\RequirePackage{fix-cm}
\documentclass[fontsize=13pt,twoside,openany,headings=normal,parskip=false,bookmarkpackage=false,DIV=calc]{scrbook}
\usepackage[no-math]{fontspec}
\setmainfont{Times New Roman}
\setsansfont{Times New Roman}
\usepackage{xeCJK}
\setCJKmainfont[AutoFakeSlant]{Microsoft YaHei}
\setCJKsansfont[AutoFakeSlant]{Microsoft YaHei}
% A semantic alternative to Chinese italics; Latin text keeps its Roman font.
\newCJKfontfamily\zhkaishu{KaiTi}
\newcommand{\textkaishu}[1]{{\zhkaishu #1}}
% The Latin font is selected outside xeCJK's character groups here, so \space
% has exactly the current Western text font's interword width and stretch.
% Additional Han-character tracking is exactly 0.07em of the active font.
\xeCJKsetup{PunctStyle=quanjiao,
  CJKglue=\hskip0.07em,CJKecglue=\space}
\usepackage{microtype}
\usepackage{amsmath,amssymb}
\allowdisplaybreaks
\raggedbottom
% Discourage isolated lines while allowing long page footnotes to split normally.
\clubpenalty=1000
\widowpenalty=1000
\displaywidowpenalty=1000
\usepackage{graphicx}
\graphicspath{{./}{../}}
% Resolve shared figures from either the main directory or a font-edition subdirectory.
% A Western word space at a Chinese-text/math boundary, in the current math style.
\newcommand{\mathwordspace}{\text{\hspace{\fontdimen2\font}}}
\renewcommand{\thechapter}{\Roman{chapter}}
\renewcommand{\thesection}{\arabic{section}}
\renewcommand*{\chapterformat}{\thechapter.\enskip}
\renewcommand*{\sectionformat}{\thesection.\enskip}
\counterwithout{figure}{chapter}
\renewcommand*{\captionformat}{}
\renewcommand{\figurename}{图}
\renewcommand{\contentsname}{目录}
\setcounter{tocdepth}{1}
\setcounter{secnumdepth}{1}
\DeclareTOCStyleEntry[dynnumwidth=true]{tocline}{chapter}
\linespread{1.35}\selectfont
\usepackage[b5paper,left=0.7cm,right=0.7cm,top=0.91cm,bottom=0.91cm,
  includehead,nofoot,headheight=\baselineskip,headsep=0.7\baselineskip]{geometry}
\setlength{\footheight}{\baselineskip}
\usepackage{scrlayer-scrpage}
\clearpairofpagestyles
\ohead{\ifnum\value{page}=1\relax\else\pagemark\fi}
\setkomafont{pageheadfoot}{\normalfont\normalsize}
\setkomafont{pagenumber}{\normalfont\normalsize}
\setkomafont{disposition}{\normalfont\bfseries}
\setkomafont{descriptionlabel}{\normalfont\bfseries}
\pagestyle{scrheadings}
\renewcommand*{\chapterpagestyle}{scrheadings}
\usepackage{indentfirst}
\setlength{\parindent}{2em}
\setlength{\emergencystretch}{2em}
\usepackage{xcolor}
\definecolor{LinkColor}{HTML}{255A83}
\usepackage[colorlinks=true,allcolors=LinkColor,linktoc=all,unicode=true,hypertexnames=false,hyperfootnotes=false]{hyperref}
\usepackage[all]{hypcap}
% Link boundaries already retain each punctuation character's elastic spacing.
% Restore only the native pair's natural-width correction, without adding a
% second source of stretch/shrink. All metrics follow the active font/style.
\ExplSyntaxOn
\NewDocumentCommand\CJKlinkpunctkern{mm}
  {\group_begin:\xeCJK_select_punct_font:
   \__xeCJK_punct_if_right:NTF #1
     {\xeCJK_get_punct_bounds:NN \c__xeCJK_right_tl #1}
     {\xeCJK_get_punct_bounds:NN \c__xeCJK_left_tl #1}
   \__xeCJK_punct_if_right:NTF #2
     {\xeCJK_get_punct_bounds:NN \c__xeCJK_right_tl #2}
     {\xeCJK_get_punct_bounds:NN \c__xeCJK_left_tl #2}
   \xeCJK_get_punct_kerning:NN #1 #2
   \skip_set:Nn \l_tmpa_skip
     {\__xeCJK_use_dim_or_skip:nNN {kern} #1 #2}
   \dim_set:Nn \l_tmpa_dim {\l_tmpa_skip}
   \nobreak\tex_kern:D \l_tmpa_dim\group_end:}
\ExplSyntaxOff
\AtBeginDocument{\pagenumbering{arabic}\thispagestyle{empty}
  \renewcommand{\contentsname}{目录}\renewcommand{\figurename}{图}}
% These references always display the page of this Chinese edition.
\newcommand{\srcpage}[1]{\hyperref[origpage-#1]{\pageref*{origpage-#1}}}
% Translator's footnotes count within chapters, with a T prefix.
\newcounter{translatornote}[chapter]
\renewcommand{\thetranslatornote}{T\arabic{chapter}.\arabic{translatornote}}
\setlength{\skip\footins}{0.8\baselineskip plus0.2\baselineskip minus0.1\baselineskip}
\renewcommand{\footnoterule}{\kern-0.4\baselineskip\hrule width0.4\columnwidth\kern0.4\baselineskip}
\hypersetup{pdftitle={Mathematische Grundlagen der Quantenmechanik / 量子力学的数学原理},pdfauthor={Johann von Neumann},pdfsubject={据1932年德文原著制作的中文译稿；所据文件为1981年重印本}}

% Font-independent ink spacing for XeLaTeX + xeCJK.
% Optional context hints: \tnote[前一字][后一字]{...},
% \origNote[前一字][后一字]{N}; the original one-argument calls remain valid.
% Hints describe neighbouring glyphs; they never print extra characters.
\makeatletter
\ExplSyntaxOn
\box_new:N \l__vn_note_mark_box
\box_new:N \l__vn_note_text_box
\box_new:N \l__vn_note_glue_box
\dim_new:N \g__vn_note_mark_left_dim
\dim_new:N \g__vn_note_mark_right_dim
\dim_new:N \g__vn_note_mark_width_dim
\dim_new:N \g__vn_note_glyph_left_dim
\dim_new:N \g__vn_note_glyph_right_dim
\dim_new:N \l__vn_note_gap_dim
\dim_new:N \l__vn_note_next_left_dim
\dim_new:N \l__vn_note_comma_right_dim
\tl_new:N \l__vn_note_string_tl
\tl_new:N \l__vn_note_first_tl
\tl_new:N \l__vn_note_last_tl
\tl_new:N \l__vn_note_footfirst_tl
\tl_set:Nn \l__vn_note_footfirst_tl { 汉 }
\cs_generate_variant:Nn \tl_item:nn { Vn }
\prg_new_conditional:Npnn \__vn_note_if_cjk:N #1 { T, F, TF }
  {
    \xeCJK_if_CJK_class:NTF #1 {\prg_return_true:}
      {
        \int_compare:nNnTF {\XeTeXcharclass `#1}
          = {\xeCJK_class_num:n {FullLeft}}
          {\prg_return_true:}
          {
            \int_compare:nNnTF {\XeTeXcharclass `#1}
              = {\xeCJK_class_num:n {FullRight}}
              {\prg_return_true:}{\prg_return_false:}
          }
      }
  }
\cs_new_protected:Npn \__vn_note_safe_footfirst:n #1
  {
    \tl_set:Nx \l__vn_note_footfirst_tl {\tl_head:n {#1}}
    \exp_after:wN \token_if_cs:NTF \l__vn_note_footfirst_tl
      {\tl_set:Nn \l__vn_note_footfirst_tl {汉}}
      {
        \exp_after:wN \token_if_math_toggle:NT \l__vn_note_footfirst_tl
          {\tl_set:Nn \l__vn_note_footfirst_tl {汉}}
      }
  }
\cs_new_protected:Npn \__vn_note_glyph:n #1
  {
    \group_begin:
      \XeTeXinterchartokenstate=0\relax
      \tl_if_blank:nTF {#1}
        { \dim_gzero:N \g__vn_note_glyph_left_dim
          \dim_gzero:N \g__vn_note_glyph_right_dim }
        {
          \__vn_note_if_cjk:NT #1 {\xeCJK_select_font:}
          \dim_gset:Nn \g__vn_note_glyph_left_dim
            { \XeTeXglyphbounds 1 \XeTeXcharglyph `#1 }
          \dim_gset:Nn \g__vn_note_glyph_right_dim
            { \XeTeXglyphbounds 3 \XeTeXcharglyph `#1 }
        }
    \group_end:
  }
\cs_new_protected:Npn \__vn_note_tight_base:nn #1#2
  {
    \tl_set:Nx \l__vn_note_string_tl {#2}
    \tl_set:Nx \l__vn_note_first_tl {\tl_head:V \l__vn_note_string_tl}
    \tl_set:Nx \l__vn_note_last_tl {\tl_item:Vn \l__vn_note_string_tl {-1}}
    \hbox_set:Nn \l__vn_note_mark_box
      {
        \bfseries
        \__vn_note_glyph:V \l__vn_note_first_tl
        \dim_gset_eq:NN \g__vn_note_mark_left_dim \g__vn_note_glyph_left_dim
        \__vn_note_glyph:V \l__vn_note_last_tl
        \dim_gset_eq:NN \g__vn_note_mark_right_dim \g__vn_note_glyph_right_dim
        #1
      }
    \kern -\g__vn_note_mark_left_dim
    \box_use:N \l__vn_note_mark_box
    \kern -\g__vn_note_mark_right_dim
  }
\cs_generate_variant:Nn \__vn_note_glyph:n { V }
\cs_new_protected:Npn \__vn_note_tight_sup:nn #1#2
  {
    \tl_set:Nx \l__vn_note_string_tl {#2}
    \tl_set:Nx \l__vn_note_first_tl { \tl_head:V \l__vn_note_string_tl }
    \tl_set:Nx \l__vn_note_last_tl { \tl_item:Vn \l__vn_note_string_tl {-1} }
    \hbox_set:Nn \l__vn_note_mark_box
      {
        \textsuperscript
          {
            \normalfont
            \hbox_set:Nn \l__vn_note_text_box {#1}
            \dim_gset:Nn \g__vn_note_mark_width_dim {\box_wd:N \l__vn_note_text_box}
            \__vn_note_glyph:V \l__vn_note_first_tl
            \dim_gset_eq:NN \g__vn_note_mark_left_dim \g__vn_note_glyph_left_dim
            \__vn_note_glyph:V \l__vn_note_last_tl
            \dim_gset_eq:NN \g__vn_note_mark_right_dim \g__vn_note_glyph_right_dim
            #1
          }
      }
    \kern -\g__vn_note_mark_left_dim
    \box_use:N \l__vn_note_mark_box
    % The math superscript also appends scriptspace; remove that dynamically.
    \kern \dim_eval:n
      { -\g__vn_note_mark_right_dim
        -\box_wd:N \l__vn_note_mark_box +\g__vn_note_mark_width_dim }
  }
\cs_new_protected:Npn \__vn_note_before:n #1
  {
    \unskip\nobreak
    \dim_set:Nn \l__vn_note_gap_dim { .14em }
    \tl_if_blank:nF {#1}
      {
        % quanjiao preserves the glyph's full advance, including its right bearing.
        \__vn_note_glyph:n {#1}
        \dim_sub:Nn \l__vn_note_gap_dim {\g__vn_note_glyph_right_dim}
      }
    \hskip \l__vn_note_gap_dim
  }
\cs_new_protected:Npn \__vn_note_after:n #1
  {
    \__vn_note_glyph:n {，}
    \dim_set_eq:NN \l__vn_note_comma_right_dim \g__vn_note_glyph_right_dim
    \__vn_note_glyph:n {汉}
    \dim_set_eq:NN \l__vn_note_next_left_dim \g__vn_note_glyph_left_dim
    \tl_if_blank:nF {#1}
      {
        \__vn_note_glyph:n {#1}
        \__vn_note_if_cjk:NT #1
          { \dim_set_eq:NN \l__vn_note_next_left_dim \g__vn_note_glyph_left_dim }
      }
    \hbox_set:Nn \l__vn_note_glue_box {\CJKglue}
    \dim_set:Nn \l__vn_note_gap_dim
      { (\l__vn_note_comma_right_dim
        +\box_wd:N \l__vn_note_glue_box +\l__vn_note_next_left_dim)/2 }
    \tl_if_blank:nF {#1}
      { \__vn_note_glyph:n {#1}
        \dim_sub:Nn \l__vn_note_gap_dim {\g__vn_note_glyph_left_dim} }
    \nobreak\hskip \l__vn_note_gap_dim
  }
\cs_generate_variant:Nn \__vn_note_after:n { V }
\cs_set_protected:Npn \@makefnmark
  {\hbox{\__vn_note_tight_sup:nn {\@thefnmark}{\@thefnmark}}}
\NewDocumentCommand \origNote { O{} O{} m }
  {
    \__vn_note_before:n {#1}
    \__vn_note_tight_sup:nn {\hyperlink{note-#3}{#3}}{#3}
    \__vn_note_after:n {#2}
  }
\NewDocumentCommand \tnote { O{} O{} m }
  {
    \stepcounter{translatornote}
    \group_begin:
      \renewcommand{\thefootnote}{\thetranslatornote}
      \__vn_note_safe_footfirst:n {#3}
      \__vn_note_before:n {#1}
      \footnote{#3}
      \__vn_note_after:n {#2}
    \group_end:
  }
\deffootnote[3em]{3em}{1em}
  {\@makefnmark\__vn_note_after:V \l__vn_note_footfirst_tl}
\NewDocumentCommand \originalnote { m +m }
  {
    \par\medskip
    \group_begin:
      \footnotesize\setlength{\parindent}{2em}\noindent
      \hypertarget{note-#1}{}
      \__vn_note_tight_base:nn {#1.}{#1.}
      \__vn_note_safe_footfirst:n {#2}
      \__vn_note_after:V \l__vn_note_footfirst_tl
      #2\par
    \group_end:
  }
\ExplSyntaxOff
\makeatother

% A nonprinting hint may be a literal glyph or semantic wrapper, e.g. \textkaishu{文}.
\ExplSyntaxOn
\bool_if_exist:NF \g__vn_note_hint_cjk_bool {\bool_new:N \g__vn_note_hint_cjk_bool}
\bool_if_exist:NF \g__vn_note_hint_fullright_bool {\bool_new:N \g__vn_note_hint_fullright_bool}
\dim_new:N \l__vn_note_removed_skip_dim
\cs_set_protected:Npn \__vn_note_glyph:n #1
  {
    \tl_if_blank:nTF {#1}
      {\dim_gzero:N \g__vn_note_glyph_left_dim
       \dim_gzero:N \g__vn_note_glyph_right_dim
       \bool_gset_false:N \g__vn_note_hint_cjk_bool
       \bool_gset_false:N \g__vn_note_hint_fullright_bool}
      {\tl_if_single_token:nTF {#1}
        {\group_begin:
           \XeTeXinterchartokenstate=0\relax
           \int_compare:nNnTF {\XeTeXcharclass `#1} = {\xeCJK_class_num:n {FullRight}}
             {\bool_gset_true:N \g__vn_note_hint_fullright_bool}
             {\bool_gset_false:N \g__vn_note_hint_fullright_bool}
           \__vn_note_if_cjk:NTF #1
             {\bool_gset_true:N \g__vn_note_hint_cjk_bool\xeCJK_select_font:}
             {\bool_gset_false:N \g__vn_note_hint_cjk_bool}
           \dim_gset:Nn \g__vn_note_glyph_left_dim
             {\XeTeXglyphbounds 1 \XeTeXcharglyph `#1}
           \dim_gset:Nn \g__vn_note_glyph_right_dim
             {\XeTeXglyphbounds 3 \XeTeXcharglyph `#1}
         \group_end:}
        {\__vn_note_glyph_styled:nn #1}}
  }
\cs_set_protected:Npn \__vn_note_glyph_styled:nn #1#2
  {#1{\__vn_note_glyph:n {#2}}}
% xeCJK FullRight emits a negative bearing rule followed by compensation glue.
% If we remove that glue, preserve its natural width in the ink-based calculation.
% Do this before any counter write nodes can block access to the trailing glue.
\cs_set_protected:Npn \__vn_note_before:n #1
  {
    \dim_set:Nn \l__vn_note_removed_skip_dim {\lastskip}
    \unskip\nobreak
    \dim_set:Nn \l__vn_note_gap_dim {.14em}
    \tl_if_blank:nF {#1}
      {\__vn_note_glyph:n {#1}
       \dim_sub:Nn \l__vn_note_gap_dim {\g__vn_note_glyph_right_dim}
       \bool_if:NT \g__vn_note_hint_fullright_bool
         {\dim_add:Nn \l__vn_note_gap_dim {\l__vn_note_removed_skip_dim}}}
    \hskip \l__vn_note_gap_dim
  }
% Restrict the reference character to a Han/CJK text character.
% Fullwidth punctuation still supplies its actual sidebearing in the final
% ink-position correction, but does not replace the Han reference glyph.
% Semantic wrappers only select the measured font; unwrap them for charclass.
\cs_new_protected:Npn \__vn_note_reference_han:n #1
  {\tl_if_single_token:nTF {#1}
     {\int_compare:nNnT {\XeTeXcharclass `#1} = {\xeCJK_class_num:n {CJK}}
        {\dim_set_eq:NN \l__vn_note_next_left_dim \g__vn_note_glyph_left_dim}}
     {\__vn_note_reference_han:nn #1}}
\cs_new_protected:Npn \__vn_note_reference_han:nn #1#2
  {\__vn_note_reference_han:n {#2}}
\cs_set_protected:Npn \__vn_note_after:n #1
  {
    \__vn_note_glyph:n {，}
    \dim_set_eq:NN \l__vn_note_comma_right_dim \g__vn_note_glyph_right_dim
    \__vn_note_glyph:n {汉}
    \dim_set_eq:NN \l__vn_note_next_left_dim \g__vn_note_glyph_left_dim
    \tl_if_blank:nF {#1}
      {\__vn_note_glyph:n {#1}
       \__vn_note_reference_han:n {#1}}
    \hbox_set:Nn \l__vn_note_glue_box {\CJKglue}
    \dim_set:Nn \l__vn_note_gap_dim
      {(\l__vn_note_comma_right_dim+\box_wd:N \l__vn_note_glue_box
        +\l__vn_note_next_left_dim)/2}
    \tl_if_blank:nF {#1}
      {\__vn_note_glyph:n {#1}
       \dim_sub:Nn \l__vn_note_gap_dim {\g__vn_note_glyph_left_dim}}
    \nobreak\hskip \l__vn_note_gap_dim
  }
\ExplSyntaxOff

% Native LaTeX page footnotes. Translator comments inside an original note
% remain visibly bracketed within that same page-footnote text.
% The preamble must set hyperfootnotes=false; links are explicit below.
\usepackage{perpage}
\MakeSorted{translatornote}
\newcounter{originalnotecounter}[chapter]
\renewcommand{\theoriginalnotecounter}{\arabic{chapter}.\arabic{originalnotecounter}}
\newcounter{translatornoteid}
\makeatletter
\ExplSyntaxOn
\tl_new:N \l__vn_note_target_tl
\tl_new:N \l__vn_note_display_tl
\bool_new:N \l__vn_inside_original_bool
\cs_set_protected:Npn \@makefnmark
  {\hbox{\tl_if_empty:NTF \l__vn_note_target_tl
    {\__vn_note_tight_sup:nn {\@thefnmark}{\@thefnmark}}
    {\hyperlink{\l__vn_note_target_tl}
      {\__vn_note_tight_sup:nn {\@thefnmark}{\@thefnmark}}}}}
\cs_new_protected:Npn \__vn_original_body:nnn #1#2#3
  {
    \tl_if_blank:nTF {#2}
      {\__vn_note_safe_footfirst:n {#3}}
      {\tl_set:Nn \l__vn_note_footfirst_tl {#2}}
    \renewcommand\thefootnote{\l__vn_note_display_tl}
    \footnote[\value{originalnotecounter}]
      {\bool_set_true:N \l__vn_inside_original_bool
       \hypertarget{note-#1}{}
       \protected@edef\@currentlabel{\l__vn_note_display_tl}
       \edef\@currentHref{note-#1}
       \label{orig-note-#1}#3}
  }
\RenewDocumentCommand \originalnote {m O{} +m}
  {\cs_gset_protected:cpn {vn_original_#1}
    {\__vn_original_body:nnn {#1}{#2}{#3}}}
\RenewDocumentCommand \origNote {O{} O{} m}
  {
    \group_begin:
      \tl_set:Nn \l__vn_note_target_tl {note-#3}
      \__vn_note_before:n {#1}
      \cs_if_exist:cTF {vn_original_seen_#3}
        {\__vn_note_tight_sup:nn
          {\hyperlink{note-#3}{\getrefnumber{orig-note-#3}}}
          {\getrefnumber{orig-note-#3}}}
        {\cs_gset:cpn {vn_original_seen_#3}{}
         \stepcounter{originalnotecounter}
         \tl_set:Nx \l__vn_note_display_tl {\theoriginalnotecounter}
         \hypertarget{note-call-#3}{}
         \use:c {vn_original_#3}}
      \__vn_note_after:n {#2}
    \group_end:
  }
\cs_new_protected:Npn \__vn_translator_anchor:
  {\hypertarget{\l__vn_note_target_tl}{}
   \protected@edef\@currentlabel{\l__vn_note_display_tl}
   \edef\@currentHref{\l__vn_note_target_tl}
   \label{translatornote-label-\arabic{translatornoteid}}}
\RenewDocumentCommand \tnote {O{} O{} O{} +m}
  {
    \group_begin:
      \bool_if:NTF \l__vn_inside_original_bool
        {\normalfont\footnotesize ［\__vn_note_before:n {［}}
        {\__vn_note_before:n {#1}}
      \stepcounter{translatornote}
      \stepcounter{translatornoteid}
      \tl_set:Nx \l__vn_note_target_tl
        {translatornote-\arabic{translatornoteid}}
      \tl_set:Nx \l__vn_note_display_tl {\thetranslatornote}
      \bool_if:NTF \l__vn_inside_original_bool
        {\hypertarget{translatornote-call-\arabic{translatornoteid}}{}
         \__vn_translator_anchor:
         \hyperlink{\l__vn_note_target_tl}
           {\__vn_note_tight_sup:nn {\l__vn_note_display_tl}{\l__vn_note_display_tl}}
         \__vn_note_after:n {译}译者注：#4］}
        {\tl_if_blank:nTF {#3}
           {\__vn_note_safe_footfirst:n {#4}}
           {\tl_set:Nn \l__vn_note_footfirst_tl {#3}}
         \hypertarget{translatornote-call-\arabic{translatornoteid}}{}
         \renewcommand\thefootnote{\l__vn_note_display_tl}
         \footnote[\value{translatornote}]
           {\__vn_translator_anchor:#4}
         \__vn_note_after:n {#2}}
    \group_end:
  }
\NewDocumentCommand \origref {m}
  {\hyperlink{note-#1}{\getrefnumber{orig-note-#1}}}
\ExplSyntaxOff
\makeatother

% Keep display spacing at its font-dependent natural size inside page footnotes.
% TeX's insertion split may accept a taller natural box by shrinking its display
% skips, but the page builder then reserves the natural box height. The mismatch
% can move the whole footnote and its call to the next page, leaving a large gap.
% Remove only this shrink component, preserving natural size and finite stretch.
% Ordinary body displays, footnote numbering, and the split penalty are unchanged.
\RequirePackage{etoolbox}
\makeatletter
\newcommand*{\vn@unshrinkdisplayskip}[1]{%
  #1=\dimexpr#1\relax plus\gluestretch#1\relax
}
\newcommand*{\vn@footnotedisplayspacing}{%
  \vn@unshrinkdisplayskip\abovedisplayskip
  \vn@unshrinkdisplayskip\belowdisplayskip
  \vn@unshrinkdisplayskip\abovedisplayshortskip
  \vn@unshrinkdisplayskip\belowdisplayshortskip
}
\AtBeginDocument{%
  \patchcmd{\@footnotetext}
    {\reset@font\footnotesize}
    {\reset@font\footnotesize\vn@footnotedisplayspacing}
    {\typeout{VN footnote display spacing: active}}
    {% KOMA-Script wraps the native insertion in a saved implementation.
     \patchcmd{\scr@saved@footnotetext}
       {\reset@font\footnotesize}
       {\reset@font\footnotesize\vn@footnotedisplayspacing}
       {\typeout{VN footnote display spacing: active (KOMA-Script)}}
       {\PackageError{vn-footnotes}{Could not set footnote display spacing}
          {The native footnote insertion changed; inspect its definition.}}}%
}
\makeatother

% Avoid phantom sorted-count entries during LaTeX's recursive child reset.
\makeatletter
\@ifundefined{vn@original@sortedcheck}{%
\let\vn@original@sortedcheck\c@schk@
\def\c@schk@#1{%
  \ifnum\value{#1}=\z@
    \expandafter\@firstoftwo
  \else
    \expandafter\@secondoftwo
  \fi
  {\pp@cl@begin\pp@cl@end}% consume reset assignment without recording a note
  {\vn@original@sortedcheck{#1}}}
}{}
\makeatother

\begin{document}

% Original notes: stored here and printed at their first call.
\originalnote{1}[德]{\hypertarget{origpage-238}{}\phantomsection\label{origpage-238}德文中已有以下综合性著作：Sommerfeld：\emph{Atombau und Spektrallinien}第4版的补卷，Braunschweig，1928。Weyl：\emph{Gruppentheorie und Quantenmechanik}，Leipzig，1928；第2版，Leipzig，1931。Fraenkel：\emph{Einführung in die Quantenmechanik}，Berlin，1929。Born与Jordan：\emph{Elementare Quantenmechanik}，Berlin，1930。Dirac：\emph{Prinzipien der Quantenmechanik}，W.~Bloch译，Leipzig，1931。}

\originalnote{2}[参]{参见\emph{Proc. Roy. Soc.}，London，第109卷（1925）及以下各卷，尤其是第113卷（1926）。P.~Jordan：\emph{Z. Physik}第40卷（1926），以及F.~London：\emph{Z. Physik}第40卷（1926），在独立于Dirac的工作中，也为这一理论提供了相近的论证。}

\originalnote{3}[参]{参见\CJKglue \hyperref[chap:IV]{第IV章}\CJKglue 及 \hyperref[sec:VI.3]{VI.3}。}

\originalnote{4}[参]{参见\CJKglue \hyperref[chap:V]{第V章}。}

\originalnote{5}[主]{主要阶段如下：Planck发现了“黑体”空腔辐射情形中的量子定律，例如可参见他的\emph{Wärmestrahlung}，Leipzig，1906中的叙述；Einstein证明了光的粒子性，即光量子理论，见\emph{Ann. Physik}〔4〕第17卷〔1905〕，由此给出了波粒二象性的第一个例子，如今我们知道，这种二象性支配着整个微观物理学；Bohr将这两组规律推广到原子模型，见\emph{Fysisk Tidskr.}第12卷（1914），以及\emph{Z. Physik}第6卷（1920）。}

\originalnote{6}[采]{采用Epstein–Sommerfeld方法，已经能够确定多重周期运动在力学定律之外还须满足的量子定律，例如参见Sommerfeld：\emph{Atombau und Spektrallinien}，Braunschweig，1924。另一方面，也已确认，自由运动的质点，或沿双曲线轨道而非椭圆轨道运行的行星，是“未量子化”的。读者若想全面了解量子理论这一发展阶段，可参阅Reiche：\emph{Die Quantentheorie, ihr Ursprung und ihre Entwicklung}，Berlin，1921，以及Landé：\emph{Fortschritte der Quantentheorie}，Dresden，1922。}

\originalnote{7}[对]{对此见Schrödinger：\emph{Ann. Physik}〔4〕第79卷（1926）。}

\originalnote{8}[\emph{Z}]{\emph{Z. Physik}第37卷（1926）。}

\originalnote{9}[参]{参见原注 \hyperlink{note-2}{\getrefnumber{orig-note-2}} 所列论文。Schrödinger的论文已汇集成书：\emph{Abhandlungen zur Wellenmechanik}，Leipzig，1928。}

\originalnote{10}[目]{目前的情况可以概括如下：理论只要涉及单个电子，或原子、分子的电子壳层，就完全成功，无论研究的是静电力的作用，还是光的产生、传播与转化中的电磁过程。相反，在处理原子核，以及尝试建立一般的、相对论性的电磁理论时，尽管在一些局部问题上已取得可观成果，似乎仍遇到了重大的困难，若没有实质性的新思想，恐怕难以克服。}

\originalnote{11}[众]{众所周知，经典力学的运动完全由哈密顿函数支配，因为它给出运动方程
\[
\dot q_l=\frac{\partial H}{\partial p_l},\qquad
\dot p_l=-\frac{\partial H}{\partial q_l}\qquad(l=1,\ldots,k).
\]\hypertarget{origpage-239}{}\phantomsection\label{origpage-239}在量子力学发现以前，人们曾保留这些运动方程，另加一些“量子条件”，以求把握量子现象，参见原注 \hyperlink{note-6}{\getrefnumber{orig-note-6}}。对于$t=0$时给定的每一组$q_1,\allowbreak \ldots,\allowbreak q_k,\allowbreak p_1,\allowbreak \ldots,\allowbreak p_k$的值，运动方程都确定其后续的时间演化，即系统在$2k$维“相空间”$q_1,\allowbreak \ldots,\allowbreak q_k,\allowbreak p_1,\allowbreak \ldots,\allowbreak p_k$中的“轨道曲线”；附加条件则从中挑选出一部分允许的初值或轨道曲线。因此，只允许少数几种轨道，相应地，能量也只能取少数几个能级。虽然量子力学彻底抛弃了这一方法，但从一开始就可以看出，哈密顿函数在其中仍须起重要作用。因为一切经验都证实了Bohr所谓对应原理：在大量子数的极限情况下，量子理论必须给出与经典力学一致的结果。}

\originalnote{12}[最]{最后这三个概念，取自前量子力学时期的量子理论，主要由N.~Bohr发展起来。后面我们将从量子力学的立场深入分析它们，参见 \hyperref[sec:III.6]{III.6} 所介绍的Dirac光的理论。若想了解它们的历史联系，可参看Bohr在1913至1916年间关于原子结构的论文，H.~Stintzing的德译本于1921年在Braunschweig出版。}

\originalnote{13}[更]{更仔细的数学分析表明，这里必然要用到无限矩阵。现在不详细考察这类矩阵的性质，后面还会深入研究。眼下只须把形式代数的演算理解为：按照熟知的矩阵加法与乘法规则进行计算。$0$及$1$尤其分别指零矩阵与单位矩阵：前者的全部元素为零；后者除对角线上为$1$外，其余元素均为零。}

\originalnote{14}[即]{即使$Q_1,\allowbreak P_1$是厄米矩阵，$Q_1P_1$及$P_1Q_1$也未必是；但$\tfrac12(Q_1P_1+P_1Q_1)$总是厄米的。以$Q_1^2P_1$为例，$\tfrac12(Q_1^2P_1+P_1Q_1^2)$与$Q_1P_1Q_1$都可以考虑；不过，当$P_1Q_1-Q_1P_1=\frac{h}{2\pi i}1$时，这两个表达式恰好相同。再以$Q_1^2P_1^2$为例，则有$\tfrac12(Q_1^2P_1^2+P_1^2Q_1^2)$、$Q_1P_1^2Q_1$、$P_1Q_1^2P_1$等等，而这些表达式即使在刚才那个特殊情形中也不再一致。这里不再详述，因为后面发展算符演算以后，这些关系会清楚得多。}

\originalnote{15}[有]{有
\[
\frac{h}{2\pi i}\frac{\partial}{\partial q_1}(q_1\psi)
=q_1\frac{h}{2\pi i}\frac{\partial\psi}{\partial q_1}
+\frac{h}{2\pi i}\psi.
\]
因此，
\[
\frac{h}{2\pi i}\frac{\partial}{\partial q_1}\cdot q_1
-q_1\cdot\frac{h}{2\pi i}\frac{\partial}{\partial q_1}
=\frac{h}{2\pi i}1,
\]
其中$1$为恒等运算，即把$\psi$映射为自身。换言之，$\frac{h}{2\pi i}\frac{\partial}{\partial q_1}$和$q_1$满足与矩阵$P_1,\allowbreak Q_1$相同的对易关系。}

\originalnote{16}[前]{前两篇论文重印于原注 \hyperlink{note-9}{\getrefnumber{orig-note-9}} 所列的书中。〔\emph{Ann. Physik}〔4〕第79卷（1926）。〕}

\originalnote{17}[参]{参见原注 \hyperlink{note-16}{\getrefnumber{orig-note-16}} 所列Schrödinger的第一篇论文。“谱”及其各个部分，将在 \hyperref[sec:II.6]{II.6} 至 \hyperref[sec:II.9]{II.9} 中准确地定义。}

\originalnote{18}[矩]{矩阵力学最初的形式，参见上面的讨论，尚未提出能够把定态作为特例涵盖在内的一般状态概念。\hypertarget{origpage-240}{}\phantomsection\label{origpage-240}理论所研究的，只是那些与能量本征值对应的定态。}

\originalnote{19}[]{$H=H(q_1,\allowbreak \ldots,\allowbreak q_k,\allowbreak p_1,\allowbreak \ldots,\allowbreak p_k)$甚至可以显含时间$t$。当然，这时一般根本不存在定态。}

\originalnote{20}[若]{若具有纯点谱，参见 \hyperref[sec:II.6]{II.6}。}

\originalnote{21}[这]{这个级数展开，以及后面的一切级数展开，都依均方收敛。\hyperref[sec:II.2]{II.2} 将详细讨论这一点。}

\originalnote{22}[定]{定态中不存在这种振动，而且只有定态中没有，这是Bohr在1913年提出的最重要的基本公设之一。经典电动力学与它相矛盾。}

\originalnote{23}[参]{参见原注 \hyperlink{note-16}{\getrefnumber{orig-note-16}} 中所列Schrödinger的第二篇论文。}

\originalnote{24}[参]{参见原注 \hyperlink{note-7}{\getrefnumber{orig-note-7}}。}

\originalnote{25}[例]{例如，参见原注 \hyperlink{note-1}{\getrefnumber{orig-note-1}} 所列Born与Jordan著作的第20、23节。}

\originalnote{26}[由]{由于
\begin{alignat*}{2}
\mathsf S^{-1}1\mathsf S&=1,&\quad  \mathsf S^{-1}(aA)\mathsf S&=a\mathsf S^{-1}A\mathsf S,\\
\mathsf S^{-1}(A+B)\mathsf S&=\mathsf S^{-1}A\mathsf S+\mathsf S^{-1}B\mathsf S,\\
\mathsf S^{-1}AB\mathsf S&=\mathsf S^{-1}A\mathsf S\cdot \mathsf S^{-1}B\mathsf S,
\end{alignat*}
对每个矩阵多项式$P(A,\allowbreak B,\allowbreak \ldots)$都有
\[
\mathsf S^{-1}P(A,B,\ldots)\mathsf S=P(\mathsf S^{-1}A\mathsf S,\mathsf S^{-1}B\mathsf S,\ldots).
\]
把对易关系的左边取作$P$，就得到这些关系的不变性；把$H$取作$P$，便得到$\mathsf S^{-1}\overline{\mathsf H}\mathsf S=\mathsf H$。}

\originalnote{27}[]{$\delta_{\mu\nu}$是熟知的Kronecker–Weierstraß记号：$\mu=\nu$时等于$1$，$\mu\ne\nu$时等于$0$。}

\originalnote{28}[对]{对所有满足$w_\varrho=\lambda$的$\varrho$，由$s_{1\varrho},\allowbreak s_{2\varrho},\allowbreak \ldots$构成的$\mathsf S$的各列形成一个完全解系；这些列还必定线性无关，因为矩阵$\mathsf S$可逆。}

\originalnote{29}[]{$\mathsf S$的列可以任意置换，只要同时对$\mathsf S^{-1}$的行作相应置换，$\mathsf H$的对角元就以相同方式置换。因此，$w_1,\allowbreak w_2,\allowbreak \ldots$的次序实际上是不确定、不可确定的。}

\originalnote{30}[积]{积分方程理论由Fredholm及Hilbert确立了最终形式。关于这一理论的详尽论述及参考文献，可参阅Courant–Hilbert：\emph{Methoden der mathematischen Physik}，Berlin，1931。}

\originalnote{31}[确]{确切地说，如果采用Lebesgue积分概念，那么，除一个测度为$0$的集合外，$q\gtrless0$时都必须有$\varphi(q)=0$；也就是说，除去这样的一个集合后，$\varphi(q)$恒等于$0$。}

\originalnote{32}[因]{因此，$\delta(q)$的图形是在$q=0$处无限细、无限高，面积为$1$的尖峰。这大致是函数$\sqrt{a/\pi}\,e^{-aq^2}$在$a\to+\infty$时的极限情形，但从严格意义上说不可能实现。}

\originalnote{33}[早]{早在量子力学产生以前，所谓“General Analysis”的创立者E.~H.~Moore就曾尝试作这种统一。对此参见Hellinger–Toeplitz在\emph{Math. Enzyklopädie}第II卷第3部第9篇中的文章，Leipzig，1927。}

\originalnote{34}[S]{Schrödinger理论中多次提到，要求波函数$\varphi$满足的条件，只是
\[
\underbrace{\int\cdots\int}_{\Omega}|\varphi(q_1,\ldots,q_k)|^2\,dq_1\cdots dq_k
\]
有限。因此，允许$\varphi$在某些地方有奇异性，例如趋于无穷，只须上述积分仍有限。一个很有启发性的例子，是Dirac的相对论性理论中的氢原子；\hypertarget{origpage-241}{}\phantomsection\label{origpage-241}参见\emph{Proc. Roy. Soc.}，London，第117卷（1928），以及W.~Gordon：\emph{Z. Physik}第48卷（1928）。}

\originalnote{35}[后]{后面对希尔伯特空间的讨论，也将给出这一定理的一个证明，参见 \hyperref[sec:II.2]{II.2}、\hyperref[sec:II.3]{II.3}，尤其是 \hyperref[thm:II.5]{II.2的定理5}。值得一提的是，这一定理中较容易证明、而且在许多问题中已经够用的一半，是$F_\Omega$与$F_Z$的某个适当部分之间的同构；这部分结果出自Hilbert，见\emph{Gött. Nachr.}，1906。Schrödinger最初的等价性证明，参见原注 \hyperlink{note-7}{\getrefnumber{orig-note-7}}，也正是建立在这一半结果之上。}

\originalnote{36}[有]{有
\[
q_mq_n\varphi(q_1,\ldots,q_k)=q_nq_m\varphi(q_1,\ldots,q_k),
\]
\[
\frac{\partial}{\partial q_m}\frac{\partial}{\partial q_n}\varphi(q_1,\ldots,q_k)
=\frac{\partial}{\partial q_n}\frac{\partial}{\partial q_m}\varphi(q_1,\ldots,q_k),
\]
\begin{align*}
&\frac{\partial}{\partial q_m}\bigl(q_n\varphi(q_1,\ldots,q_k)\bigr)
-q_n\frac{\partial}{\partial q_m}\varphi(q_1,\ldots,q_k)\\
&\qquad=\begin{cases}0,&m\ne n,\\\varphi(q_1,\ldots,q_k),&m=n,\end{cases}
\end{align*}
由此立即得到所需的算符关系。}

\originalnote{37}[用]{用条件$A.,\allowbreak B.,\allowbreak C^{(n)}$刻画$\mathfrak R_n$，出自Weyl，例如参见\emph{Raum, Zeit, Materie}，Berlin，1921。若要得到$\mathfrak R_\infty$而非$\mathfrak R_n$，只须把$C^{(n)}$换成$C^{(\infty)}$；不过，这时就必须有条件$D.,\allowbreak E.$，参见正文后面继续展开的讨论。}

\originalnote{38}[除]{除了$\mathfrak R$中的零点或零向量$0$，还有数$0$，因而同一符号用于两种东西。不过，具体语境足以区分，不会因此产生误解。}

\originalnote{39}[只]{只要求$f$属于$\mathfrak M$时$af$也属于$\mathfrak M$，以及$f,\allowbreak g$属于$\mathfrak M$时$f+g$也属于$\mathfrak M$，就已经足够。因为若$f_1,\allowbreak \ldots,\allowbreak f_k$属于$\mathfrak M$，则$a_1f_1,\allowbreak \ldots,\allowbreak a_kf_k$，继而$a_1f_1+a_2f_2$、$a_1f_1+a_2f_2+a_3f_3$、……、$a_1f_1+\cdots+a_kf_k$也都属于$\mathfrak M$。}

\originalnote{40}[]{$(f,\allowbreak f)$为实数，已经由厄米对称性得出：在该式中令$f=g$，就有$(f,\allowbreak f)=\overline{(f,\allowbreak f)}$。}

\originalnote{41}[设]{设$f$的分量为$x_1,\allowbreak \ldots,\allowbreak x_n$。按照 \hyperref[sec:II.1]{II.1} 中$\gamma)$的约定，限于有限分量的情形，有
\[
\sqrt{(f,f)}=\sqrt{\sum_{\nu=1}^{n}|x_\nu|^2},
\]
也就是通常的欧氏长度。}

\originalnote{42}[]{$(f,\allowbreak f)$为实数且非负，所以$\lVert f\rVert$为实数；这里把平方根取为非负值。对于$\lVert f-g\rVert$，结论也相同。}

\originalnote{43}[将]{将\CJKglue \hyperref[thm:II.2]{定理2} 应用于$f-g,\allowbreak g-h$，必有$f-g=0$，即$g=f$；或$g-h=0$，即$g=h$；或$g-h=c(f-g)$，其中$c$为正实数，即
\[
g=\frac{c}{c+1}f+\frac1{c+1}h.
\]
所以$g=af+(1-a)h$，其中$a=1,\allowbreak 0,\allowbreak c/(c+1)$。从几何上说，就是点$g$位于线段$f,\allowbreak h$上。}

\originalnote{44}[这]{这个聚点定义也很常见：对每个$\varepsilon>0$，都有某个属于$\mathfrak A$的$f'$，使$\lVert f-f'\rVert<\varepsilon$。两种定义等价，证明与实数分析中的完全一样。}

\originalnote{45}[为]{为简便起见，我们采用拓扑学的专门术语；例如参见Hausdorff：\emph{Mengenlehre}，Berlin，1927。正文随后会对这些术语作出解释。}

\originalnote{46}[因]{因为$\lVert f\rVert=\sqrt{(f,\allowbreak f)}=1$。}

\originalnote{47}[可]{可见，$\mathfrak R_n$中的完全正交系，对应于直角坐标系，即朝各个轴方向的单位向量。}

\originalnote{48}[作]{作为线性流形，它必须包含这样的$\{\mathfrak A\}$；又因为它是闭的，所以也包含$\{\mathfrak A\}$的聚点。\hypertarget{origpage-242}{}\phantomsection\label{origpage-242}}

\originalnote{49}[设]{设$\varphi_1,\allowbreak \varphi_2,\allowbreak \ldots$完全，则$\varphi_2,\allowbreak \varphi_3,\allowbreak \ldots$虽不再完全，却仍然是无限的！}

\originalnote{50}[注]{注意，二重序列$g_{nm}$，$n,\allowbreak m=1,\allowbreak 2,\allowbreak \ldots$，也能排成一重序列：$g_{11},\allowbreak g_{12},\allowbreak g_{21},\allowbreak g_{13},\allowbreak g_{22},\allowbreak g_{31},\allowbreak \ldots$。}

\originalnote{51}[理]{理解后面的内容，不需要本节的结果。}

\originalnote{52}[例]{例如，C.~Carathéodory：\emph{Vorlesungen über reelle Funktionen}，Leipzig，1927，尤其是第237—274页；Kamke：\emph{Das Lebesguesche Integral}，Leipzig，1925。}

\originalnote{53}[一]{一般地有
\begin{align*}
|x+y|^2&=(x+y)\overline{(x+y)}\\
&=x\bar x+y\bar y+(x\bar y+\bar xy)\\
&=|x|^2+|y|^2+2\operatorname{Re}(x\bar y).
\end{align*}}

\originalnote{54}[这]{这是Lebesgue积分理论中的惯例。}

\originalnote{55}[虽]{虽然$\alpha$连续变化，这仍然是一个和，不是积分，因为和中只会出现可数多个$\alpha$值！}

\originalnote{56}[当]{当然，把$(x(\alpha),\allowbreak y(\alpha))$定义为$\sum_\alpha x(\alpha)\overline{y(\alpha)}$。}

\originalnote{57}[这]{这是集合论中“连续统不可数”的定理。例如，参见原注 \hyperlink{note-45}{\getrefnumber{orig-note-45}} 所列Hausdorff的著作。}

\originalnote{58}[设]{设$\mathfrak R_\infty$为一个$F_\Omega$，其中$\Omega$是全体实数$x$的空间，$-\infty<x<\infty$。$d/dx$把函数映成另一个函数，也就是我们所说的算符；它只对这样的$f(x)$有定义：首先可微，其次
\[
\int_{-\infty}^{\infty}\left|\frac d{dx}f(x)\right|^2dx
\]
有限，参见 \hyperref[sec:II.8]{II.8} 对此的讨论。一般而言，$d^2f(x)/dx^2$当然未必存在；而且
\[
\int_{-\infty}^{\infty}\left|\frac{d^2}{dx^2}f(x)\right|^2dx
\]
也未必有限。例如$f(x)=|x|^{3/2}e^{-x^2}$就是如此。}

\originalnote{59}[根]{根据 \hyperref[sec:II.3]{II.3} 关于条件$E.$的讨论，只须能任意准确地逼近下述函数的所有线性组合：这类函数$f(x)$在由有限多个区间组成的集合内等于$1$，其余地方等于$0$。只要能逐一逼近其中每个这样的函数，就能做到这一点；而又只须对只有一个取值为$1$的区间的函数做到，因为其余函数是这些函数之和。设该区间为$a<x<b$。函数
\[
f(x)=\begin{cases}
0,&x<a-\varepsilon\mathwordspace{}\text{或}\mathwordspace{}x>b+\varepsilon,\\
\cos^2\!\left(\dfrac\pi2\dfrac{a-x}{\varepsilon}\right),&a-\varepsilon\le x\le a,\\
\cos^2\!\left(\dfrac\pi2\dfrac{x-b}{\varepsilon}\right),&b\le x\le b+\varepsilon,\\
1,&a<x<b
\end{cases}
\]
确实满足我们的一切正则性要求；而且只要$\varepsilon>0$足够小，就能作任意准确的逼近。}

\originalnote{60}[这]{这段考察并不严格，因为对无限和使用了线性性质。但可按以下方式补全它。设$\varphi_1,\allowbreak \varphi_2,\allowbreak \ldots$为完全规范正交系，$A,\allowbreak A^*$为彼此伴随的算符。令$f=\sum_{\nu=1}^{\infty}x_\nu\varphi_\nu$、$Af=\sum_{\nu=1}^{\infty}y_\nu\varphi_\nu$，则
\begin{align*}
y_\mu&=(Af,\varphi_\mu)=(f,A^*\varphi_\mu)\\
&=\sum_{\nu=1}^{\infty}(f,\varphi_\nu)\overline{(A^*\varphi_\mu,\varphi_\nu)}
\qquad\text{〔由\CJKglue \hyperref[thm:II.7]{定理7的$\gamma)$}〕}\\
&=\sum_{\nu=1}^{\infty}x_\nu\overline{(\varphi_\mu,A\varphi_\nu)}
=\sum_{\nu=1}^{\infty}(A\varphi_\nu,\varphi_\mu)x_\nu.
\end{align*}\hypertarget{origpage-243}{}\phantomsection\label{origpage-243}于是，若令$a_{\mu\nu}=(A\varphi_\nu,\allowbreak \varphi_\mu)$，就得到正文的公式$y_\mu=\sum_{\nu=1}^{\infty}a_{\mu\nu}x_\nu$，并保证了绝对收敛。

在序列$x_1,\allowbreak x_2,\allowbreak \ldots$的希尔伯特空间中，容易看出，序列$\varphi_1=1,\allowbreak 0,\allowbreak 0,\allowbreak \ldots$、$\varphi_2=0,\allowbreak 1,\allowbreak 0,\allowbreak \ldots$、……构成完全规范正交系。对于$f=\{x_1,\allowbreak x_2,\allowbreak \ldots\}$，有$f=\sum_{\nu=1}^{\infty}x_\nu\varphi_\nu$；对于$Af=\{y_1,\allowbreak y_2,\allowbreak \ldots\}$，有$Af=\sum_{\nu=1}^{\infty}y_\nu\varphi_\nu$。因此，与正文的联系已完全建立。

再对$A^*$构造$a^*_{\mu\nu}$，便有
\[
a^*_{\mu\nu}=(A^*\varphi_\nu,\varphi_\mu)
=(\varphi_\nu,A\varphi_\mu)
=\overline{(A\varphi_\mu,\varphi_\nu)}=\bar a_{\nu\mu}.
\]}

\originalnote{61}[]{$(Af,\allowbreak f)$必为实数，因为它等于$(A^*f,\allowbreak f)=(f,\allowbreak Af)=\overline{(Af,\allowbreak f)}$。}

\originalnote{62}[所]{所以，$U,\allowbreak U^*$必须处处有意义。另外，它们互为逆，因此每个值它们都能取到，且恰好取到一次。}

\originalnote{63}[在]{在$\int_{-\infty}^{+\infty}|\varphi(q)|^2dq$给定时，
\[
\int_{-\infty}^{+\infty}q^2|\varphi(q)|^2dq
\qquad\text{或}\qquad
\int_{-\infty}^{+\infty}\left|\frac d{dq}\varphi(q)\right|^2dq
\]
两个积分中的任意一个，都可以任意增大。例如，取$\varphi(q)=ae^{-bq^2}$，这三个积分都有限，$b>0$；它们分别正比于$a^2b^{-1/2}$、$a^2b^{-3/2}$、$a^2b^{1/2}$，所以其中任意两个都可以预先指定。}

\originalnote{64}[\emph{G}]{\emph{Gött. Nachr.}，1906。}

\originalnote{65}[在]{在变形的第三步，利用了$R$的厄米性：
\begin{align*}
&\left(R\frac{f+g}{2},\frac{f+g}{2}\right)
-\left(R\frac{f-g}{2},\frac{f-g}{2}\right)\\
&\quad=\frac{(Rf,g)+(Rg,f)}2
=\frac{(Rf,g)+\overline{(f,Rg)}}2\\
&\quad=\frac{(Rf,g)+\overline{(Rf,g)}}2
=\operatorname{Re}(Rf,g).
\end{align*}}

\originalnote{66}[参]{参见原注 \hyperlink{note-9}{\getrefnumber{orig-note-9}} 所列书中的第五篇论文，\emph{Ann. Physik}〔4〕第80卷（1926）。}

\originalnote{67}[这]{这里不考察更细的收敛问题。最初的矩阵理论与波动理论也没有准确处理这些问题；而且我们后面还会解决这些问题，例如参见 \hyperref[sec:II.9]{II.9}。}

\originalnote{68}[后]{后一个条件，只有在$\mathsf H$处处有意义且连续时才显然成立，即从$f_n\to f$还能推出$\mathsf Hf_n\to\mathsf Hf$时如此。不过，容易看出，下面这个较弱的性质就已经足够：若$f_n\to f$、$\mathsf Hf_n\to f^*$，则$\mathsf Hf=f^*$。这就是所谓$\mathsf H$的“闭性”，参见作者在\emph{Math. Ann.}第102卷〔1929〕的论文。量子力学的算符，即使不连续，也总满足这个性质。确切地说，一个不闭的厄米算符，可以唯一地扩充其定义域，使它成为仍然厄米的闭算符；例如，连续性就不具有这样的性质。参见 \hyperref[sec:II.9]{II.9}，第 \srcpage{76} 页。}

\originalnote{69}[例]{例如，参见Schrödinger对氢原子的处理，原注 \hyperlink{note-16}{\getrefnumber{orig-note-16}} 所引之处。}

\originalnote{70}[参]{参见原注 \hyperlink{note-64}{\getrefnumber{orig-note-64}} 所引之处。}

\originalnote{71}[例]{例如，参见Courant–Hilbert，原注 \hyperlink{note-30}{\getrefnumber{orig-note-30}} 所引之处。\hypertarget{origpage-244}{}\phantomsection\label{origpage-244}}

\originalnote{72}[即]{即$E(l)=(e_{\mu\nu}(l))$，
\[
E(l;\xi,\eta)=\sum_{\mu,\nu=1}^{n}e_{\mu\nu}(l)\xi_\mu\bar\eta_\nu.
\]
因此，$e_{\mu\nu}(l)=\sum_{\lambda_\varrho\le l}x_{\varrho\mu}\bar x_{\varrho\nu}$。}

\originalnote{73}[关]{关于Stieltjes积分的概念，参见Perron：\emph{Die Lehre von den Kettenbrüchen}，Leipzig，1913；另外，特别考虑算符理论需要的论述，可参见Carleman：\emph{Équations intégrales singulières}，Upsala，1923。对这些问题兴趣较少的读者，只须知道如下定义。对区间$a,\allowbreak b$的每一组分点$\Lambda_0,\allowbreak \Lambda_1,\allowbreak \ldots,\allowbreak \Lambda_k$，
\[
a\le\Lambda_0<\Lambda_1<\cdots<\Lambda_k\le b,
\]
作和$\sum_{\tau=1}^{k}f(\Lambda_\tau)\{g(\Lambda_\tau)-g(\Lambda_{\tau-1})\}$。如果随着分点$\Lambda_0,\allowbreak \ldots,\allowbreak \Lambda_k$越来越密，这个和收敛，就定义积分$\int_a^bf(x)\,dg(x)$有意义，并等于这个极限。取$g(x)=x$时，就回到熟知的Riemann积分。

因此，在这里，所推得的等式表明，$\int_{-\infty}^{+\infty}x\,dE(x;\xi,\allowbreak \eta)$存在；我们把变量$x$改记为$\lambda$，而这个积分就等于$\sum_{\mu,\allowbreak \nu=1}^{n}h_{\mu\nu}\xi_\mu\bar\eta_\nu$。}

\originalnote{74}[]{$\min(a,\allowbreak b,\allowbreak \ldots,\allowbreak e)$及$\max(a,\allowbreak b,\allowbreak \ldots,\allowbreak e)$，分别是有限个实数$a,\allowbreak b,\allowbreak \ldots,\allowbreak e$中的最小值与最大值。}

\originalnote{75}[参]{参见原注 \hyperlink{note-64}{\getrefnumber{orig-note-64}} 所引之处，以及原注 \hyperlink{note-73}{\getrefnumber{orig-note-73}} 所列Carleman的著作。后面还会多次讨论这些“连续谱”，参见 \hyperref[sec:II.8]{II.8}。}

\originalnote{76}[]{$A(\lambda)\to B$，其中$A(\lambda),\allowbreak B$为$\mathfrak R_\infty$中的算符，$\lambda$为参数，表示：对于$\mathfrak R_\infty$中的一切$f$，都有$A(\lambda)f\to Bf$。因此，它是希尔伯特空间中收敛关系的一种缩写。}

\originalnote{77}[这]{这是由原注 \hyperlink{note-73}{\getrefnumber{orig-note-73}} 中Stieltjes积分的定义得出的。证明见那里所列的著作。}

\originalnote{78}[参]{参见\emph{Math. Ann.}第102卷（1929）。}

\originalnote{79}[我]{我们是对$\lambda$序列证明这一点的。但无论考察从左侧$\lambda<\lambda_0$还是从右侧$\lambda>\lambda_0$趋于$\lambda_0$的情形，同一侧的所有$\lambda$序列都必须有相同的极限。因为同一侧的任意两个这样的序列，都可以合并成一个序列；而既然这个序列也有极限，它的两个组成部分就必须有同一极限。因此，这种收敛是朝向共同极限的收敛，对一切序列都如此；在$\lambda$连续变化时同样成立。}

\originalnote{80}[这]{这正是 \hyperref[sec:II.7]{II.7} 中$E(\lambda;\xi,\allowbreak \eta)$定义的准确转述。}

\originalnote{81}[定]{\hyperref[thm:II.10]{定理10} 证明中给出的构造表明，当$\lVert\varphi\rVert=1$时，$P_{[\varphi]}f=(f,\allowbreak \varphi)\varphi$；因此$\lVert P_{[\varphi]}f\rVert=|(f,\allowbreak \varphi)|=|(\varphi,\allowbreak f)|$。}

\originalnote{82}[由]{由\CJKglue \hyperref[thm:II.7]{定理7}，
\[
f=\sum_\varrho(f,\varphi_\varrho)\varphi_\varrho
=\sum_\varrho P_{[\varphi_\varrho]}f.
\]
这也可从 \hyperref[sec:II.4]{II.4} 末尾的讨论得出。}

\originalnote{83}[根]{根据原注 \hyperlink{note-73}{\getrefnumber{orig-note-73}}，
\[
\int_{-\infty}^{+\infty}\lambda^2\,d\left(\sum_{\lambda_\varrho\le\lambda}|x_\varrho|^2\right)
=\lim\sum_{\tau=1}^{k}\Lambda_\tau^2
\sum_{\Lambda_{\tau-1}<\lambda_\varrho\le\Lambda_\tau}|x_\varrho|^2.
\]\hypertarget{origpage-245}{}\phantomsection\label{origpage-245}如果处处有$\Lambda_\tau^2-\Lambda_{\tau-1}^2<\varepsilon$，也就是分点网格$\Lambda_0,\allowbreak \ldots,\allowbreak \Lambda_k$足够细，那么把这个和换成
\[
\sum_{\tau=1}^{k}\sum_{\Lambda_{\tau-1}<\lambda_\varrho\le\Lambda_\tau}
\lambda_\varrho^2|x_\varrho|^2
=\sum_{\Lambda_0<\lambda_\varrho\le\Lambda_k}\lambda_\varrho^2|x_\varrho|^2
\]
时，改变量小于$\varepsilon\sum_{\varrho=1}^{\infty}|x_\varrho|^2=\varepsilon\lVert f\rVert^2$；而只要$\Lambda_0$足够小、$\Lambda_k$足够大，后一式又任意接近$\sum_{\tau=1}^{\infty}\lambda_\varrho^2|x_\varrho|^2$。\tnote[。][所][原]{原书这个和使用求和指标$\tau$，却在被加项中保留$\varrho$；这里照录这一不一致。}所以，后一个和就是所求极限，即积分值。下一个积分公式也完全同样证明。}

\originalnote{84}[在]{在这里，我们采用的数学上正确的方法，与Dirac的符号方法分道扬镳，例如参见原注 \hyperlink{note-1}{\getrefnumber{orig-note-1}} 所列他的著作。后一种方法会进一步把满足$(q-\lambda)f(q)=0$的$f$看成解；这里为简便起见，取$l=j=1$、$q_j=q$。可是，对这样的普通函数，每个内积$(f,\allowbreak g)=\int f(q)\overline{g(q)}\,dq$都等于$0$；而这里又要求$f\ne0$。因此，便把$f(q)$在$q=\lambda$这一点——它唯一不为$0$的地方——设想为无限大，而且大到足以使$(f,\allowbreak g)\ne0$。由于$q\ne\lambda$时$f(q)=0$，$\int f(q)\overline{g(q)}\,dq$就只能依赖$g(\lambda)$；又因积分具有可加性，这个值显然必须正比于$\overline{g(\lambda)}$。于是它等于$c\overline{g(\lambda)}$，且应有$c\ne0$。用$f(q)/c$代替$f(q)$，便可令$c=1$。这样，我们就有了一个虚构的函数$f(q)$，满足$\int f(q)\overline{g(q)}\,dq=\overline{g(\lambda)}$。

当然，只须考察$\lambda=0$的情形，把$f(q)$记作$\delta(q)$。\tnote[。][它][这]{这里以及原注 \hyperlink{note-97}{\getrefnumber{orig-note-97}}、\hyperlink{note-98}{\getrefnumber{orig-note-98}} 中的近似函数，原书交替使用$q$与$x$；此处保留原记号。}它由下列性质定义：
\[
\text{\hypertarget{note84-Delta}{\(\Delta\).}}\qquad q\delta(q)\equiv0,\qquad
\int\delta(q)f(q)\,dq=f(0).
\]
对任意$\lambda$，$\delta(q-\lambda)$就是解。虽然具有性质 \hyperlink{note84-Delta}{$\Delta$} 的函数$\delta$并不存在，但有些函数序列会趋向这种行为，当然，极限函数并不存在。例如
\[
f_\varepsilon(q)=\begin{cases}\dfrac1{2\varepsilon},&|x|<\varepsilon,\\0,&|x|\ge\varepsilon,\end{cases}
\qquad\varepsilon\to+0,
\]
或者$f_\varepsilon(q)=\sqrt{a/\pi}\,e^{-ax^2}$，$a\to+\infty$。另见 \hyperref[sec:I.3]{I.3}，尤其是原注 \hyperlink{note-32}{\getrefnumber{orig-note-32}}。}

\originalnote{85}[这]{这个思想在这里只能看作启发式的出发点；关于它的严格展开，可见Hellinger〔\emph{J. f. Math.}第136卷〔1909〕〕及Weyl〔\emph{Math. Ann.}第68卷〔1910〕〕的论文。}

\originalnote{86}[当]{当然，只有物理应用上的成功，才能为这一立场——或者说，为在量子力学中采用这一立场——提供根据。}

\originalnote{87}[P]{Plancherel：\emph{Circ. Math. di Pal.}第30卷（1910）；Titchmarsh：\emph{London Math. Soc. Proc.}第22卷（1924）。}

\originalnote{88}[即]{即$E'(\lambda)$所对应的，并不是$A'=\frac{h}{2\pi i}\frac{\partial}{\partial q}$本身，而是某个算符，其定义域包含$A'$的定义域，并在后者上与$A'$一致。参见 \hyperref[sec:II.9]{II.9} 的讨论。}

\originalnote{89}[它]{它的全部Fourier系数都为零，所以函数本身也为零。例如，参见Courant–Hilbert，原注 \hyperlink{note-30}{\getrefnumber{orig-note-30}} 所引之处。\hypertarget{origpage-246}{}\phantomsection\label{origpage-246}}

\originalnote{90}[这]{这里虽然恒有$M^{-1}Mf(q)=f(q)$——只须对$q<0$定义$f(q)=0$，再应用前面的定理即可——却并不总有$MM^{-1}F(p)=F(p)$。因为一般有$\lVert M^{-1}F\rVert<\lVert F\rVert$，从而$\lVert MM^{-1}F\rVert<\lVert F\rVert$。因此，$M^{-1}M=1$，但$MM^{-1}\ne1$；也就是说，$M^{-1}$并不是$M$的真正逆算符。也不存在另一个逆算符，因为若有，由$M^{-1}M=1$可知，它必须等于我们这个$M^{-1}$。所以，当然可以构造$E'(\lambda)=ME(\lambda)M^{-1}$；甚至还能推出$E'(\lambda)^2=E'(\lambda)$，因为这一步只涉及$M^{-1}M$。但是，当$\lambda\to+\infty$时，$E'(\lambda)\to MM^{-1}\ne1$。}

\originalnote{91}[严]{严格说，这里不能只靠形式符号演算，而须借助严格的等式$(Af,\allowbreak g)=\int\lambda\,d(E(\lambda)f,\allowbreak g)$。这样计算如下：
\begin{align*}
(AFf,g)&=\int\lambda\,d(E(\lambda)Ff,g)
=\int\lambda\,d(FE(\lambda)f,g)\\
&=\int\lambda\,d(E(\lambda)f,Fg)
=(Af,Fg)=(FAf,g).
\end{align*}
由此得到$AF=FA$。}

\originalnote{92}[这]{这是由对Stieltjes积分普遍有效的等式得出的：
\[
\int f(\lambda)\,d\left(\int^\lambda g(\lambda')\,dh(\lambda')\right)
=\int f(\lambda)g(\lambda)\,dh(\lambda).
\]
根据$d$与$\int^\lambda$之间互为逆运算的关系，这个等式立即可见。作者曾给出一个严格证明，见\emph{Annals of Mathematics}第32卷（1931）。}

\originalnote{93}[即]{即
\[
(Bf,g)=\int_{-\infty}^{\infty}r(\lambda)\,d(E(\lambda)f,g),\qquad
(Cf,g)=\int_{-\infty}^{\infty}s(\lambda)\,d(E(\lambda)f,g).
\]}

\originalnote{94}[作]{作者在\emph{Annals of Math.}第32卷（1931）中严格建立了这个算符函数概念。更早以前，F.~Riesz曾通过多项式的极限过程，定义一般的算符函数。}

\originalnote{95}[关]{关于无界厄米算符的理论，下面除Hilbert的有界算符理论外，主要依据作者在原注 \hyperlink{note-78}{\getrefnumber{orig-note-78}} 所引之处的论述。M.~Stone独立地得到了类似结果，见\emph{Proc. Nat. Ac.}，1929及1930。}

\originalnote{96}[\emph{M}]{\emph{Math. Ann.}第69卷（1911）。}

\originalnote{97}[在]{在$-\infty<q<+\infty$上解析、并且$\int_{-\infty}^{\infty}|f(q)|^2dq$、$\int_{-\infty}^{\infty}|f'(q)|^2dq$等等有限的函数$f(q)$，也在$\mathfrak R_\infty$中处处稠密。根据 \hyperref[sec:II.3]{II.3的$D.$}，函数
\[
f_{a,b}(q)=\begin{cases}1,&a<q<b,\\0,&\text{其余情形}\end{cases}
\]
的线性组合处处稠密；因此，只须用上述函数任意准确地逼近这些$f_{a,\allowbreak b}(q)$。例如，函数
\[
f^{(\varepsilon)}_{a,b}(q)
=\frac12-\frac12\operatorname{Tg\,hyp}\frac{(x-a)(x-b)}{\varepsilon}
=\frac1{e^{2(x-a)(x-b)/\varepsilon}+1}
\]
就具有所需的性质，而且当$\varepsilon\to+0$时收敛到$f_{a,\allowbreak b}(q)$。\hypertarget{origpage-247}{}\phantomsection\label{origpage-247}}

\originalnote{98}[仍]{仍然只须用$\mathfrak A^0$中的函数逼近$f_{a,\allowbreak b}(q)$，其中$0\le a<b\le1$。例如，可以采用
\[
f^{(\varepsilon)}_{a,b}(q)=\frac12-\frac12\operatorname{Tg\,hyp}\left(
\frac1\varepsilon\frac{(x-a-\varepsilon)(x-b+\varepsilon)}{x(1-x)}\right),
\]
并令$\varepsilon\to+0$。}

\originalnote{99}[可]{可以设想用多项式逼近函数$f(\lambda)$，所以只须考察多项式，也就是组成多项式的各个幂：$f(\lambda)=\lambda^s$，$s=0,\allowbreak 1,\allowbreak 2,\allowbreak \ldots$。这里作酉变换没有影响，因此可以假定$A$为对角矩阵；对角元就是本征值，记作$\lambda_1,\allowbreak \ldots,\allowbreak \lambda_n$。于是，只须证明$A^s$也为对角矩阵，对角元为$\lambda_1^s,\allowbreak \ldots,\allowbreak \lambda_n^s$，而这显然成立。}

\originalnote{100}[这]{这些性质是厄米性或酉性的特征，只须再对对角矩阵验证即可。对于对角元为$\lambda_1,\allowbreak \ldots,\allowbreak \lambda_n$的对角矩阵$A$，其共轭转置是对角元为$\bar\lambda_1,\allowbreak \ldots,\allowbreak \bar\lambda_n$的对角矩阵$A^*$。因此，$A=A^*$意味着$\lambda_1=\bar\lambda_1,\allowbreak \ldots,\allowbreak \lambda_n=\bar\lambda_n$，即$\lambda_1,\allowbreak \ldots,\allowbreak \lambda_n$为实数；而$AA^*=A^*A=1$意味着$\lambda_1\bar\lambda_1=1,\allowbreak \ldots,\allowbreak \lambda_n\bar\lambda_n=1$，即$|\lambda_1|=\cdots=|\lambda_n|=1$。}

\originalnote{101}[这]{这些事实的证明，见作者在原注 \hyperlink{note-78}{\getrefnumber{orig-note-78}} 所引之处的论文，以及A.~Wintner：\emph{Math. Z.}第30卷（1929）。有界$f(\sigma)$的所有积分$\int_0^1 f(\sigma)\,d(E(\sigma)f,\allowbreak g)$绝对收敛，可以这样看出。只须考虑$\operatorname{Re}(E(\sigma)f,\allowbreak g)$，因为用$if,\allowbreak g$代替$f,\allowbreak g$，它就变成$\operatorname{Im}(E(\sigma)f,\allowbreak g)$。由于
\[
\operatorname{Re}(E(\sigma)f,g)
=\left(E(\sigma)\frac{f+g}{2},\frac{f+g}{2}\right)
-\left(E(\sigma)\frac{f-g}{2},\frac{f-g}{2}\right),
\]
只须研究$(E(\sigma)f,\allowbreak f)$。而$\int_0^1 f(\sigma)\,d(E(\sigma)f,\allowbreak f)$的被积函数有界，$d$后面的$\sigma$函数又单调，因此结论显然成立。}

\originalnote{102}[也]{也就是说，把Stieltjes积分应用于$\mathfrak R_\infty$中的元素，而非数。所有这些关系都应作如下理解：在$\mathfrak R_\infty$中固定一个$g$，把其中出现的每个$\mathfrak R_\infty$元素都与$g$取内积；这样得到的等式，对一切$g$都成立。与 \hyperref[sec:II.7]{II.7} 中的算符Stieltjes积分相比，这是一种半符号化的做法：当时不是选$\mathfrak R_\infty$中的一个$g$，而是选两个$f,\allowbreak g$，并且不是形成$(\ldots,\allowbreak g)$，而是形成$(\ldots f,\allowbreak g)$；后一个式子中的省略号代表算符。}

\originalnote{103}[这]{这里隐含着一个前提：对于每个给定的单位分解$F(\lambda)$，确实存在这样的算符。也就是说，当$\int_{-\infty}^{\infty}\lambda^2\,d\|F(\lambda)f\|^2$有限时，可以找到$f^*$，使对一切$g$都有
\[
(f^*,g)=\int_{-\infty}^{\infty}\lambda\,d(F(\lambda)f,g),
\]
并且这些$f$处处稠密。如此定义的算符为厄米算符，这一点由$\hyperref[cond:II.Sbar3]{\overline{\mathrm S}_3}$得出：在作为依据的等式中交换$f,\allowbreak g$，再取复共轭即可。上述两项预设，作者均在原注 \hyperlink{note-78}{\getrefnumber{orig-note-78}} 所引之处证明了。}

\originalnote{104}[要]{要使$A$的本征值问题总能求解，就必须由此得到$U$的酉性，也就是$\mathfrak E=\mathfrak F=\mathfrak R_n$或$\mathfrak R_\infty$。对于$\mathfrak R_\infty$，正如我们从非极大的$A$的存在所推得的，这并不成立。但在$\mathfrak R_n$中，它必须成立，而且可以直接看出。因为$\mathfrak R_n$的每个线性流形都是闭的，所以由全部$\varphi-U\varphi$组成的流形也闭；而它处处稠密，故等于$\mathfrak R_n$。$\mathfrak E$是所有$\varphi$组成的集合，其维数不会小于它的线性像，即所有$\varphi-U\varphi$组成的集合的维数；因此，$\mathfrak E$具有最大维数$n$。$\mathfrak F$是$\mathfrak E$的线性一一对应像，\hypertarget{origpage-248}{}\phantomsection\label{origpage-248}也必须具有同一维数。$n$有限，因而$\mathfrak E=\mathfrak F=\mathfrak R_n$。}

\originalnote{105}[不]{不过，如作者在原注 \hyperlink{note-78}{\getrefnumber{orig-note-78}} 所引之处证明的，下面这个算符虽然极大，却并非超极大。设$\mathfrak R_\infty$为所有在$0\le q<\infty$上定义、满足$\int_0^\infty|f(q)|^2\,dq$有限的函数所成的空间。令$R=i\frac{d}{dq}$，例如先对所有连续可微、满足$\int_{-\infty}^{\infty}|f'(q)|^2\,dq$有限且$f(0)=0$的$f$定义它，再将它闭化。\tnote[。][它][原]{原书在$f'$的积分中印出的下限为$-\infty$，但前面把函数的定义域规定为$[0,\allowbreak \infty)$。这里保留原印下限。}它等于$-\frac{2\pi}{h}A'$，这里$A'$是 \hyperref[sec:II.8]{II.8} 中取区间$0,\allowbreak \infty$时的算符，因此它是厄米的。这个$R$只是极大而非超极大，实际算出$\mathfrak E,\allowbreak \mathfrak F$便可验证。

这件事值得注意，因为$A'=-\frac{h}{2\pi}R$通常在物理上被解释为动量算符；这里对应的是由墙$q=0$界定的半空间。}

\originalnote{106}[这]{这个概念出自Erhard Schmidt，参见原注 \hyperlink{note-78}{\getrefnumber{orig-note-78}} 所引之处。}

\originalnote{107}[因]{因为根据 \hyperref[sec:II.5]{II.5}，$R\cdot1$、$1\cdot R$有意义，当且仅当$R$有意义；当$a\ne0$时，$R\cdot a1$、$a1\cdot R$也一样。所以，这两个乘积相等，即$R$与$a1$对易。因此，$R$与$a1$的对易性只有一种例外：$a=0$而$R$不处处有意义。这个例外颇不理想，也提示我们应当修改对易性的定义。}

\originalnote{108}[]{$a1$对应单位分解
\[
F(\mu)=\begin{cases}1,&\mu\ge a,\\0,&\mu<a,\end{cases}
\]
很容易验证。}

\originalnote{109}[对]{对于属于一个特殊类别——所谓“全连续”类，参见原注 \hyperlink{note-70}{\getrefnumber{orig-note-70}} 所引之处——的两个厄米算符$A,\allowbreak B$，Toeplitz证明了一个定理，例如参见原注 \hyperlink{note-33}{\getrefnumber{orig-note-33}} 所引之处。由这一定理可以推出正文中的上述定理。（Toeplitz的定理说，存在一个由$A,\allowbreak B$的共同本征函数组成的完全规范正交系。）关于任意$A,\allowbreak B$及$A,\allowbreak B,\allowbreak C,\allowbreak \ldots$的一般定理，作者已在原注 \hyperlink{note-94}{\getrefnumber{orig-note-94}} 所引之处证明。}

\originalnote{110}[这]{这些$\varkappa_1,\allowbreak \varkappa_2,\allowbreak \ldots$必须选得有界，例如$\varkappa_m=1/m$，使$R$连续。事实上，由$R$连续，即$\|Rf\|\le C\|f\|$，立即得到$\|R\psi_m\|=\|\varkappa_m\psi_m\|=|\varkappa_m|\le C\|\psi_m\|=C$，从而$|\varkappa_m|\le C$。反之，若$|\varkappa_m|\le C$，$m=1,\allowbreak 2,\allowbreak \ldots$，则
\begin{align*}
\|Rf\|^2&=\left\|R\left(\sum_{m=1}^{\infty}x_m\psi_m\right)\right\|^2
=\left\|\sum_{m=1}^{\infty}x_m\varkappa_m\psi_m\right\|^2
=\sum_{m=1}^{\infty}|x_m|^2|\varkappa_m|^2,\\
\|f\|^2&=\left\|\sum_{m=1}^{\infty}x_m\psi_m\right\|^2
=\sum_{m=1}^{\infty}|x_m|^2.
\end{align*}
所以$\|Rf\|^2\le C^2\|f\|^2$，$\|Rf\|\le C\|f\|$，即$R$连续。}

\originalnote{111}[允]{允许在$(q_1,\allowbreak q_2)$空间中某个Lebesgue测度为$0$的集合上有例外！}

\originalnote{112}[例]{例如，参见原注 \hyperlink{note-45}{\getrefnumber{orig-note-45}} 所引之处。}

\originalnote{113}[]{$\{a_{\mu\nu}\}$代表变换，即算符$\eta_\mu=\sum_{\nu=1}^n a_{\mu\nu}\xi_\nu$，$\mu=1,\allowbreak \ldots,\allowbreak n$。如果像 \hyperref[sec:II.7]{II.7} 的讨论那样，作变换
\[
\xi_\mu=\sum_{\nu=1}^n\varkappa_{\nu\mu}\mathfrak x_\nu,
\qquad
\eta_\mu=\sum_{\nu=1}^n\varkappa_{\nu\mu}\mathfrak y_\nu,
\qquad\mu=1,\ldots,n,
\]\hypertarget{origpage-249}{}\phantomsection\label{origpage-249}便得到$\mathfrak y_\mu=\sum_{\nu=1}^n\mathfrak a_{\mu\nu}\mathfrak x_\nu$，$\mu=1,\allowbreak \ldots,\allowbreak n$，其中
\[
\mathfrak a_{\mu\nu}=\sum_{\varrho,\sigma=1}^n
a_{\varrho\sigma}\overline{\varkappa_{\mu\varrho}}\varkappa_{\nu\sigma},
\qquad\mu,\nu=1,\ldots,n,
\]
$\{\mathfrak a_{\mu\nu}\}$就是变换后的矩阵。显然，
\begin{align*}
\sum_{\mu=1}^n\mathfrak a_{\mu\mu}
&=\sum_{\mu,\varrho,\sigma=1}^n a_{\varrho\sigma}
\overline{\varkappa_{\mu\varrho}}\varkappa_{\mu\sigma}\\
&=\sum_{\varrho,\sigma=1}^n a_{\varrho\sigma}
\left(\sum_{\mu=1}^n\overline{\varkappa_{\mu\varrho}}\varkappa_{\mu\sigma}\right)
=\sum_{\varrho=1}^n a_{\varrho\varrho},
\end{align*}
即迹不变。}

\originalnote{114}[准]{准确的表述如下：若$A$为超极大且定号的算符，则存在且只存在一个同样是超极大、定号的算符$A'$，满足$A'^2=A$。下面证明它的存在性。

设$A=\int_{-\infty}^{\infty}\lambda\,dE(\lambda)$为$A$的本征值表示。由于$A$定号，$E(\lambda)$在$\lambda<0$时为常数，从而根据$\hyperref[cond:II.Sbar1]{\overline{\mathrm S}_1}$，等于$0$。否则，可以选取适当的$\lambda_1<\lambda_2<0$，使$E(\lambda_2)-E(\lambda_1)\ne0$，于是可取$f\ne0$满足$(E(\lambda_2)-E(\lambda_1))f=f$。那么，按我们前面的推论，$\lambda\ge\lambda_2$时$E(\lambda)f=f$，$\lambda\le\lambda_1$时$E(\lambda)f=0$。所以
\begin{align*}
(Af,f)&=\int_{-\infty}^{\infty}\lambda\,d(E(\lambda)f,f)
=\int_{\lambda_1}^{\lambda_2}\lambda\,d(E(\lambda)f,f)\\
&\le\int_{\lambda_1}^{\lambda_2}\lambda_2\,d(E(\lambda)f,f)
=\lambda_2((E(\lambda_2)-E(\lambda_1))f,f)
=\lambda_2(f,f)<0.
\end{align*}
因此，
\[
A=\int_{-\infty}^{\infty}\lambda\,dE(\lambda)
=\int_0^{\infty}\lambda\,dE(\lambda)
=\int_0^{\infty}\mu^2\,dE(\mu^2),
\]
而$A'=\int_1^{\infty}\mu\,dE(\mu^2)$就满足要求。\tnote[。][][原]{原书这个$A'$公式的积分下限印作$1$；前面的计算要求下限为$0$。这里保留原印下限。}

注意，我们由定号性推出了$\lambda<0$时$E(\lambda)=0$；而由后者反过来推出定号性。因此，这一性质，也就是整个谱均非负，正是定号性的特征。}

\originalnote{115}[参]{参见原注 \hyperlink{note-64}{\getrefnumber{orig-note-64}} 所引之处。一个直接证明如下。取$\lambda_0<\lambda_1<\cdots<\lambda_n$，全部不小于$\varepsilon$，或全部不大于$-\varepsilon$，并使$E(\lambda_0)\ne E(\lambda_1)\ne E(\lambda_2)\ne\cdots\ne E(\lambda_n)$。那么$E(\lambda_\nu)-E(\lambda_{\nu-1})\ne0$，所以可选$\varphi_\nu\ne0$，使$(E(\lambda_\nu)-E(\lambda_{\nu-1}))\varphi_\nu=\varphi_\nu$。由此，$\lambda\ge\lambda_\nu$时$E(\lambda)\varphi_\nu=\varphi_\nu$，$\lambda\le\lambda_{\nu-1}$时$E(\lambda)\varphi_\nu=0$，并且还可令$\|\varphi_\nu\|=1$。由上可得$\mu\ne\nu$时$(\varphi_\mu,\allowbreak \varphi_\nu)=0$。因此，$\varphi_1,\allowbreak \ldots,\allowbreak \varphi_n$构成规范正交系，可以把它扩充为完全系$\varphi_1,\allowbreak \ldots,\allowbreak \varphi_n,\allowbreak \varphi_{n+1},\allowbreak \ldots$。对于$\nu=1,\allowbreak \ldots,\allowbreak n$，有
\begin{align*}
\|A\varphi_\nu\|^2
&=\int_{-\infty}^{\infty}\lambda^2\,d\|E(\lambda)\varphi_\nu\|^2
=\int_{\lambda_{\nu-1}}^{\lambda_\nu}\lambda^2\,d\|E(\lambda)\varphi_\nu\|^2\\
&\ge\int_{\lambda_{\nu-1}}^{\lambda_\nu}\varepsilon^2\,d\|E(\lambda)\varphi_\nu\|^2
=\varepsilon^2\bigl(\|E(\lambda_\nu)\varphi_\nu\|^2-\|E(\lambda_{\nu-1})\varphi_\nu\|^2\bigr)
=\varepsilon^2\|\varphi_\nu\|^2=\varepsilon^2.
\end{align*}
所以$\sum_{\mu=1}^{\infty}\|A\varphi_\mu\|^2$一方面不小于$\sum_{\mu=1}^n\|A\varphi_\mu\|^2\ge n\varepsilon^2$，另一方面等于$\Sigma(A)=C^2$，即$n\le C^2/\varepsilon^2$。也就是说，在$|\lambda|\ge\varepsilon$时，$E(\lambda)$至多取$2C^2/\varepsilon^2$个不同的值；因此，它只有有限次跳跃，各次跳跃之间为常值区间。\hypertarget{origpage-250}{}\phantomsection\label{origpage-250}也就是说，$|\lambda|\ge\varepsilon$的部分只有点谱。对所有$\varepsilon>0$都如此，所以整个谱是纯点谱。}

\originalnote{116}[除]{除论述过程中所引的原始论文外，这里首先应参考 Hellinger 和 Toeplitz 的百科全书文章，见\CJKglue \hyperlink{note-33}{原注 \getrefnumber{orig-note-33}}。}

\originalnote{117}[类]{类比几何学，在$\mathfrak R_\infty$中，以$\varphi_0$为中心、$r$为半径的球，定义为满足$\|f-\varphi_0\|\le r$的所有点$f$的集合；其内部为$\|f-\varphi_0\|<r$的点集，球面为$\|f-\varphi_0\|=r$的点集。对单位球，有$\varphi_0=0$、$r=1$。}

\originalnote{118}[最]{最早对系统处于状态$\varphi$时的行为作出统计断言的是 M. Born；Dirac 和 Jordan 作了更深入的论述。参看\CJKglue \hyperlink{note-8}{原注 \getrefnumber{orig-note-8}} 及\CJKglue \hyperlink{note-2}{原注 \getrefnumber{orig-note-2}} 所引文献。}

\originalnote{119}[每]{每个物理量都对应一个厄米算符，这种对应关系将在 \hyperref[sec:IV.1]{IV.1} 中详述。目前，根据 \hyperref[sec:I.2]{I.2}，我们只知道：坐标对应于算符$q_1,\allowbreak \ldots,\allowbreak q_k$，动量对应于算符$\frac{h}{2\pi i}\frac{\partial}{\partial q_1},\allowbreak \ldots,\allowbreak \frac{h}{2\pi i}\frac{\partial}{\partial q_k}$，能量对应于“能量算符”$\mathsf H$。}

\originalnote{120}[气]{气体动理论可以很好地说明这些关系。

一摩尔氧气（16克）含有$6\cdot10^{23}$个$\mathrm O_2$分子。\tnote[。][每][这]{这里的16克及$\mathrm O_2$分子数均按原文保留。}每个$\mathrm O_2$分子由两个氧原子构成。若忽略原子的内部结构，就须把每个原子看成具有三个自由度的质点。因此，这是一个具有$2\cdot3\cdot6\cdot10^{23}=36\cdot10^{23}=k$个自由度的力学系统。只要给定确定系统状态所需的$2k$个量，本可按照因果规律确定它的行为。但气体理论只使用两个量，即压力和温度；它们是上述$2k$个量的某些复杂函数。因此，气体理论只能作出统计的，也就是概率的断言。在许多情形中，这些断言几乎成为因果性的，即相应概率接近$0$或$1$；但这丝毫不改变事情在原则上的性质。}

\originalnote{121}[设]{设$F_t(\lambda)$是依赖于$t$的函数，$\frac{\partial}{\partial t}F_t(\lambda)=G_t(\lambda)$，而$\mathsf H$是厄米算符，则有$\frac{\partial}{\partial t}F_t(\mathsf H)=G_t(\mathsf H)$；因为求导是由减法、除法以及取极限构成的。对$F_t(\lambda)=e^{-\frac{2\pi i}{h}(t-t_0)\lambda}$，就得到
\[
\frac{\partial}{\partial t}e^{-\frac{2\pi i}{h}(t-t_0)\mathsf H}
=-\frac{2\pi i}{h}\mathsf H e^{-\frac{2\pi i}{h}(t-t_0)\mathsf H}.
\]
把它作用于$\varphi$，便得到所需的微分方程。

由于$|F_t(\lambda)|=1$，故$F_t(\lambda)\overline{F_t(\lambda)}=1$，因而$F_t(\mathsf H)\{F_t(\mathsf H)\}^*=1$。也就是说，我们的$F_t(\mathsf H)=e^{-\frac{2\pi i}{h}(t-t_0)\mathsf H}$是酉的。又因$t=t_0$时它显然等于$1$，所以$\varphi_{t_0}=\varphi$也成立。}

\originalnote{122}[\emph{Z}]{\emph{Z. Physik}，第37卷（1926）。此后的全部发展，参看\CJKglue \hyperlink{note-2}{原注 \getrefnumber{orig-note-2}}，都以这种看法为基础。}

\originalnote{123}[\emph{P}]{\emph{Physic. Rev.}，第26卷（1925）。另参看 W. Bothe 在\emph{Handbuch der Physik}第23卷（Quanten）中的综述，柏林，1926年，第3章，尤其第73节。}

\originalnote{124}[B]{Bohr、Kramers 和 Slater 提出的关于基本现象的统计理论，就是以这一基本思想为基础的。参看\hypertarget{origpage-251}{}\phantomsection\label{origpage-251}\emph{Z. Physik}第24卷（1924），以及\CJKglue \hyperlink{note-123}{原注 \getrefnumber{orig-note-123}}。Compton--Simon 实验可以看成对这一观点的反证。}

\originalnote{125}[J]{Jordan 认识到，这些跃变同旧 Bohr 量子理论中“量子跃迁”的观念有联系，见\emph{Z. Physik}，第40卷（1924）。}

\originalnote{126}[上]{上述一切都有一个前提：被观察系统和测量仪器的结构，也就是一切起作用的力场等等，都是确切已知的；需要确定的只是系统的状态，即各坐标的瞬时值。若这一理想化假设不成立，自然还会有其他不确定性来源。

我们对不精确测量的描述，也作了一定的理想化：我们假设，测量以绝对的确定性，判明某个值是否位于区间$I=\{\lambda',\allowbreak \lambda''\}$内，其中$\lambda'<\lambda''$。实际上，边界$\lambda',\allowbreak \lambda''$是模糊的；也就是说，判定只能以一定概率作出。不过，至少目前看来，我们采用的描述在数学上是最为合用的。}

\originalnote{127}[请]{请借助 \hyperref[post:III.1.W]{W} 核实最后这个断言！对属于$R$及$S$的单位分解，按照 \hyperref[sec:II.8]{II.8} 处理即可。}

\originalnote{128}[对]{对于能够同时测量的物理量$\mathfrak R,\allowbreak \mathfrak S$，我们已证明了这条规律：$\mathfrak R+\mathfrak S$对应的算符，是$\mathfrak R$与$\mathfrak S$所对应算符之和。这条规律对任意$\mathfrak R,\allowbreak \mathfrak S$也成立，但此时不能再证明它；它在这里是本理论的一项基本公设。参看 \hyperref[sec:IV.1]{IV.1} 末尾，以及 \hyperref[sec:IV.2]{IV.2} 的讨论。}

\originalnote{129}[关]{关于“在$\mathfrak M$中的各状态$\varphi$下的同时可测性”的详细讨论，以及对不能绝对精确测量的$\mathfrak R,\allowbreak \mathfrak S$，即连续谱等情形的讨论，留给读者。处理方法同 \hyperref[sec:III.3]{III.3} 完全类似。}

\originalnote{130}[算]{算符运算如下：
\begin{align*}
\varepsilon^2&=((R-\varrho1)^2\varphi,\varphi)
=(R^2\varphi,\varphi)-2\varrho(R\varphi,\varphi)+\varrho^2\\
&=\|R\varphi\|^2-2(R\varphi,\varphi)^2+(R\varphi,\varphi)^2
=\|R\varphi\|^2-(R\varphi,\varphi)^2;
\end{align*}
$\eta^2$同理。}

\originalnote{131}[\emph{Z}]{\emph{Z. Physik}，第43卷（1927）。Bohr 对这些考察作了深化，见\emph{Naturwiss.}第16卷（1928）。下文采用的解析方法，首先由 Kennard 着手研究，见\emph{Z. Physik}第44卷（1926）；现在的形式则由 Robertson 给出。}

\originalnote{132}[B]{Bohr 已强调这一点的根本意义，见\CJKglue \hyperlink{note-131}{原注 \getrefnumber{orig-note-131}} 所引文献。顺便说明，以下描述在某一点上并非完全经典：我们利用了光量子的存在，也就是频率为$\nu$的光，从不以小于$h\nu$的能量份额出现这一事实。}

\originalnote{133}[下]{下面的讨论依照 Heisenberg 和 Bohr，见\CJKglue \hyperlink{note-131}{原注 \getrefnumber{orig-note-131}}。}

\originalnote{134}[例]{例如可参看 Einstein 的原始论文，\emph{Ann. Physik}第14卷（1905）；或任何一本现代教科书。}

\originalnote{135}[关]{关于显微镜的理论，例如参看\emph{Handbuch der Physik}，柏林，1927年，第18卷，第2章 G。测量若非常精确，$\varepsilon$就很小，$\lambda$也很小；因此应使用$\gamma$射线，或波长更短的射线。普通透镜在这种条件下不能使用。只有那些既不被$\gamma$射线击碎、也不被它移离原位的分子，才能用来制作透镜。这类分子或粒子的存在，并不违反任何已知的自然规律，因此在思想实验中允许使用。}

\originalnote{136}[]{\hypertarget{origpage-252}{}\phantomsection\label{origpage-252}$\frac{mc\,\Delta\nu}{\nu}$比$\frac{h\,\Delta\nu}{c}$大得多，意味着$\nu$比$\frac{mc^2}{h}$小得多，也就是$\overline E=h\nu$比$mc^2$小得多。换言之，相对于粒子$\boldsymbol T$的相对论静质量所对应的能量，光量子$\boldsymbol L$的能量很小；而作非相对论计算时，这一本来就是不可避免的假设。}

\originalnote{137}[例]{例如，设$F(x)$是频率为$\nu_0$、只从$0$延伸到$\tau$的有限单色波列：
\[
F(x)=\begin{cases}\alpha\sin2\pi\nu_0x,&0\le x\le\tau,\\0,&\text{其他情形。}\end{cases}
\]
为使衔接处连续，必须有$\sin2\pi\nu_0\tau=0$，即$\nu_0=\frac n{2\tau}$，$n=1,\allowbreak 2,\allowbreak 3,\allowbreak \ldots$。根据 Fourier 积分的已知反演公式，参看\CJKglue \hyperlink{note-87}{原注 \getrefnumber{orig-note-87}}，有$a_\nu^2=b_\nu^2+c_\nu^2$，其中\tnote[中][][以]{以下原印推导中$\pi$与$2\pi$的因子并不一致，这里照原文保留。}
\begin{align*}
\left.\begin{matrix}b_\nu\\c_\nu\end{matrix}\right\}
&=2\int_{-\infty}^{+\infty}F(x)\begin{matrix}\cos\\\sin\end{matrix}2\pi\nu x\,dx
=2\alpha\int_0^\tau\sin2\pi\nu_0x\cdot\begin{matrix}\cos\\\sin\end{matrix}2\pi\nu x\,dx\\
&=\pm\alpha\int_0^\tau\left(\begin{matrix}\sin\\\cos\end{matrix}\pi(\nu+\nu_0)x
-\begin{matrix}\sin\\\cos\end{matrix}\pi(\nu-\nu_0)x\right)dx\\
&=-\alpha\left[\frac{\begin{matrix}\cos\\\sin\end{matrix}\pi(\nu+\nu_0)x}{\pi(\nu+\nu_0)}
-\frac{\begin{matrix}\cos\\\sin\end{matrix}\pi(\nu-\nu_0)x}{\pi(\nu-\nu_0)}\right]_0^\tau\\
&=\begin{cases}
-\alpha\left[\dfrac{(-1)^n\cos\pi\nu\tau-1}{\pi(\nu+\nu_0)}-\dfrac{(-1)^n\cos\pi\nu\tau-1}{\pi(\nu-\nu_0)}\right],\\[6pt]
-\alpha\left[\dfrac{(-1)^n\sin\pi\nu\tau}{\pi(\nu+\nu_0)}-\dfrac{(-1)^n\sin\pi\nu\tau}{\pi(\nu-\nu_0)}\right]
\end{cases}\\
&=\begin{cases}
-\dfrac{2\alpha\nu_0(1-(-1)^n\cos\pi\nu\tau)}{\pi(\nu^2-\nu_0^2)},\\[6pt]
\dfrac{2\alpha\nu_0(-1)^n\sin\pi\nu\tau}{\pi(\nu^2-\nu_0^2)}\text{。}
\end{cases}
\end{align*}
所以
\begin{align*}
a_\nu&=\frac{2\alpha\nu_0\sqrt{2-2(-1)^n\cos\pi\nu\tau}}{\pi(\nu^2-\nu_0^2)}\\
&=\frac{4\alpha\nu_0\left|\begin{matrix}\sin\\\cos\end{matrix}\frac12\pi\nu\tau\right|}{\pi(\nu^2-\nu_0^2)}
=\frac{4\alpha\nu_0|\sin\pi(\nu-\nu_0)\tau|}{\pi(\nu^2-\nu_0^2)}\text{。}
\end{align*}
可见，频率$\nu=\nu_0$附近的成分最强，波列能量的大部分集中在$\pi(\nu-\nu_0)\tau$取适中数值的频率范围内。所以，$\nu-\nu_0$的散布程度，或同样地$\nu$的散布程度，其数量级为$1/\tau$。精确计算有关的表达式，也得到同样的结果：
\[
\frac{\int_0^\infty a_\nu^2(\nu-\nu_0)^2\,d\nu}{\int_0^\infty a_\nu^2\,d\nu}\text{。}
\]}

\originalnote{138}[\emph{P}]{\emph{Proc. Roy. Soc., Lond.}，第114卷（1927）。另参看 Weyl 的\emph{Gruppentheorie und Quantenmechanik}，第2版，第91页及以下，莱比锡，1931年。}

\originalnote{139}[有]{有兴趣的读者，可在任何一本电动力学教科书中找到光的电磁理论，例如 Abraham--Becker 的\emph{Theorie der Elektrizität}，柏林，1930年。下文属于 Maxwell 理论范围的论述，也可参照这部书。}

\originalnote{140}[参]{参看 R. Courant 和 D. Hilbert 的\emph{Methoden der mathematischen Physik I}，第358—363页，柏林，1924年。}

\originalnote{141}[有]{有
\[
\iiint_{\boldsymbol H}[\mathfrak a,\operatorname{rot}\mathfrak b]\,dx\,dy\,dz
=\iiint_{\boldsymbol H}[\operatorname{rot}\mathfrak a,\mathfrak b]\,dx\,dy\,dz,
\]\hypertarget{origpage-253}{}\phantomsection\label{origpage-253}因为
\[
[\mathfrak a,\operatorname{rot}\mathfrak b]-[\operatorname{rot}\mathfrak a,\mathfrak b]
=\operatorname{grad}(\mathfrak a\times\mathfrak b)
\]
（$\mathfrak a\times\mathfrak b$是$\mathfrak a,\allowbreak \mathfrak b$的所谓外积，即向量积），条件是$\mathfrak a\times\mathfrak b$的法向分量在$\boldsymbol H$的边界上为零。由于$\mathfrak a\times\mathfrak b$既垂直于$\mathfrak a$，也垂直于$\mathfrak b$，所以只要$\mathfrak a$或$\mathfrak b$垂直于$\boldsymbol H$的边界，上述条件必然满足。这里$\mathfrak a=\operatorname{rot}\overline{\mathfrak A}_m$、$\mathfrak b=\overline{\mathfrak A}_n$，所以所要求的边界条件必然满足。\tnote[。][][原]{原文恒等式的右端印作$\operatorname{grad}(\mathfrak a\times\mathfrak b)$。按本节采用的通常向量微分记号，此处应为$-\operatorname{div}(\mathfrak a\times\mathfrak b)$。所给边界条件使相应曲面积分为零，因此前面的分部积分等式不受这一误印影响。原式予以保留。}}

\originalnote{142}[例]{例如参看\CJKglue \hyperlink{note-138}{原注 \getrefnumber{orig-note-138}} 所引文献。}

\originalnote{143}[由]{由于$P_\nu^x$同$Q_\nu^y,\allowbreak Q_\nu^z$对易，却不同$Q_\nu^x$对易，要说明接下来的改写成立，就须证明以下关系；这里略去多余下标，并以$F$代替$\mathfrak A$：
\[
PF(Q)-F(Q)P=\frac h{2\pi i}F'(Q),
\]
其中$P=\frac h{2\pi i}\frac\partial{\partial q}\cdots$、$Q=q\cdots$。这一关系对矩阵理论尤其重要，直接计算很容易验证。}

\originalnote{144}[要]{要看出上述指标组$k,\allowbreak M_1,\allowbreak M_2,\allowbreak \ldots$确实构成一个序列，最简单的办法如下。设$\pi_1,\allowbreak \pi_2,\allowbreak \pi_3,\allowbreak \ldots$是素数序列$2,\allowbreak 3,\allowbreak 5,\allowbreak \ldots$。乘积$\pi_1^{k-1}\pi_2^{M_1}\pi_3^{M_2}\cdots$实际上是有限乘积，因为除有限多个例外外，全部$M_n=0$，相应因子$\pi_{n+1}^{M_n}=1$。当$k,\allowbreak M_1,\allowbreak M_2,\allowbreak \ldots$遍历我们的全部指标组时，$\pi_1^{k-1}\pi_2^{M_1}\pi_3^{M_2}\cdots$遍历全部数$1,\allowbreak 2,\allowbreak 3,\allowbreak \ldots$，而且每个数恰好出现一次。因此，可以用$\pi_1^{k-1}\pi_2^{M_1}\pi_3^{M_2}\cdots$，对$\alpha_{kM_1M_2\ldots}$依次用单个整数下标编号。}

\originalnote{145}[例]{例如，可以把光量子的动量$p_x,\allowbreak p_y,\allowbreak p_z$作为坐标，再加上一个描述偏振状态的坐标$\pi$。$p_x,\allowbreak p_y,\allowbreak p_z$决定光量子的方向，即方向余弦$\alpha_x,\allowbreak \alpha_y,\allowbreak \alpha_z$，其中$\alpha_x^2+\alpha_y^2+\alpha_z^2=1$，也决定频率$\nu$、波长$\lambda$及能量。因为按照 Einstein，动量向量的大小是$h\nu/c$，参看\CJKglue \hyperlink{note-134}{原注 \getrefnumber{orig-note-134}}，所以
\[
p_x=\frac{h\nu}{c}\alpha_x,\qquad p_y=\frac{h\nu}{c}\alpha_y,\qquad p_z=\frac{h\nu}{c}\alpha_z,
\]
即
\begin{gather*}
\nu=\frac ch\sqrt{p_x^2+p_y^2+p_z^2},\qquad \lambda=\frac c\nu,\qquad \text{能量}=h\nu,\\
\alpha_x=\frac{cp_x}{h\nu},\qquad\alpha_y=\frac{cp_y}{h\nu},\qquad\alpha_z=\frac{cp_z}{h\nu}\text{。}
\end{gather*}
这里出现了一个不便：我们的本征振动$\overline{\mathfrak A}_n(xyz)\cdot\gamma\cos2\pi\varrho_n(t-\tau)$是光的驻波；由于空腔$\boldsymbol H$的壁面反射，光波必定如此。这就不能给$\overline{\mathfrak A}_n$指定一个“射线方向”$\alpha_x,\allowbreak \alpha_y,\allowbreak \alpha_z$。立即可以看出，除$\alpha_x,\allowbreak \alpha_y,\allowbreak \alpha_z$外，至少还须有相反方向$-\alpha_x,\allowbreak -\alpha_y,\allowbreak -\alpha_z$；动量也有相应的情形。因此，在$\boldsymbol H$中，必须采用$p_x,\allowbreak p_y,\allowbreak p_z,\allowbreak \pi$之外的坐标。

近来关于这一问题的一些论述，用如下办法克服了这个困难：把$\boldsymbol H$取为长方体
\[
-A<x<A,\qquad-B<y<B,\qquad-C<z<C,
\]
但边界面$x=\pm A$、$y=\pm B$、$z=\pm C$不作为壁面，而是把$x=A$同$x=-A$、$y=B$同$y=-B$、$z=C$同$z=-C$分别认作同一处。也就是说，光在$A,\allowbreak y,\allowbreak z$处到达面$x=A$时，\hypertarget{origpage-254}{}\phantomsection\label{origpage-254}便从$-A,\allowbreak y,\allowbreak z$处继续原来的路径，重新进入$\boldsymbol H$，就像什么也没有发生一样；其余各面同理。例如参看 L. Landau 和 R. Peierls 的论文，\emph{Z. Physik}第62卷（1930）。也可以说，空间在$x,\allowbreak y,\allowbreak z$方向分别以$2A,\allowbreak 2B,\allowbreak 2C$为周期。

解析处理依然相同，只是边界条件改为
\begin{gather*}
\mathfrak A(A,y,z)=\mathfrak A(-A,y,z),\qquad
\mathfrak A(x,B,z)=\mathfrak A(x,-B,z),\\
\mathfrak A(x,y,C)=\mathfrak A(x,y,-C),
\end{gather*}
而不再是边界上的$\frac\partial{\partial n}\mathfrak A=0$；用来作展开的“基本解”则为\tnote[为][][原]{原文以下相位中的因子是$c$。按前面各式，$c$为光速，所以空间项应当含$1/c$。这里保留原印相位。}
\[
\begin{matrix}\cos\\\sin\end{matrix}[2\pi\nu(t-c(\alpha_xx+\alpha_yy+\alpha_zz))],
\]
用它代替$\overline{\mathfrak A}(xyz)\cdot\widetilde\varrho(t)$。不难确定属于本征解的$\nu=\varrho_n$及$\alpha_x=\alpha_{n,\allowbreak x}$、$\alpha_y=\alpha_{n,\allowbreak y}$、$\alpha_z=\alpha_{n,\allowbreak z}$，其中$n=1,\allowbreak 2,\allowbreak \ldots$；此后理论的展开，就同正文所给的处理一致。}

\originalnote{146}[取]{取极限$S\to+\infty$，同电磁理论中所取的极限$N\to+\infty$并不相同！即使把那里出现的$M_1,\allowbreak M_2,\allowbreak \ldots$解释成光量子的数目，$N$也是不相干光量子数目的一个上界；所谓不相干，是指频率、方向及偏振并非全都相同，其中频率与方向合起来决定动量，参看\CJKglue \hyperlink{note-145}{原注 \getrefnumber{orig-note-145}}。而$S$则是不加其他区分的光量子总数的上界。}

\originalnote{147}[若]{若不借助量子力学，而只考察这一措施的后果，那么引入$\overline{\mathfrak R}^{(S)}_\infty$来代替$\mathfrak R^{(S)}_\infty$，就相当于把通常的统计法改为所谓 Bose--Einstein 统计法。例如参看\CJKglue \hyperlink{note-138}{原注 \getrefnumber{orig-note-138}} 所引 Dirac 的论文。}

\originalnote{148}[关]{关于人们如何理解这种“二重性”、又为何认为它悖谬，读者可在当时的文献中找到更详细的讨论，例如\CJKglue \hyperlink{note-6}{原注 \getrefnumber{orig-note-6}} 所引著作。

人们常说，量子力学也赋予物质同样的二重性，因为离散的粒子，如电子、质子，也用波函数描述，并表现出典型的波动性质，例如在光栅上的衍射等。［参看 Davison--Germer 的实验，\emph{Physic. Rev.}第30卷（1927），\emph{Proc. nat. Acad. Sci. U.S.A.}第14卷（1928）；另见 C. P. Thompson，\emph{Proc. Roy. Soc., Lond.}第117卷（1928），以及 Rupp，\emph{Ann. Physik}第85卷（1928）。］对此必须强调，量子力学是从一种统一的基本现象理论中，推导出这两种“本性”的。旧量子理论的悖谬之处，是为了解释经验，不得不交替诉诸两个彼此矛盾的理论：Maxwell--Hertz 的电磁光理论和 Einstein 的光量子理论。}

\originalnote{149}[W]{Weisskopf 和 Wigner 给出了这些微分方程的精确解，见\emph{Z. Physik}第63卷（1930）；从中也可以看出，这些论断是正确的。}

\originalnote{150}[众]{众所周知，N. Bohr 于1913年确定：原子从能量为$W^{(1)}$的定态，跃迁到能量为$W^{(2)}$的定态时，发出频率为$\frac{W^{(1)}-W^{(2)}}h$的光；当然，$W^{(1)}>W^{(2)}$。参看\CJKglue \hyperlink{note-5}{原注 \getrefnumber{orig-note-5}} 所引文献。在我们这里，对应的就是$\frac{W_{\bar k}-W_k}h$。}

\originalnote{151}[有]{有
\begin{align*}
|e^{ix}-1|^2&=(e^{ix}-1)\overline{(e^{ix}-1)}=(e^{ix}-1)(e^{-ix}-1)\\
&=2-e^{ix}-e^{-ix}=2-2\cos x=2(1-\cos x).
\end{align*}}

\originalnote{152}[有]{\hypertarget{origpage-255}{}\phantomsection\label{origpage-255}有如下关系，参看 Courant--Hilbert，第49页：
\[
\int_{-\infty}^{\infty}\frac{1-\cos x}{x^2}\,dx
=2\int_0^{\infty}\frac{1-\cos x}{x^2}\,dx
=\int_0^{\infty}\frac{1-\cos(2y)}{y^2}\,dy
=2\int_0^{\infty}\frac{\sin^2y}{y^2}\,dy=\pi.
\]}

\originalnote{153}[算]{算符$P_\nu^x$同全部$Q_\mu^x,\allowbreak Q_\mu^y,\allowbreak Q_\mu^z,\allowbreak P_\mu^x,\allowbreak P_\mu^y,\allowbreak P_\mu^z$对易，只有$Q_\nu^x$例外。具体地说，
\[
P_\nu^xQ_\nu^x-Q_\nu^xP_\nu^x=\frac h{2\pi i}1.
\]
因此
\[
\mathsf H_0Q_\nu^x-Q_\nu^x\mathsf H_0
=\frac1{2m_\nu}(P_\nu^x)^2Q_\nu^x-Q_\nu^x\frac1{2m_\nu}(P_\nu^x)^2
=\frac h{2\pi i}\frac1{m_\nu}P_\nu^x;
\]
参看\CJKglue \hyperlink{note-143}{原注 \getrefnumber{orig-note-143}}。}

\originalnote{154}[\emph{P}]{\emph{Physik. Z.}，第18卷（1917）。}

\originalnote{155}[必]{必须强调，在概念上要区分系统本身和处于某一状态的系统。例如，氢原子是一个系统，即一个电子和一个质子，两者之间有已知的力作用。形式上，这个系统由以下规定描述：位形空间为六维，坐标记为$q_1,\allowbreak \ldots,\allowbreak q_6$，动量记为$p_1,\allowbreak \ldots,\allowbreak p_6$，哈密顿函数为\tnote[为][][原]{原文库仑项前为正号。电子与质子相吸引时，该项应为负；这里保留原印公式。}
\begin{align*}
H(q_1\cdots q_6p_1\cdots p_6)
&=\frac{p_1^2+p_2^2+p_3^2}{2m_e}+\frac{p_4^2+p_5^2+p_6^2}{2m_p}\\
&\quad+\frac{e^2}{\sqrt{(q_1-q_4)^2+(q_2-q_5)^2+(q_3-q_6)^2}}.
\end{align*}
要确定状态，还须作进一步的规定：在经典力学中，要给坐标和动量$q_1,\allowbreak \ldots,\allowbreak q_6,\allowbreak p_1,\allowbreak \ldots,\allowbreak p_6$赋予数值$q_1^0,\allowbreak \ldots,\allowbreak q_6^0,\allowbreak p_1^0,\allowbreak \ldots,\allowbreak p_6^0$；在量子力学中，则要给定波函数$\varphi(q_1\cdots q_6)$。所需规定到此为止。系统与状态一经给定，理论便提供唯一确定的计算规则，用来回答所有问题。}

\originalnote{156}[这]{这类系综称为“集体”（Kollektive），是从频率观点建立概率论所必需的。R. von Mises 引入了这一概念，认识到它对概率论的意义，并据此建立了理论。例如参看他的著作\emph{Wahr\-schein\-lich\-keit, Statistik und ihre Wahrheit}，柏林，1928年。}

\originalnote{157}[]{$w(a')$是$a\le a'$的概率，也就是$a$属于区间$-\infty,\allowbreak a'$的概率。可见，$w(a)$具有以下性质；为了明确显出它对$\mathfrak R$的依赖，以下写作$w_{\mathfrak R}(a)$：当$a\to-\infty$时，$w_{\mathfrak R}(a)\to0$；当$a\to+\infty$时，$w_{\mathfrak R}(a)\to1$；当$a\ge a_0$且$a\to a_0$时，$w_{\mathfrak R}(a)\to w_{\mathfrak R}(a_0)$；当$a'\le a''$时，$w_{\mathfrak R}(a')\le w_{\mathfrak R}(a'')$。在量子力学中，若$E(\lambda)$是属于$R$的单位分解，则$w_{\mathfrak R}(a)=\|E(a)\varphi\|^2=(E(a)\varphi,\allowbreak \varphi)$。

如果$w_{\mathfrak R}(a)$可微，可以用通常的“概率密度”$\frac d{da}w_{\mathfrak R}(a)$来代替它。如果它在$a=a_0$处不连续，当然是从左侧而言，那么$a=a_0$这一点具有“离散概率”$w_{\mathfrak R}(a_0)-w_{\mathfrak R}(a_0-0)$。但在一般情况下都适用的概念，仍是$w_{\mathfrak R}(a)$，参看\CJKglue \hyperlink{note-156}{原注 \getrefnumber{orig-note-156}} 所引文献。}

\originalnote{158}[这]{这是所谓大数定律，即 Bernoulli 定理的结果。}

\originalnote{159}[例]{\hypertarget{origpage-256}{}\phantomsection\label{origpage-256}例如，在定义电场时，人们曾把下面这一点看作原则性的困难：为此所用的带电“试探物体”，不可能小于一个电子。}

\originalnote{160}[确]{确定值$a_0$对应于概率函数$w_{\mathfrak R}(a)$，其中
\[
w_{\mathfrak R}(a)=\begin{cases}1,&a\ge a_0,\\0,&a<a_0.\end{cases}
\]
方差$\varepsilon^2$为零，当且仅当概率函数具有上述形式。顺便指出，均值$\varrho$及方差$\varepsilon^2$按下式计算，积分为 Stieltjes 积分：
\begin{align*}
\varrho&=\int_{-\infty}^{+\infty}a\,dw_{\mathfrak R}(a),\\
\varepsilon^2&=\int_{-\infty}^{+\infty}(a-\varrho)^2\,dw_{\mathfrak R}(a)\\
&=\int_{-\infty}^{+\infty}a^2\,dw_{\mathfrak R}(a)-2\varrho\int_{-\infty}^{+\infty}a\,dw_{\mathfrak R}(a)+\varrho^2\\
&=\int_{-\infty}^{+\infty}a^2\,dw_{\mathfrak R}(a)-\varrho^2\\
&=\int_{-\infty}^{+\infty}a^2\,dw_{\mathfrak R}(a)-\left(\int_{-\infty}^{+\infty}a\,dw_{\mathfrak R}(a)\right)^2.
\end{align*}
参看 \hyperref[sec:III.4]{III.4} 及\CJKglue \hyperlink{note-130}{原注 \getrefnumber{orig-note-130}}。}

\originalnote{161}[这]{这里说的是在多个系统上作独立的测量；若在同一系统上相继测量，得到的将总是同一个值，参看 \hyperref[sec:III.3]{III.3}。}

\originalnote{162}[例]{例如，可以想一想：若把$\mathfrak R,\allowbreak \mathfrak S$取为$q$，即一个直角坐标，以及动量$p$，会发生什么？由于不确定性关系，这两个量不能同时测量。如果一个系综中$q$的散布程度很小，那么以准确度$\varepsilon$，即相应散布程度$\varepsilon$测量$p$，就会把$q$的散布程度至少增大到$h/(4\pi\varepsilon)$，参看 \hyperref[sec:III.4]{III.4}；也就是说，原先使$q$的取值集中的效果，全被破坏了。}

\originalnote{163}[]{$a\mathfrak R$、$\mathfrak S^2$、$\mathfrak R+\mathfrak S+\cdots$的意思是：按照上面给出的定义，把物理量$\mathfrak R$，或$\mathfrak S$，或$\mathfrak R,\allowbreak \mathfrak S,\allowbreak \ldots$，分别代入函数$f(x)=ax$、$f(x)=x^2$或$f(x,\allowbreak y,\allowbreak \ldots)=x+y+\cdots$。}

\originalnote{164}[例]{例如，在 Heisenberg 理论中，一个在势场$V(xyz)$中运动的电子，其能量算符为
\[
\mathsf H_0=\frac{(P^x)^2+(P^y)^2+(P^z)^2}{2m}+V(Q^x,Q^y,Q^z),
\]
例如参看 \hyperref[sec:III.6]{III.6}；它是两个不对易的算符之和：
\[
R=\frac{(P^x)^2+(P^y)^2+(P^z)^2}{2m},\qquad S=V(Q^x,Q^y,Q^z).
\]
测量属于$R$的物理量$\mathfrak R$，是动量测量；测量属于$S$的物理量$\mathfrak S$，是坐标测量。但属于$\mathsf H_0=R+S$的物理量$\mathfrak R+\mathfrak S$，则要用完全不同的方式测量。例如，可以测量这个束缚电子发出的谱线的频率，再依据 Bohr 频率关系确定能级，即$\mathfrak R+\mathfrak S$的值。尽管如此，任何情况下仍有
\[
\operatorname{Erw}(\mathfrak R+\mathfrak S)=\operatorname{Erw}(\mathfrak R)+\operatorname{Erw}(\mathfrak S).
\]}

\originalnote{165}[不]{不过，对于零方差系综，并没有理由不采用正确归一化的期望值。}

\originalnote{166}[即]{即
\[
U\varphi_m=\sum_nu_{mn}\varphi_n,
\]\hypertarget{origpage-257}{}\phantomsection\label{origpage-257}这当然要求$\sum_n|u_{mn}|^2$有限。例如，可以如下证明这一点。若$\sum_n|x_n|^2=1$，则对$\varphi=\sum_nx_n\varphi_n$，$R=P_{[\varphi]}$具有矩阵$\bar x_\mu x_\nu$，所以相应物理量$\mathfrak R$的期望值为$\sum_{mn}u_{nm}\bar x_mx_n$。由于$P_{[\varphi]}=P_{[\varphi]}^2$、$1-P_{[\varphi]}=(1-P_{[\varphi]})^2$，这个期望值$\ge0$且$\le\operatorname{Erw}(1)$；所以，至少在期望值函数$\operatorname{Erw}(\mathfrak R)$已归一化时，它$\ge0$且$\le1$。若$x_{N+1}=x_{N+2}=\cdots=0$，这意味着：当$\sum_{n=1}^N|x_n|^2=1$时，$N$维厄米型$\sum_{m,\allowbreak n=1}^Nu_{nm}\bar x_mx_n$的值$\ge0$且$\le1$；也就是说，矩阵$u_{\mu\nu}$，$\mu,\allowbreak \nu=1,\allowbreak \ldots,\allowbreak N$，的本征值都$\ge0$且$\le1$。因此，向量$y_m=\sum_{n=1}^Nu_{mn}x_n$的长度，总不大于向量$x_m$的长度。取
\[
x_m=\begin{cases}1,&m=\bar m,\\0,&\text{其他情形},\end{cases}
\]
则$y_m=u_{m\bar m}$，从而
\[
\sum_{m=1}^N|x_m|^2\ge\sum_{m=1}^N|y_m|^2,\qquad
1\ge\sum_{m=1}^N|u_{m\bar m}|^2.
\]
由于这对每个$N$都成立，所以$\sum_n|u_{\bar mn}|^2\le1$。}

\originalnote{167}[严]{严格地说，整个论证只有在$\varphi_1,\allowbreak \varphi_2,\allowbreak \ldots$都属于$R$的定义域时，才是有效的。对每个$R$，我们确实可以找到这样的完全规范正交系$\varphi_1,\allowbreak \varphi_2,\allowbreak \ldots$，参看 \hyperref[sec:II.11]{II.11}。但若$R$并非处处有定义，这个正交系就依赖于$R$。因此，实际得到的情形是：对每一个完全规范正交系$\varphi_1,\allowbreak \varphi_2,\allowbreak \ldots$，有一个依赖于该正交系的$U$，使得
\[
\operatorname{Erw}(\mathfrak R)=\operatorname{Spur}(UR)
\]
只须对定义域包含$\varphi_1,\allowbreak \varphi_2,\allowbreak \ldots$的那些$R$成立。

但所有这些$U$仍彼此相等。因为，若$U',\allowbreak U''$是其中两个，那么当$R$处处有定义时，上式对二者都成立，即$\operatorname{Spur}(U'R)=\operatorname{Spur}(U''R)$。令$R=P_{[\varphi]}$，便得到$(U'\varphi,\allowbreak \varphi)=(U''\varphi,\allowbreak \varphi)$，即$((U'-U'')\varphi,\allowbreak \varphi)=0$。由于这对全部满足$\|\varphi\|=1$的$\varphi$都成立，因而也对希尔伯特空间的全部元素成立，所以必须有$U'-U''=0$，即$U'=U''$。}

\originalnote{168}[不]{不能直接令$\mathfrak S=\sqrt{\mathfrak R}$，即$\mathfrak S=h(\mathfrak R)$、$h(x)=\sqrt x$；因为我们只考虑在所有实数$x$上有定义的实值函数，而$\sqrt x$并非如此：当$x$为负时，它是虚数。}

\originalnote{169}[也]{也可以作一个简单的直接证明。设$\varphi_1,\allowbreak \varphi_2,\allowbreak \ldots$是完全规范正交系，$a_{\mu\nu}=(A\varphi_\mu,\allowbreak \varphi_\nu)$、$b_{\mu\nu}=(B\varphi_\mu,\allowbreak \varphi_\nu)$，并有$\operatorname{Spur}(AB)=\sum_{\mu\nu}a_{\mu\nu}b_{\nu\mu}$。只要所有的和$\sum_{\mu,\allowbreak \nu=1}^Na_{\mu\nu}b_{\nu\mu}$都$\ge0$，这个量就$\ge0$。若$f=\sum_{\mu=1}^Nx_\mu\varphi_\mu$，则
\[
(Af,f)=\sum_{\mu,\nu=1}^Na_{\mu\nu}x_\mu\bar x_\nu\ge0,\qquad
(Bf,f)=\sum_{\mu,\nu=1}^Nb_{\mu\nu}x_\mu\bar x_\nu\ge0;
\]
所以有限矩阵$a_{\mu\nu},\allowbreak b_{\mu\nu}$，$\mu,\allowbreak \nu=1,\allowbreak \ldots,\allowbreak N$，同样定号。定号性以及$\sum_{\mu,\allowbreak \nu=1}^Na_{\mu\nu}b_{\nu\mu}$的值，都是$N$维空间中的正交不变量。由于$b_{\mu\nu}$是厄米的，可以用这个$N$维空间中的正交变换，把它化为对角形。因此，不妨一开始就假设它为对角矩阵，即当$\mu\ne\nu$时$b_{\mu\nu}=0$。于是
\[
\sum_{\mu,\nu=1}^Na_{\mu\nu}b_{\nu\mu}=\sum_{\mu=1}^Na_{\mu\mu}b_{\mu\mu}.
\]
但根据定号性，有$a_{\mu\mu}\ge0$、$b_{\mu\mu}\ge0$；令$\nu=\mu$时$x_\nu=1$，$\nu\ne\mu$时$x_\nu=0$即可看出。因此，上述和确实$\ge0$。}

\originalnote{170}[当]{\hypertarget{origpage-258}{}\phantomsection\label{origpage-258}当$\varphi'=\varphi''$时，这是显然的。设$\varphi'\ne\varphi''$；把$\varphi',\allowbreak \varphi''$“正交化”（见 \hyperref[sec:II.2]{II.2}），得到与$\varphi'$正交的$\varphi_1$，使$\|\varphi_1\|=1$，且$\varphi''$为$\varphi',\allowbreak \varphi_1$的线性组合。因此$\varphi''=a\varphi'+b\varphi_1$，$\|\varphi''\|^2=|a|^2+|b|^2=1$。例如，令$|a|=\cos\theta,\allowbreak |b|=\sin\theta$，则$a=e^{i\alpha}\cos\theta,\allowbreak b=e^{i\beta}\cos\theta$。\tnote[。][这][原]{原书在这里的$b$表达式中印作$\cos\theta$；前后各式要求$\sin\theta$。}这时，取$a(x)=e^{ix\alpha}\cos x\theta,\allowbreak b(x)=e^{ix\beta}\sin x\theta$，仍有$|a(x)|^2+|b(x)|^2=1$，所以$\varphi(x)=a(x)\varphi'+b(x)\varphi_1$满足$\|\varphi(x)\|=1$。而且$\varphi(x)$从$\varphi'$（$x=0$）连续变化到$\varphi''$（$x=1$）。}

\originalnote{171}[严]{严格说来，由于\CJKglue \hyperref[remark:IV.2.2]{第2项}，还应要求$V\ne0,\allowbreak W\ne0$；不过，$V=0$或$W=0$的情形，分别已包含在$c'=0,\allowbreak c''=1$或$c'=1,\allowbreak c''=0$中。}

\originalnote{172}[前]{前两节中导向纯系综的推导，由作者发表于\emph{Gött. Nachr.}，1927。纯系综的存在及其与一般系综的关系，由H.~Weyl（\emph{Z. Physik}，第46卷，1927）和作者（上述论文）各自独立发现。一般系综的一个特殊情形，即两个耦合系统的情形（见 \hyperref[sec:VI.2]{VI.2} 的讨论），更早已出现于J.~Landau，\emph{Z. Physik}，第45卷（1927）。}

\originalnote{173}[设]{设把所有“隐参数”合记为$\pi$。如果它们只取离散值$\pi_1,\allowbreak \pi_2,\allowbreak \ldots,\allowbreak \pi_n$（$n>1$），就把$\pi=\pi_1$的系统放入一个系综，把$\pi\ne\pi_1$的系统放入另一个系综，便得到两个系综，其混合正是原来的系综。如果$\pi$在区域$\Pi$中连续变化，取$\Pi$的子区域$\Pi'$；一个系综包含$\pi$属于$\Pi'$的系统，另一个包含$\pi$不属于$\Pi'$的系统。}

\originalnote{174}[相]{相空间是六维的；它的六个坐标，是质点的三个直角坐标$q_1,\allowbreak q_2,\allowbreak q_3$及相应的三个动量$p_1,\allowbreak p_2,\allowbreak p_3$。依 \hyperref[sec:III.4]{III.4}，相应的散布程度$\varepsilon_1,\allowbreak \varepsilon_2,\allowbreak \varepsilon_3,\allowbreak \eta_1,\allowbreak \eta_2,\allowbreak \eta_3$满足
\[
\varepsilon_1\eta_1\ge\frac h{4\pi},\qquad
\varepsilon_2\eta_2\ge\frac h{4\pi},\qquad
\varepsilon_3\eta_3\ge\frac h{4\pi},
\]
即
\[
\varepsilon_1\varepsilon_2\varepsilon_3\eta_1\eta_2\eta_3\ge\left(\frac h{4\pi}\right)^3.
\]
因此，经典力学相空间中的位置，无论如何都有这一体积大小的不确定范围。}

\originalnote{175}[参]{参见Schrödinger对此问题极为透彻的讨论：\emph{Naturwiss.}，第17卷（1929），第37期。}

\originalnote{176}[]{$u_\mu\to-u_\mu$（或$u_\nu\to-u_\nu$，或$v_\mu\to-v_\mu$）是$K$的对称操作，它使前面那些被积函数变号，所以其积分为零。操作$u_\mu\to v_\mu,\allowbreak v_\mu\to u_\mu$，或$u_\mu\to u_\nu,\allowbreak u_\nu\to u_\mu$，则是$K_1$的对称操作，能使后面那些积分互相变换。因此，这些积分相等，各自都等于它们的和的$1/(2s)$：
\[
\iint\cdots\iint_K(u_1^2+v_1^2+\cdots+u_s^2+v_s^2)\,dv
=\iint\cdots\iint_Kdv=\text{$K$的表面积},
\]
把这个量记为$C$。}

\originalnote{177}[若]{若$\operatorname{Spur}U$无限，最后一个由相减得到的公式或许显得可疑。不过也可如下证明：$\mathfrak R$的值为$\lambda^*$，是说在$a''<\lambda^*$或$a'\ge\lambda^*$时，有$w(a',\allowbreak a'')=0$，即$\operatorname{Spur}U(E(a'')-E(a'))=0$。这个迹总非负，随$a''$不减，随$a'$不增；所以当$a''<\lambda^*$时，只须考察$a'\to-\infty$的极限；当$a'\ge\lambda^*$时，只须考察$a''\to+\infty$的极限。因此，$a''<\lambda^*$时有$\operatorname{Spur}UE(a'')=0$，而$a'\ge\lambda^*$时有$\operatorname{Spur}U(1-E(a'))=0$。}

\originalnote{178}[也]{\hypertarget{origpage-259}{}\phantomsection\label{origpage-259}也就是说，若$\mathfrak E$成立，状态就是$\varphi$；若它不成立，则“非$\mathfrak E$”成立。此时，$E=P_{[\varphi]},\allowbreak \mathfrak M=[\varphi]$被$1-E=1-P_{[\varphi]}$、$\mathfrak R-\mathfrak M=\mathfrak R-[\varphi]$代替，而后者并不能唯一确定$U$。（$E$恰好对应于问题“状态是否为$\varphi$？”）要使测量无论得到什么结果都能唯一确定状态，就须测量这样一个物理量$\mathfrak R$：其算符$R$具有完全由单重本征值构成的纯点谱（见 \hyperref[sec:III.3]{III.3}）。测量后，系统处于状态$\varphi_1,\allowbreak \varphi_2,\allowbreak \ldots$之一，即$R$的某个本征函数所表示的状态；所以测量一般会改变$\boldsymbol S$的状态。类似地，对$\mathfrak E$的测量也改变状态：结果为肯定时，事后有$U=P_{[\varphi]}$；结果为否定时，则有$U(1-P_{[\varphi]})=U,\allowbreak UP_{[\varphi]}=0$，即$U\varphi=0$；而事前这两个条件可能都不成立。因此，量子力学对状态的这种“确定”，同时改变了状态；这本来就是可以预料的。}

\originalnote{179}[例]{例如，因为
\begin{align*}
RP_{[\varphi_n]}f&=R((f,\varphi_n)\cdot\varphi_n)
=(f,\varphi_n)\cdot R\varphi_n=\lambda_n(f,\varphi_n)\cdot\varphi_n,\\
P_{[\varphi_n]}Rf&=(Rf,\varphi_n)\cdot\varphi_n
=(f,R\varphi_n)\cdot\varphi_n=\lambda_n(f,\varphi_n)\cdot\varphi_n.
\end{align*}}

\originalnote{180}[任]{任何一本经典（哈密顿）力学教科书都给出这些关系。}

\originalnote{181}[关]{关于时间—能量这一对量的不确定性关系，已有过多次讨论。参见Heisenberg的综合论述：\emph{Die physikalischen Prinzipien der Quantentheorie}，第II.2.d节，Leipzig，1930。}

\originalnote{182}[测]{测量装置的一切其他细节，都只为把真正关心的物理量$\mathfrak R$或其算符$R$，与这里提到的$M_n$、$P^x,\allowbreak P^y,\allowbreak P^z$或$Q^x,\allowbreak Q^y,\allowbreak Q^z$耦合起来。不过，这恰好是测量技术在实践中最重要的一面。}

\originalnote{183}[例]{例如参见Nernst，\emph{Theoretische Chemie}，Stuttgart（自1893年以来多次出版），第四编关于“质量作用定律”的热力学证明的讨论。}

\originalnote{184}[读]{读者可以在许多教科书中，找到以此为基础建立的热力学唯象体系，例如Planck，\emph{Vorlesungen über Thermodynamik}，Berlin，1930。对于下文，主要需要的是基本定律的统计表述，例如以下论文中的论述：Einstein，\emph{Verh. d. dtsch. physik. Ges.}，第12卷（1914）；L.~Szilárd，\emph{Z. Physik}，第32卷（1925）。}

\originalnote{185}[这]{这是统计力学的Maxwell--Boltzmann方法（参见P.和T.~Ehrenfest在\emph{Enzykl. d. Math. Wiss.}，第II.4.D卷，Leipzig，1907中的文章所作的综述）。例如在气体理论中，“十分复杂”的系统就是由许多相互作用的分子组成的气体；被统计研究的则是分子。}

\originalnote{186}[这]{这是Gibbs的方法（见\CJKglue \hyperlink{note-185}{原注 \getrefnumber{orig-note-185}} 所引文献）。这里把整个气体作为一个单独系统，追踪同一个系统的许多份样本（即同一种气体的许多份样本）各自的演变，再对它们的性质作统计分析。}

\originalnote{187}[见]{见\CJKglue \hyperlink{note-184}{原注 \getrefnumber{orig-note-184}} 所引文献。L.~Szilárd对此又作了进一步发展。}

\originalnote{188}[众]{众所周知，气体动理论描述了容器壁如何使所包围的气体取得壁的温度。参见\CJKglue \hyperlink{note-184}{原注 \getrefnumber{orig-note-184}}、\hyperlink{note-185}{\getrefnumber{orig-note-185}} 所引文献。}

\originalnote{189}[若]{若在这个变化过程中，分别在温度$\mathbf{T}_1,\allowbreak \ldots,\allowbreak \mathbf{T}_i$消耗了热量$\mathbf{Q}_1,\allowbreak \ldots,\allowbreak \mathbf{Q}_i$，则熵差为$\frac{\mathbf{Q}_1}{\mathbf{T}_1}+\cdots+\frac{\mathbf{Q}_i}{\mathbf{T}_i}$。参见\CJKglue \hyperlink{note-184}{原注 \getrefnumber{orig-note-184}} 所引文献。}

\originalnote{190}[若]{\hypertarget{origpage-260}{}\phantomsection\label{origpage-260}若$c(\mathbf T)$为给定数量的气体在温度$\mathbf T$时的热容，则它从$\mathbf T$升温到$\mathbf T+d\mathbf T$时，吸收的热量为$c(\mathbf T)\,d\mathbf T$。因此，依\CJKglue \hyperlink{note-186}{原注 \getrefnumber{orig-note-186}}，\tnote[，][熵][原]{原书此处引用原注186；相应的熵差公式实际见\CJKglue \hyperlink{note-189}{原注 \getrefnumber{orig-note-189}}。这里保留原引用。}熵差为$\int_{\mathbf T_1}^{\mathbf T_2}\frac{c(\mathbf T)\,d\mathbf T}{\mathbf T}$。}

\originalnote{191}[对]{对理想气体，$c(\mathbf T)$为常数；但$\mathbf T$很小时，这一点一定失效。例如参见\CJKglue \hyperlink{note-6}{原注 \getrefnumber{orig-note-6}} 所引文献。}

\originalnote{192}[为]{为此，要求$\overline{\boldsymbol K}$的体积$V$比$\boldsymbol K_1,\allowbreak \ldots,\allowbreak \boldsymbol K_N$的总体积大得多；还要求“每个自由度的能量”$k\mathbf T$（$k$为Boltzmann常数）比$\frac{h^2}{\mu V^{2/3}}$大得多，其中$h$为Planck常数，$\mu$为单个分子的质量；这个量具有能量的量纲。例如参见Fermi，\emph{Z. Physik}，第36卷（1926）。}

\originalnote{193}[例]{例如参见\CJKglue \hyperlink{note-184}{原注 \getrefnumber{orig-note-184}} 所引文献。}

\originalnote{194}[参]{参见\CJKglue \hyperlink{note-185}{原注 \getrefnumber{orig-note-185}} 所引文献。关于“Maxwell妖”概念所涉及的困难，读者可在L.~Szilárd，\emph{Z. Physik}，第53卷（1929）中找到详细讨论。}

\originalnote{195}[若]{若理想气体由$M$个分子组成，其压强为$p=\frac{M\varkappa\mathbf T}{\mathcal V}$。因此，从体积$\mathcal V_1$压缩到$\mathcal V_2$时，所作的机械功为
\[
\int_{\mathcal V_1}^{\mathcal V_2}p\,d\mathcal V
=M\varkappa\mathbf T\int_{\mathcal V_1}^{\mathcal V_2}\frac{d\mathcal V}{\mathcal V}
=M\varkappa\mathbf T\ln\frac{\mathcal V_2}{\mathcal V_1}.
\]
在这里，$M=N/2,\allowbreak \mathcal V_1=\mathcal V/2,\allowbreak \mathcal V_2=\mathcal V$。}

\originalnote{196}[众]{众所周知，理想气体的能量只取决于温度。}

\originalnote{197}[直]{直接求$aP_{[\varphi]}+bP_{[\psi]}$的本征值。要求
\[
(aP_{[\varphi]}+bP_{[\psi]})f=\lambda f.
\]
左端为$\varphi,\allowbreak \psi$的线性组合，右端因而也是；所以当$\lambda\ne0$时，$f$也是这样的线性组合。$\lambda=0$当然是无限重本征值，因为每个与$\varphi,\allowbreak \psi$正交的$f$都属于它。因此，只须考察$\lambda\ne0$及$f=x\varphi+y\psi$。（设$\varphi,\allowbreak \psi$线性无关；否则$\varphi=c\psi,\allowbreak |c|=1$，两个状态就相同。）

上述方程于是成为
\[
a(x+y(\psi,\varphi))\cdot\varphi+b(x(\varphi,\psi)+y)\cdot\psi
=\lambda x\cdot\varphi+\lambda y\cdot\psi,
\]
即
\[
ax+a\overline{(\varphi,\psi)}y=\lambda x,\qquad
b(\varphi,\psi)x+by=\lambda y.
\]
这些方程的行列式必须为零：
\begin{gather*}
\begin{vmatrix}a-\lambda&a\overline{(\varphi,\psi)}\\b(\varphi,\psi)&b-\lambda\end{vmatrix}=0,
\qquad(a-\lambda)(b-\lambda)-ab|(\varphi,\psi)|^2=0,\\
\lambda^2-(a+b)\lambda+ab(1-\alpha^2)=0,\\
\lambda=\frac{a+b\pm\sqrt{(a+b)^2-4ab(1-\alpha^2)}}2
=\frac{a+b\pm\sqrt{(a-b)^2+4\alpha^2ab}}2.
\end{gather*}
分别取$a=1,\allowbreak b=0$，或$a=\frac12,\allowbreak b=\frac12$，或$a=\frac13,\allowbreak b=\frac23$，就得到正文中的公式。}

\originalnote{198}[因]{\hypertarget{origpage-261}{}\phantomsection\label{origpage-261}因为$(x\ln x)'=\ln x+1$，从而
\begin{align*}
\left(\frac{1+y}2\ln\frac{1+y}2+\frac{1-y}2\ln\frac{1-y}2\right)'
&=\frac12\left(\ln\frac{1+y}2+1\right)-\frac12\left(\ln\frac{1-y}2+1\right)\\
&=\frac12\ln\frac{1+y}{1-y},
\end{align*}
所以所考察表达式的导数为
\begin{align*}
&-3\cdot\frac12\ln\frac{3+\sqrt{1+8\alpha^2}}{3-\sqrt{1+8\alpha^2}}
\cdot\frac13\frac{8\alpha}{\sqrt{1+8\alpha^2}}+4\cdot\frac12\ln\frac{1+\alpha}{1-\alpha}\\
&\qquad=2\left(\ln\frac{1+\alpha}{1-\alpha}
-\frac{2\alpha}{\sqrt{1+8\alpha^2}}\ln\frac{3+\sqrt{1+8\alpha^2}}{3-\sqrt{1+8\alpha^2}}\right).
\end{align*}
要使它大于零，就是要有
\[
\ln\frac{1+\alpha}{1-\alpha}>
\frac{2\alpha}{\sqrt{1+8\alpha^2}}\ln\frac{3+\sqrt{1+8\alpha^2}}{3-\sqrt{1+8\alpha^2}},
\]
即
\[
\frac1{2\alpha}\ln\frac{1+\alpha}{1-\alpha}>
\frac23\cdot\frac1{2\beta}\ln\frac{1+\beta}{1-\beta},
\qquad\beta=\frac{\sqrt{1+8\alpha^2}}3.
\]
我们甚至要证明把$\frac23$换成$\frac89$后，这一不等式仍成立。由于$1-\beta^2=\frac89(1-\alpha^2)$且$\alpha<\beta$（后者由前式及$\alpha<1$得出），这也就是
\[
\frac{1-\alpha^2}{2\alpha}\ln\frac{1+\alpha}{1-\alpha}>
\frac{1-\beta^2}{2\beta}\ln\frac{1+\beta}{1-\beta}.
\]
只须证明$\frac{1-x^2}{2x}\ln\frac{1+x}{1-x}$在$0<x<1$内单调递减即可。例如，由它的幂级数就能看出：
\begin{align*}
\frac{1-x^2}{2x}\ln\frac{1+x}{1-x}
&=(1-x^2)\left(1+\frac{x^2}3+\frac{x^4}5+\cdots\right)\\
&=1-\left(1-\frac13\right)x^2-\left(\frac13-\frac15\right)x^4-\cdots.
\end{align*}}

\originalnote{199}[这]{这种技巧是唯象热力学方法所特有的；例如参见\CJKglue \hyperlink{note-184}{原注 \getrefnumber{orig-note-184}} 所引文献。}

\originalnote{200}[当]{当然，这里也可以略去因子$N\varkappa$，只考察$-\operatorname{Spur}(U\ln U)$。或者，为了保持与元素个数$N$成正比，也可取$-N\operatorname{Spur}(U\ln U)$。}

\originalnote{201}[例]{例如参见Planck，\emph{Theorie der Wärmestrahlung}，Leipzig，1913。}

\originalnote{202}[L]{L.~Szilárd在\CJKglue \hyperlink{note-194}{原注 \getrefnumber{orig-note-194}} 所引论文中证明：取得这种“知识”，至少要以$\varkappa\ln2$的熵增作为补偿。一般说来，知道两个可能情形中究竟哪一个发生，其“热力学价值”就是$\varkappa\ln2$。一切试图在不知道分子位于容器哪一半的情况下实行上述过程的方案，都可以证明不能成立；尽管这些方案有时会导向相当复杂的自动机构。}

\originalnote{203}[对]{对宏观观察者的这种刻画出自E.~Wigner。}

\originalnote{204}[若]{若$E_n$与$\mathsf H$对易，从而也与$A$对易，等式确实成立，因为
\begin{align*}
\operatorname{Spur}(A\cdot UA^{-1}E_n)
&=\operatorname{Spur}(UA^{-1}E_n\cdot A)\\
&=\operatorname{Spur}(UA^{-1}AE_n)=\operatorname{Spur}(UE_n).
\end{align*}
但是，并非一切$E_n$，也就是说，并非一切宏观可观测量，\hypertarget{origpage-262}{}\phantomsection\label{origpage-262}都与$\mathsf H$对易：其中很多量，例如扩散过程中气体的重心，会随$t$明显变化，所以$\operatorname{Spur}(UE_n)$并不恒定。由于一切宏观量彼此对易，$\mathsf H$绝不可能是宏观量；即能量不可能用宏观方法测得完全精确。这一点本来也很合理。}

\originalnote{205}[关]{关于经典$\mathsf H$定理，参见Boltzmann，\emph{Vorlesungen über Gastheorie}，Leipzig，1896，以及P.和T.~Ehrenfest在\CJKglue \hyperlink{note-185}{原注 \getrefnumber{orig-note-185}} 所引文献中极具启发性的讨论。在量子力学中，可代替Boltzmann相应假设的“无序假设”，由W.~Pauli提出（\emph{Sommerfeld-Festschrift}，1928），并用来证明$\mathsf H$定理。最近，作者也证明了经典力学的遍历定理，见\emph{Proc. Nat. Ac.}，1930年1月与3月；另见G.~D.~Birkhoff的加强结果，\emph{Proc. Nat. Ac.}，1929年12月与1930年3月。}

\originalnote{206}[\emph{Z}]{\emph{Z. Physik}，第57卷（1929）。\tnote[。][][原]{原书把本注编号印作266；根据它在原注205与207之间的位置及正文引用，可确定应为206。}}

\originalnote{207}[N]{N.~Bohr，\emph{Naturwiss.}，第17卷（1929），首先指出：量子力学在形式上使自然描述的二重性不可避免，而这种二重性也有实质上的根据；他还指出了这种二重性与心物平行原则的联系。}

\originalnote{208}[下]{下面以及 \hyperref[sec:VI.3]{VI.3} 的讨论，包含若干重要内容，作者得益于与L.~Szilárd先生的交谈。另参见Heisenberg在\CJKglue \hyperlink{note-181}{原注 \getrefnumber{orig-note-181}} 所引文献中的类似考察。}

\originalnote{209}[不]{不难证明：若$\mathsf A$厄米或超极大，则$\boldsymbol A$也具有相应性质。}

\originalnote{210}[对]{对 \hyperref[cond:IV.2.I]{I}，这一点显然；对 \hyperref[cond:IV.2.II]{II}，只要涉及的仅是多项式，也同样显然。对于一般函数，则可由下述事实看出：我们的映射$\mathsf A\to\boldsymbol A$不破坏单位分解与厄米算符之间的对应关系。}

\originalnote{211}[由]{由于指标很多而且种类不同，这里采用的矩阵记法与先前略有区别。}

\originalnote{212}[]{$I+II$中一个状态的投影，一般分别是$I$和$II$中的混合态（见下文）。这一事实由J.~Landau发现，\emph{Z. Physik}，第45卷（1927）。}

\originalnote{213}[这]{这里的数学讨论依据E.~Schmidt的一篇论文：\emph{Math. Ann.}，第63卷（1907）。}

\originalnote{214}[这]{这一试探形式还容许若干别的变体，不过出于类似理由，仍须不予考虑。}

\originalnote{215}[对]{对 \hyperref[sec:III.4]{III.4} 所讨论的位置测量情形，相应计算见Weizsäcker的一篇论文：\emph{Z. Physik}，第70卷（1931）。}

% 中文编写中间材料：frontmatter.tex
% 源PDF顺序页1按用户要求略去；下文从顺序页2开始。
\begin{center}
{\large Grundlehren der mathematischen Wissenschaften 38\par}
\medskip
数学科学基础丛书\par
\textkaishu{数学领域的系统研究丛书}\par
\vfill
{\Large Johann von Neumann\par}
\bigskip
{\huge\bfseries Mathematische Grundlagen der Quantenmechanik\par}
\bigskip
{\huge\bfseries 量子力学的数学原理\par}
\bigskip
{\large 重印本\par}
\vfill
\includegraphics[height=4\baselineskip]{figures/publisher-emblem-reprint.png}\par
\medskip
Springer 出版社\tnote[社][][本]{本页及后面扉页中的图案是出版社徽记；其中的字母S代表Springer出版社。}\par
柏林\quad 海德堡\quad 纽约\quad 1981
\end{center}
\thispagestyle{empty}
\clearpage
% 源PDF顺序页3。
\section*{出版资料}
德国图书馆在版编目简目：

Neumann, Johann von：\par
\emph{Mathematische Grundlagen der Quantenmechanik} / Johann von Neumann 著。——重印本［据柏林 Springer 1932年版］。——柏林、海德堡、纽约：Springer，1981年。\par
（Grundlehren der mathematischen Wissenschaften；第38卷）

ISBN-13：978-3-642-96049-9\par
DOI：\href{https://doi.org/10.1007/978-3-642-96048-2}{10.1007/978-3-642-96048-2}\par
电子版 ISBN-13：978-3-642-96048-2

AMS 学科分类（1970）：\par
主要分类：81 A 03，81 A 09\par
次要分类：46 N 05，47 A 05，47 A 10

本书受著作权保护。由此产生的各项权利，尤其是翻译、重印、摘取插图、广播、以照相机械或类似方式复制，以及存入数据处理设备的权利，均予保留；即使只使用本书的部分内容，也不例外。依《著作权法》第54条，为商业目的制作复制件，须向出版社支付报酬，具体数额由双方商定。

\copyright\ Springer-Verlag Berlin/Heidelberg，1971年。\par
1971年精装第一版的平装重印本。

照相制版重印：Beltz Offsetdruck，Hemsbach/Bergstrasse。\par
装订：Graphischer Betrieb Konrad Triltsch，Würzburg。\par
2140/3014-543

\clearpage
% 源PDF顺序页4。
\begin{center}
{\Large\bfseries 数学科学基础丛书\par}
\medskip
分题论述，并特别注重应用领域\par
\bigskip
协同编纂\par
W. Blaschke（汉堡）、M. Born（哥廷根）、C. Runge（哥廷根）\par
\medskip
R. Courant（哥廷根）主编\par
\bigskip
第 XXXVIII 卷\par
\bigskip
{\LARGE\bfseries 量子力学的数学原理\par}
\bigskip
Johann v. Neumann 著\par
\vfill
\includegraphics[height=4\baselineskip]{figures/publisher-emblem-1932.png}\par
\bigskip
柏林\par
Julius Springer 出版社\par
1932
\end{center}
\clearpage
% 源PDF顺序页5。
\begin{center}
{\LARGE\bfseries Mathematische Grundlagen der Quantenmechanik\par}
\bigskip
{\LARGE\bfseries 量子力学的数学原理\par}
\bigskip
{\Large Johann v. Neumann 著\par}
\bigskip
附插图4幅\par
\vfill
\includegraphics[height=4\baselineskip]{figures/publisher-emblem-1932.png}\par
\bigskip
柏林\par
Julius Springer 出版社\par
1932
\end{center}
\clearpage
% 源PDF顺序页6。
\null\vfill
\begin{center}
\large 谨以此书纪念先父
\end{center}
\vfill\null
\clearpage
% 源PDF顺序页7：由全书标题生成带动态页码的目录。
\tableofcontents


% 中文编写中间材料：introduction-I.tex
% 源PDF顺序页8，原印第1页。
\chapter*{导言}\label{chap:introduction}
\addcontentsline{toc}{chapter}{导言}\hypertarget{origpage-1}{}\phantomsection\label{origpage-1}本书要对新的量子力学作出统一的论述，并在可能而且适宜的范围内，做到数学上严密可靠。近几年来，新的量子力学已经形成了一种在主要方面看来可望定型的形式，即所谓“变换理论”。本书着重讨论与这一理论相关的一般问题和根本问题，尤其要深入考察其中困难的诠释问题；这些问题有不少至今尚未彻底澄清。其中，量子力学同统计学以及经典统计力学的关系尤为重要。至于如何把量子力学方法应用于个别问题，以及从一般理论中分化出来的各种专门理论，本书通常不作讨论；只要不妨碍读者理解各部分之间的一般联系，就尽量略去。这样处理还有一个理由：这些方面已经有好几部出色的著作，或正在陆续出版。\origNote[。][]{1}

另一方面，本书将介绍这一理论所需的数学工具，也就是希尔伯特空间及其所谓厄米算符的理论。为此，无界算符也必须详细讨论，因而需要超出 Hilbert 以及 E. Hellinger、F. Riesz、E. Schmidt、O. Toeplitz 等人所建立的传统理论范围。这里先说明所采用的方法：我们通常直接对算符本身进行运算，因为物理量正是由这些算符表示的；不以矩阵作为运算的出发点，因为只有在希尔伯特空间中任意选定一个具体坐标系之后，才得到算符的矩阵表示。这种不依赖坐标、也就是具有不变性的处理方法，带有鲜明的几何特点，并且在形式上有相当大的优点。

Dirac 也给出了不依赖坐标的量子力学表述，其简洁、优美几乎无出其右；他的数篇论文和最近出版的\hypertarget{origpage-2}{}\phantomsection\label{origpage-2}著作采用的都是这种表述。\origNote[。][因]{2}因此，我们也许应当在这里说明，为什么采用一种同他很不相同的方法。

Dirac 的上述方法清楚而优美，如今量子力学文献中已有很大一部分采用了它。然而，它完全不符合数学严格性的要求；即使把这种要求适当地、合理地降低到理论物理通常采用的程度，情况也仍然如此。例如，这种方法始终坚持一个虚构：每个自伴算符都能化为对角形。对于那些实际上不能如此处理的算符，就不得不引入具有自相矛盾性质的“广义”函数。有时，即使目的只在于数值计算一个以直观方式规定的实验的结果，也难以避免这类数学“虚构”。假如这些不能纳入现有分析学框架的概念，果真是新物理理论所不可缺少的，那么这并不构成反对它的理由。Newton 力学曾促成微积分的诞生，而微积分在当时的形式中无疑存在自相矛盾；同样，量子力学也可能要求我们重新建立“无穷多变量的分析学”。那就应当改动数学工具，而不是物理理论。但事实绝非如此。本书将表明，变换理论也可以用同样清楚、统一的方式，建立在数学上可靠的基础之上。必须强调，要正确地建立这一理论，并不是把 Dirac 的方法表述得更加精确、详尽，而是须从一开始就采用不同的途径，即以 Hilbert 的算符谱理论为依托。

在分析根本问题时，我们将特别说明，如何从几个定性的基本假设推导出量子力学的统计公式。此外，还要详细讨论：量子力学之所以具有统计性质，是否因为我们对自然界的描述还不能唯一确定结果，也就是说，这种描述还不完备？这种解释最符合一条一般原则：任何概率陈述，都来自我们知识的不完备。人们的确曾多次提出这种“隐参数”解释，也提出过与它相近的另一种解释：把“隐参数”归于观察者，而非被观察的系统。然而，结果表明，这样的解释很难令人满意；\hypertarget{origpage-3}{}\phantomsection\label{origpage-3}更准确地说，它同量子力学的某些定性基本公设并不相容。\origNote[。][]{3}

本书还考察量子统计理论同热力学的关系。深入研究表明，经典力学中那些众所周知的困难，也就是同建立热力学所需的“无序假设”有关的困难，在这里可以克服。\origNote[。][]{4}

% 源PDF顺序页11，原印第4页。
\chapter{初步考察}\label{chap:I}\hypertarget{origpage-4}{}\phantomsection\label{origpage-4}\section{变换理论的产生}\label{sec:I.1}
量子理论在1900年至1925年间取得了巨大成就；Planck、Einstein 和 Bohr 的名字标志着这段发展历程。\origNote[。][这]{5}这里无须详述这些成就。到这一时期结束时，几乎已经无可怀疑地确定：一切基本过程，也就是原子和分子尺度上发生的一切，都受量子“不连续”规律的支配。几乎各个方面也都有了定量的量子理论方法，所得结果通常同经验符合得很好，或至少尚能令人满意。还有一项具有根本意义的进展：理论物理研究中已接受这样一种观念——我们在宏观世界中所感知的连续性原则，即“自然不作跳跃”（\emph{natura non facit saltus}），不过是平均过程造成的表象，世界本质上仍是不连续的。人通常一次感知的是大量基本过程的总和，其数目为$10^{24}$的许多倍；\tnote[；][于][原]{原文的\emph{Quadrillionen}采用德语长阶命数法，单位为$10^{24}$，这里直接写出数值。}于是，抹平一切差异的大数定律，就把个别过程的真实面貌完全掩盖了。

然而，在上述时期，量子理论还没有一个数学和物理的体系，能把当时已知的一切统一起来；更谈不上具有力学、电动力学和相对论原有体系那样恢宏而严整的统一性——这种统一性已经被量子现象打破。量子理论要求普遍有效，显然有充分理由，却缺少相应的形式工具和思想工具。它把本质上不同、各自独立、彼此异质，甚至部分相互矛盾的各种方案混合在一起，难以理清。其中最突出的，一是仍有一半属于经典力学和电动力学的对应原理；不过，这一原理在后来澄清问题时将发挥决定性作用。二是光的本性中自相矛盾的二重性，即波与微粒，参看\CJKglue \hyperlink{note-5}{原注 \getrefnumber{orig-note-5}} 和\CJKglue \hyperlink{note-148}{原注 \getrefnumber{orig-note-148}}。此外，还存在不量子化的运动，即非周期运动，以及量子化的运动，即周期运动或多重周期运动。\origNote[。][]{6}

1925年，人们找到了解决办法。Heisenberg 提出的一种思路，经 Born、Heisenberg、Jordan，以及稍后的 Dirac，发展成了新的\hypertarget{origpage-5}{}\phantomsection\label{origpage-5}量子理论体系；这是物理学第一次拥有完整的量子理论体系。不久之后，Schrödinger 从完全不同的出发点建立了“波动力学”，取得了同样的成果，而且很快就证明，它同 Heisenberg--Born--Jordan 及 Dirac 的体系等价；至少在数学意义上如此，参看 \hyperref[sec:I.3]{I.3} 和 \hyperref[sec:I.4]{I.4}。\origNote[。][B]{7}Born 从统计的角度诠释了量子理论对自然界的描述。\origNote[。][在]{8}在此基础上，Dirac 和 Jordan 把两种理论融合成一种理论，即“变换理论”。\origNote[。][两]{9}两者在其中相互补充，从而能够以数学上特别简单的方式处理物理问题。

还应当提到，尽管这已不属于本书的主要论题：Goudsmit 和 Uhlenbeck 又发现了电子的磁矩和角动量，此后旧量子理论的困难几乎全部消失，我们由此有了一个几乎完全令人满意的力学体系。诚然，开头提到的同电动力学及相对论的大统一，尚未恢复；但至少已经有了一种普遍有效的力学。量子规律以不言而喻的必然性纳入其中，而我们的绝大部分经验也得到了令人满意的解释。\origNote[。][]{10}

\section{量子力学最初的表述}\label{sec:I.2}
为了先有一个大概的认识，我们来简述 Heisenberg--Born--Jordan 的“矩阵力学”和 Schrödinger 的“波动力学”各自提出的基本问题。

两种理论的出发点，都是一个由 哈密顿函数$H(q_1,\allowbreak \ldots,\allowbreak q_k,\allowbreak p_1,\allowbreak \ldots,\allowbreak p_k)$给定的经典力学问题。（其含义是力学教科书中详述的以下内容：系统具有$k$个自由度，因此只要给定$k$个坐标$q_1,\allowbreak \ldots,\allowbreak q_k$的数值，就确定了系统在该时刻的状态。能量是这些坐标及其时间导数的一个给定函数：
\[
 E=L(q_1,\ldots,q_k,\dot q_1,\ldots,\dot q_k),
\]
通常是导数$\dot q_1,\allowbreak \ldots,\allowbreak \dot q_k$的二次式。以
\[
 p_1=\frac{\partial L}{\partial\dot q_1},\ldots,
 p_k=\frac{\partial L}{\partial\dot q_k}
\]
引入坐标$q_1,\allowbreak \ldots,\allowbreak q_k$的“共轭动量”$p_1,\allowbreak \ldots,\allowbreak p_k$；在上述关于$L$的假设下，它们线性地依赖于\hypertarget{origpage-6}{}\phantomsection\label{origpage-6}$\dot q_1,\allowbreak \ldots,\allowbreak \dot q_k$。无论如何，我们可以借助$p_1,\allowbreak \ldots,\allowbreak p_k$，从$L$中消去$\dot q_1,\allowbreak \ldots,\allowbreak \dot q_k$，得到
\[
 E=L(q_1,\ldots,q_k,\dot q_1,\ldots,\dot q_k)
  =H(q_1,\ldots,q_k,p_1,\ldots,p_k).
\]
这个$H$就是 哈密顿函数。）两种理论都希望由这一 哈密顿函数，尽可能多地了解系统的真实行为，也就是量子理论意义上的行为：\origNote[：][首]{11}首先确定可能的能级，再认识与之对应的“定态”，计算各定态之间的“跃迁概率”，等等。\origNote[。][]{12}

矩阵理论给出的求解规则如下：寻找一个由$2k$个矩阵$Q_1,\allowbreak \ldots,\allowbreak Q_k,\allowbreak P_1,\allowbreak \ldots,\allowbreak P_k$组成的系统。\origNote[。][第]{13}第一，它们要满足关系
\leavevmode\phantomsection\label{rel:I.2.commutation}
\begin{align*}
 Q_mQ_n-Q_nQ_m&=0,& P_mP_n-P_nP_m&=0,\\
 P_mQ_n-Q_nP_m&=
 \begin{cases}0,&m\ne n,\\ \dfrac{h}{2\pi i}\,1,&m=n,\end{cases}
 &&(m,n=1,\ldots,k);
\end{align*}
第二，矩阵$W=H(Q_1,\allowbreak \ldots,\allowbreak Q_k,\allowbreak P_1,\allowbreak \ldots,\allowbreak P_k)$要成为对角矩阵。（这里不追究这些方程的来源，尤其不追究第一组方程，也就是所谓“对易关系”的来源；这一理论的整个非交换矩阵演算都受这些关系支配。\CJKglue \hyperlink{note-1}{原注 \getrefnumber{orig-note-1}} 所引著作中有详尽说明。这里$h$是 Planck 的作用量子。）于是，$W$的对角元，记为$w_1,\allowbreak w_2,\allowbreak \ldots$，就是系统各种可能的能级。矩阵$Q_1,\allowbreak \ldots,\allowbreak Q_k$的元素$q^{(1)}_{mn},\allowbreak \ldots,\allowbreak q^{(k)}_{mn}$，则以一定方式决定系统的跃迁概率及跃迁时发出的辐射；这里所说的跃迁，是从能量为$w_m$的第$m$态，跃迁到能量为$w_n$的第$n$态，且$w_m>w_n$。

还须指出，矩阵
\[
 W=H(Q_1,\ldots,Q_k,P_1,\ldots,P_k)
\]
并不完全由$Q_1,\allowbreak \ldots,\allowbreak Q_k,\allowbreak P_1,\allowbreak \ldots,\allowbreak P_k$及经典 哈密顿函数$H(q_1,\allowbreak \ldots,\allowbreak q_k,\allowbreak p_1,\allowbreak \ldots,\allowbreak p_k)$确定。原因是$Q_l$与$P_l$在乘法下并非全部相互交换；然而，在$H(q_1,\allowbreak \ldots,\allowbreak q_k,\allowbreak p_1,\allowbreak \ldots,\allowbreak p_k)$的经典含义中，区分$p_1q_1$项与$q_1p_1$项之类的做法，根本没有意义。因此，对$H$的规定必须超出这个表达式的经典含义，还要确定各项中$q_l$与$p_l$各因子的次序。一般情况下，这一步尚未完成；但在最重要的特殊情形中，已经知道适当的规定方式。

\hypertarget{origpage-7}{}\phantomsection\label{origpage-7}（不过，在下面这种简单情形中，并没有这样的疑问：所研究的系统由$\nu$个粒子组成，因而有$k=3\nu$个坐标$q_1,\allowbreak \ldots,\allowbreak q_{3\nu}$；例如，$q_{3\mu-2},\allowbreak q_{3\mu-1},\allowbreak q_{3\mu}$是第$\mu$个粒子的三个 直角坐标，$\mu=1,\allowbreak \ldots,\allowbreak \nu$。这些粒子之间的相互作用由势能$V(q_1,\allowbreak \ldots,\allowbreak q_{3\nu})$给定。此时，经典 哈密顿函数为
\begin{multline*}
 H(q_1,\ldots,q_{3\nu},p_1,\ldots,p_{3\nu})
 =\sum_{\mu=1}^{\nu}\frac{1}{2m_\mu}
 (p_{3\mu-2}^2+p_{3\mu-1}^2+p_{3\mu}^2)\\
 {}+V(q_1,\ldots,q_{3\nu}),
\end{multline*}
其中$m_\mu$是第$\mu$个粒子的质量，$p_{3\mu-2},\allowbreak p_{3\mu-1},\allowbreak p_{3\mu}$是它的动量分量。代入矩阵$Q_1,\allowbreak \ldots,\allowbreak Q_{3\nu},\allowbreak P_1,\allowbreak \ldots,\allowbreak P_{3\nu}$之后，这个表达式的意义十分明确；特别是，$V$项不会造成任何困难，因为全部$Q_1,\allowbreak \ldots,\allowbreak Q_{3\nu}$相互交换。）重要的一点是，我们只允许厄米矩阵，即恒满足$a_{mn}=\overline{a_{nm}}$的矩阵$A=\{a_{mn}\}$；元素$a_{mn}$可以是复数！因此，只要$Q_1,\allowbreak \ldots,\allowbreak Q_k,\allowbreak P_1,\allowbreak \ldots,\allowbreak P_k$全是厄米矩阵，$H(Q_1,\allowbreak \ldots,\allowbreak Q_k,\allowbreak P_1,\allowbreak \ldots,\allowbreak P_k)$也必须是厄米的。这使前面提到的因子次序受到一定限制，但这种限制还不够强，不能从经典$H(q_1,\allowbreak \ldots,\allowbreak q_k,\allowbreak p_1,\allowbreak \ldots,\allowbreak p_k)$唯一确定$H(Q_1,\allowbreak \ldots,\allowbreak Q_k,\allowbreak P_1,\allowbreak \ldots,\allowbreak P_k)$。\origNote[。][]{14}

与此相对，波动力学的规则如下。先写出 哈密顿函数$H(q_1,\allowbreak \ldots,\allowbreak q_k,\allowbreak p_1,\allowbreak \ldots,\allowbreak p_k)$；再对系统位形空间中的任意函数$\psi(q_1,\allowbreak \ldots,\allowbreak q_k)$，列出微分方程
\begin{multline*}
 H\left(q_1,\ldots,q_k,\frac{h}{2\pi i}\frac{\partial}{\partial q_1},
 \ldots,\frac{h}{2\pi i}\frac{\partial}{\partial q_k}\right)\psi(q_1,\ldots,q_k)\\
 =\lambda\psi(q_1,\ldots,q_k).
\end{multline*}
这里是位形空间，而非相空间；也就是说，$\psi$中不能出现$p_1,\allowbreak \ldots,\allowbreak p_k$！$H\bigl(q_1,\allowbreak \ldots,\allowbreak q_k,\allowbreak \frac{h}{2\pi i}\frac{\partial}{\partial q_1},\allowbreak \ldots,\allowbreak \frac{h}{2\pi i}\frac{\partial}{\partial q_k}\bigr)$要理解为作用于函数的一种运算，其意义不难明白。例如，在前述情形中，
\begin{multline*}
 H(q_1,\ldots,q_{3\nu},p_1,\ldots,p_{3\nu})
 =\sum_{\mu=1}^{\nu}\frac{1}{2m_\mu}(p_{3\mu-2}^2+p_{3\mu-1}^2+p_{3\mu}^2)\\
 {}+V(q_1,\ldots,q_{3\nu}),
\end{multline*}
这一运算把函数$\psi(q_1,\allowbreak \ldots,\allowbreak q_{3\nu})$变为
\[
 \sum_{\mu=1}^{\nu}\frac{1}{2m_\mu}\left(\frac{h}{2\pi i}\right)^2
 \left(\frac{\partial^2\psi}{\partial q_{3\mu-2}^2}
 +\frac{\partial^2\psi}{\partial q_{3\mu-1}^2}
 +\frac{\partial^2\psi}{\partial q_{3\mu}^2}\right)+V\psi
\]\hypertarget{origpage-8}{}\phantomsection\label{origpage-8}（这里省去了$V$和$\psi$中的变量$q_1,\allowbreak \ldots,\allowbreak q_{3\nu}$）。由于运算$q_1\frac{h}{2\pi i}\frac{\partial}{\partial q_1}$不同于运算$\frac{h}{2\pi i}\frac{\partial}{\partial q_1}q_1$，\origNote[，][因]{15}因此，$H(q_1,\allowbreak \ldots,\allowbreak q_k,\allowbreak p_1,\allowbreak \ldots,\allowbreak p_k)$中$q_m$与$p_m$各因子的次序，又带来了不确定性。不过，Schrödinger 说明了如何把问题归结为一个确定的变分原理，从而消除这种不确定性，并使所得微分方程具有自伴性。\origNote[。][]{16}

这个微分方程，即“波动方程”，完全具有本征值问题的性质：把$\lambda$看成本征值参数，并要求本征函数$\psi=\psi(q_1,\allowbreak \ldots,\allowbreak q_k)$满足某些条件，例如在位形空间，即$q_1,\allowbreak \ldots,\allowbreak q_k$空间的边界上为零，在内部正则且为单值。按照波动理论，$\lambda$的本征值，无论属于点谱还是连续谱，\origNote[，][都]{17}都是可能的能级。相应的复本征函数$\psi$，也同系统的相应状态相关联；这些状态在 Bohr 的意义下是定态。例如，考虑一个由$\nu$个电子组成的系统，$k=3\nu$，参看上文，$e$为电子的电荷。按照 Schrödinger 的看法，其中第$\mu$个电子应当设想为“弥散”在整个$x,\allowbreak y,\allowbreak z$空间中，即$q_{3\mu-2},\allowbreak q_{3\mu-1},\allowbreak q_{3\mu}$空间中。它在点$x,\allowbreak y,\allowbreak z$处测得的电荷密度为
\begin{multline*}
 e\underbrace{\int\cdots\int}_{3\nu-3\mathwordspace{}\text{重}}
 |\psi(q_1\cdots q_{3\mu-3}xyzq_{3\mu+1}\cdots q_{3\nu})|^2\\
 {}\times dq_1\cdots dq_{3\mu-3}\,dq_{3\mu+1}\cdots dq_{3\nu}.
\end{multline*}
（为了使总电荷等于$e$，必须按条件
\[
 \underbrace{\int\cdots\int}_{3\nu\mathwordspace{}\text{重}}
 |\psi(q_1\cdots q_{3\nu})|^2\,dq_1\cdots dq_{3\nu}=1
\]
将$\psi$归一化；这里要对全部$3\nu$个变量积分！对于每一个$\mu=1,\allowbreak \ldots,\allowbreak \nu$，得到的都是同一个方程。）

此外，对于不处于 Bohr 定态的系统，波动力学也能作出论断，\origNote[，][具]{18}具体如下。状态若非定态，即随时间而变化，那么波函数$\psi=\psi(q_1,\allowbreak \ldots,\allowbreak q_k;t)$就含有时间$t$，并按照微分方程\origNote[程][]{19}
\begin{multline*}
 -H\left(q_1,\ldots,q_k,\frac{h}{2\pi i}\frac{\partial}{\partial q_1},
 \ldots,\frac{h}{2\pi i}\frac{\partial}{\partial q_k}\right)\psi(q_1,\ldots,q_k;t)\\
 =\frac{h}{2\pi i}\frac{\partial}{\partial t}\psi(q_1,\ldots,q_k;t)
\end{multline*}
变化。也就是说，可以任意给定$t=t_0$时的$\psi$，此后任意时刻$t$的$\psi$就都唯一确定了。比较 Schrödinger 的两个微分方程可知，即使对应于定态的$\psi$，其实也依赖于$t$；只不过此时$t$以如下方式进入：\hypertarget{origpage-9}{}\phantomsection\label{origpage-9}\[
 \psi(q_1,\ldots,q_k;t)=e^{-\frac{2\pi i}{h}\lambda t}\psi(q_1,\ldots,q_k;0).
\]
也就是说，$t$只出现在一个因子中；这个因子不依赖于$q_1,\allowbreak \ldots,\allowbreak q_k$，即在位形空间中为常数，而且绝对值为$1$。因此，例如前面定义的电荷密度分布就不会变化。（一般说来，我们会猜想：$\psi$中一个在位形空间内为常数、绝对值又为$1$的因子，原则上无法观测。后面更细致的考察将证实这一点。）

第一个微分方程的本征函数构成完全正交系，\origNote[，][所]{20}所以每个$\psi=\psi(q_1,\allowbreak \ldots,\allowbreak q_k)$都能按它们展开。设本征函数为$\psi_1,\allowbreak \psi_2,\allowbreak \ldots$，这里仍都不依赖于$t$；相应的本征值为$\lambda_1,\allowbreak \lambda_2,\allowbreak \ldots$。展开式是\origNote[是][]{21}
\[
 \psi(q_1,\ldots,q_k)=\sum_{n=1}^{\infty}a_n\psi_n(q_1,\ldots,q_k).
\]
如果$\psi$确实依赖于$t$，那么$t$就进入展开系数$a_n$；至于本征函数$\psi_1,\allowbreak \psi_2,\allowbreak \ldots$，此处以及以下各处始终都取为不依赖于$t$。因此，如果给定的$\psi=\psi(q_1,\allowbreak \ldots,\allowbreak q_k)$其实就是$\psi(q_1,\allowbreak \ldots,\allowbreak q_k;t_0)$，那么由
\begin{align*}
 \psi=\psi(q_1,\ldots,q_k;t)&=\sum_{n=1}^{\infty}a_n(t)\psi_n,\\
 H\psi&=\sum_{n=1}^{\infty}a_n(t)H\psi_n
       =\sum_{n=1}^{\infty}\lambda_na_n(t)\psi_n,\\
 \frac{h}{2\pi i}\frac{\partial\psi}{\partial t}
   &=\sum_{n=1}^{\infty}\frac{h}{2\pi i}\dot a_n(t)\psi_n,
\end{align*}
在第二个微分方程中比较系数，便得到
\[
 \frac{h}{2\pi i}\dot a_n(t)=-\lambda_na_n(t),\qquad
 a_n(t)=c_ne^{-\frac{2\pi i}{h}\lambda_nt},
\]
也就是
\begin{align*}
 a_n(t)&=e^{-\frac{2\pi i}{h}\lambda_n(t-t_0)}a_n(t_0)
        =e^{-\frac{2\pi i}{h}\lambda_n(t-t_0)}a_n,\\
 \psi=\psi(q_1,\ldots,q_k;t)
   &=\sum_{n=1}^{\infty}e^{-\frac{2\pi i}{h}\lambda_n(t-t_0)}
     a_n\psi_n(q_1,\ldots,q_k).
\end{align*}
所以，如果$\psi$所表示的状态不是定态，也就是$a_n$并非只有一个不为零，那么当$t$变化时，$\psi$的变化就不再仅仅是乘上一个在位形空间内为常数、绝对值为$1$的因子。前面定义的电荷密度一般也将随之变化；换言之，空间中会发生真实的电振荡。\origNote[。][]{22}

\hypertarget{origpage-10}{}\phantomsection\label{origpage-10}可见，两种理论的概念与实际求解规则，表述得相当不同。尽管如此，它们从一开始就总是给出相同的结果，甚至当两者给出的细节都不同于旧量子理论时，也仍然如此。\origNote[。][正]{23}正如 \hyperref[sec:I.1]{I.1} 中所说，这一引人注目的事实，很快由 Schrödinger 关于两者数学等价的证明得到了解释。\origNote[。][下]{24}下面就来讨论这个等价性证明，同时介绍 Dirac 和 Jordan 的一般“变换理论”；这一理论把上述两种理论都包含在内。

\section{两种理论的等价性：变换理论}\label{sec:I.3}
矩阵理论的基本问题，是确定矩阵$Q_1,\allowbreak \ldots,\allowbreak Q_k,\allowbreak P_1,\allowbreak \ldots,\allowbreak P_k$，使它们首先满足 \hyperref[sec:I.2]{I.2} 的对易关系，见\CJKglue \hyperref[rel:I.2.commutation]{第 \pageref*{rel:I.2.commutation} 页}；\tnote[；][其][原]{原文此处页码误指。这里的页码和链接按对易关系的实际位置设置。}其次，使它们的一个确定函数$H(Q_1,\allowbreak \ldots,\allowbreak Q_k,\allowbreak P_1,\allowbreak \ldots,\allowbreak P_k)$成为对角矩阵。Born 和 Jordan 在最初发表的论文中，就已把这个问题分为如下两步。

第一步，任意找出一组矩阵$\overline Q_1,\allowbreak \ldots,\allowbreak \overline Q_k,\allowbreak \overline P_1,\allowbreak \ldots,\allowbreak \overline P_k$，只要求它们满足对易关系；这很容易做到。\origNote[。][此]{25}此时，一般说来，
\[
 \overline{\mathsf{H}}=H(\overline Q_1,\ldots,\overline Q_k,\overline P_1,\ldots,\overline P_k)
\]
并非对角矩阵。第二步，把所需的解设成
\begin{align*}
 Q_1&=\mathsf{S}^{-1}\overline Q_1\mathsf{S},\ldots,Q_k=\mathsf{S}^{-1}\overline Q_k\mathsf{S},\\
 P_1&=\mathsf{S}^{-1}\overline P_1\mathsf{S},\ldots,P_k=\mathsf{S}^{-1}\overline P_k\mathsf{S},
\end{align*}
其中$\mathsf{S}$本可为任意矩阵，但总须具有逆矩阵$\mathsf{S}^{-1}$，满足$\mathsf{S}^{-1}\mathsf{S}=\mathsf{S}\mathsf{S}^{-1}=1$。由$\overline Q_1,\allowbreak \ldots,\allowbreak \overline Q_k,\allowbreak \overline P_1,\allowbreak \ldots,\allowbreak \overline P_k$满足对易关系，即可推出$Q_1,\allowbreak \ldots,\allowbreak Q_k,\allowbreak P_1,\allowbreak \ldots,\allowbreak P_k$也满足这些关系，而且不论$\mathsf{S}$为何都成立！此外，$\overline{\mathsf{H}}=H(\overline Q_1,\allowbreak \ldots,\allowbreak \overline Q_k,\allowbreak \overline P_1,\allowbreak \ldots,\allowbreak \overline P_k)$变成$\mathsf{H}=H(Q_1,\allowbreak \ldots,\allowbreak Q_k,\allowbreak P_1,\allowbreak \ldots,\allowbreak P_k)$时，有$\mathsf{H}=\mathsf{S}^{-1}\overline{\mathsf{H}}\mathsf{S}$。\origNote[。][所]{26}所以，对$\mathsf{S}$只须提出一个要求：给定$\overline{\mathsf{H}}$之后，$\mathsf{S}^{-1}\overline{\mathsf{H}}\mathsf{S}$应当是对角矩阵。（当然，还应确保$\mathsf{S}^{-1}\overline Q_1\mathsf{S}$等等同原来的$\overline Q_1$等等一样，都是厄米的。不过仔细考察就会发现，对$\mathsf{S}$的这项附加要求，总能事后满足。因此，在这里的初步考察中，我们完全不予考虑。）

这样，任务就是依照$\mathsf{S}^{-1}\overline{\mathsf{H}}\mathsf{S}$的形式，把一个给定的$\overline{\mathsf{H}}$化为对角形。现在来准确说明这意味着什么。

\hypertarget{origpage-11}{}\phantomsection\label{origpage-11}设矩阵$\overline{\mathsf{H}}$的元素为$h_{\mu\nu}$，所求矩阵$\mathsf{S}$的元素为$s_{\mu\nu}$；同样未知的对角矩阵$\mathsf{H}$的对角元为$w_\mu$，即一般元素为$w_\mu\delta_{\mu\nu}$。\origNote[。][]{27}$\mathsf{H}=\mathsf{S}^{-1}\overline{\mathsf{H}}\mathsf{S}$同$\mathsf{S}\mathsf{H}=\overline{\mathsf{H}}\mathsf{S}$是一回事。按已知的矩阵乘法法则算出两边相应的元素，并令它们相等，就得到
\[
 \sum_\nu s_{\mu\nu}w_\nu\delta_{\nu\varrho}
 =\sum_\nu h_{\mu\nu}s_{\nu\varrho},
\]
即
\[
 \sum_\nu h_{\mu\nu}s_{\nu\varrho}=w_\varrho s_{\mu\varrho}.
\]
因此，矩阵$\mathsf{S}$的各列$s_{1\varrho},\allowbreak s_{2\varrho},\allowbreak \ldots$，$\varrho=1,\allowbreak 2,\allowbreak \ldots$，连同$\mathsf{H}$的相应对角元$w_\varrho$，就是所谓本征值问题的解；该问题为
\[
 \sum_\nu h_{\mu\nu}x_\nu=\lambda x_\mu\qquad(\mu=1,2,\ldots).
\]
当然，必须排除平凡解$x_1=x_2=\cdots=0$。实际上，$x_\nu=s_{\nu\varrho}$、$\lambda=w_\varrho$确为一个解；$x_\nu\equiv0$，即对所有$\nu$都有$s_{\nu\varrho}\equiv0$的情形，不可能发生。否则$\mathsf{S}$的第$\varrho$列将恒为零，这同它具有逆矩阵$\mathsf{S}^{-1}$相矛盾！值得注意的是，在本质上，除此之外再没有别的解。

上式的意思是：用矩阵$\overline{\mathsf{H}}$变换向量$x=\{x_1,\allowbreak x_2,\allowbreak \ldots\}$，等于把它乘以数$\lambda$。先用$\mathsf{S}^{-1}$变换$x=\{x_1,\allowbreak x_2,\allowbreak \ldots\}$，得到向量$y=\{y_1,\allowbreak y_2,\allowbreak \ldots\}$；再用$\mathsf{H}$变换$y$，相当于用$\mathsf{H}\mathsf{S}^{-1}=\mathsf{S}^{-1}\overline{\mathsf{H}}\mathsf{S}\cdot\mathsf{S}^{-1}=\mathsf{S}^{-1}\overline{\mathsf{H}}$变换$x$。这又相当于用$\mathsf{S}^{-1}$变换$\lambda x$，所得因而是$\lambda y$。现在，$\overline{\mathsf{H}}y$的分量是\tnote[是][][原]{原文此处写作$\overline{\mathsf{H}}y$，但这一句使用的对角矩阵分量实际属于$\mathsf{H}y$。}$\sum_\nu w_\mu\delta_{\mu\nu}y_\nu=w_\mu y_\mu$，而$\lambda y$的分量是$\lambda y_\mu$。所以，对全部$\mu=1,\allowbreak 2,\allowbreak \ldots$，都必须有$w_\mu y_\mu=\lambda y_\mu$；也就是说，只要$w_\mu\ne\lambda$，就有$y_\mu=0$。把第$\varrho$个分量为$1$、其余分量都为$0$的向量记为$\eta^\varrho$，上面的结论便是：$y$是满足$w_\varrho=\lambda$的那些$\eta^\varrho$的线性组合；特别地，若没有这样的$\varrho$，它就是零。用$\mathsf{S}$作用于$y$便得到$x$，所以$x$是上述各$\eta^\varrho$经$\mathsf{S}$变换所得向量的线性组合。$\eta^\varrho$的第$\nu$个分量为$\delta_{\nu\varrho}$，因此$\mathsf{S}\eta^\varrho$的第$\mu$个分量为$\sum_\nu s_{\mu\nu}\delta_{\nu\varrho}=s_{\mu\varrho}$。若把$\mathsf{S}$的第$\varrho$列$s_{1\varrho},\allowbreak s_{2\varrho},\allowbreak \ldots$看作向量，$x$就是所有满足$w_\varrho=\lambda$的列向量的线性组合；特别地，没有这样的列时，$x$就是零。于是，开头的断言得到证明：只有$w_1,\allowbreak w_2,\allowbreak \ldots$是本征值，而$x_\nu=s_{\nu\varrho}$、$\lambda=w_\varrho$在本质上就是全部的解。

这一点非常重要。它不仅说明，知道$\mathsf{S},\allowbreak \mathsf{H}$便确定了本征值问题的全部解；反过来，一旦完整解出了本征值问题，也可以由此确定$\mathsf{S},\allowbreak \mathsf{H}$。例如，$\mathsf{H}$可这样确定：$w_\mu$就是所有的解$\lambda$；每个$\lambda$在序列$w_1,\allowbreak w_2,\allowbreak \ldots$中出现的次数，等于属于该$\lambda$的线性无关解序列$(x_1,\allowbreak x_2,\allowbreak \ldots)$的个数。\origNote[。][这]{28}\hypertarget{origpage-12}{}\phantomsection\label{origpage-12}这样，除排列次序外，$w_1,\allowbreak w_2,\allowbreak \ldots$已经确定。\origNote[。][]{29}

所以，矩阵理论的核心问题，就是求解本征值方程
\begin{equation}
 \sum_\nu h_{\mu\nu}x_\nu=\lambda x_\mu\qquad(\mu=1,2,\ldots).
 \tag{$E_1$}\label{eq:I.3.E1}
\end{equation}
现在转向波动理论。它的基本方程是“波动方程”
\begin{equation}
 H\varphi(q_1\cdots q_k)=\lambda\varphi(q_1\cdots q_k),
 \tag{$E_2$}\label{eq:I.3.E2}
\end{equation}
其中$H$是前面讨论过的微分算符。任务是求出$\varphi(q_1\cdots q_k)$和$\lambda$的所有解，但须排除$\varphi(q_1\cdots q_k)\equiv0$、$\lambda$任意这一平凡解。这同 \hyperref[eq:I.3.E1]{$E_1$} 的要求相仿：序列$x_1,\allowbreak x_2,\allowbreak \ldots$也可看作“不连续”变量$\nu$的函数$x_\nu$，其中$\nu$取值$1,\allowbreak 2,\allowbreak \ldots$；它对应于“连续”变量$q_1,\allowbreak \ldots,\allowbreak q_k$的函数$\varphi(q_1\cdots q_k)$，而$\lambda$在两种情形中的作用相同。但是，线性变换
\[
 x_\mu\longmapsto\sum_\nu h_{\mu\nu}x_\nu
\]
同另一个变换
\[
 \varphi(q_1\cdots q_k)\longmapsto H\varphi(q_1\cdots q_k)
\]
却不大相像。怎样才能在这里建立类比呢？

我们把下标$\nu$看作变量，并使它同$k$个变量$q_1,\allowbreak \ldots,\allowbreak q_k$相对应，也就是把正整数同$k$维位形空间的一般点相对应；从现在起，这个位形空间记为$\Omega$。因此，不能指望$\sum_\nu$在$\Omega$中仍以求和的形式出现。恰当的对应物反而是积分$\int_\Omega\cdots\int\cdots\,dq_1\cdots dq_k$，或简记为$\int_\Omega\cdots\,dv$，其中$dv$是$\Omega$中的体积元$dq_1\cdots dq_k$。矩阵元$h_{\mu\nu}$依赖于两个同样取正整数值的下标，因而对应于一个函数$h(q_1,\allowbreak \ldots,\allowbreak q_k;q'_1,\allowbreak \ldots,\allowbreak q'_k)$；其中$q_1,\allowbreak \ldots,\allowbreak q_k$和$q'_1,\allowbreak \ldots,\allowbreak q'_k$各自独立地遍历$\Omega$。这样，变换
\[
 x_\mu\longmapsto\sum_\nu h_{\mu\nu}x_\nu
 \quad\text{或}\quad
 x_\nu\longmapsto\sum_{\nu'}h_{\nu\nu'}x_{\nu'}
\]
便变成
\[
 \varphi(q_1\cdots q_k)\longmapsto
 \int_\Omega\cdots\int h(q_1\cdots q_k;q'_1\cdots q'_k)
 \varphi(q'_1\cdots q'_k)\,dq'_1\cdots dq'_k;
\]
本征值问题 \hyperref[eq:I.3.E1]{$E_1$} 也可以写成
\[
 \sum_{\nu'}h_{\nu\nu'}x_{\nu'}=\lambda x_\nu,
\]
它因而变成
\begin{multline}
 \int_\Omega\cdots\int h(q_1\cdots q_k;q'_1\cdots q'_k)
 \varphi(q'_1\cdots q'_k)\,dq'_1\cdots dq'_k\\
 =\lambda\varphi(q_1\cdots q_k).
 \tag{$E_3$}\label{eq:I.3.E3}
\end{multline}\hypertarget{origpage-13}{}\phantomsection\label{origpage-13}\hyperref[eq:I.3.E3]{$E_3$} 这一类本征值问题，在数学中已经得到广泛研究；处理它们时，确实可以同 \hyperref[eq:I.3.E1]{$E_1$} 这一类问题作深入的类比。它们称为“积分方程”。\origNote[。][]{30}

遗憾的是，\hyperref[eq:I.3.E2]{$E_2$} 并不具有这种形式；确切地说，只有对微分算符
\[
 H=H\left(q_1,\ldots,q_k,\frac{h}{2\pi i}\frac{\partial}{\partial q_1},
 \ldots,\frac{h}{2\pi i}\frac{\partial}{\partial q_k}\right)
\]
找到一个函数$h(q_1\cdots q_k;q'_1\cdots q'_k)$，使下式恒成立，即对每一个$\varphi(q_1\cdots q_k)$都成立，才能把它化为这种形式：
\begin{multline}
 H\varphi(q_1\cdots q_k)\\
 =\int_\Omega\cdots\int h(q_1\cdots q_k;q'_1\cdots q'_k)
 \varphi(q'_1\cdots q'_k)\,dq'_1\cdots dq'_k.
 \tag{I}\label{eq:I.3.I}
\end{multline}
如果这样的$h(q_1\cdots q_k;q'_1\cdots q'_k)$存在，就称它为函数运算$H$的“积分核”；$H$本身则称为“积分算符”。

然而，这样的改写一般是不可能的；微分算符$H$从来不是积分算符。即使是最简单的函数算符，把每个$\varphi$映成自身、记为$1$的算符，也不是积分算符。不妨来证实这一点。为简单起见，取$k=1$，要求为
\begin{equation}
 \varphi(q)=\int_{-\infty}^{\infty}h(q,q')\varphi(q')\,dq'.
 \tag{$\Delta_1$}\label{eq:I.3.D1}
\end{equation}
把这里的$\varphi(q)$换成$\varphi(q_0+q)$，令$q=0$，并引入积分变量$q''=q'-q_0$。\tnote[。][于][原]{原文写作$q''=q'-q_0$；紧接着的公式实际需要换元$q''=q'+q_0$。这里保留原文的负号。}于是得到
\[
 \varphi(q_0)=\int_{-\infty}^{\infty}h(0,q''-q_0)\varphi(q'')\,dq''.
\]
再把$q_0,\allowbreak q''$改记为$q,\allowbreak q'$，可知除$h(q,\allowbreak q')$外，$h(0,\allowbreak q'-q)$也解决了我们的问题。因此，不妨假设$h(q,\allowbreak q')$只依赖于$q'-q$。此时要求为
\begin{equation}
 \varphi(q)=\int_{-\infty}^{\infty}h(q'-q)\varphi(q')\,dq'
 \quad\bigl(h(q,q')=h(q'-q)\bigr).
 \tag{$\Delta_2$}\label{eq:I.3.D2}
\end{equation}
再次把$\varphi(q)$换成$\varphi(q_0+q)$，就知道只考察$q=0$已经足够，即
\begin{equation}
 \varphi(0)=\int_{-\infty}^{\infty}h(q)\varphi(q)\,dq.
 \tag{$\Delta_3$}\label{eq:I.3.D3}
\end{equation}
把$\varphi(q)$换成$\varphi(-q)$，可知$h(q)$是解时，$h(-q)$也是解，从而$h_1(q)=\frac{h(q)+h(-q)}2$也是解。因此，可以假设$h(q)$是$q$的偶函数。

这些条件不可能同时满足，这一点很清楚：选择在$q\gtrless0$时满足$\varphi(q)>0$、而$\varphi(0)=0$的函数，由 \hyperref[eq:I.3.D3]{$\Delta_3$} 便推出\tnote[出][，][原]{原文此处结论写的是$\varphi(q)=0$，但这一论证所讨论的应是$h(q)=0$。这里保留原文的符号。}，在$q\gtrless0$时有$\varphi(q)=0$。\origNote[。][可]{31}\hypertarget{origpage-14}{}\phantomsection\label{origpage-14}可是，若选择$\varphi(q)\equiv1$，就得到$\int_{-\infty}^{\infty}h(q)\,dq=1$；而根据刚才的论证，又必然有$\int_{-\infty}^{\infty}h(q)\,dq=0$。

尽管如此，Dirac 仍虚设了这样一个函数的存在：
\begin{equation}
 \delta(q)=0\quad(q\gtrless0),\qquad
 \delta(q)=\delta(-q),\qquad \int_{-\infty}^{\infty}\delta(q)\,dq=1.
 \tag{$\Delta_4$}\label{eq:I.3.D4}
\end{equation}
它将满足 \hyperref[eq:I.3.D3]{$\Delta_3$}：
\begin{align*}
 \int_{-\infty}^{\infty}\delta(q)\varphi(q)\,dq
 &=\varphi(0)\int_{-\infty}^{\infty}\delta(q)\,dq
   +\int_{-\infty}^{\infty}\delta(q)\bigl(\varphi(q)-\varphi(0)\bigr)\,dq\\
 &=\varphi(0)\cdot1+\int_{-\infty}^{\infty}0\,dq=\varphi(0),
\end{align*}
从而也满足 \hyperref[eq:I.3.D1]{$\Delta_1$}、\hyperref[eq:I.3.D2]{$\Delta_2$}。我们应当设想，它除零点外处处为零，但在零点却无穷大，使$\delta(q)$的积分仍然等于$1$。\origNote[。][]{32}

一旦接受这种虚设，只要在$\delta(q)$之外也引入它的导数，就能把种种微分算符表示成积分算符。例如，
\begin{align*}
 \frac{d^n}{dq^n}\varphi(q)
 &=\frac{d^n}{dq^n}\int_{-\infty}^{\infty}\delta(q-q')\varphi(q')\,dq'\\
 &=\int_{-\infty}^{\infty}\frac{\partial^n}{\partial q^n}\delta(q-q')
    \cdot\varphi(q')\,dq'\\
 &=\int_{-\infty}^{\infty}\delta^{(n)}(q-q')\varphi(q')\,dq',\\
 q^n\varphi(q)&=\int_{-\infty}^{\infty}\delta(q-q')q^n\varphi(q')\,dq'.
\end{align*}
也就是说，$\frac{d^n}{dq^n}$及$q^n\cdots$的积分核，分别为$\delta^{(n)}(q-q')$及$\delta(q-q')q^n$。依照同样的办法，可以找出任意复杂的微分算符的积分核。有多个变量$q_1,\allowbreak \ldots,\allowbreak q_k$时，采用$\delta$的乘积即可达到目的。例如，
\begin{align*}
 &\int_\Omega\cdots\int\delta(q_1-q'_1)\delta(q_2-q'_2)\cdots
 \delta(q_k-q'_k)\varphi(q'_1\cdots q'_k)\,dq'_1\cdots dq'_k\\
 &\quad=\int_{-\infty}^{\infty}\biggl[\cdots
 \biggl[\int_{-\infty}^{\infty}
 \biggl[\int_{-\infty}^{\infty}\varphi(q'_1q'_2\cdots q'_k)
 \delta(q_1-q'_1)\,dq'_1\biggr]\\
 &\qquad{}\times\delta(q_2-q'_2)\,dq'_2\biggr]\cdots\biggr]
 \delta(q_k-q'_k)\,dq'_k\\
 &\quad=\int_{-\infty}^{\infty}\left[\cdots
 \left[\int_{-\infty}^{\infty}\varphi(q_1q'_2\cdots q'_k)
 \delta(q_2-q'_2)\,dq'_2\right]\cdots\right]
 \delta(q_k-q'_k)\,dq'_k=\cdots\\
 &\quad=\varphi(q_1q_2\cdots q_k),\\[1ex]
 &\int_\Omega\cdots\int\delta'(q_1-q'_1)\delta(q_2-q'_2)\cdots
 \delta(q_k-q'_k)\varphi(q'_1\cdots q'_k)\,dq'_1\cdots dq'_k\\
 &\quad=\frac{d}{dq_1}\int_\Omega\cdots\int
 \delta(q_1-q'_1)\delta(q_2-q'_2)\cdots\delta(q_k-q'_k)
 \varphi(q'_1\cdots q'_k)\,dq'_1\cdots dq'_k\\
 &\quad=\frac{d}{dq_1}\varphi(q_1\cdots q_k),\qquad\text{等等。}
\end{align*}

\hypertarget{origpage-15}{}\phantomsection\label{origpage-15}这样，实际上几乎所有算符都能强行采用积分表示 \hyperref[eq:I.3.I]{I}。

一旦有了这种表示，问题 \hyperref[eq:I.3.E1]{$E_1$} 与 \hyperref[eq:I.3.E3]{$E_3$} 便完全对应起来；只须作如下替换：
\begin{align*}
 \nu&\mapsto q_1\cdots q_k,& \nu'&\mapsto q'_1\cdots q'_k,\\
 \sum_{\nu'}&\mapsto\int_\Omega\cdots\int\cdots\,dq'_1\cdots dq'_k,& x&\mapsto\varphi.
\end{align*}
向量$x_\nu$对应于函数$\varphi(q_1\cdots q_k)$，同样，矩阵$h_{\nu\nu'}$就应当对应于积分核
\[
h(q_1\cdots q_k;q'_1\cdots q'_k)\text{。}
\]
更方便的办法，是直接把积分核看作矩阵，把$q_1\cdots q_k$称为行指标，把$q'_1\cdots q'_k$称为列指标，分别对应$\nu$与$\nu'$。这样，除通常的矩阵$\{h_{\nu\nu'}\}$之外，就还有另一类矩阵$\{h(q_1\cdots q_k;q'_1\cdots q'_k)\}$，也就是积分核。前者的行与列都是离散的，用数$1,\allowbreak 2,\allowbreak \ldots$标记；后者的行与列，则各由$k$个连续遍历整个$\Omega$的变量标记。

这种类比看上去也许只是形式上的，实际上并非如此。下标$\nu$和$\nu'$也可以看作某个位形空间中的坐标，只要把它们解释为量子数即可。在 Bohr 理论的意义下，量子数就是相空间中可能轨道曲线的编号；这些可能轨道受量子条件的禁限，已经变成离散的。

Dirac 和 Jordan 把这一思路发展成了关于量子过程的统一理论；这里不再继续讨论。诸如$\delta(x),\allowbreak \delta'(x),\allowbreak \ldots$之类的“广义”对象，在其中起决定性作用。它们超出了通常数学方法的范围，而本书希望借助的正是这些通常方法，来描述量子力学。因此，我们转向统一两种理论的另一种方法，即 Schrödinger 的方法。

\section{两种理论的等价性：希尔伯特空间}\label{sec:I.4}
\hyperref[sec:I.3]{I.3} 中所述方法，归根到底，是把指标值组成的“离散”空间$Z=(1,\allowbreak 2,\allowbreak \ldots)$，同力学系统的连续位形空间$\Omega$作类比；若经典力学自由度为$k$，$\Omega$就是$k$维的。这种办法若不对形式体系和数学作某种强行处理，就无法成功，并不令人奇怪：空间$Z$与$\Omega$确实十分不同，任何试图直接把两者联系起来的做法，都必然遭遇巨大的困难。\origNote[。][]{33}

但我们真正需要的，根本不是$Z$与$\Omega$之间的关系，而是两者各自的\hypertarget{origpage-16}{}\phantomsection\label{origpage-16}函数之间的关系，也就是序列$x_1,\allowbreak x_2,\allowbreak \ldots$与波函数$\varphi(q_1\cdots q_k)$之间的关系；前者是$Z$上的函数，后者是$\Omega$上的函数。只有这些函数才进入量子力学的问题。

在 Schrödinger 理论中，积分
\[
 \int_\Omega\cdots\int|\varphi(q_1\cdots q_k)|^2\,dq_1\cdots dq_k
\]
起重要作用：要让$\varphi$能够用于物理论断，就必须使它等于$1$，参看 \hyperref[sec:I.2]{I.2}。另一方面，在矩阵理论中，向量$x_1,\allowbreak x_2,\allowbreak \ldots$具有决定性作用，参看 \hyperref[sec:I.3]{I.3} 的问题 \hyperref[eq:I.3.E1]{$E_1$}。按照 Hilbert 关于这类本征值问题的理论，参看\CJKglue \hyperlink{note-30}{原注 \getrefnumber{orig-note-30}}，始终要求$\sum_\nu|x_\nu|^2$有限。甚至仅仅为了排除平凡解$x_\nu\equiv0$，通常也要作归一化$\sum_\nu|x_\nu|^2=1$。所以，在$Z$和$\Omega$中，都很自然地把允许的函数分别限制为满足
\[
 \sum_\nu|x_\nu|^2
 \quad\text{及}\quad
 \int_\Omega\cdots\int|\varphi(q_1\cdots q_k)|^2\,dq_1\cdots dq_k
\]
有限的函数；只有对这样的函数，才有可能乘上一个常数，使相应的和或积分等于$1$，也就是在通常意义下把解归一化。\origNote[。][这]{34}这两个函数类分别记为$F_Z$与$F_\Omega$。

现在有一条\textbf{定理}：$F_Z$与$F_\Omega$同构，见 Fischer 及 F. Riesz。\origNote[。][其]{35}其确切含义如下。可以在$F_Z$与$F_\Omega$之间建立一一对应：每一个满足$\sum_\nu|x_\nu|^2$有限的序列$x_1,\allowbreak x_2,\allowbreak \ldots$，对应于一个满足
\[
 \int_\Omega\cdots\int|\varphi(q_1\cdots q_k)|^2\,dq_1\cdots dq_k
\]
有限的函数$\varphi(q_1\cdots q_k)$，反之亦然；而且这一对应是线性的，并且保持长度。“线性”是说，若$x_1,\allowbreak x_2,\allowbreak \ldots$与$y_1,\allowbreak y_2,\allowbreak \ldots$分别对应于$\varphi(q_1\cdots q_k)$与$\psi(q_1\cdots q_k)$，那么$ax_1,\allowbreak ax_2,\allowbreak \ldots$及$x_1+y_1,\allowbreak x_2+y_2,\allowbreak \ldots$就分别对应于$a\varphi(q_1\cdots q_k)$及$\varphi(q_1\cdots q_k)+\psi(q_1\cdots q_k)$。对于相互对应的$x_1,\allowbreak x_2,\allowbreak \ldots$和$\varphi(q_1\cdots q_k)$，“保持长度”是说
\[
 \sum_\nu|x_\nu|^2
 =\int_\Omega\cdots\int|\varphi(q_1\cdots q_k)|^2\,dq_1\cdots dq_k.
\]
（之所以称为“保持长度”，是因为通常把$x_1,\allowbreak x_2,\allowbreak \ldots$和$\varphi(q_1\cdots q_k)$看作向量，并把
\[
 \sqrt{\sum_\nu|x_\nu|^2}
 \quad\text{及}\quad
 \sqrt{\int_\Omega\cdots\int|\varphi(q_1\cdots q_k)|^2\,dq_1\cdots dq_k}
\]
看作它们的“长度”。）事实上，还能作更强的断言：若$x_1,\allowbreak x_2,\allowbreak \ldots$与$y_1,\allowbreak y_2,\allowbreak \ldots$分别对应于$\varphi(q_1\cdots q_k)$与$\psi(q_1\cdots q_k)$，则
\[
 \sum_\nu x_\nu\overline{y_\nu}
 =\int_\Omega\cdots\int\varphi(q_1\cdots q_k)
 \overline{\psi(q_1\cdots q_k)}\,dq_1\cdots dq_k,
\]\hypertarget{origpage-17}{}\phantomsection\label{origpage-17}而且两边都绝对收敛。关于最后这一点，还应说明：我们原本也许会希望有
\[
 \sum_\nu x_\nu=\int_\Omega\cdots\int\varphi(q_1\cdots q_k)\,dq_1\cdots dq_k,
\]
或类似的关系，也就是使求和与积分完全对应。但细看就会发现，在量子力学中，求和$\sum_\nu$与积分$\int_\Omega\cdots\int\cdots\,dq_1\cdots dq_k$，始终只分别作用于$x_\nu\overline{y_\nu}$和$\varphi(q_1\cdots q_k)\overline{\psi(q_1\cdots q_k)}$这样的表达式。

如何建立这个对应，下一章还要作非常详细的讨论，所以此处不再深入研究。但必须强调，这种对应的存在具有怎样的意义。$Z$与$\Omega$是很不相同的空间，在它们之间建立关系，必然导致无法解决的数学困难。相反，$F_Z$与$F_\Omega$是同构的，即内在结构完全相同；它们在不同的数学对象中，实现了相同的抽象性质。而且，矩阵理论和波动理论各自真正的分析学基础，是这两个函数类，并不是$Z$与$\Omega$本身！因此，这一同构意味着，两种理论必定总给出相同的数值结果。确切地说，只要所述同构把矩阵
\[
 \overline{\mathsf{H}}=H(\overline Q_1,\ldots,\overline Q_k;\overline P_1,\ldots,\overline P_k)
\]
同算符
\[
 H=H\left(q_1,\ldots,q_k;\frac{h}{2\pi i}\frac{\partial}{\partial q_1},
 \ldots,\frac{h}{2\pi i}\frac{\partial}{\partial q_k}\right)
\]
相对应，情况就是如此。两者分别由矩阵$\overline Q_l,\allowbreak \overline P_l$，$l=1,\allowbreak \ldots,\allowbreak k$，及函数算符$q_l\cdots$、$\frac{h}{2\pi i}\frac{\partial}{\partial q_l}\cdots$，$l=1,\allowbreak \ldots,\allowbreak k$，经过相同的代数运算得到。因此，只须说明$q_l\cdots$对应于矩阵$\overline Q_l$，而$\frac{h}{2\pi i}\frac{\partial}{\partial q_l}\cdots$对应于矩阵$\overline P_l$。对$\overline Q_l,\allowbreak \overline P_l$，$l=1,\allowbreak \ldots,\allowbreak k$，原本只要求满足 \hyperref[sec:I.2]{I.2} 中给出的“对易关系”：\tnote[：][][以]{以下第二行的位置与动量因子，其次序同第 \hyperref[rel:I.2.commutation]{\pageref*{rel:I.2.commutation}} 页的对易关系相反，但原文未改变右端符号。这里保留原书的次序和符号。}
\begin{align*}
 \overline Q_m\overline Q_n-\overline Q_n\overline Q_m&=0,
 &\overline P_m\overline P_n-\overline P_n\overline P_m&=0,\\
 \overline Q_m\overline P_n-\overline P_n\overline Q_m
 &=\begin{cases}0,&m\ne n,\\\dfrac{h}{2\pi i}\,1,&m=n,\end{cases}
 &&(m,n=1,2,\ldots).
\end{align*}
与$q_l\cdots$、$\frac{h}{2\pi i}\frac{\partial}{\partial q_l}\cdots$相对应的矩阵，当然满足这些关系，因为函数算符$q_l\cdots$、$\frac{h}{2\pi i}\frac{\partial}{\partial q_l}\cdots$\hypertarget{origpage-18}{}\phantomsection\label{origpage-18}本身就具有上述性质，\origNote[，][而]{36}而这些性质在同构地转移到$F_Z$时不会丧失。

$F_Z$与$F_\Omega$既然同构，建立在它们之上的量子力学理论又在数学上等价，那么可以预期：只要找出函数类$F_Z$与$F_\Omega$共有的内在性质，并以此为出发点，就能统一地建立量子力学。这样建立的理论，将不依赖于所选形式框架的偶然特点，只呈现量子力学在内容上真正重要的特征。

$F_Z$通常称为“希尔伯特空间”。所以，我们首先要找出希尔伯特空间中不依赖于$F_Z$或$F_\Omega$这些具体形式的内在性质。由这些性质描述的数学对象，称为“抽象希尔伯特空间”。在具体计算时，它要同$F_Z$或$F_\Omega$等同；但在一般讨论中，处理这个抽象对象比直接处理后两者方便。

我们因此准备先描述抽象希尔伯特空间，再完全严格地证明以下两点：
\begin{enumerate}
 \item 抽象希尔伯特空间，简称 H.R.，由将要给出的性质唯一刻画；也就是说，满足这些性质的数学对象在本质上没有区别。
 \item $F_Z$和$F_\Omega$都具备这些性质。
\end{enumerate}
这样，\hyperref[sec:I.4]{I.4} 中仅作定性说明的事情，就得到严格证明。完成这些工作之后，我们再利用所得的数学工具，建立量子力学。


% 中文编写中间材料：chapter-II.tex
% 原稿顺序页25，原印页18，本章从页中开始。
\chapter{抽象希尔伯特空间的一般讨论}\label{chap:II}
\section{希尔伯特空间的刻画}\label{sec:II.1}

现在来完成 \hyperref[sec:I.4]{I.4} 末尾提出的计划。希尔伯特空间是处理量子力学的数学基础；我们要仅用以后量子力学中需要的概念来刻画这个空间。因此，这些概念必须既适用于由数列$x_\nu$（$\nu=1,\allowbreak 2,\allowbreak \ldots$）构成的“离散”函数空间$F_Z$，也同样适用于由波函数$\varphi(q_1\cdots q_k)$构成的“连续”空间$F_\Omega$，其中$q_1,\allowbreak \ldots,\allowbreak q_k$遍历位形空间$\Omega$。前面已经提到，这些概念如下：

\hypertarget{origpage-19}{}\phantomsection\label{origpage-19}$\alpha)$“数乘”，即把一个（复）数$a$与希尔伯特空间的一个元素$f$相乘，得到$af$。在$F_Z$中，这是把$x_\nu$变为$ax_\nu$；在$F_\Omega$中，则把$\varphi(q_1\cdots q_k)$变为$a\varphi(q_1\cdots q_k)$。

$\beta)$希尔伯特空间中两个元素$f,\allowbreak g$的加减：$f\pm g$。在$F_Z$中，由$x_\nu,\allowbreak y_\nu$得到$x_\nu\pm y_\nu$；在$F_\Omega$中，由$\varphi(q_1\cdots q_k),\allowbreak \psi(q_1\cdots q_k)$得到$\varphi(q_1\cdots q_k)\pm\psi(q_1\cdots q_k)$。

$\gamma)$希尔伯特空间中两个元素$f,\allowbreak g$的“内乘”。与$\alpha),\allowbreak \beta)$不同，这种运算的结果不是空间中的元素，而是一个（复）数$(f,\allowbreak g)$。在$F_Z$中，由$x_\nu,\allowbreak y_\nu$得到$\sum_\nu x_\nu\overline{y_\nu}$；在$F_\Omega$中，由$\varphi(q_1\cdots q_k),\allowbreak \psi(q_1\cdots q_k)$得到
\[
\int\!\cdots\!\int_\Omega
\varphi(q_1\cdots q_k)\overline{\psi(q_1\cdots q_k)}\,dq_1\cdots dq_k.
\]
在$F_Z,\allowbreak F_\Omega$中作出这些定义，还须补上相应的收敛性证明。我们将在 \hyperref[sec:II.3]{II.3} 中完成这些证明。

以下一律用$f,\allowbreak g,\allowbreak \ldots,\allowbreak \varphi,\allowbreak \psi,\allowbreak \ldots$表示希尔伯特空间的点，用$a,\allowbreak b,\allowbreak \ldots,\allowbreak x,\allowbreak y,\allowbreak \ldots$表示复数，用$k,\allowbreak l,\allowbreak m,\allowbreak \ldots,\allowbreak \mu,\allowbreak \nu,\allowbreak \ldots$表示正整数。必要时，也把希尔伯特空间记为$\mathfrak{R}_\infty$，这是“无限维欧几里得空间”的简写，与通常表示“$n$维欧几里得空间”的$\mathfrak{R}_n$（$n=1,\allowbreak 2,\allowbreak \ldots$）相仿。

运算$af,\allowbreak f\pm g,\allowbreak (f,\allowbreak g)$有一点值得注意：它们恰好就是向量演算的基本运算。欧几里得几何中的长度、角度计算，以及质点力学中的力、功计算等，都是以这些运算为基础的。对$F_Z$来说，这种类比最为清楚：若把$\mathfrak{R}_\infty$中的无限序列$x_1,\allowbreak x_2,\allowbreak \ldots$换成某个$\mathfrak{R}_n$中的普通点$x_1,\allowbreak \ldots,\allowbreak x_n$，运算$\alpha),\allowbreak \beta),\allowbreak \gamma)$仍有完全相同的形式。特别在$n=3$时，就得到普通空间中的情形。有时，把$x_1,\allowbreak \ldots,\allowbreak x_n$看作向量更方便，例如看作从点$0,\allowbreak \ldots,\allowbreak 0$指向点$x_1,\allowbreak \ldots,\allowbreak x_n$的向量。

因此，我们以向量的基本运算$af,\allowbreak f\pm g,\allowbreak (f,\allowbreak g)$为依据，来刻画抽象希尔伯特空间。以下讨论还将表明，我们在处理$\mathfrak{R}_\infty$的同时，也把所有$\mathfrak{R}_n$包含进来。因此，凡是暂时不必区分$\mathfrak{R}_\infty$与$\mathfrak{R}_n$的地方，都用中性的记号$\mathfrak{R}$表示空间。

首先，要求$\mathfrak{R}$具有向量所特有的性质：\origNote[：][]{37}

\leavevmode\phantomsection\textbf{A.}\label{ax:II.A} $\mathfrak{R}$是一个线性空间。

也就是说，$\mathfrak{R}$中定义了加法$f+g$和“数乘”$af$（$f,\allowbreak g$是$\mathfrak{H}$的元素，$a$是复数；$f+g,\allowbreak af$仍属于$\mathfrak{H}$），并有一个元素$0$。\origNote[。][这]{38}这些运算满足熟知的向量代数法则：

\hypertarget{origpage-20}{}\phantomsection\label{origpage-20}\begin{alignat*}{2}
f+g&=g+f &\quad &\text{（加法交换律）},\\
(f+g)+h&=f+(g+h)&\quad &\text{（加法结合律）},\\
a(f+g)&=af+ag,\quad (a+b)f=af+bf
&\quad &\text{（乘法分配律）},\\
(ab)f&=a(bf)&\quad &\text{（乘法结合律）},\\
0f&=0,\quad 1f=1&\quad &\text{（零与一的作用）}.
\end{alignat*}
这里没有列出的运算法则，也能从这些公设轻易推出。\tnote[。][例][原]{原书在上表中印作$1f=1$。这里所需的恒等式，以及下文计算实际采用的恒等式，都是$1f=f$。公式仍按原书保留。}例如，零在加法中起如下作用：
\[
f+0=1\cdot f+0\cdot f=(1+0)\cdot f=1\cdot f=f.
\]
又如，减法可以唯一定义。令
\[
-f=(-1)\cdot f,\qquad f-g=f+(-g).
\]
则
\begin{align*}
(f-g)+g&=(f+(-g))+g=f+((-g)+g),\\
(f+g)-g&=(f+g)+(-g)=f+(g+(-g)),
\end{align*}
而这两个式子都等于
\begin{align*}
f+((-1)\cdot g+1\cdot g)
&=f+((-1)+1)\cdot g\\
&=f+0\cdot g=f+0=f.
\end{align*}
再如，乘法对于减法的分配律：
\begin{align*}
a\cdot(f-g)&=a\cdot f+a\cdot(-g)=af+a\cdot((-1)\cdot g)\\
&=af+(a\cdot(-1))\cdot g=af+((-1)\cdot a)\cdot g\\
&=af+(-ag)=af-ag,\\
(a-b)\cdot f&=a\cdot f+(-b)\cdot f=af+((-1)b)\cdot f\\
&=af+(-bf)=af-bf.
\end{align*}
这些问题不必再深究下去，因为显然，线性向量演算的全部运算法则在这里仍然成立。

因此，与向量情形一样，可以定义$\mathfrak{H}$的若干元素$f_1,\allowbreak \ldots,\allowbreak f_k$的线性无关性：

\leavevmode\phantomsection\textbf{定义1.}\label{def:II.1} 如果对于复数$a_1,\allowbreak \ldots,\allowbreak a_k$，由$a_1f_1+\cdots+a_kf_k=0$必能推出$a_1=\cdots=a_k=0$，就称$f_1,\allowbreak \ldots,\allowbreak f_k$线性无关。

接着定义线性流形，它对应于向量演算中的线性图形，即经过原点的直线、平面等等。

\leavevmode\phantomsection\textbf{定义2.}\label{def:II.2} $\mathfrak{R}$的子集$\mathfrak{M}$，若包含其中任意$k$（$=1,\allowbreak 2,\allowbreak \ldots$）个元素$f_1,\allowbreak \ldots,\allowbreak f_k$的一切线性组合$a_1f_1+\cdots+a_kf_k$，就称为线性流形。\origNote[。][设]{39}设$\mathfrak{A}$是$\mathfrak{R}$的任意子集，则所有$a_1f_1+\cdots+a_kf_k$构成的集合是一个线性流形；这里$k=1,\allowbreak 2,\allowbreak \ldots$，$a_1,\allowbreak \ldots,\allowbreak a_k$为任意复数，$f_1,\allowbreak \ldots,\allowbreak f_k$为$\mathfrak{A}$的任意元素。这个流形显然包含$\mathfrak{A}$，也显然包含于任何其他包含$\mathfrak{A}$的线性流形之中。\hypertarget{origpage-21}{}\phantomsection\label{origpage-21}我们称它为“由$\mathfrak{A}$张成的线性流形”，记作$\{\mathfrak{A}\}$。

在进一步发展这些概念之前，先表述向量演算的下一条基本原理，即内积的存在。

\leavevmode\phantomsection\textbf{B.}\label{ax:II.B} $\mathfrak{R}$中定义了厄米内积。

也就是说，定义了$(f,\allowbreak g)$（$f,\allowbreak g$属于$\mathfrak{H}$，$(f,\allowbreak g)$为复数），并且它具有下列性质：\origNote[：][]{40}
\begin{alignat*}{2}
(f'+f'',g)&=(f',g)+(f'',g)&\quad &\text{（第一因子的分配律）},\\
(a\cdot f,g)&=a\cdot(f,g)&\quad &\text{（第一因子的结合律）},\\
(f,g)&=\overline{(g,f)}&\quad &\text{（厄米对称性）},\\
(f,f)&\geq0,\quad (f,f)=0\mathwordspace{}\text{仅当}\mathwordspace{}f=0
&\quad &\text{（正定性）}.
\end{alignat*}
由于厄米对称性，第一个因子的两条性质也能推出第二个因子的相应性质：交换$f,\allowbreak g$，再对等式两边取复共轭即可。于是
\[
(f,g'+g'')=(f,g')+(f,g''),\qquad
(f,a\cdot g)=\bar a\cdot(f,g).
\]
这种内积十分重要，因为有了内积就能定义距离。众所周知，在欧几里得空间中，向量$f$的大小定义为$\|f\|=\sqrt{(f,\allowbreak f)}$，\origNote[，][而]{41}而两点$f,\allowbreak g$的距离定义为$\|f-g\|$。我们也采用这种定义。

\leavevmode\phantomsection\textbf{定义3.}\label{def:II.3} $\mathfrak{R}$中元素$f$的“范数”为$\|f\|=\sqrt{(f,\allowbreak f)}$；两个元素$f,\allowbreak g$的“距离”为$\|f-g\|$。\origNote[。][]{42}

这个概念确实具备距离的一切性质，我们马上就会看到。为此先证明：

\leavevmode\phantomsection\textbf{定理1.}\label{thm:II.1} 恒有$|(f,\allowbreak g)|\leq\|f\|\cdot\|g\|$。

\textkaishu{证明。}先约定记号：若$z=u+iv$为复数，$u,\allowbreak v$为实数，则以$\operatorname{Re}z,\allowbreak \operatorname{Im}z$分别表示其实部和虚部，即$\operatorname{Re}z=u,\allowbreak \operatorname{Im}z=v$。有
\begin{align*}
\|f\|^2+\|g\|^2-2\operatorname{Re}(f,g)
&=(f,f)+(g,g)-(f,g)-(g,f)\\
&=(f-g,f-g)\geq0,\\
\operatorname{Re}(f,g)&\leq\tfrac12(\|f\|^2+\|g\|^2).
\end{align*}
把$f,\allowbreak g$分别换成$af,\allowbreak \frac1a g$，其中$a$为正实数。不难看出，左边不变，右边则成为$\frac12(a^2\|f\|^2+\frac1{a^2}\|g\|^2)$。这个量既然不小于$\operatorname{Re}(f,\allowbreak g)$，它的最小值$\|f\|\cdot\|g\|$也必不小于$\operatorname{Re}(f,\allowbreak g)$。当$f,\allowbreak g\ne0$时，取$a=\sqrt{\frac{\|g\|}{\|f\|}}$便达到此最小值；当$f=0$或$g=0$时，则分别让$a\to+\infty$或$a\to+0$来逼近它。\hypertarget{origpage-22}{}\phantomsection\label{origpage-22}因此
\[
\operatorname{Re}(f,g)\leq\|f\|\cdot\|g\|.
\]
再把这里的$f,\allowbreak g$换成$e^{i\alpha}f,\allowbreak g$，其中$\alpha$为实数。右边仍不变，因为
\[
(af,af)=a\bar a(f,f)=|a|^2(f,f),\qquad
\|af\|=|a|\cdot\|f\|,
\]
所以当$|a|=1$时，$\|af\|=\|f\|$。左边则变为
\[
\operatorname{Re}(e^{i\alpha}(f,g))
=\cos\alpha\operatorname{Re}(f,g)-\sin\alpha\operatorname{Im}(f,g).
\]
它的最大值显然为
\[
\sqrt{(\operatorname{Re}(f,g))^2+(\operatorname{Im}(f,g))^2}=|(f,g)|,
\]
于是得到所要证明的结论：
\[
|(f,g)|\leq\|f\|\cdot\|g\|.
\]

\textbf{推论。}等号成立时，$f,\allowbreak g$必只差一个常复数因子。

\textkaishu{证明。}要使$\operatorname{Re}(f,\allowbreak g)\leq\frac12(\|f\|^2+\|g\|^2)$中的等号成立，必须有$(f-g,\allowbreak f-g)=0$，即$f=g$。从这个关系过渡到$|(f,\allowbreak g)|\leq\|f\|\cdot\|g\|$时，只要$f,\allowbreak g$均不为零，我们就把它们分别换成了$e^{i\alpha}af,\allowbreak \frac1a g$（$a,\allowbreak \alpha$为实数，$a>0$）。所以等号成立时，必有$e^{i\alpha}af=\frac1a g$，即$g=a^2e^{i\alpha}f=cf$，其中$c\ne0$。反过来，若$f$或$g$为零，或$g=cf$（$c\ne0$），等号显然成立。

\leavevmode\phantomsection\textbf{定理2.}\label{thm:II.2} 恒有$\|f\|\geq0$，且仅当$f=0$时等号成立。又有$\|a\cdot f\|=|a|\cdot\|f\|$。恒有$\|f+g\|\leq\|f\|+\|g\|$，且只有当$f,\allowbreak g$只差一个非负常实数因子时，等号才成立。

\textkaishu{证明。}前两个断言已在上面证明。第三个断言中的不等式证明如下：
\begin{align*}
(f+g,f+g)&=(f,f)+(g,g)+(f,g)+(g,f)\\
&=\|f\|^2+\|g\|^2+2R(f,g)\\
&\leq\|f\|^2+\|g\|^2+2\|f\|\cdot\|g\|
=(\|f\|+\|g\|)^2,\\
\|f+g\|&\leq\|f\|+\|g\|.
\end{align*}
要使等号成立，必须有$\operatorname{Re}(f,\allowbreak g)=\|f\|\cdot\|g\|$。根据上面推论中的讨论，这就要求$f$或$g$为零，或者$g=a^2f=cf$，其中$c$为正实数。显然，此时等号也确实成立。

从\CJKglue \hyperref[thm:II.2]{定理2} 立即可知，距离$\|f-g\|$具有下列性质：$f,\allowbreak g$之间的距离为零，当且仅当$f=g$；$g,\allowbreak f$的距离与$f,\allowbreak g$的距离相同；$f,\allowbreak h$的距离不大于$f,\allowbreak g$的距离与$g,\allowbreak h$的距离之和，且只有当$g=af+(1-a)h$，其中$a$为实数、$0\leq a\leq1$时，等号才成立；\origNote[；][]{43}$af,\allowbreak ag$的距离等于$f,\allowbreak g$的距离的$|a|$倍。

\hypertarget{origpage-23}{}\phantomsection\label{origpage-23}这些恰好就是几何学和拓扑学中距离所应具备的性质；凭借这些性质，连续、有界、聚点等概念都可以归结到距离概念。我们据此作如下定义：

设$F(f)$是定义在$\mathfrak{R}$的元素$f$上的函数，函数值或者全为$\mathfrak{R}$中的点，或者全为复数。若对每个$\varepsilon>0$，都存在$\delta>0$，使$\|f-f_0\|<\delta$必能推出$\|F(f)-F(f_0)\|<\varepsilon$或$|F(f)-F(f_0)|<\varepsilon$（视$F$的值是$\mathfrak{R}$中的点还是数而定），就称$F$在$\mathfrak{R}$中的$f_0$处连续。若在$\mathfrak{R}$或它的某个给定子集上恒有$\|F(f)\|\leq C$或$|F(f)|\leq C$，其中$C$为适当选定的常数，就称$F$在该集合上有界。多变量情形的定义与此类似。若数列$\|f_1-f\|,\allowbreak \|f_2-f\|,\allowbreak \ldots$收敛于零，就说序列$f_1,\allowbreak f_2,\allowbreak \ldots$收敛于$f$，或以$f$为极限。若一个点是$\mathfrak{A}$中某个序列的极限，就称它为集合$\mathfrak{A}$（$\mathfrak{R}$的子集）的聚点。\origNote[。][特]{44}特别地，若$\mathfrak{A}$的一切聚点仍属于$\mathfrak{A}$，就称$\mathfrak{A}$为闭集；若它的聚点遍及整个$\mathfrak{R}$，就称它处处稠密。

还要指出，$af,\allowbreak f+g,\allowbreak (f,\allowbreak g)$对于各个变量都是连续的。前两项由下式即知：
\[
\|af-af'\|=|a|\cdot\|f-f'\|,
\]
\[
\|(f+g)-(f'+g')\|
=\|(f-f')+(g-g')\|\leq\|f-f'\|+\|g-g'\|.
\]
再设$\|f-f'\|<\varepsilon,\allowbreak \|g-g'\|<\varepsilon$，并令$f'-f=\varphi,\allowbreak g'-g=\psi$，则
\begin{align*}
|(f,g)-(f',g')|
&=|(f,g)-(f+\varphi,g+\psi)|\\
&=|(\varphi,g)+(f,\psi)+(\varphi,\psi)|\\
&\leq|(\varphi,g)|+|(f,\psi)|+|(\varphi,\psi)|\\
&\leq\|\varphi\|\cdot\|g\|+\|f\|\cdot\|\psi\|
+\|\varphi\|\cdot\|\psi\|\\
&\leq\varepsilon(\|f\|+\|g\|+\varepsilon).
\end{align*}
当$\varepsilon\to0$时，右端趋于零，因而能使它小于任意$\delta>0$。

可见，性质 \hyperref[ax:II.A]{A}、\hyperref[ax:II.B]{B} 已经使我们能够说明$\mathfrak{R}$的许多性质，却还不足以把各个$\mathfrak{R}_n$彼此区分，也不足以将它们与$\mathfrak{R}$区分：因为我们还没有谈到维数。众所周知，维数与线性无关向量的最大个数有关。若这样的最大个数$n=0,\allowbreak 1,\allowbreak 2,\allowbreak \ldots$存在，则对这个$n$有：

\leavevmode\phantomsection\textbf{$C^{(n)}$.}\label{ax:II.Cn} 恰有$n$个线性无关的向量。

\hypertarget{origpage-24}{}\phantomsection\label{origpage-24}也就是说，能找出$n$个这样的向量，却不存在$n+1$个。

反之，若不存在这样的最大个数，则有：

\leavevmode\phantomsection\textbf{$C^{(\infty)}$.}\label{ax:II.Cinfty} 可以有任意多个线性无关的向量。

也就是说，对每个$k=1,\allowbreak 2,\allowbreak \ldots$，都能找出$k$个这样的向量。

所以，\hyperref[ax:II.Cn]{C} 实际上并不是一条新公设：只要 \hyperref[ax:II.A]{A}、\hyperref[ax:II.B]{B} 成立，总有某个 \hyperref[ax:II.Cn]{$C^{(n)}$} 或 \hyperref[ax:II.Cinfty]{$C^{(\infty)}$} 成立。选择不同的情形，便得到不同的空间$\mathfrak{R}$。我们将看到，从 \hyperref[ax:II.Cn]{$C^{(n)}$} 可以推出$\mathfrak{R}$具有$n$维（复）欧几里得空间的一切性质。但是，仅凭 \hyperref[ax:II.Cinfty]{$C^{(\infty)}$}，还不能保证$\mathfrak{R}$与希尔伯特空间$\mathfrak{R}_\infty$在本质上相同；为此还需要两条公设 \hyperref[ax:II.D]{D}、\hyperref[ax:II.E]{E}。更确切地说，我们将证明：满足 \hyperref[ax:II.A]{A}、\hyperref[ax:II.B]{B}、\hyperref[ax:II.Cn]{$C^{(n)}$} 的$\mathfrak{R}$，具有$\mathfrak{R}_n$的一切性质，特别也具有马上要提出的 \hyperref[ax:II.D]{D}、\hyperref[ax:II.E]{E}；因此，后两条可以由 \hyperref[ax:II.A]{A}、\hyperref[ax:II.B]{B}、\hyperref[ax:II.Cn]{$C^{(n)}$} 推出。还将证明：满足 \hyperref[ax:II.A]{A}、\hyperref[ax:II.B]{B}、\hyperref[ax:II.Cinfty]{$C^{(\infty)}$}、\hyperref[ax:II.D]{D}、\hyperref[ax:II.E]{E} 的$\mathfrak{R}$，具有$\mathfrak{R}_\infty$的一切性质；而 \hyperref[ax:II.D]{D}、\hyperref[ax:II.E]{E} 在这里不可缺少，即它们不能由 \hyperref[ax:II.A]{A}、\hyperref[ax:II.B]{B}、\hyperref[ax:II.Cinfty]{$C^{(\infty)}$} 推出。现在就来表述 \hyperref[ax:II.D]{D}、\hyperref[ax:II.E]{E}。至于所有$\mathfrak{R}_n$以及$\mathfrak{R}_\infty$也具有这些性质的证明，则留到 \hyperref[sec:II.3]{II.3}。

\leavevmode\phantomsection\textbf{D.}\label{ax:II.D} $\mathfrak{R}$是完备的。\origNote[。][]{45}

也就是说，$\mathfrak{R}$中的序列$f_1,\allowbreak f_2,\allowbreak \ldots$，只要满足Cauchy收敛条件，即对每个$\varepsilon>0$都存在$N=N(\varepsilon)$，使$m,\allowbreak n\geq N$必能推出$\|f_m-f_n\|<\varepsilon$，就一定收敛；换言之，它具有上面所定义的极限$f$。

\leavevmode\phantomsection\textbf{E.}\label{ax:II.E} $\mathfrak{R}$是可分的。\origNote[。][]{45}

也就是说，$\mathfrak{R}$中存在一个处处稠密的序列$f_1,\allowbreak f_2,\allowbreak \ldots$。

按照前面的计划，我们将在 \hyperref[sec:II.2]{II.2} 中以此为基础，发展$\mathfrak{R}$的“几何学”，并看出它与$\mathfrak{R}_n$或$\mathfrak{R}_\infty$的几何学一致。

\section{希尔伯特空间的几何学}\label{sec:II.2}

先作两个定义。第一个定义只涉及几何学中角度概念的一个特例；对本书而言，只需考虑这个特例，即直角，也就是正交性。

\leavevmode\phantomsection\textbf{定义4.}\label{def:II.4} 若$(f,\allowbreak g)=0$，就称$\mathfrak{R}$的两个元素$f,\allowbreak g$正交。若线性流形$\mathfrak{M}$的每个元素都与$\mathfrak{N}$的每个元素正交，就称$\mathfrak{M},\allowbreak \mathfrak{N}$正交。若集合$\mathfrak{O}$中的任意$f,\allowbreak g$都满足
\[
(f,g)=\begin{cases}1,&f=g,\\0,&f\ne g\end{cases}
\]\hypertarget{origpage-25}{}\phantomsection\label{origpage-25}就称$\mathfrak{O}$为规范正交系。也就是说，任意两个不同元素都正交，而每个元素的范数都是$1$。\origNote[。][特]{46}特别地，若不能向$\mathfrak{O}$添入新元素而仍保持规范正交性，就称$\mathfrak{O}$是完全的。\origNote[。][]{47}

还要指出：规范正交系$\mathfrak{O}$完全，显然就是说，不存在范数为$1$且与$\mathfrak{O}$全部元素正交的$f$（参见\CJKglue \hyperlink{note-46}{原注 \getrefnumber{orig-note-46}}）。但是，只要$f\ne0$且与整个$\mathfrak{O}$正交，就能令$f'=\frac1{\|f\|}\cdot f$；由于$\|f\|>0$，$f'$便满足上述全部条件，即$\|f'\|=\frac1{\|f\|}\cdot\|f\|=1$，并与$\mathfrak{O}$正交。因此，$\mathfrak{O}$完全也就是说：每个与整个$\mathfrak{O}$正交的$f$都必须为零。

第二个定义只有在$\mathfrak{R}_\infty$中才起实质作用，因为在$\mathfrak{R}_n$中，每个线性流形本来就属于它所描述的类型（见 \hyperref[sec:II.3]{II.3} 末尾）。因此，这个定义无法用直观的几何图形来说明。

\leavevmode\phantomsection\textbf{定义5.}\label{def:II.5} 同时又是闭集的线性流形，称为闭线性流形。设$\mathfrak{A}$是$\mathfrak{R}$中的任意集合；在它张成的线性流形$\{\mathfrak{A}\}$中添入全部聚点，就得到一个包含$\mathfrak{A}$的闭线性流形。它包含于任何其他包含$\mathfrak{A}$的闭线性流形中。\origNote[。][我]{48}我们称它为$\mathfrak{A}$张成的闭线性流形，记作$[\mathfrak{A}]$。

现在更仔细地分析$\mathfrak{R}$，特别是它的完全规范正交系。某个定理若除 \hyperref[ax:II.A]{A}、\hyperref[ax:II.B]{B} 以外，还假定 \hyperref[ax:II.Cn]{$C^{(n)}$}，或假定 \hyperref[ax:II.Cinfty]{$C^{(\infty)}$}、\hyperref[ax:II.D]{D}、\hyperref[ax:II.E]{E}，我们就在定理编号上分别添上上标$(n)$或$(\infty)$；在两种情形下都成立的定理则不加上标。

\leavevmode\phantomsection\textbf{定理$3^{(n)}$.}\label{thm:II.3n} 每个规范正交系至多有$n$个元素；它完全，当且仅当它有$n$个元素。

\textkaishu{说明。}第一项断言表明，规范正交系的元素个数有最大值。按照定义，达到这个最大值的系就是完全的。因此，在 \hyperref[ax:II.Cn]{$C^{(n)}$} 情形下，完全正交系确实存在；根据本定理，它们有$n$个元素：$\varphi_1,\allowbreak \varphi_2,\allowbreak \ldots,\allowbreak \varphi_n$。

\textkaishu{证明。}每个有限规范正交系都是线性无关的。设这个系为$\varphi_1,\allowbreak \varphi_2,\allowbreak \ldots,\allowbreak \varphi_m$。若
\[
a_1\varphi_1+\cdots+a_m\varphi_m=0
\]
分别与$\varphi_\mu$（$\mu=1,\allowbreak 2,\allowbreak \ldots,\allowbreak m$）作内积，便得$a_\mu=0$。所以，根据 \hyperref[ax:II.Cn]{$C^{(n)}$}，它不能有$n+1$个元素。因而，任何规范正交系都不能含有$n+1$个元素的子集。故它必定是有限的，且元素个数不超过$n$。

\hypertarget{origpage-26}{}\phantomsection\label{origpage-26}于是，有$n$个元素的系不可能再扩充，因而是完全的。但元素个数为$m<n$的系$\varphi_1,\allowbreak \varphi_2,\allowbreak \ldots,\allowbreak \varphi_m$并不完全。因为在线性组合$a_1\varphi_1+\cdots+a_m\varphi_m$中，不可能有$n>m$个线性无关的元素；由 \hyperref[ax:II.Cn]{$C^{(n)}$}，必存在不同于所有这些组合的$f$，也就是
\[
\psi=f-a_1\varphi_1-\cdots-a_m\varphi_m
\]
总不为零。而$(\psi,\allowbreak \varphi_\mu)=0$等价于$a_\mu=(f,\allowbreak \varphi_\mu)$（$\mu=1,\allowbreak 2,\allowbreak \ldots,\allowbreak m$），所以能同时对所有$\mu$成立。因此这个系不完全。

\leavevmode\phantomsection\textbf{定理$3^{(\infty)}$.}\label{thm:II.3infty} 每个规范正交系或者是有限的，或者是一个可数无限序列；完全的规范正交系必定是无限的。

\textkaishu{说明。}因此，一切规范正交系都可以写成序列$\varphi_1,\allowbreak \varphi_2,\allowbreak \ldots$，序列可以在有限项后终止；以下就采用这种写法。须注意，与 \hyperref[ax:II.Cn]{$C^{(n)}$} 情形不同，这里元素的个数只是完全性的必要条件，并不是充分条件。\origNote[。][]{49}

\textkaishu{证明。}设$\mathfrak{O}$是规范正交系，$f,\allowbreak g$是其中两个不同的元素。则
\[
(f-g,f-g)=(f,f)+(g,g)-(f,g)-(g,f)=2,
\qquad \|f-g\|=\sqrt2.
\]
根据 \hyperref[ax:II.E]{E}，取$\mathfrak{R}$中处处稠密的序列$f_1,\allowbreak f_2,\allowbreak \ldots$。对$\mathfrak{O}$的每个$f$，这个序列中都有一个$f_m$满足$\|f-f_m\|<\frac12\sqrt2$。与$f,\allowbreak g$相应的$f_m,\allowbreak f_n$必不相同，因为若$f_m=f_n$，就会有
\begin{align*}
\|f-g\|&=\|(f-f_m)-(g-f_m)\|\\
&\leq\|f-f_m\|+\|g-f_m\|
<\tfrac12\sqrt2+\tfrac12\sqrt2=\sqrt2.
\end{align*}
因此，$\mathfrak{O}$中每个$f$都对应于序列$f_1,\allowbreak f_2,\allowbreak \ldots$的某个$f_m$，而不同的$f$对应于不同的$f_m$。所以$\mathfrak{O}$是有限的，或者可以排成一个序列。

仿照\CJKglue \hyperref[thm:II.3n]{定理$3^{(n)}$} 的证明可知：只要$\mathfrak{R}$中有多于$m$个线性无关的元素，系$\varphi_1,\allowbreak \varphi_2,\allowbreak \ldots,\allowbreak \varphi_m$就不可能完全。根据 \hyperref[ax:II.Cinfty]{$C^{(\infty)}$}，这对一切$m$成立，所以完全的系必须是无限的。

下面的定理中，涉及收敛的部分只适用于 \hyperref[ax:II.Cinfty]{$C^{(\infty)}$} 情形。不过，考虑到它们还包含其他断言，仍以一般形式表述更好。

\leavevmode\phantomsection\textbf{定理4.}\label{thm:II.4} 设$\varphi_1,\allowbreak \varphi_2,\allowbreak \ldots$是规范正交系。则一切级数$\sum_\nu(f,\allowbreak \varphi_\nu)\overline{(g,\allowbreak \varphi_\nu)}$，只要有无限项，就绝对收敛。特别地，当$f=g$时恒有$\sum_\nu|(f,\allowbreak \varphi_\nu)|^2\leq|f|^2$。

\textkaishu{证明。}令$a_\nu=(f,\allowbreak \varphi_\nu)$，$\nu=1,\allowbreak 2,\allowbreak \ldots$，则$f-\sum_{\nu=1}^Na_\nu\varphi_\nu=\psi$与所有$\varphi_\nu$（$\nu=1,\allowbreak 2,\allowbreak \ldots,\allowbreak N$）正交（见\CJKglue \hyperref[thm:II.3n]{定理$3^{(n)}$} 的证明）。\hypertarget{origpage-27}{}\phantomsection\label{origpage-27}由$f=\sum_{\nu=1}^Na_\nu\varphi_\nu+\psi$，得
\begin{align*}
(f,f)&=\sum_{\mu,\nu=1}^Na_\mu\bar a_\nu(\varphi_\mu,\varphi_\nu)
+\sum_{\nu=1}^Na_\nu(\varphi_\nu,\psi)\\
&\quad+\sum_{\nu=1}^N\bar a_\nu(\psi,\varphi_\nu)+(\psi,\psi)\\
&=\sum_{\nu=1}^N|a_\nu|^2+(\psi,\psi)
\geq\sum_{\nu=1}^N|a_\nu|^2.
\end{align*}
即$\sum_{\nu=1}^N|a_\nu|^2\leq\|f\|^2$。若系$\varphi_1,\allowbreak \varphi_2,\allowbreak \ldots$有限，立即便有$\sum_\nu|a_\nu|^2\leq\|f\|^2$；若它无限，则让$N\to\infty$，便知$\sum_\nu|a_\nu|^2$绝对收敛，且其和不超过$\|f\|^2$。第二项断言得证。

又由
\[
|(f,\varphi_\nu)\overline{(g,\varphi_\nu)}|
\leq\tfrac12\bigl\{|(f,\varphi_\nu)|^2+|(g,\varphi_\nu)|^2\bigr\},
\]
结合已得的收敛性结论，也就得到第一项断言中的一般收敛性。

\leavevmode\phantomsection\textbf{定理5.}\label{thm:II.5} 设$\varphi_1,\allowbreak \varphi_2,\allowbreak \ldots$是无限规范正交系。则级数$\sum_{\nu=1}^\infty x_\nu\varphi_\nu$收敛，当且仅当$\sum_{\nu=1}^\infty|x_\nu|^2$收敛。后一数项级数的每一项都是非负实数，所以它或者收敛，或者发散到$+\infty$。

\textkaishu{证明。}这一断言只涉及 \hyperref[ax:II.Cinfty]{$C^{(\infty)}$} 情形，因此可以应用 \hyperref[ax:II.D]{D}，即Cauchy收敛判据。于是，$\sum_{\nu=1}^\infty x_\nu\varphi_\nu$收敛，即部分和序列$\sum_{\nu=1}^Nx_\nu\varphi_\nu$在$N\to\infty$时收敛，当且仅当对每个$\varepsilon>0$都存在$N=N(\varepsilon)$，使
\[
\left\|\sum_{\nu=1}^Lx_\nu\varphi_\nu-
\sum_{\nu=1}^Mx_\nu\varphi_\nu\right\|<\varepsilon
\qquad(L,M\geq N).
\]
取$L>M\geq N$，则
\[
\left\|\sum_{\nu=1}^Lx_\nu\varphi_\nu-
\sum_{\nu=1}^Mx_\nu\varphi_\nu\right\|
=\left\|\sum_{\nu=M+1}^Lx_\nu\varphi_\nu\right\|<\varepsilon,
\]
\begin{align*}
\left\|\sum_{\nu=M+1}^Lx_\nu\varphi_\nu\right\|^2
&=\left(\sum_{\mu=M+1}^Lx_\mu\varphi_\mu,
\sum_{\nu=M+1}^Lx_\nu\varphi_\nu\right)\\
&=\sum_{\mu,\nu=M+1}^Lx_\mu\bar x_\nu(\varphi_\mu,\varphi_\nu)
=\sum_{\nu=M+1}^L|x_\nu|^2\\
&=\sum_{\nu=1}^L|x_\nu|^2-\sum_{\nu=1}^M|x_\nu|^2,
\end{align*}
所以
\[
0\leq\sum_{\nu=1}^L|x_\nu|^2-\sum_{\nu=1}^M|x_\nu|^2<\varepsilon^2.
\]
这恰好就是序列$\sum_{\nu=1}^N|x_\nu|^2$（$N\to\infty$），也就是级数$\sum_{\nu=1}^\infty|x_\nu|^2$的Cauchy收敛条件。

\textbf{推论。}若$f=\sum_\nu x_\nu\varphi_\nu$，则$(f,\allowbreak \varphi_\nu)=x_\nu$。无论正交系有限还是无限，这都成立；无限情形当然须假定级数收敛。

\hypertarget{origpage-28}{}\phantomsection\label{origpage-28}\textkaishu{证明。}当$N\geq\nu$时，内积$(\sum_{\mu=1}^Nx_\mu\varphi_\mu,\allowbreak \varphi_\nu)$等于
\[
\sum_{\mu=1}^Nx_\mu(\varphi_\mu,\varphi_\nu)=x_\nu.
\]
若$\varphi_1,\allowbreak \varphi_2,\allowbreak \ldots$有限，就令$N$等于最大指标；若它无限，就利用内积的连续性，令$N\to\infty$。两种情形都得到$(f,\allowbreak \varphi_\mu)=x_\mu$。

\leavevmode\phantomsection\textbf{定理6.}\label{thm:II.6} 设$\varphi_1,\allowbreak \varphi_2,\allowbreak \ldots$是规范正交系，$f$任意。取$x_\nu=(f,\allowbreak \varphi_\nu)$（$\nu=1,\allowbreak 2,\allowbreak \ldots$），则定义$f'=\sum_\nu x_\nu\varphi_\nu$的级数即使有无限多项，也总收敛。$f-f'$与$\varphi_1,\allowbreak \varphi_2,\allowbreak \ldots$正交。

\textkaishu{证明。}收敛性由\CJKglue \hyperref[thm:II.4]{定理4}、\hyperref[thm:II.5]{5} 得出。根据\CJKglue \hyperref[thm:II.5]{定理5} 的推论，$(f',\allowbreak \varphi_\nu)=x_\nu=(f,\allowbreak \varphi_\nu)$，所以$(f-f',\allowbreak \varphi_\nu)=0$。

有了这些准备，就可以给出规范正交系完全的一般判据；这些判据也适用于 \hyperref[ax:II.Cinfty]{$C^{(\infty)}$} 情形。

\leavevmode\phantomsection\textbf{定理7.}\label{thm:II.7} 设$\varphi_1,\allowbreak \varphi_2,\allowbreak \ldots$是规范正交系。下列每个条件都是它完全的充分必要条件：

$\alpha)$由$\varphi_1,\allowbreak \varphi_2,\allowbreak \ldots$张成的闭线性流形$[\varphi_1,\allowbreak \varphi_2,\allowbreak \ldots]$等于$\mathfrak{R}$。

$\beta)$恒有$f=\sum_\nu x_\nu\varphi_\nu$，其中$x_\nu=(f,\allowbreak \varphi_\nu)$（$\nu=1,\allowbreak 2,\allowbreak \ldots$；收敛性由\CJKglue \hyperref[thm:II.6]{定理6} 保证）。

$\gamma)$恒有
\[
(f,g)=\sum_\nu(f,\varphi_\nu)\overline{(g,\varphi_\nu)}
\]
（由\CJKglue \hyperref[thm:II.4]{定理4} 知级数绝对收敛）。

\textkaishu{证明。}若$\varphi_1,\allowbreak \varphi_2,\allowbreak \ldots$完全，则根据\CJKglue \hyperref[thm:II.6]{定理6}，$f-\sum_\nu x_\nu\varphi_\nu$（$x_\nu=(f,\allowbreak \varphi_\nu)$，$\nu=1,\allowbreak 2,\allowbreak \ldots$）与$\varphi_1,\allowbreak \varphi_2,\allowbreak \ldots$正交，所以必为零，即$\beta)$成立。若$\beta)$成立，那么在正交系有无限多项的情形下，每个$f$都是部分和$\sum_{\nu=1}^Nx_\nu\varphi_\nu$在$N\to\infty$时的极限，因而属于$[\varphi_1,\allowbreak \varphi_2,\allowbreak \ldots]$。所以这个流形等于$\mathfrak{R}$，即$\alpha)$成立。若$\alpha)$成立，则作如下推论：一个与所有$\varphi_1,\allowbreak \varphi_2,\allowbreak \ldots$正交的$f$，也与它们的一切线性组合正交；又由连续性，它还与这些线性组合的全部聚点正交，也就是与整个$[\varphi_1,\allowbreak \varphi_2,\allowbreak \ldots]$正交。因此它与整个$\mathfrak{R}$正交，特别也与自身正交：$(f,\allowbreak f)=0$，即$f=0$。所以这个系完全。

这样便得到如下逻辑关系：
\[
\text{完全性}\longrightarrow\beta)\longrightarrow\alpha)
\longrightarrow\text{完全性},
\]
可见$\alpha),\allowbreak \beta)$都是充分必要条件。

由$\gamma)$还能得出：若$f$与所有$\varphi_1,\allowbreak \varphi_2,\allowbreak \ldots$正交，就令$f=g$，得到$(f,\allowbreak f)=\sum_\nu0\cdot0=0$，所以$f=0$，即这个系完全。\hypertarget{origpage-29}{}\phantomsection\label{origpage-29}另一方面，由与完全性等价的$\beta)$可得
\begin{align*}
(f,g)&=\lim_{N\to\infty}
\left(\sum_{\nu=1}^N(f,\varphi_\nu)\varphi_\nu,
\sum_{\nu=1}^N(g,\varphi_\nu)\varphi_\nu\right)\\
&=\lim_{N\to\infty}\sum_{\mu,\nu=1}^N
(f,\varphi_\mu)\overline{(g,\varphi_\nu)}(\varphi_\mu,\varphi_\nu)\\
&=\lim_{N\to\infty}\sum_{\nu=1}^N
(f,\varphi_\nu)\overline{(g,\varphi_\nu)}\\
&=\sum_{\nu=1}^\infty(f,\varphi_\nu)\overline{(g,\varphi_\nu)}.
\end{align*}
（若系有限，则无须取极限）。这就是$\gamma)$，所以它也是充分必要条件。

\leavevmode\phantomsection\textbf{定理8.}\label{thm:II.8} 对每个序列$f_1,\allowbreak f_2,\allowbreak \ldots$，都存在一个规范正交系$\varphi_1,\allowbreak \varphi_2,\allowbreak \ldots$，与它张成同一个线性流形。这两个序列都可以在有限项后终止。

\textkaishu{证明。}先从$f_1,\allowbreak f_2,\allowbreak \ldots$中取出子序列$g_1,\allowbreak g_2,\allowbreak \ldots$，要求它与原序列张成同一个线性流形，且其中元素线性无关。做法如下：取$g_1$为第一个不等于零的$f_n$；取$g_2$为第一个不同于所有$a_1g_1$的$f_n$；取$g_3$为第一个不同于所有$a_1g_1+a_2g_2$的$f_n$；依此类推。若到某个$p$时，再也找不到不同于所有$a_1g_1+\cdots+a_pg_p$的$f_n$，就到$g_p$为止。如此得到的$g_1,\allowbreak g_2,\allowbreak \ldots$显然符合要求。

现在进行如下构造，这就是E.~Schmidt所谓的“正交化方法”：
\begin{alignat*}{2}
\gamma_1&=g_1,&\quad \varphi_1&=\frac1{\|\gamma_1\|}\gamma_1,\\
\gamma_2&=g_2-(g_2,\varphi_1)\varphi_1,
&\quad \varphi_2&=\frac1{\|\gamma_2\|}\gamma_2,\\
\gamma_3&=g_3-(g_3,\varphi_1)\varphi_1-(g_3,\varphi_2)\varphi_2,
&\quad \varphi_3&=\frac1{\|\gamma_3\|}\gamma_3,\\
&\cdots&\quad &\cdots.
\end{alignat*}
每一步构造$\varphi_p$都能进行，因为分母$\|\gamma_p\|$不为零。否则就有$\gamma_p=0$，从而$g_p$是$\varphi_1,\allowbreak \ldots,\allowbreak \varphi_{p-1}$的线性组合，也就是$g_1,\allowbreak \ldots,\allowbreak g_{p-1}$的线性组合，这与假定矛盾。此外，$g_p$显然是$\varphi_1,\allowbreak \ldots,\allowbreak \varphi_p$的线性组合，而$\varphi_p$是$g_1,\allowbreak \ldots,\allowbreak g_p$的线性组合。因此，这两个序列确定同一个线性流形。

最后，由构造可知$\|\varphi_p\|=1$；当$q<p$时，$(\gamma_p,\allowbreak \varphi_q)=0$，从而$(\varphi_p,\allowbreak \varphi_q)=0$。交换$p,\allowbreak q$可知，只要$p\ne q$，后一个等式便成立。所以$\varphi_1,\allowbreak \varphi_2,\allowbreak \ldots$是规范正交系。

\leavevmode\phantomsection\textbf{定理9.}\label{thm:II.9} 对每个闭线性流形$\mathfrak{M}$，都存在一个规范正交系，它所张成的闭线性流形恰好是$\mathfrak{M}$。

\hypertarget{origpage-30}{}\phantomsection\label{origpage-30}\textkaishu{证明。}在 \hyperref[ax:II.Cn]{$C^{(n)}$} 情形下，这一定理不言自明。因为若$\mathfrak{R}$满足 \hyperref[ax:II.A]{A}、\hyperref[ax:II.B]{B}、\hyperref[ax:II.Cn]{$C^{(n)}$}，则其中每个线性流形$\mathfrak{M}$都满足 \hyperref[ax:II.A]{A}、\hyperref[ax:II.B]{B} 和某个$m\leq n$的 \hyperref[ax:II.Cn]{$C^{(m)}$}。于是，\CJKglue \hyperref[thm:II.3n]{定理$3^{(n)}$} 的说明可应用于$\mathfrak{M}$：存在一个在$\mathfrak{M}$中完全的规范正交系$\varphi_1,\allowbreak \ldots,\allowbreak \varphi_m$。根据\CJKglue \hyperref[thm:II.7]{定理7} 的$\alpha)$，这恰好就是所要证明的结论。可见，此时甚至无须假定闭性，反而还证明了闭性；参见\CJKglue \hyperref[def:II.5]{定义5} 之前的说明。

在 \hyperref[ax:II.Cinfty]{$C^{(\infty)}$} 情形下，先回顾：根据 \hyperref[ax:II.E]{E}，$\mathfrak{R}$可分。我们要证明$\mathfrak{M}$也可分；事实上，$\mathfrak{R}$的任何子集都可分。取$\mathfrak{R}$中处处稠密的序列$f_1,\allowbreak f_2,\allowbreak \ldots$（见 \hyperref[sec:II.1]{II.1} 的 \hyperref[ax:II.E]{E}）。对每个$f_n$及每个$m=1,\allowbreak 2,\allowbreak \ldots$，作球$\mathfrak{K}_{n,\allowbreak m}$，它由一切满足$\|f-f_n\|<1/m$的$f$构成。凡$\mathfrak{K}_{n,\allowbreak m}$含有$\mathfrak{M}$的点，就从中选出一个这样的点$g_{n,\allowbreak m}$。某些$n,\allowbreak m$所对应的$g_{n,\allowbreak m}$可以没有定义，但所有已定义的点能排成$\mathfrak{M}$中的一个序列。\origNote[。][现]{50}现取$\mathfrak{M}$中任意一点$f$，并设$\varepsilon<0$。\tnote[。][可][原]{原书印作$\varepsilon<0$；下述逼近论证需要$\varepsilon>0$。这里保留原书不等式。}可取$m$使$1/m<\varepsilon/2$，再取$f_n$使$\|f_n-f\|<1/m$。于是$\mathfrak{K}_{nm}$含有$\mathfrak{M}$的一点，即$f$，所以$g_{nm}$有定义，且$\|f_n-g_{nm}\|<1/m$；从而$\|f-g_{nm}\|<2/m<\varepsilon$。因此$f$是已定义的$g_{nm}$的聚点，这个序列便符合要求。

现在把这个在$\mathfrak{M}$中处处稠密的序列记作$f_1,\allowbreak f_2,\allowbreak \ldots$。它确定的闭线性流形$[f_1,\allowbreak f_2,\allowbreak \ldots]$包含其全部聚点，因而包含整个$\mathfrak{M}$。但$\mathfrak{M}$本身是包含$f_1,\allowbreak f_2,\allowbreak \ldots$的闭线性流形，所以$[f_1,\allowbreak f_2,\allowbreak \ldots]$也包含于$\mathfrak{M}$，故两者相等。再依\CJKglue \hyperref[thm:II.8]{定理8} 取规范正交系$\varphi_1,\allowbreak \varphi_2,\allowbreak \ldots$，便有$\{\varphi_1,\allowbreak \varphi_2,\allowbreak \ldots\}=\{f_1,\allowbreak f_2,\allowbreak \ldots\}$。两边都添入聚点，就得到$[\varphi_1,\allowbreak \varphi_2,\allowbreak \ldots]=[f_1,\allowbreak f_2,\allowbreak \ldots]=\mathfrak{M}$，这正是所要证明的。

在\CJKglue \hyperref[thm:II.9]{定理9} 中取$\mathfrak{M}=\mathfrak{R}$，根据\CJKglue \hyperref[thm:II.7]{定理7} 的$\alpha)$，便得到完全规范正交系$\varphi_1,\allowbreak \varphi_2,\allowbreak \ldots$。所以完全规范正交系确实存在。有了这一点，就能够证明：视 \hyperref[ax:II.Cn]{$C^{(n)}$} 或 \hyperref[ax:II.Cinfty]{$C^{(\infty)}$} 何者成立，$\mathfrak{R}$相应地是一个$\mathfrak{R}_n$或$\mathfrak{R}_\infty$；从而它的一切性质也都确定了。

只须证明，能把$\mathfrak{R}$一一映射到所有$\{x_1,\allowbreak \ldots,\allowbreak x_n\}$构成的集合，或所有满足$\sum_{\nu=1}^\infty|x_\nu|^2$有限的$\{x_1,\allowbreak x_2,\allowbreak \ldots\}$构成的集合，并使下列条件成立：

1. 若$f\longleftrightarrow\{x_1,\allowbreak x_2,\allowbreak \ldots\}$，则$af\longleftrightarrow\{ax_1,\allowbreak ax_2,\allowbreak \ldots\}$。

2. 若$f\longleftrightarrow\{x_1,\allowbreak x_2,\allowbreak \ldots\}$且$g\longleftrightarrow\{y_1,\allowbreak y_2,\allowbreak \ldots\}$，则$f+g\longleftrightarrow\{x_1+y_1,\allowbreak x_2+y_2,\allowbreak \ldots\}$。

\hypertarget{origpage-31}{}\phantomsection\label{origpage-31}3. 若$f\longleftrightarrow\{x_1,\allowbreak x_2,\allowbreak \ldots\}$且$g\longleftrightarrow\{y_1,\allowbreak y_2,\allowbreak \ldots\}$，则
\[
(f,g)=\sum_{\nu=1}^{n\mathwordspace{}\text{或}\mathwordspace{}\infty}x_\nu\overline{y_\nu}.
\]
（无限情形下，还须证明第3项中的级数绝对收敛。）现在就来给出这个对应$f\longleftrightarrow\{x_1,\allowbreak x_2,\allowbreak \ldots\}$。

取完全规范正交系$\varphi_1,\allowbreak \varphi_2,\allowbreak \ldots$。在 \hyperref[ax:II.Cn]{$C^{(n)}$} 情形下，它到$\varphi_n$为止；在 \hyperref[ax:II.Cinfty]{$C^{(\infty)}$} 情形下，它是无限的（\hyperref[thm:II.3n]{定理$3^{(n)}$}、\hyperref[thm:II.3infty]{$3^{(\infty)}$}）。令
\[
f=\sum_{\nu=1}^{n\mathwordspace{}\text{或}\mathwordspace{}\infty}x_\nu\varphi_\nu.
\]
在无限情形下，由于$\sum_{\nu=1}^\infty|x_\nu|^2$有限，根据\CJKglue \hyperref[thm:II.5]{定理5}，这个级数也收敛；因此，$\mathfrak{R}_n$或$\mathfrak{R}_\infty$中的每个坐标序列都确实对应一个$f$。另一方面，由\CJKglue \hyperref[thm:II.7]{定理7} 的$\beta)$，以及$\sum_{\nu=1}^{n\mathwordspace{}\text{或}\mathwordspace{}\infty}|(f,\allowbreak \varphi_\nu)|^2$有限（\hyperref[thm:II.4]{定理4}），可知$\mathfrak{R}$中的每个$f$也都能由此式得到：只须令$x_\nu=(f,\allowbreak \varphi_\nu)$即可。显然，每个$\{x_1,\allowbreak x_2,\allowbreak \ldots\}$只对应于一个$f$；反向的唯一性则由\CJKglue \hyperref[thm:II.5]{定理5} 的推论得出。第1、2项显然满足，第3项由\CJKglue \hyperref[thm:II.7]{定理7} 的$\gamma)$得出。

\section[条件A—E的补充讨论]{条件A—E的补充讨论\texorpdfstring{\origNote[论][]{51}}{}}\label{sec:II.3}
还须验证 \hyperref[sec:I.4]{I.4} 末尾的第2项断言，即$F_Z,\allowbreak F_\Omega$确实满足条件 \hyperref[ax:II.A]{A}—\hyperref[ax:II.E]{E}。只考察$F_\Omega$就够了，因为在 \hyperref[sec:II.2]{II.2} 中已经证明，满足 \hyperref[ax:II.A]{A}—\hyperref[ax:II.E]{E} 的$\mathfrak{R}$，在一切性质上都必须与$\mathfrak{R}_\infty$，即$F_Z$一致；所以$F_Z$也必满足 \hyperref[ax:II.A]{A}—\hyperref[ax:II.E]{E}。此外，还要证明 \hyperref[sec:II.2]{II.2} 中提到的两点：条件 \hyperref[ax:II.D]{D}、\hyperref[ax:II.E]{E} 独立于 \hyperref[ax:II.A]{A}—\hyperref[ax:II.Cinfty]{$C^{(\infty)}$}；而它们又可由 \hyperref[ax:II.A]{A}—\hyperref[ax:II.Cn]{$C^{(n)}$} 推出，也就是说，在$\mathfrak{R}_n$中成立。这三个纯数学问题，就是本节的内容。

先验证$F_\Omega$中的 \hyperref[ax:II.A]{A}—\hyperref[ax:II.E]{E}。这里必须借助Lebesgue的积分概念；关于这一概念的基础，只能请读者参看有关专著。\origNote[。][不]{52}不过，Lebesgue积分只在这里使用；阅读后面的章节，无须掌握它。

在 \hyperref[sec:I.4]{I.4} 中，我们把$\Omega$定义为$q_1,\allowbreak \ldots,\allowbreak q_k$的$k$维空间，把$F_\Omega$定义为使
\[
\underbrace{\int\cdots\int}_{\Omega}|\varphi(q_1\cdots q_k)|^2\,dq_1\cdots dq_k
\]
有限的一切函数$f(q_1\cdots q_k)$构成的总体；这里每个$q_1,\allowbreak \ldots,\allowbreak q_k$都从$-\infty$变到$+\infty$。即使限制$q_1,\allowbreak \ldots,\allowbreak q_k$的变化范围，例如使$\Omega$成为半空间、立方体内部、球内部，或这些图形的外部等等，我们的一切推导仍然成立，而且大多数证明可以逐字沿用；甚至也可以把$\Omega$取为曲面，例如球面。

\hypertarget{origpage-32}{}\phantomsection\label{origpage-32}为了不陷入无关紧要的复杂情形，我们只讨论刚才所说的最简单情形。其他情形，读者可以参照这里的典型证明自行讨论。现在依次考察 \hyperref[ax:II.A]{A}—\hyperref[ax:II.E]{E}。

\textbf{关于 \hyperref[ax:II.A]{A}。}须证明：$f,\allowbreak g$属于$F_\Omega$时，$af,\allowbreak f\pm g$也属于$F_\Omega$。即由
\[
\int_\Omega|f|^2,\qquad\int_\Omega|g|^2
\]
有限，须能推出$\int_\Omega|af|^2=|a|^2\int_\Omega|f|^2$及$\int_\Omega|f\pm g|^2$有限。这里把
\[
\underbrace{\int\cdots\int}_{\Omega}|f(q_1\cdots q_k)|^2\,dq_1\cdots dq_k,
\quad
\underbrace{\int\cdots\int}_{\Omega}|g(q_1\cdots q_k)|^2\,dq_1\cdots dq_k
\]
简写成上述形式，不会引起误解。第一点显然成立。由于$|f\pm g|^2=|f|^2+|g|^2\pm2\operatorname{Re}(f\overline g)$，\origNote[，][只]{53}只要能保证$\int_\Omega|f\overline g|=\int_\Omega|f||g|$有限，第二点也就得证。而$|f||g|\leq\frac12(|f|^2+|g|^2)$，所以这可由假定推出。

\textbf{关于 \hyperref[ax:II.B]{B}。}把$(f,\allowbreak g)$定义为$\int_\Omega f\overline g$；刚才已知，这个积分绝对收敛。\hyperref[ax:II.B]{B} 中要求的一切性质都很明显，只有最后一条还须说明，即由$(f,\allowbreak f)=0$可推出$f\equiv0$。$(f,\allowbreak f)=0$就是说$\int_\Omega|f|^2=0$，因此，所有满足$|f|^2>0$，即$f(q_1\cdots q_k)\ne0$的点构成的集合，其Lebesgue测度必须为零。若两个函数$f,\allowbreak g$只在$q_1\cdots q_k$的一个Lebesgue零测集上取值不同，即$f(q_1\cdots q_k)\ne g(q_1\cdots q_k)$，我们就不把它们看作有本质区别。\origNote[。][作]{54}作此约定后，这里便可断定$f\equiv0$。

\textbf{关于 \hyperref[ax:II.Cn]{C}。}设$O_1,\allowbreak \ldots,\allowbreak O_n$是$\Omega$中的$n$个区域，两两没有公共点，而且每个区域的Lebesgue测度都大于零而又有限。令$f_l(q_1\cdots q_k)$在$O_l$内等于$1$，其他地方等于$0$。由于$\int_\Omega|f_l|^2$等于$O_l$的测度，$f_l$属于$F_\Omega$（$l=1,\allowbreak \ldots,\allowbreak k$）。这些$f_1,\allowbreak \ldots,\allowbreak f_n$线性无关：若$a_1f_1+\cdots+a_nf_n\equiv0$，左边函数除一个Lebesgue零测集外处处为零，所以每个$O_l$内都有零点；但它在$O_l$内恒等于$a_l$，故$a_l=0$，$l=1,\allowbreak \ldots,\allowbreak n$。对一切$n$都能作此构造，所以 \hyperref[ax:II.Cinfty]{$C^{(\infty)}$} 成立。

\textbf{关于 \hyperref[ax:II.D]{D}。}设序列$f_1,\allowbreak f_2,\allowbreak \ldots$满足Cauchy收敛条件：对每个$\varepsilon>0$，都存在$N=N(\varepsilon)$，使$m,\allowbreak n\geq N$时有$\int_\Omega|f_m-f_n|^2<\varepsilon$。取
\[
n_1=N\left(\frac18\right);\qquad
n_2\geq n_1,N\left(\frac1{8^2}\right);\qquad
n_3\geq n_1,n_2,N\left(\frac1{8^3}\right);\quad\ldots.
\]
于是$n_1\leq n_2\leq\cdots$，并且$n_\nu,\allowbreak n_{\nu+1}\geq N(1/8^\nu)$，所以
\[
\int_\Omega|f_{n_{\nu+1}}-f_{n_\nu}|^2<\frac1{8^\nu}.
\]

\hypertarget{origpage-33}{}\phantomsection\label{origpage-33}考察所有满足$|f_{n_{\nu+1}}-f_{n_\nu}|>1/2^\nu$的点构成的集合$P^{(\nu)}$。若它的Lebesgue测度为$\mu^{(\nu)}$，则
\[
\int_\Omega|f_{n_{\nu+1}}-f_{n_\nu}|^2
\geq\mu^{(\nu)}\left(\frac1{2^\nu}\right)^2
=\frac{\mu^{(\nu)}}{4^\nu},\qquad
\frac{\mu^{(\nu)}}{4^\nu}<\frac1{8^\nu},\qquad
\mu^{(\nu)}<\frac1{2^\nu}.
\]
再考察$P^{(\nu)},\allowbreak P^{(\nu+1)},\allowbreak P^{(\nu+2)},\allowbreak \ldots$的并集$Q^{(\nu)}$。它的Lebesgue测度
\[
\leq\mu^{(\nu)}+\mu^{(\nu+1)}+\mu^{(\nu+2)}+\cdots
<\frac1{2^\nu}+\frac1{2^{\nu+1}}+\frac1{2^{\nu+2}}+\cdots
=\frac1{2^{\nu-1}}.
\]
在$Q^{(\nu)}$之外有
\[
|f_{n_{\nu+1}}-f_{n_\nu}|<\frac1{2^\nu},\quad
|f_{n_{\nu+2}}-f_{n_{\nu+1}}|<\frac1{2^{\nu+1}},\quad
|f_{n_{\nu+3}}-f_{n_{\nu+2}}|<\frac1{2^{\nu+2}},\quad\ldots,
\]
所以一般地，对$\nu\leq\nu'\leq\nu''$有
\begin{align*}
|f_{n_{\nu''}}-f_{n_{\nu'}}|&<|f_{n_{\nu'+1}}-f_{n_{\nu'}}|+|f_{n_{\nu'+2}}-f_{n_{\nu'+1}}|+\cdots+|f_{n_{\nu''}}-f_{n_{\nu''-1}}|\\
&<\frac1{2^{\nu'}}+\frac1{2^{\nu'+1}}+\cdots+\frac1{2^{\nu''-1}}<\frac1{2^{\nu'-1}}.
\end{align*}
当$\nu'\to\infty$时，这个量趋于零，且这一趋近不依赖于$\nu''$。也就是说，只要$q_1\cdots q_k$不属于$Q^{(\nu)}$，序列$f_{n_1},\allowbreak f_{n_2},\allowbreak \ldots$就满足Cauchy收敛条件。对于固定的$q_1\cdots q_k$，它是一个数列，所以确实收敛。反过来说，若$f_{n_1},\allowbreak f_{n_2},\allowbreak \ldots$在某个$q_1\cdots q_k$处不收敛，该点就必属于$Q^{(\nu)}$。把所有不收敛的点$q_1\cdots q_k$构成的集合记为$Q$，则$Q$包含于$Q^{(\nu)}$，所以它的测度小于$Q^{(\nu)}$的测度，从而小于$1/2^{\nu-1}$。$Q$的定义与$\nu$无关，而这个结论对一切$\nu$成立，因此$Q$的Lebesgue测度为零。于是，例如把所有$f_n$在$Q$上的值改成零，也没有影响（见\CJKglue \hyperlink{note-54}{原注 \getrefnumber{orig-note-54}}）。这样，$f_{n_1},\allowbreak f_{n_2},\allowbreak \ldots$在$Q$上也收敛，也就处处收敛了。

这样就从$f_1,\allowbreak f_2,\allowbreak \ldots$中找到了一个在每一点$q_1\cdots q_k$处都收敛的子序列$f_{n_1},\allowbreak f_{n_2},\allowbreak \ldots$；原序列$f_1,\allowbreak f_2,\allowbreak \ldots$本身未必如此。把子序列的极限记为$f=f(q_1\cdots q_k)$。还须证明：1. $f$属于$F_\Omega$，即$\int_\Omega|f|^2$有限；2. $f$不仅是$f_{n_1},\allowbreak f_{n_2},\allowbreak \ldots$的逐点极限，也是它在希尔伯特空间中按范数收敛的极限，即$\|f-f_{n_\nu}\|\to0$，或$\int_\Omega|f-f_{n_\nu}|^2\to0$；3. 在后一意义下，$f$甚至是整个序列$f_1,\allowbreak f_2,\allowbreak \ldots$的极限，即$\|f-f_n\|\to0$，或$\int_\Omega|f-f_n|^2\to0$。

设$\varepsilon>0$，取$\nu_0$使$n_{\nu_0}\geq N(\varepsilon)$，例如令$1/8^{\nu_0}\leq\varepsilon$；再设$\nu\geq\nu_0,\allowbreak n\geq N(\varepsilon)$。于是$\int_\Omega|f_{n_\nu}-f_n|^2<\varepsilon$。当$\nu\to\infty$时，被积函数趋于$|f-f_n|^2$，根据Lebesgue积分的一个收敛定理（参见\CJKglue \hyperlink{note-53}{原注 \getrefnumber{orig-note-53}} 所引文献），\tnote[，][便][原]{原书此处引原注53；关于Lebesgue积分的文献实际列在\CJKglue \hyperlink{note-52}{原注 \getrefnumber{orig-note-52}}，原注53说明的是代数恒等式。此处保留原书引注，并给出所需文献的链接。}便得$\int_\Omega|f-f_n|^2\leq\varepsilon$。

\hypertarget{origpage-34}{}\phantomsection\label{origpage-34}于是，首先$\int_\Omega|f-f_1|^2$有限，即$f-f_1$属于$F_\Omega$；又因$f_1$属于$F_\Omega$，$f$也属于$F_\Omega$。第1点得证。其次，由上面的不等式，$n\to\infty$时有$\int_\Omega|f-f_n|^2\to0$，第2、3点也得证。

\textbf{关于 \hyperref[ax:II.E]{E}。}须找出一个在$F_\Omega$中处处稠密的函数序列$f_1,\allowbreak f_2,\allowbreak \ldots$。

设$\Omega_1,\allowbreak \Omega_2,\allowbreak \ldots$是$\Omega$中的一列区域，每个区域的测度都有限，而且它们合起来充满整个$\Omega$；例如可令$\Omega_N$为以原点为中心、半径为$N$的球。设$f=f(q_1\cdots q_k)$为$F_\Omega$的任意元素，对每个$N=1,\allowbreak 2,\allowbreak \ldots$，定义$f_N=f_N(q_1\cdots q_k)$如下：
\[
f_N(q_1\cdots q_k)=
\begin{cases}
f(q_1\cdots q_k),&q_1\cdots q_k\in\Omega_N\mathwordspace{}\text{与}\mathwordspace{}|f(q_1\cdots q_k)|\leq N,\\
0,&\text{其他情形}.
\end{cases}
\]
当$N\to\infty$时，$f_N(q_1\cdots q_k)\to f(q_1\cdots q_k)$，甚至从某个$N$起两者就相等，所以$|f-f_N|^2\to0$。此外，$f-f_N$不是$0$就是$f$，故$|f-f_N|^2\leq|f|^2$。所以这些积分$\int_\Omega|f-f_N|^2$均受同一个有限积分$\int_\Omega|f|^2$控制；而被积函数趋于零，根据刚才引用的收敛定理，积分也趋于零：$\int_\Omega|f-f_N|^2\to0$，即$\|f-f_N\|\to0$。

把满足以下条件的一切函数$g=g(q_1\cdots q_k)$构成的类记为$G$：$g\ne0$的点集测度有限，且在整个空间恒有$|g|\leq C$，其中$C$是任意选定的常数。上面的所有$f_N$都属于$G$，所以$G$在$F_\Omega$中处处稠密。

设$g$属于$G$，$\varepsilon>0$。令$M$为$g\ne0$的点集的测度，$C$为$|g|$的一个上界。取一串有理数$-C<\rho_1<\rho_2<\cdots<\rho_t<C$，使$\rho_1<-C+\varepsilon,\allowbreak \rho_2<\rho_1+\varepsilon,\allowbreak \ldots,\allowbreak \rho_t<\rho_{t-1}+\varepsilon,\allowbreak C<\rho_t+\varepsilon$；这很容易做到。把$\operatorname{Re}g(q_1\cdots q_k)$的每个值改为最接近的$\rho_s$（$s=1,\allowbreak 2,\allowbreak \ldots,\allowbreak t$），唯独$0$保持不变，便得到新函数$h_1(q_1\cdots q_k)$，它与$\operatorname{Re}g$在各处的差的绝对值都小于$\varepsilon$。对$\operatorname{Im}g$也同样构造$h_2(q_1\cdots q_k)$。令$h=h_1+ih_2$，则
\begin{align*}
\int_\Omega|g-h|^2
&=\int_\Omega|\operatorname{Re}g-h_1|^2+\int_\Omega|\operatorname{Im}g-h_2|^2\\
&\leq M\varepsilon^2+M\varepsilon^2=2M\varepsilon^2,\\
\|g-h\|&\leq\sqrt{2M}\,\varepsilon.
\end{align*}
给定$\delta>0$后，取$\varepsilon<\delta/\sqrt{2M}$，便有$\|g-h\|<\delta$。

\hypertarget{origpage-35}{}\phantomsection\label{origpage-35}把满足下列条件的一切函数$h=h(q_1\cdots q_k)$构成的类记为$H$：函数只取有限多个不同的值，每个值都形如$\rho+i\sigma$，其中$\rho,\allowbreak \sigma$为有理数；每个非零值只在测度有限的集合上取得。刚才构造的$h$属于$H$，所以$H$在$G$中处处稠密，因而在$F_\Omega$中也处处稠密。

对Lebesgue测度有限的集合$\Pi$，定义函数$f_\Pi=f_\Pi(q_1\cdots q_k)$如下：
\[
f_\Pi(q_1\cdots q_k)=\begin{cases}1,&\text{在}\mathwordspace{}\Pi,\\0,&\text{其他情形}.\end{cases}
\]
显然，$H$类恰由所有
\[
\sum_{s=1}^t(\rho_s+i\sigma_s)f_{\Pi_s}
\qquad(t=1,2,\ldots;\ \rho_s,\sigma_s\mathwordspace{}\text{为有理数})
\]
组成。现在要找一列$\Pi$型集合$\Pi^{(1)},\allowbreak \Pi^{(2)},\allowbreak \ldots$，使对每个$\Pi$型集合及每个$\varepsilon>0$，都有一个$\Pi^{(n)}$，满足下列条件：属于$\Pi$但不属于$\Pi^{(n)}$，或属于$\Pi^{(n)}$但不属于$\Pi$的全部点构成的集合，其测度小于$\varepsilon$。这个集合称为$\Pi,\allowbreak \Pi^{(n)}$的对称差。若有这样一个序列，则函数
\[
\sum_{s=1}^t(\rho_s+i\sigma_s)f_{\Pi^{(n_s)}}
\qquad(t=1,2,\ldots;\ \rho_s,\sigma_s\mathwordspace{}\text{为有理数};\ n_s=1,2,\ldots)
\]
在$H$中处处稠密。因为，对前面的每个$\Pi_s$，按上述方式选择$\Pi^{(n_s)}$，便得
\begin{align*}
&\int_\Omega\left|\sum_{s=1}^t(\rho_s+\sigma_s)f_{\Pi_s}-\sum_{s=1}^t(\rho_s+i\sigma_s)f_{\Pi^{(n_s)}}\right|^2\\
&\quad\leq\sum_{s=1}^t\int_\Omega\left|(\rho_s+i\sigma_s)f_{\Pi_s}-(\rho_s+i\sigma_s)f_{\Pi^{(n_s)}}\right|^2\\
&\quad=\sum_{s=1}^t(\rho_s^2+\sigma_s^2)\int_\Omega|f_{\Pi_s}-f_{\Pi^{(n_s)}}|^2\\
&\quad=\sum_{s=1}^t(\rho_s^2+\sigma_s^2)\,\text{下列两集合的对称差的测度：}\Pi_s,\Pi^{(n_s)}\\
&\quad<\sum_{s=1}^t(\rho_s^2+\sigma_s^2)\,\varepsilon.
\end{align*}
给定$\delta>0$，只要取$\varepsilon=\delta^2/\sum_{s=1}^t(\rho_s^2+\sigma_s^2)$，就有
\[
\left\|\sum_{s=1}^t(\rho_s+i\sigma_s)f_{\Pi_s}-\sum_{s=1}^t(\rho_s+i\sigma_s)f_{\Pi^{(n_s)}}\right\|<\delta.
\]
而这些$\sum_{s=1}^t(\rho_s+\sigma_s)f_{\Pi^{(n_s)}}$适当排列后可以构成一个序列。\tnote[。][事][原]{原书在这里，以及上一个积分估计的第一项中，都把$\rho_s+i\sigma_s$中的$i$漏去了。两处均按原书保留；前后构造使用的是复系数$\rho_s+i\sigma_s$。}事实上，设所有$\rho_1,\allowbreak \sigma_1,\allowbreak \ldots,\allowbreak \rho_t,\allowbreak \sigma_t$的公分母为$\tau$，相应分子为$\rho'_1,\allowbreak \sigma'_1,\allowbreak \ldots,\allowbreak \rho'_t,\allowbreak \sigma'_t$，则这些函数为
\[
\frac1\tau\sum_{s=1}^t(\rho'_s+i\sigma'_s)f_{\Pi^{(n_s)}},
\]

\hypertarget{origpage-36}{}\phantomsection\label{origpage-36}其中$t,\allowbreak \tau=1,\allowbreak 2,\allowbreak \ldots$；$\rho'_s,\allowbreak \sigma'_s=0,\allowbreak \pm1,\allowbreak \pm2,\allowbreak \ldots$；$n_s=1,\allowbreak 2,\allowbreak \ldots$，$s=1,\allowbreak \ldots,\allowbreak t$。把这些函数排成序列，等于把它们的数字参数组$t,\allowbreak \tau,\allowbreak \rho'_1,\allowbreak \sigma'_1,\allowbreak \ldots,\allowbreak \rho'_t,\allowbreak \sigma'_t,\allowbreak n_1,\allowbreak \ldots,\allowbreak n_t$排成序列。按下列和的递增次序排列这些参数组：
\[
t+\tau+|\rho'_1|+|\sigma'_1|+\cdots+|\rho'_t|+|\sigma'_t|+n_1+\cdots+n_t.
\]
对于这个和的每个给定值，相应的参数组只有有限多个。把每个有限集合中的参数组也按任意顺序排列，就确实得到一个单序列。

为了找出所需的集合序列$\Pi^{(1)},\allowbreak \Pi^{(2)},\allowbreak \ldots$，利用如下事实：对每个Lebesgue测度为有限值$M$的集合$\Pi$，以及每个$\delta>0$，都存在包含$\Pi$的开点集$\Pi'$，其测度比$\Pi$的测度大不到$\delta$（参见\CJKglue \hyperlink{note-52}{原注 \getrefnumber{orig-note-52}} 所引文献，以及定义开点集概念的\CJKglue \hyperlink{note-45}{原注 \getrefnumber{orig-note-45}}）。而对开集$\Pi'$和$\varepsilon>0$，显然又存在由有限多个立方体组成的集合$\Pi''$，它包含于$\Pi'$，且其测度比$\Pi''$的测度小不到$\delta$。\tnote[。][这][原]{原书这里把$\Pi''$的测度与自身比较，论证所需的比较对象应为$\Pi'$；此外，前面写$\varepsilon$，估计中却使用$\delta$。这些笔误均按原书保留。}这些立方体的边长和中心坐标都可以取为有理数。不难看出，上面定义的$\Pi$与$\Pi''$的对称差，其测度小于$\delta+\delta=2\delta$；令$\delta=\varepsilon/2$，就小于$\varepsilon$。因此，只要能把上述立方体集合排成一个序列，就达到了目的。

这些集合由所含立方体的个数$n=1,\allowbreak 2,\allowbreak \ldots$，以及各立方体的边长$\varkappa^{(\nu)}$和中心坐标$\xi_1^{(\nu)},\allowbreak \ldots,\allowbreak \xi_k^{(\nu)}$（$\nu=1,\allowbreak \ldots,\allowbreak n$）确定。$\varkappa^{(\nu)},\allowbreak \xi_1^{(\nu)},\allowbreak \ldots,\allowbreak \xi_k^{(\nu)}$都是有理数；设公分母为$\eta=1,\allowbreak 2,\allowbreak \ldots$，相应的分子为$\varkappa'^{(\nu)}=1,\allowbreak 2,\allowbreak \ldots$，$\xi_1'^{(\nu)},\allowbreak \ldots,\allowbreak \xi_k'^{(\nu)}=0,\allowbreak \pm1,\allowbreak \pm2,\allowbreak \ldots$。所以，这些立方体集合由下列数字组确定：
\[
n,\eta,\varkappa'^{(1)},\xi_1'^{(1)},\ldots,\xi_k'^{(1)},\ldots,\varkappa'^{(n)},\xi_1'^{(n)},\ldots,\xi_k'^{(n)}.
\]
按和
\[
n+\eta+\varkappa'^{(1)}+|\xi_1'^{(1)}|+\cdots+|\xi_k'^{(1)}|+\cdots+\varkappa'^{(n)}+|\xi_1'^{(n)}|+\cdots+|\xi_k'^{(n)}|
\]
的递增次序排列，便得到一个单序列，做法与前面函数线性组合的类似例子完全相同。

继续之前，先回答一个问题：设空间$\mathfrak{R}$满足 \hyperref[ax:II.A]{A}—\hyperref[ax:II.E]{E}，其中取 \hyperref[ax:II.Cinfty]{$C^{(\infty)}$}；在保持$af,\allowbreak f\pm g,\allowbreak (f,\allowbreak g)$的定义不变时，$\mathfrak{R}$的哪些子集$\mathfrak{M}$仍满足 \hyperref[ax:II.A]{A}—\hyperref[ax:II.E]{E}？

\hypertarget{origpage-37}{}\phantomsection\label{origpage-37}要使 \hyperref[ax:II.A]{A} 成立，$\mathfrak{M}$必须是线性流形。\hyperref[ax:II.B]{B} 自动成立。暂且搁下 \hyperref[ax:II.Cn]{C}：总有某个 \hyperref[ax:II.Cn]{$C^{(n)}$} 或 \hyperref[ax:II.Cinfty]{$C^{(\infty)}$} 成立。\hyperref[ax:II.D]{D} 是说，$\mathfrak{M}$中的序列只要满足Cauchy收敛条件，就在$\mathfrak{M}$中有极限。这样的序列在$\mathfrak{R}$中已经有极限，所以只需保证该极限也属于$\mathfrak{M}$；也就是说，$\mathfrak{M}$必须是闭的。至于 \hyperref[ax:II.E]{E}，我们在\CJKglue \hyperref[thm:II.9]{定理9} 的证明中已经看到，它总是成立。因此，合起来说，$\mathfrak{M}$必须是闭线性流形。设张成$\mathfrak{M}$的规范正交系为$\varphi_1,\allowbreak \varphi_2,\allowbreak \ldots$（\hyperref[thm:II.9]{定理9}）。若它无限，显然 \hyperref[ax:II.Cinfty]{$C^{(\infty)}$} 成立，$\mathfrak{M}$便与$\mathfrak{R}_\infty$同构，也就是与$\mathfrak{R}$本身同构。若它到$\varphi_n$为止，即系为$\varphi_1,\allowbreak \ldots,\allowbreak \varphi_n$，则例如由\CJKglue \hyperref[thm:II.3n]{定理$3^{(n)}$} 知 \hyperref[ax:II.Cn]{$C^{(n)}$} 成立，即$\mathfrak{M}$与$\mathfrak{R}_n$同构。

但无论哪种情形，\hyperref[ax:II.D]{D}、\hyperref[ax:II.E]{E} 在$\mathfrak{M}$中都成立，所以在每个$\mathfrak{R}_n$中都成立；因此，它们也可以由 \hyperref[ax:II.A]{A}—\hyperref[ax:II.Cn]{$C^{(n)}$} 推出。

可见，我们运用了逻辑上的技巧，避开了在$\mathfrak{R}_n$或$\mathfrak{R}_\infty$中直接验证 \hyperref[ax:II.A]{A}—\hyperref[ax:II.E]{E}（相应取 \hyperref[ax:II.Cn]{$C^{(n)}$} 或 \hyperref[ax:II.Cinfty]{$C^{(\infty)}$}）。不过，直接验证也没有本质困难，留给读者完成。

最后还须证明 \hyperref[ax:II.D]{D}、\hyperref[ax:II.E]{E} 独立于 \hyperref[ax:II.A]{A}—\hyperref[ax:II.Cinfty]{$C^{(\infty)}$}。刚才已经看到，$\mathfrak{R}_\infty$中的每个线性流形都满足 \hyperref[ax:II.A]{A}、\hyperref[ax:II.B]{B}、\hyperref[ax:II.E]{E}，以及某个 \hyperref[ax:II.Cn]{$C^{(n)}$} 或 \hyperref[ax:II.Cinfty]{$C^{(\infty)}$}；但若它不闭，就不满足 \hyperref[ax:II.D]{D}。这时必定是 \hyperref[ax:II.Cinfty]{$C^{(\infty)}$} 成立，因为 \hyperref[ax:II.Cn]{$C^{(n)}$} 可以推出 \hyperref[ax:II.D]{D}。这样的流形很容易给出：设$\varphi_1,\allowbreak \varphi_2,\allowbreak \ldots$是规范正交系，则所有$\sum_{\nu=1}^Nx_\nu\varphi_\nu$（$N=1,\allowbreak 2,\allowbreak \ldots$，$x_1,\allowbreak \ldots,\allowbreak x_N$任意）构成一个线性流形，却不是闭线性流形。因为$\sum_{\nu=1}^\infty\frac1\nu\varphi_\nu$存在——$\sum_{\nu=1}^\infty(1/\nu)^2$有限——而且是这个流形的聚点，却不是它的元素：
\[
\sum_{\nu=1}^N\frac1\nu\varphi_\nu\longrightarrow\sum_{\nu=1}^\infty\frac1\nu\varphi_\nu\qquad(N\to\infty).
\]
因此，\hyperref[ax:II.D]{D} 独立于 \hyperref[ax:II.A]{A}—\hyperref[ax:II.Cinfty]{$C^{(\infty)}$}、\hyperref[ax:II.E]{E}。

再考察所有以连续参数$\alpha$（$-\infty<\alpha<+\infty$）为自变量的复函数$x(\alpha)$。要求满足$x(\alpha)\ne0$的各个$\alpha$值能排成序列，且对这些$\alpha$求和得到的$\sum_\alpha|x(\alpha)|^2$有限。\origNote[。][全]{55}全部这些$x(\alpha)$组成空间$\mathfrak{R}_{\mathrm{kont.}}$。对于其中任意两点$x(\alpha),\allowbreak y(\alpha)$，$x(\alpha)$和$y(\alpha)$分别只在各自的一列$\alpha$值上不为零；把这两个序列合并为一个，就有某个序列$\alpha_1,\allowbreak \alpha_2,\allowbreak \ldots$，使序列以外恒有$x(\alpha)=y(\alpha)=0$。所以只须讨论$x_n=x(\alpha_n),\allowbreak y_n=y(\alpha_n)$（$n=1,\allowbreak 2,\allowbreak \ldots$）；只要涉及的仅仅是$\mathfrak{R}_{\mathrm{kont.}}$的两个点，一切就与$\mathfrak{R}_\infty$中一样。因此，\hyperref[ax:II.A]{A}、\hyperref[ax:II.B]{B} 在$\mathfrak{R}_{\mathrm{kont.}}$中也同样成立。\origNote[。][]{56}

\hypertarget{origpage-38}{}\phantomsection\label{origpage-38}涉及$k$（$=1,\allowbreak 2,\allowbreak \ldots$）个$\mathfrak{R}_{\mathrm{kont.}}$的点时也一样，所以 \hyperref[ax:II.Cinfty]{$C^{(\infty)}$} 也成立。即使涉及一列$\mathfrak{R}_{\mathrm{kont.}}$的点，情形仍然如此。设这些点为$x_1(\alpha),\allowbreak x_2(\alpha),\allowbreak \ldots$。对每个$n=1,\allowbreak 2,\allowbreak \ldots$，使$x_n(\alpha)\ne0$的$\alpha$构成一个序列$\alpha_1^{(n)},\allowbreak \alpha_2^{(n)},\allowbreak \ldots$。所有这些序列合起来组成双序列$\alpha_m^{(n)}$（$n,\allowbreak m=1,\allowbreak 2,\allowbreak \ldots$），又能写成单序列$\alpha_1^{(1)},\allowbreak \alpha_2^{(1)},\allowbreak \alpha_1^{(2)},\allowbreak \alpha_3^{(1)},\allowbreak \alpha_2^{(2)},\allowbreak \alpha_1^{(3)},\allowbreak \ldots$。因此，\hyperref[ax:II.D]{D} 在$\mathfrak{R}_{\mathrm{kont.}}$中也和$\mathfrak{R}_\infty$中一样成立。\hyperref[ax:II.E]{E} 则不同：它涉及$\mathfrak{R}$的全部点，因为每个点都必须是某个适当序列的聚点。所以不能从$\mathfrak{R}_\infty$的情形推出$\mathfrak{R}_{\mathrm{kont.}}$也满足它。事实上，它确实不成立，因为它的一项推论不成立：这里存在不能排成序列的规范正交系，这与\CJKglue \hyperref[thm:II.3infty]{定理$3^{(\infty)}$} 相违。

令
\[
x_\beta(\alpha)=\begin{cases}1,&\alpha=\beta,\\0,&\alpha\ne\beta.\end{cases}
\]
对每个$\beta$，$x_\beta(\alpha)$都属于$\mathfrak{R}_{\mathrm{kont.}}$，且这些$x_\beta(\alpha)$构成规范正交系。只有当一切$\beta$（$-\infty<\beta<+\infty$）能排成序列时，这个系才能排成序列；但众所周知，这是不可能的。\origNote[。][因]{57}因此，\hyperref[ax:II.E]{E} 也独立于 \hyperref[ax:II.A]{A}—\hyperref[ax:II.Cinfty]{$C^{(\infty)}$}、\hyperref[ax:II.D]{D}。

顺便请注意两种函数空间的根本差别：一种由使$\int_{-\infty}^{\infty}|f(x)|^2\,dx$有限的$f(x)$构成，另一种由使$\sum_\alpha|x(\alpha)|^2$有限的$x(\alpha)$构成。前一种也可以说成由使$\int_{-\infty}^{\infty}|x(\alpha)|^2\,d\alpha$有限的所有$x(\alpha)$构成的空间。两者的全部差别，只在于把$\int_{-\infty}^{\infty}\cdots\,d\alpha$换成$\sum_\alpha\cdots$；尽管如此，前一个空间是$F_\Omega$，因而满足 \hyperref[ax:II.A]{A}—\hyperref[ax:II.E]{E}，与$\mathfrak{R}_\infty$同构；后一个空间$\mathfrak{R}_{\mathrm{kont.}}$却违反 \hyperref[ax:II.E]{E}，与$\mathfrak{R}_\infty$有本质区别。而这两个空间却是同样的，只是其中范数的定义不同！

\section{闭线性流形}\label{sec:II.4}
\hyperref[sec:II.2]{II.2} 对我们很重要，不仅因为其中证明了同构，也因为其中证明了若干关于规范正交系的定理。现在进一步研究希尔伯特空间的几何结构，更仔细地考察闭线性流形。它们在$\mathfrak{R}_\infty$中的作用，类似于直线、平面等等，也就是$\mathfrak{R}_m$（$m\leq n$）在$\mathfrak{R}_n$中的作用。

先回顾\CJKglue \hyperref[def:II.2]{定义2}、\hyperref[def:II.5]{5} 中的记号：若$\mathfrak{A}$是$\mathfrak{R}$的任意子集，$\{\mathfrak{A}\}$和$[\mathfrak{A}]$分别是由$\mathfrak{A}$张成的线性流形和闭线性流形，即包含$\mathfrak{A}$的相应类型的最小对象。

\hypertarget{origpage-39}{}\phantomsection\label{origpage-39}现在扩充这个记法。若$\mathfrak{A},\allowbreak \mathfrak{B},\allowbreak \ldots$为$\mathfrak{R}$的任意子集，$f,\allowbreak g,\allowbreak \ldots$为它的元素，就用$\{\mathfrak{A},\allowbreak \mathfrak{B},\allowbreak \ldots,\allowbreak f,\allowbreak g,\allowbreak \ldots\}$和$[\mathfrak{A},\allowbreak \mathfrak{B},\allowbreak \ldots,\allowbreak f,\allowbreak g,\allowbreak \ldots]$，分别表示把各个$\mathfrak{A},\allowbreak \mathfrak{B},\allowbreak \ldots$和各个$f,\allowbreak g,\allowbreak \ldots$合并后所得集合所张成的线性流形与闭线性流形。

特别地，若$\mathfrak{M},\allowbreak \mathfrak{N},\allowbreak \ldots$是有限或无限多个闭线性流形，就把闭线性流形$[\mathfrak{M},\allowbreak \mathfrak{N},\allowbreak \ldots]$记作$\mathfrak{M}+\mathfrak{N}+\cdots$。显然，$\{\mathfrak{M},\allowbreak \mathfrak{N},\allowbreak \ldots\}$由一切和$f+g+\cdots$构成，其中$f$遍历$\mathfrak{M}$，$g$遍历$\mathfrak{N}$，依此类推；添入聚点，就得到$[\mathfrak{M},\allowbreak \mathfrak{N},\allowbreak \ldots]=\mathfrak{M}+\mathfrak{N}+\cdots$。若集合$\mathfrak{M},\allowbreak \mathfrak{N},\allowbreak \ldots$只有有限多个，且每个集合的每个元素都与其他集合的每个元素正交，那么，这两种构造的结果相同，我们很快就会看到这一点；一般情形却未必如此。

若$\mathfrak{M}$是$\mathfrak{N}$的子集，还可以考察$\mathfrak{N}$中与$\mathfrak{M}$所有元素正交的元素总体。它显然也是闭线性流形，记作$\mathfrak{N}-\mathfrak{M}$。为什么把这种构造称为减法，\CJKglue \hyperref[thm:II.14]{定理14} 将予以说明。特别重要的是$\mathfrak{R}-\mathfrak{M}$，即所有与整个$\mathfrak{M}$正交的$f$构成的集合，称为与$\mathfrak{M}$互补的闭线性流形。

最后指出三个特别简单的闭线性流形：第一是$\mathfrak{R}$本身；第二是仅由$0$构成的集合$\{0\}=[0]$；第三是一切$af$构成的集合，其中$f$为$\mathfrak{R}$的给定元素，$a$为变量。后者显然是闭线性流形，因而同时等于$\{f\}$和$[f]$。

现在引入“投影”的概念，它与欧几里得几何中的投影完全类似：

\leavevmode\phantomsection\textbf{定理10.}\label{thm:II.10} 设$\mathfrak{M}$是闭线性流形。每个$f$都能唯一分解为两项之和$f=g+h$，其中$g$属于$\mathfrak{M}$，$h$属于$\mathfrak{R}-\mathfrak{M}$。

\textkaishu{说明。}称$g$为$f$在$\mathfrak{M}$上的投影，称与整个$\mathfrak{M}$正交的$h$为从$f$到$\mathfrak{M}$的垂线向量。引入记号$P_{\mathfrak{M}}f$表示$g$。

\textkaishu{证明。}依\CJKglue \hyperref[thm:II.9]{定理9}，存在规范正交系$\varphi_1,\allowbreak \varphi_2,\allowbreak \ldots$，所张成的闭线性流形就是$\mathfrak{M}$。令$g=\sum_n(f,\allowbreak \varphi_n)\varphi_n$；根据\CJKglue \hyperref[thm:II.6]{定理6}，这个级数即使有无限项也收敛，其和$g$显然属于$\mathfrak{M}$。又由\CJKglue \hyperref[thm:II.6]{定理6}，$h=f-g$与所有$\varphi_1,\allowbreak \varphi_2,\allowbreak \ldots$正交。但与$h$正交的向量构成一个闭线性流形，所以不仅$\varphi_1,\allowbreak \varphi_2,\allowbreak \ldots$，整个$\mathfrak{M}$也都与$h$正交；即$h$属于$\mathfrak{R}-\mathfrak{M}$。

\hypertarget{origpage-40}{}\phantomsection\label{origpage-40}若还有另一个这样的分解$f=g'+h'$，其中$g'$属于$\mathfrak{M}$，$h'$属于$\mathfrak{R}-\mathfrak{M}$，则$g+h=g'+h'$，所以$g-g'=h'-h=j$。于是$j$同时属于$\mathfrak{M}$和$\mathfrak{R}-\mathfrak{M}$，因而与自身正交：$(j,\allowbreak j)=0$，即$j=0$，故$g=g',\allowbreak h=h'$。

因此，$P_{\mathfrak{M}}f$是一种运算，它给$\mathfrak{R}$中的每个$f$指定在$\mathfrak{M}$上的投影$P_{\mathfrak{M}}f$。下一节将把算符$R$定义为在$\mathfrak{R}$的某个子集上定义、以$\mathfrak{R}$中元素为值的函数，也就是给$\mathfrak{R}$中某些$f$指定$\mathfrak{R}$中的某些$Rf$的对应。这种对应不一定对所有$f$都有定义；对$\mathfrak{R}$中的其他$f$，它可以没有定义，或说“没有意义”。因此，$P_{\mathfrak{M}}$是在$\mathfrak{R}$中处处有定义的算符，称为$\mathfrak{M}$的投影算符。

\leavevmode\phantomsection\textbf{定理11.}\label{thm:II.11} 算符$P_{\mathfrak{M}}$具有下列性质：
\begin{align*}
P_{\mathfrak{M}}(a_1f_1+\cdots+a_nf_n)&=a_1P_{\mathfrak{M}}f_1+\cdots+a_nP_{\mathfrak{M}}f_n,\\
(P_{\mathfrak{M}}f,g)&=(f,P_{\mathfrak{M}}g),\\
P_{\mathfrak{M}}(P_{\mathfrak{M}}f)&=P_{\mathfrak{M}}f.
\end{align*}
$\mathfrak{M}$是$P_{\mathfrak{M}}$的全部值，即所有$P_{\mathfrak{M}}f$的集合；也可以把它刻画为$P_{\mathfrak{M}}f=f$的全部解的集合，而$\mathfrak{R}-\mathfrak{M}$则是$P_{\mathfrak{M}}f=0$的全部解的集合。

\textkaishu{说明。}下一节将看到，第一条性质刻画所谓线性算符，第二条刻画所谓厄米算符。第三条是说，$P_{\mathfrak{M}}$施行两次与施行一次的效果相同；通常将此符号化地写作$P_{\mathfrak{M}}P_{\mathfrak{M}}=P_{\mathfrak{M}}$，或$P_{\mathfrak{M}}^2=P_{\mathfrak{M}}$。

\textkaishu{证明。}由
\begin{gather*}
f_1=g_1+h_1,\quad\ldots,\quad f_n=g_n+h_n,\\
(g_1,\ldots,g_n\in\mathfrak{M};\ h_1,\ldots,h_n\in\mathfrak{R}-\mathfrak{M}),
\end{gather*}
得
\[
a_1f_1+\cdots+a_nf_n=(a_1g_1+\cdots+a_ng_n)+(a_1h_1+\cdots+a_nh_n),
\]
其中$a_1g_1+\cdots+a_ng_n$属于$\mathfrak{M}$，$a_1h_1+\cdots+a_nh_n$属于$\mathfrak{R}-\mathfrak{M}$。所以
\begin{align*}
P_{\mathfrak{M}}(a_1f_1+\cdots+a_nf_n)&=a_1g_1+\cdots+a_ng_n\\
&=a_1P_{\mathfrak{M}}f_1+\cdots+a_nP_{\mathfrak{M}}f_n.
\end{align*}
这就是第一个公式。其次，设
\[
f=g'+h',\quad g=g''+h'',\qquad
(g',g''\in\mathfrak{M};\ h',h''\in\mathfrak{R}-\mathfrak{M}).
\]
则$g',\allowbreak g''$与$h',\allowbreak h''$正交，故
\[
(g',g)=(g',g''+h'')=(g',g'')=(g'+h',g'')=(f,g''),
\]
即$(P_{\mathfrak{M}}f,\allowbreak g)=(f,\allowbreak P_{\mathfrak{M}}g)$。这就是第二个公式。最后，$P_{\mathfrak{M}}f$属于$\mathfrak{M}$，所以$P_{\mathfrak{M}}f=P_{\mathfrak{M}}f+0$就是\CJKglue \hyperref[thm:II.10]{定理10} 所保证的$P_{\mathfrak{M}}f$的分解；即$P_{\mathfrak{M}}(P_{\mathfrak{M}}f)=P_{\mathfrak{M}}f$。这就是第三个公式。

\hypertarget{origpage-41}{}\phantomsection\label{origpage-41}方程$P_{\mathfrak{M}}f=f$与$P_{\mathfrak{M}}f=0$，分别意味着：在$f=g+h$，$g$属于$\mathfrak{M}$、$h$属于$\mathfrak{R}-\mathfrak{M}$的分解（\hyperref[thm:II.10]{定理10}）中，须有$f=g,\allowbreak h=0$，或$g=0,\allowbreak f=h$；即$f$分别属于$\mathfrak{M}$或$\mathfrak{R}-\mathfrak{M}$。这是第五、第六项断言。根据定义，所有$P_{\mathfrak{M}}f$都属于$\mathfrak{M}$；而$\mathfrak{M}$的每个$f'$又等于某个$P_{\mathfrak{M}}f$，例如依刚才所述，就等于$P_{\mathfrak{M}}f'$。这就是第四项断言。

再指出，由本定理的第二、第三个公式有
\[
(P_{\mathfrak{M}}f,P_{\mathfrak{M}}g)=(f,P_{\mathfrak{M}}P_{\mathfrak{M}}g)=(f,P_{\mathfrak{M}}g)=(P_{\mathfrak{M}}f,g).
\]
现在要撇开$\mathfrak{M}$，直接定义投影算符$P_{\mathfrak{M}}$。

\leavevmode\phantomsection\textbf{定理12.}\label{thm:II.12} 处处有定义的算符$E$（算符的含义见\CJKglue \hyperref[thm:II.11]{定理11} 前的说明），是投影算符，即对某个适当的闭线性流形$\mathfrak{M}$有$E=P_{\mathfrak{M}}$，当且仅当它具有下列性质：
\[
(Ef,g)=(f,Eg),\qquad E^2=E
\]
（见\CJKglue \hyperref[thm:II.11]{定理11} 的说明）。而且此时$E$唯一确定$\mathfrak{M}$（\hyperref[thm:II.11]{定理11}）。

\textkaishu{证明。}条件的必要性，以及$E$对$\mathfrak{M}$的确定，由\CJKglue \hyperref[thm:II.11]{定理11} 即知。因此，只须证明：若$E$具有上述性质，就存在闭线性流形$\mathfrak{M}$，使$E=P_{\mathfrak{M}}$。

设$\mathfrak{M}$为所有$Ef$张成的闭线性流形。向量$g-Eg$与所有$Ef$正交，因为
\[
(Ef,g-Eg)=(Ef,g)-(Ef,Eg)=(Ef,g)-(E^2f,g)=0.
\]
$\mathfrak{R}$中与$g-Eg$正交的元素构成闭线性流形，所以既包含所有$Ef$，就也包含$\mathfrak{M}$；因此$g-Eg$属于$\mathfrak{R}-\mathfrak{M}$。这样，$g$关于$\mathfrak{M}$的分解就是$g=Eg+(g-Eg)$，故$P_{\mathfrak{M}}g=Eg$。$g$是任意的，至此全部得证。

若$\mathfrak{M}=\mathfrak{R}$或$[0]$，则$\mathfrak{R}-\mathfrak{M}$分别为$[0]$或$\mathfrak{R}$。所以$f=f+0$或$f=0+f$就是\CJKglue \hyperref[thm:II.11]{定理11} 所说的分解，此时$P_{\mathfrak{M}}f$分别等于$f$或$0$。把由$Rf=f$定义的处处有定义的算符称为算符$1$，把由$Rf=0$定义的算符称为算符$0$。于是$P_{\mathfrak{R}}=1,\allowbreak P_{[0]}=0$。此外，$f$关于$\mathfrak{M}$的分解$f=g+h$，其中$g$属于$\mathfrak{M}$、$h$属于$\mathfrak{R}-\mathfrak{M}$，显然也能改写成$f=h+g$，作为关于$\mathfrak{R}-\mathfrak{M}$的分解。因为$g$属于$\mathfrak{M}$，就与$\mathfrak{R}-\mathfrak{M}$的每个元素正交，从而属于$\mathfrak{R}-(\mathfrak{R}-\mathfrak{M})$。所以$P_{\mathfrak{M}}f=g$，$P_{\mathfrak{R}-\mathfrak{M}}f=h=f-g$，即$P_{\mathfrak{R}-\mathfrak{M}}f=f-P_{\mathfrak{M}}f$。我们把这个关系$P_{\mathfrak{R}-\mathfrak{M}}f=1f-P_{\mathfrak{M}}f$符号化地写作$P_{\mathfrak{R}-\mathfrak{M}}=1-P_{\mathfrak{M}}$。算符的加、减、乘运算，参见\CJKglue \hyperref[thm:II.14]{定理14} 的讨论。

\hypertarget{origpage-42}{}\phantomsection\label{origpage-42}请注意：前面很容易就看出$\mathfrak{M}$是$\mathfrak{R}-(\mathfrak{R}-\mathfrak{M})$的子集；但若直接证明两集合相等，大概要费些周折。然而，由
\[
P_{\mathfrak{R}-(\mathfrak{R}-\mathfrak{M})}=1-P_{\mathfrak{R}-\mathfrak{M}}=1-(1-P_{\mathfrak{M}})=P_{\mathfrak{M}}
\]
立即就可得出这个结论。上述讨论还表明，若$E$是投影算符，$1-E$也是；又因$1-(1-E)=E$，反过来也成立。

\leavevmode\phantomsection\textbf{定理13.}\label{thm:II.13} 恒有
\[
\|Ef\|^2=(Ef,f),\qquad\|Ef\|\leq\|f\|.
\]
条件$\|Ef\|=0$与$\|Ef\|=\|f\|$，分别刻画属于$\mathfrak{R}-\mathfrak{M}$和$\mathfrak{M}$的$f$。

\textkaishu{说明。}特别有
\[
\|Ef-Eg\|=\|E(f-g)\|\leq\|f-g\|;
\]
也就是说，算符$E$连续（参见 \hyperref[sec:II.1]{II.1} 中\CJKglue \hyperref[thm:II.2]{定理2} 之后的讨论）。

\textkaishu{证明。}由\CJKglue \hyperref[thm:II.11]{定理11} 之后的说明，
\[
\|Ef\|^2=(Ef,Ef)=(Ef,f).
\]
又因$1-E$也是投影算符，
\begin{align*}
\|Ef\|^2+\|f-Ef\|^2&=\|Ef\|^2+\|(1-E)f\|^2\\
&=(Ef,f)+((1-E)f,f)=(f,f)=\|f\|^2.
\end{align*}
两个加项都非负，所以也都不超过$\|f\|^2$；特别有$\|Ef\|^2\leq\|f\|^2$，$\|Ef\|\leq\|f\|$。由\CJKglue \hyperref[thm:II.11]{定理11} 已知，$\|Ef\|=0$，即$Ef=0$，表示$f$属于$\mathfrak{R}-\mathfrak{M}$。根据上面的关系，$\|Ef\|=\|f\|$意味着$\|f-Ef\|=0$，即$Ef=f$；再根据\CJKglue \hyperref[thm:II.11]{定理11}，这就是说$f$属于$\mathfrak{M}$。

设$R,\allowbreak S$为两个算符；用$R\pm S,\allowbreak aR$（$a$为复数）、$RS$表示按下式定义的算符：
\[
(R\pm S)f=Rf\pm Sf,\qquad(aR)f=a\cdot Rf,\qquad(RS)f=R(Sf),
\]
自然地，再规定
\[
R^0=1,\quad R^1=R,\quad R^2=RR,\quad R^3=RRR,\quad\ldots.
\]
这些运算满足哪些法则，很容易讨论。对于$R\pm S,\allowbreak aR$，不难验证，数的全部初等运算法则都成立；对于$RS$却不尽如此。可验证分配律$(R\pm S)T=RT\pm ST$及$R(S\pm T)=RS\pm RT$成立，不过后一条需要$R$的线性性（参见\CJKglue \hyperref[thm:II.11]{定理11} 的说明及下一节）；结合律$(RS)T=R(ST)=RST$也成立。但是交换律$RS=SR$一般并不成立，因为$(RS)f=R(Sf)$与$(SR)f=S(Rf)$未必相等。

\hypertarget{origpage-43}{}\phantomsection\label{origpage-43}若对某两个特定的$R,\allowbreak S$，交换律确实成立，就称它们可交换，或称它们对易。例如，$0$和$1$都与任何处处有定义的$R$对易：
\[
R0=0R=0,\qquad R1=1R=R;
\]
$R^m,\allowbreak R^n$也对易，因为$R^mR^n=R^{m+n}$，与$m,\allowbreak n$的次序无关。

\leavevmode\phantomsection\textbf{定理14.}\label{thm:II.14} 设$E,\allowbreak F$为闭线性流形$\mathfrak{M},\allowbreak \mathfrak{N}$各自的投影算符。$EF$也是投影算符，当且仅当$E,\allowbreak F$对易，即$EF=FE$；此时它对应于由$\mathfrak{M},\allowbreak \mathfrak{N}$的公共元素构成的闭线性流形$\mathfrak{P}$。$E+F$为投影算符，当且仅当$EF=0$，也等价于$FE=0$。这意味着整个$\mathfrak{M}$与整个$\mathfrak{N}$正交；这时$E+F$对应于$\mathfrak{M}+\mathfrak{N}=[\mathfrak{M},\allowbreak \mathfrak{N}]$，而此流形在这里等于$\{\mathfrak{M},\allowbreak \mathfrak{N}\}$。$E-F$为投影算符，当且仅当$EF=F$，也等价于$FE=F$。这意味着$\mathfrak{N}$是$\mathfrak{M}$的子集；这时$E-F$对应于$\mathfrak{M}-\mathfrak{N}$。

\textkaishu{证明。}对于$EF$，须验证\CJKglue \hyperref[thm:II.12]{定理12} 的两个条件：
\[
(EFf,g)=(f,EFg),\qquad(EF)^2=EF.
\]
由于$(EFf,\allowbreak g)=(Ff,\allowbreak Eg)=(f,\allowbreak FEg)$，第一个条件就是
\[
(f,EFg)=(f,FEg),\qquad(f,(EF-FE)g)=0.
\]
此式对一切$f$成立，所以$(EF-FE)g=0$；又对一切$g$成立，所以$EF-FE=0$，即$EF=FE$。因此，仅对第一个条件来说，对易性已是充分必要条件；而它也能推出第二个条件：
\[
(EF)^2=EFEF=EEFF=E^2F^2=EF.
\]
由于$E,\allowbreak F$满足相应关系，$E+F$总满足$((E+F)f,\allowbreak g)=(f,\allowbreak (E+F)g)$。所以这里只须验证第二个条件$(E+F)^2=E+F$。因为
\[
(E+F)^2=E^2+F^2+EF+FE=(E+F)+(EF+FE),
\]
这个条件只不过就是$EF+FE=0$。若$EF=0$，则$EF$为投影算符，由刚才的证明知$EF=FE$，故$EF+FE=0$。反过来，若$EF+FE=0$，则
\begin{align*}
E(EF+FE)&=E^2F+EFE=EF+EFE=0,\\
E(EF+FE)E&=E^2FE+EFE^2=EFE+EFE=2\cdot EFE=0.
\end{align*}
所以$EFE=0$，从而$EF=0$。因此$EF=0$是充分必要条件；由于$E,\allowbreak F$的地位相同，$FE=0$也是。

$E-F$是投影算符，当且仅当$1-(E-F)=(1-E)+F$是投影算符。$1-E,\allowbreak F$都是投影算符，所以根据刚才的证明，这等价于$(1-E)F=0$，即$F-EF=0$，即$EF=F$；也同样等价于$F(1-E)=0$，即$F-FE=0$，即$FE=F$。

\hypertarget{origpage-44}{}\phantomsection\label{origpage-44}还须证明关于$\mathfrak{M},\allowbreak \mathfrak{N}$的断言，其中$E=P_{\mathfrak{M}},\allowbreak F=P_{\mathfrak{N}}$。首先设$EF=FE$，则每个$EFf=FEf$同时属于$\mathfrak{M}$和$\mathfrak{N}$，因而属于$\mathfrak{P}$。对$\mathfrak{P}$的每个$g$，有$Eg=Fg=g$，所以$EFg=Eg=g$，即$g$具有$EFf$的形式。因此，$\mathfrak{P}$就是$EF$的值域，由\CJKglue \hyperref[thm:II.11]{定理11} 知$EF=P_{\mathfrak{P}}$。其次，设$EF=0$，从而$FE=0$。每个$(E+F)f=Ef+Ff$都属于$\{\mathfrak{M},\allowbreak \mathfrak{N}\}$；而其中的每个$g$都等于$h+j$，其中$h$属于$\mathfrak{M}$，$j$属于$\mathfrak{N}$。所以$Eh=h,\allowbreak Fh=FEh=0,\allowbreak Fj=j,\allowbreak Ej=EFj=0$，故
\[
(E+F)(h+j)=Eh+Fh+Ej+Fj=h+j,\qquad(E+F)g=g.
\]
即$g$具有$(E+F)f$的形式。因此，$\{\mathfrak{M},\allowbreak \mathfrak{N}\}$是$E+F$的值域；但$E+F$是投影算符，所以$\{\mathfrak{M},\allowbreak \mathfrak{N}\}$是它所对应的闭线性流形（\hyperref[thm:II.11]{定理11}）。既然$\{\mathfrak{M},\allowbreak \mathfrak{N}\}$是闭的，它就等于$[\mathfrak{M},\allowbreak \mathfrak{N}]=\mathfrak{M}+\mathfrak{N}$。第三，设$EF=F$，从而$FE=F$。此时$E=P_{\mathfrak{M}}$，$1-F=P_{\mathfrak{R}-\mathfrak{N}}$，所以$E-F=E-EF=E(1-F)$等于$P_{\mathfrak{P}}$，其中$\mathfrak{P}$是$\mathfrak{M}$与$\mathfrak{R}-\mathfrak{N}$的公共部分，也就是$\mathfrak{M}-\mathfrak{N}$。

最后，$EF=0$意味着恒有$(EFf,\allowbreak g)=0$，从而$(Ff,\allowbreak Eg)=0$；即整个$\mathfrak{M}$与整个$\mathfrak{N}$正交。而$EF=F$意味着$F(1-E)=0$：整个$\mathfrak{N}$与$\mathfrak{R}-\mathfrak{M}$正交；也等价于$\mathfrak{N}$是$\mathfrak{R}-(\mathfrak{R}-\mathfrak{M})=\mathfrak{M}$的子集。

若$\mathfrak{N}$是$\mathfrak{M}$的子集，也说$F=P_{\mathfrak{N}}$是$E=P_{\mathfrak{M}}$的一部分，记作$E\geq F$或$F\leq E$。这意味着$EF=F$，也等价于$FE=F$，因而包含了对易性。无论考察$\mathfrak{M},\allowbreak \mathfrak{N}$，还是直接计算，都能看出：恒有$0\leq E\leq1$；由$E\leq F,\allowbreak F\leq E$可得$E=F$；由$E\leq F,\allowbreak F\leq G$可得$E\leq G$。因此，关系$\leq$具有“大小次序”的性质。还须注意，$E\leq F$、$1-E\geq1-F$、$E$与$1-F$正交，这三者互相等价。此外，若$E,\allowbreak F$正交，且$E'\leq E,\allowbreak F'\leq F$，则$E',\allowbreak F'$也正交。若$\mathfrak{M},\allowbreak \mathfrak{N}$正交，我们也称$E,\allowbreak F$正交，即$EF=0$，或等价地$FE=0$。反过来，若$E,\allowbreak F$对易，我们也称相应的$\mathfrak{M},\allowbreak \mathfrak{N}$对易。

\leavevmode\phantomsection\textbf{定理15.}\label{thm:II.15} 关系$E\leq F$等价于对一切$f$均有$\|Ef\|\leq\|Ff\|$。

\textkaishu{证明。}由$E\leq F$得$E=EF$，再由\CJKglue \hyperref[thm:II.13]{定理13} 得$\|Ef\|=\|EFf\|\leq\|Ff\|$。反过来，这个不等式意味着：若$Ff=0$，则$\|Ef\|\leq\|Ff\|=0$，即$Ef=0$。由于$F(1-F)f=(F-F^2)f=0$，便恒有$E(1-F)f=0$，即$E(1-F)=E-EF=0$，故$E=EF$，从而$E\leq F$。

\leavevmode\phantomsection\textbf{定理16.}\label{thm:II.16} 设$E_1,\allowbreak \ldots,\allowbreak E_k$为投影算符。$E_1+\cdots+E_k$也是投影算符，当且仅当所有$E_m,\allowbreak E_l$（$m,\allowbreak l=1,\allowbreak \ldots,\allowbreak k$，$m\ne l$）两两正交。

\hypertarget{origpage-45}{}\phantomsection\label{origpage-45}另一个充分必要条件是：对一切$f$都有
\[
\|E_1f\|^2+\cdots+\|E_kf\|^2\leq\|f\|^2
\]
而且此时$E_1+\cdots+E_k$是$\mathfrak{M}_1+\cdots+\mathfrak{M}_k=[\mathfrak{M}_1,\allowbreak \ldots,\allowbreak \mathfrak{M}_k]$的投影算符；这个流形在这里等于$\{\mathfrak{M}_1,\allowbreak \ldots,\allowbreak \mathfrak{M}_k\}$，其中$E_1=P_{\mathfrak{M}_1},\allowbreak \ldots,\allowbreak E_k=P_{\mathfrak{M}_k}$。

\textkaishu{证明。}反复应用\CJKglue \hyperref[thm:II.14]{定理14}，便得最后一项断言，也得到第一个判据的充分性。若第二个判据成立，第一个也成立；因为当$m\ne l$且$E_mf=f$时，
\begin{align*}
\|f\|^2+\|E_lf\|^2&=\|E_mf\|^2+\|E_lf\|^2\\
&\leq\|E_1f\|^2+\cdots+\|E_kf\|^2\leq\|f\|^2,\\
\|E_lf\|^2&=0,\quad E_lf=0.
\end{align*}
而$E_m(E_mf)=E_mf$恒成立，所以$E_l(E_mf)=0$，即$E_lE_m=0$。最后证明第二个条件的必要性：若$E_1+\cdots+E_k$是投影算符，依\CJKglue \hyperref[thm:II.13]{定理13}，
\begin{align*}
\|E_1f\|^2+\cdots+\|E_kf\|^2&=(E_1f,f)+\cdots+(E_kf,f)\\
&=((E_1+\cdots+E_k)f,f)\\
&=\|(E_1+\cdots+E_k)f\|^2\leq\|f\|^2.
\end{align*}
因此得到如下逻辑关系：
\begin{gather*}
E_1+\cdots+E_k\mathwordspace{}\text{是投影算符}\ \Longrightarrow\ \text{第二个判据}\\
\Longrightarrow\ \text{第一个判据}\ \Longrightarrow\ E_1+\cdots+E_k\mathwordspace{}\text{是投影算符}.
\end{gather*}
所以三者互相等价。

最后再证明一个关于投影算符收敛的定理：

\leavevmode\phantomsection\textbf{定理17.}\label{thm:II.17} 设$E_1,\allowbreak E_2,\allowbreak \ldots$是单调递增或递减的投影算符序列：$E_1\leq E_2\leq\cdots$或$E_1\geq E_2\geq\cdots$。则它们收敛于某个投影算符$E$，收敛的意义是对每个$f$都有$E_nf\to Ef$；而且相应地，一切$E_n\leq E$或一切$E_n\geq E$。

\textkaishu{证明。}只讨论第二种情形就够了，因为把$E_1,\allowbreak E_2,\allowbreak \ldots,\allowbreak E$换成$1-E_1,\allowbreak 1-E_2,\allowbreak \ldots,\allowbreak 1-E$，就把第一种情形归结到第二种。因此，设$E_1\geq E_2\geq\cdots$。

根据\CJKglue \hyperref[thm:II.15]{定理15}，$\|E_1f\|^2\geq\|E_2f\|^2\geq\cdots\geq0$，所以$\lim_{m\to\infty}\|E_mf\|^2$存在。因此，对每个$\varepsilon>0$都有$N=N(\varepsilon)$，使$m,\allowbreak l\geq N$时有$|\|E_mf\|^2-\|E_lf\|^2|<\varepsilon$。当$m\leq l$时，$E_m\geq E_l$，而$E_m-E_l$是投影算符，故
\begin{align*}
\|E_mf\|^2-\|E_lf\|^2&=(E_mf,f)-(E_lf,f)=((E_m-E_l)f,f)\\
&=\|(E_m-E_l)f\|^2=\|E_mf-E_lf\|^2,
\end{align*}
从而$\|E_mf-E_lf\|<\sqrt\varepsilon$。所以序列$E_1f,\allowbreak E_2f,\allowbreak \ldots$满足Cauchy收敛条件，由 \hyperref[sec:II.1]{II.1} 的 \hyperref[ax:II.D]{D} 知，它有极限$f^*$。于是以$Ef=f^*$定义一个处处有定义的算符。

\hypertarget{origpage-46}{}\phantomsection\label{origpage-46}在$(E_nf,\allowbreak g)=(f,\allowbreak E_ng)$中取极限，得到$(Ef,\allowbreak g)=(f,\allowbreak Eg)$；在$(E_nf,\allowbreak E_ng)=(E_nf,\allowbreak g)$中取极限，得到$(Ef,\allowbreak Eg)=(Ef,\allowbreak g)$；所以$(E^2f,\allowbreak g)=(Ef,\allowbreak g)$，即$E^2=E$。因此$E$为投影算符。当$l\geq m$时，有$\|E_mf\|\geq\|E_nf\|$；令$l\to\infty$，得$\|E_mf\|\geq\|Ef\|$，故由\CJKglue \hyperref[thm:II.15]{定理15} 知$E_m\geq E$。\tnote[。][][原]{原书在令$l\to\infty$之前的不等式中写作$E_n$，其指标应为$l$。这里保留原书指标。}

若$E_1,\allowbreak E_2,\allowbreak \ldots$是两两正交的投影算符，则$E_1,\allowbreak E_1+E_2,\allowbreak E_1+E_2+E_3,\allowbreak \ldots$都是投影算符（\hyperref[thm:II.16]{定理16}），而且显然组成递增序列。根据\CJKglue \hyperref[thm:II.17]{定理17}，它们收敛于某个不小于所有这些算符的投影算符，可记作$E_1+E_2+\cdots$。设$E_1=P_{\mathfrak{M}_1},\allowbreak E_2=P_{\mathfrak{M}_2},\allowbreak \ldots,\allowbreak E_1+E_2+\cdots=P_{\mathfrak{M}}$。因为所有$E_m\leq E$，$\mathfrak{M}_m$是$\mathfrak{M}$的子集，所以$\mathfrak{M}$也包含$[\mathfrak{M}_1,\allowbreak \mathfrak{M}_2,\allowbreak \ldots]=\mathfrak{M}_1+\mathfrak{M}_2+\cdots=\mathfrak{M}'$。反过来，所有$\mathfrak{M}_m$都是$\mathfrak{M}'$的子集，所以$E_m\leq P_{\mathfrak{M}'}=E'$。利用连续性，有$E\leq E'$（参照上面证明中的推论），所以$\mathfrak{M}$又是$\mathfrak{M}'$的子集。故$\mathfrak{M}=\mathfrak{M}'$，$E=E'$，即$\mathfrak{M}=\mathfrak{M}_1+\mathfrak{M}_2+\cdots$；也可写成
\[
P_{\mathfrak{M}_1+\mathfrak{M}_2+\cdots}=P_{\mathfrak{M}_1}+P_{\mathfrak{M}_2}+\cdots.
\]
关于投影算符的讨论到此结束。

\section{希尔伯特空间中的算符}\label{sec:II.5}
现在，对于无限维空间，也就是希尔伯特空间$\mathfrak{R}_\infty$的几何关系，我们已有足够的了解，可以转而考察它的线性算符，即$\mathfrak{R}_\infty$到自身的线性映射。为此必须引入若干概念；其中一部分在上一节已经预先用到。

上一节已经涉及算符。依照其中\CJKglue \hyperref[thm:II.11]{定理11} 之前的说明，作如下定义：

\leavevmode\phantomsection\textbf{定义6.}\label{def:II.6} 算符$R$是定义在$\mathfrak{R}$的某个子集上、以$\mathfrak{R}$中的元素为值的函数；也就是给$\mathfrak{R}$中的某些元素$f$指定$\mathfrak{R}$中的某些元素$Rf$的对应。

这里除了$\mathfrak{R}_\infty$，也允许$\mathfrak{R}_n$。请注意，若$\mathfrak{R}_\infty$是一个$F_\Omega$，算符$R$就定义在$F_\Omega$的元素上，即普通的位形空间函数上，其值也同样是这类函数。因此，这些算符就是所谓“函数的函数”或“泛函”（参见 \hyperref[sec:I.2]{I.2}、\hyperref[sec:I.4]{I.4} 中的例子）。使$Rf$有意义的$f$构成的类，即$R$的定义域，不必包含整个$\mathfrak{R}$；若确实包含，就称$R$处处有定义。此外，$Rf$构成的类，即$R$的值域——定义域经$R$映射所得的像——也完全不必包含在定义域中。因此，$Rf$有意义，并不意味着$R(Rf)=R^2f$也有意义。\origNote[。][]{58}

\hypertarget{origpage-47}{}\phantomsection\label{origpage-47}$R\pm S,\allowbreak aR,\allowbreak RS,\allowbreak R^m$（$R,\allowbreak S$为算符，$a$为复数，$m=0,\allowbreak 1,\allowbreak 2,\allowbreak \ldots$）的含义，在上一节已经说明：
\begin{gather*}
(R\pm S)f=Rf\pm Sf,\qquad(aR)f=a\cdot Rf,\qquad(RS)f=R(Sf),\\
R^0=1,\quad R^1=R,\quad R^2=RR,\quad R^3=RRR,\quad\ldots.
\end{gather*}
不过，确定定义域时必须注意：只有右端有意义，左端，即算符$R\pm S,\allowbreak aR,\allowbreak RS$才有意义。例如，$R\pm S$仅在$R,\allowbreak S$定义域的公共部分上有定义，依此类推。若$Rf$对于它所能取的每个值都只取一次，就有逆算符$R^{-1}$：只要$Rg=f$有解$g$，$R^{-1}f$就有意义，且等于这个$g$。关于$R\pm S,\allowbreak aR,\allowbreak RS$所满足的运算法则，上一节已经讨论；这里仅补充关于定义域的一点说明。那里所称相等的算符，也有相同的定义域；但是$0\cdot R=0$这类算符等式，就定义域而言并不成立：$0f$总有意义，而按定义，只有$Rf$有意义，$(0\cdot R)f$才有意义；当然，只要两边都有意义，就都等于零。相反，$1\cdot R=R\cdot1=R$，以及$R^m\cdot R^l=R^{m+l}$，在定义域方面也成立。

若$R,\allowbreak S$有逆算符，$RS$也有逆算符；不难看出，$(RS)^{-1}=S^{-1}R^{-1}$。此外，当$a\ne0$时，$(aR)^{-1}=\frac1a R^{-1}$。若$R^{-1}$存在，还能构造$R$的其他负整数次幂：
\[
R^{-2}=R^{-1}R^{-1},\qquad R^{-3}=R^{-1}R^{-1}R^{-1},\quad\ldots.
\]
作了这些一般说明之后，下面更仔细地考察以后对我们特别重要的几类特殊算符。

\leavevmode\phantomsection\textbf{定义7.}\label{def:II.7} 若算符$A$的定义域是线性流形，即只要包含$f_1,\allowbreak \ldots,\allowbreak f_k$，就也包含$a_1f_1+\cdots+a_kf_k$，并且满足
\[
A(a_1f_1+\cdots+a_kf_k)=a_1Af_1+\cdots+a_kAf_k
\]
就称$A$为线性算符。以后只考察线性算符，而且只考察定义域处处稠密的线性算符。

对于量子力学中的算符，我们必须放弃处处有定义的要求；不过，在某些情况下，要求定义域处处稠密就已足够。这一点很重要，值得稍加细究。例如，考察Schr\"odinger波动力学的位形空间，为简单起见，设它是一维的：$-\infty<q<+\infty$。

\hypertarget{origpage-48}{}\phantomsection\label{origpage-48}波函数是使$\int_{-\infty}^{+\infty}|\varphi(q)|^2\,dq$有限的$\varphi(q)$，它们构成一个希尔伯特空间（见 \hyperref[sec:II.3]{II.3}）。再考察算符$q\cdots$和$\frac{h}{2\pi i}\frac{d}{dq}\cdots$。它们显然是线性算符，但其定义域都不是整个希尔伯特空间。对于$q\cdots$，这是因为即使$\int_{-\infty}^{+\infty}|\varphi(q)|^2\,dq$有限，
\[
\int_{-\infty}^{+\infty}|q\varphi(q)|^2\,dq=\int_{-\infty}^{+\infty}q^2|\varphi(q)|^2\,dq
\]
也完全可能无限，从而$q\varphi(q)$不再属于希尔伯特空间。对于$\frac{h}{2\pi i}\frac{d}{dq}$，则是因为有些函数不可微；另一些函数虽使$\int_{-\infty}^{+\infty}|\varphi(q)|^2\,dq$有限，却不能使
\[
\int_{-\infty}^{+\infty}\left|\frac{h}{2\pi i}\frac{d}{dq}\varphi(q)\right|^2\,dq
=\frac{h^2}{4\pi^2}\int_{-\infty}^{+\infty}\left|\frac{d}{dq}\varphi(q)\right|^2\,dq
\]
有限，例如$|q|^{1/2}e^{-q^2}$或$e^{-q^2}\sin(e^{q^2})$。不过，这两个定义域都是处处稠密的：凡$\varphi(q)$仅在有限区间$-C\leq q\leq C$内不为零，而且处处连续可微，这两个算符就都能作用于它；而这类函数的集合处处稠密。\origNote[。][]{59}

进一步定义：

\leavevmode\phantomsection\textbf{定义8.}\label{def:II.8} 两个算符$A,\allowbreak A^*$，若定义域相同，且在此定义域上恒有
\[
(Af,g)=(f,A^*g),\qquad(A^*f,g)=(f,Ag)
\]
就称它们互为伴随。交换$f,\allowbreak g$并对两边取复共轭，可知这两个关系中的任何一个都能推出另一个。另外，$A,\allowbreak A^*$之间的关系显然是对称的，即$A^*,\allowbreak A$也互为伴随。所以$A^{**}=A$。

还要指出，$A$至多有一个伴随算符$A^*$：若$A$既与$A_1^*$互为伴随，又与$A_2^*$互为伴随，则必须有$A_1^*=A_2^*$。因为对于使$Ag$有意义的每个$g$，有
\[
(A_1^*f,g)=(f,Ag)=(A_2^*f,g).
\]
而这些$g$处处稠密，所以$A_1^*f=A_2^*f$；此式又普遍成立，故$A_1^*=A_2^*$。于是$A$唯一确定$A^*$，同样，$A^*$也唯一确定$A$。

不难看出，$0,\allowbreak 1$，乃至一切投影算符$E$，都与自身伴随（参见\CJKglue \hyperref[thm:II.12]{定理12}）；即$0^*,\allowbreak 1^*,\allowbreak E^*$存在，且分别等于$0,\allowbreak 1,\allowbreak E$。此外，$(aA)^*=\overline aA^*$；只要$A\pm B$能够构造，即其定义域处处稠密，就有$(A\pm B)^*=A^*\pm B^*$。最后，对定义域加上容易确定的限制，还会有$(AB)^*=B^*A^*$，因为$(ABf,\allowbreak g)=(Bf,\allowbreak A^*g)=(f,\allowbreak B^*A^*g)$；以及$(A^{-1})^*=A^{*-1}$，因为
\[
(A^{-1}f,g)=(A^{-1}f,A^*A^{*-1}g)=(AA^{-1}f,A^{*-1}g)=(f,A^{*-1}g).
\]

\hypertarget{origpage-49}{}\phantomsection\label{origpage-49}特别地，在刚才考察的Schr\"odinger波动力学中，现设位形空间为$k$维，希尔伯特空间便由使$\int_{-\infty}^{\infty}\cdots\int_{-\infty}^{\infty}|\varphi(q_1\cdots q_k)|^2\,dq_1\cdots dq_k$有限的$\varphi(q_1,\allowbreak \ldots,\allowbreak q_k)$构成。对于算符$q_l\cdots$和$\frac{h}{2\pi i}\frac{\partial}{\partial q_l}\cdots$有
\[
(q_l)^*=q_l,\qquad
\left(\frac{h}{2\pi i}\frac{\partial}{\partial q_l}\right)^*
=\frac{h}{2\pi i}\frac{\partial}{\partial q_l}.
\]
第一式由
\begin{align*}
&\int_{-\infty}^{\infty}\cdots\int_{-\infty}^{\infty}q_l\varphi(q_1\cdots q_k)\overline{\psi(q_1\cdots q_k)}\,dq_1\cdots dq_k\\
&\qquad=\int_{-\infty}^{\infty}\cdots\int_{-\infty}^{\infty}\varphi(q_1\cdots q_k)\overline{q_l\psi(q_1\cdots q_k)}\,dq_1\cdots dq_k
\end{align*}
立即可知。第二式是说
\begin{align*}
&\int_{-\infty}^{\infty}\cdots\int_{-\infty}^{\infty}\frac{h}{2\pi i}\frac{\partial}{\partial q_l}\varphi(q_1\cdots q_k)\overline{\psi(q_1\cdots q_k)}\,dq_1\cdots dq_k\\
&\qquad=\int_{-\infty}^{\infty}\cdots\int_{-\infty}^{\infty}\varphi(q_1\cdots q_k)\overline{\frac{h}{2\pi i}\frac{\partial}{\partial q_l}\psi(q_1\cdots q_k)}\,dq_1\cdots dq_k,
\end{align*}
即
\begin{align*}
\int_{-\infty}^{\infty}\cdots\int_{-\infty}^{\infty}
\biggl\{&\frac{\partial}{\partial q_l}\varphi(q_1\cdots q_k)\overline{\psi(q_1\cdots q_k)}\\
&+\varphi(q_1\cdots q_k)\frac{\partial}{\partial q_l}\overline{\psi(q_1\cdots q_k)}\biggr\}\,dq_1\cdots dq_k=0,
\end{align*}
或
\begin{align*}
\lim_{A\to+\infty,\,B\to+\infty}
\int_{-\infty}^{\infty}\cdots\int_{-\infty}^{\infty}
\bigl[\varphi(q_1\cdots q_k)\overline{\psi(q_1\cdots q_k)}\bigr]_{q_l=-B}^{q_l=+A}\\
{}\times dq_1\cdots dq_{l-1}\,dq_{l+1}\cdots dq_k=0.
\end{align*}
这个极限一定存在，因为所有积分
\[
\int_{-\infty}^{\infty}\cdots\int_{-\infty}^{\infty}\cdots\,dq_1\cdots dq_k
\]
的收敛性都有保证——$\varphi,\allowbreak \psi,\allowbreak \frac{\partial}{\partial x_l}\varphi,\allowbreak \frac{\partial}{\partial x_l}\psi$都属于希尔伯特空间。现在只须证明极限为零。若极限不为零，则积分
\[
\int_{-\infty}^{\infty}\cdots\int_{-\infty}^{\infty}\varphi(q_1\cdots q_k)\overline{\psi(q_1\cdots q_k)}\,dq_1\cdots dq_{l-1}\,dq_{l+1}\cdots dq_k
\]
在$q_l\to+\infty$或$q_l\to-\infty$时必定存在的极限就不为零；这却与
\[
\int_{-\infty}^{\infty}\cdots\int_{-\infty}^{\infty}\varphi(q_1\cdots q_k)\overline{\psi(q_1\cdots q_k)}\,dq_1\cdots dq_{l-1}\,dq_l\,dq_{l+1}\cdots dq_k
\]
绝对收敛不相容，因为$\varphi,\allowbreak \psi$属于希尔伯特空间。

\hypertarget{origpage-50}{}\phantomsection\label{origpage-50}若$A$是积分算符
\[
A\varphi(q_1\cdots q_k)=\int_{-\infty}^{\infty}\cdots\int_{-\infty}^{\infty}K(q_1\cdots q_k;q'_1\cdots q'_k)\varphi(q'_1\cdots q'_k)\,dq'_1\cdots dq'_k,
\]
不难看出，$A^*$也是积分算符，只是它的核不再是$K(q_1\cdots q_k;q'_1\cdots q'_k)$，而是
\[
\overline{K(q'_1\cdots q'_k;q_1\cdots q_k)}.
\]

现在考察矩阵理论的情形；此时希尔伯特空间由满足$\sum_{n=1}^\infty|x_n|^2$有限的一切序列$x_1,\allowbreak x_2,\allowbreak \ldots$构成。线性算符$A$把$\{x_1,\allowbreak x_2,\allowbreak \ldots\}$变为$\{y_1,\allowbreak y_2,\allowbreak \ldots\}$：
\[
A\{x_1,x_2,\ldots\}=\{y_1,y_2,\ldots\}.
\]
由于$A$线性，$y_1,\allowbreak y_2,\allowbreak \ldots$必须线性地依赖于$x_1,\allowbreak x_2,\allowbreak \ldots$：\origNote[：][]{60}
\[
y_\mu=\sum_{\nu=1}^\infty a_{\mu\nu}x_\nu.
\]
因此，给出矩阵$a_{\mu\nu}$就确定了$A$。立即可见，与$A^*$相应的矩阵是$\overline{a_{\nu\mu}}$，即复共轭转置矩阵。\origNote[。][]{60}

刚才说明的、与矩阵理论的类比，自然引出下面要给出的厄米算符概念。同时再引入两个在后面很重要的概念：

\leavevmode\phantomsection\textbf{定义9.}\label{def:II.9} 若$A^*=A$，就称算符$A$为厄米算符。若此外恒有$(Af,\allowbreak f)\geq0$，还称它为定号的\tnote[的][。][原]{原书的历史术语\emph{definit}在这里允许取零；现代通常称为非负或半正定，并非严格正定。}。\origNote[。][若]{61}若$UU^*=U^*U=1$，就称算符$U$为酉算符。\origNote[。][]{62}

因此，对于酉算符有$U^*=U^{-1}$。按照定义，
\[
(Uf,Ug)=(U^*Uf,g)=(f,g),
\]
特别当$f=g$时，有$\|Uf\|=\|f\|$。反过来，如果$U$处处有定义，且能取遍所有值，后一性质就能推出酉性（见\CJKglue \hyperlink{note-62}{原注 \getrefnumber{orig-note-62}}）。证明如下：首先有
\[
\|Uf\|=\|f\|,\qquad\text{即}\mathwordspace{}(Uf,Uf)=(f,f),\quad(U^*Uf,f)=(f,f).
\]
把$f$分别换成$(f+g)/2$和$(f-g)/2$，然后相减，简单计算便得$\operatorname{Re}(Uf,\allowbreak Ug)=\operatorname{Re}(f,\allowbreak g)$。

\hypertarget{origpage-51}{}\phantomsection\label{origpage-51}再把这里的$f$换成$if$，便由$\operatorname{Re}$得到$\operatorname{Im}$。所以一般有
\[
(Uf,Ug)=(f,g),\qquad\text{即}\mathwordspace{}(U^*Uf,g)=(f,g).
\]
固定$f$，此式对一切$g$成立，故$U^*Uf=f$；又对一切$f$成立，故$U^*U=1$。还须证明$UU^*=1$。对每个$f$都有一个$g$使$Ug=f$，于是
\[
UU^*f=UU^*\cdot Ug=U\cdot U^*Ug=Ug=f,\qquad\text{故}\mathwordspace{}UU^*=1.
\]
又因由线性性有
\[
\|Uf-Ug\|=\|U(f-g)\|=\|f-g\|,
\]
所以每个酉算符$U$都连续。厄米算符却未必连续；例如，对量子力学很重要的算符$q\cdots$和$\frac{h}{2\pi i}\frac{d}{dq}\cdots$恰好就是不连续的。\origNote[。][]{63}

由$A^*$的形式运算法则，立即可得：若$U,\allowbreak V$是酉的，$U^{-1},\allowbreak UV$也酉，从而$U$的一切整数次幂也酉；若$A,\allowbreak B$是厄米的，$A\pm B$也厄米，但$aA$只有在$a$为实数时才厄米（$A=0$除外），$AB$则只有在$A,\allowbreak B$对易，即$AB=BA$时才厄米。我们还知道，一切投影算符，特别是$0,\allowbreak 1$，都是厄米的；Schr\"odinger理论中的算符$q_l\cdots,\allowbreak \frac{h}{2\pi i}\frac{\partial}{\partial q_l}\cdots$也如此。因此，$A$厄米，则它的一切幂也厄米（若$A^{-1}$存在，也包括它），它的一切实系数多项式也厄米。另有一点值得注意：若$A$厄米，$X$任意，则$XAX^*$也厄米：
\[
(XAX^*)^*=X^{**}A^*X^*=XAX^*.
\]
例如，一切$XX^*$（取$A=1$）和$X^*X$（以$X^*$代$X$）都是如此。对酉算符$U$，$UAU^{-1}$厄米，因为$U^{-1}=U^*$。

连续性对于算符至关重要，就像它对于分析学中熟悉的数值函数一样。因此，下面给出线性算符连续的若干特征条件。

\leavevmode\phantomsection\textbf{定理18.}\label{thm:II.18} 线性算符$R$只要在$f=0$处连续，就处处连续。而在$f=0$处连续的充分必要条件，是存在常数$C$，使普遍有$\|Rf\|\leq C\|f\|$。这个条件又等价于普遍有
\[
|(Rf,g)|\leq C\|f\|\|g\|.
\]
当$R$厄米时，只须对$f=g$作此要求：$|(Rf,\allowbreak f)|\leq C\|f\|^2$；又因$(Rf,\allowbreak f)$为实数（\hyperlink{note-61}{原注 \getrefnumber{orig-note-61}}），也就是
\[
-C\|f\|^2\leq(Rf,f)\leq C\|f\|^2.
\]

\hypertarget{origpage-52}{}\phantomsection\label{origpage-52}\textkaishu{说明。}算符连续的概念源于Hilbert，\origNote[，][他]{64}他称之为“有界性”，并以上面倒数第二个判据作为定义。若最后一个判据中的两个不等式只有一个普遍成立，就相应称$R$为下半有界或上半有界。例如，每个定号的$R$都是下半有界的，此时$C=0$。

\textkaishu{证明。}在$f=0$处连续，是说对每个$\varepsilon>0$都有$\delta>0$，使$\|f\|<\delta$能推出$\|Rf\|<\varepsilon$。于是
\[
\|f-f_0\|<\delta\quad\Longrightarrow\quad\|R(f-f_0)\|=\|Rf-Rf_0\|<\varepsilon;
\]
也就是说，$R$在$f=f_0$处也连续，因而处处连续。

若$\|Rf\|\leq C\|f\|$，其中当然可取$C>0$，就有连续性：令$\delta=\varepsilon/C$即可。反过来，若连续性成立，就先取与$\varepsilon=1$相应的$\delta$，再令$C=2/\delta$。因为$f=0$时，$\|Rf\|\leq C\|f\|$无论如何都成立；而$f\ne0$，即$\|f\|>0$时，令$g=\frac{\delta/2}{\|f\|}f$，则$\|g\|=\delta/2$，所以
\[
\|Rg\|=\frac{\delta/2}{\|f\|}\|Rf\|<1,
\qquad\|Rf\|<\frac{\|f\|}{\delta/2}=C\|f\|.
\]
由$\|Rf\|\leq C\|f\|$可得$|(Rf,\allowbreak g)|\leq\|Rf\|\|g\|\leq C\|f\|\|g\|$。反过来，在$|(Rf,\allowbreak g)|\leq C\|f\|\|g\|$中取$g=Rf$，得$\|Rf\|^2\leq C\|f\|\|Rf\|$，从而$\|Rf\|\leq C\|f\|$。对于厄米的$R$，还须证明
\[
|(Rf,f)|\leq C\|f\|^2\quad\Longrightarrow\quad |(Rf,g)|\leq C\|f\|\|g\|.
\]
把$f$分别换成$(f+g)/2$和$(f-g)/2$，得到\origNote[到][]{65}
\begin{align*}
|\operatorname{Re}(Rf,g)|
&=\left|\left(R\frac{f+g}{2},\frac{f+g}{2}\right)-\left(R\frac{f-g}{2},\frac{f-g}{2}\right)\right|\\
&\leq C\left(\left\|\frac{f+g}{2}\right\|^2+\left\|\frac{f-g}{2}\right\|^2\right)
=C\frac{\|f\|^2+\|g\|^2}{2}.
\end{align*}
现在仿照\CJKglue \hyperref[thm:II.1]{定理1} 的证明，把$f,\allowbreak g$换成$af,\allowbreak a^{-1}g$（$a>0$），再对右边取最小值，便得$|\operatorname{Re}(Bf,\allowbreak g)|\leq C\|f\|\|g\|$。\tnote[。][然][原]{原书此处印作$B$，前后论证中的算符则为$R$。这里保留原书字母。}然后把$f$换成$e^{i\alpha}f$（$\alpha$为实数），对左边取最大值，就有
\[
|(Rf,g)|\leq C\|f\|\|g\|.
\]
当然，这一证明只在$Rg$有意义时才成立。但这样的$g$处处稠密，而最终结果已不再含有$Rg$，所以利用连续性，结果也就普遍成立。

再证明一个关于定号算符的定理：

\hypertarget{origpage-53}{}\phantomsection\label{origpage-53}\leavevmode\phantomsection\textbf{定理19.}\label{thm:II.19} 若$R$厄米且定号，则恒有
\[
|(Rf,g)|\leq\sqrt{(Rf,f)(Rg,g)}
\]
此时由$(Rf,\allowbreak f)=0$可推出$Rf=0$。

\textkaishu{证明。}由$(Rf,\allowbreak f)\geq0$普遍成立，即由定号性，推出上述不等式的方法，与\CJKglue \hyperref[thm:II.1]{定理1} 中从$(f,\allowbreak f)\geq0$推出Schwarz不等式$|(f,\allowbreak g)|\leq\sqrt{(f,\allowbreak f)(g,\allowbreak g)}$，即$|(f,\allowbreak g)|\leq\|f\|\|g\|$的方法完全相同。若现在$(Rf,\allowbreak f)=0$，由这个不等式又得，只要$Rg$有意义，就有$(Rf,\allowbreak g)=0$。因此，这个关系在处处稠密的$g$的集合上成立，再由连续性便对一切$g$成立。所以$Rf=0$。

结束本节之前，再次提请注意两个算符$R,\allowbreak S$对易这一重要概念，即关系$RS=SR$。

由$RS=SR$有
\[
S\cdots SSR=S\cdots SRS=S\cdots RSS=\cdots=RS\cdots SS;
\]
即$R$与$S^n$对易，$n=1,\allowbreak 2,\allowbreak \ldots$。由于$R1=1R=R$且$S^0=1$，这对$n=0$也成立。若$S^{-1}$存在，则$S^{-1}\cdot SR\cdot S^{-1}=S^{-1}\cdot RS\cdot S^{-1}$。鉴于
\begin{align*}
S^{-1}\cdot SR\cdot S^{-1}&=S^{-1}S\cdot RS^{-1}=RS^{-1},\\
S^{-1}\cdot RS\cdot S^{-1}&=S^{-1}R\cdot SS^{-1}=S^{-1}R,
\end{align*}
便还有$RS^{-1}=S^{-1}R$。所以$n=-1$也允许，进而$n=-2,\allowbreak -3,\allowbreak \ldots$也都允许。即$R$与$S$的一切幂对易。再应用一次这一结果，就知$R$的每个幂都与$S$的每个幂对易。如果$R$与$S,\allowbreak T$对易，显然也与所有$aS$，以及$S\pm T,\allowbreak ST$对易。结合前一结果便知：若$R,\allowbreak S$对易，则$R$的一切多项式都与$S$的一切多项式对易。特别令$R=S$，就知$R$的一切多项式彼此对易。

\section{本征值问题}\label{sec:II.6}
在希尔伯特空间的特殊实现$F_Z,\allowbreak F_\Omega$中，量子力学的主要问题是求解 \hyperref[sec:I.3]{I.3} 的方程$\hyperref[eq:I.3.E1]{\mathbf{E}_1}$、$\hyperref[eq:I.3.E2]{\mathbf{E}_2}$。现在，我们已经可以在抽象希尔伯特空间中考察这个问题了。它称为本征值问题，我们必须重新给出统一的表述。

\hyperref[sec:I.3]{I.3} 中的$\hyperref[eq:I.3.E1]{\mathbf{E}_1}$和$\hyperref[eq:I.3.E2]{\mathbf{E}_2}$，都要求求出
\begin{equation}
\mathsf H\varphi=\lambda\varphi\tag{E}\label{eq:II.6.E}
\end{equation}
的一切非零解$\varphi\ne0$。其中$\mathsf H$是与哈密顿函数相应的厄米算符（参见该节），$\varphi$为希尔伯特空间中的元素，$\lambda$为实数；$\mathsf H$给定，$\varphi,\allowbreak \lambda$待求。

\hypertarget{origpage-54}{}\phantomsection\label{origpage-54}不过，对所求解的个数还有一定要求：所求的解必须足够多，使得

1. 在矩阵理论中，能由这些解
\[
\varphi_1=\{s_{11},s_{12},\ldots\},\quad\varphi_2=\{s_{21},s_{22},\ldots\},\quad\ldots
\]
（此时是在$F_Z$中）组成一个有逆矩阵$\mathsf S^{-1}$的矩阵$\mathsf S=\{s_{\mu\nu}\}$（见 \hyperref[sec:I.3]{I.3}）。

2. 在波动理论中，每个波函数$\varphi(q_1\cdots q_f)$（它本身不必是解）都能按解
\[
\varphi_1=\varphi_1(q_1\cdots q_f),\quad\varphi_2=\varphi_2(q_1\cdots q_f),\quad\ldots
\]
展开为级数；这里$\varphi_1,\allowbreak \varphi_2,\allowbreak \ldots$可以属于不同的$\lambda$：
\[
\varphi(q_1\cdots q_f)=\sum_{n=1}^\infty c_n\varphi_n(q_1\cdots q_f).
\]
虽说 \hyperref[sec:I.3]{I.3} 未曾提到后一点，但这一要求对波动理论的进一步发展，特别对Schr\"odinger的“微扰理论”，是不可缺少的。\origNote[。][]{66}

实际上，第1项与第2项意思相同。因为矩阵$\mathsf S$把$\{1,\allowbreak 0,\allowbreak 0,\allowbreak \ldots\},\allowbreak \{0,\allowbreak 1,\allowbreak 0,\allowbreak \ldots\},\allowbreak \ldots$分别变为
\[
\{s_{11},s_{12},s_{13},\ldots\},\quad\{s_{21},s_{22},s_{23},\ldots\},\quad\ldots,
\]
从而把整个希尔伯特空间$\mathfrak{R}_\infty$变为$\varphi_1,\allowbreak \varphi_2,\allowbreak \ldots$张成的闭线性流形。要使$\mathsf S^{-1}$存在，后一个流形也必须等于$\mathfrak{R}_\infty$。第2项恰好说的是同一件事：它同样要求，每个$\varphi$都能用$\varphi_1,\allowbreak \varphi_2,\allowbreak \ldots$的线性组合任意精确地逼近。\origNote[。][现]{67}现在来明确这个条件究竟意味着什么，并用现有的形式工具，重新考察方程 \hyperref[eq:II.6.E]{E} 的性质。

首先，由于要求$\varphi\ne0$，而且$\varphi$为解时$a\varphi$也为解，所以只考察$\|\varphi\|=1$的解即可。其次，无须另行要求$\lambda$为实数，因为由$\mathsf H\varphi=\lambda\varphi$有
\[
(\mathsf H\varphi,\varphi)=(\lambda\varphi,\varphi)=\lambda(\varphi,\varphi)=\lambda
\]
（见 \hyperref[sec:II.5]{II.5}、\CJKglue \hyperlink{note-61}{原注 \getrefnumber{orig-note-61}}）。第三，属于不同$\lambda_1,\allowbreak \lambda_2$的解$\varphi_1,\allowbreak \varphi_2$彼此正交：
\[
(\mathsf H\varphi_1,\varphi_2)=\lambda_1(\varphi_1,\varphi_2),\qquad
(\mathsf H\varphi_1,\varphi_2)=(\varphi_1,\mathsf H\varphi_2)=\lambda_2(\varphi_1,\varphi_2).
\]
于是，由$\lambda_1(\varphi_1,\allowbreak \varphi_2)=\lambda_2(\varphi_1,\allowbreak \varphi_2)$及$\lambda_1\ne\lambda_2$得到$(\varphi_1,\allowbreak \varphi_2)=0$。

现设$\lambda_1,\allowbreak \lambda_2,\allowbreak \ldots$为使 \hyperref[eq:II.6.E]{E} 可解的所有互不相同的$\lambda$。对于使$\mathsf H\varphi=\lambda\varphi$可解的每个$\lambda$，选一个范数为$1$的解$\varphi_\lambda$。根据刚才的讨论，这些$\varphi_\lambda$构成规范正交系。由 \hyperref[sec:II.2]{II.2} 的\CJKglue \hyperref[thm:II.3infty]{定理$3^{(\infty)}$} 知，这个系或者有限，或者是无限序列；所以这些$\lambda$也可以写成一个序列，并可在有限项后终止。

\hypertarget{origpage-55}{}\phantomsection\label{origpage-55}对每个$\lambda=\lambda_\rho$，$\mathsf H\varphi=\lambda\varphi$的全部解构成线性流形，而且是闭的。\origNote[。][因]{68}因此，依\CJKglue \hyperref[thm:II.9]{定理9}，存在由这些解构成的规范正交系$\varphi_{\rho,\allowbreak 1},\allowbreak \ldots,\allowbreak \varphi_{\rho,\allowbreak \nu_\rho}$，恰好张成这个闭线性流形。数$\nu_\rho$显然是$\lambda=\lambda_\rho$时线性无关解的最大个数，所以也称为本征值$\lambda_\rho$的重数。这里$\nu=1,\allowbreak 2,\allowbreak \ldots,\allowbreak \infty$；$\nu=\infty$确实可能发生，例如$\mathsf H=1,\allowbreak \nu=1$时。\tnote[。][根][原]{原书例子末尾印作$\nu=1$，按此处语义应指$\lambda=1$：恒等算符的本征值$1$具有无限重数。这里保留原书符号。}根据前面的讨论，不同$\rho$所对应的$\varphi_{\rho,\allowbreak 1},\allowbreak \ldots,\allowbreak \varphi_{\rho,\allowbreak \nu_\rho}$彼此也正交，因此所有
\[
\varphi_{\rho,\nu}\qquad(\rho=1,2,\ldots;\ \nu=1,\ldots,\nu_\rho)
\]
合起来仍构成规范正交系。从它的构造可知，它与 \hyperref[eq:II.6.E]{E} 的全部解$\varphi$张成同一个闭线性流形。

把$\varphi_{\rho,\allowbreak \nu}$按任意顺序编号为$\psi_1,\allowbreak \psi_2,\allowbreak \ldots$，把相应的$\lambda_\rho$记为$\lambda^{(1)},\allowbreak \lambda^{(2)},\allowbreak \ldots$。于是，前面提出的条件——\hyperref[eq:II.6.E]{E} 的全部解所张成的闭线性流形应为$\mathfrak{R}_\infty$——也就是说，仅由解中的一部分组成的$\psi_1,\allowbreak \psi_2,\allowbreak \ldots$就必须张成$\mathfrak{R}_\infty$。依\CJKglue \hyperref[thm:II.7]{定理7} 的$\alpha)$，这就是要求这个规范正交系完全。

因此，按照量子力学的要求解决本征值问题，就是求出 \hyperref[eq:II.6.E]{E} 的足够多个解
\[
\varphi=\psi_1,\psi_2,\ldots\qquad\text{且}\qquad\lambda=\lambda_1,\lambda_2,\ldots
\]
使它们能够组成完全规范正交系。但一般情况下，并不能做到这一点。例如在波动理论中可以看到，\hyperref[eq:II.6.E]{E}（即 \hyperref[sec:I.3]{I.3} 中的$\hyperref[eq:I.3.E2]{\mathbf{E}_2}$）的一部分解，其绝对值平方的积分并不有限，\origNote[，][因]{69}因而不属于希尔伯特空间；然而，要像上面所说的那样，把每个波函数都按解展开，又少不了所有这些解。因此，在希尔伯特空间内——而 \hyperref[eq:II.6.E]{E} 只讨论此空间——并不存在由解组成的完全规范正交系。

另一方面，Hilbert的本征值问题理论表明，这种现象对于算符的行为来说绝非例外，哪怕对于连续算符也不例外。\origNote[。][所]{70}所以，我们必须正视出现这种现象时的情形。它在物理上意味着什么，以后再看（见 \hyperref[sec:III.3]{III.3}）。若出现这种现象，即从 \hyperref[eq:II.6.E]{E} 的解中选出的规范正交系不完全，就说存在“$\mathsf H$的连续谱”。$\lambda_1,\allowbreak \lambda_2,\allowbreak \ldots$则构成“$\mathsf H$的点谱”。

\hypertarget{origpage-56}{}\phantomsection\label{origpage-56}方程 \hyperref[eq:II.6.E]{E} 既然不能满足要求，我们的下一项任务，就是为厄米算符的本征值问题找到一种表述，再将其用于量子力学。先沿着Hilbert的思路（参见\CJKglue \hyperlink{note-70}{原注 \getrefnumber{orig-note-70}} 所引文献），给出并说明这种表述。

\section{续论}\label{sec:II.7}
方程
\[
\mathsf H\varphi=\lambda\varphi,
\]
以及能由其解组成完全规范正交系的要求，都是从有限维空间$\mathfrak{R}_n$类比得来的。

在$\mathfrak{R}_n$中，$\mathsf H$是矩阵$\{h_{\mu\nu}\}$，$\mu,\allowbreak \nu=1,\allowbreak \ldots,\allowbreak n$，$h_{\mu\nu}=\overline{h_{\nu\mu}}$。代数中的一个熟知事实是，$\mathsf H\varphi=\lambda\varphi$的解$\varphi=\{x_1,\allowbreak \ldots,\allowbreak x_n\}$，即
\[
\sum_{\nu=1}^n h_{\mu\nu}x_\nu=\lambda x_\mu\qquad(\mu=1,\ldots,n)
\]
的解，包含一个完全规范正交系。\origNote[。][]{71}

我们已经看到，不能靠令$n\to\infty$，把$\mathfrak{R}_n$的这个性质推广到$\mathfrak{R}_\infty$；因此，在$\mathfrak{R}_\infty$中必须另行表述本征值问题。下面将看到，$\mathfrak{R}_n$中的本征值问题可以改写成另一种形式；新形式在$\mathfrak{R}_n$中与旧形式等价，却能够推广到$\mathfrak{R}_\infty$。也就是说，它们在每个$\mathfrak{R}_n$（$n=1,\allowbreak 2,\allowbreak \ldots$）中表达同一件事，即厄米矩阵可以作主轴变换；但其中一种能推广到$\mathfrak{R}_\infty$，另一种却不能。

设$\{x_{11},\allowbreak \ldots,\allowbreak x_{1n}\},\allowbreak \ldots,\allowbreak \{x_{n1},\allowbreak \ldots,\allowbreak x_{nn}\}$是由本征方程的解组成的完全规范正交系，$\lambda_1,\allowbreak \ldots,\allowbreak \lambda_n$为相应的$\lambda$。这些向量$\{x_{11},\allowbreak \ldots,\allowbreak x_{1n}\},\allowbreak \ldots,\allowbreak \{x_{n1},\allowbreak \ldots,\allowbreak x_{nn}\}$便在$\mathfrak{R}_n$中构成一个直角坐标系。从这个坐标系中的坐标$\mathfrak{x}_1,\allowbreak \ldots,\allowbreak \mathfrak{x}_n$到坐标$\xi_1,\allowbreak \ldots,\allowbreak \xi_n$的变换公式为
\[
\{\xi_1,\ldots,\xi_n\}=\mathfrak{x}_1\{x_{11},\ldots,x_{1n}\}+\cdots+\mathfrak{x}_n\{x_{n1},\ldots,x_{nn}\},
\]
即
\[
\xi_1=\sum_{\mu=1}^n x_{\mu1}\mathfrak{x}_\mu,\quad\ldots,\quad
\xi_n=\sum_{\mu=1}^n x_{\mu n}\mathfrak{x}_\mu,
\]
反之，
\[
\mathfrak{x}_1=\sum_{\mu=1}^n\overline{x_{1\mu}}\xi_\mu,\quad\ldots,\quad
\mathfrak{x}_n=\sum_{\mu=1}^n\overline{x_{n\mu}}\xi_\mu.
\]

\hypertarget{origpage-57}{}\phantomsection\label{origpage-57}利用变量$\mathfrak{x}_1,\allowbreak \ldots,\allowbreak \mathfrak{x}_n$，再引入一组变量$\eta_1,\allowbreak \ldots,\allowbreak \eta_n$，以及按上述公式与它们相应的$\mathfrak{y}_1,\allowbreak \ldots,\allowbreak \mathfrak{y}_n$，就能把条件$\sum_{\nu=1}^n h_{\mu\nu}x_{\rho\nu}=\lambda_\rho x_{\rho\mu}$写成
\[
\sum_{\rho,\mu=1}^n\left(\sum_{\nu=1}^n h_{\mu\nu}x_{\rho\nu}\right)\mathfrak{x}_\rho\overline{\eta_\mu}
=\sum_{\rho,\mu=1}^n\lambda_\rho x_{\rho\mu}\mathfrak{x}_\rho\overline{\eta_\mu},
\]
即
\begin{equation}
\sum_{\mu,\nu=1}^n h_{\mu\nu}\xi_\mu\overline{\eta_\nu}
=\sum_{\rho=1}^n\lambda_\rho
\left(\sum_{\mu=1}^n\overline{x_{\rho\mu}}\xi_\mu\right)
\overline{\left(\sum_{\mu=1}^m\overline{x_{\rho\mu}}\eta_\mu\right)}.
\tag{D}\label{eq:II.7.D}
\end{equation}
另一方面，这个坐标系是直角坐标系，由下式表达：
\begin{equation}
\sum_{\mu=1}^n\xi_\mu\overline{\eta_\mu}
=\sum_{\mu=1}^n
\left(\sum_{\rho=1}^n\overline{x_{\rho\mu}}\xi_\mu\right)
\overline{\left(\sum_{\rho=1}^n\overline{x_{\rho\mu}}\eta_\mu\right)}.
\tag{O}\label{eq:II.7.O}
\end{equation}
因此，求出具有性质 \hyperref[eq:II.7.D]{D}、\hyperref[eq:II.7.O]{O} 的矩阵$\{x_{\mu\nu}\}$，就等于解决$\mathfrak{R}_n$中的本征值问题；而先前无法推广到$\mathfrak{R}_\infty$的，恰是这个形式。不过，这种失败并不奇怪，原因如下。条件 \hyperref[eq:II.7.D]{D}、\hyperref[eq:II.7.O]{O} 并不能完全确定未知量$\lambda_\rho,\allowbreak x_{\mu\nu}$。诚然，这类“主轴变换”的理论表明，除排列次序外，$\lambda_\rho$是唯一确定的（参见\CJKglue \hyperlink{note-71}{原注 \getrefnumber{orig-note-71}} 所引文献）；但$x_{\mu\nu}$的情形就远没有这么好了。显然，可以给每一行$x_{\rho1},\allowbreak \ldots,\allowbreak x_{\rho n}$乘上绝对值为$1$的因子$\theta_\rho$；若若干$\lambda_\rho$重合，甚至还允许在相应的行$x_{\rho1},\allowbreak \ldots,\allowbreak x_{\rho n}$之间作任意酉变换！试图带着这些并未唯一确定的量，进行$n\to\infty$这一困难的极限过程，是没有希望的：由于未能完全确定，$\lambda_\rho,\allowbreak x_{\mu\nu}$在过程中可以任意剧烈地波动，这个过程又怎么能收敛呢！

这就提示了正确的处理方向：首先须设法替换条件 \hyperref[eq:II.7.D]{D}、\hyperref[eq:II.7.O]{O} 以及未知量$\lambda_\rho,\allowbreak x_{\mu\nu}$，让替换后的对象具有这里所缺少的唯一性。我们将看到，此时取极限的困难就小得多。

若$l$是一个或多个$\lambda_\rho$所取的值，则
\[
\sum_{\lambda_\rho=l}
\left(\sum_{\mu=1}^n\overline{x_{\rho\mu}}\xi_\mu\right)
\overline{\left(\sum_{\mu=1}^n\overline{x_{\rho\mu}}\eta_\mu\right)}
\]
在上面列举的、与 \hyperref[eq:II.7.D]{D}、\hyperref[eq:II.7.O]{O} 相容的$\lambda_\rho,\allowbreak x_{\mu\nu}$变动下保持不变。若$l$与所有$\lambda_\rho$都不同，这个式子等于零，当然更是不变。因此，对任意$l$，厄米型
\[
E(l;\xi,\eta)=\sum_{\lambda_\rho\leq l}
\left(\sum_{\rho=1}^n\overline{x_{\rho\mu}}\xi_\mu\right)
\overline{\left(\sum_{\rho=1}^n\overline{x_{\rho\mu}}\eta_\mu\right)},
\]
也保持不变；其中$\xi,\allowbreak \eta$分别代表$\xi_1,\allowbreak \ldots,\allowbreak \xi_n$和$\eta_1,\allowbreak \ldots,\allowbreak \eta_n$。\tnote[。][][原]{原书在 \hyperref[eq:II.7.D]{D} 的最后一个求和上限处印作$m$，并在 \hyperref[eq:II.7.O]{O} 及$E(l;\xi,\allowbreak \eta)$的定义式中混用了求和指标$\mu,\allowbreak \rho$。这里保留原书指标；其一致的含义可由前面的坐标变换确定。}

\hypertarget{origpage-58}{}\phantomsection\label{origpage-58}只要知道$E(l;\xi,\allowbreak \eta)$，即知道它们的系数，就很容易反求$\lambda_\rho,\allowbreak x_{\mu\nu}$。因此，如果改写本征值问题，即 \hyperref[eq:II.7.D]{D}、\hyperref[eq:II.7.O]{O}，让其中不再出现$\lambda_\rho,\allowbreak x_{\mu\nu}$，而只出现$E(l;\xi,\allowbreak \eta)$，就达到了所要求的唯一表述。

设$E(l)$为厄米型$E(l;\xi,\allowbreak \eta)$的矩阵。\origNote[。][那]{72}那么，\hyperref[eq:II.7.D]{D}、\hyperref[eq:II.7.O]{O} 对矩阵族$E(l)$意味着什么呢？

\hyperref[eq:II.7.O]{O} 意味着：当$l$充分大，即大于一切$\lambda_\rho$时，$E(l)=1$，即单位矩阵。又由$E(l)$的构造可知，当$l$充分小，即小于一切$\lambda_\rho$时，$E(l)=0$；而当$l$从$-\infty$增到$+\infty$时，除有限多个点外，$E(l)$总保持常值，只在那些点发生跳跃。这些点就是$\lambda_1,\allowbreak \ldots,\allowbreak \lambda_n$中互不相同的值，记作$l_1<l_2<\cdots<l_m$，$m\leq n$。每一处间断都发生在跳跃点的左侧，因为$\sum_{\lambda_\rho\leq l}$作为$l$的函数是右连续的；若改用$\sum_{\lambda_\rho<l}$，情形就相反。最后，我们将证明，对$l'\leq l''$有
\[
E(l')E(l'')=E(l'')E(l')=E(l'),
\]
这里是矩阵乘积。

在$\mathfrak{x}_1,\allowbreak \ldots,\allowbreak \mathfrak{x}_n$与$\mathfrak{y}_1,\allowbreak \ldots,\allowbreak \mathfrak{y}_n$的坐标系中，对$E(l';\xi,\allowbreak \eta),\allowbreak E(l'';\xi,\allowbreak \eta)$来证明更方便。引入这些变量以后，两者分别化为
\[
\sum_{\lambda_\rho\leq l'}\mathfrak{x}_\rho\overline{\mathfrak{y}_\rho}
\qquad\text{且}\qquad
\sum_{\lambda_\rho\leq l''}\mathfrak{x}_\rho\overline{\mathfrak{y}_\rho},
\]
因此，其矩阵在对角线以外全为零；在第$\rho$个对角元处，分别当$\lambda_\rho\leq l'$或$\lambda_\rho\leq l''$时为$1$，否则也为零。对这些矩阵，上述断言显然成立。

再来改写 \hyperref[eq:II.7.D]{D}。它显然是说
\[
\sum_{\mu,\nu=1}^n h_{\mu\nu}\xi_\mu\overline{\eta_\nu}
=\sum_{\tau=1}^m l_\tau\{E(l_\tau;\xi,\eta)-E(l_{\tau-1};\xi,\eta)\},
\]
其中$l_0$是任意小于$l_1$的数。但$E(l;\xi,\allowbreak \eta)$在下列每个区间上保持常值：
\[
-\infty<l<l_1;\quad l_1\leq l<l_2;\quad\ldots;\quad l_{m-1}\leq l<l_m;\quad l_m\leq l<+\infty.
\]
因此，对于任意数列$\Lambda_0<\Lambda_1<\Lambda_2<\cdots<\Lambda_k$，只要$l_1,\allowbreak \ldots,\allowbreak l_m$包含在$\Lambda_1,\allowbreak \ldots,\allowbreak \Lambda_k$中，就有
\[
\sum_{\mu,\nu=1}^n h_{\mu\nu}\xi_\mu\overline{\eta_\nu}
=\sum_{\tau=1}^k\Lambda_\tau\{E(\Lambda_\tau;\xi,\eta)-E(\Lambda_{\tau-1};\xi,\eta)\}.
\]

\hypertarget{origpage-59}{}\phantomsection\label{origpage-59}利用Stieltjes积分的概念，\origNote[，][也]{73}也可写成
\[
\sum_{\mu,\nu=1}^n h_{\mu\nu}\xi_\mu\overline{\eta_\nu}
=\int_{-\infty}^{+\infty}\lambda\,dE(\lambda;\xi,\eta).
\]
显然，$\int_{-\infty}^{+\infty}$可换成任意满足$a<l_1,\allowbreak b>l_m$的$\int_a^b$。或者，考察系数，并把对所有系数成立的等式写成矩阵本身的等式：
\[
\mathsf H=\int_{-\infty}^{+\infty}\lambda\,dE(\lambda),\qquad \mathsf H=\{h_{\mu\nu}\}.
\]
于是，到目前为止得到的问题如下：给定厄米矩阵$\mathsf H=\{h_{\mu\nu}\}$，求一族厄米矩阵$E(\lambda)$（$-\infty<\lambda<+\infty$），使它们具有下列性质：

\leavevmode\phantomsection\textbf{$S_1$.}\label{cond:II.S1} 对充分小的$\lambda$，$E(\lambda)=0$；对充分大的$\lambda$，$E(\lambda)=1$。$E(\lambda)$作为$\lambda$的函数，除有限多个点外处处为常值；在这些点发生跳跃，且跳跃总发生在该点左侧。

\leavevmode\phantomsection\textbf{$S_2$.}\label{cond:II.S2} 恒有$E(\lambda')E(\lambda'')=E(\operatorname{Min}(\lambda',\allowbreak \lambda''))$。\origNote[。][]{74}

\leavevmode\phantomsection\textbf{$S_3$.}\label{cond:II.S3} 利用Stieltjes积分，有
\[
\mathsf H=\int_{-\infty}^{+\infty}\lambda\,dE(\lambda).
\]
我们不再花时间从$\hyperref[cond:II.S1]{S_1}$—$\hyperref[cond:II.S3]{S_3}$反过来求解 \hyperref[eq:II.7.D]{D}、\hyperref[eq:II.7.O]{O}（虽然这很容易），因为量子力学中所需的正是本征值问题的这种形式。下面立即把$\hyperref[cond:II.S1]{S_1}$—$\hyperref[cond:II.S3]{S_3}$从有限多个变量推广到无限多个变量，即从$\mathfrak{R}_n$推广到$\mathfrak{R}_\infty$。

在$\mathfrak{R}_\infty$中，显然应把$\mathsf H$和$E(\lambda)$理解为厄米算符。于是，对于给定的$\mathsf H$，要仿照$\hyperref[cond:II.S1]{S_1}$—$\hyperref[cond:II.S3]{S_3}$，求与它具有某种关系的算符族$E(\lambda)$。因此，须找出$\hyperref[cond:II.S1]{S_1}$—$\hyperref[cond:II.S3]{S_3}$在$\mathfrak{R}_\infty$中的对应形式。

$\hyperref[cond:II.S2]{S_2}$可原样保留，因为$\mathfrak{R}_n$的维数在这里不起作用。不过，利用关于投影算符的结论（\hyperref[sec:II.4]{II.4}），可以稍加改写。首先，取$\lambda'=\lambda''=\lambda$，它便是说$E(\lambda)^2=E(\lambda)$，所以各个$E(\lambda)$必须是投影算符。然后，$\hyperref[cond:II.S2]{S_2}$就意味着：由$\lambda'\leq\lambda''$可得$E(\lambda')\leq E(\lambda'')$。只讨论$\lambda'\leq\lambda''$就够了，$\lambda'\geq\lambda''$的相应结论随之成立（见 \hyperref[sec:II.4]{II.4} 的\CJKglue \hyperref[thm:II.14]{定理14} 及其后的文字）。

\hypertarget{origpage-60}{}\phantomsection\label{origpage-60}对于$\hyperref[cond:II.S3]{S_3}$，则须谨慎：$A=\int_{-\infty}^{\infty}\lambda\,dE(\lambda)$没有意义，因为Stieltjes积分是对数定义的，并不是对算符定义的。不过，很容易用数代替$\mathsf H,\allowbreak E(\lambda)$，并使所得关系准确表达所要的算符关系。我们要求
\[
(\mathsf Hf,g)=\int_{-\infty}^{+\infty}\lambda\,d(E(\lambda)f,g)
\]
对$\mathfrak{R}_\infty$中的一切$f,\allowbreak g$成立，只要$\mathsf Hf$有意义；而$\mathsf H=\int_{-\infty}^{+\infty}\lambda\,dE(\lambda)$仅应理解为这一关系的符号化简写。

最后，当维数成为无限时，$\hyperref[cond:II.S1]{S_1}$会受到实质影响。在$\mathfrak{R}_n$中，使$E(\lambda)$从此等于$0$或$1$的界点，以及使它发生跳跃间断的点，都是$\mathsf H$的本征值；而它保持常值的各个区间，则是不含本征值的区间。当$n\to\infty$时，可能发生各种情况：最小、最大的本征值可能分别趋于$-\infty,\allowbreak +\infty$；其他本征值的个数越来越多，可能排列得越来越密，以至于常值区间缩成点。特别在Hilbert理论中，后一现象正是可能预示所谓连续谱出现的征兆。\origNote[。][因]{75}因此，从$\mathfrak{R}_n$过渡到$\mathfrak{R}_\infty$时，必须对$\hyperref[cond:II.S1]{S_1}$作相当大的修改：变化原有的离散、逐级跳跃的性质可能消失，这一点必须顾及。

从这个角度看，很自然可以放弃$E(\lambda)$最终实际达到$0,\allowbreak 1$的要求，改为只要求它在$\lambda\to-\infty$和$\lambda\to+\infty$时分别收敛于$0$和$1$。同样，也不再只允许区间上的常值和点处的跳跃，而允许连续增长。至于较不苛刻的要求——若有间断，只能在该点左侧发生——则不妨暂且保留。因此，把$\hyperref[cond:II.S1]{S_1}$表述如下：当$\lambda\to-\infty$时，$E(\lambda)\to0$；当$\lambda\to+\infty$时，$E(\lambda)\to1$；当$\lambda\to\lambda_0,\allowbreak \lambda\geq\lambda_0$时，$E(\lambda)\to E(\lambda_0)$。\origNote[。][]{76}

关于$\hyperref[cond:II.S3]{S_3}$还须补充一点。在有限维$\mathfrak{R}_n$中，曾有$A=\sum_{\tau=1}^ml_\tau F_\tau$，其中$F_\tau$表示矩阵$E(l_\tau)-E(l_{\tau-1})$。根据$\hyperref[cond:II.S1]{S_1}$有
\begin{align*}
\text{当}\mathwordspace{}\sigma\geq\tau:\quad
F_\tau E(l_\sigma)&=E(l_\tau)E(l_\sigma)-E(l_{\tau-1})E(l_\sigma)\\
&=E(l_\tau)-E(l_{\tau-1})=F_\tau,\\
\text{当}\mathwordspace{}\sigma\leq\tau-1:\quad
F_\tau E(l_\sigma)&=E(l_\tau)E(l_\sigma)-E(l_{\tau-1})E(l_\sigma)\\
&=E(l_\sigma)-E(l_\sigma)=0.
\end{align*}

\hypertarget{origpage-61}{}\phantomsection\label{origpage-61}于是，由$F_\sigma=E(l_\sigma)-E(l_{\sigma-1})$得
\[
F_\tau F_\sigma=\begin{cases}F_\tau,&\tau=\sigma,\\0,&\tau\ne\sigma.\end{cases}
\]
因此
\[
\mathsf H^2=\left(\sum_{\tau=1}^m l_\tau F_\tau\right)^2
=\sum_{\tau,\sigma=1}^m l_\tau l_\sigma F_\tau F_\sigma
=\sum_{\tau=1}^m l_\tau^2F_\tau.
\]
同样，$\mathsf H^p=\sum_{\tau=1}^ml_\tau^pF_\tau$。于是施行与$\mathsf H$本身相同的变换，得到
\[
\mathsf H^2=\int_{-\infty}^{+\infty}\lambda^2\,dE(\lambda).
\]
所以在$\mathfrak{R}_\infty$中，也可猜想相同形式的符号等式成立；写成数值关系就是
\[
(\mathsf H^2f,g)=\int_{-\infty}^{+\infty}\lambda^2\,d(E(\lambda)f,g).
\]
下文将证实这一点。当$f=g$时，由于
\[
(\mathsf H^2f,f)=(\mathsf Hf,\mathsf Hf)=\|\mathsf Hf\|^2,\qquad(E(\lambda)f,f)=\|E(\lambda)f\|^2,
\]
得到
\[
\|\mathsf Hf\|^2=\int_{-\infty}^{+\infty}\lambda^2\,d(\|E(\lambda)f\|^2).
\]
这个公式使我们预期：$E(\lambda)$不仅确定$\mathsf Hf$有意义时的值，也确定它何时有意义。因为积分$\int_{-\infty}^{+\infty}\lambda^2\,d(\|E(\lambda)f\|^2)$的被积函数从不为负，即$\lambda^2\geq0$；$d$之后的表达式$\|E(\lambda)f\|^2$又不递减（见$\hyperref[cond:II.S2]{S_2}$及 \hyperref[sec:II.4]{II.4} 的\CJKglue \hyperref[thm:II.15]{定理15}）。所以这个积分按其性质，或者收敛于零或有限正值，或者发散到$+\infty$。\origNote[。][这]{77}这一点与它同$\mathsf H$的关系无关，也不取决于$\mathsf Hf$是否有意义。因此，我们可以预期，$\mathsf Hf$有意义，即存在于$\mathfrak{R}_\infty$中，当且仅当拟用作$\|\mathsf Hf\|^2$的值的表达式$\int_{-\infty}^{+\infty}\lambda^2\,d\|E(\lambda)f\|^2$有限；这个表达式对每个$f$都有意义。

因此，$\hyperref[cond:II.S1]{S_1}$—$\hyperref[cond:II.S3]{S_3}$的新表述如下。对给定的厄米算符$\mathsf H$，求投影算符族$E(\lambda)$（$-\infty<\lambda<+\infty$），使它具有下列性质：

\leavevmode\phantomsection\textbf{$\overline{S}_1$.}\label{cond:II.Sbar1} 对每个$f$，当$\lambda\to-\infty$或$\lambda\to+\infty$时，$E(\lambda)f$分别趋于$0$或$f$；当$\lambda\to\lambda_0$且$\lambda\geq\lambda_0$时，$E(\lambda)f\to E(\lambda_0)f$。

\leavevmode\phantomsection\textbf{$\overline{S}_2$.}\label{cond:II.Sbar2} 由$\lambda'\leq\lambda''$可得$E(\lambda')\leq E(\lambda'')$。

\hypertarget{origpage-62}{}\phantomsection\label{origpage-62}\leavevmode\phantomsection\textbf{$\overline{S}_3$.}\label{cond:II.Sbar3} 积分
\[
\int_{-\infty}^{+\infty}\lambda^2\,d(\|E(\lambda)f\|^2),
\]
按其性质，或者收敛于零或有限正值，或者发散到$+\infty$。它确定$\mathsf H$的定义域：$\mathsf Hf$有意义，当且仅当前一种情况成立。此时，对每个$g$都有
\[
(\mathsf Hf,g)=\int_{-\infty}^{+\infty}\lambda\,d(E(\lambda)f,g).
\]
只要前一个积分有限，后一个积分就绝对收敛。\origNote[。][]{78}

性质$\hyperref[cond:II.Sbar1]{\overline{S}_1}$、$\hyperref[cond:II.Sbar2]{\overline{S}_2}$完全没有涉及$\mathsf H$本身。具有这两个性质的投影算符族$E(\lambda)$，称为单位分解。若一个单位分解与$\mathsf H$具有$\hyperref[cond:II.Sbar3]{\overline{S}_3}$的关系，就称它属于$\mathsf H$。

于是，$\mathfrak{R}_\infty$中的本征值问题就是：对于每个给定的厄米算符$\mathsf H$，是否总有属于它的单位分解？有多少个？我们希望的回答是：总是恰有一个。此外，还须考察：这里对本征值问题的定义，与量子力学中，特别是波动理论中，通常用来确定厄米算符本征值的方法，究竟有何关系。

\section{关于本征值问题的初步考察}\label{sec:II.8}
我们对本征值问题作出定义后，首先出现如下问题：条件$\hyperref[cond:II.Sbar1]{\overline{S}_1}$—$\hyperref[cond:II.Sbar3]{\overline{S}_3}$，在形式上与上一节开头的问题完全不同，已经看不出它们还有什么关系。诚然，在$\mathfrak{R}_n$中，我们从原来的条件导出了$\hyperref[cond:II.S1]{S_1}$—$\hyperref[cond:II.S3]{S_3}$；但在$\mathfrak{R}_\infty$中，情形有了本质变化：两种表述不再等价，尽管如前所说，它们在$\mathfrak{R}_n$中是等价的。因此，整个问题又须重新考察：新表述在多大程度上包含旧表述？也就是说，新的$E(\lambda)$何时、又如何确定先前的$\lambda_1,\allowbreak \lambda_2,\allowbreak \ldots$和$\varphi_1,\allowbreak \varphi_2,\allowbreak \ldots$？

设单位分解$E(\lambda)$属于厄米算符$A$。方程
\[
A\varphi=\lambda_0\varphi
\]
何时可解？$A\varphi=\lambda_0\varphi$等价于对一切$g$有$(A\varphi,\allowbreak g)-\lambda_0(\varphi,\allowbreak g)=0$，即

\hypertarget{origpage-63}{}\phantomsection\label{origpage-63}\begin{align*}
0&=\int_{-\infty}^{\infty}\lambda\,d(E(\lambda)f,g)-\lambda_0(\varphi,g)\\
&=\int_{-\infty}^{\infty}\lambda\,d(E(\lambda)f,g)-\lambda_0\int_{-\infty}^{\infty}d(E(\lambda)f,g)\\
&=\int_{-\infty}^{\infty}(\lambda-\lambda_0)\,d(E(\lambda)f,g).
\end{align*}
先取$g=E(\lambda_0)f$，便有
\begin{align*}
0&=\int_{-\infty}^{\infty}(\lambda-\lambda_0)\,d(E(\lambda)f,E(\lambda_0)f)\\
&=\int_{-\infty}^{\infty}(\lambda-\lambda_0)\,d(E(\lambda_0)E(\lambda)f,f)\\
&=\int_{-\infty}^{\infty}(\lambda-\lambda_0)\,d(E(\operatorname{Min}(\lambda,\lambda_0))f,f)\\
&=\int_{-\infty}^{\infty}(\lambda-\lambda_0)\,d(\|E(\operatorname{Min}(\lambda,\lambda_0))f\|^2).
\end{align*}
可把$\int_{-\infty}^{\infty}$分成$\int_{-\infty}^{\lambda_0}+\int_{\lambda_0}^{\infty}$。第一个积分中，$\operatorname{Min}(\lambda,\allowbreak \lambda_0)$可换成$\lambda$；第二个积分中则可换成$\lambda_0$。因此，后一个积分中$d$之后的表达式为常数，所以积分为零。第一个积分则留下\tnote[下][][原]{原书在下一积分中印出的上下限为$\lambda_0,\allowbreak \infty$，而前面的拆分要求$-\infty,\allowbreak \lambda_0$。这里保留原书积分限。}
\[
\int_{\lambda_0}^{\infty}(\lambda-\lambda_0)\,d(\|E(\lambda)f\|^2)=0.
\]
再取$g=f$，得到
\[
0=\int_{-\infty}^{\infty}(\lambda-\lambda_0)\,d(E(\lambda)f,f)
=\int_{-\infty}^{\infty}(\lambda-\lambda_0)\,d(\|E(\lambda)f\|^2).
\]
从此式减去第一个等式，并把被积函数的符号反过来，就得
\[
\int_{-\infty}^{\lambda_0}(\lambda_0-\lambda)\,d(\|E(\lambda)f\|^2)=0.
\]
现在更仔细地考察$\int_{-\infty}^{\lambda_0}(\lambda_0-\lambda)\,d(\|E(\lambda)f\|^2)$及$\int_{\lambda_0}^{\infty}(\lambda-\lambda_0)\,d(\|E(\lambda)f\|^2)$。两者的被积函数都非负，$d$之后又都是$\lambda$的单调不减函数。因此，对每个$\varepsilon>0$有
\begin{align*}
\int_{-\infty}^{\lambda_0}(\lambda_0-\lambda)\,d(\|E(\lambda)f\|^2)
&\geq\int_{-\infty}^{\lambda_0-\varepsilon}(\lambda_0-\lambda)\,d(\|E(\lambda)f\|^2)\\
&\geq\int_{-\infty}^{\lambda_0-\varepsilon}\varepsilon\,d(\|E(\lambda)f\|^2)\\
&=\varepsilon\|E(\lambda_0-\varepsilon)f\|^2,
\end{align*}
以及
\begin{align*}
\int_{\lambda_0}^{\infty}(\lambda-\lambda_0)\,d(\|E(\lambda)f\|^2)
&\geq\int_{\lambda_0+\varepsilon}^{\infty}(\lambda-\lambda_0)\,d(\|E(\lambda)f\|^2)\\
&\geq\int_{\lambda_0+\varepsilon}^{\infty}\varepsilon\,d(\|E(\lambda)f\|^2)\\
&=\varepsilon(\|f\|^2-\|E(\lambda_0+\varepsilon)f\|^2)\\
&=\varepsilon\|f-E(\lambda_0+\varepsilon)f\|^2.
\end{align*}

\hypertarget{origpage-64}{}\phantomsection\label{origpage-64}所以右端均不大于零，但它们又恒不小于零，故必为零。于是
\[
E(\lambda_0-\varepsilon)f=0,\qquad E(\lambda_0+\varepsilon)f=f.
\]
由于$E(\lambda)$右连续，可在第二式中令$\varepsilon\to0$，得到$E(\lambda_0)f=f$。当$\lambda\geq\lambda_0$时，再在第二式中取$\varepsilon=\lambda-\lambda_0\geq0$，即得$E(\lambda)f=f$；当$\lambda<\lambda_0$时，在第一式中取$\varepsilon=\lambda_0-\lambda>0$，便得$E(\lambda)f=0$。所以
\[
E(\lambda)f=\begin{cases}f,&\lambda\geq\lambda_0,\\0,&\lambda<\lambda_0.\end{cases}
\]
这个必要条件也是充分的，因为由它有
\[
(Af,g)=\int_{-\infty}^{\infty}\lambda\,d(E(\lambda)f,g)=\lambda_0(f,g),
\]
只须回顾\CJKglue \hyperlink{note-73}{原注 \getrefnumber{orig-note-73}} 中Stieltjes积分的定义即可明白。因此，对一切$g$有$(Af-\lambda_0f,\allowbreak g)=0$，即$Af=\lambda_0f$。

这个条件应如何解释？首先，它要求$E(\lambda)$在$\lambda=\lambda_0$处不连续。根据 \hyperref[sec:II.4]{II.4} 的\CJKglue \hyperref[thm:II.17]{定理17}，当$\lambda\to\lambda_0,\allowbreak \lambda<\lambda_0$以及$\lambda\to\lambda_0,\allowbreak \lambda>\lambda_0$时，$E(\lambda)$都必趋于某个投影算符；分别记为$E^{(1)}(\lambda_0),\allowbreak E^{(2)}(\lambda_0)$。\origNote[。][由]{79}由$\hyperref[cond:II.S1]{S_1}$，$E^{(2)}(\lambda_0)=E(\lambda_0)$；但在不连续时，$E^{(1)}(\lambda_0)\ne E(\lambda_0)$。又由$\hyperref[cond:II.S2]{S_2}$，当$\lambda<\lambda_0$时，$E(\lambda)\leq E(\lambda_0)$，所以恒有$E^{(1)}(\lambda_0)\leq E(\lambda_0)$。因此，$E(\lambda_0)-E^{(1)}(\lambda_0)$为投影算符；不连续的特征，就是它不等于零。

对所有$\lambda<\lambda_0$都有$E(\lambda)f=0$，能推出$E^{(1)}(\lambda_0)f=0$；反过来，由于$E(\lambda)\leq E^{(1)}(\lambda_0)$，也能推出前一条件。对所有$\lambda\geq\lambda_0$都有$E(\lambda)f=f$，则可由$E(\lambda_0)f=f$推出：因为$E(\lambda_0)\leq E(\lambda)$，$E(\lambda)E(\lambda_0)=E(\lambda_0)$，所以$E(\lambda)f=E(\lambda)E(\lambda_0)f=E(\lambda_0)f=f$。因此，$Af=\lambda_0f$的特征条件是$E^{(1)}(\lambda_0)f=0$、$E(\lambda_0)f=f$；或根据 \hyperref[sec:II.4]{II.4} 的\CJKglue \hyperref[thm:II.14]{定理14}，是$(E(\lambda_0)-E^{(1)}(\lambda_0))f=f$。若令$E(\lambda_0)-E^{(1)}(\lambda_0)=P_{\mathfrak{M}_{\lambda_0}}$，这个条件就是说$f$属于$\mathfrak{M}_{\lambda_0}$。

因此已经证明：$Af=\lambda f$只有在$E(\lambda)$的间断点处才有非零解$f$；而这些解构成刚才定义的闭线性流形$\mathfrak{M}_{\lambda_0}$。

所以，\hyperref[sec:II.6]{II.6} 所求的、由解组成的完全规范正交系（各解可属于不同的$\lambda$）存在，当且仅当全部$\mathfrak{M}_{\lambda_0}$（$-\infty<\lambda_0<\infty$）合起来所张成的闭线性流形为$\mathfrak{R}_\infty$。这种情况下如何构造所需的系，已在 \hyperref[sec:II.6]{II.6} 讨论。还可以重新看出各个$\mathfrak{M}_{\lambda_0}$彼此正交：由$\lambda_0<\mu_0$有
\[
P_{\mathfrak{M}_{\lambda_0}}P_{\mathfrak{M}_{\mu_0}}
=(E(\lambda_0)-E^{(1)}(\lambda_0))(E(\mu_0)-E^{(1)}(\mu_0))=0,
\]
因为
\begin{align*}
E(\lambda_0)-E^{(1)}(\lambda_0)&\leq E(\lambda_0)\leq E^{(1)}(\mu_0),\\
E(\mu_0)-E^{(1)}(\mu_0)&\leq1-E^{(1)}(\mu_0).
\end{align*}

\hypertarget{origpage-65}{}\phantomsection\label{origpage-65}暂不求这种情形发生的精确条件，至少能确定如下事实：若存在一个以$\mu_1,\allowbreak \mu_2$为端点的区间，使$E(\lambda)$在其中连续增长，即$\mu_1<\mu_2$，$E(\lambda)$在$\mu_1\leq\lambda\leq\mu_2$上连续，且$E(\mu_1)\ne E(\mu_2)$，那么上述完全正交系一定不存在。因为当$\lambda\leq\mu_1$时，$E(\lambda)-E^{(1)}(\lambda)\leq E(\lambda)\leq E(\mu_1)$；当$\mu_1<\lambda\leq\mu_2$时，由连续性有$E(\lambda)-E^{(1)}(\lambda)=0$；当$\mu_2<\lambda$时，$E(\lambda)-E^{(1)}(\lambda)\leq1-E^{(1)}(\lambda)\leq1-E(\mu_2)$。所以$E(\lambda)-E^{(1)}(\lambda)$总与$E(\mu_2)-E(\mu_1)$正交。设$E(\mu_2)-E(\mu_1)=P_{\mathfrak{N}}$，则所有$\mathfrak{M}_\lambda$都与$\mathfrak{N}$正交。若能从它们中选出完全规范正交系，$\mathfrak{N}$就只能含有$0$，即$E(\mu_2)-E(\mu_1)=0$，与假定矛盾。

把$E(\lambda)$的间断点称为$A$的点谱；它们就是使$Af=\lambda f,\allowbreak f\ne0$可解的那些$\lambda$。从每个非零的$\mathfrak{M}_\lambda$中选一个$\|f\|=1$的$f$，由各个$\mathfrak{M}_\lambda$的正交性，便得到规范正交系。依 \hyperref[sec:II.2]{II.2} 的\CJKglue \hyperref[thm:II.3infty]{定理3}，它是有限的，或能排成一个序列；所以点谱中的$\lambda$也至多可数。

所有在其任意邻域内$E(\lambda)$都不恒定的点，构成$A$的谱。我们已经看到，若有某些区间含有$A$的谱点，却不含点谱的点，即$E(\lambda)$在这些区间连续而不恒定，那么本征值问题按 \hyperref[sec:II.6]{II.6} 开头的意义一定不可解。对于这种不可解性的精确条件，不再进一步研究；即使每个含有谱点的区间都含有点谱的点，也可能不可解。此时，把点谱与其余部分分离就要费力得多，与本书主题也不够密切。读者可以在Hilbert所引论文中找到这些研究。

不过，我们要说明：若能从$A\varphi=\lambda\varphi$的解中选出完全规范正交系$\varphi_1,\allowbreak \varphi_2,\allowbreak \ldots$，相应的本征值为$\lambda_1,\allowbreak \lambda_2,\allowbreak \ldots$，即处于我们所谓的纯点谱情形，应如何构造$E(\lambda)$。有\origNote[有][]{80}
\[
E(\lambda)=\sum_{\lambda_\rho\leq\lambda}P_{[\varphi_\rho]}.
\]
这个和可以没有加项，此时$E(\lambda)=0$；也可以只有有限个加项，其含义显然；还可以有无限个加项，这时按 \hyperref[sec:II.4]{II.4} 末尾的讨论，它收敛。

事实上，$\hyperref[cond:II.Sbar2]{\overline{S}_2}$显然成立。因为当$\lambda'\leq\lambda''$时，
\[
E(\lambda'')-E(\lambda')=\sum_{\lambda'<\lambda_\rho\leq\lambda''}P_{[\varphi_\rho]}
\]
是投影算符，所以由\CJKglue \hyperref[thm:II.14]{定理14} 得$E(\lambda')\leq E(\lambda'')$。

\hypertarget{origpage-66}{}\phantomsection\label{origpage-66}$\hyperref[cond:II.Sbar1]{\overline{S}_1}$证明如下：对每个$f$有\origNote[有][]{81}
\[
\sum_\rho\|P_{[\varphi_\rho]}f\|^2=\sum_\rho|(f,\varphi_\rho)|^2=\|f\|^2
\]
（\hyperref[thm:II.7]{定理7}），即$\sum_\rho\|P_{[\varphi_\rho]}f\|^2$收敛。因此，对每个$\varepsilon>0$都能指定有限多个$\rho$，使仅对这些$\rho$求的和就已大于$\|f\|^2-\varepsilon$；于是，凡是不含这些$\rho$的和$\sum_\rho'$，都小于$\varepsilon$。这时还有
\[
\left\|\sum_p'P_{[\varphi_\rho]}f\right\|^2=\sum_p'\|P_{[\varphi_\rho]}f\|^2<\varepsilon.
\]
由此特别得到
\[
\left\|\sum_{\lambda_\rho\leq\lambda}P_{[\varphi_\rho]}f\right\|^2<\varepsilon,\quad
\left\|\sum_{\lambda_\rho>\lambda}P_{[\varphi_\rho]}f\right\|^2<\varepsilon,\quad
\left\|\sum_{\lambda_0<\lambda_\rho\leq\lambda}P_{[\varphi_\rho]}f\right\|^2<\varepsilon,
\]
只要分别把$\lambda$取得足够小、足够大，或足够接近$\lambda_0$且满足$\lambda\geq\lambda$。因此确实有\origNote[有][]{82}
\begin{alignat*}{2}
E(\lambda)f&=\sum_{\lambda_\rho\leq\lambda}P_{[\varphi_\rho]}f\longrightarrow0&\quad &\text{当}\mathwordspace{}\lambda\to-\infty,\\
f-E(\lambda)f&=\sum_{\lambda_\rho>\lambda}P_{[\varphi_\rho]}f\longrightarrow0&\quad &\text{当}\mathwordspace{}\lambda\to+\infty,\\
E(\lambda)f-E(\lambda_0)f&=\sum_{\lambda_0<\lambda_\rho\leq\lambda}P_{[\varphi_\rho]}f\longrightarrow0&\quad &\text{当}\mathwordspace{}\lambda\to\lambda_0,\ \lambda\geq\lambda_0.
\end{alignat*}
所以$\hyperref[cond:II.Sbar1]{\overline{S}_1}$成立。

为验证$\hyperref[cond:II.Sbar3]{\overline{S}_3}$，令$f=x_1\varphi_1+x_2\varphi_2+\cdots$，则$Af=\lambda_1x_1\varphi_1+\lambda_2x_2\varphi_2+\cdots$。要使$Af$有意义，$\sum_{\rho=1}^\infty\lambda_\rho^2|x_\rho|^2$就须有限。而\origNote[而][]{83}
\[
\begin{aligned}
\int_{-\infty}^{\infty}\lambda^2\,d\|E(\lambda)f\|^2
&=\int_{-\infty}^{\infty}\lambda^2\,d\left(\sum_{\lambda_\rho\leq\lambda}|x_\rho|^2\right)
=\sum_{\rho=1}^\infty\lambda_\rho^2|x_\rho|^2,
\end{aligned}
\]
又有\origNote[有][]{83}
\[
\begin{aligned}
\int_{-\infty}^{\infty}\lambda^2\,d(E(\lambda)f,g)
&=\int_{-\infty}^{\infty}\lambda\,d\left(\sum_{\lambda_\rho\leq\lambda}x_\rho\overline{v_\rho}\right)
=\sum_{\rho=1}^\infty\lambda_\rho x_\rho\overline{y_\rho}=(Af,g).
\end{aligned}
\]
所以$\hyperref[cond:II.Sbar3]{\overline{S}_3}$也成立。\tnote[。][][原]{原书在第二个积分恒等式左边印作$\lambda^2$，右边却要求$\lambda$；又对同一组系数先写$v_\rho$，后写$y_\rho$。上文还把本应为$\lambda\geq\lambda_0$的条件印作$\lambda\geq\lambda$。均按原书保留。}

再考察两个纯连续谱的例子，即根本不存在点谱的情形。设$\mathfrak{R}_\infty$由使
\[
\int_{-\infty}^{\infty}\cdots\int_{-\infty}^{\infty}|f(q_1\cdots q_l)|^2\,dq_1\cdots dq_l
\]
有限的一切函数$f(q_1\cdots q_l)$构成；$A$取为算符$q_j\cdots$，它显然是厄米的。

可见，$Af=\lambda f$意味着$(q_j-\lambda)f(q_1,\allowbreak \ldots,\allowbreak q_l)=0$；即$f(q_1\cdots q_l)$处处为零，至多可能在$(l-1)$维平面$q_j=0$上例外。\tnote[。][然][原]{原书把可能例外的平面印作$q_j=0$；由前面的方程应得$q_j=\lambda$，下一段也是如此表述。这里保留原书的零。}然而，这个平面的Lebesgue测度，即体积，等于零；依 \hyperref[sec:II.3]{II.3} 中关于条件 \hyperref[ax:II.B]{B} 的说明，这个平面无关紧要，所以$f\equiv0$。\origNote[。][]{84}

\hypertarget{origpage-67}{}\phantomsection\label{origpage-67}因此，$Af=\lambda f$从来没有非零解。不过，我们也能不严格地看出，应到哪里寻找解。$(q_j-\lambda)f(q_1,\allowbreak \ldots,\allowbreak q_l)=0$是说，$f$只有在$q_j=\lambda$时才允许不为零。因此，对若干$\lambda$值，例如$\lambda=\lambda',\allowbreak \lambda'',\allowbreak \ldots,\allowbreak \lambda^{(s)}$，把相应解作线性组合，得到的$f$只能在
\[
q_j=\lambda',\lambda'',\ldots,\lambda^{(s)}
\]
时不为零。由一切$\lambda\leq\lambda_0$的解作线性组合，因而可看成得到只在$\lambda\leq\lambda_0$时不为零的$f$。\tnote[。][而][原]{原文此处写作$\lambda\leq\lambda_0$；按本段讨论的函数支撑，应为$q_j\leq\lambda_0$。正文保留原文。}而在纯点谱的情形，我们曾有
\[
E(\lambda_0)=\sum_{\lambda_\rho\leq\lambda_0}P_{[\varphi_\rho]}=P_{\mathfrak{N}_{\lambda_0}},
\qquad\mathfrak{N}_{\lambda_0}=[\varphi_{\rho(\lambda_\rho\leq\lambda_0)}];
\]
即$\mathfrak{N}_\lambda$由一切满足$\lambda_\rho\leq\lambda_0$的$\varphi_\rho$的线性组合组成，也就是由一切满足$\lambda\leq\lambda_0$的方程$Af=\lambda f$的解的线性组合组成。因此可以预期——当然，这只是不严格的启发性推论——这里也有$E(\lambda_0)=P_{\mathfrak{N}_{\lambda_0}}$；其中$\mathfrak{N}_{\lambda_0}$由那些只在$q_j\leq\lambda_0$时不为零的$f$组成。于是$\mathfrak{R}_\infty-\mathfrak{N}_{\lambda_0}$显然由在$q_j\leq\lambda_0$时恒为零的$f$组成，所以
\[
E(\lambda_0)f(q_1\cdots q_l)=
\begin{cases}f(q_1\cdots q_l),&q_j\leq\lambda_0,\\0,&q_j>\lambda_0.\end{cases}
\]
我们就这样以完全不严格的方法，找到了一个投影算符族$E(\lambda)$，可以预期它对这里的$A$满足$\hyperref[cond:II.Sbar1]{\overline{S}_1}$—$\hyperref[cond:II.Sbar3]{\overline{S}_3}$。事实上，$\hyperref[cond:II.Sbar1]{\overline{S}_1}$、$\hyperref[cond:II.Sbar2]{\overline{S}_2}$显然成立；而且，在$\hyperref[cond:II.Sbar1]{\overline{S}_1}$的$\lambda\to\lambda_0$情形，连$\lambda\geq\lambda_0$的限制也不需要，也就是说$E(\lambda)$对于$\lambda$处处连续。为证明$\hyperref[cond:II.Sbar3]{\overline{S}_3}$也成立，只须验证下列等式：
\begin{align*}
&\int_{-\infty}^{\infty}\lambda^2\,d\|E(\lambda)f\|^2\\
&=\int_{-\infty}^{\infty}\lambda^2\,d\left(\begin{aligned}
&\int_{-\infty}^{\infty}\cdots\int_{-\infty}^{\lambda}\cdots\int_{-\infty}^{\infty}|f(q_1\cdots q_j\cdots q_l)|^2\\
&\qquad{}\,dq_1\cdots dq_j\cdots dq_l
\end{aligned}\right)\\
&=\int_{-\infty}^{\infty}\lambda^2\left(\begin{aligned}
&\int_{-\infty}^{\infty}\cdots\int_{-\infty}^{\infty}|f(q_1\cdots q_{j-1}\lambda q_{j+1}\cdots q_l)|^2\\
&\qquad{}\,dq_1\cdots dq_{j-1}\,dq_{j+1}\cdots dq_l
\end{aligned}\right)d\lambda\\
&=\int_{-\infty}^{\infty}\cdots\int_{-\infty}^{\infty}q_j^2|f(q_1\cdots q_{j-1}q_jq_{j+1}\cdots q_l)|^2\\
&\qquad{}\,dq_1\cdots dq_{j-1}\,dq_j\,dq_{j+1}\cdots dq_l\\
&=\|Af\|^2,
\end{align*}
以及
\begin{align*}
&\int_{-\infty}^{\infty}\lambda\,d(E(\lambda)f,g)\\
&=\int_{-\infty}^{\infty}\lambda\,d\left(\begin{aligned}
&\int_{-\infty}^{\infty}\cdots\int_{-\infty}^{\lambda}\cdots\int_{-\infty}^{\infty}f(q_1\cdots q_j\cdots q_l)\\
&\quad\times\overline{g(q_1\cdots q_j\cdots q_l)}\,dq_1\cdots dq_j\cdots dq_l
\end{aligned}\right)\\
&=\int_{-\infty}^{\infty}\lambda\left(\begin{aligned}
&\int_{-\infty}^{\infty}\cdots\int_{-\infty}^{\infty}f(q_1\cdots q_{j-1}\lambda q_{j+1}\cdots q_l)\\
&\quad\times\overline{g(q_1\cdots q_{j-1}\lambda q_{j+1}\cdots q_l)}\\
&\quad\times dq_1\cdots dq_{j-1}\,dq_{j+1}\cdots dq_l
\end{aligned}\right)d\lambda\\
&=\int_{-\infty}^{\infty}\cdots\int_{-\infty}^{\infty}q_jf(q_1\cdots q_{j-1}q_jq_{j+1}\cdots q_l)\overline{g(q_1\cdots q_{j-1}q_jq_{j+1}\cdots q_l)}\\
&\hspace{4em}\times dq_1\cdots dq_{j-1}\,dq_j\,dq_{j+1}\cdots dq_l=(Af,g).
\end{align*}

\hypertarget{origpage-68}{}\phantomsection\label{origpage-68}这里又可看出，点谱，或说本征值问题的旧定义，必然不能满足要求，因为$E(\lambda)$处处连续增长。

这个例子还提示，在一般情形下，怎样求连续谱中的$E(\lambda)$。先不作严格的数学论证，求出$Af=\lambda f$的解；由于$\lambda$属于连续谱，这些$f$根本不属于$\mathfrak{R}_\infty$。再把所有$\lambda\leq\lambda_0$的解作线性组合。其中一些组合重新属于$\mathfrak{R}_\infty$，并可能构成闭线性流形$\mathfrak{N}_{\lambda_0}$。于是令$E(\lambda_0)=P_{\mathfrak{N}_{\lambda_0}}$。如果前面的做法恰当，事后就能对$A$和这些$E(\lambda)$验证$\hyperref[cond:II.Sbar1]{\overline{S}_1}$—$\hyperref[cond:II.Sbar3]{\overline{S}_3}$，从而\textbf{事后}把先前的启发性推论补成严格论证。\origNote[。][]{85}

第二个例子，是波动力学中另一个重要算符$\frac{h}{2\pi i}\frac{\partial}{\partial q_j}$。为免无关紧要的复杂情形，取$l=j=1$；其他情形同理。因此须考察算符
\[
A'f(q)=\frac{h}{2\pi i}\frac{\partial}{\partial q}f(q).
\]
当$q$的变化范围为$-\infty<q<+\infty$时，它是厄米算符，这已在 \hyperref[sec:II.5]{II.5} 看到。但在有限区间$a\leq q\leq b$上，情况却不同：
\begin{align*}
(A'f,g)-(f,A'g)
&=\int_a^b\frac{h}{2\pi i}f'(q)\overline{g(q)}\,dq-\int_a^b f(q)\overline{\frac{h}{2\pi i}g'(q)}\,dq\\
&=\frac{h}{2\pi i}\int_a^b\{f'(q)\overline{g(q)}+f(q)\overline{g(q)}'\}\,dq\\
&=\frac{h}{2\pi i}[f(q)\overline{g(q)}]_a^b
=\frac{h}{2\pi i}[f(b)\overline{g(b)}-f(a)\overline{g(a)}].
\end{align*}
要使这个式子为零，就须缩小$\frac{h}{2\pi i}\frac{\partial}{\partial q}$的定义域，使任意取出的$f,\allowbreak g$都满足$f(a)\overline{g(a)}=f(b)\overline{g(b)}$，即$f(a):f(b)=\overline{g(b)}:\overline{g(a)}$。固定$g$而改变$f$，便知$f(a):f(b)$在整个定义域内必须为同一个数$\theta$（$\theta$也可能为$0$或$\infty$）；把$f,\allowbreak g$代入又得$\theta=1/\overline\theta$，即$|\theta|=1$。因此，要使$\frac{h}{2\pi i}\frac{\partial}{\partial q}$厄米，还必须加上一条形如
\[
f(a):f(b)=\theta
\]
的“边界条件”，其中$\theta$为任意固定的、绝对值等于$1$的数。

先取区间$-\infty<q<+\infty$。

\hypertarget{origpage-69}{}\phantomsection\label{origpage-69}$A\varphi=\lambda\varphi$，即$\frac{h}{2\pi i}\varphi'(q)=\lambda\varphi(q)$的解，是函数
\[
\varphi(q)=c e^{\frac{2\pi i}{h}\lambda q}.
\]
但这些函数不能直接拿来用于我们的目的，因为$\int_{-\infty}^{\infty}|\varphi(q)|^2\,dq=\int_{-\infty}^{\infty}|c|^2\,dq=+\infty$，除非$c=0$，即$\varphi\equiv0$。请注意：第一个例子得到的解是$\delta(q-\lambda)$，一个虚构的、并不存在的函数（见\CJKglue \hyperlink{note-84}{原注 \getrefnumber{orig-note-84}}）；现在得到的却是$e^{\frac{2\pi i}{h}\lambda q}$，一个完全正常的函数，只是由于绝对值平方的积分无限，不属于$\mathfrak{R}_\infty$。从我们的立场看，两种情形是一回事：凡不属于$\mathfrak{R}_\infty$的对象，对我们来说就不存在。\origNote[。][]{86}

与第一个例子一样，现在对$\lambda\leq\lambda_0$的解作线性组合，即构造函数
\[
f(q)=\int_{-\infty}^{\lambda_0}c(\lambda)e^{\frac{2\pi i}{h}\lambda q}\,d\lambda.
\]
希望其中有些函数属于$\mathfrak{R}_\infty$，并希望这些函数构成闭线性流形$\mathfrak{N}'_{\lambda_0}$；最后还希望投影算符$E'(\lambda_0)=P_{\mathfrak{N}'_{\lambda_0}}$构成属于$A'$的单位分解。第一个猜想可以由下例说明：令
\[
c(\lambda)=\begin{cases}1,&\lambda\geq\lambda_1,\\0,&\lambda<\lambda_1,\end{cases}
\qquad\lambda_1<\lambda_0.
\]
得到
\[
f(q)=\int_{\lambda_1}^{\lambda_0}e^{\frac{2\pi i}{h}\lambda q}\,d\lambda
=\frac{e^{\frac{2\pi i}{h}\lambda_0q}-e^{\frac{2\pi i}{h}\lambda_1q}}{\frac{2\pi i}{h}q}.
\]
这个$f(q)$在每个有限点都正则，当$q\to\pm\infty$时按$1/q$衰减，所以$\int_{-\infty}^{\infty}|f(q)|^2\,dq$有限。其他猜想也成立，这可由Fourier积分理论推出。这个理论的内容如下：\origNote[：][]{87}

设$f(x)$是任意使$\int_{-\infty}^{\infty}|f(x)|^2\,dx$有限的函数，则可构造函数
\[
Lf(x)=F(y)=\frac1{\sqrt{2\pi}}\int_{-\infty}^{\infty}e^{ixy}f(x)\,dx
\]
它的$\int_{-\infty}^{\infty}|F(y)|^2\,dy$仍有限，而且恰等于$\int_{-\infty}^{\infty}|f(x)|^2\,dx$。此外，$LLf(x)=f(-x)$。这就是所谓Laplace变换，它在微分方程理论的其他地方也起着很大作用。

\hypertarget{origpage-70}{}\phantomsection\label{origpage-70}把$x,\allowbreak y$分别换成$\sqrt{2\pi/h}\,q,\allowbreak \sqrt{2\pi/h}\,p$，便得到具有相同性质的变换
\[
Mf(q)=F(p)=\frac1{\sqrt h}\int_{-\infty}^{\infty}e^{\frac{2\pi i}{h}pq}f(q)\,dq,
\]
所以它把$\mathfrak{R}_\infty$映满自身：对$\mathfrak{R}_\infty$中的每个$g(p)$，$Mf(q)=g(p)$都有解$f(q)=Mg(-p)$。它保持$\|f\|$不变，又是线性的；因此，依 \hyperref[sec:II.5]{II.5}，这个$M$是酉的。又有$M^2f(q)=f(-q)$，所以$M^{-1}f(q)=M^*f(q)=Mf(-q)$；并且$M$与$M^2$，即与$f(q)\mapsto f(-q)$对易。

所以，对$\mathfrak{N}'_{\lambda_0}$的预期定义就是：当$F(p)=M^{-1}f(q)$对一切$p>\lambda_0$都等于零时，$f(q)$属于$\mathfrak{N}'_{\lambda_0}$。这里
\[
F(p)=\sqrt h\,c(p),
\]
其中$c(\lambda)$即前面的函数。我们知道，这些$F(p)$构成闭线性流形$\mathfrak{N}_{\lambda_0}$，所以它经$M$得到的像$\mathfrak{N}'_{\lambda_0}$也如此。由$E(\lambda_0)$得到$E'(\lambda_0)$的方式，与由$\mathfrak{N}_{\lambda_0}$得到$\mathfrak{N}'_{\lambda_0}$一样，都是对整个$\mathfrak{R}_\infty$施行变换$M$。因此$E'(\lambda_0)=ME(\lambda_0)M^{-1}$。于是，$E'(\lambda)$与$E(\lambda)$一样，具有性质$\hyperref[cond:II.Sbar1]{\overline{S}_1}$、$\hyperref[cond:II.Sbar2]{\overline{S}_2}$。还须证明$\hyperref[cond:II.Sbar3]{\overline{S}_3}$，即单位分解$E(\lambda)$属于$A'$。

这里只限于证明：若$f(q)$可微，且不存在特别的收敛困难，并有
\[
\int_{-\infty}^{\infty}|\frac{h}{2\pi i}f'(q)|^2\,dq\quad\text{有限，}
\]
则$\int_{-\infty}^{\infty}\lambda^2\,d\|E'(\lambda)f\|^2$有限，且$(A'f,\allowbreak g)=\int_{-\infty}^{\infty}\lambda\,d(E'(\lambda)f,\allowbreak g)$。\origNote[。][事]{88}事实上，设$M^{-1}f(q)=F(p)$，则
\begin{align*}
A'f(q)&=\frac{h}{2\pi i}f'(q)
=\frac{h}{2\pi i}\frac{\partial}{\partial q}(MF(p))\\
&=\frac{h}{2\pi i}\frac{\partial}{\partial q}\left(\frac1{\sqrt h}\int F(p)e^{\frac{2\pi i}{h}pq}\,dp\right)\\
&=\frac{\sqrt h}{2\pi i}\int F(p)\frac{\partial}{\partial q}\left(e^{\frac{2\pi i}{h}pq}\right)dp\\
&=\frac1{\sqrt h}\int F(p)\,p\,e^{\frac{2\pi i}{h}pq}\,dp=M(pF(p)).
\end{align*}
所以，对于所说的$f$，有$A'=MAM^{-1}$；这里$A$是算符$q\cdots$，或者说$p\cdots$，因为现在采用的变量是$p$。上述断言对$A,\allowbreak E(\lambda)$成立；用$M$变换$\mathfrak{R}_\infty$之后，它们仍成立，所以对$A'=MAM^{-1}$和$E'(\lambda)=ME(\lambda)M^{-1}$也成立。

\hypertarget{origpage-71}{}\phantomsection\label{origpage-71}对于有限区间$a\leq q\leq b$（$a<b$且两者有限）上的$\frac{h}{2\pi i}\frac{\partial}{\partial q}$，情形有本质区别。我们知道，要保持厄米性，还须加上边界条件
\[
f(a):f(b)=\theta,\qquad|\theta|=1,
\]
$\frac{h}{2\pi i}\frac{\partial}{\partial q}f(q)=\lambda f(q)$仍由
\[
f(q)=c e^{\frac{2\pi i}{h}\lambda q}
\]
求解；但现在$\int_a^b|f(q)|^2\,dq=\int_a^b|c|^2\,dq=(b-a)|c|^2$有限，因此$f(q)$总属于$\mathfrak{R}_\infty$。不过还须满足边界条件：
\[
f(a):f(b)=e^{\frac{2\pi i}{h}\lambda(a-b)}=\theta.
\]
令$\theta=e^{-i\alpha}$，$0\leq\alpha<2\pi$，则为
\begin{align*}
\frac{2\pi i}{h}\lambda(a-b)&=-i\alpha-2k\pi i\qquad(k=0,\pm1,\pm2,\ldots),\\
\lambda&=\frac h{b-a}\left(\frac\alpha{2\pi}+k\right).
\end{align*}
所以这里存在点谱。归一化的解由$(b-a)|c|^2=1$确定，例如可取$c=1/\sqrt{b-a}$：
\[
\varphi_k(q)=\frac1{\sqrt{b-a}}e^{\frac{2\pi i}{h}\lambda q}
=\frac1{\sqrt{b-a}}e^{\frac{2\pi i}{b-a}(\frac\alpha{2\pi}+k)q},
\qquad k=0,\pm1,\pm2,\ldots.
\]
这确是规范正交系，而且也是完全的。因为，若$f(q)$与所有$\varphi_k(q)$正交，$e^{\frac{\alpha i}{b-a}q}f(q)$就与所有$e^{\frac{2\pi i}{b-a}kq}$正交，从而$e^{\frac{\alpha i}{2\pi}x}f(\frac{b-a}{2\pi}x)$与所有$e^{ikx}$正交，即与$1,\allowbreak \cos x,\allowbreak \sin x,\allowbreak \cos2x,\allowbreak \sin2x,\allowbreak \ldots$正交。它定义在区间
\[
a\leq\frac{b-a}{2\pi}x\leq b
\]
上；区间长度为$2\pi$，所以根据熟知的定理，它必为零。\origNote[。][因]{89}因此$f(q)\equiv0$。

所以，这里是纯点谱情形，我们在本节开头已经一般地解决了它。请注意，“边界条件”，即$\theta$或$\alpha$，如何影响本征值和本征函数。

最后，还可以考察半无限区间，例如$0\leq q<+\infty$。

\hypertarget{origpage-72}{}\phantomsection\label{origpage-72}同样，首先必须检验厄米性。有
\begin{align*}
(A'f,g)-(f,A'g)&=\frac{h}{2\pi i}\int_0^{+\infty}\bigl(f'(q)\overline{g(q)}+f(q)\overline{g'(q)}\bigr)\,dq\\
&=\frac{h}{2\pi i}[f(q)\overline{g(q)}]_0^{+\infty}.
\end{align*}
当$q\to+\infty$时，$f(q)\overline{g(q)}\to0$，证明与 \hyperref[sec:II.5]{II.5} 中区间$-\infty,\allowbreak +\infty$的情形相同。因此须有$f(0)\overline{g(0)}=0$。令$f=g$就知，边界条件必须是$f(0)=0$。

这时出现了严重困难。$A'\varphi=\lambda\varphi$的解，与区间$-\infty<q<+\infty$的情形相同，即$ce^{\frac{2\pi i}{h}\lambda q}$；它们不属于$\mathfrak{R}_\infty$，还违反边界条件。后一现象已经提示，这种求解方法可能有问题。更奇怪的是，按前面勾画的方法，由于这些形式解，即不属于$\mathfrak{R}_\infty$的解相同，就应得到与$-\infty<q<+\infty$的情形相同的$E(\lambda)$。可是，算符已经不同，怎么会仍得到相同的结果呢？而且，所得结果也不符合我们的要求。因为，在由使$\int_0^{+\infty}|f(q)|^2\,dq$有限的$f(q)$（$0\leq q<+\infty$）构成的希尔伯特空间$F_{\Omega'}$中，若仍以
\begin{align*}
Mf(q)=F(p)&=\frac1{\sqrt h}\int_0^{\infty}e^{\frac{2\pi i}{h}pq}f(q)\,dq,\\
M^{-1}F(p)=f(q)&=\frac1{\sqrt h}\int_{-\infty}^{\infty}e^{-\frac{2\pi i}{h}pq}F(p)\,dp
\quad(=MF(-p))
\end{align*}
定义$M,\allowbreak M^{-1}$，$M$就把这个由使$\int_0^{\infty}|f(q)|^2\,dq$有限的全部$f(q)$（$0\leq q<\infty$）构成的希尔伯特空间$F_{\Omega'}$，映射到另一个希尔伯特空间：由使$\int_{-\infty}^{\infty}|F(p)|^2\,dp$有限的全部$F(p)$（$-\infty<p<\infty$）构成的$F_{\Omega''}$。

自然，仍有$\|Mf(q)\|=\|f(q)\|$，因为把$f(q)$在$-\infty<q<0$的部分延拓为零，就能使用前面提到的定理。但是$\|M^{-1}F(p)\|=\|F(p)\|$却绝非普遍成立。依那些定理，若也在$q<0$时用$f(q)=\frac1{\sqrt h}\int_{-\infty}^{\infty}e^{-\frac{2\pi i}{h}pq}F(p)\,dp$定义$f(q)$，则
\begin{align*}
\|F\|^2&=\int_{-\infty}^{\infty}|F(p)|^2\,dp\\
&=\int_{-\infty}^{\infty}|f(q)|^2\,dq,\\
\|M^{-1}F\|^2&=\|f\|^2=\int_0^{\infty}|f(q)|^2\,dq.
\end{align*}

\hypertarget{origpage-73}{}\phantomsection\label{origpage-73}所以，除非刚才这样定义出来的$f(q)$碰巧对一切$q<0$都为零，否则便有$\|M^{-1}F\|<\|F\|$。因此，$E'(\lambda)=ME(\lambda)M^{-1}$根本不是单位分解；\origNote[；][这]{90}这种方法失败了。

我们很快就会看到（\hyperlink{note-105}{原注 \getrefnumber{orig-note-105}}），这并非偶然，而是问题本身的性质所决定的，因为根本没有属于这个算符的单位分解。

结束这些初步考察之前，再给出若干形式运算法则，适用于那些已借助单位分解写成符号形式
\[
A=\int_{-\infty}^{\infty}\lambda\,dE(\lambda)
\]
的算符。首先，设$F$是与所有$E(\lambda)$对易的投影算符。对一切$\lambda'<\lambda''$有
\begin{align*}
&\|E(\lambda'')Ff-E(\lambda')Ff\|^2\\
&=\|(E(\lambda'')-E(\lambda'))Ff\|^2\\
&=\|F(E(\lambda'')-E(\lambda'))f\|^2\\
&\leq\|(E(\lambda'')-E(\lambda'))f\|^2.
\end{align*}
由于$E(\lambda''),\allowbreak E(\lambda'),\allowbreak E(\lambda'')-E(\lambda')$，以及$E(\lambda'')F,\allowbreak E(\lambda')F,\allowbreak E(\lambda'')F-E(\lambda')F=(E(\lambda'')-E(\lambda'))F$都是投影算符，因此得到
\[
\|E(\lambda'')Ff\|^2-\|E(\lambda')Ff\|^2
\leq\|E(\lambda'')f\|^2-\|E(\lambda')f\|^2.
\]
于是$\int_{-\infty}^{\infty}\lambda^2\,d\|E(\lambda)f\|^2\geq\int_{-\infty}^{\infty}\lambda^2\,d\|E(\lambda)Ff\|^2$；由$\hyperref[cond:II.Sbar3]{\overline{S}_3}$知，只要$Af$有意义，$AFf$也有意义。又有\origNote[有][]{91}
\[
AF=\int_{-\infty}^{\infty}\lambda\,d(E(\lambda)F)
=\int_{-\infty}^{\infty}\lambda\,d(FE(\lambda))=FA.
\]
即$A,\allowbreak F$也对易。特别地，由$\hyperref[cond:II.Sbar2]{\overline{S}_2}$，可以把任意一个$E(\lambda)$取作$F$。这时
\begin{align*}
AE(\lambda)&=\int_{-\infty}^{\infty}\lambda'\,d(E(\lambda')E(\lambda))
=\int_{-\infty}^{\infty}\lambda'\,d(E(\operatorname{Min}(\lambda,\lambda')))\\
&=\int_{-\infty}^{\lambda}+\int_\lambda^{\infty}
=\int_{-\infty}^{\lambda}\lambda'\,dE(\lambda')+\int_\lambda^{\infty}\lambda'\,dE(\lambda).
\end{align*}
最后一个积分中$d$之后为常数，所以积分为零，因而
\[
AE(\lambda)=E(\lambda)A=\int_{-\infty}^{\lambda}\lambda'\,dE(\lambda').
\]
进而有\origNote[有][]{92}
\[
A^2=\int_{-\infty}^{\infty}\lambda\,d(E(\lambda)A)
=\int_{-\infty}^{\infty}\lambda\,d\left(\int_{-\infty}^{\lambda}\lambda'\,dE(\lambda')\right)
=\int_{-\infty}^{\infty}\lambda^2\,dE(\lambda).
\]

\hypertarget{origpage-74}{}\phantomsection\label{origpage-74}一般有
\[
A^n=\int_{-\infty}^{\infty}\lambda^n\,dE(\lambda),
\]
因为，从$n-1$到$n$的归纳步骤为\origNote[为][]{92}
\begin{align*}
A^n=A^{n-1}A
&=\int_{-\infty}^{\infty}\lambda^{n-1}\,d(E(\lambda)A)\\
&=\int_{-\infty}^{\infty}\lambda^{n-1}\,d\left(\int_{-\infty}^{\lambda}\lambda'\,dE(\lambda')\right)\\
&=\int_{-\infty}^{\infty}\lambda^{n-1}\lambda\,dE(\lambda)
=\int_{-\infty}^{\infty}\lambda^n\,dE(\lambda).
\end{align*}
所以，对于任意多项式$p(x)=a_0+a_1x+\cdots+a_nx^n$，有
\[
p(A)=\int_{-\infty}^{\infty}p(\lambda)\,dE(\lambda).
\]
这里当然把$p(A)$理解为$a_01+a_1A+\cdots+a_nA^n$；恒等式$\int_{-\infty}^{\infty}dE(\lambda)=1$由$\hyperref[cond:II.Sbar1]{\overline{S}_1}$得出。

再指出，若$r(\lambda),\allowbreak s(\lambda)$为任意函数，并用符号关系\origNote[系][]{93}
\[
B=\int_{-\infty}^{\infty}r(\lambda)\,dE(\lambda),\qquad
C=\int_{-\infty}^{\infty}s(\lambda)\,dE(\lambda)
\]
定义两个算符$B,\allowbreak C$，则有
\[
BC=\int_{-\infty}^{\infty}r(\lambda)s(\lambda)\,dE(\lambda).
\]
证明与特殊情形$B=C=A$完全相同：
\begin{align*}
BE(\lambda)&=\int_{-\infty}^{\infty}r(\lambda')\,d(E(\lambda')E(\lambda))
=\int_{-\infty}^{\infty}r(\lambda')\,d(E(\operatorname{Min}(\lambda,\lambda')))\\
&=\int_{-\infty}^{\lambda}+\int_\lambda^{\infty}
=\int_{-\infty}^{\lambda}r(\lambda')\,dE(\lambda')+\int_\lambda^{\infty}r(\lambda')\,dE(\lambda)\\
&=\int_{-\infty}^{\lambda}r(\lambda')\,dE(\lambda'),\\
CB&=\int_{-\infty}^{\infty}s(\lambda)\,d(BE(\lambda))
=\int_{-\infty}^{\infty}s(\lambda)\,d\left(\int_{-\infty}^{\lambda}r(\lambda')\,dE(\lambda')\right)\\
&=\int_{-\infty}^{\infty}s(\lambda)\cdot r(\lambda)\,dE(\lambda)
=\int_{-\infty}^{\infty}s(\lambda)r(\lambda)\,dE(\lambda).
\end{align*}
下列关系的验证则很容易：
\begin{gather*}
B^*=\int_{-\infty}^{\infty}\overline{r(\lambda)}\,dE(\lambda),\qquad
aB=\int_{-\infty}^{\infty}ar(\lambda)\,dE(\lambda),\\
B\pm C=\int_{-\infty}^{\infty}(r(\lambda)\pm s(\lambda))\,dE(\lambda).
\end{gather*}
因此，在形式上也没有理由不能对这样的函数$r(\lambda)$写$B=r(A)$。\origNote[。][]{94}

\hypertarget{origpage-75}{}\phantomsection\label{origpage-75}特别值得注意的是不连续函数
\[
e_\lambda(\lambda')=\begin{cases}1,&\lambda'\leq\lambda,\\0,&\lambda'>\lambda.\end{cases}
\]
对这些函数，由$\hyperref[cond:II.Sbar1]{\overline{S}_1}$有
\[
e_\lambda(A)=\int_{-\infty}^{\infty}e_\lambda(\lambda')\,dE(\lambda)
=\int_{-\infty}^{\lambda}dE(\lambda')=E(\lambda).
\]
本节开头讨论过算符$A=q_j\cdots$；它的$E(\lambda)$是在$q_j\leq\lambda$或$q_j>\lambda$时，分别乘以$1$或$0$，也就是乘以$e_\lambda(q)$。因此，$e_\lambda(q_j\cdots)=e_\lambda(q_j)\cdots$；这个例子可以帮助理解这些概念。

\section{本征值问题唯一性与可解性的补充讨论}\label{sec:II.9}
上一节仅结合特殊例子，从定性上概述了哪些单位分解$E(\lambda)$属于给定的厄米算符$A$，还缺少系统研究。若在数学上完整展开，就超出本书范围；我们只能证明其中一些内容，对其余内容加以介绍。况且，要理解量子力学，并不一定需要详细了解这些问题。\origNote[。][]{95}

\hyperref[thm:II.18]{定理18} 表明，线性运算的连续性可由
\begin{equation}
\|Af\|\leq C\|f\|\tag{St}\label{eq:II.9.St}
\end{equation}
表达，其中$C$任意但固定。根据\CJKglue \hyperref[thm:II.18]{定理18}，条件 \hyperref[eq:II.9.St]{St} 还有若干等价形式：
\begin{equation}
|(Af,g)|\leq C\|f\|\|g\|,\tag{$\mathrm{St}_1$}\label{eq:II.9.St1}
\end{equation}
\begin{equation}
|(Af,f)|\leq C\|f\|^2.\tag{$\mathrm{St}_2$}\label{eq:II.9.St2}
\end{equation}
后一形式只适用于厄米的$A$。

与连续性等价的条件$\hyperref[eq:II.9.St1]{\mathrm{St}_1}$，就是Hilbert的有界性概念。Hilbert对有界，即连续的厄米算符，表述并解决了本征值问题（参见\CJKglue \hyperlink{note-70}{原注 \getrefnumber{orig-note-70}}）。不过，在讨论这一点之前，必须再引入一个概念。

厄米算符$A$，若具有下述性质，就称为闭算符：对任意点列$f_1,\allowbreak f_2,\allowbreak \ldots$，若各个$Af_n$均有意义，且$f_n\to f,\allowbreak Af_n\to f^*$，则$Af$也有意义，并等于$f^*$。须注意，连续性也可以用一种十分接近的方式定义：若所有$Af_n$及$Af$都有意义，且$f_n\to f$，就有$Af_n\to Af$。两个定义的差别在于，闭性把$Af_n$存在极限$f^*$作为前提，只在这一前提下断言$f^*=Af$；而连续性还断言$f^*$的存在。

\hypertarget{origpage-76}{}\phantomsection\label{origpage-76}举几个例子。仍设$\mathfrak{R}_\infty$由使$\int_{-\infty}^{\infty}|f(q)|^2\,dq$有限的全部$f(q)$（$-\infty<q<\infty$）构成。设$A$为算符$q\cdots$，对使$\int_{-\infty}^{\infty}|f(q)|^2\,dq$和$\int_{-\infty}^{\infty}q^2|f(q)|^2\,dq$有限的一切$f(q)$定义；$A'$为算符$\frac{h}{2\pi i}\frac{\partial}{\partial q}$，对处处可微且使$\int_{-\infty}^{\infty}|f(q)|^2\,dq$、$\int_{-\infty}^{\infty}|\frac{h}{2\pi i}f'(q)|^2\,dq$有限的一切函数定义。我们知道，它们都厄米。$A$是闭的：设$f_n\to f,\allowbreak Af_n\to f^*$，即
\[
\int_{-\infty}^{\infty}|f_n(q)-f(q)|^2\,dq\to0,\qquad
\int_{-\infty}^{\infty}|qf_n(q)-f^*(q)|^2\,dq\to0.
\]
根据 \hyperref[sec:II.3]{II.3} 中证明 \hyperref[ax:II.D]{D} 时的论证，$f_1,\allowbreak f_2,\allowbreak \ldots$有一个子序列$f_{n_1},\allowbreak f_{n_2},\allowbreak \ldots$，除$q$的一个零测集外处处收敛于极限：$f_{n_\nu}(q)\to g(q)$。所以
\[
\int_{-\infty}^{\infty}|g(q)-f(q)|^2\,dq=0,\qquad
\int_{-\infty}^{\infty}|qg(q)-f^*(q)|^2\,dq=0.
\]
因此，除一个零测集外，有$g(q)=f(q)$及$qg(q)=f^*(q)$，从而$qf(q)=f^*(q)$；也就是说，$f^*(q)$与$qf(q)$没有本质区别。而按假定，$f^*(q)$属于$\mathfrak{R}_\infty$，$qf(q)$也就属于它。所以$Af(q)$有意义，且等于$qf(q)=f^*(q)$。

相反，$A'$不闭。令$f_n(q)=e^{-\sqrt{q^2+1/n}},\allowbreak f(q)=e^{-|q|}$。显然，所有$Af_n$都有意义，但$Af$没有意义，因为$f$在$q=0$处不可微。尽管如此，不难算出：若令$f^*(q)=-\operatorname{sgn}(q)e^{-|q|}$，其中$q<0,\allowbreak q=0,\allowbreak q>0$时$\operatorname{sgn}(q)$分别为$-1,\allowbreak 0,\allowbreak +1$，则仍有$f_n\to f,\allowbreak Af_n\to f^*$。\tnote[。][][此]{此例中，原书在本应指微分算符$A'$的地方写作$A$，且$f^*$的式子漏去其因子$h/(2\pi i)$。记号和公式均按原书保留。}

现在要证明：与连续性不同，厄米算符总能不费多少力气就获得闭性，方法是扩张。所谓扩张，就是在$\mathfrak{R}_\infty$中原来已有定义的一切点上，保持算符不变；再在原来尚无定义的某些点上，补充定义。

设$A$为任意厄米算符。如下定义算符$\widetilde A$：若存在序列$f_1,\allowbreak f_2,\allowbreak \ldots$，使各个$Af_n$均有意义，$f_n$以$f$为极限，$Af_n$也有极限$f^*$，就定义$\widetilde Af=f^*$。这个定义只有在不产生歧义时才允许，即由$f_n\to f,\allowbreak g_n\to f,\allowbreak Af_n\to f^*,\allowbreak Ag_n\to g^*$必须能推出$f^*=g^*$。事实上，只要$Ag$有意义，就有
\begin{align*}
(f^*,g)&=\lim(Af_n,g)=\lim(f_n,Ag)=(f,Ag),\\
(g^*,g)&=\lim(Ag_n,g)=\lim(g_n,Ag)=(f,Ag).
\end{align*}

\hypertarget{origpage-77}{}\phantomsection\label{origpage-77}所以$(f^*,\allowbreak g)=(g^*,\allowbreak g)$。这些$g$处处稠密，故$f^*=g^*$，从而$\widetilde A$的定义成立。这个$\widetilde A$是$A$的扩张：只要$Af$有意义，$\widetilde Af$也有意义，且等于$Af$；只须令所有$f_n=f,\allowbreak f^*=Af$即可。通过取极限，由$A$的线性和厄米性，就能推出$\widetilde A$也具有这两种性质。最后，$\widetilde A$是闭的。设所有$\widetilde Af_n$有意义，且$f_n\to f,\allowbreak \widetilde Af_n\to f^*$。于是有序列$f_{n,\allowbreak 1},\allowbreak f_{n,\allowbreak 2},\allowbreak \ldots$，使$Af_{n,\allowbreak m}$有意义，$f_{n,\allowbreak m}\to f_n,\allowbreak Af_{n,\allowbreak m}\to f_n^*$，且$\widetilde Af_n=f_n^*$。对每个$n$都存在$N_n$，使$m\geq N_n$时
\[
\|f_{n,m}-f_n\|\leq\frac1n,\qquad\|Af_{n,m}-f_n^*\|\leq\frac1n.
\]
于是$f_{n,\allowbreak N_n}-f_n\to0,\allowbreak Af_{n,\allowbreak N_n}-\widetilde Af_n\to0$，从而$f_{n,\allowbreak N_n}-f\to0,\allowbreak Af_{n,\allowbreak N_n}-f^*\to0$。根据定义，就有$\widetilde Af=f^*$。

须注意，不连续的算符决不可能靠扩张变成连续的。

若算符$B$扩张算符$A$，即只要$Af$有意义，$Bf$就也有意义且等于$Af$，我们就记作$B\succ A$或$A\prec B$。上面已经证明$A\prec\widetilde A$，而$\widetilde A$厄米且闭。不难看出，对每个满足$A\prec B$的闭算符$B$，也必有$\widetilde A\prec B$。因此，$\widetilde A$是$A$的最小闭扩张，且$\widetilde{\widetilde A}=\widetilde A$。

$A$与$\widetilde A$的关系如此密切，自然使人想到，在一切讨论中都用$\widetilde A$代替$A$。因为，只有$\widetilde A$才把$A$的定义域扩充到了合理的范围；换言之，与$\widetilde A$相比，$A$的定义域受到了不必要的限制。我们就作这样的替换；此后可以假定一切厄米算符都是闭的。

再考察连续的厄米算符$A$。对它来说，算符闭等价于定义域闭。而表征连续性的条件$\|Af\|\leq C\|f\|$显然也对$\widetilde A$成立，所以$\widetilde A$也连续。它的定义域于是是闭的，又必须处处稠密，故等于$\mathfrak{R}_\infty$。因此$\widetilde A$处处有定义，每个闭的连续算符都如此。反过来也成立：若闭算符处处有定义，它就连续。这是Toeplitz定理，\origNote[，][这]{96}这里不能讨论其证明。

Hilbert的结论如下：每个连续厄米算符都有且仅有一个属于它的单位分解（参见\CJKglue \hyperlink{note-70}{原注 \getrefnumber{orig-note-70}} 所引文献）。由于连续算符总有意义，$\int\lambda^2\,d\|E(\lambda)f\|^2$必须总是有限的。它又等于$\|Af\|^2$，所以根据 \hyperref[eq:II.9.St]{St} 不超过$C^2\|f\|^2$，于是有

\hypertarget{origpage-78}{}\phantomsection\label{origpage-78}\begin{align*}
0&\geq\|Af\|^2-C^2\|f\|^2\\
&=\int_{-\infty}^{\infty}\lambda^2\,d\|E(\lambda)f\|^2-C^2\int_{-\infty}^{\infty}d\|E(\lambda)f\|^2\\
&=\int_{-\infty}^{\infty}(\lambda^2-C^2)\,d\|E(\lambda)f\|^2.
\end{align*}
现在取$f=E(-C-\varepsilon)g$，则$E(\lambda)f=E(\operatorname{Min}(\lambda,\allowbreak -C-\varepsilon))g$，在$\lambda\geq-C-\varepsilon$时为常值，所以只须考察$\int_{-\infty}^{-C-\varepsilon}$。在这个范围内，$E(\lambda)f=E(\lambda)g$，且$\lambda^2-C^2\geq(C+\varepsilon)^2-C^2>2C\varepsilon$，故
\begin{gather*}
0\geq2C\varepsilon\int_{-\infty}^{-C-\varepsilon}d\|E(\lambda)g\|^2
=2C\varepsilon\|E(-C-\varepsilon)g\|^2,\\
\|E(-C-\varepsilon)g\|^2\leq0,\qquad E(-C-\varepsilon)g=0.
\end{gather*}
同样，取$f=g-E(C+\varepsilon)g$，即可证明$g-E(C+\varepsilon)g=0$。

因此，对每个$\varepsilon>0$都有$E(-C-\varepsilon)=0,\allowbreak E(C+\varepsilon)=1$；即$\lambda<-C$时$E(\lambda)=0$，$\lambda>C$时$E(\lambda)=1$。由$\hyperref[cond:II.Sbar2]{\overline{S}_2}$，后一关系在$\lambda=C$时也成立。因此$E(\lambda)$仅在$-C\leq\lambda\leq C$内变化。反过来，由此也能推出$A$连续：
\begin{align*}
\|Af\|^2&=\int_{-\infty}^{\infty}\lambda^2\,d\|E(\lambda)f\|^2\\
&=\int_{-C}^C\lambda^2\,d\|E(\lambda)f\|^2\\
&\leq C^2\int_{-C}^C d\|Ef\|^2\\
&=C^2\int_{-\infty}^{\infty}d\|E(\lambda)f\|^2=C^2\|f\|^2,\\
\|Af\|&\leq C\|f\|.
\end{align*}
可见，只在有限$\lambda$区间内变化的单位分解，恰好给出了全部连续的$A$。那么，其余不连续的厄米算符又如何呢？还有那些在任意大的$\lambda$处仍会变化的单位分解可供使用；它们是否恰好给出了这些厄米算符？

首先，必须正确估量这些算符不能处处有定义这一事实。

可以设想，厄米算符在希尔伯特空间的某些点上没有定义，尽管在那些点定义它是合理的。例如，算符$A'=\frac{h}{2\pi i}\frac{\partial}{\partial q}$对$f(q)=e^{-|q|}$没有定义；我们本来也可以把$\frac{h}{2\pi i}\frac{\partial}{\partial q}$限制在$-\infty<q<+\infty$（$q$为实数）上的解析函数上，\origNote[，][等]{97}等等。要求定义域处处稠密，固然可防止过于任意的限制；我们也还能只考察闭算符。但这仍不够有效。

\hypertarget{origpage-79}{}\phantomsection\label{origpage-79}例如，取区间$0\leq q\leq1$上的算符$A'=\frac{h}{2\pi i}\frac{\partial}{\partial q}$。设$f(q)$处处可微，且$\int_0^1|f(q)|^2\,dq,\allowbreak \int_0^1|f'(q)|^2\,dq$有限。为使$A'$厄米，还须加上边界条件$f(0):f(1)=e^{-i\alpha}$，$0\leq\alpha<2\pi$。把这些$f(q)$构成的定义域记作$\mathfrak A_\alpha$，相应的$A'$记作$A'_\alpha$。再考察边界条件$f(1)=f(0)=0$，把相应的$f(q)$构成的定义域记作$\mathfrak A^0$，算符记作$A'{}^0$。一切$A'_\alpha$都扩张$A'{}^0$；因此$A'{}^0$厄米，且定义域处处稠密。\origNote[。][于]{98}于是它们的闭包$\widetilde A'_\alpha$也都扩张$\widetilde A'{}^0$。所有$\widetilde A'_\alpha$彼此不同，也都不同于$\widetilde A'{}^0$。因为，显然为酉的运算$f(q)\mapsto e^{i\beta q}f(q)$，把$A'$变为$A'+\frac{h\beta}{2\pi}1$，把$\mathfrak A_\alpha$变为$\mathfrak A_{\alpha+\beta}$，把$\mathfrak A^0$变为$\mathfrak A^0$；所以它把$A'_\alpha$变为$A'_{\alpha-\beta}+\frac{h\beta}{2\pi}1$，把$A'{}^0$变为$A'{}^0+\frac{h\beta}{2\pi}1$；进而把$\widetilde A'_\alpha$变为$\widetilde A'_{\alpha-\beta}+\frac{h\beta}{2\pi}1$，把$\widetilde A'{}^0$变为$\widetilde A'{}^0+\frac{h\beta}{2\pi}1$。因此，若$\widetilde A'_\alpha=\widetilde A'{}^0$，就有$\widetilde A'_{\alpha-\beta}=\widetilde A'{}^0$，从而所有$\widetilde A'_\gamma$都相等。所以只须证明$\alpha\ne\gamma$时$\widetilde A'_\alpha\ne\widetilde A'_\gamma$即可。只要$A'_\alpha,\allowbreak A'_\gamma$没有共同的厄米扩张，这便成立；也就是说，只要$A'$在$\mathfrak A_\alpha,\allowbreak \mathfrak A_\gamma$的并集上不厄米即可。由于$e^{i\alpha q}$属于$\mathfrak A_\alpha$，$e^{i\gamma q}$属于$\mathfrak A_\gamma$，而
\begin{align*}
(A'e^{i\alpha q},e^{i\gamma q})-(e^{i\alpha q},A'e^{i\gamma q})
&=i\alpha\int_0^1e^{i(\alpha-\gamma)q}\,dq-i\gamma\int_0^1e^{i(\alpha-\gamma)q}\,dq\\
&=\int_0^1e^{i(\alpha-\gamma)q}i(\alpha-\gamma)\,dq
=e^{i(\alpha-\gamma)}-1\ne0,
\end{align*}
所以确实如此。因此，闭厄米算符$\widetilde A'{}^0$的定义域过窄：它有闭厄米真扩张，即$\widetilde A'_\alpha$；所谓“真”，就是说与它本身不同。而且，这种扩张有无限多种选择，因为每个$\widetilde A'_\alpha$都可采用，并给出不同的本征值问题解答：谱总是纯点谱，但依赖于$\alpha$，即由$h(\frac{\alpha}{2\pi}+k)$，$k=0,\allowbreak \pm1,\allowbreak \pm2,\allowbreak \ldots$组成的集合。反之，对$\widetilde A'{}^0$本身，则根本不能期待本征值问题有合理解答。事实上，本节稍后将证明：具有单位分解的厄米算符，即本征值问题可解的厄米算符，没有真扩张。没有真扩张的算符，亦即在一切能不破坏厄米性而合理定义它的点上都已有定义的算符，称为极大算符。由此已经看到：单位分解只能属于极大算符。

\hypertarget{origpage-80}{}\phantomsection\label{origpage-80}另一方面，可以证明，每个厄米算符都能扩张成极大厄米算符；若算符闭而不极大，总能以无限多种不同方式扩张。因此，扩张过程唯一确定的阶段，只是从$A$到闭包$\widetilde A$（参见\CJKglue \hyperlink{note-95}{原注 \getrefnumber{orig-note-95}} 末尾）。所以，我们至多能希望得到这样一种最好的结果：每个极大厄米算符都有且仅有一个属于它的单位分解。每个闭的连续厄米算符都处处有定义，所以是极大的。

现在必须检验这一期待是否正确。其中包含两个问题：单位分解的存在性与唯一性。结论将是：对于给定的极大厄米算符，或者没有属于它的单位分解，或者恰有一个；而前一种情形确实会发生。所以，本征值问题的解若存在就是唯一的，但有时不可解；后者在某种稍后要说明的意义下是例外。下面概述导出这些结论的方法。

若$f(\lambda)$为有理函数，$A$为能够经酉变换化为对角形的有限维矩阵，则$A$的本征向量仍是$f(A)$的本征向量，而相应的本征值由$\lambda_1,\allowbreak \ldots,\allowbreak \lambda_n$变为$f(\lambda_1),\allowbreak \ldots,\allowbreak f(\lambda_n)$。\origNote[。][因]{99}因此，若$f(\lambda)$把实轴映到单位圆周，本征值全为实数的矩阵就变为本征值绝对值全为$1$的矩阵，也就是厄米矩阵变为酉矩阵。\origNote[。][例]{100}例如$f(\lambda)=\frac{\lambda+i}{\lambda-i}$就有这个性质；相应的变换称为Cayley变换：
\[
U=\frac{A-i1}{A+i1},\qquad A=-i\frac{U+1}{U-1}.
\]
我们也要对厄米算符$A$尝试这个变换。\tnote[。][其][上]{上面给出的标量函数，是所显示的$U$公式中所用函数的倒数。后面的论证始终采用$U$公式的约定；两种原书形式均予保留。}其含义是：$Uf$有定义，当且仅当$f$可写成$(A+i1)\varphi=A\varphi+i\varphi$，此时令$Uf=(A-i1)\varphi=A\varphi-i\varphi$。希望这个定义对每个$f$都给出唯一的$Uf$，且$U$为酉算符。有限维情形的证明在这里不能采用，因为它预先假定能化成对角形，即有纯点谱。不过，若这一断言成立，就能按如下方法解决本征值问题。

对于酉算符$U$，本征值问题有唯一解：存在且只存在一族投影算符$E(\sigma)$（$0\leq\sigma\leq1$），满足以下条件。

\hypertarget{origpage-81}{}\phantomsection\label{origpage-81}\begin{description}
\item[$\overline{\overline S}_1$.]\leavevmode\phantomsection\label{cond:II.Sdouble1} $E(0)=0,\allowbreak E(1)=1$；当$\sigma\to\sigma_0$且$\sigma\geq\sigma_0$时，$E(\sigma)f\to E(\sigma_0)f$。
\item[$\overline{\overline S}_2$.]\leavevmode\phantomsection\label{cond:II.Sdouble2} 由$\sigma'\leq\sigma''$得$E(\sigma')\leq E(\sigma'')$。
\item[$\overline{\overline S}_3$.]\leavevmode\phantomsection\label{cond:II.Sdouble3} $(Uf,\allowbreak g)=\int_0^1e^{2\pi i\sigma}\,d(E(\sigma)f,\allowbreak g)$。
\end{description}
这里$Uf$总有意义，积分也总绝对收敛。\origNote[。][这]{101}这可在Hilbert理论范围内、用其方法证明，因为酉算符$U$总是连续的（参见\CJKglue \hyperlink{note-70}{原注 \getrefnumber{orig-note-70}}、\hyperlink{note-101}{\getrefnumber{orig-note-101}}）。它与厄米算符的表述$\hyperref[cond:II.Sbar1]{\overline S_1}$—$\hyperref[cond:II.Sbar3]{\overline S_3}$的类似之处一望即知，二者只在两点上不同：遍历实轴的实被积函数$\lambda$换成了沿单位圆周变化的复数$e^{2\pi i\sigma}$；对定义域的条件已无必要，因为酉算符处处有定义。在有限维空间$\mathfrak R_n$中，厄米矩阵与酉矩阵的关系，也已经与实轴和单位圆周的关系十分相似（参见\CJKglue \hyperlink{note-100}{原注 \getrefnumber{orig-note-100}}）。

由$\hyperref[cond:II.Sdouble1]{\overline{\overline S}_1}$，当$\sigma\to0,\allowbreak \sigma\geq0$时有$E(\sigma)f\to0$；但当$\sigma\to1,\allowbreak \sigma\leq1$时，未必有$E(\sigma)f\to f$，即$E(\sigma)$可能在$1$处不连续。不过，依 \hyperref[sec:II.4]{II.4} 的\CJKglue \hyperref[thm:II.17]{定理17}，存在左极限投影$E'$（参见\CJKglue \hyperlink{note-79}{原注 \getrefnumber{orig-note-79}}）。若$E'\ne1$，便存在非零的$f$使$E'f=0$。由于$\sigma<1$时$E(\sigma)\leq E'$，这等价于对一切$\sigma<1$都有$E(\sigma)f=0$；像 \hyperref[sec:II.8]{II.8} 开头那样可得$Uf=f$。反过来，若$Uf=f$，则
\[
\int_0^1e^{2\pi i\sigma}\,d(E(\sigma)f,f)=(Uf,f)=(f,f),
\]
因而
\[
\mathfrak R\int_0^1e^{2\pi i\sigma}\,d(E(\sigma)f,f)=(f,f),
\qquad
\int_0^1(1-\cos2\pi\sigma)\,d(E(\sigma)f,f)=0,
\]
即
\[
\int_0^1(1-\cos2\pi\sigma)\,d\|E(\sigma)f\|^2=0.
\]
于是，仍与 \hyperref[sec:II.8]{II.8} 开头一样，得到对所有$0<\sigma<1$都有$E(\sigma)f=0$，而$\sigma=0$时本来就成立。因此，在$1$处不连续，等价于$Uf=f$有非零解。但对于Cayley变换，$\varphi=Af+if,\allowbreak U\varphi=Af-if$，所以$U\varphi=\varphi$能推出$f=0,\allowbreak \varphi=0$。因此，这时从左侧令$\sigma\to1$，也有$E(\sigma)f\to f$。

\hypertarget{origpage-82}{}\phantomsection\label{origpage-82}映射
\[
\lambda=-i\frac{e^{2\pi i\sigma}+1}{e^{2\pi i\sigma}-1}=-\cot\pi\sigma,
\qquad \sigma=-\frac1\pi\operatorname{arccot}\lambda
\]
把区间$0<\sigma<1$与整个实轴$-\infty<\lambda<+\infty$一一、单调地对应起来。因此，若令
\begin{equation}
F(\lambda)=E\left(-\frac1\pi\operatorname{arccot}\lambda\right),\qquad
E=F(-\cot\pi\sigma),
\tag{C}\label{eq:II.9.C}
\end{equation}
便从$E(\sigma)$得到满足$\hyperref[cond:II.Sbar1]{\overline S_1}$、$\hyperref[cond:II.Sbar2]{\overline S_2}$的单位分解$F(\lambda)$。我们要证明：$F(\lambda)$对$A$满足$\hyperref[cond:II.Sbar3]{\overline S_3}$，当且仅当$E(\sigma)$对$U$满足$\hyperref[cond:II.Sdouble3]{\overline{\overline S}_3}$。这样，可能不连续的厄米算符$A$的本征值问题，其存在性与唯一性就都归结到相应酉算符$U$的问题，后者已按前述最好的意义得到解决。

设$A$厄米，$U$为其Cayley变换。先讨论$U$为酉算符的情形；于是，满足$\hyperref[cond:II.Sdouble1]{\overline{\overline S}_1}$—$\hyperref[cond:II.Sdouble3]{\overline{\overline S}_3}$的$E(\sigma)$存在。依 \hyperref[eq:II.9.C]{C} 构造$F(\lambda)$，则$\hyperref[cond:II.Sbar1]{\overline S_1}$、$\hyperref[cond:II.Sbar2]{\overline S_2}$已成立。若$Af$有意义，令$\varphi=Af+if$，则$U\varphi=Af-if$，所以$f=(\varphi-U\varphi)/(2i),\allowbreak Af=(\varphi+U\varphi)/2$。下面部分采用符号运算：\origNote[：][]{102}
\begin{align*}
f&=\frac1{2i}(\varphi-U\varphi)
=\frac1{2i}\left(\varphi-\int_0^1e^{2\pi i\sigma}\,dE(\sigma)\varphi\right)\\
&=\int_0^1\frac{1-e^{2\pi i\sigma}}{2i}\,dE(\sigma)\varphi,\\
E(\sigma)f&=\int_0^1\frac{1-e^{2\pi i\sigma'}}{2i}\,d(E(\sigma)E(\sigma')\varphi)\\
&=\int_0^1\frac{1-e^{2\pi i\sigma'}}{2i}\,d(E(\operatorname{Min}(\sigma,\sigma'))\varphi)\\
&=\int_0^\sigma\frac{1-e^{2\pi i\sigma'}}{2i}\,dE(\sigma')\varphi.
\end{align*}
从而
\begin{align*}
\|E(\sigma)f\|^2&=(E(\sigma)f,f)
=\int_0^\sigma\frac{1-e^{2\pi i\sigma'}}{2i}\,d(E(\sigma')\varphi,f)\\
&=\int_0^\sigma\frac{1-e^{2\pi i\sigma'}}{2i}\,d\overline{(E(\sigma')f,\varphi)}\\
&=\int_0^\sigma\frac{1-e^{2\pi i\sigma'}}{2i}\,
d\left(\int_0^{\sigma'}\frac{1-e^{-2\pi i\sigma''}}{-2i}\,
d\overline{(E(\sigma'')\varphi,\varphi)}\right)\\
&=\int_0^\sigma\frac{1-e^{2\pi i\sigma'}}{2i}
\frac{1-e^{-2\pi i\sigma'}}{-2i}\,d\overline{(E(\sigma')\varphi,\varphi)}\\
&=\int_0^\sigma\frac{(1-e^{2\pi i\sigma'})(1-e^{-2\pi i\sigma'})}{4}\,
d\|E(\sigma')\varphi\|^2\\
&=\int_0^\sigma\sin^2\pi\sigma'\,d\|E(\sigma')\varphi\|^2.
\end{align*}

\hypertarget{origpage-83}{}\phantomsection\label{origpage-83}因此，$\hyperref[cond:II.Sbar3]{\overline S_3}$中的积分为
\begin{align*}
\int_{-\infty}^{+\infty}\lambda^2\,d\|F(\lambda)f\|^2
&=\int_0^1\cot^2\pi\sigma\,d\|E(\sigma)f\|^2\\
&=\int_0^1\cot^2\pi\sigma\,
d\left(\int_0^\sigma\sin^2\pi\sigma'\,d\|E(\sigma')\varphi\|^2\right)\\
&=\int_0^1\cot^2\pi\sigma\sin^2\pi\sigma\,d\|E(\sigma)\varphi\|^2\\
&=\int_0^1\cos^2\pi\sigma\,d\|E(\sigma)\varphi\|^2.
\end{align*}
它的绝对值不超过$\int_0^1d\|E(\sigma)\varphi\|^2=\|\varphi\|^2$，所以是有限的。此外，
\begin{align*}
Af&=\frac12(\varphi+U\varphi)
=\frac12\left(\varphi+\int_0^1e^{2\pi i\sigma}\,dE(\sigma)\varphi\right)\\
&=\int_0^1\frac{1+e^{2\pi i\sigma}}2\,dE(\sigma)\varphi\\
&=\int_0^1-i\frac{e^{2\pi i\sigma}+1}{e^{2\pi i\sigma}-1}
\frac{1-e^{2\pi i\sigma}}{2i}\,dE(\sigma)\varphi\\
&=\int_0^1-\cot\pi\sigma\,
d\left(\int_0^\sigma\frac{1-e^{2\pi i\sigma'}}{2i}\,dE(\sigma')\varphi\right)\\
&=\int_0^1-\cot\pi\sigma\,dE(\sigma)f
=\int_{-\infty}^{+\infty}\lambda\,dF(\lambda)f.
\end{align*}
即$\hyperref[cond:II.Sbar3]{\overline S_3}$中的最后一个关系也成立。因此，$A$无论如何都是按$\hyperref[cond:II.Sbar3]{\overline S_3}$属于$F(\lambda)$的那个算符的扩张；但后者是极大的，这一点还要证明，所以$A$必须与它相等。\origNote[。][]{103}\tnote[][][紧]{紧前的推导证明：对$A$定义域中的$f$，$Af$等于$F(\lambda)$的谱积分。因此直接得到的是$A$为该谱积分算符的限制。原文将扩张方向写反，这里保留原说法。}

现在讨论逆向：设依$\hyperref[cond:II.Sbar1]{\overline S_1}$—$\hyperref[cond:II.Sbar3]{\overline S_3}$，$F(\lambda)$属于$A$。仍依 \hyperref[eq:II.9.C]{C} 定义$E(\sigma)$，于是$\hyperref[cond:II.Sdouble1]{\overline{\overline S}_1}$、$\hyperref[cond:II.Sdouble2]{\overline{\overline S}_2}$成立。任取$\varphi$，令
\begin{align*}
f&=\int_{-\infty}^{+\infty}\frac1{\lambda+i}\,dF(\lambda)\varphi
=\int_0^1\frac1{-\cot\pi\sigma+i}\,dE(\sigma)\varphi\\
&=\int_0^1\frac{1-e^{2\pi i\sigma}}{2i}\,dE(\sigma)\varphi.
\end{align*}
被积函数有界，所以积分收敛。与前面一样，
\begin{align*}
F(\lambda)f=E(\sigma)f
&=\int_0^1\frac{1-e^{2\pi i\sigma'}}{2i}\,dE(\sigma)E(\sigma')\varphi\\
&=\int_0^1\frac{1-e^{2\pi i\sigma'}}{2i}\,dE(\operatorname{Min}(\sigma,\sigma'))\varphi\\
&=\int_0^\sigma\frac{1-e^{2\pi i\sigma'}}{2i}\,dE(\sigma')\varphi.
\end{align*}

\hypertarget{origpage-84}{}\phantomsection\label{origpage-84}所以
\begin{align*}
Af&=\int_{-\infty}^{+\infty}\lambda\,dF(\lambda)f
=\int_0^1-\cot(\pi\sigma)\,dE(\sigma)f\\
&=\int_0^1-i\frac{e^{2\pi i\sigma}+1}{e^{2\pi i\sigma}-1}\,
d\left(\int_0^\sigma\frac{1-e^{2\pi i\sigma'}}{2i}\,dE(\sigma')\varphi\right)\\
&=\int_0^1-i\frac{e^{2\pi i\sigma}+1}{e^{2\pi i\sigma}-1}
\frac{1-e^{2\pi i\sigma}}{2i}\,dE(\sigma)\varphi
=\int_0^1\frac{1+e^{2\pi i\sigma}}2\,dE(\sigma)\varphi,
\end{align*}
从而
\[
Af+if=\int_0^1dE(\sigma)\varphi=\varphi,\qquad
Af-if=\int_0^1e^{2\pi i\sigma}\,dE(\sigma)\varphi.
\]
因此，$U\varphi$有意义，且等于$\int_0^1e^{2\pi i\sigma}\,dE(\sigma)\varphi$。$\varphi$是任意的，所以$U$处处有意义。再与任意$\psi$作内积，然后取复共轭，便知$U^*\psi$等于$\int_0^1e^{-2\pi i\sigma}\,dE(\sigma)\psi$。由 \hyperref[sec:II.8]{II.8} 末尾的计算可得$UU^*=U^*U=1$：$U$是酉的，且属于$E(\sigma)$。

因此，$A$的本征值问题可解，当且仅当其Cayley变换$U$是酉的；而此时唯一性也有保证。最后只须判定：是否总能构造$U$，它又何时为酉的？为解决这两个问题，重新从闭厄米算符$A$出发。

$U$的定义如下：$U\varphi$有意义，当且仅当$\varphi=Af+if$，此时等于$Af-if$。首先须证明，这个定义确实可行，即给定的$\varphi$不会对应于多个$f$。也就是说，$Af+if=Ag+ig$必须能推出$f=g$；由$A$的线性性，只须由$Af+if=0$能推出$f=0$。有
\begin{align*}
\|Af\pm if\|^2&=(Af\pm if,Af\pm if)\\
&=(Af,Af)\pm(Af,if)\pm(if,Af)+(if,if)\\
&=\|Af\|^2\mp i(Af,f)\pm i(Af,f)+\|f\|^2
=\|Af\|^2+\|f\|^2.
\end{align*}
所以$Af+if=0$能推出$\|f\|^2\leq\|Af+if\|^2=0$，即$f=0$；定义由此成立。其次，$\|Af+if\|=\|Af-if\|$，即$\|U\varphi\|=\|\varphi\|$。所以$U$在有定义的地方连续。再把$U$的定义域，即所有$Af+if$的集合，记作$\mathfrak E$；把它的值域，即所有$U\varphi$，也就是所有$Af-if$的集合，记作$\mathfrak F$。由$A,\allowbreak U$的线性性，$\mathfrak E,\allowbreak \mathfrak F$是线性流形；它们还是闭的。因为，若$\varphi$为$\mathfrak E$或$\mathfrak F$的聚点，就有其中的序列$\varphi_1,\allowbreak \varphi_2,\allowbreak \ldots$使$\varphi_n\to\varphi$。

\hypertarget{origpage-85}{}\phantomsection\label{origpage-85}于是$\varphi_n=Af_n\pm if_n$。由于$\varphi_n$收敛，它们满足Cauchy收敛条件（参见 \hyperref[sec:II.1]{II.1} 的 \hyperref[ax:II.D]{D}）。由
\[
\|f_m-f_n\|\leq\|A(f_m-f_n)\pm i(f_m-f_n)\|=\|\varphi_m-\varphi_n\|,
\]
可知$f_n$也满足此条件；再由
\[
\|Af_m-Af_n\|=\|A(f_m-f_n)\|
\leq\|A(f_m-f_n)\pm i(f_m-f_n)\|=\|\varphi_m-\varphi_n\|,
\]
可知$Af_m$也满足。因此，依 \hyperref[sec:II.1]{II.1} 的 \hyperref[ax:II.D]{D}，$f_1,\allowbreak f_2,\allowbreak \ldots$和$Af_1,\allowbreak Af_2,\allowbreak \ldots$都收敛：$f_n\to f,\allowbreak Af_n\to f^*$。$A$是闭的，故$Af$有意义，且等于$f^*$。于是
\[
\varphi_n=Af_n\pm if_n\to f^*\pm if=Af\pm if,\qquad
\varphi_n\to\varphi,
\]
所以$\varphi=Af\pm if$，即$\varphi$也属于相应的$\mathfrak E$或$\mathfrak F$。

因此，$U$定义在闭线性流形$\mathfrak E$上，并把它映满闭线性流形$\mathfrak F$。它是线性的，且由于$\|Uf-Ug\|=\|U(f-g)\|=\|f-g\|$，它保持一切距离不变，我们称之为保长映射。于是$f\ne g$必有$Uf\ne Ug$，所以映射是一一的。还有$(f,\allowbreak g)=(Uf,\allowbreak Ug)$；证明与 \hyperref[sec:II.5]{II.5} 中酉算符情形完全相同。因此$U$也保持一切内积不变。不过，显然只有当$\mathfrak E=\mathfrak F=\mathfrak R_\infty$时，$U$才是酉的，反之亦然。

现设$A,\allowbreak B$为两个闭厄米算符，$U,\allowbreak V$为它们的Cayley变换，相应的上述集合分别为$\mathfrak E,\allowbreak \mathfrak F$及$\mathfrak G,\allowbreak \mathfrak H$。立即可见，若$B$是$A$的真扩张，$V$也是$U$的真扩张：$\mathfrak E$是$\mathfrak G$的真子集，$\mathfrak F$是$\mathfrak H$的真子集。因此$\mathfrak E\ne\mathfrak R_\infty,\allowbreak \mathfrak F\ne\mathfrak R_\infty$，于是$U$不酉，$A$的本征值问题不可解。这样就证明了前面多次引用的定理：$A$的本征值问题若可解，$A$就没有真扩张，即$A$为极大的。

回到单个闭厄米算符$A$及其$\mathfrak E,\allowbreak \mathfrak F,\allowbreak U$。若$Af$有意义，则对$\varphi=Af+if$，$U\varphi$有意义且等于$Af-if$。所以$f=\frac1{2i}(\varphi-U\varphi),\allowbreak Af=\frac12(\varphi+U\varphi)$；或令$\psi=\varphi/(2i)$，就有$f=\psi-U\psi,\allowbreak Af=i(\psi+U\psi)$。反过来，若$f=\psi-U\psi$，$Af$必定有意义：因为$U\psi$有意义，所以$\psi=Af'+if'$，其中$Af'$有意义，而$U\psi=Af'-if'$，故$f=\psi-U\psi=2if'$。因此，$A$的定义域是一切$\psi-U\psi$的集合，且当$f=\psi-U\psi$时，$Af=i(\psi+U\psi)$。所以，$U$及其$\mathfrak E,\allowbreak \mathfrak F$也唯一确定$A$。同时可见，一切$\psi-U\psi$必须处处稠密，因为它们构成$A$的定义域。

反过来，从两个闭线性流形$\mathfrak E,\allowbreak \mathfrak F$，以及把$\mathfrak E$映满$\mathfrak F$的线性保长映射$U$出发，是否存在一个厄米算符$A$，以这个$U$为其Cayley变换？

\hypertarget{origpage-86}{}\phantomsection\label{origpage-86}无论如何，$\psi-U\psi$必须处处稠密，所以先作此假定。根据刚才的讨论，预期的$A$便已唯一确定。剩下的问题是：这个定义是否可行？$A$是否确实厄米？$U$是否确实为它的Cayley变换？要解决第一点，只须证明在$f=\varphi-U\varphi$中，$f$唯一确定$\varphi$（若这样的$\varphi$存在）。也就是说，$\varphi-U\varphi=\psi-U\psi$应能推出$\varphi=\psi$，或者只须$\varphi-U\varphi=0$能推出$\varphi=0$。事实上，设$\varphi-U\varphi=0$；则对$g=\psi-U\psi$有
\begin{align*}
(\varphi,g)&=(\varphi,\psi)-(\varphi,U\psi)\\
&=(U\varphi,U\psi)-(\varphi,U\psi)
=(U\varphi-\varphi,U\psi)=0;
\end{align*}
由于这些$g$处处稠密，便得$\varphi=0$。

第二点须证明$(Af,\allowbreak g)=(f,\allowbreak Ag)$，即交换$f,\allowbreak g$时，$(Af,\allowbreak g)$变成它的复共轭。设$f=\varphi-U\varphi,\allowbreak g=\psi-U\psi$，则$Af=i(\varphi+U\varphi)$，而
\begin{align*}
(Af,g)&=(i(\varphi+U\varphi),\psi-U\psi)\\
&=i(\varphi,\psi)+i(U\varphi,\psi)-i(\varphi,U\psi)-i(U\varphi,U\psi)\\
&=i\big[(U\varphi,\psi)-\overline{(U\psi,\varphi)}\big]
=i(U\varphi,\psi)+\overline{i(U\psi,\varphi)}.
\end{align*}
显然，交换$f,\allowbreak g$，即交换$\varphi,\allowbreak \psi$时，这个式子具有所要求的性质。第三点如下看出：设$A$的Cayley变换为$V$，其定义域是所有
\[
Af+if=i(\varphi+U\varphi)+i(\varphi-U\varphi)=2i\varphi
\]
组成的集合，也就是$U$的定义域；在此定义域上有
\begin{align*}
V(2i\varphi)&=V(Af+if)=Af-if\\
&=i(\varphi+U\varphi)-i(\varphi-U\varphi)=2iU\varphi.
\end{align*}
即$V\varphi=U\varphi$，所以$V=U$。

因此，给每个闭厄米算符$A$指定其Cayley变换$U$，便使这些$A$，与那些使$\varphi-U\varphi$处处稠密的线性保长映射$U$，一一对应。\origNote[。][现]{104}现在很容易综观$A$的一切厄米扩张$B$，因为$U$的一切保长扩张$V$都不难找出。$\varphi-V\varphi$自然处处稠密，因为它们已包含处处稠密的子集$\varphi-U\varphi$。所以，$A$极大当且仅当$U$极大。若$U$不极大，则$\mathfrak E\ne\mathfrak R_\infty,\allowbreak \mathfrak F\ne\mathfrak R_\infty$；反过来，这两式也能推出$U$不极大。因为此时$\mathfrak R_\infty-\mathfrak E\ne0,\allowbreak \mathfrak R_\infty-\mathfrak F\ne0$，故能从$\mathfrak R_\infty-\mathfrak E$选出非零$\varphi_0$，从$\mathfrak R_\infty-\mathfrak F$选出非零$\psi_0$；再分别除以$\|\varphi_0\|,\allowbreak \|\psi_0\|$，甚至可使$\|\varphi_0\|=\|\psi_0\|=1$。在$[\mathfrak E,\allowbreak \varphi_0]$上定义算符$V$如下：若$f=\varphi+a\varphi_0$，其中$\varphi$属于$\mathfrak E$、$a$为数，就令$Vf=U\varphi+a\psi_0$。显然$V$线性；又因$\varphi$与$\varphi_0$正交，$U\varphi$与$\psi_0$正交，有

\hypertarget{origpage-87}{}\phantomsection\label{origpage-87}\[
\|f\|^2=\|\varphi\|^2+|a|^2,\qquad
\|Vf\|^2=\|U\varphi\|^2+|a|^2,
\]
所以$\|Vf\|=\|f\|$，即$V$保长。最后，$V$是$U$的真扩张。因此，$A$极大的特征条件是：$\mathfrak E=\mathfrak R_\infty$或$\mathfrak F=\mathfrak R_\infty$。

反之，若$A$不极大，则两个闭线性流形$\mathfrak R_\infty-\mathfrak E,\allowbreak \mathfrak R_\infty-\mathfrak F$都非零。设张成它们的规范正交系分别为$\varphi_1,\allowbreak \ldots,\allowbreak \varphi_p$与$\psi_1,\allowbreak \ldots,\allowbreak \psi_q$（见 \hyperref[sec:II.2]{II.2} 的\CJKglue \hyperref[thm:II.9]{定理9}）。这里$p,\allowbreak q=1,\allowbreak 2,\allowbreak \ldots,\allowbreak \infty$；$p=\infty$或$q=\infty$时，相应的$\varphi$或$\psi$序列不终止。令$r=\operatorname{Min}(p,\allowbreak q)$，在$[\mathfrak E,\allowbreak \varphi_1,\allowbreak \ldots,\allowbreak \varphi_r]$上如下定义$V$：当$f=\varphi+\sum_{\nu=1}^ra_\nu\varphi_\nu$，其中$\varphi$属于$\mathfrak E$、$a_1,\allowbreak \ldots,\allowbreak a_r$为数时，令$Vf=U\varphi+\sum_{\nu=1}^ra_\nu\psi_\nu$。仍不难看出，$V$线性、保长，并且是$U$的真扩张。它的定义域为$[\mathfrak E,\allowbreak \varphi_1,\allowbreak \ldots,\allowbreak \varphi_r]$，当$r=p$时等于$[\mathfrak E,\allowbreak \mathfrak R_\infty-\mathfrak E]=\mathfrak R_\infty$；值域为$[\mathfrak F,\allowbreak \psi_1,\allowbreak \ldots,\allowbreak \psi_r]$，当$r=q$时等于$[\mathfrak F,\allowbreak \mathfrak R_\infty-\mathfrak F]=\mathfrak R_\infty$。因此，两者中至少有一个必等于$\mathfrak R_\infty$。设$V$为厄米算符$B$的Cayley变换，依上述结论，$B$扩张$A$而且极大。须注意，$\varphi,\allowbreak \psi$的选择有无限多种，例如可把$\psi_1$换成任意$\theta\psi_1$，$|\theta|=1$；所以$V,\allowbreak B$也有无限多种选择。

至此，本征值问题已讨论完毕，结论如下：只要可解，解就只有一个；非极大算符的本征值问题则一定不可解。非极大算符总能以无限多种不同方式扩张成极大算符。这里始终说的是闭厄米算符。不过，极大性的条件与本征值问题可解的条件并不完全相同：前者是$\mathfrak E=\mathfrak R_\infty$\textbf{或}$\mathfrak F=\mathfrak R_\infty$，后者却是$\mathfrak E=\mathfrak R_\infty$\textbf{且}$\mathfrak F=\mathfrak R_\infty$。

对于满足前一个条件却不满足后一个条件的算符，这里不作进一步研究。它们的本征值问题不可解，而且又因极大性而不能作真扩张，因而也无法用扩张来改变这一状况。它们的特征是$\mathfrak E=\mathfrak R_\infty,\allowbreak \mathfrak F\ne\mathfrak R_\infty$，或$\mathfrak E\ne\mathfrak R_\infty,\allowbreak \mathfrak F=\mathfrak R_\infty$。这样的算符确实存在，且全都可由两个简单的标准形得到；所以，与本征值问题可解的极大算符相比，可以把它们看作例外。详情见\CJKglue \hyperlink{note-95}{原注 \getrefnumber{orig-note-95}} 所引作者本人的论文。无论如何，暂且必须把这些算符排除在量子力学的讨论之外。因为，如下文将看到的，属于厄米算符的单位分解，在量子力学各个概念的构造中都起着如此重要的作用，\hypertarget{origpage-88}{}\phantomsection\label{origpage-88}以至于我们不能放弃它的存在，也就不能放弃本征值问题的可解性。\origNote[。][因]{105}因此，通常只允许本征值问题可解的厄米算符。由于这个性质加强了极大性，这类算符称为超极大算符；\origNote[；][恰]{106}恰好是它们与单位分解一一对应。

最后，提到两类一定超极大的闭厄米算符。第一类是连续算符：它们处处有定义，所以极大；又因根据Hilbert理论，其本征值问题可解（参见\CJKglue \hyperlink{note-70}{原注 \getrefnumber{orig-note-70}} 所引文献），因而甚至超极大。第二类是在$\mathfrak R_\infty$的某个实现中为实的算符，前提是它们极大。因为$\mathfrak E,\allowbreak \mathfrak F$的定义只差$i$的符号，当其他一切都为实时，这个差别就无关紧要。因此，$\mathfrak E=\mathfrak R_\infty$能推出$\mathfrak F=\mathfrak R_\infty$，反之亦然；即极大性推出超极大性。不假定极大性，至少也能说$\mathfrak R_\infty-\mathfrak E$与$\mathfrak R_\infty-\mathfrak F$的维数相同。于是，用前面研究扩张时的记号，有$p=q$，所以$r=p=q$，且
\begin{align*}\relax
[\mathfrak E,\varphi_1,\ldots,\varphi_r]&=[\mathfrak E,\mathfrak R_\infty-\mathfrak E]=\mathfrak R_\infty,\\
[\mathfrak F,\psi_1,\ldots,\psi_r]&=[\mathfrak F,\mathfrak R_\infty-\mathfrak F]=\mathfrak R_\infty.
\end{align*}
即当时所得的扩张超极大。因此，实算符总有超极大扩张。\CJKglue \hyperlink{note-95}{原注 \getrefnumber{orig-note-95}} 所引论文表明，一切定号算符也有同样性质。

\section{对易算符}\label{sec:II.10}
依 \hyperref[sec:II.4]{II.4} 给出的定义，两个算符$R,\allowbreak S$对易，是说$RS=SR$；当它们不都处处有定义时，还须两边的定义域相同。先限于厄米算符，并且为了避开定义域的困难，先限于处处有定义、因而连续的算符。同时考察属于$R,\allowbreak S$的单位分解$E(\lambda),\allowbreak F(\lambda)$。

$R,\allowbreak S$对易，是说对一切$f,\allowbreak g$有$(RSf,\allowbreak g)=(SRf,\allowbreak g)$，即$(Sf,\allowbreak Rg)=(Rf,\allowbreak Sg)$。此外，由$R,\allowbreak S$对易，还能推出$R^n,\allowbreak S$（$n=0,\allowbreak 1,\allowbreak 2,\allowbreak \ldots$）对易，进而对每个多项式$p(x)=a_0+a_1x+\cdots+a_nx^n$，$p(R),\allowbreak S$也对易。按符号记法有
\[
R=\int_{-C}^C\lambda\,dE(\lambda),\qquad
s(R)=\int_{-C}^Cs(\lambda)\,dE(\lambda).
\]

\hypertarget{origpage-89}{}\phantomsection\label{origpage-89}这里$C$是 \hyperref[sec:II.9]{II.9} 为连续算符$R$（该节记作$A$）引入的常数；$s(x)$为任意函数（见 \hyperref[sec:II.8]{II.8}，特别是\CJKglue \hyperlink{note-94}{原注 \getrefnumber{orig-note-94}}）。当$s(x)$为多项式时，$(s(R)f,\allowbreak Sg)=(Sf,\allowbreak s(R)g)$，所以
\begin{equation}
\int_{-C}^Cs(\lambda)\,d(E(\lambda)f,Sg)
=\int_{-C}^Cs(\lambda)\,d(Sf,E(\lambda)g).
\tag{*}\label{eq:II.10.star}
\end{equation}
每个连续函数$s(x)$都能在$-C\leq x\leq C$上一致地用多项式任意精确逼近，因此\CJKglue \hyperref[eq:II.10.star]{（*）}\CJKglue 对连续的$s(x)$也成立。现取
\[
s(x)=\begin{cases}\lambda_0-x,&x\leq\lambda_0,\\0,&x\geq\lambda_0.\end{cases}
\]
由\CJKglue \hyperref[eq:II.10.star]{（*）}\CJKglue 得
\[
\int_{-C}^{\lambda_0}(\lambda_0-\lambda)\,d(E(\lambda)f,Sg)
=\int_{-C}^{\lambda_0}(\lambda_0-\lambda)\,d(Sf,E(\lambda)g).
\]
把$\lambda_0$换成$\lambda_0+\varepsilon$（$\varepsilon>0$），相减后除以$\varepsilon$，得到
\begin{align*}
&\int_{-C}^{\lambda_0}d(E(\lambda)f,Sg)
+\int_{\lambda_0}^{\lambda_0+\varepsilon}\frac{\lambda-\lambda_0}{\varepsilon}\,d(E(\lambda)f,Sg)\\
&\qquad=\int_{-C}^{\lambda_0}d(Sf,E(\lambda)g)
+\int_{\lambda_0}^{\lambda_0+\varepsilon}\frac{\lambda-\lambda_0}{\varepsilon}\,d(Sf,E(\lambda)g).
\end{align*}
令$\varepsilon\to0$，并注意$\hyperref[cond:II.Sbar1]{\overline S_1}$，便得
\begin{gather*}
\int_{-C}^{\lambda_0}d(E(\lambda)f,Sg)=\int_{-C}^{\lambda_0}d(Sf,E(\lambda)g),\\
(E(\lambda_0)f,Sg)=(Sf,E(\lambda_0)g).
\end{gather*}
所以，所有$-C\leq\lambda_0\leq C$的$E(\lambda_0)$都与$S$对易。其余当然更是如此，因为$\lambda_0<-C$或$\lambda_0>C$时，$E(\lambda_0)$分别为$0$或$1$。因此，若$R$与$S$对易，一切$E(\lambda)$也都与$S$对易。反过来，若所有$E(\lambda)$与$S$对易，\CJKlinkpunctkern{，}{（}\hyperref[eq:II.10.star]{（*）}\CJKglue 便对任意函数$s(x)$成立，从而一切$s(R)$也与$S$对易。由此可得两点：第一，$R$与$S$对易，当且仅当一切$E(\lambda)$与$S$对易；第二，这时$R$的一切函数，即各个$s(R)$，也都与$S$对易。

而$E(\lambda)$与$S$对易，当且仅当它与所有$F(\mu)$对易：只须把刚才的定理应用于$S,\allowbreak E(\lambda)$，以它们代替$R,\allowbreak S$即可。因此，$R,\allowbreak S$对易的另一个特征条件，就是所有$E(\lambda)$与所有$F(\mu)$对易。此外，由上述结论，$R,\allowbreak S$对易能推出$r(R),\allowbreak S$对易；再以$S,\allowbreak r(R)$代替$R,\allowbreak S$，又能推出$r(R),\allowbreak s(S)$对易。

\hypertarget{origpage-90}{}\phantomsection\label{origpage-90}若不要求厄米算符$R,\allowbreak S$连续，情形就更复杂，因为$RS,\allowbreak SR$的定义域可能很难掌握。例如，$R\cdot0$总有意义：$0f=0$，$R\cdot0f=R(0f)=R0=0$。而$0\cdot R$只有在$R$有意义时才有意义（参见 \hyperref[sec:II.5]{II.5} 的讨论）。因此，若$R$并非处处有意义，由于定义域不同，便有$R\cdot0\ne0\cdot R$；严格说来，$R$和$0$不对易。这对后面的用途很不方便：$0$应当与一切厄米算符对易，而不应仅与连续厄米算符对易。\origNote[。][所]{107}所以，对于不连续的$R,\allowbreak S$，要另行定义对易性；此处仍限于超极大的$R,\allowbreak S$，因为按照 \hyperref[sec:II.9]{II.9}，只有它们是我们关心的对象。定义如下：若属于$R,\allowbreak S$的单位分解$E(\lambda),\allowbreak F(\mu)$中，所有$E(\lambda)$都与所有$F(\mu)$按原来的意义对易，就称$R,\allowbreak S$按新的意义对易。我们知道，$R,\allowbreak S$连续时，新定义与旧定义相同；但$R,\allowbreak S$中有一个或两个不连续时，就未必相同。例如$R,\allowbreak 0$按旧定义不对易，按新定义却对易，因为对于$0$，每个$F(\mu)$都是$0$或$1$，\origNote[，][因]{108}因而与一切$E(\lambda)$对易。

上面已证明：若$R,\allowbreak S$为两个对易的连续厄米算符，$R$的每个函数$r(R)$都与$S$的每个函数$s(S)$对易。取$R=S$时前提总成立，所以同一个算符的两个函数$r(R),\allowbreak s(R)$总对易。也可由 \hyperref[sec:II.8]{II.8} 末尾的乘法公式看出这一点：$r(R)s(R)=t(R)$，其中$r(x)s(x)\equiv t(x)$。此外，若$r(x),\allowbreak s(x)$为实函数，$r(R),\allowbreak s(R)$就厄米；因为依 \hyperref[sec:II.8]{II.8}，当$r(x)$为实时有$(r(R))^*=\overline r(R)=r(R)$。

反过来也成立：若$A,\allowbreak B$为两个对易厄米算符，就存在厄米算符$R$，使两者都是它的函数，即$A=r(R),\allowbreak B=s(R)$。甚至还有更强的结论：给定任意一组有限或无限多个对易厄米算符$A,\allowbreak B,\allowbreak C,\allowbreak \ldots$，总存在一个厄米算符$R$，使所有$A,\allowbreak B,\allowbreak C,\allowbreak \ldots$都是它的函数。这里无法证明这个定理，只能请读者参阅有关文献。\origNote[。][就]{109}就本书的用途而言，只需考虑有限多个具有纯点谱的$A,\allowbreak B,\allowbreak C,\allowbreak \ldots$；下面就对这种情形加以证明，至于一般情形，只作几句初步说明。

设$A,\allowbreak B,\allowbreak C,\allowbreak \ldots$为有限多个具有纯点谱的厄米算符。对任意数$\lambda$，把$Af=\lambda f$的全部解张成的闭线性流形记作$\mathfrak L_\lambda$，其投影算符记作$E_\lambda$。

\hypertarget{origpage-91}{}\phantomsection\label{origpage-91}$\lambda$为$A$的点本征值，当且仅当存在非零解$f$，即$\mathfrak L_\lambda\ne(0)$，或$E_\lambda\ne0$。相应地，对$B$构造$\mathfrak M_\lambda,\allowbreak F_\lambda$，对$C$构造$\mathfrak N_\lambda,\allowbreak G_\lambda$，依此类推。由$Af=\lambda f$得$ABf=BAf=B(\lambda f)=\lambda Bf$，所以$f$属于$\mathfrak M_\lambda$时，$Bf$也属于它。由于$E_\lambda f$总属于$\mathfrak M_\lambda$，$BE_\lambda f$也属于它，所以$E_\lambda BE_\lambda f=BE_\lambda f$。此式恒成立，故$E_\lambda BE_\lambda=BE_\lambda$；施行$*$又得$E_\lambda BE_\lambda=E_\lambda B$，所以$E_\lambda B=BE_\lambda$。正如由$A,\allowbreak B$对易推出了$B,\allowbreak E_\lambda$对易，也能由$B,\allowbreak E_\lambda$对易推出$E_\lambda,\allowbreak F_\mu$对易。由于$A,\allowbreak B$在$A,\allowbreak B,\allowbreak C,\allowbreak \ldots$中没有任何特殊地位，可知所有$E_\lambda,\allowbreak F_\mu,\allowbreak G_\nu,\allowbreak \ldots$彼此对易。因此，$K(\lambda\mu\nu\cdots)=E_\lambda F_\mu G_\nu\cdots$也是投影算符，设相应的闭线性流形为$\mathfrak K(\lambda\mu\nu\cdots)$。依 \hyperref[sec:II.4]{II.4} 的\CJKglue \hyperref[thm:II.14]{定理14}，$\mathfrak K(\lambda\mu\nu\cdots)$是$\mathfrak L_\lambda,\allowbreak \mathfrak M_\mu,\allowbreak \mathfrak N_\nu,\allowbreak \ldots$的交集，即方程组
\[
Af=\lambda f,\qquad Bf=\mu f,\qquad Cf=\nu f,\qquad\ldots
\]
全部公共解的集合。设$\lambda,\allowbreak \mu,\allowbreak \nu,\allowbreak \ldots$与$\lambda',\allowbreak \mu',\allowbreak \nu',\allowbreak \ldots$为两组不同的数，即$\lambda\ne\lambda'$，或$\mu\ne\mu'$，或$\nu\ne\nu'$，等等。设$f$属于$\mathfrak K(\lambda\mu\nu\cdots)$，$f'$属于$\mathfrak K(\lambda'\mu'\nu'\cdots)$，则$f,\allowbreak f'$正交：当$\lambda\ne\lambda'$时，由$Af=\lambda f,\allowbreak Af'=\lambda'f'$可知；当$\mu\ne\mu'$时，由$Bf=\mu f,\allowbreak Bf'=\mu'f'$可知；其余同理。所以，整个$\mathfrak K(\lambda\mu\nu\cdots)$与整个$\mathfrak K(\lambda'\mu'\nu'\cdots)$正交。

$A$具有纯点谱，所以各$\mathfrak L_\lambda$所张成的闭线性流形是整个$\mathfrak R_\infty$。非零的$f$不可能与所有$\mathfrak L_\lambda$正交，即至少在某个$\mathfrak L_\lambda$上的投影不为零，也就是$E_\lambda f\ne0$。同样，必有某个$\mu$使$F_\mu f\ne0$，某个$\nu$使$G_\nu f\ne0$，等等。因此，对每个非零$f$，先找$\lambda$使$E_\lambda f\ne0$，再找$\mu$使$F_\mu(E_\lambda f)\ne0$，再找$\nu$使$G_\nu(F_\mu(E_\lambda f))\ne0$，依此类推。最后得到$\cdots G_\nu F_\mu E_\lambda f\ne0$，即$E_\lambda F_\mu G_\nu\cdots f\ne0$，即$K(\lambda\mu\nu\cdots)f\ne0$；所以$f$并不与$\mathfrak K(\lambda\mu\nu\cdots)$正交。换言之，一个与所有$\mathfrak K(\lambda\mu\nu\cdots)$正交的$f$，必为零。因此，这些$\mathfrak K(\lambda\mu\nu\cdots)$合起来所张成的闭线性流形是整个$\mathfrak R_\infty$。

现设$\varphi^{(1)}_{(\lambda\mu\nu\cdots)},\allowbreak \varphi^{(2)}_{(\lambda\mu\nu\cdots)},\allowbreak \ldots$为张成闭线性流形$\mathfrak K(\lambda\mu\nu\cdots)$的规范正交系。这个序列是否在有限项后终止，取决于$\mathfrak K(\lambda\mu\nu\cdots)$是有限维还是无限维；若$\mathfrak K(\lambda\mu\nu\cdots)=0$，序列就没有项。每个$\varphi^{(n)}_{(\lambda\mu\nu\cdots)}$都属于某个$\mathfrak K(\lambda\mu\nu\cdots)$，因此同时是一切$A,\allowbreak B,\allowbreak C,\allowbreak \ldots$的本征函数。任意两个不同的$\varphi^{(n)}_{(\lambda\mu\nu\cdots)}$都正交：若它们的$\lambda,\allowbreak \mu,\allowbreak \nu,\allowbreak \ldots$相同，由构造即知；若这些数不同，则因为它们属于不同的$\mathfrak K(\lambda\mu\nu\cdots)$。所有$\varphi^{(n)}_{(\lambda\mu\nu\cdots)}$与所有$\mathfrak K(\lambda\mu\nu\cdots)$张成同一个闭线性流形，即$\mathfrak R_\infty$。因此，它们构成完全规范正交系。

\hypertarget{origpage-92}{}\phantomsection\label{origpage-92}这样就构造了一个完全规范正交系，全部由$A,\allowbreak B,\allowbreak C,\allowbreak \ldots$的共同本征函数组成。以下把它改记为$\psi_1,\allowbreak \psi_2,\allowbreak \ldots$，相应的本征值方程写为
\[
A\psi_m=\lambda_m\psi_m,\qquad
B\psi_m=\mu_m\psi_m,\qquad
C\psi_m=\nu_m\psi_m,\qquad\ldots.
\]
现任取互不相同的数$\varkappa_1,\allowbreak \varkappa_2,\allowbreak \varkappa_3,\allowbreak \ldots$，构造以$\varkappa_1,\allowbreak \varkappa_2,\allowbreak \ldots$为纯点谱、以$\psi_1,\allowbreak \psi_2,\allowbreak \ldots$为相应本征函数的厄米算符$R$。即\origNote[即][]{110}
\[
R\left(\sum_{m=1}^\infty x_m\psi_m\right)
=\sum_{m=1}^\infty x_m\varkappa_m\psi_m.
\]
设$F(\varkappa)$定义在$-\infty<\varkappa<+\infty$上，满足$F(\varkappa_m)=\lambda_m$，$m=1,\allowbreak 2,\allowbreak \ldots$，其余$\varkappa$处的值任意。类似地，令$G(\varkappa_m)=\mu_m,\allowbreak H(\varkappa_m)=\nu_m$，等等。要证明
\[
A=F(R),\qquad B=G(R),\qquad C=H(R),\qquad\ldots.
\]
为此须证明：若$R$具有纯点谱$\varkappa_1,\allowbreak \varkappa_2,\allowbreak \ldots$和本征函数$\psi_1,\allowbreak \psi_2,\allowbreak \ldots$，则$F(R)$具有纯点谱$F(\varkappa_1),\allowbreak F(\varkappa_2),\allowbreak \ldots$，而本征函数仍为$\psi_1,\allowbreak \psi_2,\allowbreak \ldots$。这些函数本来已构成完全规范正交系，所以只须证明$F(R)\psi_m=F(\varkappa_m)\cdot\psi_m$。

像 \hyperref[sec:II.8]{II.8} 一样，设属于$R$的单位分解为$E(\lambda)=\sum_{\varkappa_m\leq\lambda}P_{[\psi_m]}$。按符号记法，有
\[
R=\int\lambda\,dE(\lambda),
\]
而依定义有
\[
F(R)=\int F(\lambda)\,dE(\lambda).
\]
此外，
\[
E(\lambda)\psi_m=\begin{cases}\psi_m,&\varkappa_m\leq\lambda,\\0,&\varkappa_m>\lambda.\end{cases}
\]
所以
\[
(F(R)\psi_m,g)=\int F(\lambda)\,d(E(\lambda)\psi_m,g)
=F(\varkappa_m)\cdot(\psi_m,g)
\]
对每个$g$成立，因此确有$F(R)\psi_m=F(\varkappa_m)\cdot\psi_m$。

这样，纯点谱的情形就按前述计划解决了。对于连续谱，只能请读者参阅\CJKglue \hyperlink{note-109}{原注 \getrefnumber{orig-note-109}} 所引文献；这里只指出一个特别有代表性的例子。

设$\mathfrak R_\infty$由使$\iint|f(q_1q_2)|^2\,dq_1\,dq_2$有限的一切$f(q_1q_2)$构成，$q_1,\allowbreak q_2$的变化区域为正方形$0\leq q_1,\allowbreak q_2\leq1$。构造算符$A=q_1\cdots,\allowbreak B=q_2\cdots$；它们厄米，在此$q_1,\allowbreak q_2$范围内连续（在$-\infty<q_1,\allowbreak q_2<+\infty$时则不连续），而且对易。因此，两者必定都是某个$R$的函数。

\hypertarget{origpage-93}{}\phantomsection\label{origpage-93}这个$R$于是与$A,\allowbreak B$对易，由此可推出——这里不详述——$R$必形如$s(q_1,\allowbreak q_2)\cdots$，其中$s(q_1,\allowbreak q_2)$为有界函数。因此$R^n$（$n=0,\allowbreak 1,\allowbreak 2,\allowbreak \ldots$）等于$(s(q_1,\allowbreak q_2))^n\cdots$；当$F(\varkappa)$为多项式时，$F(R)$等于$F(s_1(q_1,\allowbreak q_2))\cdots$。这个公式可以推广到所有$F(\varkappa)$，不过这里也无法详述。因此，由$F(R)=A,\allowbreak G(R)=B$得\origNote[得][]{111}
\[
F(s(q_1,q_2))=q_1,\qquad G(s(q_1,q_2))=q_2.
\]
所以，互为逆的映射$s(q_1,\allowbreak q_2)=\varkappa$与$F(\varkappa)=q_1,\allowbreak G(\varkappa)=q_2$，应当把正方形$0\leq q_1,\allowbreak q_2\leq1$与数轴上的$\varkappa$点集一一对应起来；这与几何直观相矛盾。

然而，前面引用的证明告诉我们，这必定可能；事实上，所谓Peano曲线确实给出了所需的映射。\origNote[。][仔]{112}仔细考察\CJKglue \hyperlink{note-109}{原注 \getrefnumber{orig-note-109}} 所引证明，也确实可见，在当前例子中，它导向Peano曲线或与之相近的构造。

\section{迹}\label{sec:II.11}
这里定义算符的若干重要不变量。

对于$\mathfrak R_n$中的矩阵$\{a_{\mu\nu}\}$，迹$\sum_{\mu=1}^na_{\mu\mu}$就是这样一个不变量。它具有酉不变性，即把$\{a_{\mu\nu}\}$变换到另一个直角坐标系，迹仍不变。\origNote[。][若]{113}若把矩阵$\{a_{\mu\nu}\}$换成相应的算符
\[
A\{x_1,\ldots,x_n\}=\{y_1,\ldots,y_n\},\qquad
y_\mu=\sum_{\nu=1}^n a_{\mu\nu}x_\nu,
\]
则$a_{\mu\nu}$可用$A$如下表达：向量$\varphi_1=\{1,\allowbreak 0,\allowbreak \ldots,\allowbreak 0\},\allowbreak \varphi_2=\{0,\allowbreak 1,\allowbreak \ldots,\allowbreak 0\},\allowbreak \ldots,\allowbreak \varphi_n=\{0,\allowbreak 0,\allowbreak \ldots,\allowbreak 1\}$构成完全规范正交系，显然$a_{\mu\nu}=(A\varphi_\nu,\allowbreak \varphi_\mu)$（见 \hyperref[sec:II.5]{II.5}，特别是\CJKglue \hyperlink{note-60}{原注 \getrefnumber{orig-note-60}}）。因此迹就是$\sum_{\mu=1}^n(A\varphi_\mu,\allowbreak \varphi_\mu)$，它的酉不变性意味着，对任意完全规范正交系，所得值相同。

在$\mathfrak R_\infty$中，可以立即照此定义一个类似的概念。设$A$为线性算符，任取完全规范正交系$\varphi_1,\allowbreak \varphi_2,\allowbreak \ldots$，使所有$A\varphi_\mu$都有意义。只要$A$的定义域处处稠密，这当然做得到：依 \hyperref[sec:II.2]{II.2} 的\CJKglue \hyperref[thm:II.8]{定理8}，把该定义域中的一个稠密序列$f_1,\allowbreak f_2,\allowbreak \ldots$正交化就够了。令$\operatorname{Spur}(A)=\sum_{\mu=1}^\infty(A\varphi_\mu,\allowbreak \varphi_\mu)$。须证明，它确实只依赖于$A$，不依赖于$\varphi_\mu$。

\hypertarget{origpage-94}{}\phantomsection\label{origpage-94}为此，先引入两个完全规范正交系$\varphi_1,\allowbreak \varphi_2,\allowbreak \ldots$及$\psi_1,\allowbreak \psi_2,\allowbreak \ldots$，并令
\[
\operatorname{Spur}(A;\varphi,\psi)
=\sum_{\mu,\nu=1}^\infty(A\varphi_\mu,\psi_\nu)(\psi_\nu,\varphi_\mu).
\]
依 \hyperref[sec:II.2]{II.2} 的\CJKglue \hyperref[thm:II.7]{定理7} 之$\gamma$，它等于$\sum_{\mu=1}^\infty(A\varphi_\mu,\allowbreak \varphi_\mu)$，所以只是表面上依赖$\psi_\nu$。又有
\begin{align*}
\sum_{\mu,\nu=1}^\infty(A\varphi_\mu,\psi_\nu)(\psi_\nu,\varphi_\mu)
&=\sum_{\mu,\nu=1}^\infty(\varphi_\mu,A^*\psi_\nu)(\psi_\nu,\varphi_\mu)\\
&=\overline{\sum_{\mu,\nu=1}^\infty(A^*\psi_\nu,\varphi_\mu)(\varphi_\mu,\psi_\nu)}.
\end{align*}
即$\operatorname{Spur}(A;\varphi,\allowbreak \psi)=\overline{\operatorname{Spur}(A^*;\psi,\allowbreak \varphi)}$。根据刚才所述，右边对$\varphi_\mu$的依赖只是表面的，所以左边也如此。因此，它对$\varphi_\mu$及$\psi_\nu$的依赖都只是表面的，实际上仅由$A$决定。于是可以把$\operatorname{Spur}(A;\varphi,\allowbreak \psi)$记作$\operatorname{Spur}(A)$。它又等于$\sum_{\mu=1}^\infty(A\varphi_\mu,\allowbreak \varphi_\mu)$，因此所要求的不变性得证。由最后一个等式还有$\operatorname{Spur}(A)=\overline{\operatorname{Spur}(A^*)}$。

关系
\[
\operatorname{Spur}(aA)=a\operatorname{Spur}(A),\qquad
\operatorname{Spur}(A\pm B)=\operatorname{Spur}(A)\pm\operatorname{Spur}(B)
\]
是显然的。此外，即使$A,\allowbreak B$不对易，也有
\[
\operatorname{Spur}(AB)=\operatorname{Spur}(BA).
\]
证明如下：
\begin{align*}
\operatorname{Spur}(AB)
&=\sum_{\mu=1}^\infty(AB\varphi_\mu,\varphi_\mu)
=\sum_{\mu=1}^\infty(B\varphi_\mu,A^*\varphi_\mu)\\
&=\sum_{\mu,\nu=1}^\infty(B\varphi_\mu,\psi_\nu)(\psi_\nu,A^*\varphi_\mu)\\
&=\sum_{\mu,\nu=1}^\infty(B\varphi_\mu,\psi_\nu)(A\psi_\nu,\varphi_\mu).
\end{align*}
其中$\varphi_1,\allowbreak \varphi_2,\allowbreak \ldots$和$\psi_1,\allowbreak \psi_2,\allowbreak \ldots$可以是任意两个完全规范正交系。这个式子对于$A,\allowbreak B$的对称性——同时交换$\varphi,\allowbreak \psi$——显而易见。因此，对厄米算符$A,\allowbreak B$有
\begin{align*}
\operatorname{Spur}(AB)
&=\overline{\operatorname{Spur}((AB)^*)}
=\overline{\operatorname{Spur}(B^*A^*)}\\
&=\overline{\operatorname{Spur}(BA)}
=\overline{\operatorname{Spur}(AB)}.
\end{align*}
所以$\operatorname{Spur}(AB)$为实数；$\operatorname{Spur}(A)$当然更是如此。

若$\mathfrak M$是闭线性流形，$E$是它的投影算符，则$\operatorname{Spur}(E)$可如下求出：设$\psi_1,\allowbreak \ldots,\allowbreak \psi_k$为张成$\mathfrak M$的规范正交系，$\chi_1,\allowbreak \ldots,\allowbreak \chi_l$为张成$\mathfrak R_\infty-\mathfrak M$的规范正交系。当然，$k,\allowbreak l$中至少一个是无限的，也可以两者都无限。

\hypertarget{origpage-95}{}\phantomsection\label{origpage-95}于是$\psi_1,\allowbreak \ldots,\allowbreak \psi_k,\allowbreak \chi_1,\allowbreak \ldots,\allowbreak \chi_l$合起来张成$\mathfrak R_\infty$，即依 \hyperref[sec:II.2]{II.2} 的\CJKglue \hyperref[thm:II.7]{定理7} 之$\alpha$，构成完全规范正交系。因此
\begin{align*}
\operatorname{Spur}(E)
&=\sum_{\mu=1}^k(E\psi_\mu,\psi_\mu)+\sum_{\mu=1}^l(E\chi_\mu,\chi_\mu)\\
&=\sum_{\mu=1}^k(\psi_\mu,\psi_\mu)+\sum_{\mu=1}^l(0,\chi_\mu)
=\sum_{\mu=1}^k1=k.
\end{align*}
即$\operatorname{Spur}(E)$是$\mathfrak M$的维数。

若$A$定号，所有$(A\varphi_\mu,\allowbreak \varphi_\mu)\geq0$，所以$\operatorname{Spur}(A)\geq0$。若$\operatorname{Spur}(A)=0$，所有$(A\varphi_\mu,\allowbreak \varphi_\mu)$都必须为零，从而依 \hyperref[sec:II.5]{II.5} 的\CJKglue \hyperref[thm:II.19]{定理19}，$A\varphi_\mu=0$。若$\|\varphi\|=1$，则能找到一个以$\varphi_1=\varphi$开头的完全规范正交系$\varphi_1,\allowbreak \varphi_2,\allowbreak \ldots$：取一个处处稠密的序列$f_1,\allowbreak f_2,\allowbreak \ldots$，再按 \hyperref[sec:II.2]{II.2} 的\CJKglue \hyperref[thm:II.7]{定理7} 证明中的方法，把$\varphi,\allowbreak f_1,\allowbreak f_2,\allowbreak \ldots$“正交化”，即可得到以$\varphi$开头的完全规范正交系。\tnote[。][所][原]{原书这里引定理7；正交化方法实际上在\CJKglue \hyperref[thm:II.8]{定理8} 中证明。正文保留原书的引用。}所以$A\varphi=0$。对于任意$f$，若$f=0$，$Af=0$显然成立；若$f\ne0$，就在上述结论中取$\varphi=f/\|f\|$。于是$A=0$。换言之，若$A$定号且非零，就有$\operatorname{Spur}(A)>0$。

上面关于迹的讨论虽简短明了，在数学上却并非无可非议。因为我们讨论级数
\[
\sum_{\mu,\nu=1}^\infty(A\varphi_\mu,\psi_\nu)(\psi_\nu,\varphi_\mu)
\quad\text{且}\quad
\sum_{\mu=1}^\infty(A\varphi_\mu,\varphi_\mu)
\]
时完全不顾其收敛性，又重排求和次序，把它们互相变形；也就是说，凡是严格说来不该做的事，我们几乎都做了。当然，理论物理的其他地方也不乏这样的疏忽；这里的疏忽，在量子力学的应用中也不会造成不良后果。但无论如何，必须承认这是一种疏忽。正因如此，更须强调：量子力学的基本统计断言，只对$A,\allowbreak B$均定号的算符$AB$使用迹；而迹的这个概念也完全可以建立在严格的基础上。因此，本节余下部分将汇集有关迹的那些可以用完全严格的数学论证证明的事实。

先考察$A^*A$的迹，其中$A$任意。依 \hyperref[sec:II.4]{II.4}，$A^*A$厄米，且因$(A^*Af,\allowbreak f)=(Af,\allowbreak Af)\geq0$而定号。有
\[
\operatorname{Spur}(A^*A)
=\sum_{\mu=1}^\infty(A^*A\varphi_\mu,\varphi_\mu)
=\sum_{\mu=1}^\infty(A\varphi_\mu,A\varphi_\mu)
=\sum_{\mu=1}^\infty\|A\varphi_\mu\|^2.
\]
这一数项级数的每一项都非负，所以或者收敛，或者发散至$+\infty$，无论如何都有意义。现在要不借助前面的讨论，证明它的和与$\varphi_1,\allowbreak \varphi_2,\allowbreak \ldots$的选择无关。


% 中文编写中间材料：chapter-II-s11.tex
% 第II章第11节续，承接原印95页末。
\hypertarget{origpage-96}{}\phantomsection\label{origpage-96}下面只会出现各项均不小于零的级数，因而一切都有意义，求和次序也可以任意改排。

设$\varphi_1,\allowbreak \varphi_2,\allowbreak \ldots$和$\psi_1,\allowbreak \psi_2,\allowbreak \ldots$为两个完全规范正交系，定义
\[
\Sigma(A;\varphi_\mu,\psi_\nu)
=\sum_{\mu,\nu=1}^{\infty}|(A\varphi_\mu,\psi_\nu)|^2.
\]
根据 \hyperref[sec:II.2]{II.2} 的\CJKglue \hyperref[thm:II.7]{定理7} 之$\gamma$，它等于$\sum_{\mu=1}^{\infty}\|A\varphi_\mu\|^2$，即$\Sigma(A;\varphi_\mu,\allowbreak \psi_\nu)$对$\psi_\nu$的依赖只是表面的。此外，只要$A\varphi_\mu,\allowbreak A^*\psi_\nu$都有意义，就有
\begin{align*}
\Sigma(A;\varphi_\mu,\psi_\nu)
&=\sum_{\mu,\nu=1}^{\infty}|(A\varphi_\mu,\psi_\nu)|^2
=\sum_{\mu,\nu=1}^{\infty}|(\varphi_\mu,A^*\psi_\nu)|^2\\
&=\sum_{\mu,\nu=1}^{\infty}|(A^*\psi_\nu,\varphi_\mu)|^2
=\Sigma(A^*;\psi_\nu,\varphi_\mu).
\end{align*}
因此，它对$\varphi_\mu$的依赖也只是表面的，因为按刚才的结论，等式右端与作为第二组正交系的$\varphi_\mu$无关。所以$\Sigma(A;\varphi_\mu,\allowbreak \psi_\nu)$只依赖于$A$，记作$\Sigma(A)$。根据上述证明，有
\[
\Sigma(A)=\sum_{\mu=1}^{\infty}\|A\varphi_\mu\|^2
=\sum_{\mu,\nu=1}^{\infty}|(A\varphi_\mu,\psi_\nu)|^2
\]
以及$\Sigma(A)=\Sigma(A^*)$。这样，就把$\operatorname{Spur}(A^*A)$正确地重新定义为$\Sigma(A)$。

对于$\Sigma(A)$的若干性质，我们也将另行证明；这些性质也可由前面已经导出的$\operatorname{Spur}(A)$的一般性质推出。

由定义，一般有$\Sigma(A)\geq0$；若$\Sigma(A)=0$，则所有$A\varphi_\mu$都必须为零，像前面一样可得$A=0$。因此，$A\ne0$时$\Sigma(A)>0$。

显然，$\Sigma(aA)=|a|^2\Sigma(A)$。若$A^*B=0$，则
\begin{align*}
&\|(A+B)\varphi_\mu\|^2-\|A\varphi_\mu\|^2-\|B\varphi_\mu\|^2\\
&\qquad=(A\varphi_\mu,B\varphi_\mu)+(B\varphi_\mu,A\varphi_\mu)\\
&\qquad=2\operatorname{Re}(A\varphi_\mu,B\varphi_\mu)
=2\operatorname{Re}(\varphi_\mu,A^*B\varphi_\mu)=0,
\end{align*}
所以，对$\mu=1,\allowbreak 2,\allowbreak \ldots$求和得
\[
\Sigma(A+B)=\Sigma(A)+\Sigma(B).
\]
交换$A,\allowbreak B$并不改变这个关系，所以$B^*A=0$时它也成立。还可以把$A,\allowbreak B$换成$A^*,\allowbreak B^*$；因此$AB^*=0$或$BA^*=0$同样是充分条件。于是，当$A$或$B$厄米时，可写成$AB=0$或$BA=0$。

若$E$为闭线性流形$\mathfrak M$的投影算符，则对求$\operatorname{Spur}(E)$时考察过的$\psi_1,\allowbreak \ldots,\allowbreak \psi_k,\allowbreak \chi_1,\allowbreak \ldots,\allowbreak \chi_l$有\hypertarget{origpage-97}{}\phantomsection\label{origpage-97}\begin{align*}
\Sigma(E)&=\sum_{\mu=1}^{k}\|E\psi_\mu\|^2
+\sum_{\mu=1}^{l}\|E\chi_\mu\|^2\\
&=\sum_{\mu=1}^{k}\|\psi_\mu\|^2
+\sum_{\mu=1}^{l}\|0\|^2
=\sum_{\mu=1}^{k}1=k.
\end{align*}
所以$\Sigma(E)$也是$\mathfrak M$的维数；由于$E^*E=EE=E$，这一结果本来就在意料之中。

对于两个定号的厄米算符$A,\allowbreak B$，现在可把$\operatorname{Spur}(AB)$归结到$\Sigma$。因为存在两个同样性质的算符$A',\allowbreak B'$，满足$A'^2=A,\allowbreak B'^2=B$，\origNote[，][称]{114}称为$\sqrt A,\allowbreak \sqrt B$。纯形式地有
\begin{align*}
\operatorname{Spur}(AB)
&=\operatorname{Spur}(\sqrt A\sqrt A\sqrt B\cdot\sqrt B)\\
&=\operatorname{Spur}(\sqrt B\cdot\sqrt A\sqrt A\sqrt B)\\
&=\operatorname{Spur}\bigl((\sqrt A\sqrt B)^*(\sqrt A\sqrt B)\bigr)
=\Sigma(\sqrt A\sqrt B).
\end{align*}
这个$\Sigma(\sqrt A\sqrt B)$，只依据它本身的定义，而不借助它与迹的关系，就已具有我们希望$\operatorname{Spur}(AB)$具有的全部性质，即
\begin{align*}
\Sigma(\sqrt A\sqrt B)&=\Sigma(\sqrt B\sqrt A),\\
\Sigma(\sqrt A\sqrt{B+C})&=\Sigma(\sqrt A\sqrt B)+\Sigma(\sqrt A\sqrt C),\\
\Sigma(\sqrt{A+B}\sqrt C)&=\Sigma(\sqrt A\sqrt C)+\Sigma(\sqrt B\sqrt C).
\end{align*}
第一式由$\Sigma(XY)$关于$X,\allowbreak Y$的对称性得出：
\[
\Sigma(XY)=\sum_{\mu,\nu=1}^{\infty}|(XY\varphi_\mu,\psi_\nu)|^2
=\sum_{\mu,\nu=1}^{\infty}|(Y\varphi_\mu,X\psi_\nu)|^2;
\]
利用第一式，第二式可由第三式得出。因此只须证明第三式，即$\Sigma(\sqrt A\sqrt B)$对$A$具有可加性。把它写成
\begin{align*}
\Sigma(\sqrt A\sqrt B)
&=\sum_{\mu=1}^{\infty}\|\sqrt A\sqrt B\varphi_\mu\|^2\\
&=\sum_{\mu=1}^{\infty}(\sqrt A\sqrt B\varphi_\mu,\sqrt A\sqrt B\varphi_\mu)\\
&=\sum_{\mu=1}^{\infty}(\sqrt A\cdot\sqrt A\sqrt B\varphi_\mu,\sqrt B\varphi_\mu)\\
&=\sum_{\mu=1}^{\infty}(A\sqrt B\varphi_\mu,\sqrt B\varphi_\mu)
\end{align*}
便可看出。至此，我们已在上述所需的范围内，为迹的概念建立了严格的基础。

最后一个公式还给出如下结论：若$A,\allowbreak B$定号，则
\[
\operatorname{Spur}(AB)=0\quad\text{可推出}\quad AB=0.
\]
因为前一条件是说$\Sigma(\sqrt A\sqrt B)=0$，所以$\sqrt A\sqrt B=0$（参见第 \srcpage{96} 页的说明，或下面关于$\Sigma$的进一步讨论），从而
\[
AB=\sqrt A\cdot\sqrt A\sqrt B\cdot\sqrt B=0.
\]

对于定号厄米算符$A$，即使沿用最初的形式，以迹进行计算也是正确的。设$\varphi_1,\allowbreak \varphi_2,\allowbreak \ldots$为完全规范正交系，则\hypertarget{origpage-98}{}\phantomsection\label{origpage-98}$\sum_{\mu=1}^{\infty}(A\varphi_\mu,\allowbreak \varphi_\mu)$——原本用来定义迹的那个和——各项均非负，所以或者收敛，或者发散到$+\infty$。于是有两种可能：或者，无论怎样选择$\varphi_1,\allowbreak \varphi_2,\allowbreak \ldots$，此和总是无限的；那么迹确实已不依赖于选择而定义，且等于$+\infty$。或者，至少有一种选择，例如$\bar\varphi_1,\allowbreak \bar\varphi_2,\allowbreak \ldots$，使此和有限。此时，由于
\begin{align*}
\left(\sum_{\mu=1}^{\infty}(A\bar\varphi_\mu,\bar\varphi_\mu)\right)^2
&=\sum_{\mu,\nu=1}^{\infty}(A\bar\varphi_\mu,\bar\varphi_\mu)
(A\bar\varphi_\nu,\bar\varphi_\nu)\\
&\geq\sum_{\mu,\nu=1}^{\infty}|(A\bar\varphi_\mu,\bar\varphi_\nu)|^2
=\Sigma(A),
\end{align*}
$\Sigma(A)$也有限，设等于$C^2$。若$\psi_1,\allowbreak \psi_2,\allowbreak \ldots$是任意完全规范正交系，则
\[
\Sigma(A)=\sum_{\mu=1}^{\infty}\|A\psi_\mu\|^2=C^2,
\qquad\|A\psi_1\|^2\leq C^2,
\qquad\|A\psi_1\|\leq C.
\]
每个$\|\psi\|=1$的$\psi$都可选作这种系的$\psi_1$，所以$\|\psi\|=1$能推出$\|A\psi\|\leq C$。因此普遍有$\|Af\|\leq C\|f\|$：当$f=0$时显然如此；当$f\ne0$时，令$\psi=\frac1{\|f\|}\cdot f$即可。所以$A$满足 \hyperref[sec:II.9]{II.9} 的条件 \hyperref[eq:II.9.St]{St}，是连续算符。不过，还能得出远强于此的结论。

因为$\Sigma(A)$有限，$A$属于所谓全连续算符这一类。Hilbert已证明，这类算符的本征值问题按最初的形式就可解，即存在完全规范正交系$\psi_1,\allowbreak \psi_2,\allowbreak \ldots$，满足$A\psi_\mu=\lambda_\mu\psi_\mu$，并且$\mu\to\infty$时$\lambda_\mu\to0$。\origNote[。][由]{115}由定号性，$\lambda_\mu=(A\psi_\mu,\allowbreak \psi_\mu)\geq0$；此外，
\[
\sum_{\mu=1}^{\infty}\lambda_\mu^2
=\sum_{\mu=1}^{\infty}\|A\psi_\mu\|^2=\Sigma(A)=C^2.
\]
若$\varphi_1,\allowbreak \varphi_2,\allowbreak \ldots$为任意完全规范正交系，则
\begin{align*}
\sum_{\mu=1}^{\infty}(A\varphi_\mu,\varphi_\mu)
&=\sum_{\mu=1}^{\infty}\left(\sum_{\nu=1}^{\infty}
(A\varphi_\mu,\psi_\nu)(\psi_\nu,\varphi_\mu)\right)\\
&=\sum_{\mu=1}^{\infty}\left(\sum_{\nu=1}^{\infty}
(\varphi_\mu,A\psi_\nu)(\psi_\nu,\varphi_\mu)\right)\\
&=\sum_{\mu=1}^{\infty}\left(\sum_{\nu=1}^{\infty}
\lambda_\nu(\varphi_\mu,\psi_\nu)(\psi_\nu,\varphi_\mu)\right)\\
&=\sum_{\mu=1}^{\infty}\left(\sum_{\nu=1}^{\infty}
\lambda_\mu|(\varphi_\mu,\psi_\nu)|^2\right).
\end{align*}
由于一切均非负，可以改排求和次序：\tnote[：][][原]{原书在上式最后一个内层求和中印作$\lambda_\mu$，但上一行及下面改排后的式子要求$\lambda_\nu$。这里保留原书下标。}
\begin{align*}
\sum_{\mu=1}^{\infty}(A\varphi_\mu,\varphi_\mu)
&=\sum_{\mu,\nu=1}^{\infty}\lambda_\nu|(\varphi_\mu,\psi_\nu)|^2\\
&=\sum_{\nu=1}^{\infty}\lambda_\nu
\left(\sum_{\mu=1}^{\infty}|(\varphi_\mu,\psi_\nu)|^2\right)\\
&=\sum_{\nu=1}^{\infty}\lambda_\nu\|\psi_\nu\|^2
=\sum_{\nu=1}^{\infty}\lambda_\nu.
\end{align*}

\hypertarget{origpage-99}{}\phantomsection\label{origpage-99}因此，此时$\sum_{\mu=1}^{\infty}(A\varphi_\mu,\allowbreak \varphi_\mu)$也与$\varphi_1,\allowbreak \varphi_2,\allowbreak \ldots$无关，且等于本征值之和。它对于$\bar\varphi_1,\allowbreak \bar\varphi_2,\allowbreak \ldots$这组系是有限的，所以对任何选择都有限。因此$\operatorname{Spur}(A)$仍唯一确定，不过这次是有限值。

所以，在两种情形下，以迹进行计算都有依据。

再给出几个关于$\operatorname{Spur}(A)$和$\Sigma(A)$的估计。对于一切$\Sigma(A)$有限的$A$，已有$\|Af\|\leq\sqrt{\Sigma(A)}\cdot\|f\|$；对于一切$\operatorname{Spur}(A)$有限的定号厄米算符$A$，有
\[
\|Af\|\leq\operatorname{Spur}(A)\cdot\|f\|.
\]
再设$A$定号，$\operatorname{Spur}(A)=1$，且对某个适当的$\|\varphi\|=1$的$\varphi$，有$\|A\varphi\|^2\geq1-\varepsilon$或$(A\varphi,\allowbreak \varphi)\geq1-\varepsilon$。由于$(A\varphi,\allowbreak \varphi)\leq\|A\varphi\|\cdot\|\varphi\|=\|A\varphi\|$，第二种条件可推出第一种，只是须用$(1-\varepsilon)^2\geq1-2\varepsilon$代替$1-\varepsilon$，即以$2\varepsilon$代替$\varepsilon$。因此只考察第一种就够了。

设$\psi$与$\varphi$正交，且$\|\psi\|=1$。这时能找到完全规范正交系$\chi_1,\allowbreak \chi_2,\allowbreak \ldots$，使$\chi_1=\varphi,\allowbreak \chi_2=\psi$，所以
\begin{align*}
\sum_{\mu=1}^{\infty}\|A\chi_\mu\|^2
&=\Sigma(A)\leq(\operatorname{Spur}(A))^2=1,\\
\sum_{\mu=1}^{\infty}\|A\chi_\mu\|^2
&\geq\|A\varphi\|^2+\|A\psi\|^2
\geq1-2\varepsilon+\|A\psi\|^2,\\
\|A\psi\|^2&\leq2\varepsilon,\qquad\|A\psi\|\leq\sqrt{2\varepsilon}.
\end{align*}
对于任意与$\varphi$正交的$f$，由此得$\|Af\|\leq\sqrt{2\varepsilon}\cdot\|f\|$：$f=0$时显然成立，否则令$\psi=\frac1{\|f\|}\cdot f$即可。再注意$(Af,\allowbreak g)=(f,\allowbreak Ag)$，便可说，只要$f,\allowbreak g$中有一个与$\varphi$正交，就有$|(Af,\allowbreak g)|\leq\sqrt{2\varepsilon}\cdot\|f\|\cdot\|g\|$。

现在任取$f,\allowbreak g$，有
\[
f=\alpha\varphi+f',\qquad g=\beta\varphi+g',
\]
其中$f',\allowbreak g'$与$\varphi$正交，$\alpha=(f,\allowbreak \varphi),\allowbreak \beta=(g,\allowbreak \varphi)$。所以
\[
(Af,g)=\alpha\bar\beta(A\varphi,\varphi)
+\alpha(A\varphi,g')+\bar\beta(Af',\varphi)+(Af',g'),
\]
令$(A\varphi,\allowbreak \varphi)=c$，得到
\[
|(Af,g)-c\alpha\bar\beta|
\leq|\alpha|\cdot|(A\varphi,g')|
+|\beta|\cdot|(Af',\varphi)|+|(Af',g')|,
\]
再用上述估计得
\begin{align*}
|(Af,g)-c\alpha\bar\beta|
&\leq\sqrt{2\varepsilon}\cdot
\bigl(|\alpha|\cdot\|g'\|+|\beta|\cdot\|f'\|+\|f'\|\cdot\|g'\|\bigr)\\
&\leq\sqrt{2\varepsilon}\cdot(|\alpha|+\|f'\|)(|\beta|+\|g'\|)\\
&\leq2\sqrt{2\varepsilon}\cdot
\sqrt{|\alpha|^2+\|f'\|^2}\sqrt{|\beta|^2+\|g'\|^2}\\
&=2\sqrt{2\varepsilon}\cdot\|f\|\cdot\|g\|.
\end{align*}
另一方面，
\begin{align*}
(Af,g)-c\alpha\bar\beta
&=(Af,g)-c(f,\varphi)(\varphi,g)\\
&=((A-cP_{[\varphi]})f,g).
\end{align*}
所以普遍有$|((A-cP_{[\varphi]})f,\allowbreak g)|\leq2\sqrt{2\varepsilon}\cdot\|f\|\cdot\|g\|$；由 \hyperref[sec:II.9]{II.9} 已经知道，因而也有\hypertarget{origpage-100}{}\phantomsection\label{origpage-100}\[
\|(A-cP_{[\varphi]})f\|\leq2\sqrt{2\varepsilon}\cdot\|f\|.
\]
令$f=\varphi$，便是（其中$c=(A\varphi,\allowbreak \varphi)$为非负实数）：
\[
\|A\varphi-c\varphi\|\leq2\sqrt{2\varepsilon},
\]
\begin{align*}
c=\|c\varphi\|&\leq\|A\varphi-c\varphi\|+\|A\varphi\|
\leq2\sqrt{2\varepsilon}+1,\\
c=\|c\varphi\|&\geq-\|A\varphi-c\varphi\|+\|A\varphi\|
\geq-2\sqrt{2\varepsilon}+(1-\varepsilon),
\end{align*}
\[
1-(\varepsilon+2\sqrt{2\varepsilon})\leq c\leq1+2\sqrt{2\varepsilon}.
\]
因此
\begin{align*}
\|(A-P_{[\varphi]})f\|
&\leq\|(A-cP_{[\varphi]})f\|+\|(c-1)P_{[\varphi]}f\|\\
&\leq2\sqrt{2\varepsilon}\cdot\|f\|
+(\varepsilon+2\sqrt{2\varepsilon})\|P_{[\varphi]}f\|\\
&\leq(\varepsilon+4\sqrt{2\varepsilon})\cdot\|f\|.
\end{align*}
所以当$\varepsilon\to0$时，$A$一致收敛于$P_{[\varphi]}$。

最后，考察$\mathfrak R_\infty$的实现$F_Z,\allowbreak F_\Omega$中的$\operatorname{Spur}(A),\allowbreak \Sigma(A)$（见 \hyperref[sec:I.4]{I.4} 及 \hyperref[sec:II.3]{II.3}），因为物理应用将会在这些实现中进行。

在$F_Z$中，即在由使$\sum_{\mu=1}^{\infty}|x_\mu|^2$有限的一切序列$\{x_1,\allowbreak x_2,\allowbreak \ldots\}$构成的空间中，$A$可由矩阵$\{a_{\mu\nu}\}$描述：
\[
A\{x_1,x_2,\ldots\}=\{y_1,y_2,\ldots\},
\qquad y_\mu=\sum_{\nu=1}^{\infty}a_{\mu\nu}x_\nu.
\]
序列$\{1,\allowbreak 0,\allowbreak 0,\allowbreak \ldots\},\allowbreak \{0,\allowbreak 1,\allowbreak 0,\allowbreak \ldots\},\allowbreak \ldots$构成完全规范正交系$\varphi_1,\allowbreak \varphi_2,\allowbreak \ldots$，并有
\[
A\varphi_\mu=\{a_{1\mu},a_{2\mu},\ldots\}
=\sum_{\varrho=1}^{\infty}a_{\varrho\mu}\varphi_\varrho,
\]
所以$(A\varphi_\mu,\allowbreak \varphi_\mu)=a_{\mu\mu}$，$\|A\varphi_\mu\|^2=\sum_{\varrho=1}^{\infty}|a_{\varrho\mu}|^2$。立即得到
\[
\operatorname{Spur}(A)=\sum_{\mu=1}^{\infty}a_{\mu\mu},
\qquad\Sigma(A)=\sum_{\mu,\nu=1}^{\infty}|a_{\mu\nu}|^2.
\]
在$F_\Omega$，即定义于$\Omega$且使$\int_\Omega|f(P)|^2\,dv$有限的一切$f(P)$构成的集合中，只考察积分算符
\[
Af(P)=\int_\Omega a(P,P')f(P')\,dv',
\]
其中$a(P,\allowbreak P')$是两个变量均定义在$\Omega$中的函数，称为$A$的“积分核”（见 \hyperref[sec:I.4]{I.4}）。设$\varphi_1(P),\allowbreak \varphi_2(P),\allowbreak \ldots$为任意完全规范正交系，则
\begin{align*}
\operatorname{Spur}(A)
&=\sum_{\mu=1}^{\infty}(A\varphi_\mu(P),\varphi_\mu(P))\\
&=\sum_{\mu=1}^{\infty}\int_\Omega
\left[\int_\Omega a(P,P')\varphi_\mu(P')\,dv'\right]
\overline{\varphi_\mu(P)}\,dv,
\end{align*}\hypertarget{origpage-101}{}\phantomsection\label{origpage-101}又因为一般有（把 \hyperref[sec:II.2]{II.2} 的\CJKglue \hyperref[thm:II.7]{定理7} 之$\beta$应用于$\overline{g(P)}$）
\[
\sum_{\mu=1}^{\infty}\left(\int_\Omega\overline{g(P')}\,
\overline{\varphi_\mu(P')}\,dv'\right)\varphi_\mu(P)
=\overline{g(P)},
\]
\[
\sum_{\mu=1}^{\infty}\left(\int_\Omega g(P')\varphi_\mu(P')\,dv'\right)
\overline{\varphi_\mu(P)}=g(P),
\]
所以
\[
\operatorname{Spur}(A)=\int_\Omega a(P,P)\,dv.
\]
此外，
\[
\Sigma(A)=\sum_{\mu=1}^{\infty}\int_\Omega
\left|\int_\Omega a(P,P')\varphi_\mu(P')\,dv'\right|^2\,dv,
\]
于是，由（\hyperref[sec:II.2]{II.2} 的\CJKglue \hyperref[thm:II.7]{定理7} 之$\gamma$）
\begin{align*}
\sum_{\mu=1}^{\infty}\left|\int_\Omega g(P')\varphi_\mu(P')\,dv'\right|^2
&=\sum_{\mu=1}^{\infty}\left|\int_\Omega
\overline{g(P')}\,\overline{\varphi_\mu(P')}\,dv'\right|^2\\
&=\int_\Omega|\overline{g(P')}|^2\,dv'
=\int_\Omega|g(P')|^2\,dv',
\end{align*}
进一步得到
\[
\Sigma(A)=\int_\Omega\int_\Omega|a(P,P')|^2\,dv\,dv'.
\]

可见，$\operatorname{Spur}(A),\allowbreak \Sigma(A)$实现了 \hyperref[sec:I.4]{I.4} 中借助数学上颇可疑的技巧所追求的目标：从$F_Z$过渡到$F_\Omega$时，$\sum_{\mu=1}^{\infty}\cdots$由$\int_\Omega\cdots\,dv$代替。

至此，我们关于厄米算符的数学讨论告一段落。对数学有兴趣的读者，可以从有关文献中找到进一步的内容。\origNote[。][]{116}


% 中文编写中间材料：chapter-III.tex
\chapter{量子力学的统计学}\label{chap:III}
\section{量子力学的统计断言}\label{sec:III.1}
现在回到量子力学理论的分析。前面为了展开\CJKglue \hyperref[chap:II]{第II章}\CJKglue 的数学讨论，我们暂时中断了这项分析。当时只讨论了量子力学怎样确定一个特定物理量——能量——的一切可能值：这些值就是能量算符$\mathsf H$的本征值，也就是它的谱中的数。至于其他物理量的取值，以及几个物理量的值之间有何种因果联系或统计联系，当时还没有涉及。现在就来考察理论对这些问题所作的断言。我们将以波动力学的描述方式为依据；两种理论的等价性已经得到保证。

\hypertarget{origpage-102}{}\phantomsection\label{origpage-102}采用这种描述方式，显然，关于一个系统状态的任何断言，都必须从它的波函数$\varphi(q_1,\allowbreak \ldots,\allowbreak q_k)$中推得。（设系统有$k$个自由度，$q_1,\allowbreak \ldots,\allowbreak q_k$为其位形空间的坐标。）这里不只考虑系统的定态，即量子轨道；在这种状态下，$\varphi$是$\mathsf H$的本征函数，满足$\mathsf H\varphi=\lambda\varphi$，参见 \hyperref[sec:I.3]{I.3}。我们允许同一系统处于任意状态，也就是允许任意波函数$\varphi$；它随时间的变化服从Schrödinger的含时微分方程$\mathsf H\varphi=-\frac h{2\pi i}\frac\partial{\partial t}\varphi$，参见 \hyperref[sec:I.2]{I.2}。那么，对于一个处于状态$\varphi$的系统，我们能作出哪些断言呢？

首先要注意，$\varphi$已经按 \hyperref[sec:I.3]{I.3} 中的条件归一化：
\[
\int_{-\infty}^{\infty}\cdots\int_{-\infty}^{\infty}
|\varphi(q_1\cdots q_k)|^2\,dq_1\cdots dq_k=1.
\]
用我们现在的术语说，考虑所有满足
\[
\int_{-\infty}^{\infty}\cdots\int_{-\infty}^{\infty}
|f(q_1\cdots q_k)|^2\,dq_1\cdots dq_k
\]
为有限值的函数$f(q_1\cdots q_k)$所成的希尔伯特空间$\mathfrak R_\infty$（这是一个$F_\Omega$），$\varphi$在其中的范数为$\|\varphi\|=1$。也就是说，$\varphi$应落在希尔伯特空间的单位球面上\origNote[上][。]{117}。我们已经知道，给$\varphi$乘上一个常数因子，即与$q_1,\allowbreak \ldots,\allowbreak q_k$无关的因子，并不改变其物理意义。（这就是用$a\varphi$代替$\varphi$，其中$a$为复数。由于有归一化条件$\|\varphi\|=1$，必须满足$|a|=1$。）还要指出，除了系统位形空间的坐标$q_1,\allowbreak \ldots,\allowbreak q_k$，$\varphi$也依赖于时间$t$。不过，希尔伯特空间是只针对$q_1,\allowbreak \ldots,\allowbreak q_k$构造的，因为归一化条件也只涉及这些变量。因此，必须留意$t$这一依赖关系；更恰当地说，在这里应把$t$看作一个参数。由此，$\varphi$作为$\mathfrak R_\infty$中的点，依赖于$t$，却不依赖于$q_1,\allowbreak \ldots,\allowbreak q_k$：作为$\mathfrak R_\infty$中的一个点，它表示的是这个函数随$q_1,\allowbreak \ldots,\allowbreak q_k$变化的完整情形。为此，当我们把$\varphi$看作$\mathfrak R_\infty$中的点时，有时会将参数$t$写出来，记作$\varphi_t$。

现在考察状态$\varphi=\varphi(q_1\cdots q_k)$。关于这个状态的统计断言是：系统在位形空间的点$q_1,\allowbreak \ldots,\allowbreak q_k$处的概率密度为$|\varphi(q_1\cdots q_k)|^2$。也就是说，系统位于位形空间某一区域$V$内的概率为
\[
\int\cdots\int_V|\varphi(q_1\cdots q_k)|^2\,dv.
\]
（这些是最早、也最简单的例子，人们由此认识到量子力学的统计性质\origNote[质][。]{118}。顺便说，这一断言与Schrödinger关于电荷分布的假设之间的关系也是清楚的，参见 \hyperref[sec:I.2]{I.2}。）\hypertarget{origpage-103}{}\phantomsection\label{origpage-103}此外，若系统的能量对应算符$\mathsf H$，其本征值为$\lambda_1,\allowbreak \lambda_2,\allowbreak \ldots$，本征函数为$\varphi_1,\allowbreak \varphi_2,\allowbreak \ldots$，那么，系统处于状态$\varphi$时，能量取值$\lambda_n$的概率为
\[
\left|\int\cdots\int\varphi(q_1\cdots q_k)
\overline{\varphi_n(q_1\cdots q_k)}\,dq_1\cdots dq_k\right|^2
\]
（参见原注 \hyperlink{note-118}{\getrefnumber{orig-note-118}} 中提到的论文）。下面把这些断言改写为统一的形式。

设$V$为$k$维长方体
\[
q'_1<q_1\le q''_1,\quad\ldots,\quad q'_k<q_k\le q''_k,
\]
将区间$\{q'_1,\allowbreak q''_1\},\allowbreak \ldots,\allowbreak \{q'_k,\allowbreak q''_k\}$分别记为$I_1,\allowbreak \ldots,\allowbreak I_k$。$q_1,\allowbreak \ldots,\allowbreak q_k$分别对应算符$q_1\cdots,\allowbreak \ldots,\allowbreak q_k\cdots$；相应的单位分解定义如下（参见 \hyperref[sec:II.8]{II.8}）：将$q_j$（$j=1,\allowbreak \ldots,\allowbreak k$）对应的单位分解记作$E_j(\lambda)$，则
\[
E_j(\lambda)f(q_1\cdots q_k)=
\begin{cases}f(q_1\cdots q_k),&q_j\le\lambda,\\0,&q_j>\lambda.\end{cases}
\]
我们还统一采用以下记号：若$F(\lambda)$为一个单位分解，$I$为区间$\{\lambda',\allowbreak \lambda''\}$，便记$F(I)=F(\lambda'')-F(\lambda')$。因为$\lambda'\le\lambda''$，有$F(\lambda')\le F(\lambda'')$，所以它是一个投影算符。系统位于上述$V$内，也就是$q_1$在$I_1$内、……、$q_k$在$I_k$内的概率为
\begin{align*}
&\int_{q'_1}^{q''_1}\cdots\int_{q'_k}^{q''_k}|\varphi(q_1\cdots q_k)|^2\,dq_1\cdots dq_k\\
&\qquad=\int\cdots\int|E_1(I_1)\cdots E_k(I_k)\varphi(q_1\cdots q_k)|^2\,dq_1\cdots dq_k,
\end{align*}
因为当$q_1$在$I_1$内、……、$q_k$在$I_k$内时，$E_1(I_1)\cdots E_k(I_k)\varphi(q_1\cdots q_k)$等于$\varphi(q_1\cdots q_k)$，其余情形等于$0$。因此，这个概率也就是
\[
\|E_1(I_1)\cdots E_k(I_k)\varphi\|^2.
\]
第二种情形，考察能量位于区间$I=\{\lambda',\allowbreak \lambda''\}$内的概率。$\mathsf H$对应的单位分解$E(\lambda)$按如下方式定义（参见 \hyperref[sec:II.8]{II.8}）：$E(\lambda)=\sum_{\lambda_n\le\lambda}P_{[\varphi_n]}$，于是
\[
E(I)=E(\lambda'')-E(\lambda')=\sum_{\lambda'<\lambda_n\le\lambda''}P_{[\varphi_n]}.
\]
由于能量的取值只可能是$\lambda_1,\allowbreak \lambda_2,\allowbreak \ldots$，所求概率就是所有满足$\lambda'<\lambda_n\le\lambda''$的$\lambda_n$各自概率的总和，即
\begin{align*}
&\sum_{\lambda'<\lambda_n\le\lambda''}
\left|\int\cdots\int\varphi(q_1\cdots q_k)\overline{\varphi_n(q_1\cdots q_k)}\,dq_1\cdots dq_k\right|^2\\
&\qquad=\sum_{\lambda'<\lambda_n\le\lambda''}|(\varphi,\varphi_n)|^2
=\sum_{\lambda'<\lambda_n\le\lambda''}(P_{[\varphi_n]}\varphi,\varphi)\\
&\qquad=\left(\left\{\sum_{\lambda'<\lambda_n\le\lambda''}P_{[\varphi_n]}\right\}\varphi,\varphi\right)
=(E(I)\varphi,\varphi)=\|E(I)\varphi\|^2.
\end{align*}

\hypertarget{origpage-104}{}\phantomsection\label{origpage-104}可见，这两种情形所得的结果可以统一表述如下：

\leavevmode\phantomsection\label{post:III.1.W}\textbf{（W.）}系统处于状态$\varphi$时，分别对应算符$R_1,\allowbreak \ldots,\allowbreak R_l$的物理量\origNote[量][，]{119}，其值分别落在区间$I_1,\allowbreak \ldots,\allowbreak I_l$内的概率为
\[
\|E_1(I_1)\cdots E_l(I_l)\varphi\|^2,
\]
其中$E_1(\lambda),\allowbreak \ldots,\allowbreak E_l(\lambda)$分别为$R_1,\allowbreak \ldots,\allowbreak R_l$对应的单位分解。

第一种情形对应于$l=k$、$R_1=q_1\cdots,\allowbreak \ldots,\allowbreak R_k=q_k\cdots$；第二种情形对应于$l=1$、$R_1=\mathsf H$。现在，我们把这个断言 \hyperref[post:III.1.W]{W.} 假定为普遍有效的。它确实包含了迄今为止量子力学作出的一切统计断言。

不过，必须对它的适用范围加一项限制。问题本身完全不涉及$R_1,\allowbreak \ldots,\allowbreak R_l$的先后次序，因此结果也必须与次序无关。也就是说，$E_1(I_1),\allowbreak \ldots,\allowbreak E_l(I_l)$，或者说所有$E_1(\lambda_1),\allowbreak \ldots,\allowbreak E_l(\lambda_l)$，必须彼此对易。根据 \hyperref[sec:II.10]{II.10}，这意味着$R_1,\allowbreak \ldots,\allowbreak R_l$彼此对易。$q_1\cdots,\allowbreak \ldots,\allowbreak q_k\cdots$满足这一条件；当$l=1$、$R_1=\mathsf H$时，则无需再考虑这一条件。

因此，我们对彼此对易的$R_1,\allowbreak \ldots,\allowbreak R_l$普遍提出公设 \hyperref[post:III.1.W]{W.}。此时$E_1(I_1),\allowbreak \ldots,\allowbreak E_l(I_l)$也彼此对易，所以$E_1(I_1)\cdots E_l(I_l)$为投影算符（\hyperref[thm:II.14]{II.4的定理14}），所述概率因而等于
\[
W=\|E_1(I_1)\cdots E_l(I_l)\varphi\|^2
=(E_1(I_1)\cdots E_l(I_l)\varphi,\varphi)
\]
（\hyperref[thm:II.12]{II.4的定理12}）。

继续讨论以前，必须核实 \hyperref[post:III.1.W]{W.} 具有任何合理的统计理论都应具备的几项性质。

1．各断言的先后次序无关紧要。

2．可以加入没有实质内容的断言，而不改变$W$。这指的是区间$I_j$等于$\{-\infty,\allowbreak +\infty\}$的断言，因为它只给乘积添上一个因子
\[
E_j(I)=E(+\infty)-E(-\infty)=1-0=1.
\]

3．概率的加法定理成立。也就是说，把某个区间$I_j$分成$I'_j,\allowbreak I''_j$两个区间，原来的概率就等于两个新概率之和。事实上，设$I_j,\allowbreak I'_j,\allowbreak I''_j$分别为$\{\lambda',\allowbreak \lambda''\}$、$\{\lambda',\allowbreak \lambda\}$、$\{\lambda,\allowbreak \lambda''\}$，则
\[
E(\lambda'')-E(\lambda')=(E(\lambda)-E(\lambda'))+(E(\lambda'')-E(\lambda)),
\]
即$E(I_j)=E(I'_j)+E(I''_j)$。将它代入上面$W$的第二种表达式，便得到概率的可加性，因为该式对$E_1(I_1)\cdots E_j(I_j)\cdots E_l(I_l)$是线性的。

\hypertarget{origpage-105}{}\phantomsection\label{origpage-105}4．对于不可能的断言，即某个$I_j$为空的情形，$W=0$，因为相应的$E_j(I_j)=0$。对于必然成立的断言，即所有$I_j$均为$\{-\infty,\allowbreak +\infty\}$的情形，$W=1$，因为所有$E_j(I_j)=1$，从而$W=\|\varphi\|^2=1$。一般而言，总有$0\le W\le1$，这是由 \hyperref[thm:II.13]{II.4的定理13} 得出的。

最后还要指出，\hyperref[post:III.1.W]{W.} 包含以下断言：一个物理量$R_j$只能取其本征值，也就是其谱中的数。因为当区间$I_j=\{\lambda',\allowbreak \lambda''\}$在谱之外时，$E_j(\lambda)$在此区间上为常数，故
\[
E_j(I_j)=E_j(\lambda'')-E_j(\lambda')=0,
\]
从而$W=0$。

现在令$l=1$，并将$R_1$写成$R$。$R$对应的物理量记作$\mathfrak R$（参见原注 \hyperlink{note-119}{\getrefnumber{orig-note-119}}）。设$F(\lambda)$为任意函数，来计算$F(\mathfrak R)$的期望值。

为此，将区间$\{-\infty,\allowbreak +\infty\}$分成一列子区间$\{\lambda_n,\allowbreak \lambda_{n+1}\}$，$n=0,\allowbreak \pm1,\allowbreak \pm2,\allowbreak \ldots$。$\mathfrak R$落在$\{\lambda_n,\allowbreak \lambda_{n+1}\}$内的概率为
\[
(\{E(\lambda_{n+1})-E(\lambda_n)\}\varphi,\varphi)
=(E(\lambda_{n+1})\varphi,\varphi)-(E(\lambda_n)\varphi,\varphi),
\]
因此$F(\mathfrak R)$的期望值为
\[
\sum_{n=-\infty}^{+\infty}F(\lambda'_n)
\{(E(\lambda_{n+1})\varphi,\varphi)-(E(\lambda_n)\varphi,\varphi)\},
\]
其中$\lambda'_n$为$\{\lambda_n,\allowbreak \lambda_{n+1}\}$内某个合适的中间值。当分点$\ldots,\allowbreak \lambda_{-2},\allowbreak \lambda_{-1},\allowbreak \lambda_0,\allowbreak \lambda_1,\allowbreak \lambda_2,\allowbreak \ldots$越来越密时，这些和收敛到Stieltjes积分
\[
\int_{-\infty}^{\infty}F(\lambda)\,d(E(\lambda)\varphi,\varphi).
\]
所以，所求的期望值就等于这个积分。根据 \hyperref[sec:II.8]{II.8} 中算符函数的一般定义，这个积分又等于$(F(R)\varphi,\allowbreak \varphi)$。于是有：

\leavevmode\phantomsection\label{post:III.1.E1}\textbf{（$E_1$.）}设$\mathfrak R$为任意物理量，$R$为其算符（参见原注 \hyperlink{note-119}{\getrefnumber{orig-note-119}}），$F(\lambda)$为任意函数。则在状态$\varphi$中，$f(\mathfrak R)$的期望值为
\[
\operatorname{Erw}(F(\mathfrak R);\varphi)=(F(R)\varphi,\varphi).
\]
特别地，令$F(\lambda)=\lambda$，便有：

\leavevmode\phantomsection\label{post:III.1.E2}\textbf{（$E_2$.）}设$\mathfrak R,\allowbreak R$的意义同上。则在状态$\varphi$中，$\mathfrak R$的期望值为
\[
\operatorname{Erw}(\mathfrak R;\varphi)=(R\varphi,\varphi).
\]
下面来研究 \hyperref[post:III.1.W]{W.}、\hyperref[post:III.1.E1]{$E_1$}、\hyperref[post:III.1.E2]{$E_2$} 之间的关系。

\hypertarget{origpage-106}{}\phantomsection\label{origpage-106}我们从 \hyperref[post:III.1.W]{W.} 推出了 \hyperref[post:III.1.E1]{$E_1$}，又从 \hyperref[post:III.1.E1]{$E_1$} 推出了 \hyperref[post:III.1.E2]{$E_2$}。

将$f(\mathfrak R)$的算符记作$S$，比较 \hyperref[post:III.1.E1]{$E_1$}、\hyperref[post:III.1.E2]{$E_2$}，可知
\[
(S\varphi,\varphi)=(F(R)\varphi,\varphi)
\]
对一切状态$\varphi$成立，也就是对一切满足$\|\varphi\|=1$的$\varphi$成立。因此，一般地有
\[
(Sf,f)=(F(R)f,f)
\]
当$f=0$时显然成立；其余情形取$\varphi=f/\|f\|$即可。由此还可得
\[
(Sf,g)=(F(R)f,g).
\]
只须先分别以$(f+g)/2$和$(f-g)/2$代替$f$，再将两式相减，就得到实部相等；以$if,\allowbreak g$代替$f,\allowbreak g$，便得到虚部相等。因此，必须有$S=F(R)$。将这一重要结果单列如下：

\leavevmode\phantomsection\label{post:III.1.F}\textbf{（F.）}若物理量$\mathfrak R$对应算符$R$，则物理量$F(\mathfrak R)$必须对应算符$F(R)$。

不过，在假定 \hyperref[post:III.1.F]{F.} 成立时，显然也能从 \hyperref[post:III.1.E2]{$E_2$} 推出 \hyperref[post:III.1.E1]{$E_1$}。

因此，在假定 \hyperref[post:III.1.F]{F.} 成立时，\hyperref[post:III.1.E1]{$E_1$} 与 \hyperref[post:III.1.E2]{$E_2$} 等价。现在要证明，\hyperref[post:III.1.W]{W.} 也与它们等价。由于前面已从 \hyperref[post:III.1.W]{W.} 推出了它们，只须证明能够从 \hyperref[post:III.1.E1]{$E_1$} 或 \hyperref[post:III.1.E2]{$E_2$} 推出 \hyperref[post:III.1.W]{W.}。

设$R_1,\allowbreak \ldots,\allowbreak R_l$为彼此对易的算符，分别对应物理量$\mathfrak R_1,\allowbreak \ldots,\allowbreak \mathfrak R_n$。根据 \hyperref[sec:II.10]{II.10}，它们都是某个厄米算符$R$的函数：
\[
R_1=F_1(R),\quad\ldots,\quad R_l=F_l(R).
\]
我们假定，$R$也对应一个物理量$\mathfrak R$。（这里所作的假定是：每个物理量$\mathfrak R$都对应一个〔超极大的〕厄米算符$R$，反过来也一样。参见原注 \hyperlink{note-119}{\getrefnumber{orig-note-119}} 及 \hyperref[sec:IV.2]{IV.2}。）那么，根据 \hyperref[post:III.1.F]{F.}，
\[
\mathfrak R_1=F_1(\mathfrak R),\quad\ldots,\quad\mathfrak R_l=F_l(\mathfrak R).
\]
设$I_1,\allowbreak \ldots,\allowbreak I_l$为 \hyperref[post:III.1.W]{W.} 中出现的区间，并令
\[
G_j(\lambda)=\begin{cases}1,&\lambda\mathwordspace{}\text{在}\mathwordspace{}I_j,\\0,&\text{其他情形},\end{cases}
\qquad(j=1,\ldots,l).
\]
再置
\[
H(\lambda)=G_1(F_1(\lambda))\cdots G_l(F_l(\lambda))
\]
并构造物理量
\[
\mathfrak S=H(\mathfrak R).
\]
若$\mathfrak R_j$落在$I_j$内，也就是$F_j(\mathfrak R)$落在$I_j$内，则$G_j(F_j(\mathfrak R))=1$；其余情形等于$0$。所以，只有当所有$\mathfrak R_j$都落在各自的$I_j$内（$j=1,\allowbreak \ldots,\allowbreak l$）时，$\mathfrak S=H(\mathfrak R)$才等于$1$，否则等于$0$。\hypertarget{origpage-107}{}\phantomsection\label{origpage-107}因此，$\mathfrak S$的期望值就等于$\mathfrak R_1$落在$I_1$内、……、$\mathfrak R_l$落在$I_l$内的概率$W$。于是
\begin{align*}
W&=\operatorname{Erw}(\mathfrak S,\varphi)=(H(R)\varphi,\varphi)\\
&=(G_1(F_1(R))\cdots G_l(F_l(R))\varphi,\varphi)
=(G_1(R_1)\cdots G_l(R_l)\varphi,\varphi).
\end{align*}
将$R_j$对应的单位分解仍记作$E_j(\lambda)$，并设$I_j$为区间$\{\lambda'_j,\allowbreak \lambda''_j\}$，则由 \hyperref[sec:II.8]{II.8} 末尾的讨论及其中的记号可得
\begin{align*}
G_j(\lambda)&=e_{\lambda''_j}(\lambda)-e_{\lambda'_j}(\lambda),\\
G_j(R_j)&=e_{\lambda''_j}(R_j)-e_{\lambda'_j}(R_j)
=E_j(\lambda''_j)-E_j(\lambda'_j)=E_j(I_j),
\end{align*}
故
\[
W=(E_1(I_1)\cdots E_l(I_l)\varphi,\varphi).
\]
这正是 \hyperref[post:III.1.W]{W.}。

\hyperref[post:III.1.E2]{$E_2$} 和 \hyperref[post:III.1.F]{F.} 形式简单，特别适合作为整个理论的基础。我们已经看到，最一般的概率断言 \hyperref[post:III.1.W]{W.} 可以由它们推出。\hyperref[post:III.1.W]{W.} 有两点值得注意：

1．\hyperref[post:III.1.W]{W.} 是统计性的，而非因果性的。它不说明$\mathfrak R_1,\allowbreak \ldots,\allowbreak \mathfrak R_l$在状态$\varphi$中具有什么确定的值，只说明它们以多大的概率取各个可能值。

2．\hyperref[post:III.1.W]{W.} 所提出的问题并不能对任意物理量$\mathfrak R_1,\allowbreak \ldots,\allowbreak \mathfrak R_l$求解，而只能对那些算符$R_1,\allowbreak \ldots,\allowbreak R_l$彼此对易的物理量求解。

接下来的任务，就是讨论这两件事实的意义。

\section{统计解释}\label{sec:III.2}
经典力学是一门因果性的学科。这就是说，只要准确知道系统的状态，就能用确定的数值唯一地给出每个物理量——能量、力矩等等——的值。对于一个有$k$个自由度的系统，确定状态需要$2k$个数值：位形空间的$k$个坐标$q_1,\allowbreak \ldots,\allowbreak q_k$及它们的$k$个时间导数$\frac{\partial q_1}{\partial t},\allowbreak \ldots,\allowbreak \frac{\partial q_k}{\partial t}$；也可以用$k$个动量$p_1,\allowbreak \ldots,\allowbreak p_k$代替这些导数。尽管如此，经典力学也有统计处理方法，不过这可以说是额外添上的：如果并不知道全部$2k$个确定状态的量$(q_1,\allowbreak \ldots,\allowbreak q_k,\allowbreak p_1,\allowbreak \ldots,\allowbreak p_k)$，而只知道其中一部分，甚至连这一部分也不完全准确，那么以某种方式对其余未知量求平均，至少还能对所有物理量作出统计断言。系统先前或以后的状态也是如此：\hypertarget{origpage-108}{}\phantomsection\label{origpage-108}如果知道$t=t_0$时的$q_1,\allowbreak \ldots,\allowbreak q_k,\allowbreak p_1,\allowbreak \ldots,\allowbreak p_k$，就能利用经典力学的运动方程，因果地算出它们在任意其他时刻$t$的值；如果只知道其中一部分，就必须对其余量求平均，对于其他时刻便只能作出统计断言\origNote[言][。]{120}。

量子力学中的统计断言却有另一种性质。在这里，系统若有$k$个自由度，其状态便由波函数$\varphi(q_1\cdots q_k)$描述，也就是由某种适当实现方式下的$\mathfrak R_\infty$中的一个点$\varphi$描述；$\|\varphi\|=1$，而绝对值为$1$的数因子无关紧要。虽然我们认为给定$\varphi$便完全知道了状态，但它对物理量的取值仍然只能给出统计断言。

不过，这种统计性只限于对物理量取值的预言；先前和以后的状态$\varphi_t$却都能由$\varphi_{t_0}=\varphi$因果地计算出来。含时Schrödinger微分方程（参见 \hyperref[sec:I.2]{I.2}）使这一点成为可能，因为
\[
\varphi_{t_0}=\varphi,\qquad
\frac h{2\pi i}\frac\partial{\partial t}\varphi_t=-\mathsf H\varphi_t
\]
确定了$\varphi_t$的整个变化过程。这个微分方程甚至可以显式求解：
\[
\varphi_t=e^{-\frac{2\pi i}h(t-t_0)\mathsf H}\varphi
\]
其中$e^{-\frac{2\pi i}h(t-t_0)\mathsf H}$是酉算符\origNote[符][。]{121}。（此式假定$\mathsf H$与时间无关；但即使$\mathsf H$依赖于时间，$\varphi_t$仍能唯一确定，因为这里的微分方程是一阶方程。只不过这时不再有如此简单的求解公式。）

如果仿照经典力学，要把$\varphi$与物理量取值之间的联系解释为因果性的，显然可以采取如下看法：实际上，$\varphi$并没有准确确定状态；要完全知道状态，还需要另外一些数值。也就是说，除了$\varphi$以外，系统还有其他确定状态的量，即其他坐标。假如知道了全部坐标，就能准确而确定地给出一切物理量的值；仅凭$\varphi$，则和经典力学中只知道$q_1,\allowbreak \ldots,\allowbreak q_k,\allowbreak p_1,\allowbreak \ldots,\allowbreak p_k$的一部分一样，只能作出统计断言。这一看法当然只是假说，是一种解释尝试。它的价值取决于能否真的找出那些补充$\varphi$的坐标，并借此建立一个与经验相符的因果理论；\hypertarget{origpage-109}{}\phantomsection\label{origpage-109}而在只给出$\varphi$、对其余坐标求平均时，这个理论又必须能给出量子力学的统计断言。

这些假定存在的附加坐标，除了迄今的研究所揭示的$\varphi$以外，必定还起着相当隐蔽的作用，因此通常称为“隐坐标”或“隐参数”。在经典物理学中，借助隐参数的解释已经把不少表面上的统计行为归结为力学的因果基础；气体分子运动论就是一个典型例子，参见原注 \hyperlink{note-120}{\getrefnumber{orig-note-120}}。

这种隐参数解释能否用于量子力学，是一个屡经讨论的问题。认为这个问题总有一天会得到肯定回答的人，即使在今天也有杰出的代表。如果这种观点成立，理论目前的形式就只能是暂时的，因为用$\varphi$来描述状态将是根本不完全的。

后面我们将证明（\hyperref[sec:IV.2]{IV.2}），除非对现有理论作实质性的改变，否则确实无法引入隐参数。这里先指出，$\varphi$与经典力学中$q_1,\allowbreak \ldots,\allowbreak q_k,\allowbreak p_1,\allowbreak \ldots,\allowbreak p_k$的一个子集之间，有一点根本的区别：$\varphi$对时间的依赖关系是因果性的，而非统计性的。如上所见，$\varphi_{t_0}$唯一地确定了所有$\varphi_t$。

要客观检验能否引入隐参数，还须更深入地分析量子力学的断言；这项检验将在前面提到的地方进行。在此以前，我们暂不采用这种解释方式，而采取相反的立场：承认支配基本过程的自然规律，也就是量子力学的定律，本来就是统计性的。（宏观世界所呈现的因果性，至多可能只是一种表象，来自“大数定律”的平均效应；这种效应在许多基本过程同时相互作用时便会显现。参见原注 \hyperlink{note-120}{\getrefnumber{orig-note-120}} 末尾的说明及原注 \hyperlink{note-175}{\getrefnumber{orig-note-175}}。）据此，我们把 \hyperref[post:III.1.W]{W.}，或者 \hyperref[post:III.1.E2]{$E_2$}，看作对基本过程所能作出的最充分的断言。

这种量子力学的理解方式，承认它的统计断言就是自然规律的真正形式，并放弃因果性原则，称为统计解释。这一解释由M.~Born提出\origNote[出][。]{122}。它是目前唯一能够贯彻到底的量子力学解释——这里的量子力学指我们关于基本过程的经验的总和。下面我们将采用这种解释，但不能在这里详细讨论各种解释之间的关系。

\hypertarget{origpage-110}{}\phantomsection\label{origpage-110}\section{同时可测性与一般的可测性}\label{sec:III.3}
\hyperref[sec:III.1]{III.1} 末尾所说的第二个值得注意之处，与这样一件事有关：\hyperref[post:III.1.W]{W.} 不仅给出一个物理量$\mathfrak R$取某些数值的概率，也说明几个物理量$\mathfrak R_1,\allowbreak \ldots,\allowbreak \mathfrak R_l$之间的概率联系。因为 \hyperref[post:III.1.W]{W.} 给出的是它们同时取某些给定值——更确切地说，分别落在某些区间$I_1,\allowbreak \ldots,\allowbreak I_l$内——的概率。（这一切都针对某个给定状态$\varphi$。）可是，物理量$\mathfrak R_1,\allowbreak \ldots,\allowbreak \mathfrak R_l$受到一个特别的限制：其算符$R_1,\allowbreak \ldots,\allowbreak R_l$必须彼此对易。当$R_1,\allowbreak \ldots,\allowbreak R_l$不对易时，\hyperref[post:III.1.W]{W.} 完全不能说明$\mathfrak R_1,\allowbreak \ldots,\allowbreak \mathfrak R_l$之间的相互概率依赖关系；它只能用来分别确定各个物理量的概率分布，而不顾及其他物理量。

最自然的想法，是认为这是 \hyperref[post:III.1.W]{W.} 的一项缺陷，必定还存在某个更一般的公式，能够包括这些情形。因为即使量子力学只能从统计上描述自然，我们至少仍有理由要求它既能描述单个物理量的统计规律，也能把握几个物理量之间的关联。

这种看法初看似乎合情合理，但我们很快就会发现，\hyperref[post:III.1.W]{W.} 不可能作这样的推广。除了由理论的数学工具的结构造成的形式原因之外，还有重大的物理原因，要求作出这种限制。认识到这一限制的必要性，并对它作出物理解释，甚至能使我们深入理解基本过程的本性。

要弄清楚这一点，必须进一步考察：测量一个物理量$\mathfrak R$这一过程，对于量子力学的描述方式究竟意味着什么；\hyperref[post:III.1.W]{W.} 所作的〔概率〕断言，正是关于这样的物理量。

先来看Compton和Simons在量子力学建立以前进行的一个重要实验\origNote[验][。]{123}。实验让光在电子上散射，随后截获散射光和散射电子，并测量它们的能量与动量，从而对散射过程作定量检验。也就是说，这里发生的是光量子与电子之间的碰撞，\hypertarget{origpage-111}{}\phantomsection\label{origpage-111}观察者可以根据碰撞后的轨迹，检验弹性碰撞定律是否成立。（这里不考虑其他形式的碰撞，因为当时认为，无论电子还是光量子，都不能以动能以外的形式吸收能量。根据一切已有经验，电子和光量子的结构似乎都是完全固定的！碰撞的计算当然必须采用相对论形式\origNote[式][。]{123}。）这种计算检验确实可行：碰撞前的轨迹已知，碰撞后的轨迹由实验观测，因此碰撞问题的条件是超定的。要用力学确定碰撞过程，只需四条轨迹中的两条，加上碰撞的“中心线”，即动量传递的方向；因而至多只需三条轨迹，第四条便可用作检验。实验完全证实了力学的碰撞定律。

承认碰撞定律有效，并认为碰撞前的轨迹已知，就可以把这一结果改述如下：碰撞后光量子的轨迹和碰撞后电子的轨迹，任何一条都足以算出碰撞的位置及中心线。Compton–Simons实验表明，这两种观测给出了同一个结果。

更一般地说：同一个物理量，即碰撞位置的某个坐标或中心线的方向，被用两种不同的方式测量——分别截获光量子和电子——而所得结果总是相同的。

这两次测量并不完全同时发生。光量子和电子不会同时到达；适当安排测量装置，就能先截获想先截获的那一个，当然，两者只相差$10^{-9}$至$10^{-10}$秒。将先进行的测量记作$M_1$，后进行的记作$M_2$，所测物理量记作$\mathfrak R$。关键在于：由于整个测量安排的性质，在测量之前，我们只能从统计上预言$\mathfrak R$的值，也就是$M_1,\allowbreak M_2$的结果（尤其参见原注 \hyperlink{note-123}{\getrefnumber{orig-note-123}}）；然而，$M_1,\allowbreak M_2$之间的统计关联却是完全确定的，即因果性的：$M_1$确定的$\mathfrak R$值与$M_2$确定的值必定相等。在$M_1,\allowbreak M_2$测量以前，两次测量的结果都在很大程度上尚未确定；可是一旦$M_1$完成，即使$M_2$尚未进行，$M_2$的结果也已经由因果关系唯一确定了。

从原则上说，可以设想因果性或非因果性有三个层次。第一种可能是，$\mathfrak R$的值完全是统计性的：一次测量的结果只能从统计上预言；即使紧接着第一次测量再作第二次测量，第二次的结果仍可能有一定的散布程度，不受第一次测得的值的约束，例如，散布程度可能与第一次一样大\origNote[大][。]{124}。\hypertarget{origpage-112}{}\phantomsection\label{origpage-112}第二种可能是，$\mathfrak R$在第一次测量时所得的值有一定的散布程度，但紧随其后的每次测量都必须给出与第一次相同的结果。第三种可能是，$\mathfrak R$从一开始就由因果关系确定了。

Compton–Simons实验表明，对于一种统计理论，只有第二种情形才有可能成立。因此，设系统最初所处的状态不能让我们确切预言$\mathfrak R$的值，那么，测量$\mathfrak R$的一次过程$M$——上例中的$M_1$——就会把这个状态变为另一个状态：在新状态中，$\mathfrak R$的值已经唯一确定。而且，$M$使系统进入哪一个新状态，不仅取决于$M$的装置，也取决于$M$的测量结果；这个结果在原有状态中是不能由因果关系预言的。因为$\mathfrak R$在新状态中的值，必须等于这次$M$的测量结果。

设$\mathfrak R$为一个物理量，其算符$R$具有纯点谱$\lambda_1,\allowbreak \lambda_2,\allowbreak \ldots$，相应的本征函数$\varphi_1,\allowbreak \varphi_2,\allowbreak \ldots$构成完全规范正交系。并设每个本征值都是单重的，即重数为$1$，参见 \hyperref[sec:II.6]{II.6}；也就是说，当$\mu\ne\nu$时，$\lambda_\mu\ne\lambda_\nu$。假定测量$\mathfrak R$所得的值为$\lambda^*$。测量以后，系统处于什么状态？

根据上面的讨论，这个状态$\varphi$必须满足：重新测量$\mathfrak R$，一定得到结果$\lambda^*$。（这次测量当然必须紧接着进行，因为经过$\tau$秒，$\varphi$便变成了$e^{-\frac{2\pi i}h\tau\mathsf H}\varphi$。参见 \hyperref[sec:III.2]{III.2}，$\mathsf H$为能量算符。）

对于什么状态$\varphi$，测量$\mathfrak R$必定得到$\lambda^*$？下面来回答这个问题，而且不再对算符$R$加上任何限制。

设$E(\lambda)$为$R$对应的单位分解，$I$为区间$\{\lambda',\allowbreak \lambda''\}$。上述假定也可以这样表述：如果$I$不包含$\lambda^*$，$\mathfrak R$落在$I$内的概率为$0$；如果$\lambda^*$属于$I$，即$\lambda'<\lambda^*\le\lambda''$，该概率为$1$。

根据 \hyperref[post:III.1.W]{W.}，这意味着$\|E(I)\varphi\|^2=1$。由于$\|\varphi\|=1$，便有$\|E(I)\varphi\|=\|\varphi\|$。$E(I)$为投影算符，$1-E(I)$也是，故由 \hyperref[thm:II.13]{II.4的定理13} 可得
\begin{gather*}
\|\varphi-E(I)\varphi\|^2=\|\varphi\|^2-\|E(I)\varphi\|^2=0,\\
\varphi-E(I)\varphi=0,\\
E(\lambda'')\varphi-E(\lambda')\varphi=E(I)\varphi=\varphi.
\end{gather*}
令$\lambda'\to-\infty$，得到$E(\lambda'')\varphi=\varphi$；令$\lambda''\to+\infty$，得到$E(\lambda')\varphi=0$。\hypertarget{origpage-113}{}\phantomsection\label{origpage-113}参见 \hyperref[sec:II.7]{II.7的$S_1$}。因此
\[
E(\lambda)\varphi=\begin{cases}\varphi,&\lambda\ge\lambda^*,\\0,&\lambda<\lambda^*.\end{cases}
\]
根据 \hyperref[sec:II.8]{II.8}，这正是$R\varphi=\lambda^*\varphi$的特征。

也可以从 \hyperref[post:III.1.E1]{$E_1$}，即 \hyperref[post:III.1.E2]{$E_2$}，用另一条途径证明$R\varphi=\lambda^*\varphi$。$\mathfrak R$确定取值$\lambda^*$，就意味着$(\mathfrak R-\lambda^*)^2$的期望值为$0$。于是，令$F(\lambda)=(\lambda-\lambda^*)^2$，其算符$F(R)$就是$(R-\lambda^*\cdot1)^2$，必须满足
\begin{align*}
((R-\lambda^*\cdot1)^2\varphi,\varphi)
&=((R-\lambda^*\cdot1)\varphi,(R-\lambda^*\cdot1)\varphi)\\
&=\|(R-\lambda^*\cdot1)\varphi\|^2
=\|R\varphi-\lambda^*\varphi\|^2=0,
\end{align*}
即$R\varphi=\lambda^*\varphi$。

回到刚才讨论的特殊情形，我们看到，必须有$R\varphi=\lambda^*\varphi$；按照 \hyperref[sec:II.6]{II.6} 的讨论，这意味着$\lambda^*$等于某个$\lambda_\mu$，因为$\|\varphi\|=1$，所以$\varphi\ne0$，并且$\varphi=a\varphi_\mu$。又因$\|\varphi\|=\|\varphi_\mu\|=1$，必须有$|a|=1$，所以略去$a$并不改变状态。因此，对于某个$\mu=1,\allowbreak 2,\allowbreak \ldots$，有$\lambda^*=\lambda_\mu$、$\varphi=\varphi_\mu$。（关于$\lambda^*$的断言，也可以直接从 \hyperref[post:III.1.W]{W.} 推出；但关于$\varphi$的断言却不能如此得出！）

可见，在上述关于$R$的假定下，测量$\mathfrak R$会把每个状态$\psi$变为状态$\varphi_1,\allowbreak \varphi_2,\allowbreak \ldots$中的某一个，它们分别与测量结果$\lambda_1,\allowbreak \lambda_2,\allowbreak \ldots$相联系。因此，这些状态转变的概率就等于测出$\lambda_1,\allowbreak \lambda_2,\allowbreak \ldots$的概率，可以由 \hyperref[post:III.1.W]{W.} 计算。

事实上，根据 \hyperref[post:III.1.W]{W.}，$\mathfrak R$的值落在$I$内的概率为$\|E(I)\psi\|^2$。注意到 \hyperref[sec:II.8]{II.8} 中的$E(I)=\sum_{\lambda_n\mathwordspace{}\text{在}\mathwordspace{}I}P_{[\varphi_n]}$，便有
\begin{align*}
W&=\|E(I)\psi\|^2=(E(I)\psi,\psi)
=\sum_{\lambda_n\mathwordspace{}\text{在}\mathwordspace{}I}(P_{[\varphi_n]}\psi,\psi)\\
&=\sum_{\lambda_n\mathwordspace{}\text{在}\mathwordspace{}I}|(\psi,\varphi_n)|^2.
\end{align*}
因此，可以猜想$\lambda_n$的概率为$|(\psi,\allowbreak \varphi_n)|^2$。如果能选取区间$I$，使其中唯一的$\lambda_m$就是$\lambda_n$，这便直接由上式得出。其余情形，即其他$\lambda_m$在$\lambda_n$处聚集时，可以作如下推论。取$F(\lambda)$在$\lambda=\lambda_n$处等于$1$，其余地方等于$0$，则所求概率$W_n$就是$F(\mathfrak R)$的期望值，由 \hyperref[post:III.1.E2]{$E_2$} 或 \hyperref[post:III.1.E1]{$E_1$} 可知，它等于$(F(R)\psi,\allowbreak \psi)$。根据定义（\hyperref[sec:II.8]{II.8}），
\[
(F(R)\psi,\psi)=\int_{-\infty}^{\infty}F(\lambda)\,d(\|E(\lambda)\psi\|^2).
\]
回想Stieltjes积分的定义就容易看出：如果$E(\lambda)\psi$在$\lambda=\lambda_n$处关于$\lambda$连续，此积分为$0$；\hypertarget{origpage-114}{}\phantomsection\label{origpage-114}一般而言，它等于单调不减函数$\|E(\lambda)\psi\|^2$在$\lambda=\lambda_n$处的跳跃值。该跳跃值等于$\|P_{\mathfrak M}\psi\|^2$，其中$\mathfrak M$是由$R\psi=\lambda_n\psi$的全部解张成的闭线性流形，参见 \hyperref[sec:II.8]{II.8}。当前的情形下，$\mathfrak M=[\varphi_n]$，故
\[
W_n=\|P_{[\varphi_n]}\psi\|^2=|(\psi,\varphi_n)|^2.
\]
这样，在上述关于算符$R$的假定下，测量物理量$\mathfrak R$时发生什么，已经有了回答。不过，至于这一过程“怎样”发生，目前仍不清楚：$\psi$不连续地、跳跃式地变为$\varphi_1,\allowbreak \varphi_2,\allowbreak \ldots$中的某个状态，这种变化显然不同于含时Schrödinger微分方程所描述的变化。各个$\varphi_n$本身与$\psi$无关，只有各次跃变的概率$W_n=|(\psi,\allowbreak \varphi_n)|^2$，$n=1,\allowbreak 2,\allowbreak \ldots$，才取决于$\psi$；微分方程所描述的却总是$\psi$的连续变化，终态唯一确定，并依赖于$\psi$。参见 \hyperref[sec:III.2]{III.2}。我们将在后面尝试弥合这一鸿沟，参见\CJKglue \hyperref[chap:VI]{第VI章}\origNote[章][。]{125}。

仍假定$R$具有纯点谱，但不再假定所有本征值都是单重的。这时仍能构造$\varphi_1,\allowbreak \varphi_2,\allowbreak \ldots$及$\lambda_1,\allowbreak \lambda_2,\allowbreak \ldots$，只是$\lambda_n$中允许有相等的值。测量$\mathfrak R$之后，系统必处于某个满足$R\varphi=\lambda^*\varphi$的状态$\varphi$，其中$\lambda^*$就是测量结果。因此，$\lambda^*$等于某个$\lambda_n$；但关于$\varphi$，我们只能说：设等于$\lambda^*$的那些$\lambda_n$为$\lambda_{n_1},\allowbreak \lambda_{n_2},\allowbreak \ldots$，其个数有限或无限，则
\[
\varphi=\sum_\nu a_\nu\varphi_{n_\nu}.
\]
如果$n_\nu$有无限多个，就必须有$\sum_\nu|a_\nu|^2$有限。两个这样的$\varphi$，只有在它们相差一个数因子时，也就是比值$a_1:a_2:\cdots$完全相同时，才表示同一个状态。所以，只要$n_\nu$不止一个，即本征值$\lambda^*$是多重的，即使知道了测量结果，也不能唯一确定测量后的状态$\varphi$。

$\lambda^*$的概率可以根据 \hyperref[post:III.1.W]{W.}，或 \hyperref[post:III.1.E1]{$E_1$}、\hyperref[post:III.1.E2]{$E_2$}，完全仿照上面的方式算出：
\[
W(\lambda^*)=\sum_{\lambda_n=\lambda^*}|(\psi,\varphi_n)|^2
=\sum_\nu|(\psi,\varphi_{n_\nu})|^2.
\]
如果$R$没有纯点谱，情况如下。$Rf=\lambda f$的全部解$f$张成一个闭线性流形$\mathfrak M_\lambda$。所有$\mathfrak M_\lambda$又共同张成一个闭线性流形$\overline{\mathfrak M}$；不具有纯点谱的特征正是$\overline{\mathfrak M}\ne\mathfrak R_\infty$，即$\overline{\mathfrak N}=\mathfrak R_\infty-\overline{\mathfrak M}\ne(0)$。参见 \hyperref[sec:II.8]{II.8} 及下面的讨论。至多只有可数多个$\lambda$满足$\mathfrak M_\lambda\ne(0)$，\hypertarget{origpage-115}{}\phantomsection\label{origpage-115}它们构成$R$的点谱。如果在状态$\psi$中测量$\mathfrak R$，结果为$\lambda^*$的概率为
\[
W(\lambda^*)=\|P_{\mathfrak M_{\lambda^*}}\psi\|^2
=(P_{\mathfrak M_{\lambda^*}}\psi,\psi).
\]
适宜的证明方法，仍是上面利用 \hyperref[post:III.1.E2]{$E_2$} 或 \hyperref[post:III.1.E1]{$E_1$} 以及函数
\[
F(\lambda)=\begin{cases}1,&\lambda=\lambda^*,\\0,&\lambda\ne\lambda^*\end{cases}
\]
的推论。$\mathfrak R$的值为$R$的点谱$\mathfrak P$中任一$\lambda^*$的概率，相应地为
\[
W=\sum_{\lambda^*\mathwordspace{}\text{在}\mathwordspace{}\mathfrak P}(P_{\mathfrak M_{\lambda^*}}\psi,\psi)
=(P_{\overline{\mathfrak M}}\psi,\psi)=\|P_{\overline{\mathfrak M}}\psi\|^2.
\]
借助函数
\[
F(\lambda)=\begin{cases}1,&\lambda^*\mathwordspace{}\text{在}\mathwordspace{}\mathfrak P,\\0,&\text{其他情形}\end{cases}
\]
也能直接看出这一点。

可是，假如我们准确地测量$\mathfrak R$，测量以后就必须处于满足$R\varphi=\lambda^*\varphi$的某个状态$\varphi$，因而测量结果$\lambda^*$必须属于$\mathfrak P$。所以，准确测量能够成功的概率，至多为$\|P_{\overline{\mathfrak M}}\psi\|^2$。这个数并不总等于$1$；当$\psi$属于$\overline{\mathfrak N}$时，它甚至等于$0$。因此，准确测量并非总是可能的。

我们已经看到，一个物理量$\mathfrak R$能够始终——即在每个状态$\psi$中——被准确测量，当且仅当它具有纯点谱。若不具有纯点谱，就只能作精度有限的测量：把数轴$-\infty<\lambda<+\infty$划分为区间$\ldots,\allowbreak I^{(-2)},\allowbreak I^{(-1)},\allowbreak I^{(0)},\allowbreak I^{(1)},\allowbreak I^{(2)},\allowbreak \ldots$，分点为$\ldots,\allowbreak \lambda^{(-2)},\allowbreak \lambda^{(-1)},\allowbreak \lambda^{(0)},\allowbreak \lambda^{(1)},\allowbreak \lambda^{(2)},\allowbreak \ldots$，使$I^{(n)}=\{\lambda^{(n)},\allowbreak \lambda^{(n+1)}\}$，最大区间长度，即划分的网格宽度$\varepsilon=\operatorname{Max}(\lambda^{(n+1)}-\lambda^{(n)})$，就是测量精度；然后判断$\mathfrak R$落在哪一个区间。这个过程还可以进一步计算。取如下函数$F(\lambda)$，其中$\lambda'_n$为$I^{(n)}$内的某个中间值，对每个$n=0,\allowbreak \pm1,\allowbreak \pm2,\allowbreak \ldots$都可任意选定，但在后面的讨论中固定不变：
\[
F(\lambda)=\lambda'_n\quad\text{若}\mathwordspace{}\lambda\mathwordspace{}\text{在}\mathwordspace{}I^{(n)}.
\]
那么，$\mathfrak R$的近似测量就等同于$F(\mathfrak R)$的准确测量。现在有
\begin{align*}
F(R)&=\int_{-\infty}^{\infty}F(\lambda)\,dE(\lambda)
=\sum_{n=-\infty}^{+\infty}\int_{\lambda^{(n)}}^{\lambda^{(n+1)}}F(\lambda)\,dE(\lambda)\\
&=\sum_{n=-\infty}^{+\infty}\lambda'_n\int_{\lambda^{(n)}}^{\lambda^{(n+1)}}dE(\lambda)
=\sum_{n=-\infty}^{+\infty}\lambda'_nE(I^{(n)}).
\end{align*}\hypertarget{origpage-116}{}\phantomsection\label{origpage-116}显然，$E(I^{(n)})$对应的闭线性流形中，所有$f$都满足$F(R)f=\lambda'_nf$；所以，对于$F(R)$，$\mathfrak M_{\lambda'_n}$包含这个流形。因此$P_{\mathfrak M_{\lambda'_n}}\ge E(I^{(n)})$，从而
\begin{align*}
P_{\overline{\mathfrak M}}
&\ge\sum_{n=-\infty}^{+\infty}P_{\mathfrak M_{\lambda'_n}}
\ge\sum_{n=-\infty}^{+\infty}E(I^{(n)})\\
&=\sum_{n=-\infty}^{+\infty}(E(\lambda^{(n+1)})-E(\lambda^{(n)}))=1-0=1.
\end{align*}
于是
\[
\sum_{n=-\infty}^{+\infty}P_{\mathfrak M_{\lambda'_n}}=P_{\overline{\mathfrak M}}=1,
\qquad P_{\mathfrak M_{\lambda'_n}}=E(I_n),
\]
也就是说，$F(R)$具有纯点谱，其点谱由$\lambda'_n$构成。

所以，$F(\mathfrak R)$确实可以准确测量；其值等于$\lambda'_n$，也就是$\mathfrak R$落在$I^{(n)}$内的概率为
\[
\|P_{\mathfrak M_{\lambda'_n}}\psi\|^2=\|E(I^{(n)})\psi\|^2,
\]
与 \hyperref[post:III.1.W]{W.} 对$\mathfrak R$所作的断言相符。

这个结果也可以从物理上解释，并显示出理论与通常物理直观之间的良好一致性。

按照不附加任何量子条件的经典力学观点，每个状态下的每个物理量$\mathfrak R$固然都有一个完全确定的值；但同时也承认，任何可以设想的测量仪器，都只能在某个始终不为零的误差范围内测出这个值。这是因为人的观察能力并不完善，例如读出指针或确定照相底片上感光痕迹的位置，都只能达到有限的精度。不断改进测量装置，固然可以把这个误差界限降到任意接近$0$，但它永远不会恰好等于$0$。对于那些按惯常直观图景——尤其在量子力学发现以前的图景——并未“量子化”的物理量$\mathfrak R$，人们也会预期量子理论中有同样的情形。例如，电子的直角坐标可以取$-\infty$与$+\infty$之间的每个值，其算符具有连续谱。相反，对那些按照直观图景已经“量子化”的物理量，情况却不同：它们只能取离散值，所以只要观察精确到足以消除疑问，确定究竟取了哪个“量子化”值，就能绝对准确地给出这个值。例如，知道一个氢原子的能量低于次低能级所需的能量，就绝对准确地知道了它的能量：它处于最低能级。

\hypertarget{origpage-117}{}\phantomsection\label{origpage-117}然而，正如“矩阵理论”的情况已经表明的（参见 \hyperref[sec:I.2]{I.2} 及 \hyperref[sec:II.6]{II.6}），这种“量子化”与“未量子化”的区分，正好对应于物理量$\mathfrak R$的算符$R$是否具有纯点谱。我们恰恰只对前一种物理量发现了绝对准确观测的可能；后一种物理量则只能以任意高、却从来不能绝对准确的精度加以观测\origNote[测][。]{126}。

顺便指出，\CJKglue \hyperref[chap:introduction]{引言}\CJKglue 及 \hyperref[sec:I.3]{I.3} 中提到的引入“广义”本征函数、即不属于希尔伯特空间的本征函数的做法，反倒不如我们的方法忠实于实际。参见 \hyperref[sec:II.8]{II.8}，尤其是原注 \hyperlink{note-84}{\getrefnumber{orig-note-84}}、\hyperlink{note-86}{\getrefnumber{orig-note-86}}。因为这种做法会造成假象，仿佛存在某些状态，使具有连续谱的物理量也能准确地取某些值；而这恰恰从不发生。我们认为，这类理想化不但在数学上站不住脚，出于这一物理原因也必须拒绝，尽管已有多人提出过这种做法。

至此，对于测量一个物理量时发生的过程，我们已有一个暂时的结论。现在可以转向同时测量几个物理量的问题。

先给定两个物理量$\mathfrak R,\allowbreak \mathfrak S$，其算符分别为$R,\allowbreak S$，并假定它们可以同时测量。由此能推出什么？

先要求测量绝对准确，因此$R,\allowbreak S$必须都具有纯点谱，分别为$\lambda_1,\allowbreak \lambda_2,\allowbreak \ldots$及$\mu_1,\allowbreak \mu_2,\allowbreak \ldots$。相应的本征函数构成完全规范正交系，分别记为$\varphi_1,\allowbreak \varphi_2,\allowbreak \ldots$及$\psi_1,\allowbreak \psi_2,\allowbreak \ldots$。

先从最简单的情形入手。假定两个算符之一，例如$R$，全都只有单重本征值；也就是说，当$m\ne n$时，$\lambda_m\ne\lambda_n$。

同时测量$\mathfrak R,\allowbreak \mathfrak S$以后，系统处于一个状态，其中$\mathfrak R$和$\mathfrak S$都确定地具有刚才测得的值。设这些值为$\lambda_{\bar m},\allowbreak \mu_{\bar n}$，测量后的状态为$\psi$，则必须满足$R\psi=\lambda_{\bar m}\psi$、$S\psi=\mu_{\bar n}\psi$。由前一式得到$\psi=\varphi_{\bar m}$，略去一个无关紧要的数因子；由后一式得到$\psi=\sum_\nu a_\nu\psi_{n_\nu}$，其中$\mu_{n_1},\allowbreak \mu_{n_2},\allowbreak \ldots$是所有等于$\mu_{\bar n}$的$\mu_n$。如果初态为$\varphi$，则$\lambda_{\bar m},\allowbreak \varphi_{\bar m}$的概率为$|(\varphi,\allowbreak \varphi_{\bar m})|^2$。特别是，当$\varphi=\varphi_m$时，必定有$\bar m=m$。所以，对于每个$m$，$\varphi_m$都能写成$\sum_\nu a_\nu\psi_{n_\nu}$的形式，且其中的$\mu_{n_\nu}$全都相同。即$S\varphi_m=\bar\mu\varphi_m$，其中$\bar\mu=\mu_{n_1}=\mu_{n_2}=\cdots$。因此，对$f=\varphi_m$有$RSf=SRf$，两边都等于$\lambda_m\bar\mu\cdot\varphi_m$。这对它们的线性组合也成立；而当$R,\allowbreak S$连续时，对这些线性组合的极限点也成立，即对所有$f$都成立。因此，$R,\allowbreak S$对易。

\hypertarget{origpage-118}{}\phantomsection\label{origpage-118}如果它们不连续，则作如下推论。$R,\allowbreak S$对应的单位分解$E(\lambda),\allowbreak F(\mu)$定义为
\[
E(\lambda)=\sum_{\lambda_m\le\lambda}P_{[\varphi_m]},\qquad
F(\mu)=\sum_{\mu_n\le\mu}P_{[\psi_n]}.
\]
因此，$F(\mu)\varphi_m$依$\mu\ge\bar\mu$或$\mu<\bar\mu$分别等于$\varphi_m$或$0$，这里$\bar\mu$的意义同上。类似地，$E(\lambda)\varphi_m$依$\lambda\ge\lambda_m$或$\lambda<\lambda_m$分别等于$\varphi_m$或$0$。所以，对所有$\varphi_m$总有$E(\lambda)F(\mu)\varphi_m=F(\mu)E(\lambda)\varphi_m$。仿照上面的论证，这说明$E(\lambda),\allowbreak F(\mu)$对易，从而根据 \hyperref[sec:II.10]{II.10}，$R,\allowbreak S$也对易。

根据 \hyperref[sec:II.10]{II.10}，存在一个由$R,\allowbreak S$的共同本征函数构成的完全规范正交系；也就是说，可以令$\varphi_m=\psi_m$。又因$m\ne n$时$\lambda_m\ne\lambda_n$，可以构造函数$F(\lambda)$：
\[
F(\lambda)=\begin{cases}\mu_n,&\lambda=\lambda_n,\ n=1,2,\ldots,\\
\text{任意},&\text{其他情形},\end{cases}
\]
于是$S=F(R)$，也就是$\mathfrak S=F(\mathfrak R)$。所以，$\mathfrak R,\allowbreak \mathfrak S$不但可以同时测量，甚至每次测量$\mathfrak R$都同时测量了$\mathfrak S$，因为$\mathfrak S$作为$\mathfrak R$的函数，由$\mathfrak R$因果地确定\origNote[定][。]{127}。

现在转向另一个更一般的情形，不对$R,\allowbreak S$的本征值重数作任何假定。这里采用一种根本不同的方法。

先考察物理量$\mathfrak R+\mathfrak S$。同时测量$\mathfrak R,\allowbreak \mathfrak S$，也就是测量$\mathfrak R+\mathfrak S$，因为把两个测量结果相加，就得到$\mathfrak R+\mathfrak S$的值。因此，在每个状态$\psi$中，$\mathfrak R+\mathfrak S$的期望值都等于$\mathfrak R$和$\mathfrak S$各自期望值之和。注意，这一点与$\mathfrak R,\allowbreak \mathfrak S$在统计上是否独立无关，也不取决于它们之间有无关联、有什么关联，甚至不取决于其中一个是否因果地确定另一个。因为
\[
\text{和的期望值}=\text{各期望值之和}
\]
这一规律总是成立。设$T$为$\mathfrak R+\mathfrak S$的算符，则这个期望值一方面等于$(T\psi,\allowbreak \psi)$，另一方面等于
\[
(R\psi,\psi)+(S\psi,\psi)=((R+S)\psi,\psi).
\]
即对所有$\psi$有
\[
(T\psi,\psi)=((R+S)\psi,\psi),
\]
所以$T=R+S$。因此，$\mathfrak R+\mathfrak S$对应算符$R+S$。\origNote[。][同]{128}同样可以证明，$a\mathfrak R+b\mathfrak S$（$a,\allowbreak b$为实数）对应算符$aR+bS$。（也可从第一个公式得到这一点，只须在其中用函数$F(\lambda)=a\lambda$、$G(\mu)=b\mu$分别作用于$\mathfrak R,\allowbreak \mathfrak S$及$R,\allowbreak S$。）

\hypertarget{origpage-119}{}\phantomsection\label{origpage-119}同时测量$\mathfrak R,\allowbreak \mathfrak S$，也就测量了
\begin{gather*}
\frac{\mathfrak R+\mathfrak S}2,\quad
\left(\frac{\mathfrak R+\mathfrak S}2\right)^2,\quad
\frac{\mathfrak R-\mathfrak S}2,\quad
\left(\frac{\mathfrak R-\mathfrak S}2\right)^2,\\
\left(\frac{\mathfrak R+\mathfrak S}2\right)^2-
\left(\frac{\mathfrak R-\mathfrak S}2\right)^2=\mathfrak R\cdot\mathfrak S.
\end{gather*}
所以，这些物理量对应的算符如下；这里还利用了前面的结论：如果$T$是$\mathfrak T$的算符，则$F(\mathfrak T)$的算符为$F(T)$，因而$\mathfrak T^2$的算符为$T^2$：
\begin{gather*}
\frac{R+S}2,\quad
\left(\frac{R+S}2\right)^2=\frac{R^2+S^2+RS+SR}4,\\
\frac{R-S}2,\quad
\left(\frac{R-S}2\right)^2=\frac{R^2+S^2-RS-SR}4,\\
\left(\frac{R+S}2\right)^2-\left(\frac{R-S}2\right)^2=\frac{RS+SR}2.
\end{gather*}
即$\mathfrak R\cdot\mathfrak S$对应算符$(RS+SR)/2$。对于所有$F(\mathfrak R),\allowbreak G(\mathfrak S)$，结论同样成立，因为它们也随之被测量了。因此，$F(\mathfrak R)\cdot G(\mathfrak S)$对应算符
\[
\frac{F(R)G(S)+G(S)F(R)}2.
\]
设$E(\lambda),\allowbreak F(\mu)$为$R,\allowbreak S$对应的单位分解，并令
\[
F(\lambda)=\begin{cases}1,&\lambda\le\bar\lambda,\\0,&\lambda>\bar\lambda,\end{cases}
\qquad
G(\mu)=\begin{cases}1,&\mu\le\bar\mu,\\0,&\mu>\bar\mu.\end{cases}
\]
我们知道，$F(R)=E(\bar\lambda)$、$G(S)=F(\bar\mu)$，所以$F(\mathfrak R)\cdot G(\mathfrak S)$对应算符$(EF+FE)/2$；为简便起见，将$E(\bar\lambda),\allowbreak F(\bar\mu)$写作$E,\allowbreak F$。$F(\mathfrak R)$总是$0$或$1$，所以$F(\mathfrak R)^2=F(\mathfrak R)$，从而
\[
F(\mathfrak R)\cdot(F(\mathfrak R)\cdot G(\mathfrak S))=F(\mathfrak R)\cdot G(\mathfrak S).
\]
现在对$F(\mathfrak R)$和$F(\mathfrak R)\cdot G(\mathfrak S)$应用上面的乘法公式——这些物理量都可以同时测量——便得到这个乘积的算符为
\[
\frac{E\frac{EF+FE}2+\frac{EF+FE}2E}2
=\frac{E^2F+2EFE+FE^2}4
=\frac{EF+FE+2EFE}4.
\]
它必须等于$(EF+FE)/2$，于是
\[
EF+FE=2\cdot EFE.
\]
左乘$E$，得到
\[
E^2F+EFE=2E^2FE,\quad EF+EFE=2EFE,\quad EF=EFE;
\]
右乘$E$，得到
\[
EFE+FE^2=2EFE^2,\quad EFE+FE=2EFE,\quad FE=EFE.
\]
故$EF=FE$。即所有$E(\bar\lambda),\allowbreak F(\bar\mu)$都对易，因此$R,\allowbreak S$也对易。

\hypertarget{origpage-120}{}\phantomsection\label{origpage-120}根据 \hyperref[sec:II.10]{II.10}，$R,\allowbreak S$对易，等价于存在一个厄米算符$T$，使$R,\allowbreak S$都是它的函数：$R=F(T)$、$S=G(T)$。若$T$对应物理量$\mathfrak T$，则有$\mathfrak R=F(\mathfrak T)$、$\mathfrak S=G(\mathfrak T)$。这个条件也是同时可测的充分条件，因为测量$\mathfrak T$时，就同时测量了它的函数$\mathfrak R,\allowbreak \mathfrak S$；而$\mathfrak T$确实可以绝对准确地测量，因为它具有纯点谱，参见 \hyperref[sec:II.10]{II.10}。所以，$R,\allowbreak S$对易既是必要条件，也是充分条件。

给定几个物理量$\mathfrak R,\allowbreak \mathfrak S,\allowbreak \ldots$，其个数有限，算符分别为$R,\allowbreak S,\allowbreak \ldots$，仍要求测量绝对准确，则同时可测性的情形如下。若所有$\mathfrak R,\allowbreak \mathfrak S,\allowbreak \ldots$都能同时测量，从中选出的任意一对也必定如此；所以所有算符$R,\allowbreak S,\allowbreak \ldots$必须彼此对易。反之，若$R,\allowbreak S,\allowbreak \ldots$彼此对易，则根据 \hyperref[sec:II.10]{II.10}，存在算符$T$，使它们都是$T$的函数：$R=F(T)$、$S=G(T)$、……；对相应的物理量$\mathfrak T$，便有$\mathfrak R=F(\mathfrak T)$、$\mathfrak S=G(\mathfrak T)$、……。准确测量$\mathfrak T$，也就同时测量了$\mathfrak R,\allowbreak \mathfrak S,\allowbreak \ldots$；$\mathfrak T$仍具有纯点谱，参见 \hyperref[sec:II.10]{II.10}。因此，$\mathfrak R,\allowbreak \mathfrak S,\allowbreak \ldots$同时可测的充分必要条件，就是$R,\allowbreak S,\allowbreak \ldots$彼此对易。

现在不要求测量绝对准确，而只要求达到预先给定、但可以任意高的精度。这时$R,\allowbreak S,\allowbreak \ldots$不再一定要具有纯点谱。

对$\mathfrak R,\allowbreak \mathfrak S,\allowbreak \ldots$作有限精度的测量，等效于对$F(\mathfrak R),\allowbreak G(\mathfrak S),\allowbreak \ldots$作绝对准确的测量。这里的$F(\lambda),\allowbreak G(\lambda),\allowbreak \ldots$是本节开头讨论有限精度测量时所构造的那一类函数，当时当然只写出了$F(\lambda)$。所以，只要所有$F(\mathfrak R),\allowbreak G(\mathfrak S),\allowbreak \ldots$都能绝对准确地同时测量，$\mathfrak R,\allowbreak \mathfrak S,\allowbreak \ldots$就必定能同时测量。后一个条件又等价于$F(R),\allowbreak G(S),\allowbreak \ldots$彼此对易；而$R,\allowbreak S,\allowbreak \ldots$彼此对易就能推出这一点。因此，$R,\allowbreak S,\allowbreak \ldots$对易仍然是充分条件。

反过来，假定$\mathfrak R,\allowbreak \mathfrak S,\allowbreak \ldots$可以同时测量，可以作如下推论。只要测量$\mathfrak R$足够准确，就能判断其值大于$\bar\lambda$，还是小于等于$\bar\lambda$；参见上面关于“有限精度”的定义，尤其是原注 \hyperlink{note-126}{\getrefnumber{orig-note-126}} 的讨论。因此，令$F(\lambda)$在$\lambda\le\bar\lambda$时为$1$，在$\lambda>\bar\lambda$时为$0$，则$F(\mathfrak R)$可以绝对准确地测量。同样，令$G(\mu)$在$\mu\le\bar\mu$时为$1$，在$\mu>\bar\mu$时为$0$，则$G(\mathfrak S)$也能绝对准确地测量，而且这两个物理量还能同时测量。因此$F(R),\allowbreak G(S)$对易。\hypertarget{origpage-121}{}\phantomsection\label{origpage-121}设$E(\lambda),\allowbreak F(\mu)$为$R,\allowbreak S$对应的单位分解，则$F(R)=E(\bar\lambda)$、$G(S)=F(\bar\mu)$。所以，对一切$\bar\lambda,\allowbreak \bar\mu$，$E(\bar\lambda),\allowbreak F(\bar\mu)$都对易，因而$R,\allowbreak S$对易。这必须对$R,\allowbreak S,\allowbreak \ldots$中任意一对都成立，故所有$R,\allowbreak S,\allowbreak \ldots$彼此对易。因此，这个条件也是必要的。

可见，任意有限个物理量$\mathfrak R,\allowbreak \mathfrak S,\allowbreak \ldots$同时可测的特征条件，就是其算符$R,\allowbreak S,\allowbreak \ldots$彼此对易。无论要求绝对准确，还是任意准确的测量，这一结论都成立。不过，在前一种情形，还必须加上绝对准确测量本身所要求的特征条件：这些算符应具有纯点谱。

这样，我们已经从数学上证明，\hyperref[post:III.1.W]{W.} 给出的是这个理论——即 \hyperref[post:III.1.W]{W.} 所包含的理论——所能作出的最充分的断言。因为它以算符$R_1,\allowbreak \ldots,\allowbreak R_l$对易为前提；而在没有这一条件时，本来就不能讨论$\mathfrak R_1,\allowbreak \ldots,\allowbreak \mathfrak R_l$的同时测量结果，因为这些物理量根本不可能同时测量。

\section{不确定性关系}\label{sec:III.4}
前几节已经阐明了测量一个物理量，或同时测量几个可同时测量的物理量时的测量过程。现在需要了解：若关注的是同一系统、同一状态$\varphi$中的统计规律，那些不能同时测量的物理量会表现出怎样的行为。

给定这样的两个物理量$\mathfrak R,\allowbreak \mathfrak S$及其不对易的算符$R,\allowbreak S$。尽管作了这一假定，仍可能存在某些状态$\varphi$，使两个物理量都具有确定的值，即散布程度为$0$；这些状态就是两者的共同本征函数。只是它们不可能构成完全正交系，否则$R,\allowbreak S$便会对易。参见 \hyperref[sec:II.8]{II.8} 中相应单位分解$E(\lambda),\allowbreak F(\lambda)$的构造：若$\varphi_1,\allowbreak \varphi_2,\allowbreak \ldots$就是上述完全正交系，则$E(\lambda)$及$F(\lambda)$都由$P_{[\varphi_\varrho]}$相加而成；这些$P_{[\varphi_\varrho]}$彼此对易，所以两者也对易。因此，这些$\varphi$张成的闭线性流形$\mathfrak M$必小于$\mathfrak R_\infty$；如果它就是$\mathfrak R_\infty$，则按 \hyperref[sec:II.6]{II.6} 开头对单个算符所作的讨论，便能构造出所需的完全规范正交系。

对于$\mathfrak M$中的状态，$\mathfrak R,\allowbreak \mathfrak S$仍然可以同时测量。要证明这一点，最简单的办法是给出一个同时测量的模型。\hypertarget{origpage-122}{}\phantomsection\label{origpage-122}由于$R,\allowbreak S$的共同本征函数$\varphi$张成$\mathfrak M$，其中也存在一个由这些$\varphi$构成、同样张成$\mathfrak M$的规范正交系$\varphi_1,\allowbreak \varphi_2,\allowbreak \ldots$，即它在$\mathfrak M$中是完全的。这也可以借助前面提到的 \hyperref[sec:II.6]{II.6} 中的构造得到。加入一个张成$\mathfrak R_\infty-\mathfrak M$的规范正交系$\psi_1,\allowbreak \psi_2,\allowbreak \ldots$，就能把$\varphi_1,\allowbreak \varphi_2,\allowbreak \ldots$扩充为完全系$\varphi_1,\allowbreak \varphi_2,\allowbreak \ldots,\allowbreak \psi_1,\allowbreak \psi_2,\allowbreak \ldots$。再取互不相同的数$\lambda_1,\allowbreak \lambda_2,\allowbreak \ldots,\allowbreak \mu_1,\allowbreak \mu_2,\allowbreak \ldots$，并定义$T$为
\[
T\left(\sum_nx_n\varphi_n+\sum_ny_n\psi_n\right)
=\sum_n\lambda_nx_n\varphi_n+\sum_n\mu_ny_n\psi_n,
\]
记相应的物理量为$\mathfrak T$。

根据 \hyperref[sec:III.3]{III.3}，测量$\mathfrak T$会产生状态$\varphi_1,\allowbreak \varphi_2,\allowbreak \ldots,\allowbreak \psi_1,\allowbreak \psi_2,\allowbreak \ldots$中的一个。假如所得状态是$\varphi_n$，即测量结果是某个$\lambda_n$，那么$\mathfrak R$与$\mathfrak S$的值也就知道了：按照假定，$\mathfrak R,\allowbreak \mathfrak S$在$\varphi_n$中都有确定的值，所以能够确切预言，紧接着再测量$\mathfrak R$或$\mathfrak S$，将分别得到这些值。反之，若所得状态是$\psi_n$，就不能作这种判断；$\psi_n$不在$\mathfrak M$内，所以$\mathfrak R,\allowbreak \mathfrak S$在$\psi_n$中并不都具有确定的值。不过，我们知道，得到$\psi_n$的概率为$(P_{[\psi_n]}\varphi,\allowbreak \varphi)$；得到任意一个$\psi_n$（$n=1,\allowbreak 2,\allowbreak \ldots$）的概率为
\begin{align*}
\sum_n(P_{[\psi_n]}\varphi,\varphi)
&=(P_{\mathfrak R_\infty-\mathfrak M}\varphi,\varphi)
=\|P_{\mathfrak R_\infty-\mathfrak M}\varphi\|^2\\
&=\|\varphi-P_{\mathfrak M}\varphi\|^2.
\end{align*}
因此，当$\varphi=P_{\mathfrak M}\varphi$，即$\varphi$属于$\mathfrak M$时，这个概率为$0$，也就是说，这次测量必定同时测出了$\mathfrak R,\allowbreak \mathfrak S$的值\origNote[值][。]{129}。

不过，我们现在关心的是不能同时测量的物理量，所以假定最极端的情形$\mathfrak M=(0)$：$\mathfrak R,\allowbreak \mathfrak S$在任何状态中都不能同时测量；或者说，$R,\allowbreak S$没有任何共同本征函数。

设$R,\allowbreak S$的单位分解为$E(\lambda),\allowbreak F(\lambda)$。根据 \hyperref[sec:III.1]{III.1}，在状态$\varphi$中，$R,\allowbreak S$的期望值为
\[
\varrho=(R\varphi,\varphi),\qquad\sigma=(S\varphi,\varphi),
\]
其方差，即$(\mathfrak R-\varrho)^2$、$(\mathfrak S-\sigma)^2$的期望值为
\begin{align*}
\varepsilon^2&=((R-\varrho\cdot1)^2\varphi,\varphi)
=\|(R-\varrho\cdot1)\varphi\|^2=\|R\varphi-\varrho\varphi\|^2,\\
\eta^2&=((S-\sigma\cdot1)^2\varphi,\varphi)
=\|(S-\sigma\cdot1)\varphi\|^2=\|S\varphi-\sigma\varphi\|^2.
\end{align*}
参见 \hyperref[sec:III.3]{III.3} 关于绝对准确测量的讨论。作一个熟知的变形，得到\origNote[到][]{130}
\[
\varepsilon^2=\|R\varphi\|^2-(R\varphi,\varphi)^2,\qquad
\eta^2=\|S\varphi\|^2-(S\varphi,\varphi)^2.
\]
由于$\|\varphi\|=1$，Schwarz不等式，即 \hyperref[thm:II.1]{II.1的定理1}，已能说明左边的量非负。问题在于，$\varepsilon=0$与$\eta=0$不可能同时成立，$\varepsilon$和$\eta$却又各自都可以任意小。\hypertarget{origpage-123}{}\phantomsection\label{origpage-123}这是因为$\mathfrak R,\allowbreak \mathfrak S$分别都能任意准确地测量，甚至可能绝对准确地测量。因此，$\varepsilon,\allowbreak \eta$之间必须存在某种关系，使两者不能同时变小。这种关系究竟是什么？

Heisenberg发现了这类关系的存在\origNote[在][。]{131}。它们对于认识量子力学给自然界的描述带来的不确定性，具有重大的意义，因此称为不确定性关系。下面先从数学上推导其中最重要的一条，再讨论它的根本意义，以及它与经验的关系。

在矩阵理论中，满足对易关系
\[
PQ-QP=\frac h{2\pi i}1
\]
的算符$P,\allowbreak Q$起着重要作用。例如，它们对应坐标及其共轭动量（参见 \hyperref[sec:I.2]{I.2}）；更一般地说，对应经典力学中任意一对正则共轭量，参见原注 \hyperlink{note-2}{\getrefnumber{orig-note-2}} 所列的论文。现在更一般地考察满足
\[
PQ-QP=a\cdot1
\]
的两个厄米算符$P,\allowbreak Q$。因为$(PQ-QP)^*=QP-PQ$，所以$(a\cdot1)^*=\bar a\cdot1=-a\cdot1$，即$\bar a=-a$，故$a$是纯虚数。当然，这个算符等式并不意味着定义域也相同，因为$PQ-QP$未必在所有地方都有意义。对每个$\varphi$，有
\begin{align*}
2\operatorname{Im}(P\varphi,Q\varphi)
&=-i[(P\varphi,Q\varphi)-(Q\varphi,P\varphi)]\\
&=-i[(QP\varphi,\varphi)-(PQ\varphi,\varphi)]\\
&=(i\{PQ-QP\}\varphi,\varphi)=ia\cdot\|\varphi\|^2.
\end{align*}
设$a\ne0$，根据 \hyperref[thm:II.1]{II.1的定理1}，可得
\[
\|\varphi\|^2=-\frac{2i}a\operatorname{Im}(P\varphi,Q\varphi)
\le\frac2{|a|}|(P\varphi,Q\varphi)|
\le\frac2{|a|}\|P\varphi\|\cdot\|Q\varphi\|,
\]
所以，当$\|\varphi\|=1$时，
\[
\|P\varphi\|\cdot\|Q\varphi\|\ge\frac{|a|}2.
\]
$P-\varrho\cdot1$、$Q-\sigma\cdot1$也满足上述对易关系，因而还有
\[
\|P\varphi-\varrho\varphi\|\cdot\|Q\varphi-\sigma\varphi\|\ge\frac{|a|}2.
\]
引入平均值和方差：
\begin{alignat*}{2}
\varrho&=(P\varphi,\varphi),&\quad \varepsilon^2&=\|P\varphi-\varrho\varphi\|^2,\\
\sigma&=(Q\varphi,\varphi),&\quad \eta^2&=\|Q\varphi-\sigma\varphi\|^2,
\end{alignat*}
便得到
\begin{equation}
\varepsilon\eta\ge\frac{|a|}2.\tag{U}\label{eq:III.4.U}
\end{equation}\hypertarget{origpage-124}{}\phantomsection\label{origpage-124}等号成立的充分必要条件，是推导中所用的每一个$\le$都取等号。因此，记$P'=P-\varrho\cdot1$、$Q'=Q-\sigma\cdot1$，便有
\[
-\frac{i|a|}a\operatorname{Im}(P'\varphi,Q'\varphi)
=|(P'\varphi,Q'\varphi)|=\|P'\varphi\|\cdot\|Q'\varphi\|.
\]
最后一个等式根据 \hyperref[thm:II.1]{II.1的定理1}，意味着$P'\varphi,\allowbreak Q'\varphi$只相差一个常数因子；又因为$\|P'\varphi\|\cdot\|Q'\varphi\|\ge|a|/2>0$，所以$P'\varphi\ne0$、$Q'\varphi\ne0$，必须有$P'\varphi=cQ'\varphi$，$c'\ne0$。\tnote[。][第][原]{原书此处印作$c'\ne0$，但前面比例因子记作$c$。}第一个等式则说明，$(P'\varphi,\allowbreak Q'\varphi)=c\|Q'\varphi\|^2$为纯虚数，而且它的$i$的系数与实数$-i|a|/a$同号，即与$a$的$i$的系数异号。因此，$c=i\gamma$，$\gamma$为实数，当$a/i\lessgtr0$时，$\gamma\gtrless0$。于是
\begin{equation}
P'\varphi=i\gamma Q'\varphi,\qquad
\gamma\mathwordspace{}\text{为实数，且当}\mathwordspace{}ia\lessgtr0\mathwordspace{}\text{时}\mathwordspace{}\gamma\lessgtr0.
\tag{Gl.}\label{eq:III.4.Gl}
\end{equation}
$\varrho,\allowbreak \sigma$的定义还要求$(P'\varphi,\allowbreak \varphi)=0$、$(Q'\varphi,\allowbreak \varphi)=0$。但由\CJKglue \hyperref[eq:III.4.Gl]{（Gl.）}\CJKglue 得$(P'\varphi,\allowbreak \varphi)=i\gamma(Q'\varphi,\allowbreak \varphi)$，左边是实数，右边是纯虚数，所以这两个等式自然成立。再来确定$\varepsilon,\allowbreak \eta$。有
\[
\varepsilon:\eta=\|P'\varphi\|:\|Q'\varphi\|=|c|=|\gamma|,
\qquad\varepsilon\eta=\frac{|a|}2.
\]
由于$\varepsilon,\allowbreak \eta$都为正，故
\[
\varepsilon=\sqrt{\frac{|a|\cdot|\gamma|}2},\qquad
\eta=\sqrt{\frac{|a|}{2|\gamma|}}.
\]
在量子力学的情形，$a=h/(2\pi i)$，\CJKlinkpunctkern{，}{（}\hyperref[eq:III.4.U]{（U）}\CJKglue 变为
\begin{equation}
\varepsilon\eta\ge\frac h{4\pi}.\tag{$U'$}\label{eq:III.4.Uprime}
\end{equation}
还可以进一步讨论\CJKglue \hyperref[eq:III.4.Gl]{（Gl.）}\CJKlinkpunctkern{）}{。}。例如，设$P,\allowbreak Q$为Schrödinger理论中的算符，$P=\frac h{2\pi i}\frac\partial{\partial q}\cdots$、$Q=q\cdots$，参见 \hyperref[sec:I.2]{I.2}；这里假定力学系统只有一个自由度，唯一的坐标为$q$。这时\CJKglue \hyperref[eq:III.4.Gl]{（Gl.）}\CJKglue 为
\[
\left(\frac h{2\pi i}\frac\partial{\partial q}-\varrho\right)\varphi
=i\gamma(q-\sigma)\varphi;
\]
因为$ia=h/(2\pi)>0$，所以$\gamma>0$。于是
\[
\frac\partial{\partial q}\varphi
=\left\{-\frac{2\pi}h\gamma q+\frac{2\pi}h\gamma\sigma+\frac{2\pi}h\varrho i\right\}\varphi,
\]
即
\begin{align*}
\varphi&=e^{\int^q\{-\frac{2\pi}h\gamma q+\frac{2\pi}h\gamma\sigma+\frac{2\pi}h\varrho i\}\,dq}\\
&=Ce^{-\frac{\pi\gamma}hq^2+\frac{2\pi\gamma}h\sigma q+\frac{2\pi\varrho}hiq}
=C'e^{-\frac{\pi\gamma}h(q-\sigma)^2+\frac{2\pi\varrho}hiq}.
\end{align*}\hypertarget{origpage-125}{}\phantomsection\label{origpage-125}由于$\gamma>0$，$\|\varphi\|^2=\int_{-\infty}^{+\infty}|\varphi(q)|^2\,dq$确实有限；$C'$由$\|\varphi\|=1$确定：
\begin{align*}
\|\varphi\|^2&=\int_{-\infty}^{+\infty}|\varphi(q)|^2\,dq
=|C'|^2\int_{-\infty}^{+\infty}e^{-\frac{2\pi\gamma}h(q-\sigma)^2}\,dq\\
&=|C'|^2\sqrt{\frac h{2\pi\gamma}}\int_{-\infty}^{+\infty}e^{-x^2}\,dx\\
&=|C'|^2\sqrt{\frac h{2\pi\gamma}}\sqrt\pi
=|C'|^2\sqrt{\frac h{2\gamma}}=1,\\
|C'|&=\left(\frac{2\gamma}h\right)^{1/4}.
\end{align*}
所以，略去物理上无意义、绝对值为$1$的因子，有$C'=(2\gamma/h)^{1/4}$，即
\[
\varphi=\varphi(q)=\left(\frac{2\gamma}h\right)^{1/4}
e^{-\frac{\pi\gamma}h(q-\sigma)^2+\frac{2\pi\varrho}hiq}.
\]
$\varepsilon,\allowbreak \eta$已经知道：
\[
\varepsilon=\sqrt{\frac{h\gamma}{4\pi}},\qquad
\eta=\sqrt{\frac h{4\pi\gamma}}.
\]
当$\gamma$从$0$变到$+\infty$时，除了条件$\varepsilon\eta=h/(4\pi)$以外，它们可以任意取值。也就是说，每一组满足$\varepsilon\eta=h/(4\pi)$的四个数$\varrho,\allowbreak \sigma,\allowbreak \varepsilon,\allowbreak \eta$，都恰好由一个$\varphi$实现。Heisenberg最先研究了这些$\varphi$，并用它们阐明量子力学的情形。它们的确特别适合这一用途，因为它们表示了量子力学所允许的、对经典力学情形的最大限度的近似；在经典力学中，$p,\allowbreak q$都有确定的值。而且，按需要可以给$\varepsilon$或$\eta$指定任意值。参见原注 \hyperlink{note-131}{\getrefnumber{orig-note-131}} 所引之处。

不过，上面的讨论只揭示了不确定性关系的一面，即形式的一面。要完全理解这些关系，还必须从直接的物理经验出发，换一个角度加以考察。同量子力学据以建立的许多事实相比，不确定性关系同物理现象的联系更简单、更容易看清。因此，单凭上面纯形式的推导，还不足以充分阐明这些关系。尤其有必要作直观的讨论，因为乍看之下，甚至可能觉得它们与直观相矛盾：若有足够精细的测量仪器，为什么不能以任意高的精度，同时测量一个物体的位置和速度，即坐标与动量？\hypertarget{origpage-126}{}\phantomsection\label{origpage-126}这一点并非一望便知。所以，必须仔细分析最精细、甚至只能作为思想实验实现的测量过程，才能认清这种同时测量确实不可能。实际上，熟知的波动光学、电动力学以及原子基本过程的规律，恰恰会在不确定性关系所要求的地方，给准确测量设置不可逾越的困难。而且，只须以纯经典的方式——不用量子理论——研究这些过程，就已能认识到这一点。这在原则上非常重要：即使在不依赖量子力学正确性、早已得到确证的经典领域，也就是量子现象尚未迫使旧观点作实质修正的领域——严格说来，这也是唯一能让我们直接获得直观的领域——也不会与量子力学中那些听来似乎反常的不确定性关系发生冲突\origNote[突][。]{132}。

需要说明的是：若$p,\allowbreak q$为两个正则共轭量，系统所处的状态允许以精度$\varepsilon$给出$p$的值，也就是已经作了一次误差界限为$\varepsilon$的$p$测量，那么我们对$q$的了解就不可能比$\eta=\frac h{4\pi}:\varepsilon$更准确。或者说，以精度$\varepsilon$测量$p$，必定使$q$的值产生$\eta=\frac h{4\pi}:\varepsilon$的不确定性。

当然，这种着重定性的讨论不能重现每一个细节。对于尽可能精确的测量，我们只能给出$\varepsilon\eta\sim h$，即$\varepsilon\eta$与$h$同数量级，而不能得到$\varepsilon\eta=h/(4\pi)$。下面取粒子$\boldsymbol{T}$的位置〔坐标〕与动量这一对共轭量作为典型例子\origNote[子][。]{133}。

先研究位置的确定。为此，要看见$\boldsymbol{T}$，也就是照亮它，并让眼睛吸收散射光。从光源$\boldsymbol{l}$射出的一个光量子$\boldsymbol{L}$朝向$\boldsymbol{T}$，与$\boldsymbol{T}$碰撞后，从直线路径$\beta\beta_1$偏转到$\beta\beta_2$，最后在屏$Sch$上被吸收而消失；屏代替了眼睛或照相底片，见\CJKglue \hyperref[fig:1]{图1}\tnote[1][。][图]{图中的$Sch$是德文Schirm的缩写，意为“屏”；$\boldsymbol l$为光源，$\boldsymbol L$为光量子，$\boldsymbol T$为被观测粒子。}。测量就是确定$\boldsymbol{L}$击中屏$Sch$的位置：它不在自然路径$\beta\beta_1$末端的$1$处，而在$\beta\beta_2$末端的$2$处。不过，要据此确定碰撞的位置，也就是$\boldsymbol{T}$的位置，还必须知道$\beta$及$\beta_2$的方向，即$\boldsymbol{L}$在碰撞前后的方向。\hypertarget{origpage-127}{}\phantomsection\label{origpage-127}为此，可以插入适当的光阑组$ss$及$s's'$。（这样做实际上不是测量$\boldsymbol{T}$的坐标，而是判断它是否等于某一个值；这个值就是方向$\beta$与$\beta_2$的交点，适当安排光阑便可任意选定。把几个这样的判断结合起来，也就是设置几个$s's'$光阑，才相当于测量坐标。）这种位置测量能有多高的精度呢？

光学成像定律对此规定了一个根本的界限。用波长为$\lambda$的光，确实不能使小于$\lambda$的物体成清晰的像，甚至不能将模糊程度降低到足以形成图像的地步，哪怕只是失真的图像。不过，这里并不要求光学成像，因为仅凭$\boldsymbol{L}$发生偏转的事实，就足以确定$\boldsymbol{T}$的位置。然而，光阑$ss$及$s's'$的开口不能小于$\lambda$；否则$\boldsymbol{L}$就不能简单地穿过，而会形成一组干涉条纹。这时便无法从前后设置的光阑$ss$或$s's'$的方向，推知光线的方向$\beta$或$\beta_2$。所以，即使用$\boldsymbol{L}$作为射弹，也不能以高于$\lambda$的精度瞄准并击中目标。

\begin{figure}[tbp]
\centering
\includegraphics[width=.35\linewidth,keepaspectratio]{figures/figure-1.png}
\caption{}\label{fig:1}
\end{figure}

所以，波长$\lambda$也是坐标测量误差的尺度：$\lambda\sim\varepsilon$。$\boldsymbol{L}$的其他特征量是频率$\nu$、能量$\overline E$及动量$\bar p$。它们满足熟知的关系：
\[
\nu=\frac c\lambda,\qquad
\overline E=h\nu=\frac{hc}\lambda,\qquad
\bar p=\frac{\overline E}c=\frac{h\nu}c=\frac h\lambda,
\]
其中$c$为光速\origNote[速][。]{134}。因此$\bar p\sim h/\varepsilon$。在$\boldsymbol{L}$与$\boldsymbol{T}$之间那个情况不完全清楚的碰撞中，发生了一次动量传递，显然其数量级与$\bar p$，即$h/\varepsilon$相同。因此，$\boldsymbol{T}$产生了动量不确定性$\eta\sim h/\varepsilon$。

这样似乎已经证明了$\varepsilon\eta\sim h$，但还有一点被忽略了：碰撞过程并没有那么不清楚。\hypertarget{origpage-128}{}\phantomsection\label{origpage-128}我们知道$\boldsymbol{L}$在碰撞前后的运动方向$\beta$与$\beta_2$，因而也知道其动量；由此就能算出传递给$\boldsymbol{T}$的动量。因此，$\bar p$本身并不是$\eta$的尺度，真正的尺度应是光线方向$\beta,\allowbreak \beta_2$可能具有的不确定性。为了更准确地考察瞄准小物体$\boldsymbol{T}$与伴随的方向不确定性之间的联系，宜采用比狭缝$ss$更完善的瞄准装置，也就是透镜。这样就要用到熟知的显微镜理论：要照亮一个长度为$\varepsilon$的小面元，即以精度$\varepsilon$用$\boldsymbol{L}$击中$\boldsymbol{T}$，所需波长$\lambda$和透镜孔径角$\varphi$必须满足关系$\lambda/(2\sin(\varphi/2))\sim\varepsilon$，见\CJKglue \hyperref[fig:2]{图2}\origNote[2][。]{135}。\tnote[。][][图]{图中的$\boldsymbol l$、$\boldsymbol L$、$\boldsymbol T$分别表示光源、光量子、被观测粒子；$tt$标示横向方向，$\varphi$为孔径角。}

所以，$\boldsymbol{L}$的动量沿$tt$方向的分量之所以不确定，是因为它的方向只知道在$-\varphi/2$与$+\varphi/2$之间，而无法进一步确定。因此，这个分量的不确定性为
\[
2\sin\frac\varphi2\cdot\bar p=\frac\lambda\varepsilon\cdot\frac h\lambda=\frac h\varepsilon.
\]
这才是$\eta$的真正尺度。所以仍有$\eta\sim h/\varepsilon$，即$\varepsilon\eta\sim h$。

这个例子清楚地显示了不确定性原理的机制：要瞄准得准，就需要大眼睛，即孔径角$\varphi$大的透镜，以及波长相当短的光；这又意味着光量子的动量很大，而且很不确定。光量子与所观察的物体$\boldsymbol{T}$之间，便会发生无法完全掌握的碰撞，即Compton效应，从而使$\boldsymbol{T}$的动量“模糊”起来。

\begin{figure}[tbp]
\centering
\includegraphics[width=.35\linewidth,keepaspectratio]{figures/figure-2.png}
\caption{}\label{fig:2}
\end{figure}

现在再考察相反的测量过程：测量速度，即动量。首先要说明，测量$\boldsymbol{T}$速度的自然方法，是在两个不同的时刻，例如$0$和$t$，分别测定其位置，再把坐标的变化除以$t$。但这要求速度在时间区间$0,\allowbreak t$中恒定；如果速度改变了，这一改变量就成了刚才算出的平均速度与真实速度——例如时刻$t$的瞬时速度——之间偏差的尺度，也就是测量不确定性的尺度。动量同样如此。如果坐标测量的精度为$\varepsilon$，这虽然不妨碍平均动量的测量精度，因为$t$可以取得任意大，但会引起$\sim h/\varepsilon$的动量变化，\hypertarget{origpage-129}{}\phantomsection\label{origpage-129}因而使最终动量具有$\eta\sim h/\varepsilon$的不确定性。因此，无论如何仍得到$\varepsilon\eta\sim h$。若要有新的结果，只能采用不涉及位置测量的动量测量。这种方法确实存在，例如天文学中就经常采用，依据的是Doppler效应。下面来考察这一方法。

熟知的Doppler效应是：物体$\boldsymbol{T}$以速度$v$运动时，发出在物体本身上测得频率为$\nu_0$的光，静止观察者所见的频率则变为$\nu$，其值由关系
\[
\frac{\nu-\nu_0}{\nu_0}=\frac vc\cos\theta
\]
确定，$\theta$为运动方向与发光方向之间的夹角。这个公式是非相对论的，只对很小的$v/c$有效，不过很容易加以修正。因此，只要测出$\nu$并知道$\nu_0$，就能测定速度。例如，若所观察的光属于某种已知元素的一条特定谱线，$\nu_0$便是已知的。更确切地说，所测的是沿观察方向，即发光方向的速度分量：\tnote[：][][原]{原书下一式的分母印作$\nu$；由上一式解出应为$\nu_0$。这里保留原式。}
\[
v\cos\theta=\frac{c(\nu-\nu_0)}\nu,
\]
或者相应的动量分量$p'=p\cos\theta=mc(\nu-\nu_0)/\nu_0$，其中$m$为物体$\boldsymbol{T}$的质量。$p'$的散布程度$\eta$显然与频率$\nu$的散布宽度$\Delta\nu$相联系，所以$\eta\sim mc\Delta\nu/\nu_0\sim mc\Delta\nu/\nu$。当然，$\boldsymbol{T}$因发出频率为$\nu$、动量为$\bar p=h\nu/c$的光量子，其动量也会改变；但通常情况下，这个量的不确定性$h\Delta\nu/c$与$mc\Delta\nu/\nu$相比可以忽略\origNote[略][。]{136}。

频率$\nu$通过某种干涉现象测定。就其本性而言，只有纯单色波列的干涉，才能给出绝对确定的$\nu$值。这样的波列形式为
\[
a\sin\left(2\pi\left(\frac q\lambda-\nu t\right)+\alpha\right),
\]
其中$q$为坐标，$t$为时间，$a$为振幅，$\alpha$为相位；式子代表电场或磁场强度的任意一个分量。因此，它在空间和时间上都是无限延伸的。为了避免这一点，须将上式换成另一个函数$F(q/c-t)$，使它的自变量仅在某个有限区间内时，函数才不为$0$；因为$\lambda=c/\nu$，上式也可以写作$a\sin(2\pi\nu(q/c-t)+\alpha)$。若光波具有这样的形式，就必须作Fourier分析：
\[
F(x)=\int_0^{+\infty}a_\nu\sin(2\pi\nu x+\alpha_\nu)\,d\nu,
\]\hypertarget{origpage-130}{}\phantomsection\label{origpage-130}干涉图样会显示所有满足$a_\nu\ne0$的频率$\nu$；其中，频率区间$\nu,\allowbreak \nu+d\nu$的相对强度为$a_\nu^2\,d\nu$。频率$\nu$的散布宽度，即$\Delta\nu$，应从这一分布计算。

若波列在变量$x$上的延伸范围为$\tau$，也就是说，其持续时间为$\tau$、空间长度为$c\tau$，那么可以看出，频率$\nu$的散布宽度为$\sim1/\tau$。\origNote[。][另]{137}另一方面，这种测量会造成位置不确定性，因为每次发光时，$\boldsymbol{T}$都会受到沿观察方向、大小为$h\nu/c$的反冲，即速度改变$h\nu/(mc)$。然而，发射过程需要时间$\tau$，所以这次速度改变发生的时刻，无法比$\tau$更准确地定位；因而位置不确定性为$\varepsilon\sim(h\nu/(mc))\tau$。于是
\[
\varepsilon\sim\frac{h\nu}{mc}\tau,\qquad
\eta\sim\frac{mc\Delta\nu}\nu=\frac{mc}\nu\frac1\tau,\qquad
\varepsilon\eta\sim h.
\]
所以仍有$\varepsilon\eta\sim h$。

若$\boldsymbol{T}$不像这里假定的那样自己发光，而是散射外来的光，即被照亮，计算也完全类似。

\section{作为命题的投影算符}\label{sec:III.5}
像 \hyperref[sec:III.1]{III.1} 一样，考察一个具有$k$个自由度的物理系统$\boldsymbol{S}$；因此，其位形空间由$k$个坐标$q_1,\allowbreak \ldots,\allowbreak q_k$描述，参见 \hyperref[sec:I.2]{I.2}。按经典力学的看法，在系统$\boldsymbol{S}$中可以构造的一切物理量$\mathfrak R$，都是$q_1,\allowbreak \ldots,\allowbreak q_k$及其共轭动量$p_1,\allowbreak \ldots,\allowbreak p_k$的函数：
\[
\mathfrak R=\mathfrak R(q_1,\ldots,q_k,p_1,\ldots,p_k).
\]
例如，能量就是哈密顿函数$H(q_1,\allowbreak \ldots,\allowbreak q_k,\allowbreak p_1,\allowbreak \ldots,\allowbreak p_k)$。在量子力学中，则如 \hyperref[sec:III.1]{III.1} 已经指出的，每个物理量$\mathfrak R$都唯一地对应一个超极大厄米算符$R$；特别是，$q_1,\allowbreak \ldots,\allowbreak q_k$分别对应算符$Q_1=q_1\cdots,\allowbreak \ldots,\allowbreak Q_k=q_k\cdots$，$p_1,\allowbreak \ldots,\allowbreak p_k$分别对应算符
\[
P_1=\frac h{2\pi i}\frac\partial{\partial q_1}\cdots,\quad\ldots,\quad
P_k=\frac h{2\pi i}\frac\partial{\partial q_k}\cdots.
\]
由于$Q_l,\allowbreak P_l$不对易，一般不能定义$R=\mathfrak R(Q_1,\allowbreak \ldots,\allowbreak Q_k,\allowbreak P_1,\allowbreak \ldots,\allowbreak P_k)$。在 \hyperref[sec:I.2]{I.2} 中，我们已经以哈密顿函数为例说明过这一点。虽然还不能最终说明$\mathfrak R(q_1,\allowbreak \ldots,\allowbreak q_k,\allowbreak p_1,\allowbreak \ldots,\allowbreak p_k)$与$R$之间的联系，但在 \hyperref[sec:III.1]{III.1} 及 \hyperref[sec:III.3]{III.3} 中，已得到了以下几个特殊结论：

\leavevmode\phantomsection\label{post:III.5.L}\textbf{L.}若可同时观测的物理量$\mathfrak R,\allowbreak \mathfrak S$分别对应算符$R,\allowbreak S$，则物理量$a\mathfrak R+b\mathfrak S$（$a,\allowbreak b$为实数）对应算符$aR+bS$。

\hypertarget{origpage-131}{}\phantomsection\label{origpage-131}\leavevmode\phantomsection\label{post:III.5.F}\textbf{F.}若物理量$\mathfrak R$对应算符$R$，则物理量$F(\mathfrak R)$对应算符$F(R)$，其中$F(\lambda)$为任意实函数。

\hyperref[post:III.5.L]{L.} 及 \hyperref[post:III.5.F]{F.} 还可以作一定推广。其中 \hyperref[post:III.5.F]{F.} 的推广是必然的，表述如下：

\leavevmode\phantomsection\label{post:III.5.Fstar}\textbf{$F^*$.}若可同时观测的物理量$\mathfrak R,\allowbreak \mathfrak S,\allowbreak \ldots$分别对应算符$R,\allowbreak S,\allowbreak \ldots$，它们因而彼此对易，且个数有限，则物理量$F(\mathfrak R,\allowbreak \mathfrak S,\allowbreak \ldots)$对应算符$F(R,\allowbreak S,\allowbreak \ldots)$。

这里先假定$F(\lambda,\allowbreak \mu,\allowbreak \ldots)$为$\lambda,\allowbreak \mu,\allowbreak \ldots$的实多项式，以便明确$F(R,\allowbreak S,\allowbreak \ldots)$的意义，因为$R,\allowbreak S,\allowbreak \ldots$彼此对易。不过，对任意$F(\lambda,\allowbreak \mu,\allowbreak \ldots)$，也都可以证明 \hyperref[post:III.5.Fstar]{$F^*$}；关于一般的$F(R,\allowbreak S,\allowbreak \ldots)$的定义，参见原注 \hyperlink{note-94}{\getrefnumber{orig-note-94}} 所引之处。每个多项式都可由$a\lambda$、$\lambda+\mu$、$\lambda\mu$这三种运算反复构造出来，所以只须考察这三者。而$\lambda\mu=\frac14((\lambda+\mu)^2-(\lambda-\mu)^2)$，即等于
\[
\frac14\cdot(\lambda+\mu)^2+\left(-\frac14\right)\cdot(\lambda+(-1)\cdot\mu)^2,
\]
故也可以用$a\lambda$、$\lambda+\mu$、$\lambda^2$代替这三种运算。前两种包含在 \hyperref[post:III.5.L]{L.} 中，最后一种包含在 \hyperref[post:III.5.F]{F.} 中，因此 \hyperref[post:III.5.Fstar]{$F^*$} 得证。

但在量子力学中，\hyperref[post:III.5.L]{L.} 还要推广到$\mathfrak R,\allowbreak \mathfrak S$不能同时测量的情形。这个问题留待后面 \hyperref[sec:IV.1]{IV.1} 再讨论；目前只指出，对于不能同时测量的$\mathfrak R,\allowbreak \mathfrak S$，$a\mathfrak R+b\mathfrak S$究竟应当是什么意思，本来就完全不清楚。

除了物理量$\mathfrak R$，物理学还研究别的对象：系统$\boldsymbol{S}$的状态所具有的性质。例如，某个物理量$\mathfrak R$的值为$\lambda$；或$\mathfrak R$的值为正；或两个可同时观测的物理量$\mathfrak R,\allowbreak \mathfrak S$的值分别为$\lambda,\allowbreak \mu$；或这些值的平方和大于$1$；等等。我们用$\mathfrak R,\allowbreak \mathfrak S,\allowbreak \ldots$表示物理量，用$\mathfrak E,\allowbreak \mathfrak F,\allowbreak \ldots$表示性质。物理量对应超极大厄米算符$R,\allowbreak S,\allowbreak \ldots$，刚才已经讨论过；那么，性质又对应什么呢？

给每个性质$\mathfrak E$也对应一个物理量，定义如下：凡能判定$\mathfrak E$成立或不成立的测量，都看作对这个物理量的测量；若$\mathfrak E$成立，它的值就是$1$，反之就是$0$。我们把这个对应于$\mathfrak E$的物理量，也记作$\mathfrak E$。

它只能取$0$和$1$两个值。反过来，每个只能取这两个值的物理量$\mathfrak R$，也都对应一个性质$\mathfrak E$，显然就是“$\mathfrak R$的值不为$0$”。\hypertarget{origpage-132}{}\phantomsection\label{origpage-132}所以，与性质相对应的物理量$\mathfrak E$，恰可由上述特征来刻画。

$\mathfrak E$只能取$0,\allowbreak 1$，也可以这样表述：将它代入多项式$F(\lambda)=\lambda-\lambda^2$，结果恒为零。若$\mathfrak E$对应算符$E$，则$F(\mathfrak E)$对应算符$F(E)=E-E^2$；因此，条件就是$E-E^2=0$，即$E=E^2$。换言之，$\mathfrak E$的算符$E$为投影算符。

因此，各个性质$\mathfrak E$通过我们刚定义的相应物理量$\mathfrak E$，对应于投影算符$E$。或者，在$E=P_{\mathfrak M}$以外，再考察相应的闭线性流形$\mathfrak M$，也可以说，性质对应于闭线性流形$\mathfrak M$。

下面进一步研究相互对应的$\mathfrak E,\allowbreak E,\allowbreak \mathfrak M$的演算。

若要在状态$\varphi$中判定性质$\mathfrak E$是否成立，就必须测量物理量$\mathfrak E$，确定它的值是$1$还是$0$；按定义，这正是同一件事。因此，前一种情形，即$\mathfrak E$成立的概率，等于$\mathfrak E$的期望值，也就是
\[
(E\varphi,\varphi)=\|E\varphi\|^2=\|P_{\mathfrak M}\varphi\|^2,
\]
后一种情形，即$\mathfrak E$不成立的概率，等于$1-\mathfrak E$的期望值，也就是
\[
((1-E)\varphi,\varphi)=\|(1-E)\varphi\|^2=\|\varphi-P_{\mathfrak M}\varphi\|^2.
\]
两者之和当然等于$(\varphi,\allowbreak \varphi)$，即$1$。因此，当第二个概率或第一个概率为$0$，也就是$P_{\mathfrak M}\varphi=\varphi$或$P_{\mathfrak M}\varphi=0$时，$\mathfrak E$分别必定成立或必定不成立。换言之，$\varphi$属于$\mathfrak M$时，$\mathfrak E$必定成立；$\varphi$与$\mathfrak M$正交，即属于$\mathfrak R_\infty-\mathfrak M$时，$\mathfrak E$必定不成立。

也可以把$\mathfrak M$定义为所有使性质$\mathfrak E$确定成立的$\varphi$的集合。严格说，这只给出$\mathfrak M$中满足$\|\varphi\|=1$的部分；再将其中的元素乘以正的常数，并加入$0$，才得到整个$\mathfrak M$。

若把与$\mathfrak E$相反的性质，即对$\mathfrak E$的否定，称为“非$\mathfrak E$”，由上面立即可知：若$\mathfrak E$对应$E,\allowbreak \mathfrak M$，则“非$\mathfrak E$”对应$1-E,\allowbreak \mathfrak R_\infty-\mathfrak M$。

与物理量一样，性质也有同时可测性的问题；确切地说，是同时可判定性。显然，$\mathfrak E,\allowbreak \mathfrak F$同时可判定，当且仅当相应的物理量$\mathfrak E,\allowbreak \mathfrak F$可以同时测量。由于这些物理量只能取$0,\allowbreak 1$，测量是任意准确还是绝对准确，在这里无关紧要。因此，同时可判定的充分必要条件，就是$E,\allowbreak F$对易。对更多性质$\mathfrak E,\allowbreak \mathfrak F,\allowbreak \mathfrak G,\allowbreak \ldots$，结论相同。

\hypertarget{origpage-133}{}\phantomsection\label{origpage-133}对于同时可判定的性质$\mathfrak E,\allowbreak \mathfrak F$，可以构造新性质“$\mathfrak E$且$\mathfrak F$”及“$\mathfrak E$或$\mathfrak F$”。“$\mathfrak E$且$\mathfrak F$”对应的物理量，只有在$\mathfrak E$与$\mathfrak F$所对应的两个物理量都等于$1$时才为$1$，其中任何一个为$0$时就为$0$；所以它就是两者的乘积。因此，根据 \hyperref[post:III.5.Fstar]{$F^*$}，它的算符是$\mathfrak E$与$\mathfrak F$对应算符的乘积，即$EF$。由 \hyperref[thm:II.14]{II.4的定理14}，相应的闭线性流形是$\mathfrak M,\allowbreak \mathfrak N$的公共部分$\mathfrak P$。

至于“$\mathfrak E$或$\mathfrak F$”，可以写成
\[
\text{“非（（非}\mathwordspace{}\mathfrak E\text{）且（非}\mathwordspace{}\mathfrak F\text{））”},
\]
所以它的算符为
\[
1-(1-E)(1-F)=E+F-EF,
\]
按构造可知，这当然是投影算符。$F-EF$是投影算符，所以根据 \hyperref[thm:II.14]{II.4的定理14}，$E+F-EF$对应的线性流形是$\mathfrak M+(\mathfrak N-\mathfrak P)$。它属于$\{\mathfrak M,\allowbreak \mathfrak N\}$，而且显然包含$\mathfrak M$；由对称性，也包含$\mathfrak N$，从而包含$\{\mathfrak M,\allowbreak \mathfrak N\}$。所以它就等于$\{\mathfrak M,\allowbreak \mathfrak N\}$。因此，后者是闭的，便等于$[\mathfrak M,\allowbreak \mathfrak N]$。

若$\mathfrak E$是永远成立的性质，即恒真的性质，它对应的物理量便恒等于$1$，所以$E=1$、$\mathfrak M=\mathfrak R_\infty$。反之，若$\mathfrak E$永不成立，即不可能成立，其物理量就恒等于$0$，所以$E=0$、$\mathfrak M=(0)$。若两个性质$\mathfrak E,\allowbreak \mathfrak F$互不相容，它们首先必须能同时判定，并且“$\mathfrak E$且$\mathfrak F$”必须不可能成立；即$E,\allowbreak F$对易，且$EF=0$。但$EF=0$已经蕴含对易性，参见 \hyperref[thm:II.14]{II.4的定理14}，所以单是这个条件就是特征条件。预先假定$E,\allowbreak F$对易时，$EF=0$只表示$\mathfrak M,\allowbreak \mathfrak N$的公共部分仅含$0$；但单凭这一条件不能推出$E,\allowbreak F$对易。一般而言，对$\mathfrak M,\allowbreak \mathfrak N$，$EF=0$所表达的是：整个$\mathfrak M$与整个$\mathfrak N$正交，参见 \hyperref[thm:II.14]{II.4的定理14}。

若$\mathfrak R$为一个物理量，算符为$R$，其单位分解为$E(\lambda)$，则性质“$\mathfrak R$落在区间$I=\{\lambda',\allowbreak \mu'\}$内”，$\lambda'\le\mu'$，显然对应算符$E(\mu')-E(\lambda')$。例如，该性质的概率就等于$((E(\mu')-E(\lambda'))\varphi,\allowbreak \varphi)$，参见 \hyperref[sec:III.1]{III.1} 的 \hyperref[post:III.1.W]{W.}。这是因为，与此性质对应的物理量是$\mathfrak E=F(\mathfrak R)$，其中
\[
F(\lambda)=\begin{cases}1,&\lambda'<\lambda\le\mu',\\0,&\text{其他情形},\end{cases}
\]
并且$F(R)=E(\mu')-E(\lambda')$，参见 \hyperref[sec:II.8]{II.8} 及 \hyperref[sec:III.1]{III.1}。在 \hyperref[sec:III.1]{III.1} 中，我们也把这个算符记作$E(I)$。

综上，关于性质$\mathfrak E$、投影算符$E$与相应闭线性流形$\mathfrak M$之间的联系，我们已得到以下结论：

\hypertarget{origpage-134}{}\phantomsection\label{origpage-134}\leavevmode\phantomsection\label{post:III.5.alpha}\textbf{$\alpha)$}在状态$\varphi$中，性质$\mathfrak E$成立及不成立的概率分别为
\[
(E\varphi,\varphi)=\|E\varphi\|^2=\|P_{\mathfrak M}\varphi\|^2
\]
及
\[
((1-E)\varphi,\varphi)=\|(1-E)\varphi\|^2=\|\varphi-P_{\mathfrak M}\varphi\|^2.
\]

\leavevmode\phantomsection\label{post:III.5.beta}\textbf{$\beta)$}对于属于$\mathfrak M$或$\mathfrak R_\infty-\mathfrak M$的$\varphi$，$\mathfrak E$分别确定成立或确定不成立；并且只有这些$\varphi$才具有相应的确定性。

\leavevmode\phantomsection\label{post:III.5.gamma}\textbf{$\gamma)$}几个性质$\mathfrak E,\allowbreak \mathfrak F,\allowbreak \ldots$同时可判定的特征条件，是其算符$E,\allowbreak F,\allowbreak \ldots$彼此对易。

\leavevmode\phantomsection\label{post:III.5.delta}\textbf{$\delta)$}若$\mathfrak E$对应$E,\allowbreak \mathfrak M$，则“非$\mathfrak E$”对应$1-E,\allowbreak \mathfrak R_\infty-\mathfrak M$。

\leavevmode\phantomsection\label{post:III.5.epsilon}\textbf{$\varepsilon)$}若$\mathfrak E$对应$E,\allowbreak \mathfrak M$，$\mathfrak F$对应$F,\allowbreak \mathfrak N$，且$\mathfrak E,\allowbreak \mathfrak F$可同时判定，那么“$\mathfrak E$且$\mathfrak F$”对应$EF$及$\mathfrak M,\allowbreak \mathfrak N$的公共部分；“$\mathfrak E$或$\mathfrak F$”对应$E+F-EF$及$\{\mathfrak M,\allowbreak \mathfrak N\}$，后者等于$[\mathfrak M,\allowbreak \mathfrak N]$。

\leavevmode\phantomsection\label{post:III.5.eta}\textbf{$\eta)$}$\mathfrak E$永远成立，当$E=1$，也就是$\mathfrak M=\mathfrak R_\infty$；它永不成立，当$E=0$，也就是$\mathfrak M=(0)$。

\leavevmode\phantomsection\label{post:III.5.theta}\textbf{$\vartheta)$}$\mathfrak E,\allowbreak \mathfrak F$互不相容，当$EF=0$，也就是整个$\mathfrak M$与整个$\mathfrak N$正交。

\leavevmode\phantomsection\label{post:III.5.zeta}\textbf{$\zeta)$}设$\mathfrak R$为一个物理量，$R$为其算符，$I$为区间。又设$E(\lambda)$为$R$的单位分解，$I=\{\lambda',\allowbreak \mu'\}$，$\lambda'\le\mu'$，$E(I)=E(\mu')-E(\lambda')$，参见 \hyperref[sec:III.1]{III.1}。则性质“$\mathfrak R$落在$I$内”对应算符$E(I)$。

由 \hyperref[post:III.5.alpha]{$\alpha)$} 至 \hyperref[post:III.5.zeta]{$\zeta)$}，可以推出前面的概率断言 \hyperref[post:III.1.W]{W.}、\hyperref[post:III.1.E1]{$E_1$}、\hyperref[post:III.1.E2]{$E_2$}，以及 \hyperref[sec:III.3]{III.3} 中关于同时可测性的断言。后者显然与 \hyperref[post:III.5.gamma]{$\gamma)$} 等价；由 \hyperref[post:III.5.alpha]{$\alpha)$}、\hyperref[post:III.5.epsilon]{$\varepsilon)$}、\hyperref[post:III.5.zeta]{$\zeta)$} 可推出 \hyperref[post:III.1.W]{W.}，再由它推出 \hyperref[post:III.1.E1]{$E_1$}、\hyperref[post:III.1.E2]{$E_2$}。

可见，一个物理系统的性质与投影算符之间的关系，使我们能够对这些性质建立某种逻辑演算。只是与通常的逻辑演算相比，这里多出了量子力学特有的“同时可判定性”这一概念。

此外，以投影算符为基础的命题演算，比以全部〔超极大〕厄米算符为基础的物理量演算还有一个优点：“同时可判定性”细化了“同时可测性”的概念。例如，设$\mathfrak R,\allowbreak \mathfrak S$的算符分别为$R,\allowbreak S$，单位分解分别为$E(\lambda),\allowbreak F(\mu)$，$I=\{\lambda',\allowbreak \lambda''\}$、$J=\{\mu',\allowbreak \mu''\}$。要让“$\mathfrak R$落在$I$内”与“$\mathfrak S$落在$J$内”这两个问题同时可判定，根据 \hyperref[post:III.5.gamma]{$\gamma)$} 及 \hyperref[post:III.5.zeta]{$\zeta)$}，只须算符$E(I)=E(\lambda'')-E(\lambda')$与$F(J)=F(\mu'')-F(\mu')$对易；但要让$\mathfrak R,\allowbreak \mathfrak S$同时可测，则必须所有$E(\lambda)$都与所有$F(\mu)$对易。


% 中文编写中间材料：chapter-III-6.tex
\hypertarget{origpage-135}{}\phantomsection\label{origpage-135}\section{光的理论}\label{sec:III.6}
\hyperref[sec:I.2]{I.2} 中列出的量子力学统计断言，前面已经全部重新得到了，而且还作了实质性的推广与系统阐发。只有一项例外：我们还缺少Heisenberg关于量子化系统从一个定态跃迁到另一个定态的概率表达式；而它恰恰在量子力学的形成过程中起过重要作用，参见 \hyperref[sec:I.2]{I.2} 的讨论。下面将按照Dirac的方法\origNote[法][，]{138}，说明这些跃迁概率怎样从量子力学通常的统计断言——也就是刚阐明的理论——推导出来。这一点尤其重要，因为这样既能更深入地理解定态之间的跃迁机制，也能更清楚地看出Einstein–Bohr各条能量与频率关系之间的联系。下面要介绍的Dirac光的理论，确是量子力学领域最优美的成就之一。

设$\boldsymbol{S}$为一个系统，例如一个量子化的原子，其能量对应厄米算符$\mathsf{H}_0$。用$\xi$概括表示$\boldsymbol{S}$位形空间的各个坐标；例如，若$\boldsymbol{S}$由$l$个粒子组成，便有$3l$个直角坐标$x_1=q_1,\allowbreak y_1=q_2,\allowbreak z_1=q_3,\allowbreak \ldots,\allowbreak x_l=q_{3l-2},\allowbreak y_l=q_{3l-1},\allowbreak z_l=q_{3l}$。再为简单起见，假定$\mathsf{H}_0$具有纯点谱：本征值为$W_1,\allowbreak W_2,\allowbreak \ldots$，本征函数为$\varphi_1(\xi),\allowbreak \varphi_2(\xi),\allowbreak \ldots$；有些$W_n$允许相等。$\boldsymbol{S}$的任意状态，即波函数$\varphi(\xi)$，按含时Schrödinger微分方程演化（参见 \hyperref[sec:III.2]{III.2}）：
\[
 \frac{h}{2\pi i}\frac{\partial}{\partial t}\varphi_t(\xi)=-\mathsf{H}_0\varphi_t(\xi).
\]
也就是说，若$t=t_0$时$\varphi_t(\xi)=\varphi(\xi)=\sum_{k=1}^{\infty}a_k\varphi_k(\xi)$，则一般地有
\[
 \varphi_t(\xi)=\sum_{k=1}^{\infty}a_ke^{-\frac{2\pi i}{h}W_k(t-t_0)}\varphi_k(\xi).
\]
状态$\varphi(\xi)=\varphi_k(\xi)$变为$e^{-\frac{2\pi i}{h}W_k(t-t_0)}\varphi_k(\xi)$。其中的因子$e^{-\frac{2\pi i}{h}W_k(t-t_0)}$无关紧要，因而物理状态本身没有改变：$\varphi_k(\xi)$恰是定态。这样，最初根本看不到从某个$\varphi_k(\xi)$跃迁到另一个状态的过程——那么，为什么会谈论这种跃迁呢？解释很简单：我们忽略了引起这些跃迁的因素——光。早在最初的Bohr理论中，定态的量子轨道也只有在发光时才会衰变，参见所引原注 \hyperlink{note-5}{\getrefnumber{orig-note-5}}。因此，如果像刚才的设定那样忽略光，得到绝对而永久的稳定性，\hypertarget{origpage-136}{}\phantomsection\label{origpage-136}也就十分合乎情理了。因此，必须扩大所研究的系统，把$\boldsymbol{S}$可能发出的光，也就是一切能与$\boldsymbol{S}$发生相互作用的光，都纳入其中。将光构成的系统——即经典理论的电磁场，扣除由电子与原子核电荷产生的静电场——记为$\boldsymbol{L}$；我们要研究的是$\boldsymbol{S}+\boldsymbol{L}$。

因此，需要求出以下各项：

\leavevmode\phantomsection\label{problem:III.6.1}\textbf{1．}给出$\boldsymbol{L}$的量子力学描述，也就是给出$\boldsymbol{L}$的位形空间。

\leavevmode\phantomsection\label{problem:III.6.2}\textbf{2．}求出$\boldsymbol{S}+\boldsymbol{L}$的能量算符。它分成三部分：

\leavevmode\phantomsection\label{problem:III.6.2alpha}\textbf{$\alpha)$}不依赖$\boldsymbol{L}$而存在的$\boldsymbol{S}$的能量，即不受扰动的$\boldsymbol{S}$能量。对于系统$\boldsymbol{S}$，其算符为$\mathsf{H}_0$。

\leavevmode\phantomsection\label{problem:III.6.2beta}\textbf{$\beta)$}不依赖$\boldsymbol{S}$而存在的$\boldsymbol{L}$的能量，即不受扰动的$\boldsymbol{L}$能量。其算符记作$\mathsf{H}_l$。

\leavevmode\phantomsection\label{problem:III.6.2gamma}\textbf{$\gamma)$}其余能量，即$\boldsymbol{S}$与$\boldsymbol{L}$的相互作用能，其算符记作$\mathsf{H}_w$。

可见，按照量子力学的基本原则，首先必须用经典方法回答这些问题，再把所得结果改写为算符，参见 \hyperref[sec:I.2]{I.2}。所以，我们先完全从经典观点理解光的本性：按照光的电磁理论，把光看作电磁场的一种振动状态\origNote[态][。]{139}。

为了避免不必要的复杂性，例如光散逸到无限空间中，我们设想$\boldsymbol{S}$与$\boldsymbol{L}$封闭在一个很大的空腔$\boldsymbol H$内，其体积为$\mathcal V$，腔壁具有反射性。空腔$\boldsymbol H$中电磁场的状态，如所周知，由电场强度$\mathfrak E=\{\mathfrak E_x,\allowbreak \mathfrak E_y,\allowbreak \mathfrak E_z\}$及磁场强度$\mathfrak H=\{\mathfrak H_x,\allowbreak \mathfrak H_y,\allowbreak \mathfrak H_z\}$描述。所有物理量$\mathfrak E_x,\allowbreak \ldots,\allowbreak \mathfrak H_z$都是$\boldsymbol H$内任意一点的直角坐标$x,\allowbreak y,\allowbreak z$及时间$t$的函数。还要指出，下面多次要考察实的空间向量$\mathfrak a=\{\mathfrak a_x,\allowbreak \mathfrak a_y,\allowbreak \mathfrak a_z\}$、$\mathfrak b=\{\mathfrak b_x,\allowbreak \mathfrak b_y,\allowbreak \mathfrak b_z\}$，例如$\mathfrak E,\allowbreak \mathfrak H$，以及它们的内积等运算：
\[
 [\mathfrak a,\mathfrak b]=\mathfrak a_x\mathfrak b_x+\mathfrak a_y\mathfrak b_y+\mathfrak a_z\mathfrak b_z.
\]
这应当不会与$\mathfrak R_\infty$中的内积$(\varphi,\allowbreak \psi)$混淆。用$\Delta$表示微分算符$\frac{\partial^2}{\partial x^2}+\frac{\partial^2}{\partial y^2}+\frac{\partial^2}{\partial z^2}$，用$\operatorname{div},\allowbreak \operatorname{grad},\allowbreak \operatorname{rot}$表示熟知的向量微分运算。在空的腔体$\boldsymbol H$中，向量$\mathfrak E,\allowbreak \mathfrak H$满足Maxwell方程\hypertarget{origpage-137}{}\phantomsection\label{origpage-137}\begin{align*}
 \operatorname{div}\mathfrak H&=0,&
 \operatorname{rot}\mathfrak E+\frac1c\frac{\partial}{\partial t}\mathfrak H&=0,\\
 \operatorname{div}\mathfrak E&=0,&
 \operatorname{rot}\mathfrak H-\frac1c\frac{\partial}{\partial t}\mathfrak E&=0.
\end{align*}

令$\mathfrak H=\operatorname{rot}\mathfrak A$，便满足第一行的第一个方程。$\mathfrak A=\{\mathfrak A_x,\allowbreak \mathfrak A_y,\allowbreak \mathfrak A_z\}$称为矢势，其各分量同样依赖$x,\allowbreak y,\allowbreak z,\allowbreak t$。再令$\mathfrak E=-\frac1c\frac{\partial}{\partial t}\mathfrak A$，第一行的第二个方程也得到满足。这时，第二行的方程要求：
\begin{equation}
 \operatorname{div}\mathfrak A=0,\qquad
 \Delta\mathfrak A-\frac1{c^2}\frac{\partial^2}{\partial t^2}\mathfrak A=0.
 \tag{A}\label{eq:III.6.A}
\end{equation}
为了增强空间与时间之间的对称性，矢势通常以稍有不同的方式引入。这个关于$\mathfrak A$的设定能给出Maxwell方程的最一般解，在介绍Maxwell理论的大多数著作中都有证明，这里不再详述；其中特别须注意，第二行的第一个方程实际上只给出$\frac{\partial}{\partial t}\operatorname{div}\mathfrak A=0$，即$\operatorname{div}\mathfrak A=f(x,\allowbreak y,\allowbreak z)$。参见原注 \hyperlink{note-139}{\getrefnumber{orig-note-139}} 所引之处。

以下从\CJKglue \hyperref[eq:III.6.A]{（A）}\CJKglue 出发。腔壁具有反射性，这一事实表现为：$\mathfrak A$在$\boldsymbol H$的边界上必须垂直于腔壁。求出全部这样的$\mathfrak A$，可以使用熟知的方法。因为整个问题中并不显含$t$，最一般的$\mathfrak A$就是下列乘积形式解的线性组合：一个依赖$x,\allowbreak y,\allowbreak z$的向量，乘以一个依赖$t$的标量，
\[
 \mathfrak A=\mathfrak A(xyzt)=\overline{\mathfrak A}(xyz)\,\widetilde q(t).
\]
于是\CJKglue \hyperref[eq:III.6.A]{（A）}\CJKglue 给出
\begin{align}
 \operatorname{div}\overline{\mathfrak A}&=0,\quad
 \Delta\overline{\mathfrak A}=\eta\overline{\mathfrak A},\quad
 \overline{\mathfrak A}\mathwordspace{}\text{在}\mathwordspace{}\boldsymbol H\mathwordspace{}\text{的边界上垂直于边界},
 \tag{$A_1$}\label{eq:III.6.A1}\\
 \frac{\partial^2}{\partial t^2}\widetilde q(t)&=c^2\eta\widetilde q(t),
 \tag{$A_2$}\label{eq:III.6.A2}
\end{align}
根据\CJKglue \hyperref[eq:III.6.A1]{（$A_1$）}\CJKlinkpunctkern{）}{，}，$\eta$只依赖$x,\allowbreak y,\allowbreak z$；根据\CJKglue \hyperref[eq:III.6.A2]{（$A_2$）}\CJKlinkpunctkern{）}{，}，它又只依赖$t$，所以它为常数。

因此，\CJKlinkpunctkern{，}{（}\hyperref[eq:III.6.A1]{（$A_1$）}\CJKglue 是一个本征值问题，$\eta$为本征值参数，$\overline{\mathfrak A}$为一般本征函数。这个问题的理论已得到彻底解决，这里只给出结果\origNote[果][：]{140}：\CJKlinkpunctkern{：}{（}\hyperref[eq:III.6.A1]{（$A_1$）}\CJKglue 具有纯点谱，所有本征值$\eta_1,\allowbreak \eta_2,\allowbreak \ldots$都为负，当$n\longrightarrow+\infty$时趋于$-\infty$；相应的$\overline{\mathfrak A}$记为$\overline{\mathfrak A}_1,\allowbreak \overline{\mathfrak A}_2,\allowbreak \ldots$。可以选择一个完全系$\mathfrak A_1,\allowbreak \mathfrak A_2,\allowbreak \ldots$，并按下式归一化：
\[
 \iiint_{\boldsymbol H}[\overline{\mathfrak A}_m,\overline{\mathfrak A}_n]\,dx\,dy\,dz
 =\begin{cases}4\pi c^2,&m=n,\\0,&m\ne n.\end{cases}
\]\hypertarget{origpage-138}{}\phantomsection\label{origpage-138}我们选用$4\pi c^2$，而非通常采用的$1$，因为以后这样稍方便些。将负的$\eta_n$写作$-4\pi^2\varrho_n^2/c^2$，其中$\varrho_n>0$，当$n\longrightarrow+\infty$时也有$\varrho_n\longrightarrow+\infty$。\CJKlinkpunctkern{。}{（}\hyperref[eq:III.6.A2]{（$A_2$）}\CJKglue 于是给出
\[
 \widetilde q_n(t)=\gamma\cos2\pi\varrho_n(t-\tau)
 \qquad(\gamma,\tau\mathwordspace{}\text{任意}).
\]
因此，最一般的解$\mathfrak A$为：\tnote[：][][原]{原书下一展开式的每一项都印有两个$\gamma_n$因子，此处照录。}
\begin{align*}
 \mathfrak A=\mathfrak A(xyzt)
 &=\sum_{n=1}^{\infty}\overline{\mathfrak A}_n(xyz)\widetilde q_n(t)\\
 &=\sum_{n=1}^{\infty}\gamma_n\overline{\mathfrak A}_n(xyz)\gamma_n
 \cos2\pi\varrho_n(t-\tau_n)
\end{align*}
其中$\gamma_1,\allowbreak \gamma_2,\allowbreak \ldots,\allowbreak \tau_1,\allowbreak \tau_2,\allowbreak \ldots$为任意实常数。

任意场$\mathfrak A=\sum_{n=1}^{\infty}\overline{\mathfrak A}_n(xyz)\widetilde q_n(t)$的能量如下；这里不假定$\mathfrak A$为\CJKglue \hyperref[eq:III.6.A]{（A）}\CJKglue 的解，所以$\widetilde q_n(t)$可以任取：\tnote[：][][原]{原书能量计算的第二行，在时间导数的内积前漏了因子$1/c^2$，下一行却有该因子。这里保留这一原印差异。}
\begin{align*}
 E&=\frac1{8\pi}\iiint_{\boldsymbol H}
 \bigl([\mathfrak E,\mathfrak E]+[\mathfrak H,\mathfrak H]\bigr)\,dx\,dy\,dz\\
 &=\frac1{8\pi}\iiint_{\boldsymbol H}
 \left(\left[\frac{\partial\mathfrak A}{\partial t},\frac{\partial\mathfrak A}{\partial t}\right]
 +[\operatorname{rot}\mathfrak A,\operatorname{rot}\mathfrak A]\right)\,dx\,dy\,dz\\
 &=\frac1{8\pi}\sum_{m,n=1}^{\infty}\iiint_{\boldsymbol H}
 \left(\frac1{c^2}\frac{\partial\widetilde q_m(t)}{\partial t}
 \frac{\partial\widetilde q_n(t)}{\partial t}
 [\overline{\mathfrak A}_m,\overline{\mathfrak A}_n]\right.\\
 &\qquad\left.{}+\widetilde q_m(t)\widetilde q_n(t)
 [\operatorname{rot}\overline{\mathfrak A}_m,\operatorname{rot}\overline{\mathfrak A}_n]\right)\,dx\,dy\,dz.
\end{align*}
分部积分得到\origNote[到][]{141}
\begin{align*}
 &\iiint_{\boldsymbol H}[\operatorname{rot}\overline{\mathfrak A}_m,
 \operatorname{rot}\overline{\mathfrak A}_n]\,dx\,dy\,dz
 =\iiint_{\boldsymbol H}[\operatorname{rot}\operatorname{rot}\overline{\mathfrak A}_m,
 \overline{\mathfrak A}_n]\,dx\,dy\,dz\\
 &\quad=\iiint_{\boldsymbol H}[-\Delta\overline{\mathfrak A}_m
 +\operatorname{grad}\operatorname{div}\overline{\mathfrak A}_m,
 \overline{\mathfrak A}_n]\,dx\,dy\,dz\\
 &\quad=\frac{4\pi^2\varrho_m^2}{c^2}
 \iiint_{\boldsymbol H}[\overline{\mathfrak A}_m,\overline{\mathfrak A}_n]\,dx\,dy\,dz.
\end{align*}
所以
\begin{align*}
 E&=\frac1{8\pi}\sum_{m,n=1}^{\infty}
 \left(\frac1{c^2}\frac{\partial\widetilde q_m(t)}{\partial t}
 \frac{\partial\widetilde q_n(t)}{\partial t}
 +\frac{4\pi^2\varrho_m^2}{c^2}\widetilde q_m(t)\widetilde q_n(t)\right)\\
 &\qquad{}\times\iiint_{\boldsymbol H}
 [\overline{\mathfrak A}_m,\overline{\mathfrak A}_n]\,dx\,dy\,dz\\
 &=\frac12\sum_{m=1}^{\infty}
 \left[\left(\frac{\partial\widetilde q_m(t)}{\partial t}\right)^2
 +4\pi^2\varrho_m^2(\widetilde q_m(t))^2\right].
\end{align*}
现在可以把$\widetilde q_1,\allowbreak \widetilde q_2,\allowbreak \ldots$看作描述场的瞬时状态的坐标，也就是$\boldsymbol{L}$位形空间的坐标。\hypertarget{origpage-139}{}\phantomsection\label{origpage-139}由公式
\[
 E=\frac12\sum_{n=1}^{\infty}
 \left[\left(\frac{\partial\widetilde q_n}{\partial t}\right)^2
 +4\pi^2\varrho_n^2\widetilde q_n^2\right]
\]
按经典力学求得共轭动量$\widetilde p_n$：
\[
 \widetilde p_n=\frac{\partial E}{\partial(\partial\widetilde q_n/\partial t)}
 =\frac{\partial\widetilde q_n}{\partial t},\qquad
 E=\frac12\sum_{n=1}^{\infty}(\widetilde p_n^2+4\pi^2\varrho_n^2\widetilde q_n^2)
\]
参见 \hyperref[sec:I.2]{I.2}。再由经典力学得到运动方程
\begin{align*}
 \frac{\partial\widetilde p_n}{\partial t}
 &=-\frac{\partial E}{\partial\widetilde q_n}=-4\pi^2\varrho_n^2\widetilde q_n,\\
 \frac{\partial\widetilde q_n}{\partial t}
 &=\frac{\partial E}{\partial\widetilde p_n}=\widetilde p_n,
\end{align*}
即
\[
 \frac{\partial^2\widetilde q_n}{\partial t^2}+4\pi^2\varrho_n^2\widetilde q_n=0,
\]
这恰是由Maxwell方程得到的\CJKglue \hyperref[eq:III.6.A2]{（$A_2$）}\CJKlinkpunctkern{）}{。}。于是得到Jeans的定理：

辐射场$\boldsymbol{L}$可以完全按经典力学描述，取其坐标为$\widetilde q_1,\allowbreak \widetilde q_2,\allowbreak \ldots$，能量即哈密顿函数为
\[
 E=\tfrac12\sum_{n=1}^{\infty}(\widetilde p_n^2+4\pi^2\varrho_n^2\widetilde q_n^2).
\]
这些坐标与描述场在该时刻状态的矢势$\mathfrak A$的关系为
\[
 \mathfrak A=\mathfrak A(xyz)=\sum_{n=1}^{\infty}\widetilde q_n\overline{\mathfrak A}_n(xyz)
\]

一个被约束在直线上运动的质点，质量为$1$，坐标为$\widetilde q$，处于势能为$C\widetilde q^2$、$C=2\pi^2\varrho^2$的势场中，就具有能量$\frac12[(\frac{\partial\widetilde q}{\partial t})^2+4\pi^2\varrho^2\widetilde q^2]$；由于$\widetilde p=\frac{\partial\widetilde q}{\partial t}$，它也就是$\frac12(\widetilde p^2+4\pi^2\varrho^2\widetilde q^2)$。其运动方程为$\frac{\partial^2\widetilde q}{\partial t^2}+4\pi^2\varrho^2\widetilde q=0$，解为$\widetilde q=\gamma\cos2\pi\varrho(t-\tau)$，$\gamma,\allowbreak \tau$任意。按这种运动形式，这个系统称为“频率为$\varrho$的线性振子”。因此，$\boldsymbol{L}$可以看作一列线性振子的组合，它们的频率就是空腔$\boldsymbol H$的本征频率$\varrho_1,\allowbreak \varrho_2,\allowbreak \ldots$。

这种对电磁场的“力学”描述之所以重要，是因为可以立即按照通常方法，把它重新解释为量子力学描述。$\boldsymbol{L}$的位形空间由$\widetilde q_1,\allowbreak \widetilde q_2,\allowbreak \ldots$描述；在$E$的表达式中，须分别用$\frac{h}{2\pi i}\frac{\partial}{\partial\widetilde q_n}\cdots$及$\widetilde q_n\cdots$替换$\widetilde p_n,\allowbreak \widetilde q_n$，将这些算符分别记作$\widetilde P_n,\allowbreak \widetilde Q_n$。\hypertarget{origpage-140}{}\phantomsection\label{origpage-140}这样就回答了问题 \hyperref[problem:III.6.1]{1} 及 \hyperref[problem:III.6.2beta]{2的$\beta)$}；特别是
\[
 \mathsf{H}_l=\tfrac12\sum_{n=1}^{\infty}(\widetilde P_n^2+4\pi^2\varrho_n^2\widetilde Q_n^2)
\]
就是 \hyperref[problem:III.6.2beta]{2的$\beta)$} 所要求的算符。由于假定$\mathsf{H}_0$已知，\hyperref[problem:III.6.2alpha]{2的$\alpha)$} 从一开始就已解决。因此，只剩下 \hyperref[problem:III.6.2gamma]{2的$\gamma)$}，它也不再有困难。

按经典电动力学，$\boldsymbol{S}$与$\boldsymbol{L}$的相互作用这样计算。设$\boldsymbol{S}$由$l$个粒子，例如质子和电子组成，电荷和质量分别为$e_1,\allowbreak m_1,\allowbreak \ldots,\allowbreak e_l,\allowbreak m_l$；直角坐标为$x_1=q_1,\allowbreak y_1=q_2,\allowbreak z_1=q_3,\allowbreak \ldots,\allowbreak x_l=q_{3l-2},\allowbreak y_l=q_{3l-1},\allowbreak z_l=q_{3l}$，它们就是前面用$\xi$概括表示的各个坐标；相应的动量为$p_1^x,\allowbreak p_1^y,\allowbreak p_1^z,\allowbreak \ldots,\allowbreak p_l^x,\allowbreak p_l^y,\allowbreak p_l^z$。那么，在足够好的近似下，相互作用能为\origNote[为][]{142}
\begin{multline*}
 \sum_{\nu=1}^{l}\frac{e_\nu}{cm_\nu}
 \bigl(p_\nu^x\mathfrak A_x(x_\nu,y_\nu,z_\nu)
 +p_\nu^y\mathfrak A_y(x_\nu,y_\nu,z_\nu)\\
 {}+p_\nu^z\mathfrak A_z(x_\nu,y_\nu,z_\nu)\bigr).
\end{multline*}
将$p_\nu^x,\allowbreak p_\nu^y,\allowbreak p_\nu^z,\allowbreak x_\nu,\allowbreak y_\nu,\allowbreak z_\nu$，$\nu=1,\allowbreak \ldots,\allowbreak l$，替换为算符
\[
 \frac{h}{2\pi i}\frac{\partial}{\partial x_\nu}\cdots,\quad
 \frac{h}{2\pi i}\frac{\partial}{\partial y_\nu}\cdots,\quad
 \frac{h}{2\pi i}\frac{\partial}{\partial z_\nu}\cdots,\quad
 x_\nu\cdots,\quad y_\nu\cdots,\quad z_\nu\cdots,
\]
就得到相应的量子力学算符；把这些算符分别记作$P_\nu^x,\allowbreak P_\nu^y,\allowbreak P_\nu^z,\allowbreak Q_\nu^x,\allowbreak Q_\nu^y,\allowbreak Q_\nu^z$。再注意到\tnote[到][][原]{原书下一展开式的系数印作$\varrho_n$；前面的场展开式用的是$\widetilde q_n$。这里保留原式。}
\[
 \mathfrak A(xyz)=\sum_{n=1}^{\infty}\varrho_n\overline{\mathfrak A}_n(xyz),
\]
就得到所求的$\mathsf{H}_w$：
\begin{align*}
 \mathsf{H}_w=\sum_{n=1}^{\infty}\sum_{\nu=1}^{l}\frac{e_\nu}{cm_\nu}\widetilde Q_n
 \bigl\{&P_\nu^x\overline{\mathfrak A}_{n,x}(Q_\nu^x,Q_\nu^y,Q_\nu^z)\\
 &{}+P_\nu^y\overline{\mathfrak A}_{n,y}(Q_\nu^x,Q_\nu^y,Q_\nu^z)
 +P_\nu^z\overline{\mathfrak A}_{n,z}(Q_\nu^x,Q_\nu^y,Q_\nu^z)\bigr\}.
\end{align*}
不过，这里必须注意：将乘积$p_\nu^x\overline{\mathfrak A}_{n,\allowbreak x}(x_\nu,\allowbreak y_\nu,\allowbreak z_\nu),\allowbreak \ldots$替换为算符时，我们任意选定了次序$P_\nu^x\overline{\mathfrak A}_{n,\allowbreak x}(Q_\nu^x,\allowbreak Q_\nu^y,\allowbreak Q_\nu^z),\allowbreak \ldots$。本来也完全可以采用$\overline{\mathfrak A}_{n,\allowbreak x}(Q_\nu^x,\allowbreak Q_\nu^y,\allowbreak Q_\nu^z)P_\nu^x,\allowbreak \ldots$；而且，为了保证所得算符为厄米算符，严格说本应取对称形式
\[
 \tfrac12\bigl(P_\nu^x\overline{\mathfrak A}_{n,x}(Q_\nu^x,Q_\nu^y,Q_\nu^z)
 +\overline{\mathfrak A}_{n,x}(Q_\nu^x,Q_\nu^y,Q_\nu^z)P_\nu^x\bigr),\ldots
\]
幸运的是，这些次序并不造成差别，因为\origNote[为][]{143}
\begin{align*}
 &[P_\nu^x\overline{\mathfrak A}_{n,x}(Q_\nu^x,Q_\nu^y,Q_\nu^z)+\cdots]
 -[\overline{\mathfrak A}_{n,x}(Q_\nu^x,Q_\nu^y,Q_\nu^z)P_\nu^x+\cdots]\\
 &\quad=[P_\nu^x\overline{\mathfrak A}_{n,x}(Q_\nu^x,Q_\nu^y,Q_\nu^z)
 -\overline{\mathfrak A}_{n,x}(Q_\nu^x,Q_\nu^y,Q_\nu^z)P_\nu^x]+\cdots\\
 &\quad=\frac{h}{2\pi i}\frac{\partial}{\partial x}
 \overline{\mathfrak A}_{n,x}(Q_\nu^x,Q_\nu^y,Q_\nu^z)+\cdots\\
 &\quad=\frac{h}{2\pi i}\operatorname{div}\overline{\mathfrak A}_n(Q_\nu^x,Q_\nu^y,Q_\nu^z)=0.
\end{align*}\hypertarget{origpage-141}{}\phantomsection\label{origpage-141}这样就准确给出了系统$\boldsymbol{S}+\boldsymbol{L}$的总能量，或者说它的算符
\[
 \mathsf{H}=\mathsf{H}_0+\mathsf{H}_l+\mathsf{H}_w,
\]
不过，在继续改写$\mathsf{H}$以前，还须作如下说明。$\boldsymbol{S}+\boldsymbol{L}$的位形空间由坐标$\xi$——即$q_1,\allowbreak \ldots,\allowbreak q_{3l}$，或者$x_1,\allowbreak y_1,\allowbreak z_1,\allowbreak \ldots,\allowbreak x_l,\allowbreak y_l,\allowbreak z_l$——以及$\widetilde q_1,\allowbreak \widetilde q_2,\allowbreak \ldots$描述，波函数也就依赖这些坐标。从形式上说，允许无限多自由度的系统，或者允许具有无限多个自变量的波函数，既不方便，也不妥当；因为前面的各项规则始终只涉及有限多个坐标。因此，先只考虑$\widetilde q_1,\allowbreak \widetilde q_2,\allowbreak \ldots$中的前$N$个$\widetilde q_1,\allowbreak \ldots,\allowbreak \widetilde q_N$，即将$\mathfrak A$限制为$\overline{\mathfrak A}_1,\allowbreak \ldots,\allowbreak \overline{\mathfrak A}_N$的线性组合，等到合适的时候再作实质上必需的极限过程$N\longrightarrow+\infty$。于是
\begin{align*}
 \mathsf{H}={}&\mathsf{H}_0+\tfrac12\sum_{n=1}^{N}
 (\widetilde P_n^2+4\pi^2\varrho_n^2\widetilde Q_n^2)\\
 &{}+\sum_{n=1}^{N}\sum_{\nu=1}^{l}\frac{e_\nu}{cm_\nu}\widetilde Q_n
 \bigl\{P_\nu^x\overline{\mathfrak A}_{n,x}(Q_\nu^x,Q_\nu^y,Q_\nu^z)\\
 &\qquad{}+P_\nu^y\overline{\mathfrak A}_{n,y}(Q_\nu^x,Q_\nu^y,Q_\nu^z)
 +P_\nu^z\overline{\mathfrak A}_{n,z}(Q_\nu^x,Q_\nu^y,Q_\nu^z)\bigr\}.
\end{align*}
用非厄米算符$\widetilde R_n$及其伴随算符$\widetilde R_n^*$代替$\widetilde P_n,\allowbreak \widetilde Q_n$，较为方便：
\begin{align*}
 \widetilde R_n&=\frac1{\sqrt{2h\varrho_n}}(2\pi\varrho_n\widetilde Q_n+i\widetilde P_n),\\
 \widetilde R_n^*&=\frac1{\sqrt{2h\varrho_n}}(2\pi\varrho_n\widetilde Q_n-i\widetilde P_n).
\end{align*}
那么$\widetilde Q_n=\frac1{2\pi}\sqrt{\frac{h}{2\varrho_n}}(\widetilde R_n+\widetilde R_n^*)$；又由$\widetilde P_n\widetilde Q_n-\widetilde Q_n\widetilde P_n=\frac{h}{2\pi i}1$可得
\begin{align*}
 \widetilde R_n\widetilde R_n^*
 &=\frac1{2h\varrho_n}(\widetilde P_n^2+4\pi^2\varrho_n^2\widetilde Q_n^2)+\tfrac12\,1,\\
 \widetilde R_n^*\widetilde R_n
 &=\frac1{2h\varrho_n}(\widetilde P_n^2+4\pi^2\varrho_n^2\widetilde Q_n^2)-\tfrac12\,1,
\end{align*}
特别地，$\widetilde R_n\widetilde R_n^*-\widetilde R_n^*\widetilde R_n=1$。能量公式于是为\tnote[为][][原]{原书相互作用项的平方根内印作$\varrho_\nu$；由前面的代换应得到$\varrho_n$。这里保留原下标。}
\begin{align*}
 \mathsf{H}={}&\mathsf{H}_0+\sum_{n=1}^{N}h\varrho_n\widetilde R_n^*\widetilde R_n\\
 &{}+\sum_{n=1}^{N}\sum_{\nu=1}^{l}\frac{e_\nu}{2\pi cm_\nu}
 \sqrt{\frac{h}{2\varrho_\nu}}(\widetilde R_n+\widetilde R_n^*)\\
 &\quad{}\times\bigl[P_\nu^x\overline{\mathfrak A}_{n,x}(Q_\nu^x,Q_\nu^y,Q_\nu^z)
 +P_\nu^y\overline{\mathfrak A}_{n,y}(Q_\nu^x,Q_\nu^y,Q_\nu^z)\\
 &\qquad{}+P_\nu^z\overline{\mathfrak A}_{n,z}(Q_\nu^x,Q_\nu^y,Q_\nu^z)\bigr]+C,
\end{align*}
其中$C=\text{常数}=\frac12\sum_{n=1}^{N}h\varrho_n$。能量表达式中的加法常数没有物理意义，所以可以略去$C$。况且，当$N\longrightarrow+\infty$时，$C$也会趋于无穷大，妨碍理论得到合理的完整表述，更应将它略去。

\hypertarget{origpage-142}{}\phantomsection\label{origpage-142}厄米算符$\widetilde R_n^*\widetilde R_n$是超极大的，并且具有由数$0,\allowbreak 1,\allowbreak 2,\allowbreak \ldots$构成的纯点谱。相应的本征函数记作$\psi_0^n(\widetilde q_n),\allowbreak \psi_1^n(\widetilde q_n),\allowbreak \psi_2^n(\widetilde q_n),\allowbreak \ldots$。

若用$\frac1{2\pi}\sqrt{h/\varrho_n}\,q$替代变量$\widetilde q_n$，则
\begin{gather*}
 \frac1{\sqrt{2h\varrho_n}}2\pi\varrho_n\widetilde q_n
 =2\pi\sqrt{\frac{\varrho_n}{2h}}\,\widetilde q_n,\\
 \frac1{\sqrt{2h\varrho_n}}\frac h{2\pi i}\frac\partial{\partial\widetilde q_n}
 =\frac1{2\pi}\sqrt{\frac h{2\varrho_n}}\frac1i\frac\partial{\partial\widetilde q_n}
\end{gather*}
分别变成$q/\sqrt2$和$\frac1{\sqrt2}\frac1i\frac\partial{\partial q}$，从而
\begin{align*}
 \widetilde R_n&=\frac1{\sqrt2}\left(q+\frac\partial{\partial q}\right),&
 \widetilde R_n^*&=\frac1{\sqrt2}\left(q-\frac\partial{\partial q}\right),\\
 \widetilde R_n\widetilde R_n^*&=-\frac12\frac{\partial^2}{\partial q^2}+\frac12q^2+\frac12,\\
 \widetilde R_n^*\widetilde R_n&=-\frac12\frac{\partial^2}{\partial q^2}+\frac12q^2-\frac12.
\end{align*}
读者可以在许多著作中找到这些算符的本征值理论。例如，Courant–Hilbert：第261页公式（42）、（43）及其后内容，以及第76页公式（60）、（61）；或Weyl：\emph{Gruppentheorie und Quantenmechanik}，第74页及以下。

由于$\psi_1(\xi),\allowbreak \psi_2(\xi),\allowbreak \ldots$构成$\xi$空间中的完全正交系，而$\psi_0^n(\widetilde q_n),\allowbreak \psi_1^n(\widetilde q_n),\allowbreak \ldots$构成$\widetilde q_n$空间中的完全正交系，因此函数
\[
 \Phi_{kM_1\cdots M_N}(\xi\widetilde q_1\cdots\widetilde q_N)
 =\psi_k(\xi)\psi_{M_1}^1(\widetilde q_1)\cdots\psi_{M_N}^N(\widetilde q_N),
\]
其中$k=1,\allowbreak 2,\allowbreak \ldots$，$M_1,\allowbreak \ldots,\allowbreak M_N=0,\allowbreak 1,\allowbreak 2,\allowbreak \ldots$，在由$\xi,\allowbreak \widetilde q_1,\allowbreak \ldots,\allowbreak \widetilde q_N$描述的位形空间上构成完全正交系。所以每个波函数$\Phi=\Phi(\xi,\allowbreak \widetilde q_1,\allowbreak \ldots,\allowbreak \widetilde q_N)$都能展开为
\begin{align*}
 &\varphi(\xi,\widetilde q_1,\ldots,\widetilde q_N)\\
 &\quad=\sum_{k=1}^{\infty}\sum_{M_1=0}^{\infty}\cdots\sum_{M_N=0}^{\infty}
 a_{kM_1\cdots M_N}\Phi_{kM_1\cdots M_N}(\xi\widetilde q_1\cdots\widetilde q_N)\\
 &\quad=\sum_{k=1}^{\infty}\sum_{M_1=0}^{\infty}\cdots\sum_{M_N=0}^{\infty}
 a_{kM_1\cdots M_N}\psi_k(\xi)\psi_{M_1}^1(\widetilde q_1)\cdots\psi_{M_N}^N(\widetilde q_N).
\end{align*}
这个完全正交系和展开系数用了$N+1$个指标$k,\allowbreak M_1,\allowbreak \ldots,\allowbreak M_N$，而不只是一个指标，这一点并无实质影响。根据 \hyperref[sec:II.2]{II.2}，可以作出如下判断：波函数$\Phi$的希尔伯特空间，也可以看成由$N+1$重序列$a_{kM_1\cdots M_N}$构成的空间，其中
\[
 \sum_{k=1}^{\infty}\sum_{M_1=0}^{\infty}\cdots\sum_{M_N=0}^{\infty}|a_{kM_1\cdots M_N}|^2
\]
为有限值。在希尔伯特空间的这种表示下，算符$\mathsf{H}$又具有什么形式呢？

为求出这一点，先计算$\mathsf{H}\Phi_{kM_1\cdots M_N}$。$\mathsf{H}_0$只作用于$\xi$，而$\psi_k(\xi)$是$\mathsf{H}_0$的本征函数，本征值为$W_k$；另外，$\widetilde R_n^*\widetilde R_n$只作用于$\widetilde q_n$，$\psi_{M_n}^n(\widetilde q_n)$是$\widetilde R_n^*\widetilde R_n$的本征函数，本征值为$M_n$。\hypertarget{origpage-143}{}\phantomsection\label{origpage-143}因此
\begin{align*}
 \mathsf{H}\Phi_{kM_1\cdots M_N}
 ={}&\left(W_k+\sum_{n=1}^{N}h\varrho_n M_n\right)\Phi_{kM_1\cdots M_N}\\
 &{}+\sum_{n=1}^{N}\sum_{\nu=1}^{l}\frac{e_\nu}{2\pi cm_\nu}\sqrt{\frac h{2\varrho_n}}\\
 &\quad{}\times\bigl[P_\nu^x\overline{\mathfrak A}_{n,x}(Q_\nu^x,Q_\nu^y,Q_\nu^z)
 +P_\nu^y\overline{\mathfrak A}_{n,y}(Q_\nu^x,Q_\nu^y,Q_\nu^z)\\
 &\qquad{}+P_\nu^z\overline{\mathfrak A}_{n,z}(Q_\nu^x,Q_\nu^y,Q_\nu^z)\bigr]\psi_k(\xi)\\
 &\quad{}\times\psi_{M_1}^1(\widetilde q_1)\cdots
 (\widetilde R_n+\widetilde R_n^*)\psi_{M_n}^n(\widetilde q_n)\cdots\psi_{M_N}^N(\widetilde q_N).
\end{align*}
对于像方括号中那样、只影响变量$\xi$的算符$A$，可以使用展开式
\[
 A\psi_k(\xi)=\sum_{j=1}^{\infty}(A\psi_k,\psi_j)\psi_j(\xi)
 =\sum_j(A)_{kj}\psi_j(\xi),
\]
这里定义$(A)_{kj}=(A\psi_k,\allowbreak \psi_j)$。此外，如上面所引之处证明的，
\[
 \widetilde R_n\psi_M^n(\widetilde q_n)=\sqrt M\,\psi_{M-1}^n(\widetilde q_n),\qquad
 \widetilde R_n^*\psi_M^n(\widetilde q_n)=\sqrt{M+1}\,\psi_{M+1}^n(\widetilde q_n)
\]
当$M=0$时，第一个公式的右边应置为$0$，不必理会其中出现的、无意义的$\psi_{-1}^n$。于是\tnote[是][][原]{原书下式对$j$求和，但两个$\Phi$的首个下标仍印为$k$。这里按原印指标转录。}
\begin{align*}
 &\mathsf{H}\Phi_{kM_1\cdots M_N}\\
 &\quad=\left(W_k+\sum_{n=1}^{N}h\varrho_n M_n\right)\Phi_{kM_1\cdots M_N}\\
 &\quad{}+\sum_{j=1}^{\infty}\sum_{n=1}^{N}\sqrt{\frac h{2\varrho_n}}
 \left(\sum_{\nu=1}^{l}\frac{e_\nu}{2\pi cm_\nu}
 \bigl(P_\nu^x\overline{\mathfrak A}_{n,x}(Q_\nu^x,Q_\nu^y,Q_\nu^z)+\cdots\bigr)_{kj}\right)\\
 &\quad{}\times\bigl(\sqrt{M_n+1}\,\Phi_{kM_1\cdots M_n+1\cdots M_N}\\
 &\qquad{}+\sqrt{M_n}\,\Phi_{kM_1\cdots M_n-1\cdots M_N}\bigr).
\end{align*}
现在可以确定$\mathsf{H}$作用于$a_{kM_1\cdots M_N}$时的形式：
\[
 \mathsf{H}\sum_{kM_1\cdots M_N}a_{kM_1\cdots M_N}\Phi_{kM_1\cdots M_N}
 =\sum_{kM_1\cdots M_N}a'_{kM_1\cdots M_N}\Phi_{kM_1\cdots M_N},
\]
其中
\begin{align*}
 &\mathsf{H}a_{kM_1\cdots M_N}=a'_{kM_1\cdots M_N}\\
 &\quad={}\left(W_k+\sum_{n=1}^{N}h\varrho_n M_n\right)a_{kM_1\cdots M_N}\\
 &{}+\sum_{j=1}^{\infty}\sum_{n=1}^{N}\sqrt{\frac h{2\varrho_n}}
 \left(\sum_{\nu=1}^{l}\frac{e_\nu}{2\pi cm_\nu}
 \bigl(P_\nu^x\overline{\mathfrak A}_{n,x}(Q_\nu^x,Q_\nu^y,Q_\nu^z)+\cdots\bigr)_{kj}\right)\\
 &\quad{}\times\left(\sqrt{M_n}\,a_{jM_1\cdots M_n-1\cdots M_N}
 +\sqrt{M_n+1}\,a_{jM_1\cdots M_n+1\cdots M_N}\right).
\end{align*}
对$\mathsf{H}$的讨论进行到这里，便可以取极限$N\longrightarrow\infty$了。由于$a_{kM_1\cdots M_N}$的指标结构随之改变，这将产生一种全新的$\mathsf{H}$算符。必须引入具有无限多个指标$M_1,\allowbreak M_2,\allowbreak \ldots$的分量$a_{kM_1M_2\cdots}$；但仅仅为了保证$\mathsf{H}$中出现的和$\sum_{n=1}^{\infty}h\varrho_n M_n$有限，\hypertarget{origpage-144}{}\phantomsection\label{origpage-144}就必须把$M_1,\allowbreak M_2,\allowbreak \ldots$限制为仅有限项不等于$0$的序列。由此，希尔伯特空间由所有满足$\sum_{k=1}^{\infty}\sum_{M_1=0}^{\infty}\sum_{M_2=0}^{\infty}\cdots|a_{kM_1M_2\cdots}|^2$有限的序列$a_{kM_1M_2\cdots}$构成，其中指标$k,\allowbreak M_1,\allowbreak M_2,\allowbreak \ldots$的取值范围为：$k=1,\allowbreak 2,\allowbreak \ldots$，$M_1,\allowbreak M_2,\allowbreak \ldots=0,\allowbreak 1,\allowbreak 2,\allowbreak \ldots$，但只允许有限个$M_n\ne0$，个数本身可以任意\origNote[意][。]{144}。$\mathsf{H}$的最终形式为：\tnote[：][][从]{从前面的有限$N$公式过渡到这个极限时，原书将两项系数所对应的$M_n+1$与$M_n-1$互换了。这里分别保留两式的原印形式。}
\begin{align*}
 &\mathsf{H}a_{kM_1M_2\cdots}=a'_{kM_1M_2\cdots}\\
 &\quad={}\left(W_k+\sum_{n=1}^{\infty}h\varrho_n M_n\right)a_{kM_1M_2\cdots}\\
 &{}+\sum_{j=1}^{\infty}\sum_{n=1}^{\infty}w_{kj}^n
 \left(\sqrt{M_n+1}\,a_{jM_1M_2\cdots M_n+1\cdots}
 +\sqrt{M_n}\,a_{jM_1M_2\cdots M_n-1\cdots}\right),
\end{align*}
其中$w_{kj}^n$定义为
\[
 w_{kj}^n=\sqrt{\frac h{2\varrho_n}}\sum_{\nu=1}^{l}\frac{e_\nu}{2\pi cm_\nu}
 \bigl(P_\nu^x\overline{\mathfrak A}_{n,x}(Q_\nu^x,Q_\nu^y,Q_\nu^z)+\cdots\bigr)_{kj}.
\]

在从这个结果得出所需的物理结论以前，须再次提醒：这个结果是从电动力学的光理论得到的。首先要确定，我们按固定规则作出的量子力学改写，是否足以说明光偏离波动模型的性质，尤其是它实际上具有的离散粒子性。原本也有这样一种可能：为了把握这种性质，必须直接从光的粒子模型出发，而不是像这里这样，将电磁场“量子化”。

从$\mathsf{H}$的表达式，立即可以看出，它确实包含了某种光粒子。如果略去第二项，就只剩下
\[
 \mathsf{H}_1a_{kM_1M_2\cdots}=\left(W_k+\sum_{n=1}^{\infty}h\varrho_n M_n\right)a_{kM_1M_2\cdots}.
\]
第二项代表某种扰动，如我们很快将看到的，系统$\boldsymbol{S}$从一个“定态”向另一个定态的量子跃迁正由它引起。这才是我们真正关心的现象。但这种扰动的作用远弱于物质系统$\boldsymbol{S}$与已有辐射的作用；如后面将看到的，这两者体现在第一项中。上面的能量表达式可以解释为：系统$\boldsymbol{S}$的能量$W_k$，加上各项能量$h\varrho_n M_n$，$n=1,\allowbreak 2,\allowbreak \ldots$。很自然，可以把$M_n=0,\allowbreak 1,\allowbreak 2,\allowbreak \ldots$理解为能量分别等于$h\varrho_n$的粒子数；而$h\varrho_n$恰是按照Einstein的观点应赋予频率为$\varrho_n$的光量子的能量，参见原注 \hyperlink{note-134}{\getrefnumber{orig-note-134}}。因此，$\mathsf{H}_1$的结构使我们猜想，\hypertarget{origpage-145}{}\phantomsection\label{origpage-145}在$\boldsymbol H$中的电磁场，扣除静电部分后，也就是辐射$\boldsymbol{L}$，可以理解为以下情形：频率为$\varrho_1,\allowbreak \varrho_2,\allowbreak \ldots$、因而能量为$h\varrho_1,\allowbreak h\varrho_2,\allowbreak \ldots$的光量子，其数目分别为$M_1,\allowbreak M_2,\allowbreak \ldots$，各为$0,\allowbreak 1,\allowbreak 2,\allowbreak \ldots$。为什么只计入$\varrho_1,\allowbreak \varrho_2,\allowbreak \ldots$这些频率，也有合理的解释：它们正是空腔$\boldsymbol H$的本征频率，因为矢势$\overline{\mathfrak A}_n(xyz)\gamma\cos2\pi\varrho_n(t-\tau)$表示的是$\boldsymbol H$中仅有的那些定态电磁振动。

不过，这些推测和解释还只有启发性的价值。只有当我们能从辐射$\boldsymbol{L}$的光量子模型出发，同样得到上述$\mathsf{H}$的能量表达式时，问题才算有了满意而最终的答案。前面先作经典考察，是因为光量子与物质相互作用能的表达式尚无从设定；在这一点上，经典电动力学无法直接改写为量子力学。现在却可以了：只要从一般的相互作用能推导出的结果，在形式上与$\mathsf{H}$相符，就可以比较系数，确定相互作用能。

以光量子假说为依据，$\boldsymbol{L}$的位形空间是什么样的，即问题 \hyperref[problem:III.6.1]{1} 应如何回答？设空腔$\boldsymbol H$中的单个光量子由某些坐标表征，用符号$u$概括表示这些坐标\origNote[标][。]{145}。它在$\boldsymbol H$中的定态，设具有波函数$\psi_1(u),\allowbreak \psi_2(u),\allowbreak \ldots$及能量$E_1,\allowbreak E_2,\allowbreak \ldots$，这些波函数构成完全规范正交系；它们对应于频率为$\varrho_1,\allowbreak \varrho_2,\allowbreak \ldots$的电磁本征振动$\overline{\mathfrak A}_1,\allowbreak \overline{\mathfrak A}_2,\allowbreak \ldots$。按照Einstein的观点，必须有$E_n=h\varrho_n$，这一点后面也将证明。此外还须注意：在上面的电磁学讨论中，我们已经把光能量的零点选在它的最小值处，对应指标$M_1=M_2=\cdots=0$。因此，我们已把光不存在的情形也承认为光的一种可能状态，这在物理上也是有根据的；光量子确能被发射和吸收，也就是被产生和消灭。但这种想法对量子力学来说完全陌生：每个粒子都以其坐标参与系统的位形空间，并与整个系统的形式描述紧密相联，以致绝对不可消灭。因此，即使在一个粒子被消灭以后，也必须赋予它一种潜在的存在，使它的坐标继续属于位形空间。于是，能量为$E_n=0$的某个状态$\psi_n(u)$，就必须对应于光量子的不存在。这里选用记号$\psi_0(u)$，$E_0=0$。这样，$\psi_1(u),\allowbreak \psi_2(u),\allowbreak \ldots$对应存在着的光量子，\hypertarget{origpage-146}{}\phantomsection\label{origpage-146}但只有$\psi_0(u),\allowbreak \psi_1(u),\allowbreak \psi_2(u),\allowbreak \ldots$才构成完全正交系。

现在转向全部光量子构成的系统$\boldsymbol{L}$。由于还把不存在的光量子也计入，$\boldsymbol{L}$必须含有足够多的光量子，才能保证以后产生的数目永远不超过这一数量，也就是必须有无限多个。不过，从一开始就对$\boldsymbol{L}$中无限多组成部分进行运算并不方便，所以先假设只有$S$个光量子，$S=1,\allowbreak 2,\allowbreak \ldots$，最后再作极限过程$S\longrightarrow+\infty$。\origNote[。][将]{146}将这$S$个光量子编号为$1,\allowbreak \ldots,\allowbreak S$，其坐标分别记作$u_1,\allowbreak \ldots,\allowbreak u_S$。于是，$\boldsymbol{L}$的位形空间由$u_1,\allowbreak \ldots,\allowbreak u_S$描述，$\boldsymbol{S}+\boldsymbol{L}$的位形空间由$\xi,\allowbreak u_1,\allowbreak \ldots,\allowbreak u_S$描述。因此，$\boldsymbol{S}+\boldsymbol{L}$的一般波函数为$f(\xi,\allowbreak u_1,\allowbreak \ldots,\allowbreak u_S)$，而$\varphi_k(\xi)\allowbreak \psi_{n_1}(u_1)\cdots\psi_{n_S}(u_S)$，$k=1,\allowbreak 2,\allowbreak \ldots$、$u_1,\allowbreak \ldots,\allowbreak u_S=0,\allowbreak 1,\allowbreak 2,\allowbreak \ldots$，构成完全正交系。\tnote[。][][原]{原书这处指标范围印作$u_1,\allowbreak \ldots,\allowbreak u_S$；所列波函数的指标应为$n_1,\allowbreak \ldots,\allowbreak n_S$。这里保留原印记号。}

光量子还有一个根本性质：彼此完全相同。也就是说，若两个光量子的坐标$u$相同，世界上就没有任何办法将它们区分开来。换句话说，编号为$m,\allowbreak n$的光量子分别具有$u_m=u'$、$u_n=u''$的状态，与具有$u_m=u''$、$u_n=u'$的状态，没有区别。这是经典的描述方式，而非量子力学描述，因为这里给出的是$u$的值，不是波函数$\varphi(u)$！用量子力学来说，就是波函数$f(\xi,\allowbreak u_1,\allowbreak \ldots,\allowbreak u_m,\allowbreak \ldots,\allowbreak u_n,\allowbreak \ldots,\allowbreak u_S)$与$f(\xi,\allowbreak u_1,\allowbreak \ldots,\allowbreak u_n,\allowbreak \ldots,\allowbreak u_m,\allowbreak \ldots,\allowbreak u_S)$对应的状态不可区分。也就是说，每个物理量$\mathfrak R$在两者中的期望值相同；既然对$F(\mathfrak R)$也如此，其统计规律便完全相同，参见 \hyperref[sec:III.1]{III.1} 中关于 \hyperref[post:III.1.E1]{$E_1$}、\hyperref[post:III.1.E2]{$E_2$} 的讨论。将交换$u_m,\allowbreak u_n$的函数运算记作$O_{mn}$；容易看出，$O_{mn}$既是厄米算符又是酉算符，且$O_{mn}^2=1$。于是上述要求意味着，$\mathfrak R$在$f$中的期望值与在$O_{mn}f$中的相同，即
\[
 (Rf,f)=(RO_{mn}f,O_{mn}f)=(O_{mn}RO_{mn}f,f),
\]
故$R=O_{mn}RO_{mn}$，也就是$O_{mn}R=RO_{mn}$。这意味着，在目前情形，只允许那些与所有$O_{mn}$，$m,\allowbreak n=1,\allowbreak \ldots,\allowbreak S$、$m\ne n$，对易的算符$R$；考虑到$O_{mn}$的定义，也就是说，$R$对一切坐标$u_1,\allowbreak \ldots,\allowbreak u_S$都以对称的方式发生作用。

一个对所有变量$u_1,\allowbreak \ldots,\allowbreak u_S$都对称的波函数$f$，即对$m,\allowbreak n=1,\allowbreak \ldots,\allowbreak S$、$m\ne n$满足$O_{mn}f=f$的波函数，经这样的算符$R$作用后，仍是对称的：$O_{mn}Rf=RO_{mn}f=Rf$。这些$f$构成闭线性流形，即全部$f$所成的希尔伯特空间$\mathfrak R_{\infty}^{(S)}$中的希尔伯特子空间$\overline{\mathfrak R}_{\infty}^{(S)}$。\hypertarget{origpage-147}{}\phantomsection\label{origpage-147}各个$R$把$\overline{\mathfrak R}_{\infty}^{(S)}$中的元素映射为同一空间中的元素，因而可以看作希尔伯特空间$\overline{\mathfrak R}_{\infty}^{(S)}$上的算符。所以，就量子力学的用途而言，$\overline{\mathfrak R}_{\infty}^{(S)}$与原先考虑的$\mathfrak R_{\infty}^{(S)}$同样可用。于是出现一个问题：既然$\boldsymbol{L}$对于光量子的交换具有对称性，是否应当把波函数限制为对称的，也就是用$\overline{\mathfrak R}_{\infty}^{(S)}$替代$\mathfrak R_{\infty}^{(S)}$？我们就这样做；后面得到的成功，即与电磁学所推导的$\mathsf{H}$表达式完全吻合，将为这一做法提供依据\origNote[据][。]{147}。

$\varphi_k(\xi)\allowbreak \psi_{n_1}(u_1)\cdots\psi_{n_S}(u_S)$构成$\mathfrak R_{\infty}^{(S)}$中的完全规范正交系，下面用它们构造$\overline{\mathfrak R}_{\infty}^{(S)}$中的一个这样的系。设$M_0,\allowbreak M_1,\allowbreak \ldots$为任意取自$0,\allowbreak 1,\allowbreak 2,\allowbreak \ldots$的数，满足$M_0+M_1+\cdots=S$，因而只有有限个不为$0$。用$[M_1,\allowbreak M_2,\allowbreak \ldots]$表示所有这样的指标组$u_1,\allowbreak \ldots,\allowbreak u_S$：其中$0$出现$M_0$次，$1$出现$M_1$次，等等；恰有$M_0!M_1!\cdots$个不同的指标组。\tnote[。][令][本]{本页的计数与归一化因子均按原印转录。对于具有所述重数的不同指标排列，其数目应为$S!/(M_0!M_1!\cdots)$，而非此处所印的阶乘乘积。原书还交替使用$[M_1,\allowbreak M_2,\allowbreak \ldots]$和$[M_0,\allowbreak M_1,\allowbreak \ldots]$，并在指称指标时交替使用$u_j$和$n_j$。}令
\begin{align*}
 &\Phi_{M_0M_1\cdots}(u_1\cdots u_S)\\
 &\quad=\sum_{n_1\cdots n_S\mathwordspace{}\text{在}\mathwordspace{}[M_0M_1\cdots]}
 \psi_{n_1}(u_1)\cdots\psi_{n_S}(u_S).
\end{align*}
$\Phi_{M_0M_1\cdots}$由$M_0!M_1!\cdots$个两两正交、范数为$1$的项相加而成，所以它的范数平方等于$M_0!M_1!\cdots$个$1$相加，其范数为$\sqrt{M_0!M_1!\cdots}$。两个不同的$\Phi_{M_0M_1\cdots}$的各项彼此正交，所以它们也正交。因此，
\begin{align*}
 &\psi_{M_0M_1\cdots}(u_1\cdots u_S)\\
 &\quad=\frac1{\sqrt{M_0!M_1!\cdots}}\Phi_{M_0M_1\cdots}(u_1\cdots u_S)
\end{align*}
构成规范正交系。对$u_1,\allowbreak \ldots,\allowbreak u_S$对称的$f(\xi,\allowbreak u_1,\allowbreak \ldots,\allowbreak u_S)$，与$\varphi_k(\xi)\allowbreak \Phi_{M_0M_1\cdots}(u_1\cdots u_S)$的所有组成项都有相同的内积。所以，如果它与$\varphi_k(\xi)\allowbreak \Phi_{M_0M_1\cdots}(u_1\cdots u_S)$正交，也就是与$\varphi_k(\xi)\allowbreak \psi_{M_0M_1\cdots}(u_1\cdots u_S)$正交，它就与这些组成项中的每一项都正交。也就是说，如果它与所有$\varphi_k(\xi)\allowbreak \psi_{M_0M_1\cdots}(u_1\cdots u_S)$都正交，它也与所有$\varphi_k(\xi)\allowbreak \psi_{n_1}(u_1)\cdots\psi_{n_S}(u_S)$都正交，因而等于$0$。因此，$\varphi_k(\xi)\allowbreak \psi_{M_0M_1\cdots}(u_1\cdots u_S)$本身属于$\overline{\mathfrak R}_{\infty}^{(S)}$，并在其中构成完全规范正交系。

现在系统地考察$\boldsymbol{S}+\boldsymbol{L}$中出现的各种能量。第一，问题 \hyperref[problem:III.6.2alpha]{2的$\alpha)$}，是$\boldsymbol{S}$的能量。对$\boldsymbol{S}$，其算符由$\mathsf{H}_0\varphi_k(\xi)=W_k\varphi_k(\xi)$定义；对$\boldsymbol{S}+\boldsymbol{L}$，则由
\[
 \mathsf{H}_0\varphi_k(\xi)\psi_{M_0M_1\cdots}(u_1\cdots u_S)
 =W_k\varphi_k(\xi)\psi_{M_0M_1\cdots}(u_1\cdots u_S)
\]
定义。第二，问题 \hyperref[problem:III.6.2beta]{2的$\beta)$}，是每个光量子$l'$的能量。对$l'$，有$\mathsf{H}_{l'}\psi_n(u)=E_n\psi_n(u)$，\hypertarget{origpage-148}{}\phantomsection\label{origpage-148}因此，对$\boldsymbol{S}+\boldsymbol{L}$中的第$m$个光量子，$m=1,\allowbreak \ldots,\allowbreak S$，有
\begin{align*}
 &\mathsf{H}_{l_m}\varphi_k(\xi)\psi_{n_1}(u_1)\cdots\psi_{n_m}(u_m)\cdots\psi_{n_S}(u_S)\\
 &\quad=E_{n_m}\varphi_k(\xi)\psi_{n_1}(u_1)\cdots\psi_{n_m}(u_m)\cdots\psi_{n_S}(u_S),
\end{align*}
并须构造$\mathsf{H}_l=\mathsf{H}_{l_1}+\cdots+\mathsf{H}_{l_S}$。最后，问题 \hyperref[problem:III.6.2gamma]{2的$\gamma)$}，设单个光量子$l'$与$\boldsymbol{S}$的相互作用能由某个尚不详知的算符$V$描述，先用它的矩阵来表征：
\[
 V_{l'}\varphi_k(\xi)\psi_n(u)=\sum_{j=1}^{\infty}\sum_{p=0}^{\infty}V_{kn/jp}\varphi_j(\xi)\psi_p(u).
\]
在$\boldsymbol{S}+\boldsymbol{L}$中，对于第$m$个光量子，有
\begin{align*}
 &V_{l_m}\varphi_k(\xi)\psi_{n_1}(u_1)\cdots\psi_{n_m}(u_m)\cdots\psi_{n_S}(u_S)\\
 &\quad=\sum_{j=1}^{\infty}\sum_{p=0}^{\infty}V_{kn_m/jp}\varphi_j(\xi)
 \psi_{n_1}(u_1)\cdots\psi_p(u_m)\cdots\psi_{n_S}(u_S)\\
 &\quad=\sum_{j=1}^{\infty}\sum_{p_1,\ldots,p_S=0}^{\infty}
 \delta(n_1-p_1)\cdots V_{kn_m/jp_m}\cdots\delta(n_S-p_S)\\
 &\qquad{}\times\varphi_j(\xi)\psi_{p_1}(u_1)\cdots\psi_{p_m}(u_m)\cdots\psi_{p_S}(u_S)
\end{align*}
其中$\delta(n)$在$n=0$时为$1$，在$n\ne0$时为$0$；并须构造$\mathsf{H}_w=V_{l_1}+\cdots+V_{l_S}$。

综上，得到
\begin{align*}
 &\mathsf{H}\varphi_k(\xi)\psi_{n_1}(u_1)\cdots\psi_{n_S}(u_S)\\
 &\quad=(W_k+E_{n_1}+\cdots+E_{n_S})\varphi_k(\xi)\psi_{n_1}(u_1)\cdots\psi_{n_S}(u_S)\\
 &\qquad{}+\sum_{j=1}^{\infty}\sum_{p_1,\ldots,p_S=0}^{\infty}\sum_{m=1}^{S}
 \delta(n_1-p_1)\cdots V_{kn_m/jp_m}\cdots\delta(n_S-p_S)\\
 &\qquad{}\times\varphi_j(\xi)\psi_{p_1}(u_1)\cdots\psi_{p_S}(u_S).
\end{align*}
稍作变形便得
\begin{align*}
 &\mathsf{H}\varphi_k(\xi)\Phi_{M_0M_1\cdots}(u_1\cdots u_S)\\
 &\quad=\left(W_k+\sum_{n=0}^{\infty}M_nE_n\right)\varphi_k(\xi)\Phi_{M_0M_1\cdots}(u_1\cdots u_S)\\
 &\qquad{}+\sum_{j=1}^{\infty}\sum_{n,p=0}^{\infty}M_nV_{kn/jp}\varphi_j(\xi)
 \Phi_{M_0M_1\cdots M_n-1\cdots M_p+1\cdots}(u_1\cdots u_S)
\end{align*}
当$n=p$时，$\cdots M_n-1\cdots M_p+1\cdots$应换为$\cdots M_n\cdots$。所以，对规范正交函数，有
\begin{align*}
 &\mathsf{H}\varphi_k(\xi)\psi_{M_0M_1\cdots}(u_1\cdots u_S)\\
 &\quad=\left(W_k+\sum_{n=0}^{\infty}M_nE_n\right)\varphi_k(\xi)\psi_{M_0M_1\cdots}(u_1\cdots u_S)\\
 &\qquad{}+\sum_{j=1}^{\infty}\sum_{n,p=0}^{\infty}\sqrt{M_n(M_p+1-\delta(n-p))}\,V_{kn/jp}\varphi_j(\xi)\\
 &\qquad{}\times\psi_{M_0M_1\cdots M_n-1\cdots M_p+1\cdots}(u_1\cdots u_S).
\end{align*}
$\overline{\mathfrak R}_{\infty}^{(S)}$中的一般$f(\xi,\allowbreak u_1\cdots u_S)$可以按这些正交函数\hypertarget{origpage-149}{}\phantomsection\label{origpage-149}展开：\tnote[：][][原]{原书本页用大写$\Psi$表示前两页中记为小写$\psi$的归一化函数。这里保留其大小写变化。}
\[
 f(\xi u_1\cdots u_S)=\sum_{k=1}^{\infty}
 \sum_{\substack{M_0,M_1,\ldots\ge0\\M_0+M_1+\cdots=S}}
 a_{kM_0M_1\cdots}\varphi_k(\xi)\Psi_{M_0M_1\cdots}(u_1\cdots u_S).
\]
因此，也可以将$\overline{\mathfrak R}_{\infty}^{(S)}$理解为序列$a_{kM_0M_1\cdots}$所成的希尔伯特空间，其中$k=1,\allowbreak 2,\allowbreak \ldots$，$M_0,\allowbreak M_1,\allowbreak \ldots=0,\allowbreak 1,\allowbreak 2,\allowbreak \ldots$，$M_0+M_1+\cdots=S$，并且$\sum_{kM_0M_1\cdots}|a_{kM_0M_1\cdots}|^2$有限。这时
\[
 \mathsf{H}a_{kM_0M_1\cdots}=a'_{kM_0M_1\cdots}
\]
由下式定义：
\begin{align*}
 &\mathsf{H}\sum_{k=1}^{\infty}\sum_{\substack{M_0,M_1,\ldots\ge0\\M_0+M_1+\cdots=S}}
 a_{kM_0M_1\cdots}\varphi_k(\xi)\Psi_{M_0M_1\cdots}(u_1\cdots u_S)\\
 &\quad=\sum_{k=1}^{\infty}\sum_{\substack{M_1,M_2,\ldots\ge0\\M_1+M_2+\cdots=S}}
 a'_{kM_0M_1\cdots}\varphi_k(\xi)\Psi_{M_0M_1\cdots}(u_1\cdots u_S),
\end{align*}
故
\begin{align*}
 &\mathsf{H}a_{kM_0M_1\cdots}=a'_{kM_0M_1\cdots}\\
 &\quad={}\left(W_k+\sum_{n=0}^{\infty}M_nE_n\right)a_{kM_0M_1\cdots}\\
 &{}+\sum_{j=1}^{\infty}\sum_{n,p=0}^{\infty}\sqrt{M_n(M_p+1-\delta(n-p))}\,
 \overline{V_{kn/jp}}\,a_{jM_0M_1\cdots M_n-1\cdots M_p+1\cdots}.
\end{align*}
与$\varphi_k(\xi)\Psi_{M_0M_1\cdots}(u_1\cdots u_S)$的公式相比，$k,\allowbreak j$及$n,\allowbreak p$的作用已互换。考虑到$V$是厄米算符，我们把$V_{jp/kn}$写成了$\overline{V_{kn/jp}}$。

现在为极限过程$S\longrightarrow+\infty$作准备。$M_0$由$M_1,\allowbreak M_2,\allowbreak \ldots$同时确定，因为$M_0=S-M_1-M_2-\cdots$，所以可将$a_{kM_0M_1\cdots}$改记作$a_{kM_1M_2\cdots}$。此时，指标须满足$k=1,\allowbreak 2,\allowbreak \ldots$、$M_1,\allowbreak M_2,\allowbreak \ldots=0,\allowbreak 1,\allowbreak 2,\allowbreak \ldots$、$M_1+M_2+\cdots\le S$。注意到$E_0=0$，并引入记号
\[
 SV_{k0/j0}=V_{k/j},\qquad \sqrt S\,V_{k0/jn}=V_{k/jn},\qquad
 \sqrt S\,V_{kn/j0}=\overline{V_{j/kn}}
\]
这里$V_{kn/jp}$是厄米的，便得到
\begin{align*}
 &\mathsf{H}a_{kM_1M_2\cdots}=a'_{kM_1M_2\cdots}\\
 &\quad={}\left(W_k+\sum_{n=1}^{\infty}M_nE_n\right)a_{kM_1M_2\cdots}\\
 &{}+\sum_{j=1}^{\infty}V_{k/j}a_{jM_1M_2\cdots}\\
 &{}+\sum_{j=1}^{\infty}\sum_{n=1}^{\infty}\sqrt{M_n}
 \sqrt{\frac{S-M_1-M_2-\cdots+1}{S}}\,V_{j/kn}a_{jM_1M_2\cdots M_n-1\cdots}\\
 &{}+\sum_{j=1}^{\infty}\sum_{n=1}^{\infty}\sqrt{M_n+1}
 \sqrt{\frac{S-M_1-M_2-\cdots}{S}}\,\overline{V_{k/jn}}\,a_{jM_1M_2\cdots M_n+1\cdots}\\
 &{}+\sum_{j=1}^{\infty}\sum_{n,p=1}^{\infty}\sqrt{M_n(M_p+1)}\,
 \overline{V_{kn/jp}}\,a_{jM_1M_2\cdots M_n-1\cdots M_p+1\cdots}.
\end{align*}
现在可以令$S\longrightarrow+\infty$。这时，$a_{kM_1M_2\cdots}$重新定义在所有满足$k=1,\allowbreak 2,\allowbreak \ldots$、$M_1,\allowbreak M_2,\allowbreak \ldots=0,\allowbreak 1,\allowbreak 2,\allowbreak \ldots$，\hypertarget{origpage-150}{}\phantomsection\label{origpage-150}且只有有限个$M_n\ne0$的指标序列$kM_1M_2\cdots$上；非零项个数可以任意，参见原注 \hyperlink{note-144}{\getrefnumber{orig-note-144}}。$\mathsf{H}$变成\tnote[成][][从]{从前面的有限$S$公式过渡到此极限式时，原书将两项系数所对应的$M_n+1$与$M_n-1$互换了。这里保留各式的原印形式。}
\begin{align*}
 &\mathsf{H}a_{kM_1M_2\cdots}=a'_{kM_1M_2\cdots}\\
 &\quad={}\left(W_k+\sum_{n=1}^{\infty}M_nE_n\right)a_{kM_1M_2\cdots}
 +\sum_{j=1}^{\infty}V_{k/j}a_{jM_1M_2\cdots}\\
 &{}+\sum_{j=1}^{\infty}\sum_{n=1}^{\infty}
 \bigl(V_{j/kn}\sqrt{M_n+1}\,a_{jM_1M_2\cdots M_n+1\cdots}\\
 &\qquad{}+\overline{V_{k/jn}}\sqrt{M_n}\,a_{jM_1M_2\cdots M_n-1\cdots}\bigr)\\
 &{}+\sum_{j=1}^{\infty}\sum_{n,p=1}^{\infty}\overline{V_{kn/jp}}\sqrt{M_n(M_p+1)}\,
 a_{jM_1M_2\cdots M_n-1\cdots M_p+1\cdots}.
\end{align*}
它与从光的电磁理论推导出的方程之间的相似性一目了然。要让两者完全吻合，只须令
\[
 E_n=h\varrho_n,\qquad V_{k/j}=0,\qquad V_{k/jn}=w_{jk}^n=\overline{w_{kj}^n},\qquad V_{kn/jp}=0.
\]
可见，光量子的观点与经典电磁学的观点一致，只要：

1．按照一般的量子力学规则改写后者；

2．按照Einstein的规则，把光量子的能量取为频率的$h$倍；

3．正确设定光量子与物质之间的相互作用，参见上面$V$的各个表达式。

这样，早期量子理论最严重的一个悖论，即光的双重本性——电磁波或离散粒子、光量子——便得到了十分出色的解决\origNote[决][。]{148}。当然，要为刚算出的光与物质之间的相互作用能$V$找到直接而直观的意义，仍很困难。尤其是，它仅有的不为零的矩阵元$V_{kn/jp}$，即$n=0,\allowbreak p\ne0$或$n\ne0,\allowbreak p=0$的那些矩阵元，竟依赖于所有可能光量子的总数$S$：它们与$1/\sqrt S$成正比，而最终又必须令$S\longrightarrow\infty$。不过，模型描述终究只是一种近似；理论的准确内容是由$\mathsf{H}$算符的表达式给出的，考虑到这一点，上述情形尚可接受。

现在回到真正的任务：确定跃迁概率。按照含时Schrödinger微分方程，$a_{kM_1M_2\cdots}=a_{kM_1M_2\cdots}(t)$的变化由下式确定：
\begin{align*}
 &\frac h{2\pi i}\frac\partial{\partial t}a_{kM_1M_2\cdots}=-\mathsf{H}a_{kM_1M_2\cdots}\\
 &\quad=-\left(W_k+\sum_{n=1}^{\infty}h\varrho_n M_n\right)a_{kM_1M_2\cdots}\\
 &{}-\sum_{j=1}^{\infty}\sum_{n=1}^{\infty}w_{kj}^n
 \left(\sqrt{M_n+1}\,a_{jM_1M_2\cdots M_n+1\cdots}
 +\sqrt{M_n}\,a_{jM_1M_2\cdots M_n-1\cdots}\right).
\end{align*}\hypertarget{origpage-151}{}\phantomsection\label{origpage-151}$a_{kM_1M_2\cdots}$的主要变化来自这个表达式的第一项，宜用下列设定将它分离出来：
\[
 a_{kM_1M_2\cdots}(t)=e^{-\frac{2\pi i}{h}(W_k+\sum_{n=1}^{\infty}h\varrho_n M_n)t}\,b_{kM_1M_2\cdots}(t).
\]
于是\tnote[是][][原]{原书下列微分方程的两项之间印作减号，前面的方程却为加号。这里照录。}
\begin{align*}
 &\frac\partial{\partial t}b_{kM_1M_2\cdots}\\
 &\quad=-\frac{2\pi i}{h}\sum_{j=1}^{\infty}\sum_{n=1}^{\infty}w_{kj}^n
 \bigl(e^{-\frac{2\pi i}{h}(W_j-W_k+h\varrho_n)t}\sqrt{M_n+1}\,b_{jM_1M_2\cdots M_n+1\cdots}\\
 &\qquad{}-e^{-\frac{2\pi i}{h}(W_j-W_k-h\varrho_n)t}\sqrt{M_n}\,b_{jM_1M_2\cdots M_n-1\cdots}\bigr).
\end{align*}
$a_{kM_1M_2\cdots}$及$b_{kM_1M_2\cdots}$的物理意义，由它们的构造已经确定。在$\overline M_0+\overline M_1+\overline M_2+\cdots=S$时，$\varphi_{\bar k}(\xi)\Psi_{\overline M_0\overline M_1\cdots}(u_1\cdots u_S)$表示以下状态：$\boldsymbol{S}$处于第$\bar k$条量子轨道，分别有$\overline M_0,\allowbreak \overline M_1,\allowbreak \overline M_2,\allowbreak \ldots$个光量子处于$\psi_0,\allowbreak \psi_1,\allowbreak \psi_2,\allowbreak \ldots$状态，即$\overline M_0$个处于“不存在”的状态，而$\overline M_1,\allowbreak \overline M_2,\allowbreak \ldots$个分别处于对应于本征振动$\overline{\mathfrak A}_1,\allowbreak \overline{\mathfrak A}_2,\allowbreak \ldots$的状态。属于这个波函数的$a_{kM_1M_2\cdots}$为：\tnote[：][][原]{原书下式第一个占据数因子印作$M_0$。下文的概率公式及初始条件中，同一构型用$M_1$表示。此处保留原下标。}
\[
 a_{kM_1M_2\cdots}=\delta(k-\bar k)\delta(M_0-\overline M_1)\delta(M_2-\overline M_2)\cdots.
\]
只有有限个因子不等于$1$，因为除有限个例外外，总有$M_n=\overline M_n=0$。即使令$S\longrightarrow\infty$，这也当然仍然成立。对于$\boldsymbol{S}+\boldsymbol{L}$的任意状态$a_{kM_1M_2\cdots}$，某个给定构型在测量时出现的概率，按照 \hyperref[sec:III.3]{III.3} 关于单重纯点谱的论述，为
\begin{align*}
 &\left|\sum_{k,M_1,M_2,\ldots}a_{kM_1M_2\cdots}
 \delta(k-\bar k)\delta(M_1-\overline M_1)\delta(M_2-\overline M_2)\cdots\right|^2\\
 &\qquad=|a_{\bar k\overline M_1\overline M_2\cdots}|^2
 =|b_{\bar k\overline M_1\overline M_2\cdots}|^2.
\end{align*}
特别地，$\boldsymbol{S}$处于第$\bar k$条量子轨道的总概率为
\[
 \vartheta_{\bar k}=\sum_{\overline M_1,\overline M_2,\ldots}
 |b_{\bar k\overline M_1\overline M_2\cdots}|^2.
\]
设某一时刻$t=0$，原子处于第$\bar k$个状态，并分别存在$\overline M_1,\allowbreak \overline M_2,\allowbreak \ldots$个处于$\overline{\mathfrak A}_1,\allowbreak \overline{\mathfrak A}_2,\allowbreak \ldots$状态的光量子，即
\[
 b_{kM_1M_2\cdots}=a_{kM_1M_2\cdots}
 =\delta(k-\bar k)\delta(M_1-\overline M_1)\delta(M_2-\overline M_2)\cdots.
\]
按上述微分方程，在第一近似下，也就是$t$短到仍可把方程右边看作常量时，\hypertarget{origpage-152}{}\phantomsection\label{origpage-152}只有那些$M_1,\allowbreak M_2,\allowbreak \ldots,\allowbreak M_n+1,\allowbreak \ldots$或$M_1,\allowbreak M_2,\allowbreak \ldots,\allowbreak M_n-1,\allowbreak \ldots$与$\overline M_1,\allowbreak \overline M_2,\allowbreak \ldots$相同的$\frac\partial{\partial t}b_{kM_1M_2\cdots}$才可能不为$0$。也就是说，只有$k,\allowbreak \overline M_1,\allowbreak \overline M_2,\allowbreak \ldots,\allowbreak \overline M_n\pm1,\allowbreak \ldots$这些指标组合有此可能。对此积分得到\tnote[到][][原]{原书第二式的分母含$-h\varrho_n$，指数中却含$+h\varrho_n$；后面的吸收概率采用了后一个符号。这里保留原式。}
\begin{align*}
 b_{k\overline M_1\overline M_2\cdots\overline M_n+1\cdots}
 &=w_{\bar k k}^n\frac{1-e^{-\frac{2\pi i}{h}(W_{\bar k}-W_k-h\varrho_n)t}}
 {W_{\bar k}-W_k-h\varrho_n}\sqrt{\overline M_n+1},\\
 b_{k\overline M_1\overline M_2\cdots\overline M_n-1\cdots}
 &=w_{\bar k k}^n\frac{1-e^{-\frac{2\pi i}{h}(W_{\bar k}-W_k+h\varrho_n)t}}
 {W_{\bar k}-W_k-h\varrho_n}\sqrt{\overline M_n}.
\end{align*}
所有其他$b_{kM_1M_2\cdots}$在这个近似下均为$0$，唯独$b_{\bar k\overline M_1\overline M_2\cdots}$例外。在此近似中，即略去$t^2$项后，它应等于初值$1$。不过，断言$\frac\partial{\partial t}b_{\bar k\overline M_1\overline M_2\cdots}=0$并不稳妥，因为在这种情况下，微分方程右边包含无限多个在本近似下不消失的$b_{k\overline M_1\overline M_2\cdots\overline M_n\pm1\cdots}$。所以，即使$t$很小时每一项都小，也还不能断定它们的总和小。实际上，计算下一阶近似会表明，$b_{\bar k\overline M_1\overline M_2\cdots}$偏离$1$的不是$t^2$项，而是$t$项\origNote[项][。]{149}。不过，由于
\[
 \sum_{k,M_1,M_2,\ldots}|b_{kM_1M_2\cdots}|^2
 =\sum_{k,M_1,M_2,\ldots}|a_{kM_1M_2\cdots}|^2=1,
\]
所以
\[
 |b_{\bar k\overline M_1\overline M_2\cdots}|^2
 =1-\sum_{kM_1M_2\cdots\ne\bar k\overline M_1\overline M_2\cdots}|b_{kM_1M_2\cdots}|^2,
\]
因而根本不必直接求出这个$b_{\bar k\overline M_1\overline M_2\cdots}$。

这里已能清楚看出过程的定性特征。$b_{k\overline M_1\overline M_2\cdots\overline M_n+1\cdots}$对应于发射一个频率为$\varrho_n$、属于$\overline{\mathfrak A}_n$振动的光量子；分母$W_{\bar k}-W_k-h\varrho_n$越小，它增长得越多，也就是说，光的本征频率$\varrho_n$越接近“Bohr频率”$(W_{\bar k}-W_k)/h$，\origNote[，][增]{150}增长就越多。相应地，代表吸收的$b_{k\overline M_1\overline M_2\cdots\overline M_n-1\cdots}$，则在$\varrho_n$越接近$(W_k-W_{\bar k})/h$时增长越多。可见，Bohr频率关系并不严格成立，因为$\varrho_n$并不包含一切频率；但当时间$t$很短、$\varrho_n$又排列得很密时，即空腔$\boldsymbol H$很大时，它却以压倒性的概率成立。此外，系数$w_{\bar k k}^n$也会提高这些过程的发生频率；稍后我们就能用跃迁概率来解释这些系数。

\hypertarget{origpage-153}{}\phantomsection\label{origpage-153}由上面的$b_{k\overline M_1\overline M_2\cdots\overline M_n\pm1\cdots}$公式，还可得发射概率\origNote[率][、]{151}、吸收概率\origNote[率][及]{151}及其余构型的概率：
\begin{align*}
 &|b_{k\overline M_1\overline M_2\cdots\overline M_n+1\cdots}|^2\\
 &\quad=\frac2{h^2}(\overline M_n+1)|w_{\bar k k}^n|^2
 \frac{1-\cos 2\pi\left(\varrho_n-\frac{W_{\bar k}-W_k}{h}\right)t}
 {\left(\varrho_n-\frac{W_{\bar k}-W_k}{h}\right)^2},\\
 &|b_{k\overline M_1\overline M_2\cdots\overline M_n-1\cdots}|^2\\
 &\quad=\frac2{h^2}\overline M_n|w_{\bar k k}^n|^2
 \frac{1-\cos 2\pi\left(\varrho_n-\frac{W_k-W_{\bar k}}{h}\right)t}
 {\left(\varrho_n-\frac{W_k-W_{\bar k}}{h}\right)^2},\\
 &|b_{kM_1M_2\cdots}|^2=0\\
 &\quad\text{当}\mathwordspace{}kM_1M_2\cdots\ne\bar k\overline M_1\overline M_2\cdots,\;\\
 &\qquad k\overline M_1\overline M_2\cdots\overline M_n\pm1\cdots.
\end{align*}
从而对于$k\ne\bar k$，有
\begin{align*}
 \vartheta_k={}&\sum_{n=1}^{\infty}\frac2{h^2}(\overline M_n+1)|w_{\bar k k}^n|^2
 \frac{1-\cos 2\pi\left(\varrho_n-\frac{W_{\bar k}-W_k}{h}\right)t}
 {\left(\varrho_n-\frac{W_{\bar k}-W_k}{h}\right)^2}\\
 &{}+\sum_{n=1}^{\infty}\frac2{h^2}\overline M_n|w_{\bar k k}^n|^2
 \frac{1-\cos 2\pi\left(\varrho_n-\frac{W_k-W_{\bar k}}{h}\right)t}
 {\left(\varrho_n-\frac{W_k-W_{\bar k}}{h}\right)^2}.
\end{align*}
第一个对$n$的和对应发射，第二个对应吸收。为了给出这些$\vartheta_k$的闭式表达式，下面作一些简化假定：一方面把$\boldsymbol H$取得很大，即令其体积$\mathcal V\longrightarrow\infty$；另一方面，对$\boldsymbol H$中的本征振动$\overline{\mathfrak A}_n$作统计处理。为此，在上述两个和中，分别把$\varrho_n$处于$\varrho$与$\varrho+d\varrho$之间的各项合并，把$w_{\bar k k}^n$的值代入，并假定$d\varrho\ll\varrho$：
\begin{multline*}
 \frac1{4\pi^2c^2h\varrho}
 \left[\sum_{\varrho\le\varrho_n<\varrho+d\varrho}
 \left|\sum_{\nu=1}^{l}\frac{e_\nu}{m_\nu}
 \bigl(P_\nu^x\overline{\mathfrak A}_{n,x}(Q_\nu^x,Q_\nu^y,Q_\nu^z)+\cdots\bigr)_{k\bar k}\right|^2
 (\overline M_n+1)\right]\\
 {}\times\frac{1-\cos2\pi\left(\varrho-\frac{W_{\bar k}-W_k}{h}\right)t}
 {\left(\varrho-\frac{W_{\bar k}-W_k}{h}\right)^2},
\end{multline*}
另一项形式相同，但用$\overline M_n$代替$\overline M_n+1$，用$(W_k-W_{\bar k})/h$代替$(W_{\bar k}-W_k)/h$。现在需要确定方括号中的量。

通常描述$M_1,\allowbreak M_2,\allowbreak \ldots$时，并不逐一给出它们的值，而只提供少得多的信息，即各频段的能量密度。也就是说，频率区间$\varrho$至\hypertarget{origpage-154}{}\phantomsection\label{origpage-154}$\varrho+d\varrho$内，单位体积的光能为$I(\varrho)d\varrho$。这意味着
\begin{gather*}
 \sum_{\varrho\le\varrho_n<\varrho+d\varrho}h\varrho_n\overline M_n
 \approx h\varrho\sum_{\varrho\le\varrho_n<\varrho+d\varrho}\overline M_n
 =\mathcal V I(\varrho)d\varrho,\\
 \sum_{\varrho\le\varrho_n<\varrho+d\varrho}\overline M_n
 =\frac{\mathcal V I(\varrho)}{h\varrho}\,d\varrho.
\end{gather*}
根据Weyl的一个普遍有效的渐近公式，参见原注 \hyperlink{note-140}{\getrefnumber{orig-note-140}} 所引之处，区间$\varrho\le\varrho_n<\varrho+d\varrho$中$\varrho_n$的数目为$\frac{8\pi\mathcal V\varrho^2}{c^2}d\varrho$，因此\tnote[此][][原]{原书此处Weyl模式计数公式印作$c^2$，紧随其后的表达式却用$c^3$。这里保留这两种原印幂次。}
\[
 \sum_{\varrho\le\varrho_n<\varrho+d\varrho}(\overline M_n+1)
 \approx\frac{\mathcal V\left(I(\varrho)+\frac{8\pi h\varrho^3}{c^3}\right)}{h\varrho}\,d\varrho.
\]
对于上面的方括号，如果
\[
 \left|\sum_{\nu=1}^{l}\frac{e_\nu}{m_\nu}
 \bigl(P_\nu^x\overline{\mathfrak A}_{n,x}(Q_\nu^x,Q_\nu^y,Q_\nu^z)+\cdots\bigr)_{k\bar k}\right|^2
\]
在区间$\varrho\le\varrho_n<\varrho+d\varrho$内足够迅速地围绕某个平均值波动，将此平均值记作$w_{k\bar k}(\varrho)$，那么相应的两个表达式就是
\[
 w_{k\bar k}(\varrho)\frac{\mathcal V\left(I(\varrho)+\frac{8\pi h\varrho^3}{c^3}\right)}{h\varrho}\,d\varrho
 \quad\text{及}\quad
 w_{k\bar k}(\varrho)\frac{\mathcal V I(\varrho)}{h\varrho}\,d\varrho.
\]
再用$\nu_{\bar k k}$表示$(W_{\bar k}-W_k)/h$，用$\nu_{k\bar k}$表示$(W_k-W_{\bar k})/h$，上述各个和便给出
\begin{align*}
 \vartheta_k=\frac{\mathcal V}{4\pi^2c^2h^2}\int_0^{\infty}
 \biggl\{&\left(I(\varrho)+\frac{8\pi h}{c^3}\varrho^3\right)
 \frac{1-\cos2\pi(\varrho-\nu_{\bar k k})t}{(\varrho-\nu_{\bar k k})^2}\\
 &{}+I(\varrho)\frac{1-\cos2\pi(\varrho-\nu_{k\bar k})t}{(\varrho-\nu_{k\bar k})^2}
 \biggr\}\frac{w_{k\bar k}(\varrho)}{\varrho^2}\,d\varrho.
\end{align*}
对于小的$t$，这个积分显然一般为$t^2$的数量级，因为$1-\cos2\pi ct$就是如此；但积分区域内分母$(\varrho-\nu_{\bar k k})^2$或$(\varrho-\nu_{k\bar k})^2$很小的部分例外。那里可能产生比$t^2$大的贡献，而若果真如此，这些贡献就给出$\vartheta_k$的渐近表达式。实际上，我们将看到确实如此，因为会得到$t$数量级的贡献；现在就来求它。

由于$\nu_{\bar k k}=-\nu_{k\bar k}=(W_{\bar k}-W_k)/h$，在$W_{\bar k}>W_k$时，只有第一项的分母可能变小；在$W_{\bar k}<W_k$时，则只有第二项如此。因此，前一种情形只保留第一项，后一种情形只保留第二项。

远离$\nu_{\bar k k}$或$\nu_{k\bar k}$的$\varrho$，只给积分带来$t^2$的贡献。把正的那个频率简记为$\bar\nu_{k\bar k}=|W_{\bar k}-W_k|/h$，\hypertarget{origpage-155}{}\phantomsection\label{origpage-155}便可把被积函数中其余的因子换成在$\varrho=\bar\nu_{k\bar k}$处的值，即$I w_{k\bar k}(\bar\nu_{k\bar k})/\bar\nu_{k\bar k}^{2}$；这里$I$在两种情形中分别为$I(\bar\nu_{k\bar k})+\frac{8\pi h}{c^3}\bar\nu_{k\bar k}^{3}$及$I(\bar\nu_{k\bar k})$。于是
\[
 \vartheta_k=\frac{\mathcal V I w_{k\bar k}(\bar\nu_{k\bar k})}{4\pi^2c^2h^2\bar\nu_{k\bar k}^{2}}
 \int_0^{\infty}\frac{1-\cos2\pi(\varrho-\bar\nu_{k\bar k})t}{(\varrho-\bar\nu_{k\bar k})^2}\,d\varrho.
\]
同样，由于改变积分下限也只带来$t^2$的贡献，可以把$\int_0^{\infty}$换成$\int_{-\infty}^{\infty}$，再引入新变量$x=2\pi(\varrho-\bar\nu_{k\bar k})t$。由\origNote[由][]{152}
\[
 \int_{-\infty}^{\infty}\frac{1-\cos2\pi(\varrho-\bar\nu_{k\bar k})t}{(\varrho-\bar\nu_{k\bar k})^2}\,d\varrho
 =2\pi t\int_{-\infty}^{\infty}\frac{1-\cos x}{x^2}\,dx=2\pi^2t,
\]
最终得到\tnote[到][][原]{原书下一式略去了前面分母中的因子$c^2$。这里保留原式。}
\[
 \vartheta_k=\frac{\mathcal V I w_{k\bar k}(\bar\nu_{k\bar k})}{2h^2\bar\nu_{k\bar k}^{2}}\,t,
\]
这也就证明了$\vartheta_k$的数量级为$t$。

为了确定$w_{k\bar k}(\bar\nu_{k\bar k})$，必须把
\[
 \left|\sum_{\nu=1}^{l}\frac{e_\nu}{m_\nu}
 \bigl(P_\nu^x\overline{\mathfrak A}_{n,x}(Q_\nu^x,Q_\nu^y,Q_\nu^z)+\cdots\bigr)_{k\bar k}\right|^2
\]
改写成不含$\overline{\mathfrak A}_n$的表达式。考虑到$\overline{\mathfrak A}_n$变化很快，可以用一个长度恒定、方向随机的向量代替它。在这种近似中，向量不再依赖$Q_\nu^x,\allowbreak Q_\nu^y,\allowbreak Q_\nu^z$，其分量因而只相当于一个数因子乘以单位矩阵$1$。其恒定长度$\gamma_n$须从归一化条件
\[
 \iiint_{\boldsymbol H}[\overline{\mathfrak A}_n,\overline{\mathfrak A}_n]\,dx\,dy\,dz=4\pi c^2
\]
求得。因此
\[
 \mathcal V\gamma_n^2=4\pi c^2,\qquad \gamma_n^2=\frac{4\pi c^2}{\mathcal V}.
\]
平均而言，$x$分量$\overline{\mathfrak A}_{n,\allowbreak x}^{2}$占$[\overline{\mathfrak A}_n,\allowbreak \overline{\mathfrak A}_n]=\overline{\mathfrak A}_{n,\allowbreak x}^{2}+\overline{\mathfrak A}_{n,\allowbreak y}^{2}+\overline{\mathfrak A}_{n,\allowbreak z}^{2}=\gamma_n^2$的三分之一，所以等于$\frac13\gamma_n^2=\frac{4\pi c^2}{3\mathcal V}$；$\overline{\mathfrak A}_{n,\allowbreak y}^{2}$及$\overline{\mathfrak A}_{n,\allowbreak z}^{2}$也相同。因此
\begin{align*}
 w_{k\bar k}(\varrho)
 &=\mathop{\text{平均}}_{\varrho\le\varrho_n<\varrho+d\varrho}
 \left|\sum_{\nu=1}^{l}\frac{e_\nu}{m_\nu}
 \bigl(P_\nu^x\overline{\mathfrak A}_{n,x}(Q_\nu^x,Q_\nu^y,Q_\nu^z)+\cdots\bigr)_{k\bar k}\right|^2\\
 &\approx\frac{4\pi c^2}{3\mathcal V}
 \left(\left|\left(\sum_{\nu=1}^{l}\frac{e_\nu}{m_\nu}P_\nu^x\right)_{k\bar k}\right|^2+\cdots\right).
\end{align*}

\hypertarget{origpage-156}{}\phantomsection\label{origpage-156}单独考察系统$\boldsymbol{S}$时，其能量$\mathsf{H}_0$等于动能加势能，即具有形式
\[
 \mathsf{H}_0=\sum_{\nu=1}^{l}\frac1{2m_\nu}\bigl((P_\nu^x)^2+(P_\nu^y)^2+(P_\nu^z)^2\bigr)
 +V(Q_1^x,Q_1^y,Q_1^z,\ldots,Q_l^x,Q_l^y,Q_l^z),
\]
故有\origNote[有][]{153}
\[
 \mathsf{H}_0Q_\nu^x-Q_\nu^x\mathsf{H}_0=\frac h{2\pi i}\frac1{m_\nu}P_\nu^x.
\]
又因$\mathsf{H}_0$是对角元为$W_1,\allowbreak W_2,\allowbreak \ldots$的对角矩阵，即$(\mathsf{H}_0)_{kj}=W_k\delta_{kj}$，于是矩阵元满足
\begin{align*}
 (P_\nu^x)_{k\bar k}
 &=\frac{2\pi i m_\nu}{h}(\mathsf{H}_0Q_\nu^x-Q_\nu^x\mathsf{H}_0)_{k\bar k}
 =\frac{2\pi i m_\nu}{h}(W_k-W_{\bar k})(Q_\nu^x)_{k\bar k}\\
 &=\pm i\,2\pi m_\nu\bar\nu_{k\bar k}(Q_\nu^x)_{k\bar k}.
\end{align*}
所以\tnote[以][][前]{前面关于动量的恒等式中，$h$已约去，但原书下列$w_{k\bar k}(\varrho)$表达式的分母仍多保留了$h^2$。此处照录。}
\[
 w_{k\bar k}(\varrho)=\frac{16\pi^3c^2}{3\mathcal Vh^2}\bar\nu_{k\bar k}^{2}
 \left(\left|\left(\sum_{\nu=1}^{l}e_\nu Q_\nu^x\right)_{k\bar k}\right|^2+\cdots\right).
\]
代入$\vartheta_k$公式，得到
\[
 \vartheta_k=\frac{8\pi^3}{3h^2}
 \left(\left|\left(\sum_{\nu=1}^{l}e_\nu Q_\nu^x\right)_{k\bar k}\right|^2+\cdots\right)It.
\]
再置$w_{k\bar k}=\left|\left(\sum_{\nu=1}^{l}e_\nu Q_\nu^x\right)_{k\bar k}\right|^2+\cdots$，这个结果显然可解释如下。处于第$k$个状态的原子$\boldsymbol{S}$，会发生以下跃迁，即量子跃迁：

1．跃迁到较高的状态$\bar k$，$W_{\bar k}>W_k$，每秒次数为
\[
 \frac{8\pi^3}{3h^2}w_{k\bar k}I\left(\frac{W_{\bar k}-W_k}{h}\right)
 \]
即正比于辐射场在相应Bohr频率$(W_{\bar k}-W_k)/h$处的光谱能量密度。

2．跃迁到较低的状态$\bar k$，$W_{\bar k}<W_k$，每秒次数为
\[
 \frac{8\pi^3}{3h^2}w_{k\bar k}I\left(\frac{W_k-W_{\bar k}}{h}\right)
 \]
即正比于辐射场在相应Bohr频率$(W_k-W_{\bar k})/h$处的光谱能量密度。

3．此外，还会跃迁到较低的状态$\bar k$，$W_{\bar k}<W_k$，每秒次数为
\[
 \frac{64\pi^4}{3hc^3}w_{k\bar k}\left(\frac{W_k-W_{\bar k}}{h}\right)^3
 \]
这完全不依赖于已有的辐射场。

第1项对应从辐射场吸收的过程；第2项对应由辐射场引起的受激发射；第3项则是自发发射，只要原子尚未最终稳定在最低定态，即$W_k$最小的状态，它便总会发生。

\hypertarget{origpage-157}{}\phantomsection\label{origpage-157}早在量子力学发现以前，Einstein已经从热力学推得了这三种跃迁机制\origNote[制][，]{154}，只是还不知道“跃迁概率”$w_{k\bar k}$的值。上面得到的
\[
 w_{k\bar k}=\left|\left(\sum_{\nu=1}^{l}e_\nu Q_\nu^x\right)_{k\bar k}\right|^2
 +\left|\left(\sum_{\nu=1}^{l}e_\nu Q_\nu^y\right)_{k\bar k}\right|^2
 +\left|\left(\sum_{\nu=1}^{l}e_\nu Q_\nu^z\right)_{k\bar k}\right|^2
\]
正如前面说过的，已包含在Heisenberg最初的解释方案中。现在，我们按照Dirac的方法，已从一般理论为它建立了依据。


% 中文编写中间材料：chapters-IV-V.tex
\chapter{理论的演绎构建}\label{chap:IV}
\section{统计理论的原理依据}\label{sec:IV.1}

在第 \hyperref[chap:III]{III} 章中，我们已把量子力学的一切论断归结为一个统计公式（在那里记作 \hyperref[post:III.1.E2]{$E_2$}）：
\begin{equation}\tag{$\overline E$}\label{eq:IV.1.E}
\operatorname{Erw}(\mathfrak R,\varphi)=(R\varphi,\varphi).
\end{equation}
其中，$\operatorname{Erw}(\mathfrak R,\allowbreak \varphi)$是物理量$\mathfrak R$在状态$\varphi$中的期望值，$R$是与$\mathfrak R$对应的算符。下面要说明，这个公式本身也可以从少数几个一般性的定性假设推导出来；同时，还要重新检验第 \hyperref[chap:III]{III} 章以它为基础构建量子力学的整个过程。不过，在此之前还须补充一点。

在状态$\varphi$中，物理量$\mathfrak R$的期望值为$\varrho=(R\varphi,\allowbreak \varphi)$，其方差$\varepsilon^2$则是物理量$(\mathfrak R-\varrho)^2$的期望值，即
\[
((R-\varrho\mathbf1)^2\varphi,\varphi)=\|R\varphi\|^2-(R\varphi,\varphi)^2.
\]
参见原注 \hyperlink{note-130}{\getrefnumber{orig-note-130}}；这里的一切都由 \hyperref[eq:IV.1.E]{$\overline E$} 算得。这一方差通常大于零，只有在$R\varphi=\varrho\varphi$时才等于零，参见 \hyperref[sec:III.3]{III.3}。所以，正如我们多次指出的，即使状态$\varphi$是单一的，也只能作统计描述。如果连实际处于什么状态都不知道，统计性质还会进一步加强。例如，可能存在若干状态$\varphi_1,\allowbreak \varphi_2,\allowbreak \ldots$，它们各自的概率为$w_1,\allowbreak w_2,\allowbreak \ldots$，其中$w_1\geq0,\allowbreak w_2\geq0,\allowbreak \ldots$，且$w_1+w_2+\cdots=1$。按照普遍成立的概率法则，此时$\mathfrak R$的期望值为
\[
\varrho'=\sum_nw_n(R\varphi_n,\varphi_n).
\]
一般说来，$(R\varphi,\allowbreak \varphi)=\operatorname{Spur}(P_{[\varphi]}R)$。事实上，取一个完全的规范正交系$\psi_1,\allowbreak \psi_2,\allowbreak \ldots$，使$\psi_1=\varphi$，因而$\psi_2,\allowbreak \psi_3,\allowbreak \ldots$都与$\varphi$正交，就有
\[
P_{[\varphi]}\psi_n=\begin{cases}\varphi,&n=1,\\0,&\text{其他情形},\end{cases}
\]
于是
\begin{align*}
\operatorname{Spur}(P_{[\varphi]}R)
&=\sum_{m,n}(P_{[\varphi]}\psi_n,\psi_m)(R\psi_m,\psi_n)\\
&=\sum_n(\varphi,\psi_m)(R\psi_m,\varphi)=(R\varphi,\varphi).
\end{align*}

\hypertarget{origpage-158}{}\phantomsection\label{origpage-158}因此，上述$\varrho'$等于$\operatorname{Spur}(\{\sum_nw_nP_{[\varphi_n]}\}R)$。\tnote[。][算][原]{原书上一式第二行的求和下标印作$n$，但被加项中的下标为$m$。这里保留原印。}算符
\[
U=\sum_nw_nP_{[\varphi_n]}
\]
是定号的，因为所有$P_{[\varphi_n]}$都定号，且$w_n\geq0$；又因为$\operatorname{Spur}P_{[\varphi_n]}=1$，它的迹等于$\sum_nw_n=1$。此外，它完全确定了刚才所说的状态混合的统计性质：
\[
\varrho'=\operatorname{Spur}(UR).
\]
可见，除了状态之外，还须研究状态的混合。记住这一点，下面转入一般性的讨论。

现在暂且撇开整个量子力学，只保留下面的设定。给定一个系统$\boldsymbol{S}$；\origNote[；][对]{155}对实验者来说，只要列出在该系统上实际能够测量的所有物理量，以及这些量之间的函数关系，就确定了这个系统。严格地说，一个物理量就是一套测量规程：怎样测量，以及怎样根据测量仪器指针的位置读出或算出它的数值。若$\mathfrak R$是一个物理量，$f(x)$是任意函数，那么物理量$f(\mathfrak R)$定义如下：测量$f(\mathfrak R)$时，先测量$\mathfrak R$；如果后者的测得值为$a$，前者的值就是$f(a)$。可见，固定$\mathfrak R$而任取$f(x)$，这样得到的所有物理量$f(\mathfrak R)$，都能在测量$\mathfrak R$的同时一并测得。这就是同时可测物理量的第一个例子。

一般地说，如果存在一种装置，能够在同一系统上同时测量两个（或更多）物理量$\mathfrak R,\allowbreak \mathfrak S$，只是要以不同方法根据读数算出各自的值，我们就称这些量同时可测。众所周知，经典力学中的所有物理量都同时可测；量子力学则不然，参见 \hyperref[sec:III.3]{III.3}。对于同时可测的量，还可以由二元函数$f(x,\allowbreak y)$定义物理量$f(\mathfrak R,\allowbreak \mathfrak S)$：同时测量$\mathfrak R,\allowbreak \mathfrak S$，若测得$a,\allowbreak b$，则$f(\mathfrak R,\allowbreak \mathfrak S)$的值为$f(a,\allowbreak b)$。但须明确，若$\mathfrak R,\allowbreak \mathfrak S$不能同时测量，试图组成$f(\mathfrak R,\allowbreak \mathfrak S)$就是全无意义的，因为无法给出相应的测量装置。

然而，我们并不只能在单个对象$\boldsymbol{S}$上研究物理量。特别是，若几个量能否同时测量尚有疑问，自然还应考虑大的统计系综：它由许多个系统$\boldsymbol{S}_1,\allowbreak \ldots,\allowbreak \boldsymbol{S}_N$组成，也就是$\boldsymbol{S}$的$N$个副本，其中$N$很大\origNote[大][。]{156}。在这种系综$[\boldsymbol{S}_1,\allowbreak \ldots,\allowbreak \boldsymbol{S}_N]$上测量的，自然不是物理量$\mathfrak R$的某个“值”，而是它的数值分布。也就是说，对每一个区间$a'<a\leq a''$，\hypertarget{origpage-159}{}\phantomsection\label{origpage-159}其中$a',\allowbreak a''$给定且$a'\leq a''$，确定$\boldsymbol{S}_1,\allowbreak \ldots,\allowbreak \boldsymbol{S}_N$中有多少个系统的$\mathfrak R$值落在其中；把这个数除以$N$，就得到概率函数$w(a',\allowbreak a'')=w(a'')-w(a')$。\origNote[。][采]{157}采用这种系综的重大好处在于：

\leavevmode\phantomsection\textbf{1.}\label{advantage:IV.1.1}即使测量物理量$\mathfrak R$会使被测系统$\boldsymbol{S}$发生很大变化，只要$N$足够大，对系综$[\boldsymbol{S}_1,\allowbreak \ldots,\allowbreak \boldsymbol{S}_N]$中$\mathfrak R$的概率分布作统计测定，仍可使系综受到的改变任意小。量子力学中的测量确实会造成上述变化；而且，我们在 \hyperref[sec:III.4]{III.4} 已经看到，在基本过程的物理学中，这种变化原则上不可避免，因为测量的干预与系统本身、或与其中被观察的部分，属于同一数量级。

\leavevmode\phantomsection\textbf{2.}\label{advantage:IV.1.2}即使两个（或更多）物理量$\mathfrak R,\allowbreak \mathfrak S$在单个系统$\boldsymbol{S}$上不能同时测量，只要$N$足够大，仍能在同一个系综$[\boldsymbol{S}_1,\allowbreak \ldots,\allowbreak \boldsymbol{S}_N]$中同时确定它们的概率分布，而且可以任意精确。

对一个含$N$个元素的系综，要统计考察$\mathfrak R$的数值分布，并不必测量全部$N$个元素$\boldsymbol{S}_1,\allowbreak \ldots,\allowbreak \boldsymbol{S}_N$。任选一个含$M\leq N$个元素的子系综，例如$[\boldsymbol{S}_1,\allowbreak \ldots,\allowbreak \boldsymbol{S}_M]$，就已足够；条件是$M,\allowbreak N$都很大，但$M$相对于$N$仍可很小\origNote[小][。]{158}。这样，测量引起的改变至多只涉及整个系综的$M/N$；只要把$M/N$取得足够小，这个比例便可任意小。$N$足够大时，即使$M$很大也做得到，这就证明了 \hyperref[advantage:IV.1.1]{\textbf{1}}。若要同时测量两个（或更多）物理量$\mathfrak R,\allowbreak \mathfrak S$，就需要两个子系综，例如$[\boldsymbol{S}_1,\allowbreak \ldots,\allowbreak \boldsymbol{S}_M]$和$[\boldsymbol{S}_{M+1},\allowbreak \ldots,\allowbreak \boldsymbol{S}_{2M}]$，其中$2M\leq N$。前者用来统计$\mathfrak R$，后者用来统计$\mathfrak S$。这两次测量虽然针对同一个系综$[\boldsymbol{S}_1,\allowbreak \ldots,\allowbreak \boldsymbol{S}_N]$，却不会互相干扰；只要$2M/N$足够小，甚至只会改变其中任意小的一部分。$N$足够大时，即使$M$很大也能做到，这就证明了 \hyperref[advantage:IV.1.2]{\textbf{2}}。

可见，正因为测量可能影响单个系统，几个物理量也可能不能同时测量，我们才需要引入统计系综，也就是采用概率论的方法。一般的理论必须顾及这两种情况：基本过程中可能出现这些情况，本就是应当预料到的\origNote[的][；]{159}；而今天，精确的分析已使这一点不容置疑，参见 \hyperref[sec:III.4]{III.4}。统计系综消除了这些困难，\hypertarget{origpage-160}{}\phantomsection\label{origpage-160}从而使客观描述重新成为可能。这种描述既不依赖偶然的结果，也不依赖在给定状态下，究竟选择测量两个不能同时测量的量中的哪一个。

在这样的系综中，一个物理量$\mathfrak R$没有确定值，并不值得惊讶。它的数值分布可以不是只集中在单个值$a_0$上\origNote[上][，]{160}，而是允许几个值或若干数值区间，因而有正的方差\origNote[差][。]{160}。不过，这种情形可能有两种不同原因：

\leavevmode\phantomsection\textbf{I.}\label{case:IV.1.I}系综中的各个系统$\boldsymbol{S}_1,\allowbreak \ldots,\allowbreak \boldsymbol{S}_N$可能处于不同状态，而系综$[\boldsymbol{S}_1,\allowbreak \ldots,\allowbreak \boldsymbol{S}_N]$由这些状态的相对频率来规定。此时，物理量不能给出确定值，是由于我们有所不知：不知道测量的是哪一种状态，自然就无法说出将得到什么结果。

\leavevmode\phantomsection\textbf{II.}\label{case:IV.1.II}所有单个系统$\boldsymbol{S}_1,\allowbreak \ldots,\allowbreak \boldsymbol{S}_N$都处于同一状态，但自然规律不具有因果性。此时，方差并非来自我们的无知，而是自然本身不遵守“充足理由原则”。

情形 \hyperref[case:IV.1.I]{\textbf{I}} 人人熟悉，情形 \hyperref[case:IV.1.II]{\textbf{II}} 则既重要又新颖。起初，人们固然会怀疑后者是否可能存在，但我们将找到一个客观判据，用来判断它究竟是否出现。乍看之下，似乎可以提出重大异议，否认这种情形能够设想、或否认它有意义。然而，我们认为这些异议站不住脚；要摆脱某些困难（例如量子力学中的困难），\hyperref[case:IV.1.II]{\textbf{II}} 还是唯一的出路。因此，下面讨论 \hyperref[case:IV.1.II]{\textbf{II}} 在概念上遇到的困难。

对 \hyperref[case:IV.1.II]{\textbf{II}} 可以提出如下反驳：自然根本不可能违背“充足理由原则”，即因果性，因为这实际上更像是在定义“相同”。也就是说，两个相同的对象$\boldsymbol{S}_1,\allowbreak \boldsymbol{S}_2$——即系统$\boldsymbol{S}$的两个副本，处于同一状态——在一切可以设想的干预下都将作出相同反应；这一命题之所以为真，只是因为它没有说出任何实质内容。假如$\boldsymbol{S}_1,\allowbreak \boldsymbol{S}_2$受到相同干预却表现不同，例如测量物理量$\mathfrak R$时得到不同的值，人们就不会称它们“相同”。于是，若系综$[\boldsymbol{S}_1,\allowbreak \ldots,\allowbreak \boldsymbol{S}_N]$对某个物理量$\mathfrak R$有非零方差，那么按定义，单个系统$\boldsymbol{S}_1,\allowbreak \ldots,\allowbreak \boldsymbol{S}_N$就不可能全处于同一状态。把这个说法用于量子力学，就是说：只要$\varphi$不是$\mathfrak R$对应算符$R$的本征函数\origNote[数][，]{161}，在若干个都具有波函数$\varphi$的系统上测量同一个量$\mathfrak R$，\hypertarget{origpage-161}{}\phantomsection\label{origpage-161}就会得到不同的值。所以这些系统本来就不“相同”，也就是说，波函数给出的描述并不完备。因而还应存在其他决定因素，即 \hyperref[sec:III.2]{III.2} 提到的“隐参数”。我们很快就会看到，事情并没有这样简单。

按照这种看法，只要一个大的统计系综中还有某个量$\mathfrak R$的方差不为零，就必定能够按其元素所处的不同状态，把系综分成构成各异的几部分。这个说法似乎更有道理，因为看起来确实有一个简单的分法：按$\mathfrak R$在其中取得的不同值来划分。对所有存在的量$\mathfrak R,\allowbreak \mathfrak S,\allowbreak \mathfrak T,\allowbreak \ldots$都完成这样的细分或分解之后，才应得到真正的纯系综。最后得到的每一个子系综中，这些量的方差就都为零了。

首先，刚才这几句话中的断言是错误的，因为它们忽略了测量会改变被测系统。为简单起见，设$\mathfrak R$只取$a_1,\allowbreak a_2$两个值；在全部对象$\boldsymbol{S}_1,\allowbreak \ldots,\allowbreak \boldsymbol{S}_N$上测量它，设在$\boldsymbol{S}'_1,\allowbreak \ldots,\allowbreak \boldsymbol{S}'_{N_1}$上得到$a_1$，在$\boldsymbol{S}''_1,\allowbreak \ldots,\allowbreak \boldsymbol{S}''_{N-N_1}$上得到$a_2$。测量之后，$\mathfrak R$在$[\boldsymbol{S}'_1,\allowbreak \ldots,\allowbreak \boldsymbol{S}'_{N_1}]$和$[\boldsymbol{S}''_1,\allowbreak \ldots,\allowbreak \boldsymbol{S}''_{N-N_1}]$中的方差确实都为零，在前者中恒为$a_1$，在后者中恒为$a_2$。但这不只是把$[\boldsymbol{S}_1,\allowbreak \ldots,\allowbreak \boldsymbol{S}_N]$分成了上述两部分，因为测量$\mathfrak R$还改变了各个系统。

我们在 \hyperref[advantage:IV.1.1]{\textbf{1}} 中确实找到一种办法，能在只略微改变$[\boldsymbol{S}_1,\allowbreak \ldots,\allowbreak \boldsymbol{S}_N]$的情况下确定$\mathfrak R$的数值分布：只测量$\boldsymbol{S}_1,\allowbreak \ldots,\allowbreak \boldsymbol{S}_M$，其中$M$很大，而$M/N$很小。但这不能给出分解，因为对大多数系统$\boldsymbol{S}_1,\allowbreak \ldots,\allowbreak \boldsymbol{S}_N$，也就是$\boldsymbol{S}_{M+1},\allowbreak \ldots,\allowbreak \boldsymbol{S}_N$，并没有查明每个系统各自的$\mathfrak R$值。现在可以看出，上述产生纯系综的方法为什么失效。接着在$[\boldsymbol{S}'_1,\allowbreak \ldots,\allowbreak \boldsymbol{S}'_{N_1}]$和$[\boldsymbol{S}''_1,\allowbreak \ldots,\allowbreak \boldsymbol{S}''_{N-N_1}]$中测量第二个量$\mathfrak S$；它也只取$b_1,\allowbreak b_2$两个值。设在$\boldsymbol{S}'''_1,\allowbreak \ldots,\allowbreak \boldsymbol{S}'''_{N_{11}}$与$\boldsymbol{S}^{\mathrm V}_1,\allowbreak \ldots,\allowbreak \boldsymbol{S}^{\mathrm V}_{N_{12}}$中测得$b_1$，在$\boldsymbol{S}^{\mathrm{IV}}_1,\allowbreak \ldots,\allowbreak \boldsymbol{S}^{\mathrm{IV}}_{N_1-N_{11}}$与$\boldsymbol{S}^{\mathrm{VI}}_1,\allowbreak \ldots,\allowbreak \boldsymbol{S}^{\mathrm{VI}}_{N-N_1-N_{12}}$中测得$b_2$。于是$\mathfrak S$在
\begin{gather*}\relax
[\boldsymbol{S}'''_1,\ldots,\boldsymbol{S}'''_{N_{11}}],\quad
[\boldsymbol{S}^{\mathrm{IV}}_1,\ldots,\boldsymbol{S}^{\mathrm{IV}}_{N_1-N_{11}}],\\
[\boldsymbol{S}^{\mathrm V}_1,\ldots,\boldsymbol{S}^{\mathrm V}_{N_{12}}],\quad
[\boldsymbol{S}^{\mathrm{VI}}_1,\ldots,\boldsymbol{S}^{\mathrm{VI}}_{N-N_1-N_{12}}]
\end{gather*}
中的方差都为零，它的值依次恒为$b_1,\allowbreak b_2,\allowbreak b_1,\allowbreak b_2$。然而，前两个系综虽然是$[\boldsymbol{S}'_1,\allowbreak \ldots,\allowbreak \boldsymbol{S}'_{N_1}]$的部分，后两个虽然是$[\boldsymbol{S}''_1,\allowbreak \ldots,\allowbreak \boldsymbol{S}''_{N-N_1}]$的部分，而$\mathfrak R$在这两个原系综中的方差都为零，但在新得到的四个系综中，方差却可能个个都不为零：测量$\mathfrak S$改变了组成它们的单个系统！因此，这样做并没有进展，每一步都破坏了前一步的结果\origNote[果][。]{162}。无论怎样累积接连进行的\hypertarget{origpage-162}{}\phantomsection\label{origpage-162}测量，也不能在这团混乱中建立因果秩序。在原子领域，我们已经来到物理世界的边缘；每一次测量都是与被测对象同一数量级的干预，因而会显著影响它。这实质上是不确定性关系的结果。

所以，未必存在一种方法，总能在不改变系综各个元素的条件下，把存在非零方差的系综继续分解；更不能保证最终达到所有物理量的方差都为零的纯系综。我们习惯于把这种纯系综看成由彼此相同、按因果律确定的单个对象组成。尽管如此，人们仍可试着保留如下设想：任何存在非零方差的系综，都能在不改变其元素的条件下分成两个（或更多）部分，各部分彼此不同，也不同于原系综；也就是说，将两个分解所得的系综混合，便重新得到原来的系综。

可见，原先想把因果性作为“相同”的定义来确立，最后却变成了一个事实问题。这个问题能够而且必须回答，答案甚至可能是否定的：凡是其中某些物理量$\mathfrak R$的方差不为零的系综$[\boldsymbol{S}_1,\allowbreak \ldots,\allowbreak \boldsymbol{S}_N]$，是否真的都能表示成两个（或更多）系综的混合，而这些系综彼此不同，也不同于它？若多于两个，例如$n=3,\allowbreak 4,\allowbreak \ldots$，只须把第一个与其余$n-1$个的混合态分开考虑，就可归结为两个的情形。

例如，若$[\boldsymbol{S}_1,\allowbreak \ldots,\allowbreak \boldsymbol{S}_N]$是$[\boldsymbol{S}'_1,\allowbreak \ldots,\allowbreak \boldsymbol{S}'_P]$与$[\boldsymbol{S}''_1,\allowbreak \ldots,\allowbreak \boldsymbol{S}''_Q]$的混合，那么对于每个物理量$\mathfrak R$，它的概率函数$w_{\mathfrak R}(a)$（参见原注 \hyperlink{note-157}{\getrefnumber{orig-note-157}}）都可以用后两个系综的概率函数$w'_{\mathfrak R}(a),\allowbreak w''_{\mathfrak R}(a)$表示为
\begin{equation}\tag{$M_1$}\label{eq:IV.1.M1}
w_{\mathfrak R}(a)=\alpha w'_{\mathfrak R}(a)+\beta w''_{\mathfrak R}(a),\qquad
\alpha>0,\quad\beta>0,\quad\alpha+\beta=1.
\end{equation}
其中$\alpha=P/N$，$\beta=Q/N$，$N=P+Q$，而$\alpha,\allowbreak \beta$不依赖$\mathfrak R$。因此，归根结底，这是一个数学问题：给定一个概率函数为$w_{\mathfrak R}(a)$的系综，其中有些量$\mathfrak R$的方差不为零（$w_{\mathfrak R}(a)$的相应性质见原注 \hyperlink{note-160}{\getrefnumber{orig-note-160}}），是否总能找到另外两个系综，其概率函数为$w'_{\mathfrak R}(a)$与$w''_{\mathfrak R}(a)$，使 \hyperref[eq:IV.1.M1]{$M_1$} 对一切$\mathfrak R$成立？也可以换一种表述：不用各个量$\mathfrak R$的概率函数$w_{\mathfrak R}(a)$，而用它们的期望值
\[
\operatorname{Erw}(\mathfrak R)=\int_{-\infty}^{+\infty}a\,dw_{\mathfrak R}(a)
\]
来确定系综。这样，问题可表述如下：如果一个系综\hypertarget{origpage-163}{}\phantomsection\label{origpage-163}对每个$\mathfrak R$的方差
\[
\operatorname{Erw}\bigl([\mathfrak R-\operatorname{Erw}(\mathfrak R)]^2\bigr)
=\operatorname{Erw}(\mathfrak R^2)-[\operatorname{Erw}(\mathfrak R)]^2
\]
都等于零（参见原注 \hyperlink{note-160}{\getrefnumber{orig-note-160}}），即满足
\begin{equation}\tag{$\mathrm{Str}_1$}\label{eq:IV.1.Str1}
\operatorname{Erw}(\mathfrak R^2)=[\operatorname{Erw}(\mathfrak R)]^2
\end{equation}
就称它为零方差系综。如果不是这样，是否总能找到另外两个系综，其期望值分别为$\operatorname{Erw}'(\mathfrak R)$与$\operatorname{Erw}''(\mathfrak R)$，且
\[
\operatorname{Erw}(\mathfrak R)\not\equiv\operatorname{Erw}'(\mathfrak R)
\not\equiv\operatorname{Erw}''(\mathfrak R),
\]
使
\begin{equation}\tag{$M_2$}\label{eq:IV.1.M2}
\begin{gathered}
\operatorname{Erw}(\mathfrak R)=\alpha\operatorname{Erw}'(\mathfrak R)+\beta\operatorname{Erw}''(\mathfrak R),\\
\alpha>0,\quad\beta>0,\quad\alpha+\beta=1
\end{gathered}
\end{equation}
恒成立，其中$\alpha,\allowbreak \beta$不依赖$\mathfrak R$？［须注意：对于一个给定的量$\mathfrak R$，数$\operatorname{Erw}(\mathfrak R)$不能代替函数$w_{\mathfrak R}(a)$；但知道所有$\operatorname{Erw}(\mathfrak R)$，就等于知道所有$w_{\mathfrak R}(a)$。因为若以
\[
f_a(x)=\begin{cases}1,&x\leq a,\\0,&x>a,\end{cases}
\]
定义$f_a(x)$，就有$w_{\mathfrak R}(a)=\operatorname{Erw}(f_a(\mathfrak R))$。］

要从数学上处理这个问题，考察相应的$\operatorname{Erw}(\mathfrak R)$，比考察系综$[\boldsymbol{S}_1,\allowbreak \ldots,\allowbreak \boldsymbol{S}_N]$本身更方便。每个系综都有这样一个函数：它对$\boldsymbol{S}$中的所有物理量$\mathfrak R$都有定义，取值为实数；反过来，它也完全确定系综的一切统计性质。这里可参照上面对$\operatorname{Erw}(\mathfrak R)$与$w_{\mathfrak R}(a)$之间关系的说明。当然，我们还须弄清，一个以$\mathfrak R$为自变量的函数必须具备哪些性质，才能成为某个适当系综的$\operatorname{Erw}(\mathfrak R)$。一旦弄清，就可以作出定义：

\noindent\phantomsection$\alpha)$\label{def:IV.1.alpha}若一个以$\mathfrak R$为自变量的函数是某个$\operatorname{Erw}(\mathfrak R)$，且满足条件 \hyperref[eq:IV.1.Str1]{$\mathrm{Str}_1$}，就称它\textbf{零方差}。

\noindent\phantomsection$\beta)$\label{def:IV.1.beta}若一个以$\mathfrak R$为自变量的函数是某个$\operatorname{Erw}(\mathfrak R)$，并且对它而言，\hyperref[eq:IV.1.M2]{$M_2$} 蕴含
\[
\operatorname{Erw}(\mathfrak R)\equiv\operatorname{Erw}'(\mathfrak R)
\equiv\operatorname{Erw}''(\mathfrak R)
\]
就称它\textbf{纯}的（或\textbf{均一}的）。从这些概念的含义看，每一个零方差的$\operatorname{Erw}(\mathfrak R)$函数都是纯的；这一点我们稍后还要证明。因此，我们的问题是：每一个纯的$\operatorname{Erw}(\mathfrak R)$函数是否都具有零方差性质？

显然，每一个$\operatorname{Erw}(\mathfrak R)$函数都必须具备以下性质：

\noindent\phantomsection\textbf{A.}\label{cond:IV.1.A}若物理量$\mathfrak R$恒等于$1$，则$\operatorname{Erw}(\mathfrak R)=1$。这里所谓恒等于$1$，是指$\mathfrak R$的“测量规程”为：什么都不必测量，因为$\mathfrak R$永远取值$1$。

\noindent\phantomsection\textbf{B.}\label{cond:IV.1.B}对每个$\mathfrak R$和每个实数$a$，都有$\operatorname{Erw}(a\mathfrak R)=a\operatorname{Erw}(\mathfrak R)$。\origNote[。][]{163}

\hypertarget{origpage-164}{}\phantomsection\label{origpage-164}\noindent\phantomsection\textbf{C.}\label{cond:IV.1.C}若物理量$\mathfrak R$按其本性绝不取负值，例如它是另一个物理量$\mathfrak S$的平方\origNote[方][，]{163}，则$\operatorname{Erw}(\mathfrak R)\geq0$。

\noindent\phantomsection\textbf{D.}\label{cond:IV.1.D}若物理量$\mathfrak R,\allowbreak \mathfrak S,\allowbreak \ldots$同时可测，则
\[
\operatorname{Erw}(\mathfrak R+\mathfrak S+\cdots)=\operatorname{Erw}(\mathfrak R)+\operatorname{Erw}(\mathfrak S)+\cdots.
\quad\text{\textsuperscript{\hyperlink{note-163}{\getrefnumber{orig-note-163}}}}
\]
若$\mathfrak R,\allowbreak \mathfrak S,\allowbreak \ldots$不能同时测量，$\mathfrak R+\mathfrak S+\cdots$便没有定义，参见上文。

这一切都直接来自所考察物理量的定义（即测量规程），以及期望值的定义：在一个足够大的统计系综中，把所有测量结果取算术平均。关于 \hyperref[cond:IV.1.D]{D}，应注意它所依据的是概率论中的一条定理：一个和的期望值总等于各个加项的期望值之和，不论这些加项之间是否存在概率上的依赖关系；例如乘积就没有这样的性质。我们只对同时可测的$\mathfrak R,\allowbreak \mathfrak S,\allowbreak \ldots$提出这一条件，自然是因为否则$\mathfrak R+\mathfrak S+\cdots$便没有意义。

不过，量子力学中还存在一种超出上述讨论的运算：任意两个物理量都可以相加，不要求它们同时可观测。这是因为，两个厄米算符$R,\allowbreak S$的和$R+S$仍是厄米的，即使$R,\allowbreak S$不对易也是如此；例如，它们的乘积$RS$就只有在二者对易时才是厄米的，参见 \hyperref[sec:II.5]{II.5}。在每个状态$\varphi$中，期望值相加：
\[
(R\varphi,\varphi)+(S\varphi,\varphi)=((R+S)\varphi,\varphi);
\]
参见 \hyperref[sec:III.1]{III.1} 的 \hyperref[post:III.1.E2]{$E_2$}。多个加项也一样。我们把这一事实纳入目前的一般性设定；这个设定到现在还完全没有专门限定为量子力学：

\noindent\phantomsection\textbf{E.}\label{cond:IV.1.E}若$\mathfrak R,\allowbreak \mathfrak S,\allowbreak \ldots$是任意物理量，就存在另一个物理量$\mathfrak R+\mathfrak S+\cdots$，它不依赖于$\operatorname{Erw}(\mathfrak R)$函数的选取，而且满足
\[
\operatorname{Erw}(\mathfrak R+\mathfrak S+\cdots)=\operatorname{Erw}(\mathfrak R)+\operatorname{Erw}(\mathfrak S)+\cdots.
\]
若$\mathfrak R,\allowbreak \mathfrak S,\allowbreak \ldots$同时可测，根据 \hyperref[cond:IV.1.D]{D}，这个$\mathfrak R+\mathfrak S+\cdots$必定就是通常的和。但在一般情形中，它只由 \hyperref[cond:IV.1.E]{E} 间接规定；我们很难把$\mathfrak R,\allowbreak \mathfrak S,\allowbreak \ldots$的测量规程合成为$\mathfrak R+\mathfrak S+\cdots$的测量规程。\origNote[。][]{164}

在以上说明之外，再补充一点。我们不只容许代表期望值的$\operatorname{Erw}(\mathfrak R)$函数，也要容许代表相对期望值的函数；也就是说，撤去归一化条件 \hyperref[cond:IV.1.A]{A}。根据 \hyperref[cond:IV.1.C]{C}，$\operatorname{Erw}(1)\geq0$；如果它有限且不为零，便没有本质区别，因为改用$\operatorname{Erw}(\mathfrak R)/\operatorname{Erw}(1)$，一切就与原来相同。但$\operatorname{Erw}(1)=\infty$是一种本质上不同的可能性。我们作这种\hypertarget{origpage-165}{}\phantomsection\label{origpage-165}推广，正是为了包括它。用一个简单例子最容易说明：有时采用相对概率比采用真正的概率更好，尤其是相对概率的总和为无穷大的情形；$\operatorname{Erw}(1)$对应的正是总概率。设所考察的系统是一个能沿一维运动的粒子，其统计分布使它在一条无限直线上各处出现的可能性相等。这时，直线上每个有限区间的概率都是零。但各处等概率并不由这一点体现，而是由如下事实体现：两个有限区间的概率之比，等于它们长度之比。由于$0/0$没有意义，我们只能把区间长度引入为相对概率，借此表达上述关系；不过，总相对概率也就成为$\infty$。

综合以上各点，最后把条件写成如下形式（\hyperref[cond:IV.Aprime]{A$'$} 对应 \hyperref[cond:IV.1.C]{C}，\hyperref[cond:IV.Bprime]{B$'$} 对应 \hyperref[cond:IV.1.B]{B}、\hyperref[cond:IV.1.D]{D}、\hyperref[cond:IV.1.E]{E}）：

\noindent\phantomsection\textbf{A$'$.}\label{cond:IV.Aprime}若物理量$\mathfrak R$按其本性绝不取负值，例如它是另一个物理量$\mathfrak S$的平方，则$\operatorname{Erw}(\mathfrak R)\geq0$。

\noindent\phantomsection\textbf{B$'$.}\label{cond:IV.Bprime}若$\mathfrak R,\allowbreak \mathfrak S,\allowbreak \ldots$是任意物理量，$a,\allowbreak b,\allowbreak \ldots$是实数，则
\[
\operatorname{Erw}(a\mathfrak R+b\mathfrak S+\cdots)=a\operatorname{Erw}(\mathfrak R)+b\operatorname{Erw}(\mathfrak S)+\cdots.
\]
还须强调：

\noindent\phantomsection1.\label{remark:IV.1.1}由于只考察相对期望值，函数$\operatorname{Erw}(\mathfrak R)$与$c\operatorname{Erw}(\mathfrak R)$（$c>0$为常数）应视为没有本质区别。

\noindent\phantomsection2.\label{remark:IV.1.2}对所有$\mathfrak R$都有$\operatorname{Erw}(\mathfrak R)\equiv0$时，不能给出任何论断，故排除这一函数。

\noindent\phantomsection3.\label{remark:IV.1.3}当$\operatorname{Erw}(1)=1$时，得到的是绝对期望值，即正确归一化的期望值。根据 \hyperref[cond:IV.Aprime]{A$'$}，$\operatorname{Erw}(1)$总是$\geq0$；只要它有限且不为零，便可在 \hyperref[remark:IV.1.1]{1} 中取$c=1/\operatorname{Erw}(1)$，恢复正确的归一化。若$\operatorname{Erw}(1)=0$，则如下文所证，属于 \hyperref[remark:IV.1.2]{2} 的情形，应予排除。若$\operatorname{Erw}(1)=\infty$，则统计本质上未归一化，即只能是相对的。

还须回到定义 \hyperref[def:IV.1.alpha]{$\alpha)$}、\hyperref[def:IV.1.beta]{$\beta)$}。根据 \hyperref[remark:IV.1.1]{1}，可把 \hyperref[eq:IV.1.M2]{$M_2$} 换成较简单的条件
\begin{equation}\tag{$M_3$}\label{eq:IV.1.M3}
\operatorname{Erw}(\mathfrak R)=\operatorname{Erw}'(\mathfrak R)+\operatorname{Erw}''(\mathfrak R).
\end{equation}
关于 \hyperref[eq:IV.1.Str1]{$\mathrm{Str}_1$}，应注意先前的计算假设了$\operatorname{Erw}(1)=1$。当$\operatorname{Erw}(1)=\infty$时，根本无法定义零方差性，因为它所说的是$\operatorname{Erw}((\mathfrak R-\varrho)^2)=0$，其中$\varrho$为$\mathfrak R$的绝对\hypertarget{origpage-166}{}\phantomsection\label{origpage-166}期望值，也就是$\operatorname{Erw}(\mathfrak R)/\operatorname{Erw}(1)$，而这个比值在此可能是$\infty/\infty$，从而没有意义。\origNote[。][因]{165}因此，\hyperref[def:IV.1.alpha]{$\alpha)$}、\hyperref[def:IV.1.beta]{$\beta)$} 现在改述如下：

\noindent\phantomsection$\alpha')$\label{def:IV.1.alphaprime}一个以$\mathfrak R$为自变量、可作为$\operatorname{Erw}(\mathfrak R)$的函数，若$\operatorname{Erw}(1)$有限且不为零，就可以依 \hyperref[remark:IV.1.1]{1} 设$\operatorname{Erw}(1)=1$；在此条件下，以 \hyperref[eq:IV.1.Str1]{$\mathrm{Str}_1$} 为判据，满足它的称为\textbf{零方差}。

\noindent\phantomsection$\beta')$\label{def:IV.1.betaprime}一个以$\mathfrak R$为自变量、可作为$\operatorname{Erw}(\mathfrak R)$的函数，若对它而言，\hyperref[eq:IV.1.M3]{$M_3$} 必定蕴含
\[
\operatorname{Erw}'(\mathfrak R)=c'\operatorname{Erw}(\mathfrak R),\qquad
\operatorname{Erw}''(\mathfrak R)=c''\operatorname{Erw}(\mathfrak R),
\]
就称它\textbf{纯}的（或\textbf{均一}的）。其中$c',\allowbreak c''$是常数，当然有$c'+c''=1$，而且根据 \hyperref[cond:IV.Aprime]{A$'$} 及 \hyperref[remark:IV.1.1]{1}、\hyperref[remark:IV.1.2]{2}，有$c'>0$、$c''>0$。

根据 \hyperref[cond:IV.Aprime]{A$'$}、\hyperref[cond:IV.Bprime]{B$'$} 及 \hyperref[def:IV.1.alphaprime]{$\alpha')$}、\hyperref[def:IV.1.betaprime]{$\beta')$}，只要知道$\boldsymbol{S}$中的物理量$\mathfrak R,\allowbreak \mathfrak S,\allowbreak \ldots$以及它们之间的函数关系，就能对因果性问题作出确切判断。下一节将在量子力学中进行这一工作。

本节最后再补充两点。

第一点涉及$\operatorname{Erw}(1)=0$的情形。此时，由 \hyperref[cond:IV.Bprime]{B$'$} 得$\operatorname{Erw}(c)=0$。如果一个物理量$\mathfrak R$恒满足$\mathfrak R\geq c'$、$\mathfrak R\leq c''$，由 \hyperref[cond:IV.Aprime]{A$'$} 得$\operatorname{Erw}(c''-\mathfrak R)\geq0$、$\operatorname{Erw}(\mathfrak R-c')\geq0$；再由 \hyperref[cond:IV.Bprime]{B$'$} 得$\operatorname{Erw}(c')\leq\operatorname{Erw}(\mathfrak R)\leq\operatorname{Erw}(c'')$，即$\operatorname{Erw}(\mathfrak R)=0$。现在任取$\mathfrak R$，并取一列有界函数$f_1(x),\allowbreak f_2(x),\allowbreak \ldots$，使
\[
f_1(x)+f_2(x)+\cdots=x
\]
例如
\[
f_1(x)=\frac{\sin x}{x},\qquad
f_n(x)=\frac{\sin nx}{nx}-\frac{\sin(n-1)x}{(n-1)x},\quad n=2,3,\ldots.
\]
于是对$n=1,\allowbreak 2,\allowbreak \ldots$有$\operatorname{Erw}(f_n(\mathfrak R))=0$，再由 \hyperref[cond:IV.Bprime]{B$'$} 得$\operatorname{Erw}(\mathfrak R)=0$。所以，正如前面所说，$\operatorname{Erw}(1)=0$属于 \hyperref[remark:IV.1.2]{2} 所排除的情形。\tnote[。][][原]{原书所印的例子不满足所述求和关系：其前$N$项和为$\sin(Nx)/(Nx)$，当$x\ne0$时趋于$0$，并非$x$。公式仍按原书保留。}

第二点是：\hyperref[eq:IV.1.Str1]{$\mathrm{Str}_1$} 用$\operatorname{Erw}(\mathfrak R^2)=[\operatorname{Erw}(\mathfrak R)]^2$来刻画零方差性，颇值得注意。因为在零方差的情形，$\operatorname{Erw}(\mathfrak R)$就是$\mathfrak R$的值，$\operatorname{Erw}(f(\mathfrak R))$就是$f(\mathfrak R)$的值，所以对每个函数$f(x)$都应有
\begin{equation}\tag{$\mathrm{Str}_2$}\label{eq:IV.1.Str2}
\operatorname{Erw}(f(\mathfrak R))=f(\operatorname{Erw}(\mathfrak R)).
\end{equation}
\hyperref[eq:IV.1.Str1]{$\mathrm{Str}_1$} 只是 \hyperref[eq:IV.1.Str2]{$\mathrm{Str}_2$} 中$f(x)=x^2$这一特例，为什么就够了呢？答案是：只要 \hyperref[eq:IV.1.Str2]{$\mathrm{Str}_2$} 对$f(x)=x^2$成立，它便自动对所有$f(x)$成立。甚至可以用$x$的任意其他连续的严格凸函数来代替$x^2$，也就是满足下式的函数：
\[
f\left(\frac{x+y}{2}\right)<\frac{f(x)+f(y)}2\qquad(x\ne y).
\]
这里不再证明。

\hypertarget{origpage-167}{}\phantomsection\label{origpage-167}\section{统计公式的证明}\label{sec:IV.2}
我们知道，量子力学系统中的每个物理量，都唯一对应一个超极大厄米算符，参见 \hyperref[sec:III.5]{III.5} 中的讨论。为便于讨论，我们假定这种对应是一一对应，即每一个超极大厄米算符确实都对应一个物理量；在 \hyperref[sec:III.3]{III.3} 中，我们也曾偶尔用到这一点。此外，还有如下规律成立（参见 \hyperref[sec:III.5]{III.5} 中的 \hyperref[post:III.5.F]{F}、\hyperref[post:III.5.L]{L} 及 \hyperref[sec:IV.1]{IV.1} 末尾的说明）：

\noindent\phantomsection\textbf{I.}\label{cond:IV.2.I}若物理量$\mathfrak R$对应算符$R$，则物理量$f(\mathfrak R)$对应算符$f(R)$。

\noindent\phantomsection\textbf{II.}\label{cond:IV.2.II}若物理量$\mathfrak R,\allowbreak \mathfrak S,\allowbreak \ldots$分别对应算符$R,\allowbreak S,\allowbreak \ldots$，则物理量$\mathfrak R+\mathfrak S+\cdots$对应算符$R+S+\cdots$。这里不假设$\mathfrak R,\allowbreak \mathfrak S,\allowbreak \ldots$同时可测，参见前述位置的说明。

\hyperref[cond:IV.Aprime]{A$'$}、\hyperref[cond:IV.Bprime]{B$'$}、\hyperref[def:IV.1.alphaprime]{$\alpha')$}、\hyperref[def:IV.1.betaprime]{$\beta')$} 以及 \hyperref[cond:IV.2.I]{I}、\hyperref[cond:IV.2.II]{II}，构成了下面分析的数学出发点。

设$\varphi_1,\allowbreak \varphi_2,\allowbreak \ldots$是一个完全的规范正交系。对于每个算符$R$，我们考察矩阵$a_{\mu\nu}=(R\varphi_\mu,\allowbreak \varphi_\nu)$，以它代替$R$。再构造如下厄米算符，其矩阵分别为
\[
e_{\mu\nu}^{(n)}=\begin{cases}1,&\mu=\nu=n,\\0,&\text{其他情形},\end{cases}
\]
\[
f_{\mu\nu}^{(mn)}=\begin{cases}1,&\mu=m,\ \nu=n,\\1,&\mu=n,\ \nu=m,\\0,&\text{其他情形},\end{cases}
\qquad
g_{\mu\nu}^{(mn)}=\begin{cases}i,&\mu=m,\ \nu=n,\\-i,&\mu=n,\ \nu=m,\\0,&\text{其他情形}.\end{cases}
\]
顺便指出，这些算符就是
\[
P_{[\varphi_n]},\qquad
P_{[(\varphi_m+\varphi_n)/\sqrt2]}-P_{[(\varphi_m-\varphi_n)/\sqrt2]},\qquad
P_{[(\varphi_m+i\varphi_n)/\sqrt2]}-P_{[(\varphi_m-i\varphi_n)/\sqrt2]}.
\]
相应的物理量记作$\mathfrak U^{(n)},\allowbreak \mathfrak V^{(mn)},\allowbreak \mathfrak W^{(mn)}$。由$a_{nm}=\overline{a}_{mn}$，显然有
\[
a_{\mu\nu}=\sum_n a_{nn}e_{\mu\nu}^{(n)}
+\sum_{m<n}\operatorname{Re}a_{mn}f_{\mu\nu}^{(mn)}
+\sum_{m<n}\operatorname{Im}a_{mn}g_{\mu\nu}^{(mn)},
\]
以及
\[
\mathfrak R=\sum_n a_{nn}\mathfrak U^{(n)}
+\sum_{m<n}\operatorname{Re}a_{mn}\mathfrak V^{(mn)}
+\sum_{m<n}\operatorname{Im}a_{mn}\mathfrak W^{(mn)}.
\]
根据 \hyperref[cond:IV.2.II]{II} 及 \hyperref[cond:IV.Bprime]{B$'$}，
\begin{align*}
\operatorname{Erw}(\mathfrak R)={}&\sum_n a_{nn}\operatorname{Erw}(\mathfrak U^{(n)})
+\sum_{m<n}\operatorname{Re}a_{mn}\operatorname{Erw}(\mathfrak V^{(mn)})\\
&+\sum_{m<n}\operatorname{Im}a_{mn}\operatorname{Erw}(\mathfrak W^{(mn)}).
\end{align*}\hypertarget{origpage-168}{}\phantomsection\label{origpage-168}所以，令
\begin{alignat*}{2}
u_{nn}&=\operatorname{Erw}(\mathfrak U^{(n)}),\\
u_{mn}&=\tfrac12\operatorname{Erw}(\mathfrak V^{(mn)})+\tfrac i2\operatorname{Erw}(\mathfrak W^{(mn)}),&\quad &m<n,\\
u_{nm}&=\tfrac12\operatorname{Erw}(\mathfrak V^{(mn)})-\tfrac i2\operatorname{Erw}(\mathfrak W^{(mn)}),&\quad &m<n,
\end{alignat*}
便得
\[
\operatorname{Erw}(\mathfrak R)=\sum_{m,n}u_{nm}a_{mn}.
\]
由于$u_{nm}=\overline{u}_{mn}$，可以用$(U\varphi_m,\allowbreak \varphi_n)=u_{mn}$定义一个厄米算符$U$；\origNote[；][而]{166}而右端就是$\operatorname{Spur}(UR)$，参见 \hyperref[sec:II.11]{II.11}。因此，最后的公式为
\begin{equation}\tag{Sp.}\label{eq:IV.2.Sp}
\operatorname{Erw}(\mathfrak R)=\operatorname{Spur}(UR),
\end{equation}
$U$是一个不依赖于$R$的厄米算符\origNote[符][，]{167}，因而由系综本身决定。

根据 \hyperref[cond:IV.2.II]{II}，无论怎样选取$U$，\hyperref[eq:IV.2.Sp]{Sp.} 都满足 \hyperref[cond:IV.Bprime]{B$'$}。所以只须确定 \hyperref[cond:IV.Aprime]{A$'$} 对$U$施加了什么限制。任取$\|\varphi\|=1$的$\varphi$。由$P_{[\varphi]}^2=P_{[\varphi]}$及 \hyperref[cond:IV.2.I]{I}，$P_{[\varphi]}$对应的物理量$\mathfrak R$满足$\mathfrak R^2=\mathfrak R$，故由 \hyperref[cond:IV.Aprime]{A$'$} 得$\operatorname{Erw}(\mathfrak R)\geq0$。于是
\[
\operatorname{Spur}(UP_{[\varphi]})=(U\varphi,\varphi)\geq0.
\]
对于任意$f\ne0$，可取$\varphi=f/\|f\|$，于是$(U\varphi,\allowbreak \varphi)=(Uf,\allowbreak f)/\|f\|^2$，仍有$(Uf,\allowbreak f)\geq0$；而$f=0$时本来就成立。所以$U$定号。这样，\hyperref[cond:IV.Aprime]{A$'$} 蕴含$U$的定号性；反过来，定号性也足以保证 \hyperref[cond:IV.Aprime]{A$'$} 成立。

\hyperref[cond:IV.Aprime]{A$'$} 所要求的，无非是每个$\operatorname{Erw}(\mathfrak S^2)$都$\geq0$。因为，若$\mathfrak R$只能取非负值，对于$f(x)=|x|$就有$f(\mathfrak R)=\mathfrak R$；又因为$g(x)=\sqrt{|x|}$恒满足$(g(x))^2=f(x)$，便有$(g(\mathfrak R))^2=f(\mathfrak R)$，所以$\mathfrak R=\mathfrak S^2$，其中$\mathfrak S=g(\mathfrak R)$。\origNote[。][若]{168}若$S$是$\mathfrak S$的算符，所需条件就是$\operatorname{Spur}(US^2)\geq0$。而$S^2$定号，因为
\[
(S^2f,f)=(Sf,Sf)\geq0.
\]
因此，用$A,\allowbreak B$分别代替$U,\allowbreak S^2$，需要证明的是：若$A,\allowbreak B$厄米且定号，则$\operatorname{Spur}(AB)\geq0$。这一点已在 \hyperref[sec:II.11]{II.11} 中借助定号算符的一般定理（参见原注 \hyperlink{note-114}{\getrefnumber{orig-note-114}}）证明过。\origNote[。][]{169}

至此，函数$\operatorname{Erw}(\mathfrak R)$已完全确定：它们对应于定号厄米算符$U$，相互关系由 \hyperref[eq:IV.2.Sp]{Sp.} 给出。我们把$U$称为相应系综的\textbf{统计算符}。

现在很容易讨论 \hyperref[sec:IV.1]{IV.1} 中的附注 \hyperref[remark:IV.1.1]{1}、\hyperref[remark:IV.1.2]{2}、\hyperref[remark:IV.1.3]{3}。与之对应，有：

\noindent\phantomsection1.\label{remark:IV.2.1}从相对概率和相对期望值来看，$U$与$cU$没有本质区别，其中$c>0$为常数。

\hypertarget{origpage-169}{}\phantomsection\label{origpage-169}\noindent\phantomsection2.\label{remark:IV.2.2}$U=0$不能给出任何论断，故须排除。

\noindent\phantomsection3.\label{remark:IV.2.3}当$\operatorname{Spur}U=1$时，得到绝对的、即正确归一化的概率和期望值。只要$\operatorname{Spur}U$有限，就可以根据 \hyperref[remark:IV.2.1]{1}，把$U$乘以$c=1/\operatorname{Spur}U$而归一化。由于$U$定号，$\operatorname{Spur}U\geq0$；实际上还有$\operatorname{Spur}U>0$，因为$\operatorname{Spur}U=0$会导致$U=0$，而这已被 \hyperref[remark:IV.2.2]{2} 排除。这一点在 \hyperref[sec:IV.1]{IV.1} 末尾已作一般证明，在目前情形中也可直接由 \hyperref[sec:II.11]{II.11} 得到。只有$\operatorname{Spur}U$为无穷大时，概率和期望值才本质上只能是相对的。

最后还须研究 \hyperref[sec:IV.1]{IV.1} 的 \hyperref[def:IV.1.alpha]{$\alpha)$}、\hyperref[def:IV.1.beta]{$\beta)$}，也就是在这些$U$中找出零方差的和纯的。

先看零方差的。此时须假设$U$已正确归一化，参见 \hyperref[sec:IV.1]{IV.1}；而且应恒有$\operatorname{Erw}(\mathfrak R^2)=[\operatorname{Erw}(\mathfrak R)]^2$，即$\operatorname{Spur}(UR^2)=[\operatorname{Spur}(UR)]^2$。对$R=P_{[\varphi]}$，有$R^2=R=P_{[\varphi]}$以及$\operatorname{Spur}(UP_{[\varphi]})=(U\varphi,\allowbreak \varphi)$，所以应有$(U\varphi,\allowbreak \varphi)=(U\varphi,\allowbreak \varphi)^2$，即$(U\varphi,\allowbreak \varphi)=0$或$1$。若$\|\varphi'\|=1$、$\|\varphi''\|=1$，总能使$\varphi$从$\varphi'$连续变化到$\varphi''$，并始终保持$\|\varphi\|=1$。\origNote[。][此]{170}此时$(U\varphi,\allowbreak \varphi)$也连续变化，但它只能取$0$或$1$，因而必为常数。所以$(U\varphi',\allowbreak \varphi')=(U\varphi'',\allowbreak \varphi'')$。于是$(U\varphi,\allowbreak \varphi)$或恒等于$0$，或恒等于$1$；立刻可得$U=0$或$U=1$。$U=0$已被 \hyperref[remark:IV.2.2]{2} 排除；$U=1$则不能归一化，因为$\operatorname{Spur}1$等于空间的维数，即$\infty$。而且不难直接看出，它也不是零方差的。因此，零方差系综不存在。

再看纯的。根据 \hyperref[def:IV.1.beta]{$\beta)$} 及 \hyperref[eq:IV.2.Sp]{Sp.}，称$U$为纯的，是指若
\[
U=V+W,
\]
其中$V,\allowbreak W$与$U$一样厄米且定号，就必有$V=c'U$、$W=c''U$。\origNote[。][我]{171}我们断言：具有这一性质的恰好是$U=P_{[\varphi]}$，$\|\varphi\|=1$。

先设$U$具有上述性质。由于$U\ne0$，存在$f_0$使$Uf_0\ne0$；于是$f_0\ne0$，从而$(Uf_0,\allowbreak f_0)>0$（参见 \hyperref[sec:II.5]{II.5}，\CJKglue \hyperref[thm:II.19]{定理19}）。构造两个厄米算符
\[
Vf=\frac{(f,Uf_0)}{(Uf_0,f_0)}Uf_0,\qquad
Wf=Uf-\frac{(f,Uf_0)}{(Uf_0,f_0)}Uf_0
\]
则
\begin{align*}
(Vf,f)&=\frac{|(f,Uf_0)|^2}{(Uf_0,f_0)}\geq0,\\
(Wf,f)&=\frac{(Uf,f)(Uf_0,f_0)-|(f,Uf_0)|^2}{(Uf_0,f_0)}\geq0;
\end{align*}
参见 \hyperref[sec:II.5]{II.5}，\CJKglue \hyperref[thm:II.19]{定理19}。因此$V,\allowbreak W$定号，且显然$U=V+W$。于是$V=c'U$；又因$Vf_0=Uf_0\ne0$，必有$c'=1$，即$U=V$。令
\[
\varphi=\frac{Uf_0}{\|Uf_0\|},\qquad\|\varphi\|=1,\qquad
c=\frac{\|Uf_0\|^2}{(Uf_0,f_0)}>0,
\]\hypertarget{origpage-170}{}\phantomsection\label{origpage-170}就得$Uf=Vf=c(f,\allowbreak \varphi)\varphi=cP_{[\varphi]}f$，即$U=cP_{[\varphi]}$。根据 \hyperref[remark:IV.2.1]{1}，$U$本质上就是$P_{[\varphi]}$。

反过来，设$U=P_{[\varphi]}$，$\|\varphi\|=1$。若$U=V+W$，$V,\allowbreak W$定号，便可作如下推论。由$Uf=0$得
\[
0\leq(Vf,f)\leq(Vf,f)+(Wf,f)=(Uf,f)=0,
\]
所以$(Vf,\allowbreak f)=0$，从而$Vf=0$（参见刚才所引处）。而$(f,\allowbreak \varphi)=0$会给出$Uf=P_{[\varphi]}f=0$，故也有$Vf=0$。于是对任意$g$，有$(f,\allowbreak Vg)=(Vf,\allowbreak g)=0$。也就是说，凡与$\varphi$正交的，也与$Vg$正交，所以$Vg=c_g\varphi$，其中$c_g$是依赖于$g$的数。我们只须用到$g=\varphi$的情形，即$V\varphi=c'\varphi$。每个$f$均可写成$(f,\allowbreak \varphi)\varphi+f'$，其中$f'$与$\varphi$正交，因而
\[
Vf=(f,\varphi)V\varphi+Vf'=(f,\varphi)c'\varphi=c'P_{[\varphi]}f=c'Uf.
\]
所以$V=c'U$，$W=U-V=(1-c')U$，至此全部证毕。

因此，纯系综对应于$U=P_{[\varphi]}$，$\|\varphi\|=1$；在这种情形中，\hyperref[eq:IV.2.Sp]{Sp.} 就成为 \hyperref[sec:III.1]{III.1} 的公式 \hyperref[post:III.1.E2]{$E_2$}：
\begin{equation}\tag{$E_2$}\label{eq:IV.2.E2}
\operatorname{Erw}(\mathfrak R)=\operatorname{Spur}(P_{[\varphi]}R)=(R\varphi,\varphi).
\end{equation}
注意，$\operatorname{Erw}(1)=\operatorname{Spur}(P_{[\varphi]})=1$，因为$P_{[\varphi]}$对应于一维空间$[\varphi]$，也可以直接由 \hyperref[eq:IV.2.E2]{$E_2$} 看出。所以上述形式的$U$已正确归一化。最后，再问$P_{[\varphi]}$与$P_{[\psi]}$何时给出相同统计，也就是何时有$P_{[\varphi]}=cP_{[\psi]}$，其中$c>0$为常数（参见 \hyperref[remark:IV.2.1]{1}）。由于$\operatorname{Spur}(P_{[\varphi]})=\operatorname{Spur}(P_{[\psi]})=1$，必有$c=1$，于是$P_{[\varphi]}=P_{[\psi]}$，$[\varphi]=[\psi]$，即$\varphi=a\psi$。再由$\|\varphi\|=\|\psi\|=1$，得常数$a$满足$|a|=1$；这显然也是充分条件。

归结起来：零方差系综不存在；纯系综则存在，它们恰好对应于$U=P_{[\varphi]}$，$\|\varphi\|=1$。对这些$U$，\hyperref[eq:IV.2.Sp]{Sp.} 化为 \hyperref[eq:IV.2.E2]{$E_2$}，归一化也已正确。若把$\varphi$换成$a\varphi$，其中$a$为常数且$|a|=1$，则$U$不变；除此以外，对$\varphi$的任何改变都会使$U$发生本质变化，参见 \hyperref[remark:IV.2.1]{1}。因此，纯系综正对应于我们先前所刻画的量子力学状态：它由希尔伯特空间中一个满足$\|\varphi\|=1$的$\varphi$表示，模为$1$的常数因子无关紧要（参见 \hyperref[sec:III.2]{III.2}），而统计断言由 \hyperref[eq:IV.2.E2]{$E_2$} 给出。\origNote[。][]{172}

这一切都是从纯定性的条件 \hyperref[cond:IV.Aprime]{A$'$}、\hyperref[cond:IV.Bprime]{B$'$}、\hyperref[def:IV.1.alpha]{$\alpha)$}、\hyperref[def:IV.1.beta]{$\beta)$}、\hyperref[cond:IV.2.I]{I}、\hyperref[cond:IV.2.II]{II} 推导出来的。

在我们的条件范围内，判断至此已经作出，而且结论不支持因果性：所有系综中都存在方差不为零的物理量，纯系综也不例外。

还须讨论 \hyperref[sec:III.2]{III.2} 提出的“隐参数”问题。也就是说，波函数$\varphi$（即 \hyperref[eq:IV.2.E2]{$E_2$}）所刻画的纯系综之所以有非零\hypertarget{origpage-171}{}\phantomsection\label{origpage-171}方差，会不会是因为它们并非真正的状态，而只是若干状态的混合？要刻画真正的状态，除了给出波函数$\varphi$，还须给出其他数据，即“隐参数”；这些数据合在一起，便按因果律确定一切，也就是导向零方差系综。这样，纯系综（$U=P_{[\varphi]}$，$\|\varphi\|=1$）的统计，就来自对组成它的所有真正状态取平均，即对这些状态中实际实现的“隐参数”取值范围取平均。但有两个理由使这不可能：第一，这样一来，该纯系综就能表示成两个不同系综的混合\origNote[合][，]{173}，与其定义矛盾；第二，与“真正”状态对应的零方差系综——即全由处于同一“真正”状态的系统组成的系综——根本不存在。须注意，我们完全不必深入“隐参数”机制的细节，就能作出这一判断：无论如何，不能借助它们重新得到量子力学那些已确立的结果。甚至，只要波函数之外还存在其他决定因素，即“隐参数”，就不可能仍保有所述的那些物理量及其运算关系，也就是说，\hyperref[cond:IV.2.I]{I}、\hyperref[cond:IV.2.II]{II} 不可能仍成立。

在量子力学用算符表示的那些已知物理量之外，即使还有尚未发现的其他物理量，也不足以解决问题；因为即便对前面那些已知量，量子力学所规定的运算关系，即 \hyperref[cond:IV.2.I]{I}、\hyperref[cond:IV.2.II]{II}，也必须失效。因此，这并不像人们常以为的那样，只是怎样诠释量子力学的问题。要使基本过程可能表现出统计行为以外的行为，量子力学本身就必须在客观上是错误的。

另有一点也值得提出。初看起来，不确定性关系与相对论的基本公设颇为相似。相对论说：原则上不可能确定两个相距$r$的点上发生的事件是否同时，其精度不能超过长度为$r/c$的时间区间（$c$为光速）。\hypertarget{origpage-172}{}\phantomsection\label{origpage-172}不确定性关系则说：原则上不可能在相空间中指定一个质点的位置，使其精度超过体积为$(h/4\pi)^3$的区域。\origNote[。][然]{174}然而，二者有根本区别。相对论固然否认客观而精确地测量同时性的可能，却仍可通过引入一个伽利略参考系，在世界中设置一个坐标系，由此定义一种与我们通常的同时性概念足够相符的同时性。之所以不赋予这种定义客观意义，只是因为坐标系可以有无穷多种不同选择，从而产生无穷多种同样好的同时性定义。也就是说，这里不能测量，是因为可采用的理论定义有无穷多种。量子力学则不同：即使引入新的、假设的、未被观察到的坐标，即“隐参数”，也根本无法用相空间中的点来描述波函数为$\varphi$的系统，因为这将导向零方差系综。因此，不仅测量不可能，任何合理的理论定义也不可能。这里说的合理定义，是指像相对论中的定义那样，虽不能由经验来证明，却至少也不能由经验来驳倒的定义。因此，测量之所以原则上不可能，在一种情形中，是因为有关概念有无穷多种定义方式，都不违背直接经验或理论的一般基本假设；在另一种情形中，则是因为这样的定义根本不存在。

据此，可以这样概括因果性在当代物理学中的处境。在宏观领域，没有支持它的经验，也不可能有。世界在大尺度上（即对肉眼可见的对象）表现出的因果秩序，原因无疑只是“大数定律”，与支配基本过程的自然规律——即真正的自然规律——是否具有因果性毫不相干。\origNote[。][宏]{175}宏观上相同的对象在宏观上表现相同，与因果性并没有多大关系。它们其实并不相同：精确规定各原子状态的那些坐标，几乎从不相等；宏观描述则对这些坐标取了平均，在这里，它们就是“隐参数”。然而，这些坐标的数量极大（$1\,\mathrm g$物质中约有$10^{25}$个），所以按照概率论的已知定理，取平均会使一切方差大幅减小。当然，这只是一般情况；在某些特殊场合，例如布朗运动、不稳定状态等，这种表面上的宏观因果性就不成立。只有在原子领域、在基本过程本身中，才能真正检验\hypertarget{origpage-173}{}\phantomsection\label{origpage-173}因果性问题；但以我们目前的知识，一切都不利于它。现有的形式理论中，唯一能把经验相当令人满意地整理、归纳起来的，就是量子力学，而它与因果性有着不可回避的逻辑矛盾。固然，若说因果性就此被彻底否定，未免言过其实。量子力学目前的形式无疑还有缺漏，甚至可能是错误的；不过，考虑到它在理解一般问题和计算具体问题时表现出的惊人效能，后一种可能相当小。它虽然与经验极为吻合，也让我们看到了世界在性质上全新的一面，但对任何理论，我们都不能说它已被经验证实，只能说它是目前已知对经验最好的概括。即使充分保留这些限制，我们仍可以说：目前既没有理由，也没有借口，主张自然中存在因果性。没有任何经验支持它的存在，因为宏观经验原则上就不适合做这件事；而与我们关于基本过程的经验相容的唯一已知理论——量子力学——又与它矛盾。

因果性固然是人人根深蒂固的一种看法，却绝不是思维上的必然要求；统计理论能够建立起来，这件事本身就已说明了这一点。没有成见地面对问题，就没有理由坚持它。在这样的情况下，为它牺牲一种合理的物理理论，难道有根据吗？

\section{由实验作出的推论}\label{sec:IV.3}
上一节说明，与我们的定性基本假设相容的最一般统计系综，可按规律 \hyperref[eq:IV.2.Sp]{Sp.} 由一个定号算符$U$刻画。其中我们称为纯系综的特殊情形，由$U=P_{[\varphi]}$、$\|\varphi\|=1$刻画；由于它们就是系统$\boldsymbol{S}$真正的、不能再分解的状态，我们也称它们为状态，并把$U=P_{[\varphi]}$称为状态$\varphi$。

如果$U$有纯点谱，例如本征值为$w_1,\allowbreak w_2,\allowbreak \ldots$，本征函数为$\varphi_1,\allowbreak \varphi_2,\allowbreak \ldots$，且构成完全规范正交系，那么由 \hyperref[sec:II.8]{II.8}，
\[
U=\sum_n w_nP_{[\varphi_n]}.
\]\hypertarget{origpage-174}{}\phantomsection\label{origpage-174}因为$U$定号，所有$w_n\geq0$：事实上，$U\varphi_n=w_n\varphi_n$，故$(U\varphi_n,\allowbreak \varphi_n)=w_n\geq0$。此外，
\[
\sum_nw_n=\sum_n(U\varphi_n,\varphi_n)=\operatorname{Spur}U;
\]
也参见 \hyperref[sec:IV.1]{IV.1} 开头。因此，$U$已正确归一化时，$\sum_nw_n=1$。按照 \hyperref[sec:IV.1]{IV.1} 开头的说明，$U$便可看成状态$\varphi_1,\allowbreak \varphi_2,\allowbreak \ldots$以相对权重$w_1,\allowbreak w_2,\allowbreak \ldots$混合的结果；若$U$已正确归一化，这些权重在绝对意义上也正确。

但由 \hyperref[sec:II.11]{II.11}（特别参见原注 \hyperlink{note-115}{\getrefnumber{orig-note-115}}），正确归一化的$U$，即$\operatorname{Spur}U=1$，是全连续的，因而有纯点谱；当然，只要$\operatorname{Spur}U$有限也是如此。$\operatorname{Spur}U$为无穷大的情形可以视为极限，这里不深入讨论。因此，在真正值得研究的情形中，所考察的系综总能表示为状态的混合，而且我们还把这些状态选成了两两正交的。以下就把一般系综称为\textbf{混合态}，以区别于纯系综，即状态。

若$U$的所有本征值都是单重的，即$w_1,\allowbreak w_2,\allowbreak \ldots$彼此不同，那么我们知道，$\varphi_1,\allowbreak \varphi_2,\allowbreak \ldots$只差模为$1$的常数因子便唯一确定；因此相应的状态以及$P_{[\varphi_1]},\allowbreak P_{[\varphi_2]},\allowbreak \ldots$都完全确定。权重$w_1,\allowbreak w_2,\allowbreak \ldots$也唯一确定，当然不计排列次序。所以，此时能唯一指出构成混合态$U$的那些两两正交的状态。但若$U$还有重本征值（“简并”），情形便本质不同。在 \hyperref[sec:II.8]{II.8} 中，我们已详细讨论$\varphi_1,\allowbreak \varphi_2,\allowbreak \ldots$的选法：它们有无穷多种本质不同的选择，而$w_1,\allowbreak w_2,\allowbreak \ldots$仍然唯一确定。只须把$w_1,\allowbreak w_2,\allowbreak \ldots$中互不相同的值列为$w',\allowbreak w'',\allowbreak \ldots$，对每一个写出其本征函数构成的闭线性流形，即$Uf=wf$的解空间$\mathfrak M_{w'},\allowbreak \mathfrak M_{w''},\allowbreak \ldots$。然后，在各个$\mathfrak M_{w'},\allowbreak \mathfrak M_{w''},\allowbreak \ldots$中任意选择张成它的规范正交系，分别记为$\chi'_1,\allowbreak \chi'_2,\allowbreak \ldots$、$\chi''_1,\allowbreak \chi''_2,\allowbreak \ldots$，等等。$\chi'_1,\allowbreak \chi'_2,\allowbreak \ldots,\allowbreak \chi''_1,\allowbreak \chi''_2,\allowbreak \ldots,\allowbreak \ldots$就是所求$\varphi_1,\allowbreak \varphi_2,\allowbreak \ldots$，相应的本征值$w',\allowbreak w',\allowbreak \ldots,\allowbreak w'',\allowbreak w'',\allowbreak \ldots,\allowbreak \ldots$就是$w_1,\allowbreak w_2,\allowbreak \ldots$。只要有一个$\mathfrak M_w$的维数超过一，即存在重本征值，其中的$\chi_1,\allowbreak \chi_2,\allowbreak \ldots$就不再只是相差模为$1$的常数因子；例如，$\chi_1$可以是$\mathfrak M_w$中的任何归一化元素。也就是说，它们的状态本身也不唯一。

这一现象也可表述如下。如果状态$\chi_1,\allowbreak \chi_2,\allowbreak \ldots$两两正交，即构成一个有限或无限的规范正交系，而我们以完全相同的\hypertarget{origpage-175}{}\phantomsection\label{origpage-175}权重，例如相对权重$1:1:\cdots$，把它们混合，那么所得混合态只依赖于它们张成的闭线性流形$\mathfrak M$。事实上，
\[
U=P_{[\chi_1]}+P_{[\chi_2]}+\cdots=P_{\mathfrak M}.
\]
如果$\chi_1,\allowbreak \chi_2,\allowbreak \ldots$只有有限个，比如$s$个$\chi_1,\allowbreak \ldots,\allowbreak \chi_s$，还可把这个$U$看成$\mathfrak M$中所有归一化元素，即所有状态的混合。这些元素是
\[
\chi=x_1\chi_1+\cdots+x_s\chi_s,\qquad |x_1|^2+\cdots+|x_s|^2=1.
\]
确实，令$x_1=u_1+iv_1,\allowbreak \ldots,\allowbreak x_s=u_s+iv_s$，把条件$|x_1|^2+\cdots+|x_s|^2=1$对应的$(2s-1)$维单位球面
\[
u_1^2+v_1^2+\cdots+u_s^2+v_s^2=1,
\]
记为$K$，面元记为$do$。对于
\[
U'=\int\!\!\int_K\cdots\int\!\!\int P_{[\chi]}\,do,
\]
有
\begin{align*}
(U'f,g)
&=\int\!\!\int_K\cdots\int\!\!\int(P_{[\chi]}f,g)\,do\\
&=\int\!\!\int_K\cdots\int\!\!\int(f,\chi)\overline{(g,\chi)}\,do\\
&=\int\!\!\int_K\cdots\int\!\!\int
\left(f,\sum_{\mu=1}^s(u_\mu+iv_\mu)\chi_\mu\right)\\
&\qquad{}\cdot\overline{\left(g,\sum_{\mu=1}^s(u_\mu+iv_\mu)\chi_\mu\right)}\,do\\
&=\int\!\!\int_K\cdots\int\!\!\int\sum_{\mu,\nu=1}^s
(f,\chi_\mu)\overline{(g,\chi_\nu)}\\
&\qquad{}\cdot(u_\mu-iv_\mu)(u_\nu+iv_\nu)\,do\\
&=\sum_{\mu,\nu=1}^s(f,\chi_\mu)\overline{(g,\chi_\nu)}\\
&\qquad{}\cdot\int\!\!\int_K\cdots\int\!\!\int
[(u_\mu u_\nu+v_\mu v_\nu)+i(u_\mu v_\nu-u_\nu v_\mu)]\,do.
\end{align*}
由于对称性，所有$u_\mu v_\nu$、$v_\nu u_\mu$的积分，以及$\mu\ne\nu$时$u_\mu u_\nu$、$v_\mu v_\nu$的积分都等于零\origNote[零][；]{176}；当$\mu=\nu$时，后两种积分都等于$C/(2s)$，$C>0$。\origNote[。][所]{176}所以
\begin{align*}
(U'f,g)&=\frac Cs\sum_{\mu=1}^s(f,\chi_\mu)\overline{(g,\chi_\mu)}
=\frac Cs\sum_{\mu=1}^s(P_{[\chi_\mu]}f,g)\\
&=\left(\left\{\frac Cs\sum_{\mu=1}^sP_{[\chi_\mu]}\right\}f,g\right).
\end{align*}
于是
\[
U'=\frac Cs\sum_{\mu=1}^sP_{[\chi_\mu]}=\frac CsU,
\]
即$U',\allowbreak U$没有本质区别。

这些事实对量子力学统计的性质十分重要，因此我们再说一遍：

\noindent1.若一个混合态由彼此正交的状态以完全相等的权重组成，就无法再确定这些状态究竟是哪些。换言之，把不同组彼此正交的组分以完全相等的\hypertarget{origpage-176}{}\phantomsection\label{origpage-176}比例混合，可以得到同一个混合态。

\noindent2.若组分只有有限个，这样得到的混合态，与这些组分的一切线性组合所代表的状态的混合相同。

最简单的例子是：将正交的$\varphi,\allowbreak \psi$按$1:1$混合，所得结果，与把$(\varphi+\psi)/\sqrt2$和$(\varphi-\psi)/\sqrt2$按$1:1$混合相同，也与把所有满足$|x|^2+|y|^2=1$的$x\varphi+y\psi$混合相同。如果混合的是两个不正交的$\varphi,\allowbreak \psi$（比例甚至不必是$1:1$），那么从所得混合态中就更加无法认出它们，因为这个混合态肯定也能视为正交状态的混合。

关于混合态的性质，留到 \hyperref[sec:V.2]{V.2} 及以后讨论热力学时再深入研究。

\hyperref[sec:IV.2]{IV.2} 的公式 \hyperref[eq:IV.2.Sp]{Sp.} 已经说明：在统计算符为$U$的混合态中，算符为$R$的物理量$\mathfrak R$的期望值为$\operatorname{Spur}(UR)$。若问$R$的值$a$落在区间$a'<a\leq a''$内的概率是多少（$a',\allowbreak a''$给定，且$a'\leq a''$），便可照 \hyperref[sec:III.1]{III.1} 或 \hyperref[sec:III.5]{III.5} 的方法计算。用函数
\[
F(x)=\begin{cases}1,&a'<x\leq a'',\\0,&\text{其他情形}\end{cases}
\]
构造物理量$F(\mathfrak R)$，其期望值就是所求概率。根据 \hyperref[sec:IV.2]{IV.2} 的 \hyperref[cond:IV.2.I]{I}，$F(\mathfrak R)$对应算符$F(R)$；若$E(\lambda)$是$R$的单位分解，便如我们多次算过的那样，有$F(R)=E(a'')-E(a')$。因此所求概率为
\[
w(a',a'')=\operatorname{Spur}\bigl(U(E(a'')-E(a'))\bigr).
\]
所以，描述$\mathfrak R$的统计的概率函数为$w(a)=\operatorname{Spur}(UE(a))$，参见 \hyperref[sec:IV.1]{IV.1}、原注 \hyperlink{note-175}{\getrefnumber{orig-note-175}}。\tnote[。][对][原]{原书引作原注 \hyperlink{note-175}{\getrefnumber{orig-note-175}}；概率函数的定义实际见原注 \hyperlink{note-157}{\getrefnumber{orig-note-157}}。}对于状态，即$U=P_{[\varphi]}$，又得到$w(a)=\operatorname{Spur}(P_{[\varphi]}E(a))=(E(a)\varphi,\allowbreak \varphi)$。当然，若$U$未正确归一化，这些概率就只是相对概率。

算符为$R$的物理量$\mathfrak R$，在统计算符为$U$的混合态中，何时必定取值$\lambda^*$？这个问题可以直接借助$w(a)$回答：对$a<\lambda^*$，须有$w(a)=0$；对$a\geq\lambda^*$，须有$w(a)=1$。若$U$未正确归一化，后一条件只改为$w(a)=\operatorname{Erw}(1)=\operatorname{Spur}U$。所以\origNote[以][]{177}
\begin{align*}
\operatorname{Spur}(UE(a))&=0\quad(a<\lambda^*),\\
\operatorname{Spur}(U(1-E(a)))&=0\quad(a\geq\lambda^*).
\end{align*}
但对定号算符$A,\allowbreak B$，$\operatorname{Spur}(AB)=0$蕴含$AB=0$（参见 \hyperref[sec:II.11]{II.11}）。因此$a<\lambda^*$时$UE(a)=0$，$a\geq\lambda^*$时$UE(a)=U$；等价地，分别有$E(a)U=0$与$E(a)U=U$，因为乘积为厄米算符时，两因子必对易。也就是说，对$f=Ug$有
\[
E(a)f=\begin{cases}f,&a\geq\lambda^*,\\0,&a<\lambda^*;\end{cases}
\]
由 \hyperref[sec:II.8]{II.8} 的讨论，\hypertarget{origpage-177}{}\phantomsection\label{origpage-177}这就表示$Rf=\lambda^*f$，即对$g$恒有$RUg=\lambda^*Ug$。所以最后的条件就是$RU=\lambda^*U$。或者，把$Rh=\lambda^*h$的所有解$h$构成的闭线性流形记为$\mathfrak M$，则条件是：$Uf$始终落在$\mathfrak M$中。

顺便说，同一结果也可以从方差为零导出，即要求$(\mathfrak R-\lambda^*)^2$的期望值（也可以是相对期望值）为零。

在 \hyperref[sec:III.3]{III.3} 中，我们已回答以下问题（$\mathfrak R,\allowbreak \mathfrak S,\allowbreak \ldots$为物理量，$R,\allowbreak S,\allowbreak \ldots$为相应算符）：

\noindent\phantomsection1.\label{criterion:IV.3.1}$\mathfrak R$何时能绝对精确地测量？答：$R$具有纯点谱时。

\noindent\phantomsection2.\label{criterion:IV.3.2}$\mathfrak R,\allowbreak \mathfrak S$何时能同时绝对精确地测量？答：$R,\allowbreak S$都有纯点谱且互相对易时。

\noindent\phantomsection3.\label{criterion:IV.3.3}几个物理量$\mathfrak R,\allowbreak \mathfrak S,\allowbreak \ldots$何时能同时绝对精确地测量？答：$R,\allowbreak S,\allowbreak \ldots$都有纯点谱且彼此对易时。

\noindent\phantomsection4.\label{criterion:IV.3.4}几个物理量$\mathfrak R,\allowbreak \mathfrak S,\allowbreak \ldots$何时能同时任意精确地测量？答：$R,\allowbreak S,\allowbreak \ldots$彼此对易时。

当时采用的是从Compton--Simon实验结果中抽出的如下原理：

\noindent\phantomsection\textbf{(M.)}\label{post:IV.3.M}在一个系统$\boldsymbol{S}$上紧接着测量物理量$\mathfrak R$两次，两次得到相同的值。即使$\mathfrak R$在$\boldsymbol{S}$的原状态中的方差不为零，而且测量$\mathfrak R$还会改变$\boldsymbol{S}$的状态，这一原理仍然成立。

我们在 \hyperref[sec:III.3]{III.3} 已详细讨论过 \hyperref[post:IV.3.M]{M} 的物理意义。回答 \hyperref[criterion:IV.3.1]{1}—\hyperref[criterion:IV.3.4]{4} 时还用到了三个假设：\hyperref[sec:III.3]{III.3} 关于状态的统计公式 \hyperref[post:III.1.E2]{$E_2$}；那里假设 \hyperref[post:III.1.F]{F} 所说的，若$\mathfrak R$的算符为$R$，则$F(\mathfrak R)$的算符为$F(R)$；以及如下假设：若同时可观测的量$\mathfrak R,\allowbreak \mathfrak S$分别对应算符$R,\allowbreak S$，则$R+S$对应算符$\mathfrak R+\mathfrak S$。\tnote[。][][这]{这一句中原书把物理量与算符的位置颠倒了：$\mathfrak R+\mathfrak S$是物理量，$R+S$是其算符。这里保留原书的次序。}

现在这三个假设仍然可用：第一个由 \hyperref[sec:IV.2]{IV.2} 的 \hyperref[eq:IV.2.Sp]{Sp.} 得出，另两个分别对应那里的 \hyperref[cond:IV.2.I]{I}、\hyperref[cond:IV.2.II]{II}。\hyperref[post:IV.3.M]{M} 也应视为成立，因为我们已经认识到，它对量子力学的概念构建不可缺少。因此，\hyperref[sec:III.3]{III.3} 对 \hyperref[criterion:IV.3.1]{1}—\hyperref[criterion:IV.3.4]{4} 的证明在这里仍然适用，以上答案也就仍然正确。

我们在 \hyperref[sec:III.5]{III.5} 考察了只取$0,\allowbreak 1$两个值的物理量，它们与性质$\mathfrak E$一一对应。给定性质$\mathfrak E$时，可以如下定义一个物理量：测量它，就是判断$\mathfrak E$\hypertarget{origpage-178}{}\phantomsection\label{origpage-178}是否具备；具备时取值$1$，不具备时取值$0$。反之，若给定这个量，$\mathfrak E$就是“此量取值$1$而非$0$”这一性质。由那里的 \hyperref[post:III.5.F]{F}，即 \hyperref[sec:IV.2]{IV.2} 的 \hyperref[cond:IV.2.I]{I}，可知，与这些量对应的算符$E$都是投影算符，而且每个投影算符都对应这样的量。具备$\mathfrak E$的概率，等于上述物理量的期望值。在 \hyperref[sec:III.5]{III.5} 中，我们只对状态$\varphi$，即$U=P_{[\varphi]}$、$\|\varphi\|=1$，计算了这一概率；现在可由 \hyperref[eq:IV.2.Sp]{Sp.} 一般地求得它为$\operatorname{Spur}(UE)$。这是相对概率；只有$U$正确归一化，即$\operatorname{Spur}U=1$时，才是绝对概率。

既然 \hyperref[criterion:IV.3.1]{1}—\hyperref[criterion:IV.3.4]{4} 已经确立，\hyperref[sec:III.5]{III.5} 由它们导出的 \hyperref[post:III.5.alpha]{$\alpha)$}—\hyperref[post:III.5.zeta]{$\zeta)$} 也就确立了。不过须注意，那里 \hyperref[post:III.5.alpha]{$\alpha)$} 只涉及状态；我们刚才已把它推广到一切混合态：

\noindent\phantomsection$\alpha')$\label{prop:IV.3.alphaprime}在统计算符为$U$的混合态中，具备与不具备性质$\mathfrak E$的概率分别为
\[
\operatorname{Spur}(UE)\qquad\text{及}\qquad\operatorname{Spur}(U(1-E)),
\]
这些是相对概率；仅当$U$正确归一化，即$\operatorname{Spur}U=1$时，才是绝对概率。

若要研究几个物理量$\mathfrak R_1,\allowbreak \ldots,\allowbreak \mathfrak R_l$，设它们的算符分别为$R_1,\allowbreak \ldots,\allowbreak R_l$，相应的单位分解为$E_1(\lambda),\allowbreak \ldots,\allowbreak E_l(\lambda)$。给定$l$个区间
\[
I_1:\lambda'_1<\lambda\leq\lambda''_1,\quad\ldots,\quad
I_l:\lambda'_l<\lambda\leq\lambda''_l
\]
并令$E_1(I_1)=E_1(\lambda''_1)-E_1(\lambda'_1),\allowbreak \ldots,\allowbreak E_l(I_l)=E_l(\lambda''_l)-E_l(\lambda'_l)$。那么，性质“$\mathfrak R_1$落在$I_1$中”、……、“$\mathfrak R_l$落在$I_l$中”分别对应投影算符$E_1(I_1),\allowbreak \ldots,\allowbreak E_l(I_l)$，参见 \hyperref[post:III.5.zeta]{$\zeta)$}。这些性质同时可判定，当且仅当$E_1(I_1),\allowbreak \ldots,\allowbreak E_l(I_l)$彼此对易，参见 \hyperref[post:III.5.gamma]{$\gamma)$}；它们同时成立这一性质的投影算符为$E=E_1(I_1)\cdots E_l(I_l)$，参见 \hyperref[post:III.5.epsilon]{$\varepsilon)$}。所以其概率为$\operatorname{Spur}(UE)$，参见 \hyperref[prop:IV.3.alphaprime]{$\alpha')$}。

现在反过来走。假设我们不知道系统$\boldsymbol{S}$的状态，但对它作过某些测量，并知道结果。实际情况总是如此，因为只有测量结果能使我们获知$\boldsymbol{S}$的状态。严格说来，状态只是一种理论构造；实际存在的只有测量结果，而物理学的任务，是找出过去与未来的测量结果之间的联系。固然，人们总是引入“状态”这个辅助概念来完成这一任务；但这样一来，物理理论一方面必须说明怎样由过去的测量结果推断现在的状态，另一方面又必须说明怎样由现在的状态推断未来的测量结果。我们此前只讨论了\hypertarget{origpage-179}{}\phantomsection\label{origpage-179}第二个问题，现在须转向第一个。

如果过去的测量不足以唯一确定现在的状态，或许仍能由它们得知各个状态出现的概率。在经典力学等因果性理论中如此，在量子力学中也如此。因此，真正的问题是：给定测量结果，找出一个混合态，使它的统计与如下系统$\boldsymbol{S}$所应有的统计相同——关于这个系统，我们只知道对它进行过所述测量，并得到所述结果。不过还须更确切地说明，关于$\boldsymbol{S}$“只知道这些”、此外一无所知，究竟是什么意思，以及怎样由此得到统计。

与统计的联系无论如何必须是这样：如果对许多个系统$\boldsymbol{S}'_1,\allowbreak \ldots,\allowbreak \boldsymbol{S}'_M$（都是$\boldsymbol{S}$的副本）作所述测量，都得到所述结果，那么系综$[\boldsymbol{S}'_1,\allowbreak \ldots,\allowbreak \boldsymbol{S}'_M]$的一切统计性质，都与上面配给这些测量结果的混合态一致。按照 \hyperref[post:IV.3.M]{M}，可把所有$\boldsymbol{S}'_1,\allowbreak \ldots,\allowbreak \boldsymbol{S}'_M$得到相同结果的过程设想为：原先有一个大系综$[\boldsymbol{S}_1,\allowbreak \ldots,\allowbreak \boldsymbol{S}_N]$，对它作所述测量，把得到所需结果的那些元素汇集为新系综，即$[\boldsymbol{S}'_1,\allowbreak \ldots,\allowbreak \boldsymbol{S}'_M]$。当然，一切取决于初始系综$[\boldsymbol{S}_1,\allowbreak \ldots,\allowbreak \boldsymbol{S}_N]$如何选取。这个系综可以说规定了系统$\boldsymbol{S}$各个状态的先验概率。一般概率论中，这种情况本来就很熟悉：要由测量结果推断状态，也就是由果推因、计算后验概率，必须先知道先验概率。一般说来，先验概率有多种选法，因此我们的问题也没有唯一答案。但我们将看到，在量子力学的特殊情况下，初始系综$[\boldsymbol{S}_1,\allowbreak \ldots,\allowbreak \boldsymbol{S}_N]$——也就是先验概率——有一种特别突出的规定方式。

如果测量结果已经充分，能够完全确定$\boldsymbol{S}$的状态，情况就不同了：此时每个问题都必须有唯一答案。稍后将看到这一点怎样起作用。

最后还须指出：与其说已知关于$\boldsymbol{S}$的若干测量结果，也可以说，已检验$\boldsymbol{S}$是否具备某一性质$\mathfrak E$，并已确定它\hypertarget{origpage-180}{}\phantomsection\label{origpage-180}具备这一性质。我们由 \hyperref[post:III.5.alpha]{$\alpha)$}—\hyperref[post:III.5.zeta]{$\zeta)$} 已知道这些说法如何联系。例如，若能同时确定的测量结果是各量$\mathfrak R_1,\allowbreak \ldots,\allowbreak \mathfrak R_l$分别落在区间$I_1,\allowbreak \ldots,\allowbreak I_l$内，那么用前面的记号，$\mathfrak E$的投影算符就是$E=E_1(I_1)\cdots E_l(I_l)$。

因此，我们关于$\boldsymbol{S}$的知识，总归结为某个性质$\mathfrak E$成立；形式上，给出投影算符$E$即可刻画它。所求的是等价系综$[\boldsymbol{S}'_1,\allowbreak \ldots,\allowbreak \boldsymbol{S}'_M]$的统计算符$U$，以及一般初始系综$[\boldsymbol{S}_1,\allowbreak \ldots,\allowbreak \boldsymbol{S}_N]$的统计算符$U_0$。$E,\allowbreak U,\allowbreak U_0$之间有怎样的数学关系？

由 \hyperref[post:IV.3.M]{M}，$[\boldsymbol{S}'_1,\allowbreak \ldots,\allowbreak \boldsymbol{S}'_M]$也必定具备$\mathfrak E$，即对应于$\mathfrak E$的量取值$1$。如本节开头所见，这表示$EU=U$，也就是$Uf$恒落在$\mathfrak M$中；$\mathfrak M$是所有满足$Ef=f$的$f$所成的集合，即对应于$E$的闭线性流形。

不用$EU=U$，也可以写成$UE=U$、$U(1-E)=0$。于是对所有$g=(1-E)f$都有$Ug=0$，也就是对$1-E$所对应的闭线性流形中的所有$g$成立，亦即对$\mathfrak R-\mathfrak M$中的所有$g$成立。因此，$f$属于$\mathfrak R-\mathfrak M$时$Uf=0$，而$f$属于$\mathfrak M$时$Uf$也在$\mathfrak M$中。由这一路径，关于$U$只能说到这里。

当且仅当$\mathfrak M$为零维或一维时，上述条件才能在本质上确定$U$，即只剩一个常数因子未定。事实上，$\mathfrak M=[0]$时得$U=0$，根据 \hyperref[sec:III.2]{III.2} 的附注1，这是不可能的。\tnote[。][若][原]{原书引作III.2附注1；排除$U=0$的论述实际见 \hyperref[remark:IV.2.2]{IV.2附注2}。}若$\mathfrak M=[\varphi]$，$\varphi\ne0$，便可设$\|\varphi\|=1$；此时$U\varphi=c\varphi$，所以对$\mathfrak M$中一切$f$都有$Uf=cf$，因为它们均形如$a\varphi$。于是一般地有$Uf=UEf=cEf$，即$U=cE=cP_{[\varphi]}$。因$U$定号且非零，有$c>0$，所以$U$本质上就是$E=P_{[\varphi]}$。若$\mathfrak M$至少二维，就能从中选取规范正交的$\varphi,\allowbreak \psi$；这时$P_{[\varphi]},\allowbreak P_{[\psi]}$就是满足条件的两个本质不同的$U$。所以：$E=0$不可能；$E=P_{[\varphi]}$、$\|\varphi\|=1$时，有$U=E=P_{[\varphi]}$；其余情况下$U$不唯一。

如果系统$\boldsymbol{S}$可能具备$E=0$所对应的性质$\mathfrak E$，那么根本找不到$U$就会成为麻烦。然而，\hyperref[post:III.5.eta]{$\eta)$} 排除了这一可能：$\mathfrak E$永不成立，其概率总为$0$。一维的$\mathfrak M$，即$E=P_{[\varphi]}$、$\|\varphi\|=1$，则唯一确定$U$，它就是状态$\varphi$。因此，这正是一种能够完全确定$\boldsymbol{S}$状态的测量：如果结果为肯定，就确定为状态$\varphi$。\origNote[。][其]{178}其余测量都是不完全的，不足以确定状态。

一般情形可如下处理。把对应于$\mathfrak E$的量也记为$\mathfrak E$，则$U$的产生方式是：在\hypertarget{origpage-181}{}\phantomsection\label{origpage-181}$U_0$所对应的系综$[\boldsymbol{S}_1,\allowbreak \ldots,\allowbreak \boldsymbol{S}_N]$上测量$\mathfrak E$，把取值为$1$的所有元素汇集成$U$的系综$[\boldsymbol{S}'_1,\allowbreak \ldots,\allowbreak \boldsymbol{S}'_M]$。测量$\mathfrak E$的方法可以有很多种，例如测量另一个量$\mathfrak R$，只须$\mathfrak E$是它的一个已知函数，即$\mathfrak E=F(\mathfrak R)$。比如，设$\varphi_1,\allowbreak \varphi_2,\allowbreak \ldots$是张成$\mathfrak M$的规范正交系，$\psi_1,\allowbreak \psi_2,\allowbreak \ldots$是$\mathfrak R-\mathfrak M$中的相应正交系。于是$\varphi_1,\allowbreak \varphi_2,\allowbreak \ldots,\allowbreak \psi_1,\allowbreak \psi_2,\allowbreak \ldots$张成$\mathfrak M+(\mathfrak R-\mathfrak M)=\mathfrak R$，即为完全系。取互不相同的实数$\lambda_1,\allowbreak \lambda_2,\allowbreak \ldots,\allowbreak \mu_1,\allowbreak \mu_2,\allowbreak \ldots$，并以
\[
R\left(\sum_nx_n\varphi_n+\sum_nx_n\psi_n\right)
=\sum_n\lambda_nx_n\varphi_n+\sum_n\mu_ny_n\psi_n
\]
定义算符$R$。显然，$R$的纯点谱为$\lambda_1,\allowbreak \lambda_2,\allowbreak \ldots,\allowbreak \mu_1,\allowbreak \mu_2,\allowbreak \ldots$，相应本征函数为$\varphi_1,\allowbreak \varphi_2,\allowbreak \ldots,\allowbreak \psi_1,\allowbreak \psi_2,\allowbreak \ldots$，且所有本征值都是单重的。\tnote[。][若][原]{原书左端第二个和式的系数印作$x_n$，右端相应处则为$y_n$。公式保留原印。}若$F(x)$是满足
\[
F(\lambda_n)=1,\qquad F(\mu_n)=0
\]
的任意函数，则对$\varphi_1,\allowbreak \varphi_2,\allowbreak \ldots$，进而对$\mathfrak M$中每个$f$，$F(R)$的本征值都为$1$；对$\psi_1,\allowbreak \psi_2,\allowbreak \ldots$，进而对$\mathfrak R-\mathfrak M$中每个$f$，本征值都为$0$。所以$E=F(R)$。若$R$对应于$\mathfrak R$，便有$\mathfrak E=F(\mathfrak R)$。因此，测量$\mathfrak E$可以通过测量$\mathfrak R$来实现。

此时可以算出$U_0,\allowbreak U$的关系。测量$\mathfrak R$以后，每个系统都处于$\varphi_1,\allowbreak \varphi_2,\allowbreak \ldots,\allowbreak \psi_1,\allowbreak \psi_2,\allowbreak \ldots$中的某一个状态，分别对应测得值$\lambda_1,\allowbreak \lambda_2,\allowbreak \ldots,\allowbreak \mu_1,\allowbreak \mu_2,\allowbreak \ldots$。其概率依次为
\begin{gather*}
\operatorname{Spur}(U_0P_{[\varphi_1]})=(U_0\varphi_1,\varphi_1),\quad
\operatorname{Spur}(U_0P_{[\varphi_2]})=(U_0\varphi_2,\varphi_2),\quad\ldots,\\
\operatorname{Spur}(U_0P_{[\psi_1]})=(U_0\psi_1,\psi_1),\quad
\operatorname{Spur}(U_0P_{[\psi_2]})=(U_0\psi_2,\psi_2),\quad\ldots;
\end{gather*}
参见 \hyperref[sec:III.3]{III.3} 中已经确认有效的论述。也就是说，$U_0$系综中相应比例的部分，分别转入系综$P_{[\varphi_1]},\allowbreak P_{[\varphi_2]},\allowbreak \ldots,\allowbreak P_{[\psi_1]},\allowbreak P_{[\psi_2]},\allowbreak \ldots$。由于$\mathfrak E=1$对应$\mathfrak R=\lambda_1,\allowbreak \lambda_2,\allowbreak \ldots$，把第一组汇集起来，就得到$U$系综。因此
\[
U=\sum_n(U_0\varphi_n,\varphi_n)P_{[\varphi_n]}.
\]
每个$P_{[\varphi_n]}$都与$R$对易\origNote[易][，]{179}，所以$U$也必须如此。也就是说，如果$U$不能与上述方式构成的每个$R$对易，就有一些测量方法——即以相应$\mathfrak R$为依据的方法——不能从$U_0$产生它。这样，我们对$U$的了解就不只是“它由测量$\mathfrak E$产生”而已。但$U$所要表示的，正是我们仅有所述知识的情况，因此应尽量坚持这一条件：如果存在某种$U$，对它无需排除任何一种$\mathfrak E$的测量方法，我们就优先选取这样的$U$。下面考察这样的$U$是否存在，又究竟是哪些。

\hypertarget{origpage-182}{}\phantomsection\label{origpage-182}如前所见，$U$必须与上述方式产生的所有$R$对易。于是
\[
RU\varphi_n=UR\varphi_n=U(\lambda_n\varphi_n)=\lambda_nU\varphi_n,
\]
即$U\varphi_n$是$R$对应于本征值$\lambda_n$的本征函数，所以$U\varphi_n=a_n\varphi_n$；特别有$U\varphi_1=a_1\varphi_1$。任给$\mathfrak M$中满足$\|\varphi\|=1$的$\varphi$，总可以把$\varphi_1,\allowbreak \varphi_2,\allowbreak \ldots,\allowbreak \psi_1,\allowbreak \psi_2,\allowbreak \ldots$选成$\varphi_1=\varphi$。所以每个这样的$\varphi$都是$U$的本征函数，而且它们全属于同一个本征值。否则，若$\varphi,\allowbreak \psi$属于不同本征值，就必定正交；但$(\varphi+\psi)/\sqrt2$也是本征函数，又因为
\[
\left(\frac{\varphi+\psi}{\sqrt2},\varphi\right)=\frac{(\varphi,\varphi)}{\sqrt2}=\frac1{\sqrt2},\qquad
\left(\frac{\varphi+\psi}{\sqrt2},\psi\right)=\frac{(\psi,\psi)}{\sqrt2}=\frac1{\sqrt2},
\]
它与二者都不正交，因而又须与二者都属于同一本征值，这与原设矛盾。因此$U\varphi=a\varphi$，$a$为常数；显然可撤去$\|\varphi\|=1$的限制，故对$\mathfrak M$中所有$f$都有$Uf=af$。于是恒有$UEg=aEg$，即$UE=aE$；但$U=UE$，故$U=aE$。$U,\allowbreak E$都定号且非零，所以$a>0$；不改变$U$的本质，就可以令$U=E$。

只要适当选取$U_0$，这个$U$确实能对每个$R$，即每一组$\varphi_1,\allowbreak \varphi_2,\allowbreak \ldots,\allowbreak \psi_1,\allowbreak \psi_2,\allowbreak \ldots$，解决上述问题。取$U_0=1$便有
\[
\sum_n(U_0\varphi_n,\varphi_n)P_{[\varphi_n]}
=\sum_n(\varphi_n,\varphi_n)P_{[\varphi_n]}
=\sum_nP_{[\varphi_n]}=P_{\mathfrak M}=E=U.
\]
于是，按照上面提出的思路，$U=E$得到了依据。如果再要求$U_0$具有普遍性，即不依赖于$E,\allowbreak R$，那么$U_0$本身也可确定：$U_0=1$满足要求，而且只有它满足。因为
\begin{align*}
(U\varphi_m,\varphi_m)&=(E\varphi_m,\varphi_m)=(\varphi_m,\varphi_m)=1,\\
(U\varphi_m,\varphi_m)&=\sum_n(U_0\varphi_n,\varphi_n)(P_{[\varphi_n]}\varphi_m,\varphi_m)\\
&=\sum_n(U_0\varphi_n,\varphi_n)|(\varphi_n,\varphi_m)|^2
=(U_0\varphi_m,\varphi_m).
\end{align*}
所以$(U_0\varphi_m,\allowbreak \varphi_m)=1$。$\mathfrak M$中每个满足$\|\varphi\|=1$的$\varphi$都可以选为$\varphi_1$，故$(U_0\varphi,\allowbreak \varphi)=1$，由此对$\mathfrak M$中所有$f$得$(U_0f,\allowbreak f)=(f,\allowbreak f)$。由于$\mathfrak M$任意，这对一切$f$都成立，所以$U_0=1$。

给定两个未必同时可判定的性质$\mathfrak E,\allowbreak \mathfrak F$。如果刚刚确认系统$\boldsymbol{S}$具备$\mathfrak E$，随即检验时又具备$\mathfrak F$的概率，根据上述结果为
\[
\operatorname{Spur}(EF)=\sum(EF),
\]
其中$E,\allowbreak F$是$\mathfrak E,\allowbreak \mathfrak F$的算符。第一个公式来自$U=E$，第二个来自$E^2=E$、$F^2=F$及 \hyperref[sec:II.11]{II.11}。附带指出，这些概率\hypertarget{origpage-183}{}\phantomsection\label{origpage-183}是相对的，其中视$\mathfrak E$为固定，$\mathfrak F$为可变。如果$\operatorname{Spur}(E)=\sum(E)$，即$\mathfrak M$的维数，是有限的，就可除以它而归一化。

也可以用物理量代替$\mathfrak E,\allowbreak \mathfrak F$。设$\mathfrak R_1,\allowbreak \ldots,\allowbreak \mathfrak R_j$是一组同时可测的量，$\mathfrak S_1,\allowbreak \ldots,\allowbreak \mathfrak S_l$是另一组；相应的算符为$R_1,\allowbreak \ldots,\allowbreak R_j,\allowbreak S_1,\allowbreak \ldots,\allowbreak S_l$，相应的单位分解为$E_1(\lambda),\allowbreak \ldots,\allowbreak E_j(\lambda),\allowbreak F_1(\lambda),\allowbreak \ldots,\allowbreak F_l(\lambda)$。设
\begin{gather*}
I_1:\lambda'_1<\lambda\leq\lambda''_1,\quad\ldots,\quad I_j:\lambda'_j<\lambda\leq\lambda''_j,\\
J_1:\mu'_1<\lambda\leq\mu''_1,\quad\ldots,\quad J_l:\mu'_l<\lambda\leq\mu''_l
\end{gather*}
为若干区间，并令
\begin{gather*}
E_1(I_1)=E_1(\lambda''_1)-E_1(\lambda'_1),\quad\ldots,\quad E_j(I_j)=E_j(\lambda''_j)-E_j(\lambda'_j),\\
F_1(J_1)=F_1(\mu''_1)-F_1(\mu'_1),\quad\ldots,\quad F_l(J_l)=F_l(\mu''_l)-F_l(\mu'_l).
\end{gather*}
问题是：在$\boldsymbol{S}$上测量了$\mathfrak R_1,\allowbreak \ldots,\allowbreak \mathfrak R_j$，测得值分别落在$I_1,\allowbreak \ldots,\allowbreak I_j$内；紧接着测量$\mathfrak S_1,\allowbreak \ldots,\allowbreak \mathfrak S_l$，值分别落在$J_1,\allowbreak \ldots,\allowbreak J_l$内的概率是多少？显然应令$E=E_1(I_1)\cdots E_j(I_j)$、$F=F_1(J_1)\cdots F_l(J_l)$（参见 \hyperref[post:III.5.epsilon]{$\varepsilon)$}、\hyperref[post:III.5.zeta]{$\zeta)$}），所求概率就是
\begin{align*}
&\operatorname{Spur}\bigl(E_1(I_1)\cdots E_j(I_j)F_1(J_1)\cdots F_l(J_l)\bigr)\\
&\qquad=\sum\bigl(E_1(I_1)\cdots E_j(I_j)F_1(J_1)\cdots F_l(J_l)\bigr).
\end{align*}

最后再强调一般初始系综$U_0=1$的意义。我们在测量$\mathfrak R$时将它分为两部分，从而得到$U$。如果不把它分开，也就是在它的每个元素上都测量$\mathfrak R$，然后把全部元素重新汇成一个系综，结果就仍是$U_0=1$。这一点可以直接验证，也可以取$E=1$来证明：此时$\mu_1,\allowbreak \mu_2,\allowbreak \ldots$和$\psi_1,\allowbreak \psi_2,\allowbreak \ldots$消失，$\lambda_1,\allowbreak \lambda_2,\allowbreak \ldots$与$\varphi_1,\allowbreak \varphi_2,\allowbreak \ldots$构成完全系。因此，测量$\mathfrak R$虽然可能改变单个元素，所有这些改变却必定恰好互相抵消，因为整个系综没有变化。而且，这一性质恰能刻画$U_0=1$。事实上，若对每个完全规范正交系$\varphi_1,\allowbreak \varphi_2,\allowbreak \ldots$都有
\[
U_0=\sum_{n=1}^{\infty}(U_0\varphi_n,\varphi_n)P_{[\varphi_n]}
\]
则$U_0$与$P_{[\varphi_1]}$对易，因为每个$P_{[\varphi_n]}$都与它对易。于是$U_0$与所有满足$\|\varphi\|=1$的$P_{[\varphi]}$都对易。因而
\[
U_0\varphi=U_0P_{[\varphi]}\varphi=P_{[\varphi]}U_0\varphi=(U_0\varphi,\varphi)\varphi,
\]
即$\varphi$是$U_0$的本征函数。由此推出$U_0=1$，与前面由相应关系推出$U=E$完全相同，只须把$\mathfrak M,\allowbreak E$换成$\mathfrak R,\allowbreak 1$。

因此，在$U_0=1$中，所有可能状态达到最大限度的平衡，任何测量干预都不能改变这种平衡。对每个\hypertarget{origpage-184}{}\phantomsection\label{origpage-184}完全规范正交系$\varphi_1,\allowbreak \varphi_2,\allowbreak \ldots$都有
\[
U_0=1=\sum_{n=1}^{\infty}P_{[\varphi_n]},
\]
即所有状态$\varphi_1,\allowbreak \varphi_2,\allowbreak \ldots$的$1:1:\cdots$混合。由此可见，$U_0=1$对应于旧量子论中惯用的热力学假设：“\textkaishu{所有非简并量子轨道的先验等概率性}”。在下面几节的热力学讨论中，它也将起重要作用。

\chapter{一般性讨论}\label{chap:V}
\section{测量与可逆性}\label{sec:V.1}
在统计算符为$U$的混合态中测量一个算符为$R$的物理量$\mathfrak R$，会发生什么？只要这个问题本身容许唯一答案，我们就能够作出预言。这里是说，在系综的每个元素上都作测量，再把经过测量的各个系统重新汇成一个系综。

第一，设$R$有单重的纯点谱，$\varphi_1,\allowbreak \varphi_2,\allowbreak \ldots$是其本征函数构成的完全规范正交系，$\lambda_1,\allowbreak \lambda_2,\allowbreak \ldots$为相应本征值，按假设彼此不同。测量之后，原系综中占$(U\varphi_m,\allowbreak \varphi_m)$的那一部分，$\mathfrak R$取值$\lambda_n$（$n=1,\allowbreak 2,\allowbreak \ldots$）。\tnote[。][因][原]{原文此处写作$(U\varphi_m,\allowbreak \varphi_m)$；按上下文，下标$m$应为$n$。正文保留原印。}因此，这一部分构成一个系综，其中$\mathfrak R$必定取值$\lambda_n$，参见 \hyperref[sec:IV.3]{IV.3} 的 \hyperref[post:IV.3.M]{M}；它就是状态$\varphi_n$，其正确归一化的统计算符为$P_{[\varphi_n]}$。把它们汇集起来，所得混合态的统计算符便是
\[
U'=\sum_{n=1}^{\infty}(U\varphi_n,\varphi_n)P_{[\varphi_n]}.
\]

第二，设$R$有纯点谱，$\varphi_1,\allowbreak \varphi_2,\allowbreak \ldots$与$\lambda_1,\allowbreak \lambda_2,\allowbreak \ldots$的意义仍如上，但谱不是单重的，即有些$\lambda_n$相等。此时$\mathfrak R$的测量过程并不唯一，正如 \hyperref[sec:IV.3]{IV.3} 中$\mathfrak E$的情形。事实上，取互不相同的实数$\mu_1,\allowbreak \mu_2,\allowbreak \ldots$，以$S$表示本征函数为$\varphi_1,\allowbreak \varphi_2,\allowbreak \ldots$、相应本征值为$\mu_1,\allowbreak \mu_2,\allowbreak \ldots$的算符，以$\mathfrak S$表示相应物理量。若函数$F(x)$满足
\[
F(\lambda_n)=\mu_n\qquad(n=1,2,\ldots),
\]
就有$F(S)=R$，即$F(\mathfrak S)=\mathfrak R$。\tnote[。][因][原]{原书印作$F(\lambda_n)=\mu_n$。若要使下述$F(S)=R$成立，鉴于$S$的本征值为$\mu_n$、$R$的本征值为$\lambda_n$，应为$F(\mu_n)=\lambda_n$。这里仍保留原印公式。}因此，测量$\mathfrak S$也应算作测量$\mathfrak R$，它把$U$变为前述$U'$。\hypertarget{origpage-185}{}\phantomsection\label{origpage-185}这个$U'$虽不依赖完全任意的$\mu_1,\allowbreak \mu_2,\allowbreak \ldots$，却依赖$\varphi_1,\allowbreak \varphi_2,\allowbreak \ldots$；而后者由于$R$的本征值有重数，并非唯一确定。我们在 \hyperref[sec:IV.2]{IV.2} 接续 \hyperref[sec:II.8]{II.8} 的论述，已说明$\varphi_1,\allowbreak \varphi_2,\allowbreak \ldots$所能确定到的程度。把$\lambda_1,\allowbreak \lambda_2,\allowbreak \ldots$中的不同值记为$\lambda',\allowbreak \lambda'',\allowbreak \ldots$，把满足$Rf=\lambda'f$、$Rf=\lambda''f$等的$f$所成集合分别记为$\mathfrak M_{\lambda'},\allowbreak \mathfrak M_{\lambda''},\allowbreak \ldots$。任取分别张成这些流形的规范正交系$\chi'_1,\allowbreak \chi'_2,\allowbreak \ldots$、$\chi''_1,\allowbreak \chi''_2,\allowbreak \ldots$等，则$\chi'_1,\allowbreak \chi'_2,\allowbreak \ldots,\allowbreak \chi''_1,\allowbreak \chi''_2,\allowbreak \ldots,\allowbreak \ldots$就是最一般的$\varphi_1,\allowbreak \varphi_2,\allowbreak \ldots$。所以，随$\mathfrak S$的选择，即随实际测量装置的选择，任何形如
\[
U'=\sum_n(U\chi'_n,\chi'_n)P_{[\chi'_n]}
+\sum_n(U\chi''_n,\chi''_n)P_{[\chi''_n]}+\cdots
\]
的$U'$都有可能出现。只有特殊情况下它才唯一。

再确定一下这种特殊情况。此时，每个加项必须分别唯一。也就是说，对于每个本征值$\lambda$，设$\mathfrak M_\lambda$为满足$Rf=\lambda f$的$f$所成集合，那么对任何张成$\mathfrak M_\lambda$的规范正交系，和$\sum_n(U\chi_n,\allowbreak \chi_n)P_{[\chi_n]}$都必须相同。把这个和记为$V$。沿用 \hyperref[sec:IV.3]{IV.3} 的论证，只须把那里的$U_0,\allowbreak U,\allowbreak \mathfrak M$换成$U,\allowbreak V,\allowbreak \mathfrak M_\lambda$，便知必有$V=c_\lambda P_{\mathfrak M_\lambda}$，$c_\lambda$是正常数；这等价于对$\mathfrak M_\lambda$中的一切$f$有$(Uf,\allowbreak f)=c_\lambda(f,\allowbreak f)$。这些$f$恰好是一切$g$所给出的$P_{\mathfrak M_\lambda}g$，故所需条件为
\[
(UP_{\mathfrak M_\lambda}g,P_{\mathfrak M_\lambda}g)
=c_\lambda(P_{\mathfrak M_\lambda}g,P_{\mathfrak M_\lambda}g),
\]
或$(P_{\mathfrak M_\lambda}UP_{\mathfrak M_\lambda}g,\allowbreak g)=c_\lambda(P_{\mathfrak M_\lambda}g,\allowbreak g)$，即
\[
P_{\mathfrak M_\lambda}UP_{\mathfrak M_\lambda}=c_\lambda P_{\mathfrak M_\lambda}
\]
对$R$的每个本征值$\lambda$成立。这一条件显然对$U$有很强的限制。如果它不成立，不同的$\mathfrak R$测量装置确实会把$U$变成不同的$U'$。不过，在 \hyperref[sec:V.4]{V.4} 中，我们还能从热力学出发，对一般的$\mathfrak R$测量结果作进一步说明。

第三，设$R$没有纯点谱。由 \hyperref[sec:III.3]{III.3}，或 \hyperref[sec:IV.3]{IV.3} 的判据 \hyperref[criterion:IV.3.1]{1}，它根本不能绝对精确地测量。如在那里所讨论的，有限精度的$\mathfrak R$测量，等价于对具有纯点谱的物理量作测量。

对物质系统还有另一种干预方式，由含时Schr\"odinger微分方程给出。实验或测量是不连续、非因果、瞬时起作用的，而这个方程描述的是：已知总能量时，系统怎样随时间连续地、依因果律变化。对于状态$\varphi$，若能量算符为$\mathsf{H}$，方程为
\begin{equation}\tag{$Z_1$}\label{eq:V.1.Z1}
\frac{\partial}{\partial t}\varphi_t=-\frac{2\pi i}{h}\mathsf{H}\varphi_t.
\end{equation}

\hypertarget{origpage-186}{}\phantomsection\label{origpage-186}对状态$\varphi_t$的统计算符$U_t=P_{[\varphi_t]}$，这表示
\begin{align*}
\left(\frac{\partial U_t}{\partial t}\right)f
&=\frac{\partial}{\partial t}(U_tf)
=\frac{\partial}{\partial t}\bigl((f,\varphi_t)\varphi_t\bigr)\\
&=\left(f,\frac{\partial\varphi_t}{\partial t}\right)\varphi_t
+(f,\varphi_t)\frac{\partial\varphi_t}{\partial t}\\
&=-\left(f,\frac{2\pi i}{h}\mathsf{H}\varphi_t\right)\varphi_t
-(f,\varphi_t)\frac{2\pi i}{h}\mathsf{H}\varphi_t\\
&=\frac{2\pi i}{h}\bigl((\mathsf{H}f,\varphi_t)\varphi_t-(f,\varphi_t)\mathsf{H}\varphi_t\bigr)
=\frac{2\pi i}{h}(U_t\mathsf{H}-\mathsf{H}U_t)f.
\end{align*}
\begin{equation}\tag{$Z_2$}\label{eq:V.1.Z2}
\frac{\partial U_t}{\partial t}=\frac{2\pi i}{h}(U_t\mathsf{H}-\mathsf{H}U_t).
\end{equation}
若$U_t$不是一个状态，而是若干状态（例如$P_{[\varphi_t^{(1)}]},\allowbreak P_{[\varphi_t^{(2)}]},\allowbreak \ldots$分别以权重$w_1,\allowbreak w_2,\allowbreak \ldots$）的混合，则其变化必须与这些$P_{[\varphi_t^{(1)}]},\allowbreak P_{[\varphi_t^{(2)}]},\allowbreak \ldots$的变化相符。把相应的 \hyperref[eq:V.1.Z2]{$Z_2$} 方程相加，便知 \hyperref[eq:V.1.Z2]{$Z_2$} 对这些$U_t$也成立。而所有$U$都是这样的混合态，或是其极限情形（例如，每个迹有限的$U$都是这样的混合态），所以 \hyperref[eq:V.1.Z2]{$Z_2$} 具有一般有效性。

顺便说，在 \hyperref[eq:V.1.Z2]{$Z_2$} 中，正如在Schr\"odinger微分方程 \hyperref[eq:V.1.Z1]{$Z_1$} 中一样，$\mathsf{H}$也可以依赖$t$。若它不依赖$t$，就能把解显式写出。对于 \hyperref[eq:V.1.Z1]{$Z_1$}，我们已知
\begin{equation}\tag{$Z'_1$}\label{eq:V.1.Z1prime}
\varphi_t=e^{-\frac{2\pi i}{h}t\mathsf{H}}\varphi_0,
\end{equation}
对于 \hyperref[eq:V.1.Z2]{$Z_2$}，则为
\begin{equation}\tag{$Z'_2$}\label{eq:V.1.Z2prime}
U_t=e^{-\frac{2\pi i}{h}t\mathsf{H}}U_0e^{\frac{2\pi i}{h}t\mathsf{H}}.
\end{equation}
很容易验证它们确实是解，也很容易由其中一个推出另一个。给定初值$\varphi_0$或$U_0$后解唯一，这也很清楚，因为微分方程 \hyperref[eq:V.1.Z1]{$Z_1$}、\hyperref[eq:V.1.Z2]{$Z_2$} 都是关于$t$的一阶方程。

因此，对系统$\boldsymbol{S}$或系综$[\boldsymbol{S}_1,\allowbreak \ldots,\allowbreak \boldsymbol{S}_N]$，有两种根本不同的干预。第一种是测量造成的人为变化，用公式
\begin{equation}\tag{1}\label{eq:V.1.1}
U\longrightarrow U'=\sum_{n=1}^{\infty}(U\varphi_n,\varphi_n)P_{[\varphi_n]}
\end{equation}
表示，其中$\varphi_1,\allowbreak \varphi_2,\allowbreak \ldots$为完全规范正交系，参见上文。第二种是随时间自动发生的变化，用公式
\begin{equation}\tag{2}\label{eq:V.1.2}
U\longrightarrow U_t=e^{-\frac{2\pi i}{h}t\mathsf{H}}Ue^{\frac{2\pi i}{h}t\mathsf{H}}
\end{equation}\hypertarget{origpage-187}{}\phantomsection\label{origpage-187}表示，其中$\mathsf{H}$为能量算符，$t$为时间，$\mathsf{H}$不依赖$t$。若$\mathsf{H}$依赖$t$，就把所考察的时间区间分成小区间，使$\mathsf{H}$在每个小区间内不变、或只略有变化，再分别使用 \hyperref[eq:V.1.2]{2}；叠合起来便得最终结果。

现在必须仔细研究这两种干预，弄清它们的性质和相互关系。

首先值得注意的是，\hyperref[eq:V.1.2]{2} 容许$\mathsf{H}$按上述方式依赖时间。因此，按理说，只用 \hyperref[eq:V.1.2]{2} 就应足以描述测量所造成的干预。物理干预无非是在一段时间内，把某些能量耦合引入被考察的系统；也就是由观察者规定，使$\mathsf{H}$具有适当的时间依赖性。为什么测量还需要特殊过程 \hyperref[eq:V.1.1]{1} 呢？原因在于：测量时不能只考察系统$\boldsymbol{S}$本身，要通过计算追踪它与测量装置$\boldsymbol{M}$的相互作用，必须研究系统$\boldsymbol{S}+\boldsymbol{M}$。测量理论所陈述的是关于$\boldsymbol{S}+\boldsymbol{M}$的事情，因为它要说明$\boldsymbol{S}$的状态怎样与$\boldsymbol{M}$状态的某些性质相联系，也就是怎样与观察者读出的某些指针位置相联系。是否把观察者也算进$\boldsymbol{M}$，其实相当任意；还可以不用$\boldsymbol{S}$的状态与$\boldsymbol{M}$中指针位置的联系，而用它与观察者眼中、乃至脑中化学变化的联系，亦即与他“看见”或“统觉到”的东西的联系。这一点将在 \hyperref[sec:VI.1]{VI.1} 进一步研究。无论如何，\hyperref[eq:V.1.2]{2} 只能用于$\boldsymbol{S}+\boldsymbol{M}$；但必须证明，它对$\boldsymbol{S}$给出的结果，与直接把 \hyperref[eq:V.1.1]{1} 用于$\boldsymbol{S}$相同。只有做到这一点，才算在量子力学基础上获得了统一的物理世界图景。此问题留待 \hyperref[sec:VI.3]{VI.3} 讨论。

其次，关于 \hyperref[eq:V.1.1]{1}，我们曾反复强调，其中所指的测量必须是瞬时的，即在足够短的时间内完成，使 \hyperref[eq:V.1.2]{2} 造成的$U$的变化尚不可察觉。即使想据 \hyperref[eq:V.1.2]{2} 计算变化后的$U_t$来作修正，也无济于事。因为要使用任何一个$U_t$，都必须准确知道$t$，即测量发生的时刻；所以测量持续时间仍须很短。原则上说，这颇成问题：众所周知，在经典力学中，有一个与时间\hypertarget{origpage-188}{}\phantomsection\label{origpage-188}正则共轭的量，就是能量。\origNote[。][因]{180}因此，时间—能量这对正则共轭量，也应像直角坐标—动量这一对一样，满足不确定性关系。\origNote[。][狭]{181}狭义相对论也表明，二者之间必有深刻的类比：三个空间坐标与时间组成“四维矢量”，三个空间动量与能量也一样。如果有这样的不确定性关系，就意味着不能在很短时间内完成十分精确的能量测量；具体说来，测量误差$\varepsilon$与持续时间$\tau$应有形如
\[
\varepsilon\tau\sim h
\]
的关系。对这一问题作与 \hyperref[sec:III.4]{III.4} 中$p,\allowbreak q$情形类似的物理讨论，确实得到这一结果\origNote[果][。]{181}。不详述细节，只举光量子为例。根据Bohr频率关系，它的能量不确定度$\varepsilon$等于频率不确定度的$h$倍，即$h\Delta\nu$；而如原注 \hyperlink{note-137}{\getrefnumber{orig-note-137}} 所论，频率不确定度$\Delta\nu$至多只能缩小到持续时间倒数$1/\tau$的量级，所以$\varepsilon\gtrsim h/\tau$。要确定光量子在整个$\tau$区间内的单色性，测量就必须遍及这一整段时间。光量子的例子具有代表性，因为原子能级通常是根据相应谱线的频率确定的。能量既然如此，其他量$\mathfrak R$的测量精度也可能与持续时间相联系。那么，怎样为瞬时测量的假设辩护呢？

首先必须承认，这一反对意见触及量子力学的一项重大弱点，甚至是主要弱点：它具有非相对论性质，把时间$t$与三个空间坐标$x,\allowbreak y,\allowbreak z$区别对待，并预设客观的同时性概念。其他一切量，特别是由Lorentz变换与$t$紧密联系的$x,\allowbreak y,\allowbreak z$，都用算符表示；时间却像在经典力学中一样，只用普通的数值参数$t$表示。又如，由两个粒子组成的系统，其波函数依赖二者的$2\times3=6$个空间坐标，却只有一个时间$t$；但从Lorentz变换来看，本应希望有两个时间。因此，我们能够忽略自然规律对测量规定的最短持续时间，或许正与量子力学的非相对论性质有关。这算是一种解释，却不是令人满意的解释！

\hypertarget{origpage-189}{}\phantomsection\label{origpage-189}不过，细究之后可知，情况并没有那么糟。我们真正需要的不是$t$很小，而只是计算概率$(U\varphi_n,\allowbreak \varphi_n)$，进而构成$U'=\sum_{n=1}^{\infty}(U\varphi_n,\allowbreak \varphi_n)P_{[\varphi_n]}$时，从$U$出发或从某个$U_t=e^{-\frac{2\pi i}{h}t\mathsf{H}}Ue^{\frac{2\pi i}{h}t\mathsf{H}}$出发，差别很小。由于
\begin{align*}
(U_t\varphi_n,\varphi_n)
&=\left(e^{-\frac{2\pi i}{h}t\mathsf{H}}Ue^{\frac{2\pi i}{h}t\mathsf{H}}\varphi_n,\varphi_n\right)\\
&=\left(Ue^{\frac{2\pi i}{h}t\mathsf{H}}\varphi_n,e^{\frac{2\pi i}{h}t\mathsf{H}}\varphi_n\right),
\end{align*}
例如，可以用适当的扰动能来改变$\mathsf{H}$，使$e^{\frac{2\pi i}{h}t\mathsf{H}}\varphi_n$与$\varphi_n$只差一个模为$1$的常数因子。也就是说，让状态$\varphi_n$在 \hyperref[eq:V.1.2]{2} 作用下本质不变，成为定态；等价地，$\mathsf{H}\varphi_n$等于$\varphi_n$乘以一个实常数，即$\varphi_n$是$\mathsf{H}$的本征函数。乍看之下，这样改变能量算符$\mathsf{H}$，使$R$的本征函数成为定态，也就是成为$\mathsf{H}$的本征函数，从而使$R,\allowbreak \mathsf{H}$对易，似乎不大可信。但实际上并非如此，甚至可以说明，典型测量装置所求的正是对$\mathsf{H}$的这种改变。

每次测量的结尾，总是一个光量子或质点以某种能量沿某个方向飞走，由这些特征（即动量）显示测量结果；或者一个质点（例如刻度盘上的指针）静止下来，由它的直角坐标显示结果。对光量子，用 \hyperref[sec:III.6]{III.6} 的说法，指出哪一个$M_n$为$1$、其余为$0$，即给出所有$M_1,\allowbreak M_2,\allowbreak \ldots$的值，就等于给出待测量的值；对飞走的质点，则是给出三个动量分量$P^x,\allowbreak P^y,\allowbreak P^z$；对静止的质点，则是给出三个直角坐标$x,\allowbreak y,\allowbreak z$，或用其算符的名称，$Q^x,\allowbreak Q^y,\allowbreak Q^z$。但测量只有在光量子或质点真正“飞走”时才算完成：光不再有被吸收的危险，质点不再受势能偏转；或者，静止的质点确实保持静止，这要求很大的质量。\origNote[。][仅]{182}仅由不确定性关系也可得出后一要求：速度须接近零，因而速度的散布程度也须很小；但速度与质量的乘积——动量——却因坐标的散布很小而必须有很大的散布。指针确实是宏观对象，也就是极其庞大。能量算符$\mathsf{H}$对\hypertarget{origpage-190}{}\phantomsection\label{origpage-190}光量子而言，由 \hyperref[sec:III.6]{III.6} 为
\begin{align*}
&\sum_{n=1}^{\infty}h\varrho_nM_n\\
&\quad+\sum_{j=1}^{\infty}\sum_{n=1}^{\infty}w_{kj}^{n}
\Biggl(\sqrt{M_n+1}\begin{pmatrix}k\to j\\M_n\to M_n+1\end{pmatrix}\\
&\qquad\qquad{}+\sqrt{M_n}\begin{pmatrix}k\to j\\M_n\to M_n-1\end{pmatrix}\Biggr);
\end{align*}
对两个质点例子则为
\[
\frac{(P^x)^2+(P^y)^2+(P^z)^2}{2m}+V(Q^x,Q^y,Q^z),
\]
其中$m$为质量，$V$为势能。我们的判据分别要求：$w_{kj}^n$为零，$V$为常数，或$m$极大。而这些条件恰好使$M_n$，或$P^x,\allowbreak P^y,\allowbreak P^z$，或$Q^x,\allowbreak Q^y,\allowbreak Q^z$，与上述$\mathsf{H}$对易。

这一讨论最后还应提到：在理论物理学的其他问题中，也会用到把所研究的状态（这里是$\varphi_1,\allowbreak \varphi_2,\allowbreak \ldots$）变为定态的办法。物理化学的“思想实验”常常不得不假定，可以通过“毒化”中止化学反应；这种假设正属于此类。\origNote[。][]{183}

两种干预 \hyperref[eq:V.1.1]{1}、\hyperref[eq:V.1.2]{2} 根本不同。它们在形式上都唯一确定，因而都是因果的，这一点无关紧要。因为我们所考察的是系综的统计性质，任何变化，即使是统计性的变化，也会使概率和期望值按因果律改变，并不足为奇；引入统计系综和概率，本来就是为此。关键在于：\hyperref[eq:V.1.2]{2} 不会增加$U$中已有的统计不确定性，\hyperref[eq:V.1.1]{1} 却会。前者把状态变为状态：
\[
P_{[\varphi]}\quad\text{变为}\quad P_{[e^{-\frac{2\pi i}{h}t\mathsf{H}}\varphi]},
\]
而 \hyperref[eq:V.1.1]{1} 完全可能把状态变成混合态。从这个意义说，状态按 \hyperref[eq:V.1.1]{1} 的演变是统计性的，按 \hyperref[eq:V.1.2]{2} 的演变则是因果性的。

此外，固定$\mathsf{H},\allowbreak t$后，\hyperref[eq:V.1.2]{2} 无非是对所有$U$作酉变换：
\[
U_t=AUA^{-1},\qquad A=e^{-\frac{2\pi i}{h}t\mathsf{H}}\quad\text{为酉算符}.
\]
于是$Uf=g$就意味着$U_t(Af)=Ag$，即$U_t$由$U$经过希尔伯特空间的酉变换$A$而来。这是一种使我们一切几何基本概念保持不变的同构，参见 \hyperref[sec:I.4]{I.4} 所述原理。因此它可逆：只须把$A$换为$A^{-1}$。这是做得到的，因为$\mathsf{H},\allowbreak t$有很大的选择自由，$A,\allowbreak A^{-1}$可以视为完全任意的酉算符。因此，过程 \hyperref[eq:V.1.2]{2} 和经典\hypertarget{origpage-191}{}\phantomsection\label{origpage-191}力学一样，未能再现现实世界最本质、最显著的一项性质：不可逆性，即“未来”与“过去”两个时间方向的根本差别。

\hyperref[eq:V.1.1]{1} 的情况则根本不同。转变
\[
U\longrightarrow U'=\sum_{n=1}^{\infty}(U\varphi_n,\varphi_n)P_{[\varphi_n]}
\]
不能直接倒转。我们很快还会看到，它根本不可逆：一般说来，即使任意反复运用 \hyperref[eq:V.1.1]{1}、\hyperref[eq:V.1.2]{2} 两种过程，也无法从给定的$U'$回到它原来的$U$。

至此，就有理由引入热力学的观点了。\hyperref[eq:V.1.1]{1} 与 \hyperref[eq:V.1.2]{2} 的差别显然涉及可逆性问题，只有这种观点才能使我们正确理解二者的差别。

\section{热力学讨论}\label{sec:V.2}
我们将从两个不同角度研究量子力学系综的热力学。首先，假设热力学两条基本定律成立，也就是第一类和第二类永动机不可能存在，即能量原理与熵原理\origNote[理][，]{184}，再据此计算每个系综的熵。这里采用唯象热力学的通常方法；量子力学只在如下意义上起作用：热力学所考察的对象，即我们的系综及其统计算符$U$，遵守量子力学规律。两条基本定律本身则是假设成立，并不加以证明。随后，我们再证明它们在量子力学中成立。由于能量原理本来就成立，实际只剩熵原理：我们要说明，干预 \hyperref[eq:V.1.1]{1}、\hyperref[eq:V.1.2]{2} 绝不减少按第一种方法算出的熵。这个次序看似有些不自然，却自有理由：只有经过唯象热力学的讨论，才能掌握问题的全貌，而这是第二种讨论所必需的。

因此，先从唯象讨论开始；它还将使我们能够消解经典热力学中一个著名的悖论。这里先指出，我们的“思想实验”虽然很大胆，实际上不能实行，却不因此\hypertarget{origpage-192}{}\phantomsection\label{origpage-192}失去证明效力。在唯象热力学中，任何可以设想的过程，只要不违背基本定律，都能用来论证。

目标是确定统计算符为$U$的系综$[\boldsymbol{S}_1,\allowbreak \ldots,\allowbreak \boldsymbol{S}_N]$的熵，并假设$U$已正确归一化，即$\operatorname{Spur}U=1$。用经典统计力学的说法，这是一个Gibbs系综：统计与热力学所涉及的，不是单个极其复杂的力学系统中彼此作用的各组成部分——这种系统的自由度很多，而我们对这些自由度所知有限\origNote[限][—]{185}——而是由大量力学系统组成的系综。每个系统本身可以有（但不必有）多得无法把握的自由度，但各系统彼此完全分离，互不作用。\origNote[。][由]{186}由于$\boldsymbol{S}_1,\allowbreak \ldots,\allowbreak \boldsymbol{S}_N$完全分离，而且我们要用的是概率论中纯粹计数的方法，因此毫无疑问，应采用通常的概率算法。Bose--Einstein与Fermi--Dirac的不同计数法，在这里都不适用；它们适用于某些由不可区分而彼此作用的个体组成的系综，即分别适用于光量子，以及电子与质子，参见 \hyperref[sec:III.6]{III.6}，特别是原注 \hyperlink{note-147}{\getrefnumber{orig-note-147}}。

Einstein提出的、从热力学处理这类系综$[\boldsymbol{S}_1,\allowbreak \ldots,\allowbreak \boldsymbol{S}_N]$的方法如下。\origNote[。][把]{187}把每个系统$\boldsymbol{S}_1,\allowbreak \ldots,\allowbreak \boldsymbol{S}_N$分别装进一个盒子$\boldsymbol{K}_1,\allowbreak \ldots,\allowbreak \boldsymbol{K}_N$，盒壁隔绝一切作用的传递；各系统既然互不作用，这样做是允许的。每个盒子还须有很大的质量，使$\boldsymbol{S}_1,\allowbreak \ldots,\allowbreak \boldsymbol{S}_N$可能发生的状态变化，以及随之而来的能量即质量变化，对盒子的质量只产生很小影响。这也使下面思想实验中的盒子速度足够小，可以用非相对论方法计算。再把$\boldsymbol{K}_1,\allowbreak \ldots,\allowbreak \boldsymbol{K}_N$放进一个极大的盒子$\overline{\boldsymbol{K}}$中，其体积$\mathcal V$远大于小盒子总体积。为简单起见，设$\overline{\boldsymbol{K}}$中没有力场，特别是远离一切引力场，并大到可以忽略小盒子的质量所产生的作用。于是，可以把内含$\boldsymbol{S}_1,\allowbreak \ldots,\allowbreak \boldsymbol{S}_N$的盒子$\boldsymbol{K}_1,\allowbreak \ldots,\allowbreak \boldsymbol{K}_N$，视为封在大盒子$\overline{\boldsymbol{K}}$中的气体分子。若使$\overline{\boldsymbol{K}}$接触一个温度为$\mathbf T$的巨大热库，盒壁也会达到这一温度，其真正的分子便会作相应的布朗运动。因此，它们\hypertarget{origpage-193}{}\phantomsection\label{origpage-193}会把动量传给邻近的$\boldsymbol{K}_1,\allowbreak \ldots,\allowbreak \boldsymbol{K}_N$，使这些盒子运动起来，继而把动量传给其他盒子。不久，全部小盒子都会运动，在$\overline{\boldsymbol{K}}$的壁上与壁分子交换动量，并在内部互相碰撞、交换动量。只有当小盒子取得与温度为$\mathbf T$的壁分子布朗运动相平衡的速度分布时，运动才达到稳定的平衡；这就是温度为$\mathbf T$、以$\boldsymbol{K}_1,\allowbreak \ldots,\allowbreak \boldsymbol{K}_N$为“分子”的气体的Maxwell速度分布。\origNote[。][所]{188}所以可以说，$[\boldsymbol{S}_1,\allowbreak \ldots,\allowbreak \boldsymbol{S}_N]$气体已达到温度$\mathbf T$。为简便起见，以下把统计算符为$U$的系综称为$U$系综，相应气体称为$U$气体。

我们引入这种气体，是因为按定义，$U$系综与$V$系综的熵差，必须通过把前者可逆地变成后者来确定。\origNote[。][这]{189}这里$U,\allowbreak V$都是迹为$1$的定号算符，相应系综分别为$[\boldsymbol{S}_1,\allowbreak \ldots,\allowbreak \boldsymbol{S}_N]$与$[\boldsymbol{S}'_1,\allowbreak \ldots,\allowbreak \boldsymbol{S}'_N]$。借道$U$气体、$V$气体，最容易实现这种转变。我们断言：只要把两种气体放在同一温度$\mathbf T$下，它们的熵差就恰好等于两系综的熵差，而$\mathbf T$除此以外可以任意。如果$\mathbf T$极接近零，这显然可任意精确地成立，因为$U$系综与零温的$U$气体之间的区别消失了：后者的小盒子没有自身运动，静止的$\boldsymbol{K}_1,\allowbreak \ldots,\allowbreak \boldsymbol{K}_N,\allowbreak \overline{\boldsymbol{K}}$在场，与热力学无关（对$V$亦然）。所以，只须证明：温度$\mathbf T$改变一定数值时，$U$气体与$V$气体的熵变化相同。气体从$\mathbf T_1$升温到$\mathbf T_2$的熵变，只取决于它的量热状态方程，更确切地说，取决于热容。\origNote[。][当]{190}当然，当$\mathbf T_1$如这里一样必须取得接近零时，不能假定气体是理想气体。\origNote[。][不]{191}不过，两种气体（$U$与$V$）的状态方程和热容肯定相同，因为从气体动理论看，起主导作用的是盒子$\boldsymbol{K}_1,\allowbreak \ldots,\allowbreak \boldsymbol{K}_N$，它们把内部的$\boldsymbol{S}_1,\allowbreak \ldots,\allowbreak \boldsymbol{S}_N$或$\boldsymbol{S}'_1,\allowbreak \ldots,\allowbreak \boldsymbol{S}'_N$完全遮蔽了。因此，在这一升温过程中，$U,\allowbreak V$之别并不显现，两种熵差便如所断言的那样相等。以下\hypertarget{origpage-194}{}\phantomsection\label{origpage-194}只比较$U$气体与$V$气体，并把$\mathbf T$取得足够高，使二者可视为理想气体。\origNote[。][这]{192}这样，气体的动理学行为就完全掌握在我们手中，可以着手真正的任务：把$U$气体可逆地变为$V$气体。与此前的过程不同，这时还须对盒内的$\boldsymbol{S}_1,\allowbreak \ldots,\allowbreak \boldsymbol{S}_N$施加作用，也就是必须“打开”盒子。

先证明所有状态$U=P_{[\varphi]}$都有相同的熵：可以不消耗也不产生热能，而把$P_{[\varphi]}$系综可逆地变成$P_{[\psi]}$系综。当然，如果二者的能量期望值不同，就必须消耗或产生机械能，参见原注 \hyperlink{note-185}{\getrefnumber{orig-note-185}}。为此甚至不必引入刚才的气体；在零温下，直接对系综本身也能完成这一转变。证明之后，就可以而且将要规定$U$系综的熵的零点，使所有状态的熵均为零。

顺便说，所述$P_{[\varphi]}$到$P_{[\psi]}$的转变甚至不必可逆。若不可逆，熵差须不小于原注 \hyperlink{note-185}{\getrefnumber{orig-note-185}} 给出的表达式（参见该注），因而须$\geq0$；交换$P_{[\varphi]},\allowbreak P_{[\psi]}$，同样得熵差$\leq0$，所以仍然只能为零。

最简单的办法本来是采用含时Schr\"odinger微分方程，也就是过程 \hyperref[eq:V.1.2]{2}；为此须找到能量算符$\mathsf{H}$及一个时间$t$，使酉算符$e^{-\frac{2\pi i}{h}t\mathsf{H}}$把$\varphi$变成$\psi$。那么经过$t$秒，$P_{[\varphi]}$就自动变为$P_{[\psi]}$，过程甚至可逆，也完全不涉及热，参见 \hyperref[sec:V.1]{V.1}。不过，我们宁愿避免假定能量算符$\mathsf{H}$可能具有什么形式，而只用过程 \hyperref[eq:V.1.1]{1}，即类似测量的干预。最简单的设想是，在系综$P_{[\varphi]}$上测量物理量$\mathfrak R$，其算符$R$有纯点谱，所有本征值$\lambda_1,\allowbreak \lambda_2,\allowbreak \ldots$都是单重的，而本征函数$\psi_1,\allowbreak \psi_2,\allowbreak \ldots$中包含$\psi$，例如$\psi_1=\psi$。测量把状态$\varphi$变成这些状态的混合，因此也产生了$\psi$。然而，这并不合适：$\psi_1=\psi$只以概率$|(\varphi,\allowbreak \psi)|^2$出现，另有$1-|(\varphi,\allowbreak \psi)|^2$的部分转入其他状态；若$\varphi,\allowbreak \psi$正交，这就是整个系综。不过，与之类似的\hypertarget{origpage-195}{}\phantomsection\label{origpage-195}方法仍能达到目的：接连作大量不同的测量，就能把$P_{[\varphi]}$变为一个与$P_{[\psi]}$相差任意小的系综。如上所述，这些操作都是（或至少可能是）不可逆的，并不重要。

设$\varphi,\allowbreak \psi$正交；否则可取同时正交于二者的$\chi$，使$\|\chi\|=1$，先由$\varphi$转到$\chi$，再由$\chi$转到$\psi$。设$k=1,\allowbreak 2,\allowbreak \ldots$是一个留待选定的数，令
\[
\psi^{(\nu)}=\cos\frac{\pi\nu}{2k}\,\varphi+\sin\frac{\pi\nu}{2k}\,\psi
\qquad(\nu=0,1,\ldots,k).
\]
显然$\psi^{(0)}=\varphi$、$\psi^{(k)}=\psi$，且恒有$\|\psi^{(\nu)}\|=1$。把每个$\psi^{(\nu)}$（$\nu=1,\allowbreak \ldots,\allowbreak k$）扩充为完全规范正交系$\psi_1^{(\nu)},\allowbreak \psi_2^{(\nu)},\allowbreak \ldots$，其中$\psi_1^{(\nu)}=\psi^{(\nu)}$。取一个纯点谱算符$R^{(\nu)}$，其互不相同的本征值为$\lambda_1^{(\nu)},\allowbreak \lambda_2^{(\nu)},\allowbreak \ldots$，本征函数为$\psi_1^{(\nu)},\allowbreak \psi_2^{(\nu)},\allowbreak \ldots$；相应物理量为$\mathfrak R^{(\nu)}$。还应注意
\begin{align*}
(\psi^{(\nu-1)},\psi^{(\nu)})
&=\cos\frac{\pi(\nu-1)}{2k}\cos\frac{\pi\nu}{2k}
+\sin\frac{\pi(\nu-1)}{2k}\sin\frac{\pi\nu}{2k}\\
&=\cos\left(\frac{\pi\nu}{2k}-\frac{\pi(\nu-1)}{2k}\right)=\cos\frac\pi{2k}.
\end{align*}

现在，在$U^{(0)}=P_{[\psi^{(0)}]}=P_{[\varphi]}$系综上测量$\mathfrak R^{(1)}$，得到$U^{(1)}$；再在$U^{(1)}$上测量$\mathfrak R^{(2)}$，得到$U^{(2)}$；依此类推，最后在$U^{(k-1)}$上测量$\mathfrak R^{(k)}$，得到$U^{(k)}$。$k$足够大时，$U^{(k)}$任意接近$P_{[\psi^{(k)}]}=P_{[\psi]}$；直观解释如下：在$\psi^{(\nu-1)}$上测量$\mathfrak R^{(\nu)}$，有
\[
|(\psi^{(\nu-1)},\psi^{(\nu)})|^2=\left(\cos\frac\pi{2k}\right)^2
\]
这一比例转入$\psi^{(\nu)}$。所以接连测量$\mathfrak R^{(1)},\allowbreak \mathfrak R^{(2)},\allowbreak \ldots,\allowbreak \mathfrak R^{(k)}$，至少有$(\cos(\pi/2k))^{2k}$的部分，由$\psi^{(0)}=\varphi$经过$\psi^{(1)},\allowbreak \psi^{(2)},\allowbreak \ldots,\allowbreak \psi^{(k-1)}$而转入$\psi=\psi^{(k)}$。当$k\to\infty$时，$(\cos(\pi/2k))^{2k}\to1$，所以只要$k$足够大，所得几乎全是$\psi$。严格证明如下。过程 \hyperref[eq:V.1.1]{1} 不改变迹，故由$\operatorname{Spur}U^{(0)}=\operatorname{Spur}P_{[\varphi]}=1$，有$\operatorname{Spur}U^{(1)}=\operatorname{Spur}U^{(2)}=\cdots=\operatorname{Spur}U^{(j)}=1$。另一方面，
\begin{align*}
(U^{(\nu)}f,f)
&=\sum_n(U^{(\nu-1)}\psi_n^{(\nu)},\psi_n^{(\nu)})(P_{[\psi_n^{(\nu)}]}f,f)\\
&=\sum_n(U^{(\nu-1)}\psi_n^{(\nu)},\psi_n^{(\nu)})|(\psi_n^{(\nu)},f)|^2.
\end{align*}
对$\nu=1,\allowbreak \ldots,\allowbreak k-1$取$f=\psi_1^{(\nu+1)}=\psi^{(\nu+1)}$，对$\nu=k$取$f=\psi_1^{(k)}=\psi^{(k)}=\psi$，就有
\begin{align*}
(U^{(\nu)}\psi^{(\nu+1)},\psi^{(\nu+1)})
&\geq(U^{(\nu-1)}\psi^{(\nu)},\psi^{(\nu)})|(\psi^{(\nu)},\psi^{(\nu+1)})|^2\\
&=\left(\cos\frac\pi{2k}\right)^2(U^{(\nu-1)}\psi^{(\nu)},\psi^{(\nu)}),\\
(U^{(k)}\psi^{(k)},\psi^{(k)})&=(U^{(k-1)}\psi^{(k)},\psi^{(k)}).
\end{align*}\hypertarget{origpage-196}{}\phantomsection\label{origpage-196}再结合
\[
(U^{(0)}\psi^{(1)},\psi^{(1)})=(P_{[\psi^{(0)}]}\psi^{(1)},\psi^{(1)})
=|(\psi^{(0)},\psi^{(1)})|^2=\left(\cos\frac\pi{2k}\right)^2,
\]
得到\tnote[到][][这]{这里原书印作$U^{(j)}$；前面的论证和后面的极限涉及的则是$U^{(k)}$。}
\[
(U^{(j)}\psi,\psi)\geq\left(\cos\frac\pi{2k}\right)^{2k}.
\]
由于$\operatorname{Spur}U^{(k)}=1$，且$k\to\infty$时$(\cos(\pi/2k))^{2k}\to1$，可用 \hyperref[sec:II.11]{II.11} 所得结果，得知$U^{(k)}$收敛到$P_{[\psi]}$。目的至此达到。

接着须讨论：唯象热力学“思想实验”的一项主要工具——所谓半透壁——在多大程度上也能用于量子力学系统？唯象热力学容许如下命题：若I、II是同一系统$\boldsymbol{S}$的两个不同状态，就可以假定存在一种壁，对I完全透过，对II完全不透。\origNote[。][可]{193}可以说，这从热力学上定义了两个状态何为不同，因而也定义了何为相同。在量子力学中，这种假定有多大适用范围呢？

先证明：若$\varphi_1,\allowbreak \varphi_2,\allowbreak \ldots,\allowbreak \psi_1,\allowbreak \psi_2,\allowbreak \ldots$为规范正交系，就存在一种半透壁，让处于任一$\varphi_1,\allowbreak \varphi_2,\allowbreak \ldots$状态的系统$\boldsymbol{S}$无阻碍地透过，让处于任一$\psi_1,\allowbreak \psi_2,\allowbreak \ldots$状态的系统原状反射。其他状态的系统则可能在碰撞壁时改变。

可设$\varphi_1,\allowbreak \varphi_2,\allowbreak \ldots,\allowbreak \psi_1,\allowbreak \psi_2,\allowbreak \ldots$完全；否则可补上$\chi_1,\allowbreak \chi_2,\allowbreak \ldots$使其完全，再把它们归入$\varphi_1,\allowbreak \varphi_2,\allowbreak \ldots$一组。选取纯点谱算符$R$，所有本征值$\lambda_1,\allowbreak \lambda_2,\allowbreak \ldots,\allowbreak \mu_1,\allowbreak \mu_2,\allowbreak \ldots$都是单重的，相应本征函数为$\varphi_1,\allowbreak \varphi_2,\allowbreak \ldots,\allowbreak \psi_1,\allowbreak \psi_2,\allowbreak \ldots$，且$\lambda_n<0$、$\mu_n>0$。设相应物理量为$\mathfrak R$。在壁上安装许多窗口，各作如下安排：我们气体的每个“分子”$\boldsymbol{K}_1,\allowbreak \ldots,\allowbreak \boldsymbol{K}_N$（现在仍考察温度$\mathbf T>0$的$U$气体）到达窗口时被暂时留住、打开；测量其中$\boldsymbol{S}_1$、……或$\boldsymbol{S}_N$的$\mathfrak R$，再关上盒子。若测得值$<0$，让整个盒子连同内容物保持动量穿过窗口；若$>0$，则使其反射。这一装置显然达到了预期目的。剩下只须讨论：碰撞后装置中留下什么变化，以及它在\hypertarget{origpage-197}{}\phantomsection\label{origpage-197}多大程度上与热力学中所谓“Maxwell妖”类似或不同。\origNote[。][]{194}

第一点是：测量$\mathfrak R$可能改变$\boldsymbol{S}$的状态，进而改变其能量期望值。因此，按第一定律，这一机械能差额必须由测量装置供给或吸收。例如，装置中可装一根弹簧，让它松开、拉紧，或采用类似办法。由于这是纯自动运行的测量机构，而且转换的只是机械能而非热能，所以这里肯定没有熵变；目前重要的只有这一点。若$\boldsymbol{S}$处于$\varphi_1,\allowbreak \varphi_2,\allowbreak \ldots,\allowbreak \psi_1,\allowbreak \psi_2,\allowbreak \ldots$之一，测量$\mathfrak R$根本不改变它，测量装置中便也不留下补偿性的变化。

第二点更棘手。这个装置与“Maxwell妖”颇为相似；后者是一种让右边来的分子透过、左边来的分子反射的半透壁。把它插入装有气体的容器中央，不久全部气体就会聚在左半边，即体积减半而不付出熵的代价。这意味着气体的熵无补偿地减少，所以按第二定律，这种壁不可能存在。不过，我们的半透壁与这种热力学所不容的壁，有着本质差别：我们只依据“分子”$\boldsymbol{K}_1,\allowbreak \ldots,\allowbreak \boldsymbol{K}_N$的内部性质，即盒内$\boldsymbol{S}_1$、……或$\boldsymbol{S}_N$的状态，而不依据外部性质，例如从左来还是从右来。关键恰恰在此。L.~Szilard的研究澄清了半透壁、“Maxwell妖”，以及一般而言“\textkaishu{有智力的生物干预热力学系统}”的作用，据此可以作详尽分析。这里不能深入这些问题，何况读者在原注 \hyperlink{note-194}{\getrefnumber{orig-note-194}} 所引著作中就能找到详细论述。

以上特别表明：系统$\boldsymbol{S}$的两个状态$\varphi,\allowbreak \psi$，若彼此正交，肯定能用半透壁分开。现在证明反命题：若它们不正交，假定存在这样的半透壁，就与第二定律矛盾。因此，能用半透壁分离的充要条件是$(\varphi,\allowbreak \psi)=0$，而不是经典理论中的$\varphi\ne\psi$（这里以$\varphi,\allowbreak \psi$代替上面的I、II）。这样就消解了经典形式热力学中的一个老悖论，即使用\hypertarget{origpage-198}{}\phantomsection\label{origpage-198}半透壁时令人不快的不连续性：状态不论相差多么微小，仍能百分之百分开；而完全相同的状态，却全然不能分开。现在则有连续过渡：只有$(\varphi,\allowbreak \psi)=0$时才能完全分开，随$(\varphi,\allowbreak \psi)$增大，分离就越来越不完全；直到其模达到最大值$|(\varphi,\allowbreak \psi)|=1$时，状态$\varphi,\allowbreak \psi$才相同，分离才完全不可能。因为$\|\varphi\|=\|\psi\|=1$，故$|(\varphi,\allowbreak \psi)|=1$蕴含$\varphi=c\psi$，其中$c$为常数，$|c|=1$。

要进行上述论证，须先借用本节最后才会得到的结果，即$U$系综的熵。当然，推导这一结果时，不会用到眼下所证的结论。

假设存在能分开$\varphi,\allowbreak \psi$的半透壁，我们要由此推出$(\varphi,\allowbreak \psi)=0$。考察$\tfrac12(P_{[\varphi]}+P_{[\psi]})$气体，即由$N/2$个处于$\varphi$状态的系统和$N/2$个处于$\psi$状态的系统组成的气体；该算符的迹为$1$。选取$\mathcal V$（即$\overline{\boldsymbol{K}}$）和$\mathbf T$，使气体理想。设$\overline{\boldsymbol{K}}$的纵截面为图 \ref{fig:3} 所示的$12341$。\tnote[。][从][图]{图中$W$标记热库；$a,\allowbreak b$是隔板位置的字母标记，箭头表示隔板移动的方向。}从一端$aa$插入半透壁，再推进到中央$bb$；让另一端$23$接触温度为$\mathbf T$的大热库$\boldsymbol{W}$，保持气体温度不变。
\begin{figure}[tbp]
\centering
\includegraphics[width=0.55\linewidth]{figures/figure-3.png}
\caption{}\label{fig:3}
\end{figure}
在这个过程中，$\varphi$分子不受影响，$\psi$分子却被推到$\overline{\boldsymbol{K}}$的右半部，即$bb$与$23$之间。也就是说，$\tfrac12(P_{[\varphi]}+P_{[\psi]})$气体是$P_{[\varphi]}$气体与$P_{[\psi]}$气体的$1:1$混合；前者不变，后者则等温压缩到一半体积。由理想气体状态方程知，所作机械功为$(N/2)\varkappa\mathbf T\ln2$，其中$N/2$是$P_{[\psi]}$气体的分子数，$\varkappa$为Boltzmann常数。\origNote[。][因]{195}因为等温过程不改变气体能量\origNote[量][，]{196}，这部分能量被热库$\boldsymbol{W}$吸收，所以热库的熵增加$\mathbf Q/\mathbf T=N\varkappa\cdot\tfrac12\ln2$，参见原注 \hyperlink{note-186}{\getrefnumber{orig-note-186}}。

过程结束后，$bb$左侧是原$P_{[\varphi]}$气体的一半，即$N/4$个分子；右侧则是\hypertarget{origpage-199}{}\phantomsection\label{origpage-199}另一半$P_{[\varphi]}$气体（$N/4$个分子），以及全部$P_{[\psi]}$气体（$N/2$个分子），合计是含$3N/4$个分子的$\tfrac13P_{[\varphi]}+\tfrac23P_{[\psi]}$气体。把前者压缩至体积$\mathcal V/4$，后者膨胀至$3\mathcal V/4$，相应的机械功仍由热库$\boldsymbol{W}$吸收或供给，分别为$(N/4)\varkappa\mathbf T\ln2$与$(3N/4)\varkappa\mathbf T\ln(3/2)$，参见原注 \hyperlink{note-191}{\getrefnumber{orig-note-191}}。因此，热库的熵变化量分别为$N\varkappa\cdot\tfrac14\ln2$与$-N\varkappa\cdot\tfrac34\ln(3/2)$，合计
\[
N\varkappa\left(\tfrac12\ln2+\tfrac14\ln2-\tfrac34\ln\tfrac32\right)
=N\varkappa\cdot\tfrac34\ln\tfrac43.
\]

最后，体积$\mathcal V/4$内有$N/4$个分子的$P_{[\varphi]}$气体，体积$3\mathcal V/4$内有$3N/4$个分子的$\tfrac13P_{[\varphi]}+\tfrac23P_{[\psi]}$气体。原来则是体积$\mathcal V$内有$N$个分子的$\tfrac12P_{[\varphi]}+\tfrac12P_{[\psi]}$气体；也可把它看成两份同种气体，分别有$N/4$和$3N/4$个分子，占据$\mathcal V/4$和$3\mathcal V/4$。因此，整个过程的改变是：体积$\mathcal V/4$内的$N/4$个分子，从$\tfrac12P_{[\varphi]}+\tfrac12P_{[\psi]}$气体变成$P_{[\varphi]}$气体；体积$3\mathcal V/4$内的$3N/4$个分子，从同一种原气体变成$\tfrac13P_{[\varphi]}+\tfrac23P_{[\psi]}$气体；而$\boldsymbol{W}$的熵增加$N\varkappa\cdot\tfrac34\ln(4/3)$。过程可逆，所以总熵增必须为零，即两份气体的熵变必须恰好抵消$\boldsymbol{W}$的熵变。因此还须知道气体的熵变。

我们稍后将看到：若把同体积、同温度的$P_{[\chi]}$气体的熵算作零，则含$M$个分子的$U$气体的熵为$-M\varkappa\operatorname{Spur}(U\ln U)$，参见前文。若$U$有纯点谱，本征值为$w_1,\allowbreak w_2,\allowbreak \ldots$，这就是$-M\varkappa\sum_{n=1}^{\infty}w_n\ln w_n$，其中$x=0$时约定$x\ln x=0$。容易算出，算符
\[
P_{[\varphi]},\qquad\tfrac12P_{[\varphi]}+\tfrac12P_{[\psi]},\qquad
\tfrac13P_{[\varphi]}+\tfrac23P_{[\psi]}
\]
的本征值依次为
\[
1,0;\qquad \frac{1+\alpha}2,\frac{1-\alpha}2,0;\qquad
\frac{3+\sqrt{1+8\alpha^2}}6,\frac{3-\sqrt{1+8\alpha^2}}6,0,
\]
其中$\alpha=|(\varphi,\allowbreak \psi)|$，故$0\leq\alpha\leq1$；零本征值总是无穷重，其他本征值则是单重的。\origNote[。][因]{197}因此，气体的熵变化量为
\begin{align*}
&-\frac N4\varkappa\cdot0
-\frac{3N}4\varkappa\Biggl[
\frac{3+\sqrt{1+8\alpha^2}}6\ln\frac{3+\sqrt{1+8\alpha^2}}6\\
&\qquad\qquad\qquad+\frac{3-\sqrt{1+8\alpha^2}}6\ln\frac{3-\sqrt{1+8\alpha^2}}6\Biggr]\\
&\qquad+N\varkappa\left(\frac{1+\alpha}2\ln\frac{1+\alpha}2
+\frac{1-\alpha}2\ln\frac{1-\alpha}2\right).
\end{align*}\hypertarget{origpage-200}{}\phantomsection\label{origpage-200}这与$\boldsymbol{W}$的熵增$N\varkappa\cdot\tfrac34\ln(4/3)$相加应得零。除以$N\varkappa/4$，有
\begin{align*}
&-\frac{3+\sqrt{1+8\alpha^2}}2\ln\frac{3+\sqrt{1+8\alpha^2}}6
-\frac{3-\sqrt{1+8\alpha^2}}2\ln\frac{3-\sqrt{1+8\alpha^2}}6\\
&\qquad+2(1+\alpha)\ln\frac{1+\alpha}2
+2(1-\alpha)\ln\frac{1-\alpha}2+3\ln\frac43=0,
\end{align*}
其中$0\leq\alpha\leq1$。

容易看出，$\alpha$从$0$增到$1$时，左端不断增加\origNote[加][，]{198}，从$0$增至$3\ln(4/3)$。故必有$\alpha=0$；否则，所述过程的逆过程就会减少熵，违反第二定律。于是$(\varphi,\allowbreak \psi)=0$得证。

完成这些准备后，现在确定体积为$\mathcal V$、温度为$\mathbf T$、含$N$个分子的$U$气体的熵；更确切地说，确定它比同条件的$P_{[\varphi]}$气体多出的熵。根据此前的说明及上述零点规定，这就是含$N$个单个系统的$U$系综的熵。如上，设$\operatorname{Spur}U=1$。
\begin{figure}[tbp]
\centering
\includegraphics[width=0.55\linewidth]{figures/figure-4.png}
\caption{}\label{fig:4}
\end{figure}
我们知道，$U$有纯点谱$w_1,\allowbreak w_2,\allowbreak \ldots$，其中$w_1\geq0,\allowbreak w_2\geq0,\allowbreak \ldots$，$w_1+w_2+\cdots=1$。相应本征函数设为$\varphi_1,\allowbreak \varphi_2,\allowbreak \ldots$，则$U=\sum_{n=1}^{\infty}w_nP_{[\varphi_n]}$，参见 \hyperref[sec:IV.3]{IV.3}。所以$U$气体是$P_{[\varphi_1]},\allowbreak P_{[\varphi_2]},\allowbreak \ldots$气体的混合，各分别含$w_1N,\allowbreak w_2N,\allowbreak \ldots$个分子，均占同一体积$\mathcal V$。仍选取$\mathbf T,\allowbreak \mathcal V$使所有气体都理想，并设$\overline{\boldsymbol{K}}$为矩形截面。为把$\varphi_1,\allowbreak \varphi_2,\allowbreak \ldots$分子互相分离，作如下可逆操作，参见图 \ref{fig:4}。\tnote[。][在][图]{图中字母$a,\allowbreak b,\allowbreak c,\allowbreak d$分别标记各隔板的位置；箭头表示移动方向。}在盒子$\overline{\boldsymbol{K}}$（$23452$）旁接上同样大小的矩形盒$\overline{\boldsymbol{K}}'$（$12561$）。把公共壁$25$换成两片紧贴的壁：一片固定在$25$处，半透，对$\varphi_1$透过，对$\varphi_2,\allowbreak \varphi_3,\allowbreak \ldots$反射；另一片$bb$可移动，是普通的完全不透壁。再在紧邻$34$的$dd$处插入另一半透壁，对$\varphi_2,\allowbreak \varphi_3,\allowbreak \ldots$透过，对$\varphi_1$反射。然后保持$bb,\allowbreak dd$间距不变，分别把它们移至$aa,\allowbreak cc$，即紧邻$16,\allowbreak 25$处。这样，\hypertarget{origpage-201}{}\phantomsection\label{origpage-201}$\varphi_2,\allowbreak \varphi_3,\allowbreak \ldots$分子全不受影响，$\varphi_1$分子却被迫留在两片移动壁$bb,\allowbreak dd$之间。由于两壁间距恒定，不必克服气体压力作功，因而也没有热的产生。最后，把$25,\allowbreak cc$换成一片固定的完全不透壁$25$，并撤去$aa$，便恢复了盒子$\overline{\boldsymbol{K}},\allowbreak \overline{\boldsymbol{K}}'$。现在全部$\varphi_1$分子都在$\overline{\boldsymbol{K}}'$中；我们就这样把它们从$\overline{\boldsymbol{K}}$可逆地“抽取”到同样大的$\overline{\boldsymbol{K}}'$，既不作功，也不产生热，更不改变温度。\origNote[。][]{199}

同样，把$\varphi_2,\allowbreak \varphi_3,\allowbreak \ldots$分子分别“抽取”到等大的盒子$\overline{\boldsymbol{K}}'',\allowbreak \overline{\boldsymbol{K}}''',\allowbreak \ldots$中，最后得到$P_{[\varphi_1]},\allowbreak P_{[\varphi_2]},\allowbreak \ldots$气体，各含$w_1N,\allowbreak w_2N,\allowbreak \ldots$个分子，每份体积都是$\mathcal V$。再分别等温压缩至$w_1\mathcal V,\allowbreak w_2\mathcal V,\allowbreak \ldots$。为了补偿，须从温度为$\mathbf T$的大热库向它们供给热量$w_1N\varkappa\mathbf T\ln w_1$、$w_2N\varkappa\mathbf T\ln w_2$等，以保证过程可逆。这些热量全是负的；对各气体所作压缩机械功恰好等于它们的相反数，参见原注 \hyperlink{note-191}{\getrefnumber{orig-note-191}}。所以，气体在这些过程中的熵变化量为$\sum_{n=1}^{\infty}w_nN\varkappa\ln w_n$。最后，把$P_{[\varphi_1]},\allowbreak P_{[\varphi_2]},\allowbreak \ldots$气体全变成$P_{[\varphi]}$气体，$\varphi$为任取的状态。现在只有$P_{[\varphi]}$气体，分别有$w_1N,\allowbreak w_2N,\allowbreak \ldots$个分子，占据体积$w_1\mathcal V,\allowbreak w_2\mathcal V,\allowbreak \ldots$。由于它们完全相同，密度也同为$N:\mathcal V$，便可把它们混合，且仍是可逆的；由$\sum_{n=1}^{\infty}w_n=1$，得到体积$\mathcal V$内含$N$个分子的$P_{[\varphi]}$气体。

这样就完成了所需的可逆转变。气体的熵变化量为
\[
N\varkappa\sum_{n=1}^{\infty}w_n\ln w_n\text{，}
\]
而最终的熵按规定为零，所以原先的熵是
\[
-N\varkappa\sum_{n=1}^{\infty}w_n\ln w_n.
\]
由于$U$的本征函数$\varphi_1,\allowbreak \varphi_2,\allowbreak \ldots$分别对应本征值$w_1,\allowbreak w_2,\allowbreak \ldots$，$U\ln U$便有相同本征函数，而本征值为$w_1\ln w_1,\allowbreak w_2\ln w_2,\allowbreak \ldots$。所以$\operatorname{Spur}(U\ln U)=\sum_{n=1}^{\infty}w_n\ln w_n$。注意$0\leq w_n\leq1$，因此$w_n\ln w_n\leq0$，仅当$w_n=0,\allowbreak 1$时为零。$w_n=0$时必须规定$w_n\ln w_n=0$，因为我们论证中无需考虑为零的$w_n$；连续性也给出同样的规定。

\hypertarget{origpage-202}{}\phantomsection\label{origpage-202}因此，由$N$个单个系统组成的$U$系综，其熵已确定为
\[
-N\varkappa\operatorname{Spur}(U\ln U).
\]
由刚才关于$w_n\ln w_n$的说明，它恒$\geq0$。要使它等于零，全部$w_n$必须是$0$或$1$。由于$\operatorname{Spur}U=1$，这时恰有一个$w_n=1$，其余全零，即$U=P_{[\varphi_n]}$。所以：$U=P_{[\varphi]}$，也就是状态，熵为零；其他一切混合态的熵均大于零。

\section{可逆性与平衡问题}\label{sec:V.3}
现在可以证明 \hyperref[sec:V.1]{V.1} 所断言的测量过程不可逆性。例如，设$U$为一个状态，即$U=P_{[\varphi]}$，测量物理量$\mathfrak R$，其算符$R$的本征函数为$\varphi_1,\allowbreak \varphi_2,\allowbreak \ldots$；它使$U$变为系综
\[
U'=\sum_{n=1}^{\infty}(P_{[\varphi]}\varphi_n,\varphi_n)P_{[\varphi_n]}
=\sum_{n=1}^{\infty}|(\varphi,\varphi_n)|^2P_{[\varphi_n]}.
\]
如果$U'$不是状态，熵便增加了，因为$U$的熵为零，而$U'$的熵大于零，所以过程不可逆。如果$U'$是状态，由于$\varphi_n$是它的本征函数，它必为某个$P_{[\varphi_{\bar n}]}$；这要求所有$|(\varphi,\allowbreak \varphi_n)|^2$除一个以外全为零，而这一个等于$1$。即$\varphi$与一切$n\ne\bar n$的$\varphi_n$正交。于是$\varphi=c\varphi_{\bar n}$，$|c|=1$，故$P_{[\varphi]}=P_{[\varphi_{\bar n}]}$，$U=U'$。因此，对一个状态所作的测量都不可逆，唯一例外是该状态为被测物理量的本征函数，即此量在给定状态中有确定值；而这时测量根本不改变状态。可见，非因果行为也与某些热力学伴随现象有确定的联系。

下面在完全一般的情形下讨论：过程 \hyperref[eq:V.1.1]{1}，即$U\to U'=\sum_{n=1}^{\infty}(U\varphi_n,\allowbreak \varphi_n)P_{[\varphi_n]}$，何时使熵增加。

$U$的熵为$-N\varkappa\operatorname{Spur}(U\ln U)$。若它的本征值为$w_1,\allowbreak w_2,\allowbreak \ldots$，本征函数为$\psi_1,\allowbreak \psi_2,\allowbreak \ldots$，则熵等于
\[
-N\varkappa\sum_{n=1}^{\infty}w_n\ln w_n
=-N\varkappa\sum_{n=1}^{\infty}(U\psi_n,\psi_n)\ln(U\psi_n,\psi_n).
\]
$U'$的本征值为$(U\varphi_1,\allowbreak \varphi_1),\allowbreak (U\varphi_2,\allowbreak \varphi_2),\allowbreak \ldots$，故其熵为
\[
-N\varkappa\sum_{n=1}^{\infty}(U\varphi_n,\varphi_n)\ln(U\varphi_n,\varphi_n).
\]
因此，$U$的熵与$U'$的熵的大小关系$\lesseqgtr$，取决于
\begin{equation}\tag{*}\label{eq:V.3.star}
\sum_{n=1}^{\infty}(U\psi_n,\psi_n)\ln(U\psi_n,\psi_n)
\gtreqless\sum_{n=1}^{\infty}(U\varphi_n,\varphi_n)\ln(U\varphi_n,\varphi_n)
\end{equation}
中的相应关系。先证明 \hyperref[eq:V.3.star]{*} 中总有$\geq$，即过程$U\to U'$不会减少熵。热力学上这\hypertarget{origpage-203}{}\phantomsection\label{origpage-203}当然明白，但为后面的用途，还须有纯数学证明。做法是固定$U$，从而固定$\psi_1,\allowbreak \psi_2,\allowbreak \ldots$，让$\varphi_1,\allowbreak \varphi_2,\allowbreak \ldots$遍历一切完全规范正交系。

由于连续性，可先只考察如下$\varphi_1,\allowbreak \varphi_2,\allowbreak \ldots$：其中只有有限个$\varphi_n$不同于相应$\psi_n$。例如，设$n>M$时恒有$\varphi_n=\psi_n$，则$n\leq M$的$\varphi_n$都是$n\leq M$的$\psi_n$的线性组合，反过来也一样。即
\[
\psi_m=\sum_{n=1}^Mx_{mn}\varphi_n\qquad(m=1,\ldots,M),
\]
其中$M$维矩阵$\{x_{mn}\}$显然是酉矩阵。\tnote[。][有][从]{从这里起，原书在本证明中把$\varphi$与$\psi$的作用同前文互换，改以$\varphi_n$为$U$的本征函数。以下保留原书记号。后一句恒等式还把求和上限印为$N$，并印作$w=1,\allowbreak \ldots,\allowbreak M$，而非$m=1,\allowbreak \ldots,\allowbreak M$。}有$(U\varphi_m,\allowbreak \varphi_m)=w_m$，并且容易算出$(U\psi_m,\allowbreak \psi_m)=\sum_{n=1}^Nw_n|x_{mn}|^2$，$w=1,\allowbreak \ldots,\allowbreak M$。所以须证明
\[
\sum_{m=1}^Mw_m\ln w_m\geq
\sum_{m=1}^M\left(\sum_{n=1}^Mw_n|x_{mn}|^2\right)
\ln\left(\sum_{n=1}^Mw_n|x_{mn}|^2\right).
\]
右端是$M^2$个有界变量$x_{mn}$的连续函数，因此在$\{x_{mn}\}$为酉矩阵的约束下有最大值，而且能取到。左端恰是$x_{mn}=1$（$m=n$）、$x_{mn}=0$（$m\ne n$）时的值，所以须证明，在这一组$x_{mn}$处取得最大值。

设$x^0_{mn}$（$m,\allowbreak n=1,\allowbreak \ldots,\allowbreak M$）是一组使最大值达到的数值。把矩阵$\{x^0_{mn}\}$乘以酉矩阵
\[
\begin{pmatrix}
\alpha&\beta&0&\cdots&0\\
-\overline\beta&\overline\alpha&0&\cdots&0\\
0&0&1&\cdots&0\\
\vdots&\vdots&\vdots&\ddots&\vdots\\
0&0&0&\cdots&1
\end{pmatrix},\qquad |\alpha|^2+|\beta|^2=1,
\]
便得另一酉矩阵$\{x'_{mn}\}$，也是允许的数值组。取$\alpha=\sqrt{1-\varepsilon^2}$，$\beta=\theta\varepsilon$，其中$\varepsilon$为实数、$|\theta|=1$。令$\varepsilon$很小，以下只精确保留常数项、$\varepsilon$项和$\varepsilon^2$项，略去$\varepsilon^3,\allowbreak \varepsilon^4,\allowbreak \ldots$。于是$\alpha\approx1-\tfrac12\varepsilon^2$，新矩阵中有
\begin{align*}
 x'_{1n}&\approx(1-\tfrac12\varepsilon^2)x^0_{1n}+\theta\varepsilon x^0_{2n},\\
 x'_{2n}&\approx-\overline\theta\varepsilon x^0_{1n}+(1-\tfrac12\varepsilon^2)x^0_{2n},\\
 x'_{mn}&=x^0_{mn}\qquad(m\geq3).
\end{align*}\hypertarget{origpage-204}{}\phantomsection\label{origpage-204}进而
\begin{align*}
\sum_{n=1}^Mw_n|x'_{1n}|^2
&\approx\sum_{n=1}^Mw_n|x^0_{1n}|^2
+\sum_{n=1}^M2w_n\operatorname{Re}(\overline\theta x^0_{1n}\overline{x^0_{2n}})\varepsilon\\
&\quad+\sum_{n=1}^Mw_n(-|x^0_{1n}|^2+|x^0_{2n}|^2)\varepsilon^2,\\
\sum_{n=1}^Mw_n|x'_{2n}|^2
&\approx\sum_{n=1}^Mw_n|x^0_{2n}|^2
-\sum_{n=1}^M2w_n\operatorname{Re}(\overline\theta x^0_{1n}\overline{x^0_{2n}})\varepsilon\\
&\quad-\sum_{n=1}^Mw_n(-|x^0_{1n}|^2+|x^0_{2n}|^2)\varepsilon^2,\\
\sum_{n=1}^Mw_n|x'_{mn}|^2&=\sum_{n=1}^Mw_n|x^0_{mn}|^2\qquad(m\geq3).
\end{align*}
把这些表达式代入$f(x)=x\ln x$，注意$f'(x)=\ln x+1$、$f''(x)=1/x$，再全部相加，得到
\begin{align*}
&\sum_{m=1}^M\left(\sum_{n=1}^Mw_n|x'_{mn}|^2\right)\ln\left(\sum_{n=1}^Mw_n|x'_{mn}|^2\right)\\
&\quad\approx\sum_{m=1}^M\left(\sum_{n=1}^Mw_n|x^0_{mn}|^2\right)\ln\left(\sum_{n=1}^Mw_n|x^0_{mn}|^2\right)\\
&\qquad+\left[\ln\left(\sum_{n=1}^Mw_n|x^0_{1n}|^2\right)-\ln\left(\sum_{n=1}^Mw_n|x^0_{2n}|^2\right)\right]\\
&\qquad\qquad\cdot\sum_{n=1}^M2w_n\operatorname{Re}(\overline\theta x^0_{1n}\overline{x^0_{2n}})\varepsilon\\
&\qquad+\Biggl\{-\left[\ln\left(\sum_{n=1}^Mw_n|x^0_{1n}|^2\right)-\ln\left(\sum_{n=1}^Mw_n|x^0_{2n}|^2\right)\right]\\
&\qquad\qquad\cdot\left[\sum_{n=1}^Mw_n|x^0_{1n}|^2-\sum_{n=1}^Mw_n|x^0_{2n}|^2\right]\\
&\qquad\qquad+\frac12\left[\frac1{\sum_{n=1}^Mw_n|x^0_{1n}|^2}+\frac1{\sum_{n=1}^Mw_n|x^0_{2n}|^2}\right]\\
&\qquad\qquad\cdot\left[\sum_{n=1}^M2\operatorname{Re}(\overline\theta x^0_{1n}\overline{x^0_{2n}})\right]^2\Biggr\}\varepsilon^2.
\end{align*}
右端第一项要是最大值，$\varepsilon$的系数就必须为零，$\varepsilon^2$的系数必须$\leq0$。\tnote[。][前][本]{本展开式最后一个平方和中，原书漏掉了前面一阶项中出现的权重$w_n$。公式仍按原书保留。}前一个系数有两个因子：
\[
\ln\left(\sum_{n=1}^Mw_n|x^0_{1n}|^2\right)-\ln\left(\sum_{n=1}^Mw_n|x^0_{2n}|^2\right),
\qquad\sum_{n=1}^M2w_n\operatorname{Re}(\overline\theta x^0_{1n}\overline{x^0_{2n}}).
\]
若第一个因子为零，$\varepsilon^2$系数中的第一项便为零；这一项总是$\leq0$。第二项显然总是$\geq0$，所以要使整个系数$\leq0$，第二项也必须为零。这就是说，$\sum_{n=1}^M2w_n\operatorname{Re}(\overline\theta x^0_{1n}\overline{x^0_{2n}})=0$。因此，$\varepsilon$系数的第二个因子无论如何都等于零，亦可写成
\[
2\operatorname{Re}\left(\overline\theta\sum_{n=1}^Mw_nx^0_{1n}\overline{x^0_{2n}}\right)=0.
\]
适当选取$\theta$，可把它变成和的绝对值，所以这个和必须为零：
\[
\sum_{n=1}^Mw_nx^0_{1n}\overline{x^0_{2n}}=0.
\]

\hypertarget{origpage-205}{}\phantomsection\label{origpage-205}由于$1,\allowbreak 2$可换成$1,\allowbreak \ldots,\allowbreak M$中任意两个不同的$k,\allowbreak j$，有
\[
\sum_{n=1}^Mw_nx^0_{kn}\overline{x^0_{jn}}=0\qquad(k\ne j).
\]
即矩阵为$\{x^0_{mn}\}$的酉坐标变换，把对角元为$w_1,\allowbreak \ldots,\allowbreak w_n$的对角矩阵再变成一个对角矩阵。对角元就是矩阵的乘子（或本征值），坐标变换不改变它们，至多重新排列。变换前它们是$w_m$（$m=1,\allowbreak \ldots,\allowbreak M$），变换后是$\sum_{n=1}^Mw_n|x^0_{mn}|^2$（$m=1,\allowbreak \ldots,\allowbreak N$），所以两个和
\[
\sum_{n=1}^Mw_n\ln w_n,\qquad
\sum_{m=1}^M\left(\sum_{n=1}^Mw_n|x^0_{mn}|^2\right)\ln\left(\sum_{n=1}^Mw_n|x^0_{mn}|^2\right)
\]
相等。也就是说，在$x_{mn}=1$（$m=n$）、$x_{mn}=0$（$m\ne n$）处也取得最大值，所言得证。

再确定 \hyperref[eq:V.3.star]{*} 中何时成立等号。如果等号成立，那么令$\chi_1,\allowbreak \chi_2,\allowbreak \ldots$遍历一切完全规范正交系时，
\[
\sum_{n=1}^{\infty}(U\chi_n,\chi_n)\ln(U\chi_n,\chi_n)
\]
不只在$\chi_n=\varphi_n$（$n=1,\allowbreak 2,\allowbreak \ldots$）处达到最大值，也在$\chi_n=\psi_n$（$n=1,\allowbreak 2,\allowbreak \ldots$）处达到。特别地，只把前$M$个$\psi_n$相互作酉变换，而$n>M$时仍取$\chi_n=\psi_n$，也有这一性质。令$u_{mn}=(U\psi_m,\allowbreak \psi_n)$（$m,\allowbreak n=1,\allowbreak \ldots,\allowbreak M$），设这个有限的、同样厄米且定号的矩阵$\{u_{mn}\}$的本征值为$v_1,\allowbreak \ldots,\allowbreak v_N$，$\{\alpha_{mn}\}$（$m,\allowbreak n=1,\allowbreak \ldots,\allowbreak M$）为将它化成对角形的矩阵。这一变换把$\psi_1,\allowbreak \ldots,\allowbreak \psi_M$变成$\omega_1,\allowbreak \ldots,\allowbreak \omega_M$，即
\[
\psi_m=\sum_{n=1}^M\alpha_{mn}\omega_n\qquad(m=1,\ldots,M).
\]
于是$U\omega_n=v_n\omega_n$，故
\[
(U\omega_m,\omega_n)=\begin{cases}v_n,&m=n,\\0,&m\ne n.\end{cases}
\]
若$\xi_m=\sum_{n=1}^Mx_{mn}\omega_n$（$m=1,\allowbreak \ldots,\allowbreak M$），其中$\{x_{mn}\}$也为酉矩阵，则$(U\xi_k,\allowbreak \xi_j)=\sum_{n=1}^Mv_nx_{kn}\overline{x_{jn}}$。根据对$\psi_1,\allowbreak \ldots,\allowbreak \psi_N$的假设，
\[
\sum_{n=1}^M\left(\sum_{n=1}^Mv_n|x_{mn}|^2\right)\ln\left(\sum_{n=1}^Mv_n|x_{mn}|^2\right)
\]
在$x_{mn}=\alpha_{mn}$处取得最大值。于是，由上面的证明得$\sum_{n=1}^Mv_n\alpha_{kn}\overline{\alpha_{jn}}=0$（$k\ne j$），即$(U\psi_k,\allowbreak \psi_j)=0$（$k\ne j$，$k,\allowbreak j=1,\allowbreak \ldots,\allowbreak M$）。

\hypertarget{origpage-206}{}\phantomsection\label{origpage-206}这对所有$M$都须成立，因此$U\psi_k$与所有$k\ne j$的$\psi_j$正交，故等于$w'_k\psi_k$，$w'_k$为常数。于是$\psi_1,\allowbreak \psi_2,\allowbreak \ldots$是$U$的本征函数，相应本征值为$w'_1,\allowbreak w'_2,\allowbreak \ldots$（即$w_1,\allowbreak w_2,\allowbreak \ldots$的一种排列）。在这种情况下，
\[
U'=\sum_{n=1}^{\infty}(U\psi_n,\psi_n)P_{[\psi_n]}
=\sum_{n=1}^{\infty}w'_nP_{[\psi_n]}=U.
\]
所以我们得到：

过程 \hyperref[eq:V.1.1]{1}，即$U\to U'=\sum_{n=1}^{\infty}(U\psi_n,\allowbreak \psi_n)P_{[\psi_n]}$（$\psi_1,\allowbreak \psi_2,\allowbreak \ldots$为被测物理量$\mathfrak R$的算符$R$的本征函数），从不减少熵；甚至总会增加熵，唯一例外是所有$\psi_1,\allowbreak \psi_2,\allowbreak \ldots$都是$U$的本征函数，而这时正有$U=U'$。此时$U$也与$R$对易，这同样是刻画上述情形的条件，因为它等价于存在公共本征函数系$\psi_1,\allowbreak \psi_2,\allowbreak \ldots$，参见 \hyperref[sec:II.10]{II.10}。

因此，只要过程 \hyperref[eq:V.1.1]{1} 造成任何改变，它就是不可逆的。

现在按照 \hyperref[sec:V.2]{V.2} 中计划的第二步，不借助唯象热力学，来处理过程 \hyperref[eq:V.1.1]{1}、\hyperref[eq:V.1.2]{2} 的可逆性。所需的数学方法我们已经有了：如果第二定律成立，熵必为$-N\varkappa\operatorname{Spur}(U\ln U)$，而它在任何 \hyperref[eq:V.1.1]{1}、\hyperref[eq:V.1.2]{2} 过程中都不能减小。因此，现在把$-N\varkappa\operatorname{Spur}(U\ln U)$只作为一个计算量，撇开熵的解释，求它在 \hyperref[eq:V.1.1]{1}、\hyperref[eq:V.1.2]{2} 中的变化。\origNote[。][]{200}

在 \hyperref[eq:V.1.2]{2} 中，$U$变为$U_t=e^{-\frac{2\pi i}{h}t\mathsf{H}}Ue^{\frac{2\pi i}{h}t\mathsf{H}}$；若以$A$记酉算符$e^{-\frac{2\pi i}{h}t\mathsf{H}}$，就是$U\to U_t=AUA^{-1}$。由于$A$的酉性，$f\to Af$是希尔伯特空间到自身的同构，把任一算符$P$变为$APA^{-1}$，所以一般有$F(APA^{-1})=AF(P)A^{-1}$。于是$U_t\ln U_t=A\cdot U\ln U\cdot A^{-1}$，故$\operatorname{Spur}(U_t\ln U_t)=\operatorname{Spur}(U\ln U)$。因此，我们的量$-N\varkappa\operatorname{Spur}(U\ln U)$在 \hyperref[eq:V.1.2]{2} 中不变。在 \hyperref[eq:V.1.1]{1} 中的情况刚才也已查明，而且未援引第二定律：如果$U$发生变化，即$U\ne U'$，这个量就增加；如果$U$不变，即$U=U'$，或$\psi_1,\allowbreak \psi_2,\allowbreak \ldots$是$U$的本征函数，或$U,\allowbreak R$对易，它当然也不变。因此，对任意多个、任意次序的 \hyperref[eq:V.1.1]{1}、\hyperref[eq:V.1.2]{2} 过程组成的干预，若每个 \hyperref[eq:V.1.1]{1} 都无效，即没有造成改变，$-N\varkappa\operatorname{Spur}(U\ln U)$就不变；在其余一切情形中，它都增加。\hypertarget{origpage-207}{}\phantomsection\label{origpage-207}所以，若只考虑干预 \hyperref[eq:V.1.1]{1}、\hyperref[eq:V.1.2]{2}，每个造成了改变的过程 \hyperref[eq:V.1.1]{1} 都不可逆。

顺便说，也有比$-\operatorname{Spur}(U\ln U)$更简单的量，在 \hyperref[eq:V.1.1]{1} 中不增加，在 \hyperref[eq:V.1.2]{2} 中不变，例如$U$的最大本征值。在 \hyperref[eq:V.1.2]{2} 中，它同$U$的所有本征值一样不变；在 \hyperref[eq:V.1.1]{1} 中，$U$的本征值$w_1,\allowbreak w_2,\allowbreak \ldots$变为$U'$的本征值$\sum_{n=1}^{\infty}w_n|x_{1n}|^2$、$\sum_{n=1}^{\infty}w_n|x_{2n}|^2$等，参见本节前面的讨论。由矩阵$\{x_{mn}\}$的酉性，$\sum_{n=1}^{\infty}|x_{1n}|^2=1$、$\sum_{n=1}^{\infty}|x_{2n}|^2=1$等等，因此这些数全不超过最大的$w_n$。最大的$w_n$确实存在，因为各$w_n\geq0$，而$\sum_{n=1}^{\infty}w_n=1$蕴含$w_n\to0$。现在，完全可以改变$U$，使$-\operatorname{Spur}(U\ln U)=-\sum_{n=1}^{\infty}w_n\ln w_n$不变，而最大的$w_n$减小。这说明，有些转变为唯象热力学所容许，因而确实能用我们的气体过程实现，却绝不可能仅靠接连运用 \hyperref[eq:V.1.1]{1}、\hyperref[eq:V.1.2]{2} 产生。引入气体方法，可见不可避免。

也可以不用$-\operatorname{Spur}(U\ln U)$，而考察适当函数$F(x)$所给出的$\operatorname{Spur}(F(U))$。只要它具备前面实际用到的$F(x)=-x\ln x$的那些性质，就能像那里一样证明：$U\ne U'$时，它在 \hyperref[eq:V.1.1]{1} 中增加；$U=U'$时，以及在 \hyperref[eq:V.1.2]{2} 中，它自然不变。这些性质是$F''(x)<0$，以及$F'(x)$单调递减，而后者又从前者推出。所以，在不依靠热力学的不可逆性讨论中，只要$F(x)$是严格凹函数，即在$0\leq x\leq1$上$F''(x)<0$（$U$的本征值全在其中），任何$\operatorname{Spur}(F(U))$都可以使用。

最后还须证明：混合两个系综$U,\allowbreak V$，例如按$\alpha:\beta$（$\alpha>0$、$\beta>0$、$\alpha+\beta=1$）混合，也不会使熵减小，即
\[
-\operatorname{Spur}\bigl((\alpha U+\beta V)\ln(\alpha U+\beta V)\bigr)
\geq-\alpha\operatorname{Spur}(U\ln U)-\beta\operatorname{Spur}(V\ln V).
\]
把$-x\ln x$换成任意严格凹函数$F(x)$，此式仍成立。\tnote[。][证][原]{原书称这些函数为“向上凸”，并明确要求$F''(x)<0$；按现今常用术语，这就是凹函数。}证明留给读者。

现在要寻找能量给定时的定常平衡混合态，也就是熵最大的混合态。所谓能量给定，自然是说能量的期望值\hypertarget{origpage-208}{}\phantomsection\label{origpage-208}给定；按照原注 \hyperlink{note-184}{\getrefnumber{orig-note-184}} 所述研究统计系综热力学的方法，只能采用这一表述。因此，只容许$U$满足$\operatorname{Spur}U=1$、$\operatorname{Spur}(U\mathsf{H})=\mathbf E$的混合态，其中$\mathsf{H}$为能量算符，$\mathbf E$为给定的能量期望值；在这些约束下，使$-N\varkappa\operatorname{Spur}(U\ln U)$最大。再作一项简化假设：$\mathsf{H}$有纯点谱，例如本征值$W_1,\allowbreak W_2,\allowbreak \ldots$，本征函数$\varphi_1,\allowbreak \varphi_2,\allowbreak \ldots$；本征值可以有重数。

取一个物理量$\mathfrak R$，其算符$R$的本征函数为$\varphi_1,\allowbreak \varphi_2,\allowbreak \ldots$，本征值则彼此不同。测量$\mathfrak R$，根据 \hyperref[eq:V.1.2]{2}\tnote[2][，][这]{这里原书引用过程2，但下述测量变换实际是过程 \hyperref[eq:V.1.1]{1}。}，会使$U$变成$U'=\sum_{n=1}^{\infty}(U\varphi_n,\allowbreak \varphi_n)P_{[\varphi_n]}$。除非$U=U'$，$-N\varkappa\operatorname{Spur}(U\ln U)$都会增加，而$\operatorname{Spur}(U)$与$\operatorname{Spur}(U\mathsf{H})$不变。后者是因为$\varphi_n$为$\mathsf{H}$的本征函数，$m\ne n$时$(\mathsf{H}\varphi_m,\allowbreak \varphi_n)$为零：
\begin{align*}
\operatorname{Spur}(U'\mathsf{H})
&=\sum_{n=1}^{\infty}(U\varphi_n,\varphi_n)\operatorname{Spur}(P_{[\varphi_n]}\mathsf{H})\\
&=\sum_{n=1}^{\infty}(U\varphi_n,\varphi_n)(\mathsf{H}\varphi_n,\varphi_n)\\
&=\sum_{m,n=1}^{\infty}(U\varphi_m,\varphi_n)(\mathsf{H}\varphi_n,\varphi_m)
=\operatorname{Spur}(U\mathsf{H}).
\end{align*}
从$R,\allowbreak \mathsf{H}$对易，即$\mathfrak R$与能量同时可测，也必须得出这一点。因此，即使只限于$U'$，也就是以$\varphi_1,\allowbreak \varphi_2,\allowbreak \ldots$为本征函数的统计算符，所求最大值也不变；而且最大值只会在其中取得。

于是$U=\sum_{n=1}^{\infty}w_nP_{[\varphi_n]}$。$U,\allowbreak U\mathsf{H},\allowbreak U\ln U$均以$\varphi_n$为本征函数，相应本征值依次为$w_n,\allowbreak W_nw_n,\allowbreak w_n\ln w_n$，故问题是，在约束$\sum_{n=1}^{\infty}w_n=1$、$\sum_{n=1}^{\infty}W_nw_n=\mathbf E$下，使$-N\varkappa\sum_{n=1}^{\infty}w_n\ln w_n$最大。这恰好就是通常气体理论中相应平衡问题所遇到的极值问题\origNote[题][，]{201}，解法也相同。按照熟知的极值计算规则，使极值达到的一组$w_1,\allowbreak w_2,\allowbreak \ldots$必须满足
\[
\frac{\partial}{\partial w_n}\left(\sum_{m=1}^{\infty}w_m\ln w_m\right)
+\alpha\frac{\partial}{\partial w_n}\left(\sum_{m=1}^{\infty}w_m\right)
+\beta\frac{\partial}{\partial w_m}\left(\sum_{m=1}^{\infty}W_mw_m\right)=0,
\]
其中$\alpha,\allowbreak \beta$为适当常数，$n=1,\allowbreak 2,\allowbreak \ldots$。即
\[
(\ln w_n+1)+\alpha+\beta W_n=0,\qquad
w_n=e^{-1-\alpha-\beta W_n}=ae^{-\beta W_n},
\]
这里以常数$a=e^{-1-\alpha}$代替了$\alpha$。

\hypertarget{origpage-209}{}\phantomsection\label{origpage-209}由$\sum_{n=1}^{\infty}w_n=1$，得$a=1/\sum_{n=1}^{\infty}e^{-\beta W_n}$，所以
\[
w_n=\frac{e^{-\beta W_n}}{\sum_{m=1}^{\infty}e^{-\beta W_m}}.
\]
再由$\sum_{n=1}^{\infty}W_nw_n=\mathbf E$，须有
\[
\frac{\sum_{n=1}^{\infty}W_ne^{-\beta W_n}}{\sum_{n=1}^{\infty}e^{-\beta W_n}}=\mathbf E,
\]
此式确定$\beta$。按惯例引入“配分函数”
\[
Z(\beta)=\sum_{n=1}^{\infty}e^{-\beta W_n}=\operatorname{Spur}(e^{-\beta \mathsf{H}});
\]
关于这一点及以下内容，参见原注 \hyperlink{note-197}{\getrefnumber{orig-note-197}}。\tnote[。][于][下]{下式$Z'(\beta)$等式最后的迹项前应有负号。这里保留原印公式。}于是
\[
Z'(\beta)=-\sum_{n=1}^{\infty}W_ne^{-\beta W_n}=\operatorname{Spur}(\mathsf{H}e^{-\beta \mathsf{H}}),
\]
因此，$\beta$的条件为
\[
-\frac{Z'(\beta)}{Z(\beta)}=\mathbf E.
\]
还假设对所有$\beta>0$，$\sum_{n=1}^{\infty}e^{-\beta W_n}$及$\sum_{n=1}^{\infty}W_ne^{-\beta W_n}$都收敛，即$n\to\infty$时$W_n\to\infty$足够快，例如$W_n/\ln n\to\infty$就够了。对$U$本身，得到
\[
U=\sum_{n=1}^{\infty}ae^{-\beta W_n}P_{[\varphi_n]}
=ae^{-\beta \mathsf{H}}=\frac{e^{-\beta \mathsf{H}}}{\operatorname{Spur}(e^{-\beta \mathsf{H}})}
=\frac{e^{-\beta \mathsf{H}}}{Z(\beta)}.
\]
给定$\mathbf E$或$\beta$，平衡系综$U$就确定了；正如所应有的，它的性质只依赖一个参数，现在可按气体理论的通常方法求出。

我们系综的熵为
\begin{align*}
\mathbf S&=-N\varkappa\operatorname{Spur}(U\ln U)
=-N\varkappa\operatorname{Spur}\left(\frac{e^{-\beta \mathsf{H}}}{Z(\beta)}\ln\frac{e^{-\beta \mathsf{H}}}{Z(\beta)}\right)\\
&=-\frac{N\varkappa}{Z(\beta)}\operatorname{Spur}\bigl(e^{-\beta \mathsf{H}}(-\beta \mathsf{H}-\ln Z(\beta))\bigr)\\
&=\frac{\beta N\varkappa}{Z(\beta)}\operatorname{Spur}(\mathsf{H}e^{-\beta \mathsf{H}})
+\frac{\ln Z(\beta)N\varkappa}{Z(\beta)}\operatorname{Spur}(e^{-\beta \mathsf{H}})\\
&=N\varkappa\left[-\beta\frac{Z'(\beta)}{Z(\beta)}+\ln Z(\beta)\right],
\end{align*}\hypertarget{origpage-210}{}\phantomsection\label{origpage-210}总能量为
\[
N\mathbf E=-N\frac{Z'(\beta)}{Z(\beta)}.
\]
应当与$\mathbf S$并列的是这个量，而不是$\mathbf E$本身。至此已把$U,\allowbreak \mathbf S,\allowbreak N\mathbf E$用$\beta$表示出来。与其只把$\beta$表示成$\mathbf E$的函数，更实用的是求平衡混合态的温度$\mathbf T$，再把一切归结为它。方法如下：使平衡混合态接触温度为$\mathbf T'$的热库，并从中取得能量$N\,d\mathbf E$。按照两条基本定律，总能量不变，熵不能减少。因此，热库损失能量$N\,d\mathbf E$，其熵变化量为$-N\,d\mathbf E/\mathbf T'$，同时须有
\[
d\mathbf S-\frac{N\,d\mathbf E}{\mathbf T'}
=\left(\frac{d\mathbf S}{N\,d\mathbf E}-\frac1{\mathbf T'}\right)N\,d\mathbf E\geq0.
\]
另一方面，$\mathbf T'\gtreqless\mathbf T$时显然有$N\,d\mathbf E\gtreqless0$，因为较冷物体从较热物体吸收能量。所以$\mathbf T'\gtreqless\mathbf T$意味着$d\mathbf S/(N\,d\mathbf E)-1/\mathbf T'\gtreqless0$，即
\[
\mathbf T'\gtreqless\frac{N\,d\mathbf E}{d\mathbf S}
=\frac{N\,d\mathbf E/d\beta}{d\mathbf S/d\beta}.
\]
因此
\begin{align*}
\mathbf T&=\frac{N\,d\mathbf E/d\beta}{d\mathbf S/d\beta}
=-\frac1\varkappa\frac{(Z'(\beta)/Z(\beta))'}{(\ln Z(\beta)-\beta Z'(\beta)/Z(\beta))'}\\
&=-\frac1\varkappa\frac{(Z'(\beta)/Z(\beta))'}{-\beta(Z'(\beta)/Z(\beta))'}
=\frac1{\varkappa\beta},
\end{align*}
也就是
\[
\beta=\frac1{\varkappa\mathbf T}.
\]
这样，$U,\allowbreak \mathbf S,\allowbreak N\mathbf E$全都表示成温度的函数。

上面得到的熵、平衡系综等表达式，与以经典力学为基础的热力学理论所给出的相应结果，具有显著的相似性。先看熵$-N\varkappa\operatorname{Spur}(U\ln U)$。$U=\sum_{n=1}^{\infty}w_nP_{[\varphi_n]}$是系综$P_{[\varphi_1]},\allowbreak P_{[\varphi_2]},\allowbreak \ldots$按$w_1:w_2:\cdots$的混合，即有$Nw_1$个$\varphi_1$系统、$Nw_2$个$\varphi_2$系统等等。对“微观排列数”
\[
\frac{N!}{(Nw_1)!(Nw_2)!\cdots}
\]
取对数，再乘以$\varkappa$，就得到这个系综的Boltzmann熵\origNote[熵][。]{201}。$N$很大，故可用Stirling公式$x!\approx\sqrt{2\pi x}\,e^{-x}x^x$近似阶乘。这样，
\[
\varkappa\ln[N!/((Nw_1)!(Nw_2)!\cdots)]
\]
在主要项上就变为
\[
-N\varkappa\sum_{n=1}^{\infty}w_n\ln w_n,
\]
而这正是$-\operatorname{Spur}(U\ln U)$。\tnote[。][再][原]{原书最后这一表达式省略了前述熵公式中的因子$N\varkappa$。}\hypertarget{origpage-211}{}\phantomsection\label{origpage-211}再看所得平衡系综$U=e^{-\mathsf{H}/(\varkappa\mathbf T)}$，这里略去归一化因子$1/Z(\beta)$。它等于$\sum_{n=1}^{\infty}e^{-W_n/(\varkappa\mathbf T)}P_{[\varphi_n]}$，即状态$P_{[\varphi_1]},\allowbreak P_{[\varphi_2]},\allowbreak \ldots$（能量为$W_1,\allowbreak W_2,\allowbreak \ldots$的定态）按相对权重$e^{-W_1/(\varkappa\mathbf T)}:e^{-W_2/(\varkappa\mathbf T)}:\cdots$的混合。若某个能量本征值有重数，例如$W_{n_1}=\cdots=W_{n_\nu}=W$，那么$P_{[\varphi_{n_1}]}+\cdots+P_{[\varphi_{n_\nu}]}$在混合态中的权重为$e^{-W/(\varkappa\mathbf T)}$，也就是说，正确归一化的系综$\frac1\nu(P_{[\varphi_{n_1}]}+\cdots+P_{[\varphi_{n_\nu}]})$（参见 \hyperref[sec:IV.3]{IV.3} 开头）的权重为$\nu e^{-W/(\varkappa\mathbf T)}$。除去出现了量子力学特有的构造$\frac1\nu(P_{[\varphi_{n_1}]}+\cdots+P_{[\varphi_{n_\nu}]})$外，这正是经典“正则”系综的定义。这就是所谓Boltzmann定理\origNote[理][。]{201}。

当$\mathbf T\to0$时，权重$e^{-W_n/(\varkappa\mathbf T)}$趋于$1$，所以$U$趋于$\sum_{n=1}^{\infty}P_{[\varphi_n]}=1$。\tnote[。][因][原]{原书印作$\mathbf T\to0$；固定$W_n$时，所列权重在$\mathbf T\to\infty$时才趋于$1$。这里保留原印极限。}因此，无能量限制时，$U=1$为绝对平衡态；这一结果我们在 \hyperref[sec:IV.3]{IV.3} 已得到。可见，“\textkaishu{量子轨道的先验等概率性}”（指单重、非简并的轨道；一般说来，本征值的重数是“先验权重”，参见上文）从这一理论中自行得出。

还值得查明：不借助热力学，单靠以下事实，对给定能量的平衡系综$U$究竟能说出多少？这些事实是：$U$为定常的，即随时间不变（过程 \hyperref[eq:V.1.1]{1}）；以及它在所有不影响能量的测量中也不变，即测量与能量同时可测的物理量时不变（过程 \hyperref[eq:V.1.2]{2}，其中$R,\allowbreak \mathsf{H}$对易，即$\varphi_1,\allowbreak \varphi_2,\allowbreak \ldots$为$\mathsf{H}$的本征函数）。\tnote[。][][本]{本段原书把两个过程的编号互换了：时间演化是过程 \hyperref[eq:V.1.2]{2}，测量是过程 \hyperref[eq:V.1.1]{1}。}

由微分方程$\partial U/\partial t=(2\pi i/h)(U\mathsf{H}-\mathsf{H}U)$，第一个条件仅表示$\mathsf{H},\allowbreak U$对易。第二个条件则是说：凡能作$\mathsf{H}$的完全本征函数系的$\varphi_1,\allowbreak \varphi_2,\allowbreak \ldots$，都有$U=U'$，也就是这些$\varphi_1,\allowbreak \varphi_2,\allowbreak \ldots$也为$U$的本征函数。设其相应的$\mathsf{H}$本征值为$W_1,\allowbreak W_2,\allowbreak \ldots$，$U$本征值为$w_1,\allowbreak w_2,\allowbreak \ldots$。若$W_j=W_k$，可把$\mathsf{H}$的$\varphi_j,\allowbreak \varphi_k$换成$(\varphi_j+\varphi_k)/\sqrt2$、$(\varphi_j-\varphi_k)/\sqrt2$，因而它们也必须是$U$的本征函数，遂得$w_j=w_k$。所以能构造满足$F(W_n)=w_n$\hypertarget{origpage-212}{}\phantomsection\label{origpage-212}（$n=1,\allowbreak 2,\allowbreak \ldots$）的函数$F(x)$，并有$F(\mathsf{H})=U$。这一条件充分是显然的，它蕴含$\mathsf{H},\allowbreak U$对易也一样清楚。

因此，这种方法固然得到$U=F(\mathsf{H})$，却无法确定$F(x)$；我们已知它应为$F(x)=e^{-\beta x}/Z(\beta)$，$\beta=1/(\varkappa\mathbf T)$。由$\operatorname{Spur}U=1$、$\operatorname{Spur}(U\mathsf{H})=\mathbf E$，还得到
\[
\sum_{n=1}^{\infty}F(W_n)=1,\qquad
\sum_{n=1}^{\infty}W_nF(W_n)=\mathbf E,
\]
但这一方法至此已用尽。

\section{宏观测量}\label{sec:V.4}
我们已经看到，所得熵表达式与经典熵完全类似；然而，它在系统随时间演变（过程 \hyperref[eq:V.1.2]{2}）时不变，只在测量（过程 \hyperref[eq:V.1.1]{1}）时增加，这仍令人诧异。经典理论中，测量全不起作用，熵通常仅在系统普通的力学时间演化中就会增加。因此，必须澄清这一表面上矛盾的行为。

通常的经典热力学论证如下：取体积为$\mathcal V$的容器，用隔板把左右两半分开，在右半部（体积$\mathcal V/2$）放入温度为$\mathbf T$、含$M$个分子的气体，为简单起见设为理想气体。如果让它等温、可逆地膨胀到体积$\mathcal V$，即让气体压力推回隔板，利用所得机械功，并以温度为$\mathbf T$的大热库保持气体温度恒定，那么外部（热库）的熵会减少$M\varkappa\ln2$（参见原注 \hyperlink{note-195}{\getrefnumber{orig-note-195}}），所以气体的熵就增加同样多。反之，若直接抽掉隔板，气体便扩散到空的左半部，体积增至$\mathcal V$；熵因而增加$M\varkappa\ln2$，却没有任何补偿，所以过程不可逆。熵在系统简单的力学时间演化——扩散——中增加了。为什么我们的理论得不到类似结果？

取$M=1$时，情况最清楚。对这种单分子气体，热力学仍然适用；体积加倍，熵确实增加$\varkappa\ln2$。但这个$\varkappa\ln2$的差别，只有当我们对分子的了解确实仅限于它位于体积为\hypertarget{origpage-213}{}\phantomsection\label{origpage-213}$\mathcal V/2$或$\mathcal V$的区域内时，才真正存在。例如，分子在体积$\mathcal V$中，但已知它在容器中点的左边还是右边，那么只须在中央插入隔板，让分子等温、可逆地把隔板推到左端或右端，就能作机械功$\varkappa\mathbf T\ln2$；这份能量取自热库。最后，分子仍在体积$\mathcal V$中，只是我们不再知道它位于中点左侧还是右侧；作为补偿，热库的熵减少了$\varkappa\ln2$。我们用自己的知识换得了$\varkappa\ln2$的熵减。\origNote[。][因]{202}因此，只要知道分子在哪一半，体积$\mathcal V$中的熵就与体积$\mathcal V/2$中的相同。若扩散前已准确知道分子的情况（位置与动量），就能计算扩散后每一时刻它位于右半部还是左半部，熵根本没有增加。只有在仅掌握宏观信息，即先前体积为$\mathcal V/2$时，扩散才真正使熵增加。

因此，对知道一切坐标和动量的经典观察者，熵恒定，而且为零，因为Boltzmann“微观排列数”为$1$，参见原注 \hyperlink{note-201}{\getrefnumber{orig-note-201}} 所引处。这恰如我们理论中的状态$U=P_{[\varphi]}$；它同样对应于观察者对系统所能拥有的最多知识。

所以，熵随时间变化，是因为观察者并非无所不知，或者不能把原则上可测的一切都确定、测量出来。他的感官只能感知所谓宏观量。不过，既然这样解释了最初那个表面矛盾，就有义务在量子力学系综中找到经典宏观熵的确切对应物；也就是站在某种观察者的角度来看熵：他不能测量一切量，只能测量某些特殊的量，即宏观量，而且可能连这些量也只能测到有限精度。

我们在 \hyperref[sec:III.3]{III.3} 得知，任何有限精度的测量，都可替换为对其他物理量作绝对精确的测量；后者是原物理量的函数，并有纯点谱。设$\mathfrak R$是这种量，算符为$R$，$\lambda^{(1)},\allowbreak \lambda^{(2)},\allowbreak \ldots$是其本征值中的不同值，那么测量$\mathfrak R$，等于回答：“是否\hypertarget{origpage-214}{}\phantomsection\label{origpage-214}$\mathfrak R=\lambda^{(1)}$？”“是否$\mathfrak R=\lambda^{(2)}$？”等等。也可以直接说：要以有限精度测量算符为$S$的$\mathfrak S$，例如判断它落在哪个区间$c_{n-1}<\lambda\leq c_n$，其中
\[
\cdots<c_{-2}<c_{-1}<c_0<c_1<c_2<\cdots,
\qquad c_n\to\pm\infty\mathwordspace{}\text{当}\mathwordspace{}n\to\pm\infty,
\]
就是要回答全部问题：“$\mathfrak S$是否落在$c_{n-1}<\lambda\leq c_n$中？”（$n=0,\allowbreak \pm1,\allowbreak \pm2,\allowbreak \ldots$）。

按 \hyperref[sec:III.5]{III.5}，这些问题对应于投影算符$E$，真正测量的是相应的、只取$0,\allowbreak 1$两个值的量$\mathfrak E$。在上述例子中，$\mathfrak E$分别是函数$F_n(\mathfrak R)$（$n=1,\allowbreak 2,\allowbreak \ldots$），其中
\[
F(\lambda)=\begin{cases}1,&\lambda=\lambda^{(n)},\\0,&\text{其他情形},\end{cases}
\]
或函数$G_n(\mathfrak S)$（$n=0,\allowbreak \pm1,\allowbreak \pm2,\allowbreak \ldots$），其中
\[
G(\lambda)=\begin{cases}1,&c_{n-1}<\lambda\leq c_n,\\0,&\text{其他情形}.\end{cases}
\]
相应的$E$则分别为$F_n(R)$或$G_n(S)$。因此，不必给出宏观可测量$\mathfrak S$及其可达测量精度，而可以给出宏观测量所能回答的问题$\mathfrak E$，或其投影算符$E$，参见 \hyperref[sec:III.5]{III.5}。给出观察者的这些$E$，正是对宏观观察者的刻画；在经典情形中，例如可以规定，他能测量气体每立方厘米中的温度和压力，或许还受某种精度限制，但不能测量别的。\origNote[。][]{203}

宏观测量的本意是：凡是这样可测的，就都同时可测；即所有宏观上可回答的问题$\mathfrak E$都能同时回答，也就是所有$E$彼此对易。量子力学物理量不能同时测量，最初之所以显得如此矛盾，恰是因为宏观观点中没有这一概念。这一点在原理上十分重要，宜再详细讨论。

回想怎样以有限精度同时测量两个公认不能同时测量的量，例如坐标$q$和动量$p$（参见 \hyperref[sec:III.4]{III.4}）。设测量的均方根误差分别为$\varepsilon,\allowbreak \eta$，由不确定性关系$\varepsilon\eta\sim h$。\hyperref[sec:III.4]{III.4} 的讨论说明，在这种精度要求下，同时测量确实可能：$q$（位置）测量用不太短波长的光，$p$（动量）测量用不太长的光波列。一切安排就绪后，真正的测量在于以某种方法，例如照相，探测两个光量子：\hypertarget{origpage-215}{}\phantomsection\label{origpage-215}一个是在$q$测量中经Compton效应偏转的光量子；另一个是在$p$测量中经Doppler效应反射、频率改变，又被棱镜或衍射光栅等光学装置偏转，以便测定该频率的光量子。因此，实验结束时留下两个光量子，甚至两张照相底片；我们须根据光量子的方向，或底片上感光斑的位置，计算$q,\allowbreak p$。但应强调：我们完全可以任意精确地确定这两个方向或感光斑位置；它们显然是同时可测的量，是两个不同对象的动量或坐标。然而，这对测量$q,\allowbreak p$帮助不大。正如 \hyperref[sec:III.4]{III.4} 所示，这些量与$q,\allowbreak p$的耦合使后者恰好仍留下$\varepsilon,\allowbreak \eta$的不确定度，即使前者已完全精确地知道也是如此；不可能安排到$\varepsilon\eta\ll h$。若把这两个方向或感光斑位置本身引入为物理量，算符记为$Q',\allowbreak P'$，便知$Q',\allowbreak P'$确实精确对易，但$q,\allowbreak p$的算符$Q,\allowbreak P$不能借助它们表示到分别优于$\varepsilon,\allowbreak \eta$的精度。设$Q',\allowbreak P'$对应的量为$q',\allowbreak p'$，那么真正宏观可测的是$q',\allowbreak p'$而非$q,\allowbreak p$本身，这种解释就十分自然（实际测得的确实是$q',\allowbreak p'$），也与一切宏观量同时可测的公设相容。

宜把这一事实作一般解释，视为宏观观点的特征：把各种一般不对易的算符$A,\allowbreak B,\allowbreak C,\allowbreak \ldots$，换成另一组彼此对易的算符$A',\allowbreak B',\allowbreak C',\allowbreak \ldots$，而前者近似为后者的函数。由于这些函数也可直接记为$A',\allowbreak B',\allowbreak C',\allowbreak \ldots$，便可说：$A',\allowbreak B',\allowbreak C',\allowbreak \ldots$分别近似$A,\allowbreak B,\allowbreak C,\allowbreak \ldots$，却精确地彼此对易。若以$\varepsilon_A,\allowbreak \varepsilon_B,\allowbreak \varepsilon_C,\allowbreak \ldots$衡量$A'-A,\allowbreak B'-B,\allowbreak C'-C,\allowbreak \ldots$的大小，就知$\varepsilon_A\varepsilon_B$应与$AB-BA$（一般非零）同数量级，其他亦然。这给出了所能达到的近似限度。当然，列举$A,\allowbreak B,\allowbreak C,\allowbreak \ldots$时，合理的做法是只考虑这样一些算符：其对应的物理量至少能够以尚可接受的精度作宏观测量。

这些定性说明，如果不能证明其要求在数学上\hypertarget{origpage-216}{}\phantomsection\label{origpage-216}可行，就仍是空洞的计划。因此，对典型情形$Q,\allowbreak P$，再从纯数学角度讨论上述$Q',\allowbreak P'$是否存在。给定两个正数$\varepsilon,\allowbreak \eta$，$\varepsilon\eta=h/(4\pi)$，要求两个对易的$Q',\allowbreak P'$，使$Q'-Q,\allowbreak P'-P$分别具有$\varepsilon,\allowbreak \eta$的数量级；这里的含义还须进一步确定。

我们采用完全精确可测的量$q',\allowbreak p'$，即要求$Q',\allowbreak P'$具有纯点谱。由于二者对易，有一个由公共本征函数组成的完全规范正交系$\varphi_1,\allowbreak \varphi_2,\allowbreak \ldots$，参见 \hyperref[sec:II.10]{II.10}。设$Q',\allowbreak P'$的相应本征值分别为$a_1,\allowbreak a_2,\allowbreak \ldots$和$b_1,\allowbreak b_2,\allowbreak \ldots$，即
\[
Q'=\sum_{n=1}^{\infty}a_nP_{[\varphi_n]},\qquad
P'=\sum_{n=1}^{\infty}b_nP_{[\varphi_n]}.
\]
测量应安排成测后得到$\varphi_1,\allowbreak \varphi_2,\allowbreak \ldots$之一；例如，可测量物理量$\mathfrak R$，其算符$R$有本征函数$\varphi_1,\allowbreak \varphi_2,\allowbreak \ldots$和互不相同的本征值$c_1,\allowbreak c_2,\allowbreak \ldots$，于是$Q',\allowbreak P'$为$R$的函数。说这样的测量近似包含了对$Q,\allowbreak P$的测量，显然是指：在状态$\varphi_n$中，$Q,\allowbreak P$可近似地用$Q',\allowbreak P'$的相应数值$a_n,\allowbreak b_n$来表示；也就是说，测得值集中在这些数值附近，散布很小。衡量这种散布的量是$(q-a_n)^2,\allowbreak (p-b_n)^2$的期望值，即
\begin{align*}
((Q-a_n1)^2\varphi_n,\varphi_n)&=\|(Q-a_n1)\varphi_n\|^2=\|Q\varphi_n-a_n\varphi_n\|^2,\\
((P-b_n1)^2\varphi_n,\varphi_n)&=\|(P-b_n1)\varphi_n\|^2=\|P\varphi_n-b_n\varphi_n\|^2.
\end{align*}
它们衡量$Q'$与$Q$、$P'$与$P$的差的平方，所以应分别约为$\varepsilon^2,\allowbreak \eta^2$。因此要求
\[
\|Q\varphi_n-a_n\varphi_n\|\lesssim\varepsilon,\qquad
\|P\varphi_n-b_n\varphi_n\|\lesssim\eta.
\]
所以，与其谈$Q',\allowbreak P'$，不如直接寻找一个完全规范正交系$\varphi_1,\allowbreak \varphi_2,\allowbreak \ldots$，使适当选取$a_1,\allowbreak a_2,\allowbreak \ldots$与$b_1,\allowbreak b_2,\allowbreak \ldots$后，上述估计成立。

在 \hyperref[sec:III.4]{III.4} 中，我们已找到满足$\|\varphi\|=1$的单个$\varphi$，对适当的$a,\allowbreak b$还有
\[
\|Q\varphi-a\varphi\|=\varepsilon,\qquad\|P\varphi-b\varphi\|=\eta:
\]
具体为
\[
\varphi_{\varrho,\sigma,\gamma}
=\varphi_{\varrho,\sigma,\gamma}(q)
\equiv\left(\frac{2\gamma}{h}\right)^{1/4}
 e^{-\frac{\pi\gamma}{h}(q-\sigma)^2+\frac{2\pi\varrho}{h}iq}.
\]
由于$\varepsilon\eta=h/(4\pi)$，这里又置$\varepsilon=\sqrt{h\gamma/(4\pi)}$、$\eta=\sqrt{h/(4\pi\gamma)}$，即$\gamma=\varepsilon/\eta$，并须选$a=\sigma$、$b=\varrho$。\tnote[。][现][对]{对上列高斯函数，$Q,\allowbreak P$的期望值分别为$\sigma,\allowbreak \varrho$，标准差分别为$\sqrt{h/(4\pi\gamma)}$、$\sqrt{h\gamma/(4\pi)}$。原书这里颠倒了两个宽度，下一段又颠倒了两个期望值。此处保留原书记号。}现在要从这些$\varphi_{\varrho,\allowbreak \sigma,\allowbreak \gamma}$组成一个完全规范正交系。\hypertarget{origpage-217}{}\phantomsection\label{origpage-217}由于$\varrho$是$Q$的期望值、$\sigma$是$P$的期望值，让$\varrho,\allowbreak \sigma$彼此独立地分别遍历一组数，前者间距约为$\varepsilon$，后者约为$\eta$，看来是合理的。实际上，采用单位$2\sqrt\pi\,\varepsilon=\sqrt{h\gamma}$与$2\sqrt\pi\,\eta=\sqrt{h/\gamma}$很方便，即
\[
\varrho=\sqrt{h\gamma}\,\mu,\qquad\sigma=\sqrt{h/\gamma}\,\nu
\qquad(\mu,\nu=0,\pm1,\pm2,\ldots).
\]
因此，以$\psi_{\mu,\allowbreak \nu}=\varphi_{\sqrt{h\gamma}\,\mu,\allowbreak \sqrt{h/\gamma}\,\nu,\allowbreak \gamma}$（$\mu,\allowbreak \nu=0,\allowbreak \pm1,\allowbreak \pm2,\allowbreak \ldots$）对应$\varphi_n$（$n=1,\allowbreak 2,\allowbreak \ldots$）。用两个下标$\mu,\allowbreak \nu$代替一个$n$显然无关紧要。

不过，这些$\psi_{\mu,\allowbreak \nu}$还不正交。它们已归一化，且
\[
\|Q\psi_{\mu,\nu}-\sqrt{h\gamma}\,\mu\|=\varepsilon,\qquad
\|P\psi_{\mu,\nu}-\sqrt{h/\gamma}\,\nu\|=\eta.
\]
如果按E.~Schmidt的方法，以任意次序将它们“正交化”（参见 \hyperref[sec:II.2]{II.2}，\CJKglue \hyperref[thm:II.8]{定理8} 的证明），便可不太困难地证明所得规范正交系$\psi'_{\mu,\allowbreak \nu}$的完全性，并保证估计
\[
\|Q\psi_{\mu,\nu}-\sqrt{h\gamma}\,\mu\|\leq C\varepsilon,\qquad
\|P\psi_{\mu,\nu}-\sqrt{h/\gamma}\,\nu\|\leq C\eta
\]
成立，其中$C$约为$60$；这个数值大概还能大幅降低。证明需要相当冗长、但不需要任何新观点的计算，这里略去。$C\sim60$这样的因子影响不大，因为$\varepsilon\eta=h/(4\pi)$用宏观的CGS单位衡量极其微小，约为$10^{-28}$！

总之，有理由假设一切宏观算符彼此对易，特别是前面引入的宏观投影算符$E$也如此。

$E$对应于一切宏观可回答的问题$\mathfrak E$，即关于所考察系统、在宏观上能够作出的一切分类判断。如前所见，它们彼此对易；由 \hyperref[sec:III.5]{III.5}，若有$E$，也就有$1-E$，若有$E,\allowbreak F$，也有$EF$、$E+F-EF$、$E-EF$。很自然可以假设它们只有有限个$E_1,\allowbreak \ldots,\allowbreak E_n$。暂时引入记号$E^{(+)}=E$、$E^{(-)}=1-E$，考察全部$2^n$个乘积$E_1^{(s_1)}\cdots E_n^{(s_n)}$（$s_1,\allowbreak \ldots,\allowbreak s_n=\pm$）。任意两个不同乘积相乘都为零：因为若$E_1^{(s_1)}\cdots E_n^{(s_n)}$与$E_1^{(t_1)}\cdots E_n^{(t_n)}$不同，例如$s_\nu\ne t_\nu$，则它们之积包含因子$E_\nu^{(s_\nu)},\allowbreak E_\nu^{(t_\nu)}$，即$E_\nu^{(+)}=E$\hypertarget{origpage-218}{}\phantomsection\label{origpage-218}与$E_\nu^{(-)}=1-E$，其乘积为零。每个$E_\nu$都是若干这样的乘积之和，具体为
\[
E_\nu=\sum_{\substack{s_1,\ldots,s_{\nu-1},s_{\nu+1},\ldots,s_n=\pm}}
 E_1^{(s_1)}\cdots E_{\nu-1}^{(s_{\nu-1})}E_\nu^{(+)}E_{\nu+1}^{(s_{\nu+1})}\cdots E_n^{(s_n)}.
\]
把这些乘积中非零者记为$E'_1,\allowbreak \ldots,\allowbreak E'_m$。显然$m\leq2^n$，甚至$m\leq n-1$，因为它们必须在$E_1,\allowbreak \ldots,\allowbreak E_n$中出现，而且非零。于是$E'_\mu\ne0$，$\mu\ne\nu$时$E'_\mu E'_\nu=0$，每个$E_\mu$都是若干$E'_\kappa$之和；后一事实还给出$n=2^m$。注意，不可能有$E_\mu+E_\nu=E'_\varrho$，除非$E_\mu=0,\allowbreak E_\nu=E'_\varrho$，或$E_\mu=E'_\varrho,\allowbreak E_\nu=0$。否则，$E_\mu,\allowbreak E_\nu$各是若干$E'_\kappa$之和，$E'_\varrho$便成为至少两个$E'_\kappa$的和（或有重复）。由 \hyperref[sec:II.4]{II.4} 的\CJKglue \hyperref[thm:II.15]{定理15}、\hyperref[thm:II.16]{16}，这些项彼此不同；且既然至少有两项，也都不同于$E'_\varrho$，因此与$E'_\varrho$的乘积都为零，其和与$E'_\varrho$相乘也为零；但这个和原应等于$E'_\varrho$，矛盾。

因此，$E'_1,\allowbreak \ldots,\allowbreak E'_m$所对应的性质$\mathfrak E'_1,\allowbreak \ldots,\allowbreak \mathfrak E'_m$，是具有以下特征的宏观性质：每一个都有可能成立，任意两个彼此排斥，每个宏观性质都是其中若干个的析取，而它们中的任何一个都不能再分解成两个更细的宏观性质的析取。因此，$\mathfrak E'_1,\allowbreak \ldots,\allowbreak \mathfrak E'_m$就是宏观上所能作出的最细分类，不能再作宏观分解。

以下不再要求它们的数目有限，只要求存在宏观不可分的性质$\mathfrak E'_1,\allowbreak \mathfrak E'_2,\allowbreak \ldots$。相应投影算符为$E'_1,\allowbreak E'_2,\allowbreak \ldots$，仍然全非零、两两正交，而每个宏观$E$都是其中若干个的和。

所以$1$也是其中若干个的和。如果某个$E'_\nu$不在其中，它就与所有入选的算符正交，因而与$1$正交，即$E'_\nu=E'_\nu\cdot1=0$，不可能。所以$E'_1+E'_2+\cdots=1$。以下去掉撇号，写为$\mathfrak E_1,\allowbreak \mathfrak E_2,\allowbreak \ldots$及$E_1,\allowbreak E_2,\allowbreak \ldots$。相应的闭线性流形记为$\mathfrak M_1,\allowbreak \mathfrak M_2,\allowbreak \ldots$，维数为$s_1,\allowbreak s_2,\allowbreak \ldots$。

若$s_n=1$，即$\mathfrak M_n$为一维，就有$\mathfrak M_n=[\varphi_n]$、$E_n=P_{[\varphi_n]}$；又由$E_1+E_2+\cdots=1$，$\varphi_1,\allowbreak \varphi_2,\allowbreak \ldots$构成完全规范正交系。这将意味着，仅靠宏观测量就能完全确定所考察系统的状态。通常并非如此，故一般有$s_n>1$，甚至$s_n\gg1$。

此外，$E_n$作为宏观世界描述的基本构件，在某种\hypertarget{origpage-219}{}\phantomsection\label{origpage-219}意义上对应于经典理论惯用的相空间分格。我们已经看到，它们能够近似再现不对易算符的行为，特别是对相空间至关重要的$Q,\allowbreak P$的行为。

对于不可分投影算符为$E_1,\allowbreak E_2,\allowbreak \ldots$的宏观观察者，混合态$U$具有什么熵？更确切地说，这样的观察者把$U$变为$V$时，最多能够获得多少熵？也就是说，作为转变$U\to V$的补偿，他最多能使外部对象的熵减少多少（这个熵减自然也可能为零或为负）？

首先，若两个系综$U,\allowbreak U'$对全部$E_1,\allowbreak E_2,\allowbreak \ldots$都有相同的期望值，即
\[
\operatorname{Spur}(UE_n)=\operatorname{Spur}(U'E_n)\qquad(n=1,2,\ldots)\text{，}
\]
观察者就根本无法区分二者。过一段时间也许可以区分，因为$U,\allowbreak U'$按 \hyperref[eq:V.1.2]{2} 改变，而
\[
\operatorname{Spur}(AUA^{-1}E_n)=\operatorname{Spur}(AU'A^{-1}E_n),\qquad
A=e^{-\frac{2\pi i}{h}t\mathsf{H}},
\]
未必仍成立\origNote[立][；]{204}；但这里考察的只是能够立即进行的测量，所以在上述条件下，可以认为$U,\allowbreak U'$不可区分。此外，观察者也只能用这样的半透壁：让落在某些$\mathfrak M_n$中的状态$\varphi$透过，让其余状态原状反射。按照 \hyperref[sec:V.2]{V.2} 的方法，不难确认，这已经足以把$U'=\sum_{n=1}^{\infty}x_nE_n$可逆地变成$V'=\sum_{n=1}^{\infty}y_nE_n$。所以熵差仍为$\varkappa\operatorname{Spur}(U'\ln U')-\varkappa\operatorname{Spur}(V'\ln V')$，而$U'$的熵为$-\varkappa\operatorname{Spur}(U'\ln U')$。不过，若要存在这种满足$\operatorname{Spur}U'=1$的$U'$，各$\operatorname{Spur}E_n$，即$s_n$，必须有限。以下假设每个$s_n$都有限。$U'$的本征值$x_1$有$s_1$重，$x_2$有$s_2$重，等等；于是$-U'\ln U'$的本征值$-x_1\ln x_1$有$s_1$重，$-x_2\ln x_2$有$s_2$重，等等。因此，$\operatorname{Spur}U'=1$表示$\sum_{n=1}^{\infty}s_nx_n=1$，熵则为$-\varkappa\sum_{n=1}^{\infty}s_nx_n\ln x_n$。由于
\[
U'E_m=\sum_{n=1}^{\infty}x_nE_nE_m=x_mE_m,\qquad
\operatorname{Spur}(U'E_m)=x_m\operatorname{Spur}E_m=s_mx_m,
\]
有$x_m=\operatorname{Spur}(U'E_m)/s_m$，故所说的熵等于
\[
-\varkappa\sum_{n=1}^{\infty}\operatorname{Spur}(U'E_n)
\ln\frac{\operatorname{Spur}(U'E_n)}{s_n}.
\]
对任意$U$（$\operatorname{Spur}U=1$），熵也必须为
\[
-\varkappa\sum_{n=1}^{\infty}\operatorname{Spur}(UE_n)
\ln\frac{\operatorname{Spur}(UE_n)}{s_n}.
\]\hypertarget{origpage-220}{}\phantomsection\label{origpage-220}因为若置
\[
x_n=\frac{\operatorname{Spur}(UE_n)}{s_n},\qquad U'=\sum_{n=1}^{\infty}x_nE_n,
\]
就有$\operatorname{Spur}(UE_n)=\operatorname{Spur}(U'E_n)$，而$U,\allowbreak U'$不可区分，故有相同的熵。

还须指出，这个熵总不小于通常的熵：
\[
-\varkappa\sum_{n=1}^{\infty}\operatorname{Spur}(UE_n)
\ln\frac{\operatorname{Spur}(UE_n)}{s_n}
\geq-\varkappa\operatorname{Spur}(U\ln U),
\]
仅当$U=\sum_{n=1}^{\infty}x_nE_n$时成立等号。由 \hyperref[sec:V.3]{V.3} 的结果，只要能通过若干次（未必宏观的）过程 \hyperref[eq:V.1.1]{1} 把$U$变为
\[
U'=\sum_{n=1}^{\infty}[\operatorname{Spur}(UE_n)/s_n]E_n,
\]
上述不等式就必然成立；因为左端就是$-\varkappa\operatorname{Spur}(U'\ln U')$，而$U=\sum_{n=1}^{\infty}x_nE_n$与$U=U'$同义。取规范正交系$\varphi_1^{(n)},\allowbreak \ldots,\allowbreak \varphi_{s_n}^{(n)}$，张成$E_n$对应的闭线性流形$\mathfrak M_n$；由$\sum_{n=1}^{\infty}E_n=1$，全部$\varphi_\nu^{(n)}$（$n=1,\allowbreak 2,\allowbreak \ldots$，$\nu=1,\allowbreak \ldots,\allowbreak s_n$）构成完全规范正交系。设算符$R$以它们为本征函数，并有互不相同的本征值，相应物理量为$\mathfrak R$。测量$\mathfrak R$，按 \hyperref[eq:V.1.1]{1}，使$U$变为
\[
U''=\sum_{n=1}^{\infty}\sum_{\nu=1}^{s_n}
(U\varphi_\nu^{(n)},\varphi_\nu^{(n)})P_{[\varphi_\nu^{(n)}]}.
\]
然后令
\[
\psi_\mu^{(n)}=\frac1{\sqrt{s_n}}\sum_{\nu=1}^{s_n}
 e^{\frac{2\pi i}{s_n}\mu\nu}\varphi_\nu^{(n)}
\qquad(\mu=1,\ldots,s_n).
\]
$\psi_1^{(n)},\allowbreak \ldots,\allowbreak \psi_{s_n}^{(n)}$构成规范正交系，张成与$\varphi_1^{(n)},\allowbreak \ldots,\allowbreak \varphi_{s_n}^{(n)}$相同的闭线性流形$\mathfrak M_n$。所以全部$\psi_\nu^{(n)}$（$n=1,\allowbreak 2,\allowbreak \ldots$，$\nu=1,\allowbreak \ldots,\allowbreak s_n$）也构成完全规范正交系。以它们为本征函数构造算符$S$，相应物理量为$\mathfrak S$。还记录下列公式：
\[
(P_{[\varphi_\nu^{(n)}]}\psi_\mu^{(m)},\psi_\mu^{(m)})
=\begin{cases}0,&m\ne n,\\1/s_n,&m=n,\end{cases}
\]
\[
\sum_{\nu=1}^{s_n}P_{[\varphi_\nu^{(n)}]}
=\sum_{\nu=1}^{s_n}P_{[\psi_\nu^{(n)}]}=E_n.
\]\hypertarget{origpage-221}{}\phantomsection\label{origpage-221}因此，测量$\mathfrak S$，按 \hyperref[eq:V.1.1]{1}，就把$U''$变成
\begin{align*}
&\sum_{m=1}^{\infty}\sum_{\mu=1}^{s_m}(U''\psi_\mu^{(m)},\psi_\mu^{(m)})P_{[\psi_\mu^{(m)}]}\\
&\quad=\sum_{m=1}^{\infty}\sum_{\mu=1}^{s_m}
\left[\sum_{n=1}^{\infty}\sum_{\nu=1}^{s_n}(U\varphi_\nu^{(n)},\varphi_\nu^{(n)})
(P_{[\varphi_\nu^{(n)}]}\psi_\mu^{(m)},\psi_\mu^{(m)})\right]P_{[\psi_\mu^{(m)}]}\\
&\quad=\sum_{m=1}^{\infty}\sum_{\mu=1}^{s_m}
\left[\sum_{\nu=1}^{s_m}\frac{(U\varphi_\nu^{(m)},\varphi_\nu^{(m)})}{s_m}\right]P_{[\psi_\mu^{(m)}]}\\
&\quad=\sum_{m=1}^{\infty}\sum_{\mu=1}^{s_m}
\frac{\operatorname{Spur}(UE_m)}{s_m}P_{[\psi_\mu^{(m)}]}
=\sum_{m=1}^{\infty}\frac{\operatorname{Spur}(UE_m)}{s_m}E_m=U'.
\end{align*}
所以，两个过程 \hyperref[eq:V.1.1]{1} 已足以把$U$变成$U'$；证明所需的正是这一点。

对于状态$U=P_{[\varphi]}$，有$\operatorname{Spur}(UE_n)=(E_n\varphi,\allowbreak \varphi)=\|E_n\varphi\|^2$，所以这个熵为
\[
-\varkappa\sum_{n=1}^{\infty}\|E_n\varphi\|^2\ln\frac{\|E_n\varphi\|^2}{s_n}.
\]
它已没有“微观”熵的那些缺陷：一般说来，它不随时间（即过程 \hyperref[eq:V.1.2]{2}）保持恒定，也并非对所有状态$U=P_{[\varphi]}$都为零。事实上，构成这个熵的$\operatorname{Spur}(UE_n)$一般不随时间恒定，原注 \hyperlink{note-202}{\getrefnumber{orig-note-202}} 中已讨论过。很容易指出状态$U=P_{[\varphi]}$何时熵为零：由于$0\leq\|E_n\varphi\|^2/s_n\leq1$，熵表达式中的各个加项$\|E_n\varphi\|^2\ln(\|E_n\varphi\|^2/s_n)$都$\leq0$，故必须全为零。这就要求$\|E_n\varphi\|^2/s_n=0$或$1$。前者表示$E_n\varphi=0$；后者表示$\|E_n\varphi\|=\sqrt{s_n}$，但由
\[
\|E_n\varphi\|\leq1,\qquad s_n\geq1,
\]
得$s_n=1$、$\|E_n\varphi\|=\|\varphi\|$，即$E_n\varphi=\varphi$；也就是$s_n=1$且$\varphi$在$\mathfrak M_n$中。后一情况肯定不能对两个不同的$n$同时成立，也不能对任何$n$都不成立，否则就恒有$E_n\varphi=0$，再由$\sum_{n=1}^{\infty}E_n=1$得$\varphi=0$。所以，必须恰有一个$n$使$\varphi$属于$\mathfrak M_n$，而且$s_n=1$。我们已知一般所有$s_n\gg1$，故这不可能。因此，我们的熵始终大于零。

宏观熵既然随时间变化，接着的问题就是：它是否像现实世界中唯象热力学的熵那样，主要趋于增加？在经典力学理论中，所谓Boltzmann的H定理给出肯定回答，\hypertarget{origpage-222}{}\phantomsection\label{origpage-222}但必须作某些统计假设，即所谓“无序假设”。\origNote[。][在]{205}在量子力学中，作者已成功地不依靠这类假设证明了相应定理。\origNote[。][如]{206}如果深入讨论这个问题，以及与之密切相连的遍历定理（在原注 \hyperlink{note-206}{\getrefnumber{orig-note-206}} 所引论文中也已证明），就超出这里的范围了，因此不再复述这些研究。有兴趣的读者可参阅所引论文。


% 中文编写中间材料：chapter-VI.tex
\chapter{测量过程}\label{chap:VI}
\section{问题的表述}\label{sec:VI.1}
此前讨论量子力学与各种因果性、统计性的自然描述方法之间的关系时，我们发现量子力学的描述方式存在一种奇特的两重性，尚未得到充分解释。一方面，状态$\varphi$在能量算符$\mathsf H$的影响下，在时间区间$0\leq\tau\leq t$内，按如下方式变为状态$\varphi'$：
\begin{align*}
 \frac{\partial}{\partial\tau}\varphi_\tau
 &=-\frac{2\pi i}{h}\mathsf H\varphi_\tau\qquad(0\leq\tau\leq t),\\
 \varphi_0&=\varphi,\qquad\varphi_t=\varphi',
\end{align*}
即
\[
 \varphi'=e^{-\frac{2\pi i}{h}t\mathsf H}\varphi,
\]
所以完全遵循因果律。混合态$U$相应地变为
\[
 U'=e^{-\frac{2\pi i}{h}t\mathsf H}Ue^{\frac{2\pi i}{h}t\mathsf H};
\]
$\varphi$按照因果律变为$\varphi'$时，状态$U=P_{[\varphi]}$也变为状态$U'=P_{[\varphi']}$，这是 \hyperref[sec:V.1]{V.1} 的过程 \hyperref[eq:V.1.2]{\textbf{2}}。另一方面，测量时——例如测量一个本征值全为单重、本征函数为$\varphi_1,\allowbreak \varphi_2,\allowbreak \ldots$的量——状态$\varphi$经历非因果变化：可能产生$\varphi_1,\allowbreak \varphi_2,\allowbreak \ldots$中的任何一个状态，概率分别为$|(\varphi,\allowbreak \varphi_1)|^2,\allowbreak |(\varphi,\allowbreak \varphi_2)|^2,\allowbreak \ldots$。因此得到混合态
\[
 U'=\sum_{n=1}^{\infty}|(\varphi,\varphi_n)|^2P_{[\varphi_n]}
\]
一般说来，混合态$U$变为
\[
 U'=\sum_{n=1}^{\infty}(U\varphi_n,\varphi_n)P_{[\varphi_n]}
\]
这是 \hyperref[sec:V.1]{V.1} 的过程 \hyperref[eq:V.1.1]{\textbf{1}}。这里状态变成混合态，所以过程不遵循因果律。

这两个过程$U\longrightarrow U'$的差别是根本性的：除了对因果性原则表现不同，前者在热力学上可逆，后者则不可逆，参见 \hyperref[sec:V.3]{V.3}。

\hypertarget{origpage-223}{}\phantomsection\label{origpage-223}现在把这些情况，与自然界及其观察中实际存在的情况比较。首先，测量，或者与测量相伴的主观统觉过程，相对于物理环境确实是一种全新的现象，不能还原为物理环境中的过程。因为它把我们引向物理环境之外；更确切地说，是引向个体内在的思想生活。这种内在生活无法检验，因为任何检验的尝试都已经以它为前提，参见下文。不过，自然科学的世界观有一个基本要求，即所谓心物平行原则（\textit{psychophysikalischer Parallelismus}）：主观统觉其实是物理范围之外的过程，但必须能像它发生于物理世界中那样描述它。也就是说，必须能把这一主观过程的各个环节，分别与普通空间中、客观环境内的物理过程对应起来。当然，在作这种对应时，总会需要把这些过程定位于我们身体所占据的空间区域内；但这不改变它们属于环境这一点。

用一个简单例子说明这种看法。要测量某个温度，可以只计算到温度计水银球周围的温度，然后说：温度计测量了这一温度。但也可继续计算，根据分子动理论所解释的水银性质，算出它如何受热、膨胀，以及由此形成的水银柱长度，再说：观察者看到了这个长度。再往前，还可把他的光源包括进来，求光量子在不透明水银柱上的反射、其余光量子进入眼睛的路径、在晶状体中的折射，以及视网膜上的成像，然后才说：观察者的视网膜接收了这个像。假如生理学知识比今天更精确，还可继续追踪这个像在视网膜、神经通路、脑中引起的化学反应，最后才说：观察者统觉到自己脑细胞中的这些化学变化。但不论计算到哪里——水银球、温度计刻度、视网膜，还是脑内——总有一处必须说：然后观察者感知到它。也就是说，总须把世界分成两部分，一部分是被观察的系统，另一部分是观察者。\hypertarget{origpage-224}{}\phantomsection\label{origpage-224}前一部分中的一切物理过程，至少原则上都可以任意精确地追踪；对后一部分，这样做则没有意义。两者的界限有很大的任意性；上例中就有四种不同划法。特别是，这种意义上的观察者，绝不必等同于实际观察者的身体：上例中，有一次甚至把温度计也归到观察者一侧，另一次却连观察者的眼睛和神经通路也不包括在内。界限可以任意深入实际观察者身体内部，这就是心物平行原则的内容。但这不改变如下事实：任何一种描述，要避免成为空谈，也就是要能够与经验比较，就必须在某处划出这条界限。因为经验只作这样的陈述：“\textkaishu{某个观察者有了某种主观感知}”，从不作这样的陈述：“\textkaishu{某个物理量具有某个数值}”。

量子力学用过程 \hyperref[eq:V.1.2]{\textbf{2}}（\hyperref[sec:V.1]{V.1}）描述世界中被观察部分的事件，条件是它们不与观察者部分相互作用；一旦发生这种相互作用，即发生测量，就规定用过程 \hyperref[eq:V.1.1]{\textbf{1}}。这样，两重性便有了理由。\origNote[。][不]{207}不过，在证明被观察系统与观察者的界限能按上述意义任意移动之前，仍有违背心物平行原则的危险。

为讨论这一点，把世界分为$I,\allowbreak II,\allowbreak III$三部分：$I$为被观察的系统本身，$II$为测量仪器，$III$为观察者本身。\origNote[。][须]{208}须证明：界限既可划在$I$与$II+III$之间，也可划在$I+II$与$III$之间。在上例中，比较第一、第二种划法时，$I$为待研究系统，$II$为温度计，$III$为光加观察者；比较第二、第三种时，$I$为待研究系统加温度计，$II$为光加观察者的眼睛，$III$为从视网膜开始的观察者；比较第三、第四种时，$I$为直到观察者视网膜之前的一切，$II$为其视网膜、神经通路和脑，$III$为抽象的“我”。因此，第一种划法把过程 \hyperref[eq:V.1.2]{\textbf{2}} 用于$I$，把过程 \hyperref[eq:V.1.1]{\textbf{1}} 用于$I$与$II+III$间的相互作用；第二种则把 \hyperref[eq:V.1.2]{\textbf{2}} 用于$I+II$，把 \hyperref[eq:V.1.1]{\textbf{1}} 用于$I+II$与$III$间的相互作用。$III$本身在两种情况下都不纳入计算。我们真正的任务，是证明两种方法对$I$给出相同论断；两种情况下，只有它始终属于世界中被观察的部分。

\hypertarget{origpage-225}{}\phantomsection\label{origpage-225}但要顺利开展这一工作，须先详细研究两个物理系统的合并过程，即从$I$、$II$到$I+II$的过程。

\section{复合系统}\label{sec:VI.2}
如前节结尾所说，现在考察两个物理系统$I,\allowbreak II$及其合并所得$I+II$；这里$I,\allowbreak II$不必具有前节中的意义。在经典力学描述中，设$I$有$k$个自由度，坐标为$q_1,\allowbreak \ldots,\allowbreak q_k$，以后总用一个符号$q$简写；相应地，$II$有$l$个自由度，坐标为$r_1,\allowbreak \ldots,\allowbreak r_l$，简写为$r$。于是$I+II$有$k+l$个自由度，坐标为$q_1,\allowbreak \ldots,\allowbreak q_k,\allowbreak r_1,\allowbreak \ldots,\allowbreak r_l$，简写为$q,\allowbreak r$。因此，在量子力学中，$I$的波函数形如$\varphi(q)$，$II$的形如$\xi(r)$，$I+II$的形如$\Phi(q,\allowbreak r)$。在相应希尔伯特空间$\mathfrak R^I,\allowbreak \mathfrak R^{II},\allowbreak \mathfrak R^{I+II}$中，内积分别定义为
\[
 \int\varphi(q)\overline{\psi(q)}\,dq,\qquad
 \int\xi(r)\overline{\eta(r)}\,dr,\qquad
 \iint\Phi(q,r)\overline{\Psi(q,r)}\,dq\,dr,
\]
$I,\allowbreak II,\allowbreak I+II$的物理量，分别对应$\mathfrak R^I,\allowbreak \mathfrak R^{II},\allowbreak \mathfrak R^{I+II}$中的超极大厄米算符$\mathsf A,\allowbreak A,\allowbreak \boldsymbol A$。

$I$中的每个物理量当然也是$I+II$中的物理量；应从其$\mathsf A$按如下方式确定$\boldsymbol A$：求$\boldsymbol A\Phi(q,\allowbreak r)$时，将$r$视为常数，把$\mathsf A$作用于以$q$为变量的函数$\Phi(q,\allowbreak r)$。\origNote[。][这]{209}这一对应规则对坐标、动量算符肯定正确，即对$Q_1,\allowbreak \ldots,\allowbreak Q_k$与$P_1,\allowbreak \ldots,\allowbreak P_k$，亦即$q_1,\allowbreak \ldots,\allowbreak q_k$与$\frac{h}{2\pi i}\frac{\partial}{\partial q_1},\allowbreak \ldots,\allowbreak \frac{h}{2\pi i}\frac{\partial}{\partial q_k}$正确（参见 \hyperref[sec:I.2]{I.2}）；它也符合 \hyperref[sec:IV.2]{IV.2} 的对应原则 \hyperref[cond:IV.2.I]{\textbf{I}}、\hyperref[cond:IV.2.II]{\textbf{II}}。\origNote[。][因]{210}因此，把它作为普遍公设；量子力学中通常也正如此使用。

同样，$II$中的每个物理量也是$I+II$中的物理量，其$A$也按同样规则给出$\boldsymbol A$：$\boldsymbol A\Phi(q,\allowbreak r)$等于把$q$视为常数、把$\Phi(q,\allowbreak r)$视为$r$的函数时所得的$A\Phi(q,\allowbreak r)$。

若$\varphi_m(q)$（$m=1,\allowbreak 2,\allowbreak \ldots$）是$\mathfrak R^I$中的完全规范正交系，$\xi_n(r)$（$n=1,\allowbreak 2,\allowbreak \ldots$）是$\mathfrak R^{II}$中的完全规范正交系，那么$\Phi_{mn}(q,\allowbreak r)=\varphi_m(q)\xi_n(r)$（$m,\allowbreak n=1,\allowbreak 2,\allowbreak \ldots$）显然是$\mathfrak R^{I+II}$中的完全规范正交系。因此，算符$\mathsf A,\allowbreak A,\allowbreak \boldsymbol A$分别可用矩阵$\{a_{m/m'}\}$、$\{a_{n/n'}\}$、$\{\alpha_{mn/m'n'}\}$表示（$m,\allowbreak m',\allowbreak n,\allowbreak n'=1,\allowbreak 2,\allowbreak \ldots$）。\origNote[。][下]{211}下面将多次用到这一点。矩阵表示的含义为
\begin{align*}
 \mathsf A\varphi_m(q)&=\sum_{m'=1}^{\infty}a_{m/m'}\varphi_{m'}(q),\\
 A\xi_n(r)&=\sum_{n'=1}^{\infty}a_{n/n'}\xi_{n'}(r),
\end{align*}\hypertarget{origpage-226}{}\phantomsection\label{origpage-226}以及
\[
 \boldsymbol A\Phi_{mn}(q,r)
 =\sum_{m',n'=1}^{\infty}\alpha_{mn/m'n'}\Phi_{m'n'}(q,r),
\]
即
\[
 \boldsymbol A\varphi_m(q)\xi_n(r)
 =\sum_{m',n'=1}^{\infty}\alpha_{mn/m'n'}\varphi_{m'}(q)\xi_{n'}(r).
\]
特别地，对应$\mathsf A\longrightarrow\boldsymbol A$表示
\[
 \boldsymbol A\varphi_m(q)\xi_n(r)
 =(\mathsf A\varphi_m(q))\xi_n(r)
 =\sum_{m'=1}^{\infty}a_{m/m'}\varphi_{m'}(q)\xi_n(r),
\]
即
\[
 \alpha_{mn/m'n'}=a_{m/m'}\delta_{n/n'},\qquad
 \delta_{n/n'}=\begin{cases}1,&n=n',\\0,&n\ne n'.\end{cases}
\]
类似地，对应$A\longrightarrow\boldsymbol A$表示$\alpha_{mn/m'n'}=a_{n/n'}\delta_{m/m'}$。

$I+II$中的统计系综由其统计算符$\boldsymbol U$，或其矩阵$\{v_{mn/m'n'}\}$刻画。它确定$I+II$中一切物理量的统计性质，特别也确定$I$中物理量的统计性质。因此，它也对应于单独$I$中的一个统计系综：一个只能感知$I$、不能感知$II$的观察者，会把$I+II$系统的系综看作$I$系统的系综。这个$I$系综的统计算符$\mathsf U$，或其矩阵$\{u_{m/m'}\}$是什么？求法如下。矩阵为$\{a_{m/m'}\}$的$I$物理量，当作$I+II$物理量时，矩阵为$\{a_{m/m'}\delta_{n/n'}\}$；所以在$I$中计算，期望值为$\sum_{m,\allowbreak m'=1}^{\infty}u_{m/m'}a_{m'/m}$，而在$I+II$中计算，得到
\begin{align*}
 \sum_{m,n,m',n'=1}^{\infty}v_{mn/m'n'}a_{m'/m}\delta_{n'/n}
 &=\sum_{m,m',n=1}^{\infty}v_{mn/m'n}a_{m'/m}\\
 &=\sum_{m,m'=1}^{\infty}\left(\sum_{n=1}^{\infty}v_{mn/m'n}\right)a_{m'/m}.
\end{align*}
要使二者相等，必须有
\[
 u_{m/m'}=\sum_{n=1}^{\infty}v_{mn/m'n}.
\]

完全类似地，如果只考虑$II$而忽略$I$，这个$I+II$系综也确定一个$II$系综，其统计算符为$U$，矩阵为$\{u_{n/n'}\}$。同样得到
\[
 u_{n/n'}=\sum_{m=1}^{\infty}v_{mn/mn'}.
\]

这样也得到了$I,\allowbreak II,\allowbreak I+II$的统计算符$\mathsf U,\allowbreak U,\allowbreak \boldsymbol U$之间的对应规则；但它们与物理量算符$\mathsf A,\allowbreak A,\allowbreak \boldsymbol A$之间的对应规则有本质差别。

\hypertarget{origpage-227}{}\phantomsection\label{origpage-227}还须指出，$\mathsf U,\allowbreak U,\allowbreak \boldsymbol U$的这一对应，只是表面上依赖于完全规范正交系$\varphi_m(q)$、$\xi_n(q)$的选择。因为它是从一个不变条件导出的，而且只有这一对应满足该条件：$\mathsf A$与$\boldsymbol A$的期望值相等，$A$与$\boldsymbol A$的期望值也相等。

$\boldsymbol U$表示$I+II$中的统计，$\mathsf U,\allowbreak U$分别表示只限于$I$、$II$的统计。问题是：$\mathsf U,\allowbreak U$能否唯一确定$\boldsymbol U$？一般应预期不能，因为仅知道$\mathsf U,\allowbreak U$，也就是只分别了解$I$、$II$，会丢失两系统之间可能存在的一切“概率依赖关系”。但若$I$和$II$的状态都准确已知，就不再涉及“概率依赖关系”，$I+II$也就准确已知。不过，精确的数学讨论优于这种定性考虑，下面就来进行。

问题是：给定两个定号矩阵$\{u_{m/m'}\}$、$\{u_{n/n'}\}$，求第三个定号矩阵$\{v_{mn/m'n'}\}$，使
\[
 \sum_{n=1}^{\infty}v_{mn/m'n}=u_{m/m'},\qquad
 \sum_{m=1}^{\infty}v_{mn/mn'}=u_{n/n'}.
\]
由$\sum_{m=1}^{\infty}u_{m/m}=1$、$\sum_{n=1}^{\infty}u_{n/n}=1$，便自动得$\sum_{m,\allowbreak n=1}^{\infty}v_{mn/mn}=1$，因此仍满足所需的归一化条件。问题总是有解，例如$v_{mn/m'n'}=u_{m/m'}u_{n/n'}$总是一个解；不难看出这个矩阵定号。不过，要问的是，它何时是唯一解。

我们证明：当且仅当两个矩阵$\{u_{m/m'}\}$、$\{u_{n/n'}\}$中至少有一个对应于状态时，解才唯一。先证必要性，即两个矩阵都对应混合态时，存在多个解。这时（参见 \hyperref[sec:IV.2]{IV.2}）
\[
 u_{m/m'}=\alpha v_{m/m'}+\beta w_{m/m'},\qquad
 u_{n/n'}=\gamma v_{n/n'}+\delta w_{n/n'},
\]
其中$v_{m/m'},\allowbreak w_{m/m'}$定号，且不只是相差一个常数因子；$v_{n/n'},\allowbreak w_{n/n'}$也一样；还有
\[
 \sum_{m=1}^{\infty}v_{m/m}=\sum_{m=1}^{\infty}w_{m/m}
 =\sum_{n=1}^{\infty}v_{n/n}=\sum_{n=1}^{\infty}w_{n/n}=1,
\]
及$\alpha,\allowbreak \beta,\allowbreak \gamma,\allowbreak \delta>0$、$\alpha+\beta=1$、$\gamma+\delta=1$。容易验证，任一\tnote[一][][下]{下式左端原书最后的$n$漏印了撇号；由右端及后续用法，应为$v_{mn/m'n'}$。这里保留原印下标。}
\begin{multline*}
 v_{mn/m'n}=\pi v_{m/m'}v_{n/n'}+\varrho w_{m/m'}v_{n/n'}\\
 {}+\sigma v_{m/m'}w_{n/n'}+\tau w_{m/m'}w_{n/n'}
\end{multline*}\hypertarget{origpage-228}{}\phantomsection\label{origpage-228}只要满足
\begin{gather*}
 \pi+\sigma=\alpha,\quad\varrho+\tau=\beta,\quad
 \pi+\varrho=\gamma,\quad\sigma+\tau=\delta,\\
 \pi,\varrho,\sigma,\tau>0,
\end{gather*}
就是解。$\pi,\allowbreak \varrho,\allowbreak \sigma,\allowbreak \tau$有无穷多种选法：由于$\alpha+\beta=\gamma+\delta$，四个方程只有三个独立，例如$\varrho=\gamma-\pi$、$\sigma=\alpha-\pi$、$\tau=(\delta-\alpha)+\pi$；要全部为正，只须$\alpha-\delta=\gamma-\beta<\pi<\alpha,\allowbreak \gamma$，而满足它的$\pi$有无穷多个。不同的$\pi,\allowbreak \varrho,\allowbreak \sigma,\allowbreak \tau$给出不同的$v_{mn/m'n'}$，因为$v_{m/m'}v_{n/n'},\allowbreak \ldots,\allowbreak w_{m/m'}w_{n/n'}$线性无关；这是由$v_{m/m'},\allowbreak w_{m/m'}$线性无关、$v_{n/n'},\allowbreak w_{n/n'}$也线性无关而来。

再证充分性。可设$u_{m/m'}$对应一个状态，另一种情形同样处理。于是$\mathsf U=P_{[\varphi]}$；完全规范正交系$\varphi_1,\allowbreak \varphi_2,\allowbreak \ldots$原本任意，故可取$\varphi_1=\varphi$。显然，$\mathsf U=P_{[\varphi_1]}$的矩阵为
\[
 u_{m/m'}=\begin{cases}1,&m=m'=1,\\0,&\text{其他情形}.\end{cases}
\]
因此
\[
 \sum_{n=1}^{\infty}v_{mn/m'n}
 =\begin{cases}1,&m=m'=1,\\0,&\text{其他情形}.\end{cases}
\]
特别地，$m\ne1$时$\sum_{n=1}^{\infty}v_{mn/mn}=0$。但由$v_{mn/m'n'}$定号，各$v_{mn/mn}\geq0$［因为$v_{mn/mn}=(\boldsymbol U\Phi_{mn},\allowbreak \Phi_{mn})$］，所以此时$v_{mn/mn}=0$。即$(\boldsymbol U\Phi_{mn},\allowbreak \Phi_{mn})=0$，再由$\boldsymbol U$定号，有$(\boldsymbol U\Phi_{mn},\allowbreak \Phi_{m'n'})=0$（参见 \hyperref[sec:II.5]{II.5}，\CJKglue \hyperref[thm:II.19]{定理19}），其中$m',\allowbreak n'$任意。因此，$m\ne1$蕴含$v_{mn/m'n'}=0$；由厄米性，$m'\ne1$也给出同样结论。而$m=m'=1$时，
\[
 v_{1n/1n'}=\sum_{m=1}^{\infty}v_{mn/mn'}=u_{n/n'}.
\]
所以解$v_{mn/m'n'}$确如所言，唯一确定。

把结果概括起来，也可这样说：$I+II$中的统计系综，算符为$\boldsymbol U=\{v_{mn/m'n'}\}$；它在单独$I$、单独$II$中分别确定统计系综，算符为$\mathsf U=\{u_{m/m'}\}$、$U=\{u_{n/n'}\}$。当且仅当满足以下条件时，后二者能唯一确定前者：

\textbf{1.}有$v_{mn/m'n'}=v_{m/m'}v_{n/n'}$。由
\[
 \operatorname{Spur}\boldsymbol U=\sum_{m,n=1}^{\infty}v_{mn/mn}
 =\sum_{m=1}^{\infty}v_{m/m}\sum_{n=1}^{\infty}v_{n/n}=1
\]
可知，把$v_{m/m'}$与$v_{n/n'}$分别乘上互为倒数的两个常数因子，就能使$\sum_{m=1}^{\infty}v_{m/m}=1$、$\sum_{n=1}^{\infty}v_{n/n}=1$。\hypertarget{origpage-229}{}\phantomsection\label{origpage-229}此时便知$u_{m/m'}=v_{m/m'}$、$u_{n/n'}=v_{n/n'}$。

\textbf{2.}或者$v_{m/m'}=\overline{x_m}x_{m'}$，或者$v_{n/n'}=\overline{x_n}x_{n'}$。因为$\mathsf U=P_{[\varphi]}$意味着$\varphi=\sum_{m=1}^{\infty}y_m\varphi_m$，从而$u_{m/m'}=\overline{y_m}y_{m'}$，对$v_{m/m'}$也相应成立；$U=P_{[\xi]}$的情形类似。

我们把$\mathsf U,\allowbreak U$分别称为$\boldsymbol U$在$I,\allowbreak II$中的投影。\origNote[。][]{212}

现在转向$I+II$的状态$\boldsymbol U=P_\Phi$。相应的波函数$\Phi(q,\allowbreak r)$可按完全规范正交系$\Phi_{mn}(q,\allowbreak r)=\varphi_m(q)\xi_n(r)$展开：
\[
 \Phi(q,r)=\sum_{m,n=1}^{\infty}f_{mn}\varphi_m(q)\xi_n(r).
\]

因此，也可以改用系数$f_{mn}$（$m,\allowbreak n=1,\allowbreak 2,\allowbreak \ldots$）来表示波函数；系数只受一个约束：$\sum_{m,\allowbreak n=1}^{\infty}|f_{mn}|^2=\|\Phi\|^2$有限。

用
\begin{align*}
 F\varphi(q)&=\int\overline{\Phi(qr)}\varphi(q)\,dq,\\
 F^*\xi(r)&=\int\Phi(qr)\xi(r)\,dr
\end{align*}
定义两个算符$F,\allowbreak F^*$。这两个算符都是线性的，但有一点特殊：定义域分别是$\mathfrak R^I$、$\mathfrak R^{II}$，值却分别取在$\mathfrak R^{II}$、$\mathfrak R^I$中。二者互为伴随，因为显然$(F\varphi,\allowbreak \xi)=(\varphi,\allowbreak F^*\xi)$；左边内积在$\mathfrak R^{II}$中取，右边在$\mathfrak R^I$中取！$\mathfrak R^I$与$\mathfrak R^{II}$的差别在数学上无关紧要，因而可用 \hyperref[sec:II.11]{II.11} 的结果：由于它们是积分算符，$\Sigma(F),\allowbreak \Sigma(F^*)$都等于
\[
 \iint|\Phi(qr)|^2\,dq\,dr=\|\Phi\|^2=1
 \qquad(\|\Phi\|\mathwordspace{}\text{在}\mathwordspace{}\mathfrak R^{I+II}\mathwordspace{}\text{中取}!),
\]
所以有限。因此$F,\allowbreak F^*$连续，甚至全连续，而$F^*F,\allowbreak FF^*$为定号算符，$\operatorname{Spur}(F^*F)=\Sigma(F)=1$、$\operatorname{Spur}(FF^*)=\Sigma(F^*)=1$。再顾及$\mathfrak R^I,\allowbreak \mathfrak R^{II}$的差别，就知$F^*F$定义在$\mathfrak R^I$中，$FF^*$定义在$\mathfrak R^{II}$中。

由于$F\varphi_m(q)=\sum_{n=1}^{\infty}\overline{f_{mn}}\xi_n(r)$，$F$的矩阵为$\{\overline{f_{mn}}\}$［分别采用完全规范正交系$\varphi_m(q)$与$\xi_n(r)$］；同样，$F^*$的矩阵为$\{f_{mn}\}$。所以$F^*F,\allowbreak FF^*$的矩阵分别为
\[
 \left\{\sum_{n=1}^{\infty}\overline{f_{mn}}f_{m'n}\right\},\qquad
 \left\{\sum_{m=1}^{\infty}\overline{f_{mn}}f_{mn'}\right\},
\]

\hypertarget{origpage-230}{}\phantomsection\label{origpage-230}另一方面，在正交系$\Phi_{mn}(q,\allowbreak r)=\varphi_m(q)\xi_n(r)$中，$\boldsymbol U=P_{[\Phi]}$的矩阵为$\{\overline{f_{mn}}f_{m'n'}\}$，因此它在$I,\allowbreak II$中的投影$\mathsf U,\allowbreak U$的矩阵分别如下。\origNote[。][]{213}
\[
 \left\{\sum_{n=1}^{\infty}\overline{f_{mn}}f_{m'n}\right\},\qquad
 \left\{\sum_{m=1}^{\infty}\overline{f_{mn}}f_{mn'}\right\},
\]
所以
\[
 \mathsf U=F^*F,\qquad U=FF^*,
\]
而且这一论断不依赖于$\varphi_m,\allowbreak \xi_n$的选择，因为$\mathsf U,\allowbreak U,\allowbreak F$都不依赖这种选择。

算符$\mathsf U,\allowbreak U$全连续，根据 \hyperref[sec:II.11]{II.11} 及 \hyperref[sec:IV.3]{IV.3}，可写成
\[
 \mathsf U=\sum_{k=1}^{\infty}w'_kP_{[\psi_k]},\qquad
 U=\sum_{k=1}^{\infty}w''_kP_{[\eta_k]},
\]
其中$\psi_k$构成$\mathfrak R^I$中的完全规范正交系，$\eta_k$构成$\mathfrak R^{II}$中的完全规范正交系，且各$w'_k,\allowbreak w''_k\geq0$。现在分别略去两个公式中$w'_k=0$或$w''_k=0$的项，把余项重新按$k=1,\allowbreak 2,\allowbreak \ldots$连续编号。此时$\psi_k$与$\eta_k$仍是规范正交系，却未必完全；和式$\sum_{k=1}^{\infty}$分别换成$\sum_{k=1}^{M'}$、$\sum_{k=1}^{M''}$，$M',\allowbreak M''$可以有限，也可以无穷，而所有$w'_k,\allowbreak w''_k$现在都大于零。

考察一个$\psi_k$。有$\mathsf U\psi_k=w'_k\psi_k$，故$F^*F\psi_k=w'_k\psi_k$、$FF^*F\psi_k=w'_kF\psi_k$、$UF\psi_k=w'_kF\psi_k$。此外，
\begin{align*}
 (F\psi_k,F\psi_l)&=(F^*F\psi_k,\psi_l)=(\mathsf U\psi_k,\psi_l)\\
 &=w'_k(\psi_k,\psi_l)
 =\begin{cases}w'_k,&k=l,\\0,&k\ne l.\end{cases}
\end{align*}
特别有$\|F\psi_k\|^2=w'_k$。因此，$\frac{1}{\sqrt{w'_k}}F\psi_k$在$\mathfrak R^{II}$中构成规范正交系，且为$U$的本征函数，其本征值与$\psi_k$对$\mathsf U$的本征值相同，即$w'_k$。所以$\mathsf U$的每个本征值也是$U$的本征值，重数至少相同。交换$\mathsf U,\allowbreak U$便知，它们具有相同本征值，重数也相同。因此，$w'_k$与$w''_k$除次序外相同，$M'=M''=M$；把$w''_k$重新编号，就能使$w''_k=w'_k=w_k$。做到这一点后，显然一般可选$\eta_k=\frac{1}{\sqrt{w_k}}F\psi_k$。于是
\[
 \frac{1}{\sqrt{w_k}}F^*\eta_k
 =\frac{1}{w_k}F^*F\psi_k=\frac{1}{w_k}\mathsf U\psi_k=\psi_k.
\]
所以：\origNote[：][]{212}
\[
 \eta_k=\frac{1}{\sqrt{w_k}}F\psi_k,\qquad
 \psi_k=\frac{1}{\sqrt{w_k}}F^*\eta_k.
\]

\hypertarget{origpage-231}{}\phantomsection\label{origpage-231}把规范正交系$\psi_1,\allowbreak \psi_2,\allowbreak \ldots$扩充成完全系$\psi_1,\allowbreak \psi_2,\allowbreak \ldots,\allowbreak \psi'_1,\allowbreak \psi'_2,\allowbreak \ldots$，同样把$\eta_1,\allowbreak \eta_2,\allowbreak \ldots$扩充成$\eta_1,\allowbreak \eta_2,\allowbreak \ldots,\allowbreak \eta'_1,\allowbreak \eta'_2,\allowbreak \ldots$；补入的$\psi'_1,\allowbreak \psi'_2,\allowbreak \ldots$与$\eta'_1,\allowbreak \eta'_2,\allowbreak \ldots$各自都可以为空、有限或无限，彼此无关。前面的讨论不依赖于完全规范正交系$\varphi_1,\allowbreak \varphi_2,\allowbreak \ldots$、$\xi_1,\allowbreak \xi_2,\allowbreak \ldots$的选择，因而可以让它们分别与刚扩充的两个正交系重合。设$\psi_k$对应$\varphi_{\mu_k}$，$\eta_k$对应$\xi_{\nu_k}$，其中$\mu_1,\allowbreak \mu_2,\allowbreak \ldots$互不相同，$\nu_1,\allowbreak \nu_2,\allowbreak \ldots$也互不相同。于是
\[
 F\varphi_{\mu_k}=\sqrt{w_k}\xi_{\nu_k},\qquad
 F\varphi_m=0\quad\text{当}\mathwordspace{}m\ne\mu_1,\mu_2,\ldots.
\]
所以
\[
 f_{mn}=\begin{cases}\sqrt{w_k},&m=\mu_k,\ n=\nu_k,\ k=1,2,\ldots,\\
 0,&\text{其他情形},\end{cases}
\]
或等价地，
\[
 \Phi(q,r)=\sum_{k=1}^{M}\sqrt{w_k}\varphi_{\mu_k}(q)\xi_{\nu_k}(q).
\]

这样，适当选取正交系$\varphi_m(q)$、$\xi_n(q)$后，矩阵$\{f_{mn}\}$的每行、每列至多含一个非零元素；这个元素甚至是实的、正的，等于$\sqrt{w_k}$，但对下文并不重要。这一形式陈述的物理意义是什么？\tnote[？][][这]{这里原书把$\xi_{\nu_k}$及$\xi_n$的自变量印作$q$；但它们属于第二个系统，其他地方的自变量是$r$。此处保留原印。}

设算符$A$的本征函数为$\varphi_1,\allowbreak \varphi_2,\allowbreak \ldots$，互不相同的本征值为$a_1,\allowbreak a_2,\allowbreak \ldots$；同样，$B$有$\xi_1,\allowbreak \xi_2,\allowbreak \ldots$及$b_1,\allowbreak b_2,\allowbreak \ldots$。$A$对应$I$中的物理量，$B$对应$II$中的物理量，因而二者同时可测。容易看出，“$A$的值为$a_m$，$B$的值为$b_n$”两句话合起来，确定状态$\Phi_{mn}(q,\allowbreak r)=\varphi_m(q)\xi_n(r)$；在状态$\Phi(q,\allowbreak r)$中，它出现的概率为$(P_{[\Phi_{mn}]}\Phi,\allowbreak \Phi)=|(\Phi,\allowbreak \Phi_{mn})|^2=|f_{mn}|^2$。所以我们的结论意味着：$A,\allowbreak B$同时可测，在$\Phi$中测出其中一个，就唯一确定另一个的值。如果某个$a_m$的全部$f_{mn}$都为零，它就不可能出现，因为只要测到了$a_m$，其总概率$\sum_{n=1}^{\infty}|f_{mn}|^2$就不能为零；所以恰有一个$n$使$f_{mn}\ne0$。对$b_n$也一样。也就是说，状态$\Phi$中可能有多个$A$值：凡$\sum_{n=1}^{\infty}|f_{mn}|^2>0$，即存在某个$n$使$f_{mn}\ne0$的$a_m$都可能，通常是全部$a_m$；也有同样多个$B$值：凡$\sum_{n=1}^{\infty}|f_{mn}|^2>0$，即存在某个$m$使$f_{mn}\ne0$的$b_n$都可能。但$\Phi$在可能的$A$值与$B$值之间建立了一一对应。\tnote[。][][谈]{谈$B$值的括号中，原书重复使用了求和下标$n$；相应的边缘概率应对$m$求和。这里保留原印下标。}

\hypertarget{origpage-232}{}\phantomsection\label{origpage-232}把可能的$m$记为$\mu_1,\allowbreak \mu_2,\allowbreak \ldots$，相应可能的$n$记为$\nu_1,\allowbreak \nu_2,\allowbreak \ldots$，则
\[
 f_{mn}=\begin{cases}c_k\ne0,&m=\mu_k,\ n=\nu_k,\ k=1,2,\ldots,\\
 0,&\text{其他情形},\end{cases}
\]
故有（$M$有限或无限）
\[
 \Phi(q,r)=\sum_{k=1}^{M}c_k\varphi_{\mu_k}(q)\xi_{\nu_k}(r).
\]
于是
\begin{align*}
 u_{m/m'}&=\sum_{n=1}^{\infty}\overline{f_{mn}}f_{m'n}
 =\begin{cases}|c_k|^2,&m=m'=\mu_k,\ k=1,2,\ldots,\\0,&\text{其他情形},\end{cases}\\
 u_{n/n'}&=\sum_{m=1}^{\infty}\overline{f_{mn}}f_{mn'}
 =\begin{cases}|c_k|^2,&n=n'=\nu_k,\ k=1,2,\ldots,\\0,&\text{其他情形},\end{cases}
\end{align*}
因此
\[
 \mathsf U=\sum_{k=1}^{M}|c_k|^2P_{[\varphi_{\mu_k}]},\qquad
 U=\sum_{k=1}^{M}|c_k|^2P_{[\xi_{\nu_k}]}.
\]
所以，投影到$I$或$II$中，$\Phi$一般会成为混合态；它只有在$I+II$中才是状态，因为它包含一项关于$I+II$的论断，无法在单独$I$或单独$II$中利用：即$A$值与$B$值之间的一一对应。

因此，对每个$\Phi$都能选取$A,\allowbreak B$，亦即$\varphi_m,\allowbreak \xi_n$，使我们的条件成立；若$A,\allowbreak B$任意选定，当然不能指望这一条件成立。每个状态$\Phi$都在$I$与$II$之间建立一种特殊联系，其中相关联的量$A,\allowbreak B$依赖于$\Phi$。$\Phi$能把它们（即$\varphi_m,\allowbreak \xi_n$）确定到什么程度，不难判断：若所有$|c_k|$彼此不同且非零，$\mathsf U,\allowbreak U$（由$\Phi$确定）就唯一确定$\varphi_m,\allowbreak \xi_n$，参见 \hyperref[sec:IV.3]{IV.3}。一般讨论留给读者。

最后指出：因为$|c_k|^2>0$，$M\ne1$时，$\mathsf U,\allowbreak U$都不是状态。$M=1$时则二者都是状态：$\mathsf U=P_{[\varphi_{\mu_1}]}$，$U=P_{[\xi_{\nu_1}]}$；此时$\Phi(q,\allowbreak r)=c_1\varphi_{\mu_1}(q)\xi_{\nu_1}(r)$，可把$c_1$吸收到$\varphi_{\mu_1}(q)$中。因此，当且仅当$\Phi(q,\allowbreak r)$形如$\varphi(q)\xi(r)$时，$\mathsf U,\allowbreak U$是状态，且分别等于$P_{[\varphi]}$、$P_{[\xi]}$。

根据以上结果，再强调一点：若$I$处于状态$\varphi(q)$，$II$处于状态$\xi(r)$，则$I+II$处于状态$\Phi(q,\allowbreak r)=\varphi(q)\xi(r)$。反之，若$I+II$处于一个不能写成乘积$\varphi(q)\xi(r)$的状态$\Phi(q,\allowbreak r)$，那么$I,\allowbreak II$都是混合态，但$\Phi$使$I$、$II$中某些量的可能值之间建立一一对应。

\hypertarget{origpage-233}{}\phantomsection\label{origpage-233}\section{测量过程的讨论}\label{sec:VI.3}
在借助 \hyperref[sec:VI.2]{VI.2} 的形式工具，按 \hyperref[sec:VI.1]{VI.1} 的思路完成测量过程的讨论之前，先用 \hyperref[sec:VI.2]{VI.2} 的结果，排除一种曾多次被提出的解释；它试图说明过程 \hyperref[eq:V.1.1]{\textbf{1}}（\hyperref[sec:V.1]{V.1}）的统计性质。其思路如下：设$I$为被观察系统，$II$为观察者。测量前，$I$处于状态$\mathsf U=P_{[\varphi]}$，$II$却是混合态$U=\sum_{n=1}^{\infty}w_nP_{[\xi_n]}$；那么$I+II$就是一个唯一确定的混合态$\boldsymbol U$。确实，由 \hyperref[sec:VI.2]{VI.2} 容易算出
\[
 \boldsymbol U=\sum_{n=1}^{\infty}w_nP_{[\Phi_n]},\qquad
 \Phi_n(q,r)=\varphi(q)\xi_n(r).
\]
现在测量$I$中的量$A$，应把它理解为$I,\allowbreak II$之间的相互作用；这是 \hyperref[sec:V.1]{V.1} 的过程 \hyperref[eq:V.1.2]{\textbf{2}}，能量算符为$\boldsymbol H$。若持续时间为$t$，$\boldsymbol U$就变为
\[
 \boldsymbol U'=e^{-\frac{2\pi i}{h}t\boldsymbol H}\boldsymbol Ue^{\frac{2\pi i}{h}t\boldsymbol H},
\]
显然
\[
 \boldsymbol U'=\sum_{n=1}^{\infty}w_nP_{[e^{-\frac{2\pi i}{h}t\boldsymbol H}\Phi_n]}.
\]
如果每个$\Phi_n(q,\allowbreak r)$都形如$\psi_n(q)\eta_n(r)$，其中$\psi_n$是$A$的本征函数，$\eta_n$是某个固定的完全规范正交系，那么这一干预就具有测量的性质\tnote[质][：][原]{原书此句仍用测量前的记号$\Phi_n$。按上文的演化式及此处论证，应指演化后的$e^{-\frac{2\pi i}{h}t\boldsymbol H}\Phi_n$具有$\psi_n(q)\eta_n(r)$的形式。正文保留原印记号。}：把$I$的每个状态$\varphi$都变成$A$的本征函数$\psi_n$的混合。统计性质是这样产生的：测量前$I$虽处于单一状态，$II$却是混合态；相互作用期间，$II$的混合性质“传染”了$I+II$，特别使它在$I$中的投影成为混合态。也就是说，测量结果不确定，是因为测量前不知道观察者的确切状态。这样的机制似乎可能存在，因为自然规律也许限制了观察者对自身状态的了解。这样的限制会体现在$w_n$的取值上；而这些数值必须完全由观察者决定，因而与$\varphi$无关！

但这一解释尝试恰在此失败：量子力学要求$w_n=(P_{\psi_n}\varphi,\allowbreak \varphi)=|(\varphi,\allowbreak \psi_n)|^2$，即$w_n$依赖$\varphi$！即使可能另有一种分解$\boldsymbol U'=\sum_{n=1}^{\infty}w'_nP_{[\Phi'_n]}$，也无济于事；这里$\Phi'_n(q,\allowbreak r)=\psi_n(q)\eta_n(r)$是规范正交的。因为$w'_n$除次序外已由$\boldsymbol U'$唯一确定（\hyperref[sec:IV.3]{IV.3}），故与$w_n$相等。\origNote[。][]{214}

因此，过程 \hyperref[eq:V.1.1]{\textbf{1}} 的非因果性，并非由于对观察者状态了解不足；以下便始终假定这个状态准确已知。

\hypertarget{origpage-234}{}\phantomsection\label{origpage-234}回到 \hyperref[sec:VI.1]{VI.1} 末尾提出的任务。系统$I,\allowbreak II,\allowbreak III$仍取那里的意义；对$I,\allowbreak II$的量子力学研究，采用 \hyperref[sec:VI.2]{VI.2} 的记号，$III$仍不纳入计算，参见 \hyperref[sec:VI.1]{VI.1} 的讨论。设$A$是真正待测的量（在$I$中），$\varphi_1(q),\allowbreak \varphi_2(q),\allowbreak \ldots$是其本征函数，$I$处于状态$\varphi(q)$。

若$I$是被观察系统，$II+III$是观察者，就须采用过程 \hyperref[eq:V.1.1]{\textbf{2}}\tnote[\textbf{2}][，][原]{原书这里印作“过程2”；但按 \hyperref[sec:V.1]{V.1} 及 \hyperref[sec:VI.1]{VI.1} 的术语，同句所述的测量变换是过程1。这里保留原印数字。}，得到：测量使$I$从状态$\varphi$转入$\varphi_n$之一（$n=1,\allowbreak 2,\allowbreak \ldots$），概率分别为$|(\varphi,\allowbreak \varphi_n)|^2$（$n=1,\allowbreak 2,\allowbreak \ldots$）。那么，若以$I+II$为被观察系统，只有$III$为观察者，应怎样描述呢？

此时必须说：$II$是一台测量仪器，在刻度上显示$A$（在$I$中）的值；刻度上的指针位置是物理量$B$（在$II$中），这才是$III$真正观察的量。若$II$已经在观察者体内，刻度和指针位置就须换成相应生理概念，例如视网膜及其上的像等。$B$的值与$A$的值一一对应。设$A$取值$a_1,\allowbreak a_2,\allowbreak \ldots$，$B$取值$b_1,\allowbreak b_2,\allowbreak \ldots$，编号使$a_n$与$b_n$对应。

起初，$I$处于未知状态$\varphi(q)$，$II$处于已知状态$\xi(r)$，因此$I+II$处于状态$\Phi(q,\allowbreak r)=\varphi(q)\xi(r)$。由$II$对$I$实行的测量，如前例一样，是在时间$t$内由$I+II$中的能量算符$\boldsymbol H$实现的；这是过程 \hyperref[eq:V.1.2]{\textbf{2}}，把$\Phi$变为$\Phi'=e^{-\frac{2\pi i}{h}t\boldsymbol H}\Phi$。从观察者$III$看来，只有满足以下条件才能称为测量：若$III$用过程 \hyperref[eq:V.1.1]{\textbf{1}} 测量同时可测的$A,\allowbreak B$（分别在$I,\allowbreak II$中，或都在$I+II$中），$m\ne n$的值对$a_m,\allowbreak b_n$概率为零，$m=n$的值对则有某些概率$w_n$。也就是说，只须“看一眼”$II$，便已测量了$I$中的$A$。量子力学还要求$w_n=|(\varphi,\allowbreak \varphi_n)|^2$。

做到这一点，测量过程在$II$中进行的部分，就在理论上得到“解释”；也就是说，\hyperref[sec:VI.1]{VI.1} 所述的界限从$I\mid II+III$移到了$I+II\mid III$。

因此，数学任务如下：给定$I$中的完全规范正交系$\varphi_1,\allowbreak \varphi_2,\allowbreak \ldots$，求$\mathfrak R^{II}$中的同类正交系$\xi_1,\allowbreak \xi_2,\allowbreak \ldots$、状态$\xi$，以及$\mathfrak R^{I+II}$中的能量算符$\boldsymbol H$和时间$t$，使下述性质成立：若$\varphi$是$\mathfrak R^I$中的任意状态，并令$\Phi(q,\allowbreak r)=\varphi(q)\xi(r)$、$\Phi'(q,\allowbreak r)=e^{-\frac{2\pi i}{h}t\boldsymbol H}\Phi(q,\allowbreak r)$，则$\Phi'(q,\allowbreak r)$应形如$\sum_{n=1}^{\infty}c_n\varphi_n(q)\xi_n(r)$（$c_n$自然依赖$\varphi$），并有$|c_n|^2=|(\varphi,\allowbreak \varphi_n)|^2$。\hypertarget{origpage-235}{}\phantomsection\label{origpage-235}最后这一要求与上面物理要求等价，已在 \hyperref[sec:VI.2]{VI.2} 讨论。

以下除$\varphi_1,\allowbreak \varphi_2,\allowbreak \ldots$外，也把$\xi_1,\allowbreak \xi_2,\allowbreak \ldots$及$\xi$固定给定；不求$\boldsymbol H$，而求酉算符$\Delta=e^{-\frac{2\pi i}{h}t\boldsymbol H}$。

这个数学问题把我们带回 \hyperref[sec:VI.2]{VI.2} 中已解的问题：那里给定$\Phi'$，证明$c_n,\allowbreak \varphi_n,\allowbreak \xi_n$存在；这里$\varphi_n,\allowbreak \xi_n$固定，$\Phi,\allowbreak c_n$作为$\varphi$的函数给定，要确定$\Delta$，使$\Phi'=\Delta\Phi$恰好给出这些$c_n,\allowbreak \varphi_n,\allowbreak \xi_n$。

下面证明，这样的$\Delta$确实能确定。我们只关心原理问题，即某种这样的$\Delta$是否存在；至于最简单、最直观的实验装置（例如 \hyperref[sec:III.4]{III.4} 中的装置）的$\Delta=e^{-\frac{2\pi i}{h}t\boldsymbol H}$是否也有这一性质，则不再研究。因为已经看到，我们的条件事实上等价于判断一个干预是否具有测量性质的直观判据，而上述实验装置无疑具有这种直观的测量性质；如果它们的$\Delta$竟不满足我们的条件，哪怕近似也不满足，量子力学就会给出与直观严重矛盾的结果。\origNote[。][因]{215}因此，以下只给出一个严格满足条件的抽象$\Delta$。

设$\varphi_m$（$m=0,\allowbreak \pm1,\allowbreak \pm2,\allowbreak \ldots$）、$\xi_n$（$n=0,\allowbreak \pm1,\allowbreak \pm2,\allowbreak \ldots$）分别是$\mathfrak R^I$、$\mathfrak R^{II}$中给定的完全规范正交系。让$m,\allowbreak n$不遍历$1,\allowbreak 2,\allowbreak \ldots$而遍历$0,\allowbreak \pm1,\allowbreak \pm2,\allowbreak \ldots$，纯为技术方便，在原理上没有区别。为简单起见，取状态$\xi=\xi_0$。以
\[
 \Delta\sum_{m,n=-\infty}^{\infty}x_{mn}\varphi_m(q)\xi_n(r)
 =\sum_{m,n=-\infty}^{\infty}x_{mn}\varphi_m(q)\xi_{m+n}(r)
\]
定义$\Delta$。$\varphi_m(q)\xi_n(r)$与$\varphi_m(q)\xi_{m+n}(r)$都构成$\mathfrak R^{I+II}$中的完全规范正交系，故$\Delta$为酉算符。现在
\[
 \varphi(q)=\sum_{m=-\infty}^{\infty}(\varphi,\varphi_m)\varphi_m(q),\qquad
 \xi(r)=\xi_0(r),
\]
所以\tnote[以][][下]{下列两行的第一行中，原书把$\xi(r)$印作$\xi(q)$。这里保留原印自变量。}
\begin{align*}
 \Phi(q,r)&=\varphi(q)\xi(q)
 =\sum_{m=-\infty}^{\infty}(\varphi,\varphi_m)\varphi_m(q)\xi_0(r),\\
 \Phi'(q,r)&=\Delta\Phi(q,r)
 =\sum_{m=-\infty}^{\infty}(\varphi,\varphi_m)\varphi_m(q)\xi_m(r).
\end{align*}
目的至此达到，甚至有$c_n=(\varphi,\allowbreak \varphi_n)$。

\hypertarget{origpage-236}{}\phantomsection\label{origpage-236}如果用具体的Schrödinger波函数举例，并且不给$\Delta$而直接给出$\boldsymbol H$，就能更清楚地看出这一过程的机制。

设被观察对象与观察者，即$I$与$II$，分别只用一个连续遍历$-\infty$至$+\infty$的变量$q$、$r$刻画；也就是说，把二者都视为能沿直线运动的点，波函数分别总形如$\psi(q)$、$\eta(r)$。假设其质量$m_1,\allowbreak m_2$很大，使能量算符中的动能部分
\[
 \frac{1}{2m_1}\left(\frac{h}{2\pi i}\frac{\partial}{\partial q}\right)^2
 +\frac{1}{2m_2}\left(\frac{h}{2\pi i}\frac{\partial}{\partial r}\right)^2,
\]
可以忽略。于是$\boldsymbol H$只剩对测量起决定作用的相互作用能部分，取其特殊形式为$\frac{h}{2\pi i}q\frac{\partial}{\partial r}$。

此时，$I+II$的波函数$\psi_t=\psi_t(q,\allowbreak r)$所满足的含时Schrödinger微分方程为
\begin{align*}
 \frac{h}{2\pi i}\frac{\partial}{\partial t}\psi_t(q,r)
 &=-\frac{h}{2\pi i}q\frac{\partial}{\partial r}\psi_t(q,r),\\
 \left(\frac{\partial}{\partial t}+q\frac{\partial}{\partial r}\right)\psi_t(q,r)&=0,
\end{align*}
即
\[
 \psi_t(q,r)=f(q,r-tq).
\]
若$t=0$时$\psi_0(q,\allowbreak r)=\Phi(q,\allowbreak r)$，则$f(q,\allowbreak r)=\Phi(q,\allowbreak r)$，所以
\[
 \psi_t(q,r)=\Phi(q,r-tq).
\]
特别地，若$I,\allowbreak II$的初始状态分别由$\varphi(q)$与$\xi(r)$表示，在我们的计算框架中（把其中的时间$t$取为$1$），便有
\begin{align*}
 \Phi(q,r)&=\varphi(q)\xi(r),\\
 \Phi'(q,r)&=\psi_1(q,r)=\varphi(q)\xi(r-q).
\end{align*}
现在要证明，这可以用来由$II$测量$I$的位置，也就是说，坐标$q,\allowbreak r$的值相互对应。由于$q,\allowbreak r$有连续谱，只能任意精确而不能绝对精确地测量，这一目标也只能近似达到。

为此，设$\xi(r)$只在很窄的区间$-\varepsilon<r<\varepsilon$内非零，即测量前观察者的坐标$r$已十分精确地知道。此外，$\xi$当然须归一化：
\[
 \|\xi\|=1,\qquad\text{即}\qquad\int|\xi(r)|^2\,dr=1.
\]

\hypertarget{origpage-237}{}\phantomsection\label{origpage-237}$q$落在区间$q_0-\delta<q<q_0+\delta$、$r$落在区间$r_0-\delta'<r<r_0+\delta'$的概率为
\begin{multline*}
 \int_{q_0-\delta}^{q_0+\delta}\int_{r_0-\delta'}^{r_0+\delta'}
 |\Phi'(qr)|^2\,dq\,dr\\
 =\int_{q_0-\delta}^{q_0+\delta}\int_{r_0-\delta'}^{r_0+\delta'}
 |\varphi(q)|^2|\xi(r-q)|^2\,dq\,dr.
\end{multline*}
如果$q_0,\allowbreak r_0$相差超过$\delta+\delta'+\varepsilon$，此式为零；也就是说，$q,\allowbreak r$的取值彼此紧密对应，差别不可能超过$\delta+\delta'+\varepsilon$。而当$r_0=q_0$时，只要选$\delta'\geq\delta+\varepsilon$，根据对$\xi$的假设，此式便等于$\int_{q_0-\delta}^{q_0+\delta}|\varphi(q)|^2\,dq$。由于$\delta,\allowbreak \delta',\allowbreak \varepsilon$可以任意小，只须大于零，这就说明：$q,\allowbreak r$的值可以在任意高的精度下相互对应，概率密度也具有量子力学所要求的值$|\varphi(q)|^2$。

这样，\hyperref[sec:VI.1]{VI.1} 及本节所讨论的测量情况就实现了。

更复杂的例子，例如与 \hyperref[sec:VI.1]{VI.1} 那个四层次例子类似的情况，或由第二个观察者$III$检验$II$对$I$所作测量是否正确，也可类似讨论。这留给读者。


\end{document}
